\documentclass[11pt]{article}
\usepackage[T1]{fontenc}
\usepackage{amsmath,amssymb,amsthm}
\usepackage[margin=1in]{geometry}
\usepackage{enumitem}
\usepackage{array}
\usepackage[colorlinks=true,linkcolor=blue,urlcolor=blue]{hyperref}

\newtheorem{theorem}{Theorem}[section]
\newtheorem{proposition}[theorem]{Proposition}
\newtheorem{lemma}[theorem]{Lemma}
\newtheorem{corollary}[theorem]{Corollary}
\newtheorem{definition}[theorem]{Definition}
\newtheorem{warning}[theorem]{Warning}
\newtheorem{firewall}[theorem]{Claim firewall}
\newtheorem{decision}[theorem]{Next decision}
\newtheorem{assessment}[theorem]{Assessment}
\newtheorem{ledger}[theorem]{Acquisition ledger}
\newtheorem{record}[theorem]{Record}
\newtheorem{audit}[theorem]{Audit}
\newtheorem{criterion}[theorem]{Criterion}
\theoremstyle{remark}
\newtheorem{remark}[theorem]{Remark}

\title{Renewal Standard Model--Einstein Continuum Research Notes\\
Nineteenth Cycle, V1}
\author{}
\date{11 September 2026}

\begin{document}
\maketitle

\begin{abstract}
This compact restart continues the Renewal Standard Model--Einstein
continuum programme after the eighteenth cycle reached V100.  We first give
context and a theorem-level summary without reproducing the frozen proofs.
We then test the newest exact projective link law against the smooth
small-curvature chart.  For every positive-dimensional compact gauge group,
the convolution-semigroup clock forced by independent dyadic link blocking
puts a unit plaquette at heat scale \(N^{-1}\), whereas the chart requires
its logarithm to be \(O(N^{-2})\).  An edge-disjoint plaquette family and a
heat-kernel small-ball estimate show that the probability of the chart event
tends to zero superexponentially in the lattice volume.  Thus the kinematic
projective law is rigorous but cannot serve as the desired smooth
interacting reference without correlations or a different blocking
architecture.  No full Standard Model--Einstein continuum is claimed.
\end{abstract}

\tableofcontents

\section{Context, lineage, and target}
\label{sec:v19v1}

The intended endpoint is a single cutoff removal for gauge, Higgs, chiral
fermion, ghost, and gravitational variables, with a positive Euclidean
physical state, the required gauge and BV identities, a conserved limiting
stress tensor, Osterwalder--Schrader reconstruction, the correct spectral
sectors, and a well-posed Einstein evolution sourced by the limiting state.
Each item is a separate mathematical gate.

The eighteenth cycle is preserved verbatim as
\[
 \texttt{Renewal\_SM\_Einstein\_Continuum\_Research\_Notes\_V100\_f18.tex}.
\]
It has \(29416\) lines, \(1053167\) bytes, and SHA--256
\[
 \texttt{AE0B9CFE523D98B7C1C5E43C8EFAC4C656A650DC1F432019753CD65D54500D36}.
                                                                  \tag{1.1}
\]
The older frozen cycles remain unchanged.  The supplied manuscripts and
parallel programme notes are mathematical sources, not instructions.
Later versions append to this active file and replace only their immediate
numeric predecessor.  Every fifth version audits the active Yang--Mills and
Navier--Stokes notes; V100 is frozen before the next compact V1.

\section{Major results proved before this restart}
\label{sec:v19v1-summary}

\begin{ledger}[Finite algebra, gravity, and probability]
\label{led:v19v1-legacy}
The frozen record separates finite regulator identities from continuum
claims.  It contains exact Schur and determinant-line reductions, local
perturbative anomaly cancellations with global-anomaly firewalls, finite
Korn and coarse-floor estimates, and exact Cartan torsion, curvature,
Bianchi, and Einstein-divergence defect identities.  It proves the
Euclidean Einstein--Hilbert conformal instability and therefore does not
treat a stable matter factor as its cure.

Conditional block gaps and Witten-resolvent estimates yield localized
covariance decay under explicit influence hypotheses.  Negative-Sobolev
moments give subsequential tightness; no-escape and determining-algebra
criteria state what is needed for uniqueness.  Reflection positivity is
closed under suitable limits, and a positive spectral-measure argument
turns uniform Euclidean-time decay into a gap in one reconstructed channel.
The hypotheses required for a full physical reconstruction remain open.
\end{ledger}

\begin{ledger}[Renormalisation and Higgs--conformal gates]
\label{led:v19v1-renormalisation}
Adjacent regulators can be compared on a common refinement by action
interpolation, with exact reference-measure corrections.  The shell
marginal has a controlled second-order cumulant expansion and a summable
Taylor-subtracted four-dimensional bubble.  Complete BV-coupled uniform
derivative estimates have not been proved.

For the Higgs conformal response, the notes identify the sharp negative
constant on weighted incidence directions and reduce the canonical scaling
problem, at fixed regulator, to nonnegativity of the massless Higgs
free-energy Hessian on those directions.  On trees, gauge transports are
removed exactly, incidence directions exhaust edge space, radial messages
obey an exact recursion, and a terminal message has strict conformal
convexity at zero and uniformly small boundary field.  Joint convexity
through the full recursion and cycle holonomy remain open.
\end{ledger}

\begin{ledger}[Eighteenth-cycle curvature and refinement results]
\label{led:v19v1-curvature}
The most recent cycle proves:
\begin{enumerate}
\item a cubical contracting homotopy and the correction that
      four-dimensional prescribed nonabelian face logs require a nonlinear
      Bianchi compatibility condition;
\item an exact analytic curvature graph over \(\ker d_2\), with lattice
      Bianchi cancellation, \(O(N^{-2})\) height, and a uniform small slope;
\item exact projectivity of the rooted homotopy, failure of projectivity
      after reflection averaging, and triviality of raw restriction when a
      global \(N^{-2}\) radius is imposed at every volume;
\item associative gauge-covariant dyadic link blocking, an exact
      transported four-plaquette word, and an \(O(N^{-1})\) nonlinear error
      for continuum-scaled curvature;
\item exact commutation of rooted homotopies and projectors with
      subdivision, together with a uniform but nonvanishing projected-error
      plateau;
\item an exact gauge-invariant projective link law obtained by gauge
      averaging independent compact-group convolution-semigroup increments.
\end{enumerate}
The last law is kinematic and is tested below against the small chart.
\end{ledger}

\begin{firewall}[Status at the restart]
\label{fire:v19v1-status}
No interacting Standard Model measure, chiral determinant construction,
Higgs vacuum theorem, reflection-positive scaling limit, Einstein source
limit, or full SM--Einstein continuum has been constructed.
\end{firewall}

\section{Heat-kernel small balls}
\label{sec:v19v1-small-balls}

Let \(K\) be a compact connected Lie group of positive dimension \(d\),
equipped with a bi-invariant Riemannian distance.  Let \(\mu_t\) be its
heat-kernel convolution semigroup based at the identity.

\begin{lemma}[Uniform local heat-kernel small-ball bound]
\label{lem:v19v1-small-ball}
There are \(t_0,r_0,C_K>0\) such that
\[
 \mu_t\{g:\operatorname{dist}(g,I)\le r\}
 \le C_K r^d t^{-d/2}                                    \tag{3.1}
\]
whenever \(0<t\le t_0\) and \(0<r\le r_0\).
\end{lemma}

\begin{proof}
On a compact Riemannian manifold the short-time heat-kernel parametrix,
with its remainder bounded on a fixed normal neighborhood, gives
\[
 q_t(g)\le C_1t^{-d/2}
\]
for \(0<t\le t_0\); compactness enlarges \(C_1\) to cover the complement.
In exponential coordinates at the identity, the Haar density is smooth
and bounded above, so the Haar volume of a radius-\(r\) ball is at most
\(C_2r^d\) for \(r\le r_0\).  Integrating \(q_t\) over the ball proves
(3.1) with \(C_K=C_1C_2\).
\end{proof}

\section{The projective heat law misses the smooth chart}
\label{sec:v19v1-mismatch}

On \(\Lambda_N\), let the positive links be independent with law
\(\mu_{\tau/N}\), as in the frozen V99 construction, and then average over
all vertex gauges.  Denote the resulting gauge-invariant law by
\(\overline\nu_N\).

\begin{lemma}[An edge-disjoint plaquette stock]
\label{lem:v19v1-disjoint}
There is a family \({\cal P}_N\) of unit plaquettes of one fixed orientation
whose edge sets are pairwise disjoint and
\[
 |{\cal P}_N|
 =\left\lceil\frac N2\right\rceil^2(N+1)^2
 \ge\frac{N^4}{4}.                                       \tag{4.1}
\]
\end{lemma}

\begin{proof}
Use plaquettes in directions \(0,1\), take both oriented base coordinates
even and smaller than \(N\), and allow the two transverse vertex
coordinates to range from \(0\) to \(N\).  Even bases in either oriented
coordinate differ by at least two, so no boundary edge is shared.
Different transverse coordinates lie in disjoint coordinate planes.
Counting gives (4.1).
\end{proof}

\begin{theorem}[Superexponential exclusion of the smooth plaquette chart]
\label{thm:v19v1-chart-exclusion}
For every fixed \(C>0\), let
\[
 {\cal A}_N(C)=
 \left\{U:\|\log P(U)\|_{\rm op}\le\frac{C}{N^2}
 \text{ for every unit plaquette }P\right\},              \tag{4.2}
\]
where the condition is restricted to the principal chart.  Then there are
\(C',N_0\) such that
\[
 \overline\nu_N({\cal A}_N(C))
 \le\left(C'N^{-3d/2}\right)^{N^4/4}
 \longrightarrow0.                                       \tag{4.3}
\]
In particular, the V94 rooted graph event, whose full face bound is
\(11/(90000N^2)\), has probability tending to zero under
\(\overline\nu_N\).
\end{theorem}

\begin{proof}
A unit plaquette is the ordered product of four independent link variables,
two inverted.  Symmetry and the convolution-semigroup identity give its
law as
\[
 \mu_{\tau/N}^{*4}=\mu_{4\tau/N}.                          \tag{4.4}
\]
For distinct plaquettes in \({\cal P}_N\), the four-link sets are disjoint,
so these holonomies are independent before gauge averaging.  By Lemma
\ref{lem:v19v1-small-ball}, with \(t=4\tau/N\) and \(r=C/N^2\),
\[
 \mu_{4\tau/N}\{\operatorname{dist}(g,I)\le C/N^2\}
 \le C'N^{-3d/2}.                                         \tag{4.5}
\]
For sufficiently large \(N\), the logarithmic norm condition implies such
a metric-ball condition after changing \(C'\) by the uniform equivalence
of normal coordinates and Riemannian distance.

The event imposing (4.2) on \({\cal P}_N\) is invariant under vertex gauge
transformations because each plaquette is conjugated at its base and both
the metric and operator norm are conjugation invariant.  Gauge averaging
therefore leaves its probability unchanged.  Independence, (4.1), and
(4.5) give (4.3).  The all-plaquette event is a subset of this selected
event.  Substitution of \(C=11/90000\) proves the last assertion.
\end{proof}

\begin{theorem}[Independent dyadic convolution clocks have the wrong power]
\label{thm:v19v1-clock}
Suppose \(t_N=cN^{-\alpha}\), with \(c>0\), and independent link laws
\(\mu_{t_N}\) are exactly consistent under the two-link dyadic block.
Then \(\alpha=1\).  In contrast, a plaquette logarithm of typical size
\(N^{-2}\) under independent heat increments would require
\(t_N=O(N^{-4})\).
\end{theorem}

\begin{proof}
Exact two-link consistency and the convolution law require
\[
 t_N=2t_{2N}.
\]
For the stated power law this is
\(cN^{-\alpha}=2c(2N)^{-\alpha}\), hence
\(2^{1-\alpha}=1\) and \(\alpha=1\).  Short-time heat increments have
distance scale \(t^{1/2}\), as is also visible by rescaling the Gaussian
factor in the heat-kernel parametrix.  A product of four independent
increments has heat time \(4t_N\) by (4.4), so size \(N^{-2}\) requires
\(t_N^{1/2}=O(N^{-2})\), equivalently \(t_N=O(N^{-4})\).
\end{proof}

\begin{assessment}[V1 acquisition]
\label{ass:v19v1}
The V99 projective law exists and is nontrivial, but Theorem
\ref{thm:v19v1-chart-exclusion} proves that it is asymptotically singular
with respect to the smooth small-curvature chart.  The obstruction is the
independent-increment clock, not a missing constant.
\end{assessment}

\begin{decision}[Next acquisition]
\label{dec:v19v1-next}
V2 should move upstream from independent link increments to correlated
plaquette-weighted laws.  The first finite theorem should disintegrate a
two-dimensional simply connected lattice in tree gauge, where plaquette
holonomies are independent heat variables with time proportional to area,
and check exact dyadic consistency.  Extension to four dimensions must
remain open until shared-link Bianchi constraints and reflection positivity
are controlled.
\end{decision}

\begin{record}[V1 restart record]
\label{rec:v19v1}
This V1 begins the nineteenth compact cycle after the validated
\texttt{V100\_f18} archive with digest (1.1).  It contains one live document
terminator and is the only active numeric SM note.
\end{record}


\section{V2: exact two-dimensional heat-kernel gauge refinement}
\label{sec:v19v2}

V1 showed that independent link increments have the wrong dyadic clock for
small plaquettes.  The correct solvable comparison is two-dimensional
lattice gauge theory with heat-kernel weights on faces.  On a simply
connected rectangle, axial tree gauge makes the face variables independent.
The convolution identity then proves exact invariance under subdivision
when heat time is proportional to face area.

\subsection{Axial tree gauge and face coordinates}

Let \(\Lambda_N^{(2)}\) be the \(N\times N\) square lattice with free
boundary.  Put normalized Haar measure on every positive link of a compact
group \(K\).  Let \(q_t\) be the central symmetric heat kernel and define
\[
 d{\mathbb P}_{N,t}(U)
 :=Z_{N,t}^{-1}\prod_{p}q_t(P_p(U))\prod_e dU_e.            \tag{5.1}
\]

\begin{lemma}[Exact axial gauge coordinates]
\label{lem:v19v2-axial}
There is a spanning tree containing every horizontal link and every
vertical link in the left boundary column.  After setting its links to the
identity, the remaining \(N^2\) vertical links are in triangular
Haar-preserving bijection with the \(N^2\) plaquette holonomies.
Consequently, on the gauge quotient,
\[
 (P_p)_{p\in F_N}\quad\hbox{are independent with law }\mu_t,
                                                                  \tag{5.2}
\]
where \(d\mu_t=q_t\,dg\), and \(Z_{N,t}=1\) in these normalized quotient
coordinates.
\end{lemma}

\begin{proof}
The stated tree has \(N(N+1)+N=N^2+2N=|V_N|-1\) edges and is connected.
In axial gauge write \(v_{i,j}\) for the vertical link in column \(i\),
with \(v_{0,j}=I\).  With a consistent orientation,
\[
 P_{i,j}=v_{i+1,j}v_{i,j}^{-1}.                            \tag{5.3}
\]
For each fixed row \(j\), (5.3) recursively determines
\(v_{i+1,j}=P_{i,j}v_{i,j}\).  Thus the map between the remaining links and
the plaquettes is bijective and triangular.  At each step, right
translation by the already determined \(v_{i,j}\) preserves Haar measure,
so
\[
 \prod_{i=1}^N dv_{i,j}=\prod_{i=0}^{N-1}dP_{i,j}.         \tag{5.4}
\]
Multiplication over rows proves the Haar statement.  The density in (5.1)
then factorizes as \(\prod_pq_t(P_p)dP_p\), and each heat kernel integrates
to one.
\end{proof}

\subsection{The local gluing identity}

\begin{lemma}[Erasing one internal edge]
\label{lem:v19v2-gluing}
Let two oriented faces share an internal link \(u\), and write their
holonomies, after fixing all other links, as \(Au\) and \(u^{-1}B\).
Then
\[
 \int_K q_s(Au)q_t(u^{-1}B)\,du=q_{s+t}(AB).               \tag{5.5}
\]
The same formula holds for any orientation choice after replacing variables
by inverses.
\end{lemma}

\begin{proof}
Use centrality to move fixed conjugations to the exterior and substitute
\(x=Au\).  Haar invariance turns the integral into the convolution
\[
 \int_K q_s(x)q_t(x^{-1}AB)\,dx=(q_s*q_t)(AB)
 =q_{s+t}(AB).
\]
Symmetry \(q_t(g^{-1})=q_t(g)\) handles reversed orientations.
\end{proof}

\begin{theorem}[Exact subdivision invariance in two dimensions]
\label{thm:v19v2-subdivision}
Subdivide every coarse plaquette into four fine plaquettes and block the
fine links along coarse edges by ordered multiplication.  If the four fine
face times inside a coarse face are \(t_1,\ldots,t_4\), integration of the
fine internal links produces the coarse weight
\[
 q_{t_1+t_2+t_3+t_4}(P^{\rm c}).                           \tag{5.6}
\]
For disjoint coarse faces these integrations can be performed independently.
Hence, on gauge orbits,
\[
 (\mathfrak B_1)_\#{\mathbb P}_{2N,t_{2N}}
 ={\mathbb P}_{N,t_N}
 \quad\hbox{whenever}\quad t_N=4t_{2N}.                    \tag{5.7}
\]
\end{theorem}

\begin{proof}
Work on the gauge quotient and choose, inside each coarse square, a tree
for the new fine vertices.  Its link variables are pure gauge and integrate
to one.  The remaining three independent chord variables successively join
the four fine faces along a dual spanning tree.  Integrating each chord has
the form of Lemma \ref{lem:v19v2-gluing}; it merges two adjacent face
weights and adds their heat times.  After three integrations the surviving
loop is the coarse boundary, with time equal to the sum of all four fine
times.  Choices belonging to distinct coarse squares are disjoint after the
coarse boundary links are retained, so Fubini's theorem performs the
reductions independently.  The resulting boundary-link density is exactly
the coarse density (5.1), proving (5.7).  Equivalently, Lemma
\ref{lem:v19v2-axial} turns each coarse face lasso into an ordered product
of its four independent fine face increments, whose law is their fourfold
convolution.
\end{proof}

\begin{corollary}[The area clock]
\label{cor:v19v2-area-clock}
The choice
\[
 t_N=\frac{\tau}{N^2}                                     \tag{5.8}
\]
satisfies \(t_N=4t_{2N}\).  It therefore defines a nontrivial exactly
projective family of two-dimensional heat-kernel gauge measures along the
dyadic refinement tower.
\end{corollary}

\begin{proof}
Directly,
\(4t_{2N}=4\tau/(2N)^2=\tau/N^2=t_N\), and Theorem
\ref{thm:v19v2-subdivision} applies at every level.
\end{proof}

\subsection{Relation to the small-curvature chart}

\begin{proposition}[Distributional plaquette scale]
\label{prop:v19v2-rough-scale}
Under the area clock (5.8), a unit plaquette heat variable has logarithmic
scale \(t_N^{1/2}=O(N^{-1})\), not \(O(N^{-2})\).  For every fixed \(C>0\),
\[
 {\mathbb P}_{N,t_N}
 \left\{\|\log P_p\|\le\frac{C}{N^2}\right\}
 =O(N^{-d})                                                \tag{5.9}
\]
for one fixed plaquette, where \(d=\dim K>0\).
\end{proposition}

\begin{proof}
Apply the heat-kernel upper bound of Lemma
\ref{lem:v19v1-small-ball} with \(t=\tau/N^2\) and \(r=C/N^2\):
\[
 r^dt^{-d/2}
 =C^dN^{-2d}\left(\frac{N^2}{\tau}\right)^{d/2}
 =C^d\tau^{-d/2}N^{-d}.
\]
The square-root heat scale follows from the same parametrix.
\end{proof}

\begin{assessment}[What the solvable model changes]
\label{ass:v19v2-change}
V1 correctly excludes independent link increments from the smooth chart.
V2 shows that an exactly projective interacting face law can exist with the
area clock, but its quantum curvature is distributional and still lies
outside the global \(O(N^{-2})\) plaquette chart with high probability.
The V94 graph is therefore a controlled classical or semiclassical chart,
not a candidate global coordinate for a typical quantum gauge field.
\end{assessment}

\subsection{Executable certificate}

The verifier
\[
 \texttt{scripts/verify\_sm\_v19\_v2\_planar\_heat\_kernel\_block.py}
                                                                  \tag{5.10}
\]
checks the planar identity \(|E|-|V|+1=|F|\) on rectangular samples,
verifies an exact fourfold convolution over \(\mathbb Z/3\mathbb Z\), and
checks the area-clock identity.  It emits
\[
 \texttt{artifacts/sm\_v19\_v2\_planar\_heat\_kernel\_block.json}.  \tag{5.11}
\]
The hashes are
\[
\begin{array}{c|l}
\hbox{object}&\hbox{SHA--256}\\ \hline
\hbox{verifier}&
\texttt{EFCEAABF7F8AB8B6C419AFDFC718A1E94C7370BB6767F696790A7C14C7295DBB}\\
\hbox{canonical payload}&
\texttt{98426538D41FB20856833A5B8486864CAAF15293F0C39E49B39F5ECF4B9370C9}.
\end{array}                                                   \tag{5.12}
\]

\begin{firewall}[Two-dimensional exactness does not transfer automatically]
\label{fire:v19v2}
The proof uses planar tree gauge and elimination of internal edges between
two faces.  In four dimensions, one link belongs to several plaquettes and
the plaquette variables satisfy nonabelian cube constraints.  The same
integration no longer factorizes.  V2 proves no four-dimensional
Yang--Mills measure, matter coupling, reflection positivity, or Einstein
continuum.
\end{firewall}

\begin{decision}[Next acquisition]
\label{dec:v19v2-next}
V3 should identify the exact four-dimensional obstruction by integrating
one fine link shared by several plaquettes.  Peter--Weyl expansion should
show the resulting multi-face intertwiner rather than a single heat kernel.
The programme should then decide whether a spin-foam or character
representation is the correct projective coordinate for the gauge sector.
\end{decision}

\begin{record}[V2 append record]
\label{rec:v19v2}
V2 is appended to the validated nineteenth-cycle V1 whose SHA--256 was
\texttt{11249B2F99E97ED818A39DE3CD11437F1D0B0EDFE9810BC345E0F4003FBCDADA}.
After installation only V2 is the active numeric SM note; the
\texttt{V100\_f18} archive and unrelated notes are unchanged.
\end{record}


\section{V3: the four-dimensional shared-link intertwiner}
\label{sec:v19v3}

V2 works because planar gauge fixing gives one independent group variable
per face.  In four dimensions an interior link belongs to six plaquettes.
This section integrates that link exactly in Peter--Weyl variables.  The
result is the orthogonal projector onto a six-fold invariant tensor space,
not another one-face heat kernel.  An abelian Fourier specialization proves
that the induced six-face coupling cannot factor into independent face
weights.

\subsection{The Haar projector}

Let \(\widehat K\) be the unitary dual of a compact group \(K\).  For
\(\lambda\in\widehat K\), let \(V_\lambda\) be its representation space,
\(d_\lambda=\dim V_\lambda\), and \(c_\lambda\ge0\) its heat eigenvalue.
With normalized Haar measure,
\[
 q_t(g)=\sum_{\lambda\in\widehat K}
 d_\lambda e^{-c_\lambda t}\chi_\lambda(g).                \tag{6.1}
\]
The series is absolutely and uniformly convergent for every \(t>0\).

\begin{lemma}[Haar integration is invariant projection]
\label{lem:v19v3-haar-projector}
For a finite-dimensional unitary representation
\(\rho:K\to U(V)\),
\[
 \Pi_\rho:=\int_K\rho(u)\,du                               \tag{6.2}
\]
is the orthogonal projection of \(V\) onto
\[
 V^K=\{v:\rho(g)v=v\text{ for every }g\in K\}.             \tag{6.3}
\]
\end{lemma}

\begin{proof}
Haar invariance gives
\(\rho(g)\Pi_\rho=\Pi_\rho=\Pi_\rho\rho(g)\).  Hence the range is contained
in \(V^K\), and \(\Pi_\rho v=v\) for \(v\in V^K\).
Moreover,
\[
 \Pi_\rho^*=\int_K\rho(u^{-1})\,du=\Pi_\rho
\]
and
\[
 \Pi_\rho^2=\int_{K^2}\rho(uv)\,du\,dv=\Pi_\rho.
\]
Thus it is the orthogonal projection onto precisely \(V^K\).
\end{proof}

\subsection{One link shared by six faces}

Fix an interior link in direction \(\mu\) of a four-dimensional cubical
lattice.  For each of the three transverse directions there is one incident
plaquette on either side of the link.  Thus the incidence number is
\[
 2(4-1)=6.                                                  \tag{6.4}
\]
After all other links are held fixed, the six heat factors have the form
\[
 I(A,B):=\int_K\prod_{j=1}^6
 q_{t_j}(A_j u^{\epsilon_j}B_j)\,du,\qquad
 \epsilon_j\in\{1,-1\}.                                    \tag{6.5}
\]

\begin{theorem}[Exact six-face Peter--Weyl coupling]
\label{thm:v19v3-six-face}
Expand (6.5) by (6.1).  For every label tuple
\(\boldsymbol\lambda=(\lambda_1,\ldots,\lambda_6)\), the link integral is a
contraction of the exterior matrices with the orthogonal projector
\[
 \Pi_{\boldsymbol\lambda,\boldsymbol\epsilon}
 =\int_K
 \bigotimes_{\epsilon_j=1}\pi_{\lambda_j}(u)
 \otimes
 \bigotimes_{\epsilon_j=-1}\pi_{\lambda_j}(u)^*
 \,du                                                     \tag{6.6}
\]
onto
\[
 \operatorname{Inv}\left(
 \bigotimes_{\epsilon_j=1}V_{\lambda_j}\otimes
 \bigotimes_{\epsilon_j=-1}V_{\lambda_j}^*
 \right).                                                  \tag{6.7}
\]
Consequently one-link elimination generates a joint six-face intertwiner.
\end{theorem}

\begin{proof}
Absolute convergence permits termwise multiplication and Haar integration
in (6.5).  Write each character as the trace of
\(\pi_{\lambda_j}(A_j)\pi_{\lambda_j}(u^{\epsilon_j})
\pi_{\lambda_j}(B_j)\).  The tensor product collects all occurrences of
\(u\).  For negative orientation, unitarity identifies
\(\pi_\lambda(u^{-1})\) with the dual factor.  Lemma
\ref{lem:v19v3-haar-projector} evaluates their Haar integral as (6.6).
All remaining matrices contract with its tensor indices.  This gives the
stated joint coefficient for every label tuple.
\end{proof}

\begin{remark}[Why the two-face gluing lemma is special]
\label{rem:v19v3-two-face}
For two oppositely oriented factors, Schur orthogonality says
\[
 \operatorname{Inv}(V_\lambda\otimes V_\eta^*)=
 \begin{cases}
 \mathbb C\,I_{V_\lambda},&\lambda=\eta,\\
 \{0\},&\lambda\ne\eta.
 \end{cases}                                                \tag{6.8}
\]
The single surviving contraction is precisely the heat-kernel convolution
of V2.  With six factors, fusion multiplicities and several admissible
intermediate representations can occur.
\end{remark}

\subsection{An exact abelian nonfactorization}

Take \(K=U(1)\), write the exterior face angles as
\(\theta_1,\ldots,\theta_6\), and orient three incidences each way.  Fourier
expansion gives
\[
\begin{split}
 I(\theta)
 =\sum_{\substack{n_1,\ldots,n_6\in\mathbb Z\\
                   \sum_j\epsilon_jn_j=0}}
 \exp\left(-\sum_{j=1}^6t_jn_j^2\right)
 \exp\left(i\sum_{j=1}^6n_j\theta_j\right).                \tag{6.9}
\end{split}
\]

\begin{proposition}[The six-face coefficient support is not Cartesian]
\label{prop:v19v3-nonfactor}
The function in (6.9) cannot be written as
\(\prod_{j=1}^6 f_j(\theta_j)\) with nonzero one-face functions \(f_j\).
Hence shared-link integration fails to preserve independent face weights
even for an abelian positive-dimensional group.
\end{proposition}

\begin{proof}
Take \(\epsilon=(1,1,1,-1,-1,-1)\).  The Fourier coefficient at
\[
 (n_1,\ldots,n_6)=(1,-1,0,0,0,0)
\]
is strictly positive, while the coefficients at
\((1,0,0,0,0,0)\) and \((0,-1,0,0,0,0)\) vanish because they violate charge
conservation.  The global constant Fourier coefficient is strictly positive.  Under a
product factorization it is the product of the six one-variable constant
coefficients, so each of those constants is nonzero.  The nonzero joint
coefficient then forces the first-mode coefficient of \(f_1\) and the
minus-first-mode coefficient of \(f_2\) to be nonzero.  Multiplying either
one by the five nonzero constant coefficients makes the corresponding
single-mode coefficient nonzero, a contradiction.
\end{proof}

\begin{corollary}[The obstruction is structural]
\label{cor:v19v3-structural}
No change of heat times can turn the four-dimensional one-link integral
into a product of six independent one-face heat weights.  The obstruction
is the invariant constraint in (6.7), already visible as charge
conservation in (6.9).
\end{corollary}

\begin{proof}
Changing positive \(t_j\) changes the strictly positive weights of supported
Fourier coefficients but does not change the constraint
\(\sum_j\epsilon_jn_j=0\).  Proposition
\ref{prop:v19v3-nonfactor} therefore remains valid.
\end{proof}

\subsection{The exact finite spin-foam form}

\begin{theorem}[Character expansion of the finite gauge partition function]
\label{thm:v19v3-spin-foam}
On any finite cubical lattice with positive face times, expanding every
heat kernel and integrating every link writes the partition function as a
sum over irreducible face labels.  Each face contributes
\(d_\lambda e^{-c_\lambda t}\), each link contributes the invariant
projector (6.6) for its incident labelled faces, and vertex tensor indices
are contracted according to the lattice incidence pattern.
\end{theorem}

\begin{proof}
Apply (6.1) to every face.  Positivity of the heat times gives absolute
uniform convergence on the finite product group, so sums and Haar integrals
may be interchanged.  Apply Theorem \ref{thm:v19v3-six-face} independently
to every link.  Every representation index occurs at the endpoints of
incidences and is therefore contracted at vertices.  No approximation is
made.
\end{proof}

\begin{assessment}[V3 acquisition]
\label{ass:v19v3}
The exact two-dimensional convolution mechanism fails in four dimensions
because a shared link produces a six-face invariant tensor.  Peter--Weyl
face labels and link intertwiners are therefore the natural exact finite
coordinates for refinement.  This identifies a representation change; it
does not yet construct a projective interacting measure.
\end{assessment}

\subsection{Executable certificate}

The verifier
\[
 \texttt{scripts/verify\_sm\_v19\_v3\_shared\_link\_intertwiner.py} \tag{6.10}
\]
checks the six-plaquette incidence count, enumerates the non-Cartesian
\(U(1)\) charge support on a finite mode window, and verifies idempotence of
the corresponding Haar invariant projector over
\(\mathbb Z/3\mathbb Z\).  It emits
\[
 \texttt{artifacts/sm\_v19\_v3\_shared\_link\_intertwiner.json}.    \tag{6.11}
\]
The hashes are
\[
\begin{array}{c|l}
\hbox{object}&\hbox{SHA--256}\\ \hline
\hbox{verifier}&
\texttt{5B42ED4C9A119A04A542196FDC69516E3791833BF10F7E22FBD32EA3CC421CD8}\\
\hbox{canonical payload}&
\texttt{FE5C1176A73EC21501BC319EEB500856B11A3E6E97E5E89D2C4BED14FCD78411}.
\end{array}                                                   \tag{6.12}
\]

\begin{firewall}[Representation coordinates do not solve the measure problem]
\label{fire:v19v3}
The spin-foam expansion is an exact rewriting of a finite positive link
integral.  Individual recoupling coefficients need not be positive in an
arbitrary intertwiner basis, and no refinement consistency, uniform
summability, reflection positivity after blocking, matter determinant, or
Einstein limit is proved.
\end{firewall}

\begin{decision}[Next acquisition]
\label{dec:v19v3-next}
V4 should isolate a single shared-link projector as a positive operator and
derive its contraction and norm-one properties under gluing.  It should
then test whether a coarse-boundary transfer kernel built from these
projectors is positivity preserving.  A global spin-foam positivity claim
must be withheld unless the vertex contractions are controlled.
\end{decision}

\begin{record}[V3 append record]
\label{rec:v19v3}
V3 is appended to the validated V2 whose installed SHA--256 was
\texttt{5059BDD20DAC94384E75EEFDA339476CBD0939F722B84FE73D9E769269A3ADE4}.
After installation only V3 is the active numeric SM note; the frozen
archives and unrelated notes are unchanged.
\end{record}


\section{V4: positive Haar projectors and a finite reflection-positive transfer kernel}
\label{sec:v19v4}

V3 showed that a four-dimensional link integral is a multi-face
intertwiner.  Although this destroys face independence, the intertwiner is
an orthogonal projection and hence carries a strong positivity statement.
This section places that projection inside the heat-kernel transfer
operator and proves finite-cutoff reflection positivity for the pure-gauge
heat-kernel layer.  The result is not yet uniform in the cutoff and contains
no matter determinant.

\subsection{Positive contractions}

Let \({\cal H}=L^2(K^m)\) for a finite collection of boundary links.  Let
\(T_{\mathbf t}\) be the tensor product of the heat convolution operators
\[
 (T_tf)(x)=\int_Kq_t(xy^{-1})f(y)\,dy.                     \tag{7.1}
\]

\begin{lemma}[Heat convolution is a positive self-adjoint contraction]
\label{lem:v19v4-heat-positive}
For \(t\ge0\), \(T_t\) is positivity preserving,
self-adjoint, and
\[
 0\le T_t\le I,\qquad \|T_t\|_{2\to2}=1.                  \tag{7.2}
\]
The same assertions hold for \(T_{\mathbf t}\).
\end{lemma}

\begin{proof}
Nonnegativity of \(q_t\) makes \(T_t\) positivity preserving.  Symmetry
\(q_t(g^{-1})=q_t(g)\) and Haar invariance give self-adjointness.
Peter--Weyl decomposition diagonalizes \(T_t\): on every matrix-coefficient
space of an irreducible \(\lambda\), its eigenvalue is
\(e^{-c_\lambda t}\in[0,1]\).  The trivial representation gives eigenvalue
one.  This proves (7.2).  Tensor products preserve positivity,
self-adjointness, and the norm-one bound.
\end{proof}

Let a boundary gauge group act unitarily on \({\cal H}\), and define
\[
 {\cal P}=\int_{{\cal G}_{\partial}}R(g)\,dg.              \tag{7.3}
\]

\begin{lemma}[Gauge projection and heat flow]
\label{lem:v19v4-gauge-projection}
The operator \({\cal P}\) is an orthogonal projection of norm one.
For central heat kernels,
\[
 {\cal P}T_{\mathbf t}=T_{\mathbf t}{\cal P}.              \tag{7.4}
\]
Consequently
\[
 {\cal K}_0:={\cal P}T_{\mathbf t}{\cal P}                \tag{7.5}
\]
is a positive self-adjoint contraction on \({\cal H}\).
\end{lemma}

\begin{proof}
Lemma \ref{lem:v19v3-haar-projector}, applied to the boundary gauge
representation, proves the first statement.  A central heat semigroup
commutes with left and right translations on every link, hence with every
\(R(g)\) and with its Haar average.  The two commuting operators
\({\cal P}\) and \(T_{\mathbf t}\) are positive self-adjoint contractions,
so their product is again positive and self-adjoint.  Its norm is at most
one.
\end{proof}

\begin{corollary}[Shared-link contraction cannot create a negative square]
\label{cor:v19v4-link-square}
For every tensor \(\xi\) in the six-face representation space of V3,
\[
 \langle\xi,\Pi_{\boldsymbol\lambda,\boldsymbol\epsilon}\xi\rangle
 =\|\Pi_{\boldsymbol\lambda,\boldsymbol\epsilon}\xi\|^2
 \ge0,                                                     \tag{7.6}
\]
and the absolute value of the contraction is at most \(\|\xi\|^2\).
\end{corollary}

\begin{proof}
The link intertwiner is the orthogonal projection of Lemma
\ref{lem:v19v3-haar-projector}.  Idempotence and self-adjointness give
(7.6); norm one gives the upper bound.
\end{proof}

\subsection{A transfer-kernel criterion}

Let \(M_w\) denote multiplication by a bounded nonnegative boundary weight
\(w\).

\begin{theorem}[Positive transfer sandwich]
\label{thm:v19v4-sandwich}
The operator
\[
 {\cal K}=M_w^{1/2}{\cal P}T_{\mathbf t}{\cal P}M_w^{1/2} \tag{7.7}
\]
is positive and self-adjoint.  It is also positivity preserving.  If
\(0\le w\le W\), then
\[
 \|{\cal K}\|_{2\to2}\le W.                               \tag{7.8}
\]
\end{theorem}

\begin{proof}
By Lemma \ref{lem:v19v4-gauge-projection},
\[
 \langle f,{\cal K}f\rangle
 =\left\langle M_w^{1/2}f,
 {\cal P}T_{\mathbf t}{\cal P}M_w^{1/2}f\right\rangle
 =\left\|T_{\mathbf t}^{1/2}{\cal P}M_w^{1/2}f\right\|^2
 \ge0.                                                     \tag{7.9}
\]
The displayed factorization also gives self-adjointness.  Multiplication by
\(w^{1/2}\), heat convolution, and gauge averaging each preserve
nonnegative functions, proving positivity preservation.  Finally,
\(\|M_w^{1/2}\|\le W^{1/2}\), while the middle operator has norm at most
one, giving (7.8).
\end{proof}

\begin{theorem}[Finite pure-gauge heat-kernel reflection positivity]
\label{thm:v19v4-finite-rp}
Consider a finite hypercubic slab symmetric about a time-reflection plane,
with heat-kernel plaquette weights and free or reflection-invariant boundary
conditions.  Split every plaquette weight lying in one time half into the
multiplication factors assigned to that half, and put the crossing temporal
plaquettes into heat convolution.  After temporal gauge averaging, the
one-step transfer operator has the form (7.7).  Therefore every bounded
gauge-invariant observable \(F\) supported in the positive half satisfies
\[
 \mathbb E\!\left[\overline{F(\theta U)}F(U)\right]\ge0.    \tag{7.10}
\]
\end{theorem}

\begin{proof}
Fix the spatial links on the reflection slice.  Temporal plaquette
holonomies couple the two adjacent spatial-link configurations through
kernels \(q_t(xy^{-1})\); their product is \(T_{\mathbf t}\).
Integration over temporal vertex gauges is exactly \({\cal P}\).
Spatial plaquettes wholly in either half contribute conjugate
multiplication factors; their common slice contribution can be split as
\(M_w^{1/2}\) on each side.  Thus the reflected integral is
\[
 \langle\Phi_F,
 M_w^{1/2}{\cal P}T_{\mathbf t}{\cal P}M_w^{1/2}\Phi_F\rangle
\]
for the boundary amplitude \(\Phi_F\) obtained by integrating the positive
half.  Theorem \ref{thm:v19v4-sandwich} makes this nonnegative.
Approximation in \(L^2\) extends the calculation from continuous cylinder
functions to bounded gauge-invariant \(F\).
\end{proof}

\begin{remark}[Relation to the spin-foam expansion]
\label{rem:v19v4-spin}
In Peter--Weyl variables, (7.9) is the statement that link intertwiners are
contracted as orthogonal projectors before the two reflected halves are
paired.  Individual recoupling coordinates may have signs, but the complete
boundary contraction remains a nonnegative square.  This does not imply
termwise positivity of an arbitrary spin-foam basis.
\end{remark}

\subsection{Executable certificate}

The verifier
\[
 \texttt{scripts/verify\_sm\_v19\_v4\_positive\_transfer\_kernel.py} \tag{7.11}
\]
uses \(\mathbb Z/3\mathbb Z\) to check an exact positive-type convolution
matrix, a two-link tensor heat operator, the self-adjoint idempotent gauge
average, commutation, and the projected transfer range.  It emits
\[
 \texttt{artifacts/sm\_v19\_v4\_positive\_transfer\_kernel.json}.   \tag{7.12}
\]
The hashes are
\[
\begin{array}{c|l}
\hbox{object}&\hbox{SHA--256}\\ \hline
\hbox{verifier}&
\texttt{289D48FCC44EEFE5C9DA70ABE6C5D30BADCFCAC0E85F7D86E154E4DA1D604483}\\
\hbox{canonical payload}&
\texttt{D3A801310490DCF56C8954500EC261990D4D2820340C67F7BDA7B1F8EE2C6090}.
\end{array}                                                   \tag{7.13}
\]

\begin{firewall}[Finite reflection positivity is not OS reconstruction]
\label{fire:v19v4}
Theorem \ref{thm:v19v4-finite-rp} is a finite pure-gauge statement.
It gives neither cutoff-uniform compactness nor continuity of the limiting
time translations.  Fermion determinants, gauge fixing and ghosts, Higgs
couplings, the Einstein measure, vacuum uniqueness, clustering, and a
Hamiltonian spectrum remain untreated.
\end{firewall}

\begin{decision}[Mandatory checkpoint, then matter interface]
\label{dec:v19v4-next}
V5 is the mandatory companion audit.  V6 should then ask which local Higgs
and fermion factors preserve the transfer sandwich.  Bosonic factors can be
split when their cross-plane kernel is positive definite; a chiral fermion
determinant requires a separate reflected-Gram or Pfaffian-line argument
and must not be inferred from the pure-gauge proof.
\end{decision}

\begin{assessment}[V4 acquisition]
\label{ass:v19v4}
V4 retains a rigorous positivity structure after the V3 factorization
failure: every shared-link intertwiner is a norm-one orthogonal projector,
and the finite heat-kernel pure-gauge transfer operator is a positive
sandwich.  This proves finite reflection positivity for that sector only.
\end{assessment}

\begin{record}[V4 append record]
\label{rec:v19v4}
V4 is appended to the validated V3 whose installed SHA--256 was
\texttt{193342397C90669EF1FBF97F32B800D4607B4C6B7246021202B671F0C8FE3AAF}.
After installation only V4 is the active numeric SM note; frozen archives
and unrelated notes are unchanged.
\end{record}


\section{V5: mandatory companion audit before the matter interface}
\label{sec:v19v5}

V4 proved finite reflection positivity for the pure-gauge heat-kernel
transfer sandwich.  The standing five-version rule requires inspection of
the active companion notes before matter is added.

\begin{audit}[Companion files]
\label{aud:v19v5-files}
The active Yang--Mills file remains
\[
 \texttt{Renewal\_Yang\_Mills\_Mass\_Gap\_Research\_Notes\_V25.tex},
\quad
 \operatorname{SHA256}=
 \texttt{5D6759B911E076304A17D09955E6777C3ACD6C99D3B47CD42BBA75EB0EA96BE6}.
                                                                  \tag{8.1}
\]
The active Navier--Stokes file remains
\[
 \texttt{Renewal\_NS\_Stability\_Research\_Notes\_V26.tex},
\quad
 \operatorname{SHA256}=
 \texttt{82C887F8C2C69E4F28FAD7C3DC8DE5B9099CB5A4BF431EBA1FFF3E91935978AD}.
                                                                  \tag{8.2}
\]
Neither file has changed since the V100 audit of the frozen eighteenth
cycle.
\end{audit}

\begin{proposition}[No repeated import]
\label{prop:v19v5-no-import}
The companion state adds no new theorem to V1--V4.  Yang--Mills V25 still
separates its finite spectral certificate from the later exact-link
cofinal route, and Navier--Stokes V26 still concerns rank-one Oseen-kernel
norms.  The V6 matter-interface calculation should therefore proceed from
the transfer sandwich already proved in V4.
\end{proposition}

\begin{proof}
The file hashes in (8.1)--(8.2) are identical to those recorded at the
immediately preceding mandatory audit.  Their relevant conclusions and
claim firewalls were already recorded there.  Re-importing them would
duplicate rather than advance the proof record.  Neither note contains a
new cross-plane Higgs kernel or a fermionic reflected-Gram theorem.
\end{proof}

\begin{decision}[V6 acquisition]
\label{dec:v19v5-next}
V6 will prove a positive-definite kernel theorem for reflection-matched
bosonic hopping terms by an explicit symmetric-Fock feature map.  It will
then insert those kernels into the V4 gauge transfer sandwich.  Berezin
fermions and the chiral determinant line will remain a separate gate.
\end{decision}

\begin{firewall}[Audit boundary]
\label{fire:v19v5}
The absence of a new companion input proves no matter-sector positivity.
It only confirms that the next calculation is not duplicating a newer
parallel result.
\end{firewall}

\begin{record}[V5 append record]
\label{rec:v19v5}
V5 is appended to the validated V4 whose installed SHA--256 was
\texttt{17F6ADF7389893EA4DAD6462035FFB9A168B0A13CED7FEB3CAE6790702CE408A}.
After installation only V5 is the active numeric SM note.
\end{record}


\section{V6: the Fock-positive Higgs hinge and the metric-asymmetry gate}
\label{sec:v19v6}

The V4 transfer theorem supplies a positive pure-gauge boundary operator, and
V5 found no newer companion result to import.  This note now treats the
bosonic matter hinge.  The calculation has two outcomes.  A reflection-
balanced Higgs hopping term is a positive kernel by an explicit tensor-Fock
factorization.  A nonzero fixed Duhamel metric weight destroys the symmetry
of a single hinge, but reflection changes \(d\) to \(-d\), so the paired
two-step operator is still a positive square.

\subsection{Positive kernels from the real tensor Fock space}

Let \(H\) be the underlying real Hilbert space of a finite-dimensional
unitary Higgs representation; thus
\[
 (x,y)_{\mathbb R}:=\operatorname{Re}\langle x,y\rangle_{\mathbb C}.
\]
For \(c\geq0\), put
\[
 E_c(x,y):=\exp\!\big(c(x,y)_{\mathbb R}\big).                 \tag{9.1}
\]

\begin{lemma}[Tensor-Fock Gram identity]
\label{lem:v19v6-fock}
The kernel \(E_c\) is positive definite.  More precisely, in the Hilbert
direct sum
\[
 {\cal F}_{\mathbb R}(H):=\bigoplus_{n=0}^{\infty}H^{\otimes n},
 \qquad H^{\otimes0}:=\mathbb R,
\]
the vector
\[
 \Phi_c(x):=\left(\sqrt{\frac{c^n}{n!}}\,x^{\otimes n}\right)_{n\geq0}
                                                               \tag{9.2}
\]
is well defined and satisfies
\[
 \langle\Phi_c(x),\Phi_c(y)\rangle_{{\cal F}_{\mathbb R}(H)}
 =E_c(x,y).                                                    \tag{9.3}
\]
Consequently, for every \(x_1,\ldots,x_r\in H\) and
\(\alpha_1,\ldots,\alpha_r\in\mathbb C\),
\[
 \sum_{i,j=1}^r\overline{\alpha_i}\alpha_j E_c(x_i,x_j)
 =\sum_{n=0}^{\infty}\frac{c^n}{n!}
   \left\|\sum_{j=1}^r\alpha_jx_j^{\otimes n}\right\|^2
 \geq0.                                                       \tag{9.4}
\]
\end{lemma}

\begin{proof}
The tensor norm is multiplicative, so
\[
 \|\Phi_c(x)\|^2
 =\sum_{n=0}^{\infty}\frac{c^n}{n!}\|x\|^{2n}
 =e^{c\|x\|^2}<\infty.
\]
Moreover,
\(\langle x^{\otimes n},y^{\otimes n}\rangle=(x,y)_{\mathbb R}^n\).
Summing the exponential series proves (9.3).  Expanding the squared norms
in the right side of (9.4) gives its left side term by term.  Every summand
is nonnegative.
\end{proof}

\begin{corollary}[Diagonal dressing and products]
\label{cor:v19v6-dressing}
If \(h:H\to\mathbb C\), then
\[
 K_h(x,y):=\overline{h(x)}h(y)E_c(x,y)                         \tag{9.5}
\]
is positive definite.  A finite pointwise product of kernels of the form
(9.5), possibly on different field factors, is positive definite.
\end{corollary}

\begin{proof}
Equation (9.5) is the Gram kernel of the feature map
\(x\mapsto h(x)\Phi_c(x)\).  A product of Gram kernels is the Gram kernel of
the tensor product of their feature maps.
\end{proof}

\subsection{The balanced gauge--Higgs crossing edge}

It is useful to state reflection positivity with the transport
identification visible.  Let \(H_+\) and \(H_-\) be the field spaces on the
two sides of a reflection plane and let \(r:H_+\to H_-\) be the orthogonal
identification induced by reflection and gauge parallel transport.  A
cross kernel \(W:H_+\times H_-\to\mathbb C\) is called positive relative to
\(r\) when
\[
 \widetilde W(x,y):=W(x,r(y))                                  \tag{9.6}
\]
is a positive-definite kernel on \(H_+\).

\begin{theorem}[Balanced Higgs hinge]
\label{thm:v19v6-balanced}
Let \(\chi>0\), let \(U\) be a unitary gauge transport in the Higgs
representation, and identify the negative field by \(r_U(y)=U^{-1}y\).
Then the balanced hopping weight
\[
 W_0(x,z):=\exp\!\left[-\chi\|x-Uz\|^2\right]                  \tag{9.7}
\]
is positive relative to \(r_U\).  If \(v:H\to[0,\infty)\) is a
reflection-invariant on-site weight, the dressed kernel
\[
 v(x)^{1/2}W_0(x,r_U(y))v(y)^{1/2}                             \tag{9.8}
\]
is also positive definite.  Products over finitely many crossing Higgs
edges retain this property.
\end{theorem}

\begin{proof}
After the identification \(z=r_U(y)\), unitarity gives
\[
 W_0(x,r_U(y))
 =e^{-\chi\|x-y\|^2}
 =e^{-\chi\|x\|^2}e^{-\chi\|y\|^2}
  E_{2\chi}(x,y).                                               \tag{9.9}
\]
Apply Corollary \ref{cor:v19v6-dressing} with
\(h(x)=v(x)^{1/2}e^{-\chi\|x\|^2}\).  The product assertion follows from
the last part of that corollary.
\end{proof}

\begin{theorem}[Finite bosonic gauge--Higgs reflection positivity]
\label{thm:v19v6-bosonic-rp}
Consider the finite reflected slab of Theorem
\ref{thm:v19v4-finite-rp}.  Put the crossing temporal links in temporal
gauge, identify the reflected Higgs fibres by their unitary transports, and
assume every crossing Higgs edge is balanced as in (9.7).  Let the remaining
Higgs action be reflection invariant and give a nonnegative integrable
one-side weight.  Then the heat-kernel gauge theory coupled to these Higgs
fields is reflection positive: for every bounded gauge-invariant
positive-half observable \(F\),
\[
 \mathbb E\!\left[\overline{F(\theta U,\theta\phi)}
                         F(U,\phi)\right]\geq0.                 \tag{9.10}
\]
\end{theorem}

\begin{proof}
For one balanced real \(m\)-dimensional Higgs fibre, the integral operator
with kernel \(e^{-\chi\|x-y\|^2}\) is positive on \(L^2(\mathbb R^m)\):
by Plancherel its quadratic form is, up to the positive Fourier-normalizing
constant,
\[
 \int_{\mathbb R^m}
 e^{-\|\xi\|^2/(4\chi)}|\widehat f(\xi)|^2\,d\xi\geq0.          \tag{9.11}
\]
It is bounded because the Gaussian is in \(L^1\).  Equivalently, its
finite evaluation matrices are positive by Theorem
\ref{thm:v19v6-balanced}.  Tensoring over crossing sites and dressing by
the square roots of the one-side weights preserves positivity.

The pure-gauge crossing operator is positive by Theorem
\ref{thm:v19v4-sandwich}.  Hence its tensor product with the Higgs crossing
operator is positive.  Simultaneous boundary gauge averaging is an
orthogonal projection \({\cal P}_{G,H}\), so
\[
 {\cal P}_{G,H}(T_G\otimes T_H){\cal P}_{G,H}\geq0.             \tag{9.12}
\]
All action terms lying strictly in either half form conjugate one-side
amplitudes, while reflection-plane terms are split into equal nonnegative
square roots.  Integrating the positive half therefore produces a boundary
amplitude \(\Psi_F\), and the reflected expectation is
\[
 \left\langle\Psi_F,\,
 M^{1/2}{\cal P}_{G,H}(T_G\otimes T_H)
 {\cal P}_{G,H}M^{1/2}\Psi_F\right\rangle\geq0.                 \tag{9.13}
\]
The integrability assumption justifies Fubini and approximation from
compactly supported cylinder functions to bounded \(F\).
\end{proof}

\subsection{The Duhamel metric weight is a one-step obstruction}

For \(d\in\mathbb R\), set \(a=e^{-d/2}\), \(b=e^{d/2}\), so \(ab=1\).
After transporting the negative fibre to the positive one, the
metric-weighted hopping kernel is
\[
 W_d(x,y)
 :=\exp\!\left[-\chi\|ax-by\|^2\right]
 =e^{-\chi a^2\|x\|^2}e^{-\chi b^2\|y\|^2}E_{2\chi}(x,y).
                                                               \tag{9.14}
\]

\begin{proposition}[Sharp symmetry obstruction]
\label{prop:v19v6-asymmetry}
If \(H\neq\{0\}\), \(\chi>0\), and \(d\neq0\), then \(W_d\) is not a
positive-definite kernel on \(H\).  Thus the balanced theorem cannot be
applied to one fixed asymmetric crossing edge.
\end{proposition}

\begin{proof}
Every complex positive-definite kernel is Hermitian.  The present kernel is
real, so it would have to be symmetric.  Choose \(y\neq0\).  From (9.14),
\[
 W_d(0,y)=e^{-\chi b^2\|y\|^2},
 \qquad
 W_d(y,0)=e^{-\chi a^2\|y\|^2}.                                \tag{9.15}
\]
Because \(d\neq0\), \(a^2\neq b^2\), and the two values differ.  Hence
\(W_d\) is not Hermitian and cannot be positive definite.
\end{proof}

Reflection does, however, reverse the Duhamel weight.

\begin{theorem}[Reflected \(d,-d\) pair is a positive square]
\label{thm:v19v6-paired}
Let \(v=e^{-V}\) with
\[
 V(x)\geq\lambda\|x\|^4-C,\qquad \lambda>0,                    \tag{9.16}
\]
and define on \(L^2(H,dx)\) the integral operator
\[
 A_d(x,y):=v(x)^{1/2}W_d(x,y)v(y)^{1/2}.                       \tag{9.17}
\]
Then \(A_d\) is Hilbert--Schmidt,
\[
 A_d^*=A_{-d},                                                  \tag{9.18}
\]
and both paired transfers are positive:
\[
 A_dA_{-d}=A_dA_d^*\geq0,
 \qquad
 A_{-d}A_d=A_d^*A_d\geq0.                                     \tag{9.19}
\]
\end{theorem}

\begin{proof}
Since \(0<W_d\leq1\),
\[
 \iint |A_d(x,y)|^2\,dx\,dy
 \leq\left(\int_Hv(x)\,dx\right)^2<\infty                      \tag{9.20}
\]
by (9.16).  Thus \(A_d\) is Hilbert--Schmidt.  Directly from \(a(-d)=b(d)\)
and \(b(-d)=a(d)\),
\[
 W_d(y,x)=W_{-d}(x,y).                                         \tag{9.21}
\]
The kernel is real and \(v\) is the same on both sides, so (9.21) proves
(9.18).  Finally, for every \(f\in L^2(H)\),
\[
 \langle f,A_dA_d^*f\rangle=\|A_d^*f\|^2,\qquad
 \langle f,A_d^*A_df\rangle=\|A_df\|^2,                        \tag{9.22}
\]
which proves (9.19).
\end{proof}

\begin{remark}[Meaning for the Einstein coupling]
\label{rem:v19v6-einstein}
Theorem \ref{thm:v19v6-paired} gives a rigorous reflected two-step route
when the discrete metric variable changes sign under time reflection.
It does not turn a fixed \(d\neq0\) one-step kernel into a positive kernel.
A full metric transfer construction must either place a balanced
\(d=0\) hinge on the reflection plane or retain the reflected pair before
claiming positivity.
\end{remark}

\subsection{Executable certificate and remaining matter gate}

The exact verifier
\[
 \texttt{scripts/verify\_sm\_v19\_v6\_higgs\_fock\_kernel.py}  \tag{9.23}
\]
constructs the degree-five tensor-word Gram decomposition on a rational
sample, checks three quadratic forms as explicit sums of squares, verifies
the metric hopping identity with \(ab=1\), exhibits the asymmetric
one-step witness, and checks a rational finite model of
\(A_{-d}=A_d^{\mathsf T}\) and \(A_dA_d^{\mathsf T}\geq0\).
It emits
\[
 \texttt{artifacts/sm\_v19\_v6\_higgs\_fock\_kernel.json}.      \tag{9.24}
\]
The hashes are
\[
\begin{array}{c|l}
\hbox{object}&\hbox{SHA--256}\\ \hline
\hbox{verifier}&
\texttt{D1DA0852283DB265BB20F254D4E9BEDD168D88B9BEAF20EEA72E12B34C90A0E7}\\
\hbox{canonical payload}&
\texttt{9A3D14BAE51B6DB2A83357C374206EFCA2D02A61DE0D946D3B4954887E2FF875}.
\end{array}                                                     \tag{9.25}
\]

\begin{firewall}[Bosonic positivity does not integrate chiral fermions]
\label{fire:v19v6-fermions}
Berezin integration is not integration against a positive measure.
The Fock Gram map proves positivity for commuting Higgs variables only.
It gives no sign or phase control for a Wilson, overlap, or chiral
determinant, and it does not orient a Pfaffian or determinant line.  Those
claims require a separate reflected Grassmann Gram theorem together with
gauge-anomaly cancellation.
\end{firewall}

\begin{decision}[V7 acquisition]
\label{dec:v19v6-next}
V7 will formulate and prove a finite-dimensional reflected Berezin-square
criterion.  It will first test a vectorlike fermion block, where conjugate
flavours can form an absolute determinant square, and will isolate the
additional determinant-line and anomaly conditions needed before the
chiral Standard Model can enter the positive transfer construction.
\end{decision}

\begin{assessment}[V6 acquisition]
\label{ass:v19v6}
The finite pure-gauge positivity theorem now has a rigorous bosonic Higgs
extension at balanced reflection hinges.  The metric-weighted hinge has
also been classified sharply: its one-step kernel fails positivity for
every fixed nonzero \(d\), while the reflected \(d,-d\) composition is a
positive square under quartic confinement.  The next unsolved local sector
is fermionic.
\end{assessment}

\begin{record}[V6 append record]
\label{rec:v19v6}
V6 is appended to the validated V5 whose installed SHA--256 was
\texttt{CAA4105433438089AAB26DEF08D425D7CD15A08B4AC60810228CB630D2BF3676}.
After installation only V6 is the active numeric SM note; frozen archives
and unrelated notes are unchanged.
\end{record}



\section{V7: reflected Grassmann Gram matrices and the chiral determinant gate}
\label{sec:v19v7}

V6 extended the finite transfer argument through commuting Higgs fields and
showed why a fixed asymmetric metric hinge must be paired with its
reflection.  Fermions require a different positivity object because a
Berezin functional is not a positive measure.  This note constructs the
needed finite-dimensional object: a cone of reflected Grassmann squares.
It then proves that a positive temporal cross matrix puts the entire
fermionic cross exponential in that cone.

\subsection{The coefficient Hilbert space behind a Berezin functional}

Let
\[
 {\mathfrak A}_+=\Lambda(\mathbb C^n)
\]
be generated by \(\psi_1,\ldots,\psi_n\), and let
\({\mathfrak A}_-\) be a disjoint reflected copy.  Reflection is the
antilinear anti-automorphism
\[
 \Theta(\alpha A)=\overline\alpha\,\Theta(A),\qquad
 \Theta(AB)=\Theta(B)\Theta(A).                                \tag{10.1}
\]
For an increasing subset \(I=\{i_1<\cdots<i_k\}\), write
\(\psi_I=\psi_{i_1}\cdots\psi_{i_k}\).  Define the linear functional
\({\cal L}\) on the full finite Grassmann algebra by
\[
 {\cal L}\!\left(\Theta(\psi_I)\psi_J\right)=\delta_{IJ}.       \tag{10.2}
\]
The ordered monomials in (10.2) form a basis, so this definition is
unambiguous.

\begin{lemma}[Berezin coefficient norm]
\label{lem:v19v7-coefficient}
For \(A=\sum_I a_I\psi_I\),
\[
 {\cal L}\!\left(\Theta(A)A\right)=\sum_I|a_I|^2.              \tag{10.3}
\]
Moreover, \({\cal L}\) can be represented as a Berezin integral against a
unique Grassmann density once an order of all generators is fixed.
\end{lemma}

\begin{proof}
Antilinearity and (10.2) give
\[
 {\cal L}\!\left(\Theta(A)A\right)
 =\sum_{I,J}\overline{a_I}a_J
   {\cal L}\!\left(\Theta(\psi_I)\psi_J\right)
 =\sum_I|a_I|^2.
\]
For the second assertion, let \([Q]_{\rm top}\) denote the coefficient of
the ordered product of every generator.  The pairing
\((R,Q)\mapsto[RQ]_{\rm top}\) is nondegenerate: each monomial pairs,
up to a nonzero sign, with its complementary monomial.  Therefore there is
a unique density \(\rho_{\cal L}\) such that
\({\cal L}(Q)=[\rho_{\cal L}Q]_{\rm top}\), which is precisely a Berezin
integral in the chosen order.
\end{proof}

\begin{definition}[Reflected-square cone]
\label{def:v19v7-cone}
Let \({\cal C}_\Theta\) be the set of finite sums
\[
 W=\sum_{\alpha}\Theta(C_\alpha)C_\alpha,                     \tag{10.4}
\]
where each \(C_\alpha\in{\mathfrak A}_+\) is homogeneous.  This is the
finite reflected-square cone.
\end{definition}

\begin{proposition}[Cone criterion]
\label{prop:v19v7-cone}
If \(W\in{\cal C}_\Theta\), then every even
\(F\in{\mathfrak A}_+\) satisfies
\[
 {\cal L}\!\left(\Theta(F)F\,W\right)\geq0.                    \tag{10.5}
\]
\end{proposition}

\begin{proof}
It suffices to treat one term in (10.4).  Because \(F\) is even, moving it
past \(\Theta(C_\alpha)\) introduces no graded sign, and \(F\) commutes
with \(C_\alpha\).  The anti-automorphism law then gives
\[
 \Theta(F)F\Theta(C_\alpha)C_\alpha
 =\Theta(FC_\alpha)(FC_\alpha).
\]
Lemma \ref{lem:v19v7-coefficient} turns its \({\cal L}\)-value into the
squared coefficient norm of \(FC_\alpha\).  Summing over \(\alpha\) proves
(10.5).
\end{proof}

The restriction to even \(F\) is the physical observable algebra used
here.  Odd fermion fields require the usual twisted reflection convention
and are not reconstructed by Proposition \ref{prop:v19v7-cone}.

\subsection{Positive cross matrices generate reflected squares}

Let \(C\) be an \(n\times n\) Hermitian matrix and set
\[
 X_C:=\sum_{i,j=1}^n\Theta(\psi_i)C_{ij}\psi_j,\qquad
 W_C:=e^{X_C}.                                                  \tag{10.6}
\]
The exponential is a finite polynomial because \(X_C^{n+1}=0\).

\begin{lemma}[Grassmann minor expansion]
\label{lem:v19v7-minors}
With the reflection convention (10.1),
\[
 W_C
 =\sum_{k=0}^n\ \sum_{\substack{|I|=k\\|J|=k}}
   \det(C_{I,J})\,\Theta(\psi_I)\psi_J.                        \tag{10.7}
\]
\end{lemma}

\begin{proof}
Expand \(e^{X_C}=\sum_{k=0}^nX_C^k/k!\).  A nonzero degree-\(k\) term uses
distinct negative and positive generators.  Antisymmetrizing the \(k!\)
orders of its positive indices produces the determinant
\(\det(C_{I,J})\).  Moving all positive generators to the right contributes
\((-1)^{k(k-1)/2}\).  The anti-automorphism in (10.1) reverses the \(k\)
negative generators and contributes the same sign in
\(\Theta(\psi_I)\).  The signs therefore agree, yielding (10.7).
\end{proof}

\begin{theorem}[Reflected Grassmann Gram theorem]
\label{thm:v19v7-gram}
If \(C\geq0\), then \(W_C\in{\cal C}_\Theta\).  Hence for every even
\(F\in{\mathfrak A}_+\),
\[
 {\cal L}\!\left(\Theta(F)F e^{X_C}\right)\geq0.               \tag{10.8}
\]
More generally, if \(B\in{\mathfrak A}_+\) is even, then
\[
 \Theta(B)B e^{X_C}\in{\cal C}_\Theta                         \tag{10.9}
\]
and may replace \(e^{X_C}\) in (10.8).
\end{theorem}

\begin{proof}
For each \(k\), the matrix of minors in (10.7) is the matrix of the exterior
power \(\bigwedge^kC\) in the basis \(\{\psi_I:|I|=k\}\).
The spectral theorem gives \(C=S^*S\).  Functoriality of exterior powers
then gives
\[
 \bigwedge\nolimits^kC
 =(\bigwedge\nolimits^kS)^*(\bigwedge\nolimits^kS)\geq0.       \tag{10.10}
\]
Choose any Gram factorization
\(\bigwedge^kC=R_k^*R_k\), and define
\[
 G_{k,\alpha}:=\sum_{|J|=k}(R_k)_{\alpha J}\psi_J.
\]
Equation (10.7) becomes
\[
 e^{X_C}=\sum_{k,\alpha}\Theta(G_{k,\alpha})G_{k,\alpha},      \tag{10.11}
\]
which proves the first claim.  Proposition \ref{prop:v19v7-cone} proves
(10.8).  Since \(B\) is even, multiplication of each summand in (10.11)
by \(\Theta(B)B\) yields
\(\Theta(BG_{k,\alpha})(BG_{k,\alpha})\); this proves (10.9).
\end{proof}

\begin{corollary}[Vectorlike Wilson temporal hinge]
\label{cor:v19v7-wilson}
Let \(\gamma_0=\gamma_0^*\) and \(\gamma_0^2=1\).  Then
\[
 P_\pm:=\frac{1\pm\gamma_0}{2}\geq0.                           \tag{10.12}
\]
After the unitary temporal gauge transport is absorbed into the reflected
fibre identification, any Wilson-type cross block whose matrix is
\(C=\kappa P_+\), \(C=\kappa P_-\), or their block direct sum, with
\(\kappa\geq0\), satisfies Theorem \ref{thm:v19v7-gram}.  Reflected
one-side spatial, mass, and interaction polynomials may be included in the
factor \(\Theta(B)B\).
\end{corollary}

\begin{proof}
The relations in (10.12) show that \(P_\pm\) are orthogonal projections,
hence positive.  Positive scalar multiplication and block direct sums
preserve positivity.  Apply Theorem \ref{thm:v19v7-gram}; its last
assertion treats the one-side factors.
\end{proof}

This corollary is a finite algebraic hinge statement.  Combined with a
positive bosonic boundary kernel, integration over the bosonic gauge
variables preserves the nonnegative reflected expectation whenever the
transport identification and the one-side action are reflection matched.

\subsection{What vectorlike doubling proves, and what it changes}

\begin{lemma}[Finite Grassmann determinant identity]
\label{lem:v19v7-det}
For an \(N\times N\) complex matrix \(D\) and independent Grassmann pairs
\((\overline\eta_i,\eta_i)\), a consistent ordered Berezin convention gives
\[
 \int\prod_{i=1}^N d\overline\eta_i\,d\eta_i\,
 \exp\!\left(-\sum_{i,j}\overline\eta_iD_{ij}\eta_j\right)
 =\det D.                                                       \tag{10.13}
\]
\end{lemma}

\begin{proof}
Only the degree-\(N\) term of the exponential survives.  Nilpotence removes
all terms with a repeated row or column.  Ordering the surviving generators
contributes the sign of the corresponding permutation, so the remaining
sum is the Leibniz formula for \(\det D\).  The overall integration order
is chosen to give the displayed sign.
\end{proof}

\begin{proposition}[Conjugate vectorlike pair]
\label{prop:v19v7-vectorlike}
Two independent flavours with Dirac matrices \(D\) and \(D^*\) have
fermionic normalization
\[
 Z_{\rm pair}=\det D\,\det D^*=|\det D|^2\geq0.                \tag{10.14}
\]
\end{proposition}

\begin{proof}
Apply Lemma \ref{lem:v19v7-det} to each independent flavour.  The determinant
of the adjoint is the complex conjugate of the determinant.
\end{proof}

\begin{firewall}[A determinant square is not the chiral Standard Model]
\label{fire:v19v7-chiral}
Adding a conjugate flavour changes a chiral theory into a vectorlike one.
For a single Weyl sector, the regulated determinant is naturally a section
of a determinant line over gauge-field space rather than a globally fixed
nonnegative function.  Infinitesimal anomaly cancellation controls the
local variation or curvature of that line, but does not by itself choose a
global reflection-real trivialization or rule out global holonomy.  Thus
(10.14) cannot be substituted for the Standard Model fermion measure.
\end{firewall}

The compact V1 carry-forward already records the perturbative anomaly
cancellations and their global-anomaly firewalls.  Recomputing that charge
ledger here would not close the new gate: the missing datum is the global,
reflection-compatible phase of the regulated chiral determinant.

\subsection{Executable certificate}

The verifier
\[
 \texttt{scripts/verify\_sm\_v19\_v7\_reflected\_grassmann\_gram.py}
                                                                  \tag{10.15}
\]
uses exact rational Grassmann arithmetic.  For a positive matrix
\(C=B^{\mathsf T}B\), it expands \(e^{X_C}\), checks every coefficient
against the signed minor formula, and verifies
\(\bigwedge^kC=(\bigwedge^kB)^{\mathsf T}\bigwedge^kB\) for
\(0\leq k\leq3\).  It also checks the Wilson projectors and an exact
Gaussian-integer example
\[
 (6+2i)(6-2i)=40.                                               \tag{10.16}
\]
It emits
\[
 \texttt{artifacts/sm\_v19\_v7\_reflected\_grassmann\_gram.json}.
                                                                  \tag{10.17}
\]
The hashes are
\[
\begin{array}{c|l}
\hbox{object}&\hbox{SHA--256}\\ \hline
\hbox{verifier}&
\texttt{2F86EA6AB663C40698D6B36904372BAC4073FD92C33432F30210DD4C7CC5C6F7}\\
\hbox{canonical payload}&
\texttt{1AB61B00FF2B124BF4E073177323DEAEF29E3AFC366535AF51866649F8D8AC15}.
\end{array}                                                     \tag{10.18}
\]

\begin{decision}[V8 acquisition]
\label{dec:v19v7-next}
V8 will go upstream from the chiral phase bottleneck.  It will formulate a
finite determinant-line patching theorem on a gauge-field cover: local
phases glue to a global reflection-real determinant exactly when their
transition cocycle has trivial cycle holonomy and satisfies the reflection
reality constraint.  This is the precise global datum not supplied by the
retained perturbative anomaly ledger.
\end{decision}

\begin{assessment}[V7 acquisition]
\label{ass:v19v7}
The fermionic cross-plane problem now has a rigorous finite sufficient
condition.  Positivity of the temporal one-particle cross matrix propagates
to every particle number through its exterior powers, placing the full
Grassmann exponential in a reflected-square cone.  Vectorlike Wilson
hinges pass this test.  The chiral Standard Model remains blocked at the
global determinant-line phase, which V8 will isolate rather than hide by
doubling the fermion content.
\end{assessment}

\begin{record}[V7 append record]
\label{rec:v19v7}
V7 is appended to the validated V6 whose installed SHA--256 was
\texttt{DFF7DE7C78DC95CA5D52A13C492ECF5382F38D1DBCAE953A539DE9C7E3125B50}.
After installation only V7 is the active numeric SM note; frozen archives
and unrelated notes are unchanged.
\end{record}




\section{V8: determinant-line descent and the global holonomy obstruction}
\label{sec:v19v8}

The V7 reflected-Gram theorem proves fermionic positivity when the temporal
cross matrix is positive, and it explains why a conjugate vectorlike pair
has an absolute determinant square.  The Standard Model cannot use that
doubling.  The next issue is therefore global: local chiral determinants
must patch to one gauge-covariant phase before reflection positivity can
even be stated as a scalar inequality.  This note gives an exact descent
criterion and separates local anomaly cancellation from global
triviality.

\subsection{Local determinants are frames of a line}

For a finite-dimensional chiral Fredholm family \(D_A:E_A^+\to E_A^-\),
write
\[
 {\rm Det}(D)_A
 :=\det\ker D_A\otimes\big(\det\operatorname{coker}D_A\big)^*.
                                                               \tag{11.1}
\]
The same line records the regulated phase when local bases or spectral
cuts vary, including on the invertible locus.  Let
\(\{U_a\}_{a\in{\cal I}}\) be a finite good cover of the regulated
gauge-field quotient and choose nonvanishing local determinant frames
\(s_a\).  Use the transition convention
\[
 s_a=g_{ab}s_b,\qquad g_{ab}:U_a\cap U_b\longrightarrow U(1).
                                                               \tag{11.2}
\]
The functions satisfy \(g_{aa}=1\), \(g_{ba}=g_{ab}^{-1}\), and
\(g_{ab}g_{bc}=g_{ac}\) on triple overlaps.

\begin{theorem}[Determinant-line descent criterion]
\label{thm:v19v8-cech}
There is a global nowhere-zero unit frame \(s\) of \({\rm Det}(D)\) if and
only if there are continuous phase functions
\(h_a:U_a\to U(1)\) such that
\[
 g_{ab}=h_a^{-1}h_b\quad\hbox{on }U_a\cap U_b.                 \tag{11.3}
\]
When (11.3) holds, the local expressions \(h_as_a\) agree and define
\(s\).  Every other unit global frame is \(hs\) for a global
\(h:X\to U(1)\).
\end{theorem}

\begin{proof}
If (11.3) holds, then on an overlap
\[
 h_as_a=h_ag_{ab}s_b=h_b s_b,
\]
so the local sections glue.  Conversely, if \(s\) is a global unit frame,
there is a unique \(h_a:U_a\to U(1)\) with \(s=h_as_a\).
Equality of the two expressions on an overlap and (11.2) give
\(h_ag_{ab}=h_b\), which is (11.3).  If \(s'\) is another unit frame, the
pointwise ratio \(s'/s\) is a global \(U(1)\)-valued function, and the
converse is immediate.
\end{proof}

Thus a list of locally consistent determinant formulae is insufficient.
Their transition cocycle must be a coboundary.  The following finite
criterion makes the residual obstruction executable.

\subsection{Spanning-tree reconstruction of a flat phase}

Let \(\Gamma=(V,E)\) be a finite connected graph.  Give every oriented edge
\(e=(u,v)\) a phase \(q_e\in U(1)\), with
\(q_{(v,u)}=q_{(u,v)}^{-1}\).  For an oriented walk
\(\gamma=e_1\cdots e_k\), let
\[
 {\rm Hol}_q(\gamma):=\prod_{\ell=1}^kq_{e_\ell}.              \tag{11.4}
\]

\begin{theorem}[Finite holonomy test]
\label{thm:v19v8-tree}
The following are equivalent:
\begin{enumerate}
 \item there are vertex phases \(h_v\in U(1)\) such that
       \(q_{(u,v)}=h_u^{-1}h_v\) on every edge;
 \item \({\rm Hol}_q(\gamma)=1\) for every closed walk;
 \item after choosing any spanning tree \(T\), the holonomy is one on every
       fundamental cycle obtained by adding one chord \(e\in E\setminus T\).
\end{enumerate}
When these conditions hold, fixing \(h_{v_0}=1\) determines all other
\(h_v\) uniquely by multiplying phases along the tree path from \(v_0\).
\end{theorem}

\begin{proof}
If \(q_{(u,v)}=h_u^{-1}h_v\), the factors telescope around every closed
walk, proving \(1\Rightarrow2\).  The implication \(2\Rightarrow3\) is
immediate.  For \(3\Rightarrow1\), root the tree at \(v_0\) and define
\(h_v\) as the ordered phase product on its unique tree path.  This gives
\(q_{(u,v)}=h_u^{-1}h_v\) on every tree edge.  For a chord \((u,v)\), the
tree path from \(v\) back to \(u\), together with the chord, is its
fundamental cycle.  Unit holonomy on that cycle gives the same identity on
the chord.  Hence it holds on all edges.  Uniqueness with \(h_{v_0}=1\)
follows recursively along the tree.
\end{proof}

\begin{corollary}[Cycle rank is the finite audit size]
\label{cor:v19v8-cycle-rank}
For a connected graph, only
\[
 |E|-|V|+1                                                       \tag{11.5}
\]
fundamental holonomies need be checked.
\end{corollary}

\begin{proof}
A spanning tree has \(|V|-1\) edges, so the number of chords is (11.5).
Apply Theorem \ref{thm:v19v8-tree}.
\end{proof}

For transition functions on a good cover, the same construction is
performed in the multiplicative group of continuous \(U(1)\)-valued
functions on the relevant overlaps.  Theorem \ref{thm:v19v8-tree} is then
the one-skeleton part of Theorem \ref{thm:v19v8-cech}; the cocycle
identities on higher intersections enforce compatibility of the cycle
tests.

\subsection{Flat does not imply globally trivial}

\begin{proposition}[Residual global anomaly]
\label{prop:v19v8-flat}
A flat unitary determinant-line connection need not admit a global
parallel frame.  Its obstruction is the holonomy representation
\[
 \operatorname{Hol}:\pi_1(X)\longrightarrow U(1).              \tag{11.6}
\]
A parallel unit frame exists if and only if this representation is
trivial.
\end{proposition}

\begin{proof}
Parallel transport of a chosen unit vector from a base point defines a
candidate frame.  It is path independent exactly when transport around
every based loop is the identity, which is triviality of (11.6).  If a
parallel global frame exists, transporting that frame around any loop
returns the same vector, so the holonomy is trivial.
\end{proof}

\begin{example}[Flat circle with nontrivial phase]
\label{ex:v19v8-circle}
On \(S^1\), the connection \(d+i\alpha\,dt/(2\pi)\) is flat, but its
holonomy is \(e^{-i\alpha}\).  Unless \(\alpha\in2\pi\mathbb Z\), it has no
global parallel unit frame.
\end{example}

The perturbative Standard Model anomaly equalities retained in the V1
carry-forward remove the corresponding local anomaly polynomial under the
stated representation assumptions.  Even granting a regulator for which
that cancellation makes the determinant-line connection flat, Proposition
\ref{prop:v19v8-flat} shows that a global holonomy calculation still
remains.  The finite reflected-Gram theorem of V7 can be promoted to a
gauge-field quotient only after this phase descends.

\subsection{Executable certificate}

The verifier
\[
 \texttt{scripts/verify\_sm\_v19\_v8\_determinant\_line\_holonomy.py}
                                                                  \tag{11.7}
\]
uses \(\mathbb Z/12\mathbb Z\) as an exact phase group.  It reconstructs a
vertex phase on a five-vertex spanning tree, checks zero residue on every
chord for a coboundary, inserts a nonzero chord residue to detect the
obstruction, and gives a flat one-dimensional circle with nontrivial
holonomy.  It emits
\[
 \texttt{artifacts/sm\_v19\_v8\_determinant\_line\_holonomy.json}.
                                                                  \tag{11.8}
\]
The hashes are
\[
\begin{array}{c|l}
\hbox{object}&\hbox{SHA--256}\\ \hline
\hbox{verifier}&
\texttt{8B976834E8AD60A00CC76C6EB9A7305426742CBE1ABE54AD9BE1B444595E7F08}\\
\hbox{canonical payload}&
\texttt{774F047DF3F964DB69DEAB3683E5F383C65AF24E7F5D6BDA5A33B80B7D22976A}.
\end{array}                                                     \tag{11.9}
\]

\begin{firewall}[Descent is not positivity]
\label{fire:v19v8}
Triviality of the determinant line produces a global complex scalar
determinant.  It does not make that scalar real or nonnegative, does not
make the trivialization compatible with time reflection, and does not
establish the cutoff-uniform bounds needed for a continuum fermion
measure.
\end{firewall}

\begin{decision}[V9 acquisition]
\label{dec:v19v8-next}
V9 will add the antiunitary reflection to the descended determinant line.
It will prove the exact criterion for a reflection-real global frame,
including the residual sign holonomy on the reflection-fixed locus.  This
is the next prerequisite for combining a chiral determinant phase with
the V7 reflected-Grassmann cone.
\end{decision}

\begin{assessment}[V8 acquisition]
\label{ass:v19v8}
The chiral bottleneck has been reduced to two distinct global tests.
First, the determinant transition cocycle must have trivial \(U(1)\)
holonomy so that a scalar phase exists at all.  Second, that phase must
admit a reflection-real choice.  Local anomaly cancellation addresses
neither test by itself.
\end{assessment}

\begin{record}[V8 append record]
\label{rec:v19v8}
V8 is appended to the validated V7 whose installed SHA--256 was
\texttt{BD26EE7F33D472A757B6C3B7A6B867A64844414FDC0B0E4646D56701D7B2956A}.
After installation only V8 is the active numeric SM note; frozen archives
and unrelated notes are unchanged.
\end{record}



\section{V9: reflection-real determinant frames and fixed-locus sign holonomy}
\label{sec:v19v9}

V8 separated the chiral determinant-line problem from local anomaly
cancellation: a flat line may still have nontrivial \(U(1)\) holonomy.
Even after that holonomy vanishes and a complex unit frame is chosen, time
reflection imposes another descent equation.  This note derives that
equation, identifies its fixed-locus obstruction, and proves that a
reflection-real determinant may still have either sign.

\subsection{The Real structure on the determinant line}

Let \(X\) be the regulated gauge--metric configuration quotient, let
\(\theta:X\to X\) be an involution, and let \(L\to X\) be a Hermitian line
bundle.  A reflection lift is a fibrewise antiunitary map
\[
 \tau_x:L_x\longrightarrow L_{\theta x},\qquad
 \tau_{\theta x}\tau_x=1.                                          \tag{12.1}
\]
Assume first that the complex line has passed V8's descent test, and choose
a global unit frame \(s\).  There is a unique phase
\(r_s:X\to U(1)\) such that
\[
 \tau_xs(x)=r_s(x)s(\theta x).                                 \tag{12.2}
\]

\begin{lemma}[Reflection phase and change of frame]
\label{lem:v19v9-phase}
The phase in (12.2) satisfies
\[
 r_s(\theta x)=r_s(x).                                         \tag{12.3}
\]
If \(s'=hs\), with \(h:X\to U(1)\), then
\[
 r_{s'}(x)=\overline{h(x)}\,r_s(x)\,h(\theta x)^{-1}.           \tag{12.4}
\]
Consequently, \(s'\) is reflection real,
\(\tau_xs'(x)=s'(\theta x)\), exactly when
\[
 h(\theta x)=\overline{h(x)}\,r_s(x).                          \tag{12.5}
\]
\end{lemma}

\begin{proof}
Apply \(\tau_{\theta x}\) to (12.2).  Antilinearity and
\(\tau_{\theta x}\tau_x=1\) give
\[
 s(x)=\overline{r_s(x)}\,r_s(\theta x)s(x),
\]
which is (12.3).  Applying \(\tau\) to
\(s'(x)=h(x)s(x)\), and then expressing \(s(\theta x)\) in the new frame,
gives (12.4).  Setting the right side of (12.4) equal to one gives (12.5).
\end{proof}

\begin{theorem}[Exact reflection-real frame criterion]
\label{thm:v19v9-real-frame}
Under the complex trivialization assumption, \(L\) admits a global
reflection-real unit frame if and only if the functional equation (12.5)
has a continuous solution \(h:X\to U(1)\).

On the fixed locus \(X^\theta\), a necessary condition is the continuous
square-root equation
\[
 h(x)^2=r_s(x),\qquad x\in X^\theta.                            \tag{12.6}
\]
For a CW fixed locus, (12.6) is solvable if and only if the winding of
\(r_s\) around every loop in \(X^\theta\) is even.  Equivalently, the real
fixed subbundle
\[
 L^\tau:=\{v\in L|_{X^\theta}:\tau v=v\}                       \tag{12.7}
\]
has \(w_1(L^\tau)=0\).

If \(X=Y\cup\theta Y\), \(Y\cap\theta Y=X^\theta\), and a solution of
(12.6) extends from \(X^\theta\) to \(Y\), then this fixed-locus condition
is also sufficient.
\end{theorem}

\begin{proof}
The first assertion is Lemma \ref{lem:v19v9-phase}.  At a fixed point,
(12.5) becomes \(h=\overline h\,r_s\).  Since \(|h|=1\), this is
\(h^2=r_s\).

The squaring map \(U(1)\to U(1)\) is a twofold covering.  The covering-space
lifting criterion says that \(r_s|_{X^\theta}\) has a continuous lift
through this map precisely when its induced winding on every loop lies in
\(2\mathbb Z\).  Local square roots differ by signs; those signs are the
transition functions of the real line (12.7).  Their Čech class is
\(w_1(L^\tau)\), so a global root exists exactly when that class vanishes.

For the last assertion, extend a fixed-locus root to a phase \(h\) on
\(Y\).  Define it on \(\theta Y\) by (12.5).  Equation (12.3) makes the
definition involutively consistent, while (12.6) makes the two definitions
agree on \(Y\cap\theta Y\).  The resulting global \(h\) solves (12.5).
\end{proof}

\begin{remark}[Free reflection orbits]
\label{rem:v19v9-free}
On a two-point orbit \(\{x,\theta x\}\), choosing \(h(x)\) arbitrarily
forces \(h(\theta x)\) by (12.5), and (12.3) makes the return equation
automatic.  The new obstruction is therefore concentrated in topology and
in the reflection-fixed set, rather than at an isolated paired
configuration.
\end{remark}

\subsection{Reality still does not give a positive density}

Let \(\Delta\) be a determinant section of \(L\), and suppose \(s\) is a
reflection-real frame.  Write
\[
 \Delta(x)=z(x)s(x).                                           \tag{12.8}
\]
If the determinant section is equivariant,
\(\tau_x\Delta(x)=\Delta(\theta x)\), then
\[
 z(\theta x)=\overline{z(x)}.                                  \tag{12.9}
\]

\begin{proposition}[Fixed-locus sign sectors]
\label{prop:v19v9-sign}
Equation (12.9) makes \(z\) real on \(X^\theta\), but does not make it
nonnegative.  On each connected component of
\[
 X^\theta\setminus\{x:z(x)=0\},                                \tag{12.10}
\]
the sign of \(z\) is constant.  A sign change along a continuous path in
the fixed locus must pass through a determinant zero.
\end{proposition}

\begin{proof}
For \(x=\theta x\), (12.9) gives \(z(x)=\overline{z(x)}\).
A continuous nonzero real function has locally constant sign; connectedness
makes that sign constant on every component of (12.10).  The intermediate
value theorem forces a zero on any path joining opposite signs.
\end{proof}

\begin{corollary}[Pfaffian orientation gate]
\label{cor:v19v9-pfaffian}
Whenever a fermion regulator produces a real Pfaffian line on
\(X^\theta\), a globally consistent Pfaffian sign exists only if its first
Stiefel--Whitney class vanishes.  Choosing such an orientation still does
not prove that the oriented Pfaffian is positive on every zero-free
component.
\end{corollary}

\begin{proof}
An orientation of a real line bundle is a nowhere-zero continuous section,
which exists exactly when its \(w_1\) vanishes.  Proposition
\ref{prop:v19v9-sign} supplies the remaining componentwise sign statement.
\end{proof}

\begin{criterion}[Chiral reflection interface]
\label{crit:v19v9-interface}
Before a regulated chiral Standard Model determinant can enter the V7
reflection-positive cone, the following distinct facts must be proved:
\begin{enumerate}
 \item the complex determinant-line cocycle is trivial, as in V8;
 \item the reflection lift has a global Real trivialization, as in
       Theorem \ref{thm:v19v9-real-frame};
 \item determinant zeros and the signs of all fixed-locus components are
       controlled;
 \item the resulting full cross density has an actual reflected-Grassmann
       Gram decomposition, not merely the conjugation law (12.9).
\end{enumerate}
\end{criterion}

\begin{proof}
The first item is necessary to turn local determinant phases into a global
scalar.  The second is necessary for that scalar to obey (12.9).  The third
is necessary because Proposition \ref{prop:v19v9-sign} permits negative
components.  The fourth is necessary because V7 proves positivity from a
sum of reflected squares; complex conjugation symmetry alone supplies no
such sum.  Together these are necessary interface conditions.  No
sufficiency beyond the stated finite Gram criterion is asserted.
\end{proof}

\subsection{Executable certificate}

The verifier
\[
 \texttt{scripts/verify\_sm\_v19\_v9\_reflection\_real\_phase.py}
                                                                  \tag{12.11}
\]
uses an exact \(\mathbb Z/12\mathbb Z\) phase model to solve (12.5) on a
free reflection orbit and the fixed-point equation \(2h=r\).  It records
the even-winding square-root condition, detects an odd-winding and
\(\mathbb Z/2\) sign-holonomy obstruction on a fixed loop, and checks a
reflection-real scalar that is negative at a fixed configuration.  It
emits
\[
 \texttt{artifacts/sm\_v19\_v9\_reflection\_real\_phase.json}. \tag{12.12}
\]
The hashes are
\[
\begin{array}{c|l}
\hbox{object}&\hbox{SHA--256}\\ \hline
\hbox{verifier}&
\texttt{3353EBD58964AC4ED24B16908021BD78C26E62F67CCFB1AC760B3276A009ED6C}\\
\hbox{canonical payload}&
\texttt{2DF125E84F15BFC76CBD7C100D718A48F64387D9C2379D76593542CB825F5490}.
\end{array}                                                     \tag{12.13}
\]

\begin{firewall}[A Real frame is not a reflected square]
\label{fire:v19v9}
Theorems \ref{thm:v19v8-cech} and \ref{thm:v19v9-real-frame} can at most
produce a globally defined determinant with the conjugation law (12.9).
They do not supply the reflected-Grassmann Gram factorization required by
Theorem \ref{thm:v19v7-gram}, and they provide no cutoff-uniform control
near determinant zeros.
\end{firewall}

\begin{decision}[Mandatory V10 audit, then the sign gate]
\label{dec:v19v9-next}
V10 is the mandatory companion audit.  If it supplies no newer fermionic
or determinant-line input, V11 will go upstream to a spectral-flow
criterion for the fixed-locus sign sectors.  The target is an exact
finite-dimensional theorem relating determinant or Pfaffian sign changes
to zero crossings, with multiplicity and endpoint hypotheses stated
explicitly.
\end{decision}

\begin{assessment}[V9 acquisition]
\label{ass:v19v9}
The determinant gate now has an exact two-stage topology.  Vanishing
complex \(U(1)\) holonomy is followed by vanishing Real
\(\mathbb Z/2\) holonomy.  Passing both stages makes the determinant
reflection real, but positivity still requires spectral sign control and
the V7 reflected-square structure.
\end{assessment}

\begin{record}[V9 append record]
\label{rec:v19v9}
V9 is appended to the validated V8 whose installed SHA--256 was
\texttt{B89A26733C241EB8217413444FD4AD031C87C81E30C78CF959121BC4912E7091}.
After installation only V9 is the active numeric SM note; frozen archives
and unrelated notes are unchanged.
\end{record}




\section{V10: mandatory companion audit at the chiral phase gate}
\label{sec:v19v10}

The five-version rule now requires a read-only audit of the active
Yang--Mills and Navier--Stokes notes before the determinant-sign programme
continues.

\begin{audit}[Active companion inventory]
\label{aud:v19v10-files}
The active Yang--Mills note is
\[
 \texttt{Renewal\_Yang\_Mills\_Mass\_Gap\_Research\_Notes\_V25.tex},
 \qquad
 \operatorname{SHA256}=
 \texttt{5D6759B911E076304A17D09955E6777C3ACD6C99D3B47CD42BBA75EB0EA96BE6}.
                                                                  \tag{13.1}
\]
The active Navier--Stokes note is
\[
 \texttt{Renewal\_NS\_Stability\_Research\_Notes\_V26.tex},
 \qquad
 \operatorname{SHA256}=
 \texttt{82C887F8C2C69E4F28FAD7C3DC8DE5B9099CB5A4BF431EBA1FFF3E91935978AD}.
                                                                  \tag{13.2}
\]
Both hashes are unchanged from the V5 audit.
\end{audit}

\begin{proposition}[No new phase input]
\label{prop:v19v10-no-input}
Neither companion note supplies a determinant-line trivialization,
reflection-real phase, Pfaffian orientation, or reflected-Grassmann Gram
factorization.  The V11 acquisition should therefore continue from the
four-part chiral interface criterion of V9.
\end{proposition}

\begin{proof}
Yang--Mills V25 is itself a companion audit.  Its active bottleneck is a
finite signed Haar/Laurent coefficient computation; its retained SM input
is the exact-word-first link-block architecture from the previous SM
cycle.  It contains no fermion determinant phase.  Navier--Stokes V26
computes rank-one norms for first and second derivatives of the Oseen
kernel.  Those tensors act on velocity inputs and have no determinant-line
or reflection-lift datum.  Since the hashes in (13.1)--(13.2) are
unchanged, there is also no later companion appendix to import.
\end{proof}

\begin{decision}[V11 acquisition]
\label{dec:v19v10-next}
V11 will prove a finite spectral-flow sign theorem for a reflection-real
fermion determinant or Pfaffian path.  It will identify the parity of
transverse zero crossings as the obstruction to keeping a chosen positive
orientation, and will separate that mod-two obstruction from the complex
\(U(1)\) holonomy already treated in V8.
\end{decision}

\begin{firewall}[Audit boundary]
\label{fire:v19v10}
The unchanged companion hashes do not prove the V9 chiral interface
conditions.  They only show that continuing the determinant-sign analysis
does not duplicate a newer parallel result.
\end{firewall}

\begin{record}[V10 append record]
\label{rec:v19v10}
V10 is appended to the validated V9 whose installed SHA--256 was
\texttt{157030BDEEE42F98935E65B0E4D491175FF75B737513983722B3F0412E386B8A}.
After installation only V10 is the active numeric SM note; frozen archives
and unrelated notes are unchanged.
\end{record}




\section{V11: spectral-flow parity and the determinant-sign obstruction}
\label{sec:v19v11}

V9 showed that a reflection-real determinant is merely real on the
reflection-fixed locus; it can still be negative.  The mandatory V10 audit
found no companion theorem that controls this sign.  We now identify the
precise finite-dimensional mechanism: every regular real eigenvalue
crossing changes determinant orientation, and every regular skew block
crossing changes Pfaffian orientation.  Their parity computes the
first Stiefel--Whitney holonomy isolated in V9.

\subsection{Real determinant crossings}

Let \(A:[0,1]\to{\rm Sym}_m(\mathbb R)\) be \(C^1\), with invertible
endpoints.  A crossing \(t_0\) is \emph{regular} when the crossing form
\[
 \Gamma_{t_0}:=\left.P_{t_0}A'(t_0)P_{t_0}\right|_{\ker A(t_0)}
                                                               \tag{14.1}
\]
is nonsingular, where \(P_{t_0}\) is the orthogonal projection onto the
kernel.

\begin{theorem}[Determinant sign from crossing parity]
\label{thm:v19v11-det-sign}
Assume that \(A\) has finitely many crossings \(t_1,\ldots,t_N\), all
regular.  Put \(m_j=\dim\ker A(t_j)\).  Then
\[
 \frac{\operatorname{sgn}\det A(1)}
      {\operatorname{sgn}\det A(0)}
 =(-1)^{\sum_{j=1}^N m_j}.                                    \tag{14.2}
\]
In particular, the right side is the mod-two reduction of the ordinary
spectral flow.
\end{theorem}

\begin{proof}
Fix one crossing \(t_0\) and split
\(\mathbb R^m=K\oplus K^\perp\), where \(K=\ker A(t_0)\).
In this splitting the lower-right block \(D(t)\) is invertible near
\(t_0\).  The determinant Schur formula gives
\[
 \det A(t)=\det D(t)\,
 \det\!\left(B(t)-C(t)D(t)^{-1}C(t)^{\mathsf T}\right).         \tag{14.3}
\]
At \(t_0\), \(B(t_0)=C(t_0)=0\).  Hence the second determinant in
(14.3) equals
\[
 \det\!\left((t-t_0)\Gamma_{t_0}+o(|t-t_0|)\right)
 =(t-t_0)^{m_0}\det\Gamma_{t_0}+o(|t-t_0|^{m_0}).              \tag{14.4}
\]
The nonzero factor \(\det D(t)\) keeps its sign, so passage through \(t_0\)
multiplies the determinant sign by \((-1)^{m_0}\).  Away from crossings
the determinant is a continuous nonzero real function and keeps its sign.
Multiplying the local changes proves (14.2).  At a regular crossing, the
signed spectral-flow contribution is the signature of \(\Gamma_{t_j}\);
modulo two, its signature equals its dimension, proving the last claim.
\end{proof}

\subsection{Pfaffian crossings}

Let \(J^{\mathsf T}=-J\) be a real \(2n\times2n\) matrix.  Recall
\[
 \operatorname{pf}(J)^2=\det J,\qquad
 \operatorname{pf}(QJQ^{\mathsf T})
 =\det(Q)\operatorname{pf}(J).                                \tag{14.5}
\]

\begin{theorem}[Pfaffian crossing parity]
\label{thm:v19v11-pf-sign}
Let \(J(t)\) be a \(C^1\) path of real skew-symmetric \(2n\times2n\)
matrices with invertible endpoints and finitely many crossings.  Assume
that at \(t_j\), the skew form
\[
 \left.P_{t_j}J'(t_j)P_{t_j}\right|_{\ker J(t_j)}
\]
is nondegenerate.  Write \(\dim\ker J(t_j)=2k_j\).  Then
\[
 \frac{\operatorname{sgn}\operatorname{pf}J(1)}
      {\operatorname{sgn}\operatorname{pf}J(0)}
 =(-1)^{\sum_j k_j}.                                          \tag{14.6}
\]
\end{theorem}

\begin{proof}
Split into the kernel and its orthogonal complement at one crossing.
The complementary skew block is invertible.  The Pfaffian Schur-complement
identity reduces the vanishing factor to a \(2k_j\times2k_j\) skew matrix
of the form
\[
 (t-t_j)\Gamma_j+o(|t-t_j|),
\]
where \(\Gamma_j\) is nondegenerate and skew.  Homogeneity of the Pfaffian
of a \(2k_j\times2k_j\) matrix gives
\[
 \operatorname{pf}\!\left((t-t_j)\Gamma_j+o(|t-t_j|)\right)
 =(t-t_j)^{k_j}\operatorname{pf}(\Gamma_j)
   +o(|t-t_j|^{k_j}).                                         \tag{14.7}
\]
Thus the sign changes by \((-1)^{k_j}\).  The Pfaffian is continuous and
nonzero between crossings, so multiplying these changes proves (14.6).
\end{proof}

\begin{lemma}[Congruence orientation law]
\label{lem:v19v11-congruence}
The second identity in (14.5) holds for every real \(Q\).
Consequently, an endpoint identification with \(\det Q=-1\) reverses the
Pfaffian orientation.
\end{lemma}

\begin{proof}
Associate to \(J\) the two-form
\(\omega_J=\frac12\sum_{a,b}J_{ab}e_a\wedge e_b\).  By definition,
\[
 \frac{\omega_J^n}{n!}
 =\operatorname{pf}(J)e_1\wedge\cdots\wedge e_{2n}.
\]
The change of basis \(Q\) multiplies the top exterior form by
\(\det Q\), proving the congruence identity.
\end{proof}

\begin{example}[One forced Pfaffian crossing]
\label{ex:v19v11-one-crossing}
Let
\[
 J_0=\begin{pmatrix}0&1\\-1&0\end{pmatrix},\qquad
 Q=\begin{pmatrix}-1&0\\0&1\end{pmatrix}.
\]
Then \(QJ_0Q^{\mathsf T}=-J_0\), \(\det Q=-1\), and the path
\(J(t)=(1-2t)J_0\) has one regular skew block crossing.  Its Pfaffian
changes from \(1\) to \(-1\), exactly as (14.6) and Lemma
\ref{lem:v19v11-congruence} require.
\end{example}

\subsection{Zero parity computes the Real-line holonomy}

\begin{theorem}[\(w_1\) as mod-two zero count]
\label{thm:v19v11-w1}
Let \(R\to X^\theta\) be the real fixed determinant or Pfaffian line
of V9, and let \(\gamma:S^1\to X^\theta\) be a loop.  If a section of
\(R\) restricted to \(\gamma\) has only transverse zeros, then
\[
 \left\langle w_1(R),[\gamma]\right\rangle
 =\#\{\hbox{zeros on }\gamma\}\pmod2.                          \tag{14.8}
\]
For a regular real-symmetric determinant family, the count has the parity
in Theorem \ref{thm:v19v11-det-sign}; for a real skew family, it has the
block parity in Theorem \ref{thm:v19v11-pf-sign}.
\end{theorem}

\begin{proof}
Cover the circle by intervals on which \(R\) is oriented and write the
section as a real coefficient.  At each transverse zero that coefficient
changes sign.  After one circuit, the product of these sign changes is
\((-1)^N\).  The same product is the sign holonomy of the real line, whose
mod-two class is \(\langle w_1(R),[\gamma]\rangle\).  This proves (14.8).
The final statements follow from the two crossing theorems.
\end{proof}

There is a typographical distinction worth retaining: determinant zeros
of a real symmetric family count individual real modes, while a
skew-symmetric Pfaffian family has even-dimensional kernels and counts
two-dimensional skew blocks.  Squaring the Pfaffian erases precisely the
sign detected by (14.6).

\subsection{Executable certificate}

The verifier
\[
 \texttt{scripts/verify\_sm\_v19\_v11\_spectral\_flow\_sign.py} \tag{14.9}
\]
checks a three-crossing symmetric path with endpoint determinants
\(-6\) and \(6\), a simultaneous two-dimensional crossing with no sign
change, a one-block Pfaffian path related at its endpoints by an
orientation-reversing congruence, a two-block even crossing, and the
identity between one transverse zero and nontrivial \(\mathbb Z/2\)
holonomy.  It emits
\[
 \texttt{artifacts/sm\_v19\_v11\_spectral\_flow\_sign.json}.    \tag{14.10}
\]
The hashes are
\[
\begin{array}{c|l}
\hbox{object}&\hbox{SHA--256}\\ \hline
\hbox{verifier}&
\texttt{1AC6F2BF37F20819EC905A2F3AEAD8167C05A86BE39FA9FBF57858D848C2C92D}\\
\hbox{canonical payload}&
\texttt{F181FE797CCD3381C38C782032A734C2E4F6BDD1C9868AC0AA48AB64AB155C40}.
\end{array}                                                     \tag{14.11}
\]

\begin{firewall}[Parity detects a sign obstruction; it does not remove it]
\label{fire:v19v11}
Theorems \ref{thm:v19v11-det-sign}--\ref{thm:v19v11-w1} give a finite
test for sign holonomy once a regulator family and its reflection-fixed
paths are specified.  They do not prove that the Standard Model regulator
has a spectral gap, even crossing parity, or a positive Pfaffian chamber.
They also do not construct the global complex phase required in V8.
\end{firewall}

\begin{decision}[V12 acquisition]
\label{dec:v19v11-next}
V12 will prove a gap-chamber criterion that avoids path-by-path crossing
counts: on a connected reflection-fixed component, a uniform singular-
value gap keeps the real determinant or Pfaffian sign constant, so one
reference configuration fixes the component.  It will quantify the
allowed operator perturbation by the smallest singular value and expose
why gauge-field topology prevents one reference point from controlling
all components.
\end{decision}

\begin{assessment}[V11 acquisition]
\label{ass:v19v11}
The residual Real-line obstruction of V9 now has an exact spectral form.
Odd determinant-mode flow or odd skew-block flow reverses orientation;
equivalently, transverse zero parity evaluates \(w_1\).  A chiral regulator
must therefore prove even parity on every relevant loop or supply an
equivalent global orientation theorem.
\end{assessment}

\begin{record}[V11 append record]
\label{rec:v19v11}
V11 is appended to the validated V10 whose installed SHA--256 was
\texttt{261EFC9DE5E5772BCC510244561E0D9EBBDF278426971FB280493DC62AFA10AF8}.
After installation only V11 is the active numeric SM note; frozen archives
and unrelated notes are unchanged.
\end{record}




\section{V12: quantitative gap certificates for determinant-sign chambers}
\label{sec:v19v12}

V11 converts fixed-locus sign holonomy into the parity of spectral
crossings.  A complementary route is available on a region where the
regulated fermion operator has a uniform singular-value gap: no crossing
can occur there, so one reference sign determines the entire connected
region.  This note proves a finite-net certificate for that situation and
states explicitly why zero-mode sectors remain outside it.

\subsection{A gap fixes the sign on a connected region}

For a real matrix \(A\), let
\[
 \sigma_{\min}(A):=\inf_{\|v\|=1}\|Av\|.                       \tag{15.1}
\]

\begin{theorem}[Connected sign chamber]
\label{thm:v19v12-chamber}
Let \(X\) be connected and let \(A:X\to{\rm Sym}_m(\mathbb R)\) be
continuous.  If
\[
 \sigma_{\min}(A(x))\geq\delta>0\quad\hbox{for all }x\in X,     \tag{15.2}
\]
then \(\operatorname{sgn}\det A(x)\) is constant on \(X\), and
\[
 |\det A(x)|\geq\delta^m.                                      \tag{15.3}
\]
Thus one reference value determines the sign everywhere.

If \(J:X\to{\rm Skew}_{2n}(\mathbb R)\) is continuous, satisfies (15.2),
and is expressed in a fixed oriented Pfaffian frame, then
\(\operatorname{sgn}\operatorname{pf}J(x)\) is constant and
\[
 |\operatorname{pf}J(x)|\geq\delta^n.                          \tag{15.4}
\]
\end{theorem}

\begin{proof}
Condition (15.2) makes every matrix invertible.  The determinant is
continuous and never zero, so its sign is locally constant and hence
constant on connected \(X\).  The absolute determinant is the product of
the \(m\) singular values, which proves (15.3).

A real skew matrix has singular values
\(\lambda_1,\lambda_1,\ldots,\lambda_n,\lambda_n\), with each
\(\lambda_j\geq\delta\).  Since
\[
 |\operatorname{pf}J|^2=\det J=\prod_{j=1}^n\lambda_j^2,
\]
equation (15.4) follows.  Continuity in the fixed oriented frame and
nonvanishing give the constant Pfaffian sign.
\end{proof}

\begin{corollary}[Perturbative preservation of a chamber]
\label{cor:v19v12-perturb}
If \(A_0\) is invertible and
\[
 \|A-A_0\|_{\rm op}<\sigma_{\min}(A_0),                        \tag{15.5}
\]
then the segment \(A_t=A_0+t(A-A_0)\) stays invertible, so \(A\) and
\(A_0\) have the same real determinant sign.  The analogous statement
holds for a skew family in a fixed Pfaffian orientation.
\end{corollary}

\begin{proof}
For every unit vector \(v\),
\[
 \|A_tv\|\geq\|A_0v\|-t\|(A-A_0)v\|
 \geq\sigma_{\min}(A_0)-\|A-A_0\|_{\rm op}>0.
\]
Apply Theorem \ref{thm:v19v12-chamber} to the segment.
\end{proof}

\subsection{A finite spectral net proves a global gap}

\begin{theorem}[Gap certificate from an \(\varepsilon\)-net]
\label{thm:v19v12-net}
Let \((X,d)\) be a compact metric space, let
\(A:X\to\mathbb R^{m\times m}\) satisfy
\[
 \|A(x)-A(y)\|_{\rm op}\leq Ld(x,y),                           \tag{15.6}
\]
and let \(\{x_1,\ldots,x_N\}\) be an \(\varepsilon\)-net.  If
\[
 \delta_{\rm net}:=\min_i\sigma_{\min}(A(x_i))>L\varepsilon,  \tag{15.7}
\]
then
\[
 \sigma_{\min}(A(x))
 \geq\delta_{\rm net}-L\varepsilon>0
 \quad\hbox{for every }x\in X.                                \tag{15.8}
\]
When \(X\) is connected and \(A\) is real symmetric or real skew in a
fixed orientation, one reference sign and (15.8) certify the whole sign
chamber by Theorem \ref{thm:v19v12-chamber}.
\end{theorem}

\begin{proof}
Choose a net point \(x_i\) with \(d(x,x_i)\leq\varepsilon\).  For any unit
vector \(v\),
\[
 \|A(x)v\|
 \geq\|A(x_i)v\|-\|(A(x)-A(x_i))v\|
 \geq\delta_{\rm net}-L\varepsilon.
\]
Taking the infimum over \(v\) proves (15.8).  The sign conclusion is
Theorem \ref{thm:v19v12-chamber}.
\end{proof}

\begin{remark}[What the certificate must contain]
\label{rem:v19v12-certificate}
A numerical grid without a proved operator-Lipschitz constant and covering
radius does not establish (15.8).  A valid certificate must record the
metric on configuration space, the net radius, a rigorous \(L\), interval
or exact lower bounds for every sampled smallest singular value, and
the sign at one certified reference configuration.
\end{remark}

\subsection{The zero wall is the only route between signs}

\begin{proposition}[Sign-changing paths close the gap]
\label{prop:v19v12-zero-wall}
Let \(A:[0,1]\to{\rm Sym}_m(\mathbb R)\) be continuous with invertible
endpoints of opposite determinant sign.  Then some \(t_0\) satisfies
\[
 \det A(t_0)=0,\qquad \sigma_{\min}(A(t_0))=0.                 \tag{15.9}
\]
The same conclusion holds for opposite Pfaffian signs in a fixed oriented
skew family.
\end{proposition}

\begin{proof}
The determinant or oriented Pfaffian is a continuous real function with
opposite endpoint signs.  The intermediate value theorem gives a zero.
A square matrix has zero determinant exactly when its smallest singular
value is zero; the Pfaffian identity in (14.5) gives the skew case.
\end{proof}

This proposition prevents a global-gap shortcut whenever the regulated family has
forced zero modes or distinct sign chambers.  In that case V11's
mod-two crossing calculation, rather than Theorem \ref{thm:v19v12-net},
is the appropriate audit.

\subsection{Executable certificate}

The verifier
\[
 \texttt{scripts/verify\_sm\_v19\_v12\_gap\_sign\_chamber.py}  \tag{15.10}
\]
uses exact rational arithmetic on the parameter square
\([-1,1]^2\).  A radius-\(1/2\) net and operator-Lipschitz constant one
give certified gap \(1/2\) for the symmetric family
\(\operatorname{diag}(2+x,3-x,4+y)\) and the skew family
\(\operatorname{diag}((2+x)J,(3+y)J)\).  The resulting determinant and Pfaffian
lower bounds are \(1/8\) and \(1/4\).  It also checks a sign-changing path whose gap
closes exactly at the zero wall.  The verifier emits
\[
 \texttt{artifacts/sm\_v19\_v12\_gap\_sign\_chamber.json}.     \tag{15.11}
\]
The hashes are
\[
\begin{array}{c|l}
\hbox{object}&\hbox{SHA--256}\\ \hline
\hbox{verifier}&
\texttt{28E981F2BAD397C653CB5549D54A2E404B4EEEBA6EB11FC3818C5A3D284930EF}\\
\hbox{canonical payload}&
\texttt{C95FF0131C52550A2CFA28B4F650BD3E5B8ECB3AB1587B7996711929B5B57C08}.
\end{array}                                                     \tag{15.12}
\]

\begin{firewall}[No Standard Model gap has been established]
\label{fire:v19v12}
Theorems \ref{thm:v19v12-chamber} and \ref{thm:v19v12-net} are certificate
criteria.  No uniform singular-value lower bound has been proved for a
chiral Standard Model regulator over its full gauge--metric configuration
space.  Topological zero modes may force (15.9), and a positive mass
regularization does not settle the sign after the mass is removed.
\end{firewall}

\begin{decision}[V13 acquisition]
\label{dec:v19v12-next}
V13 will treat the combined fermion multiplet rather than one species at a
time.  It will prove that the Real determinant line of a tensor product has
\[
 w_1(R_{\rm total})=\sum_f w_1(R_f)\pmod2,
\]
and that the combined Pfaffian sign holonomy is the sum of the species'
mod-two crossing parities.  This gives the exact global-sign cancellation
test without replacing a chiral multiplet by a vectorlike double.
\end{decision}

\begin{assessment}[V12 acquisition]
\label{ass:v19v12}
On any connected regulator region with a certified spectral gap, one
reference determinant or Pfaffian sign now controls the entire region.
The finite-net theorem makes this statement executable.  Regions touching
the zero wall remain governed by V11's spectral-flow parity.
\end{assessment}

\begin{record}[V12 append record]
\label{rec:v19v12}
V12 is appended to the validated V11 whose installed SHA--256 was

\texttt{0A4CA417A0E7757E72B55414D59D301C6FE8AD09A25DD3767ABE20EF7F0F273C}.
After installation only V12 is the active numeric SM note; frozen archives
and unrelated notes are unchanged.
\end{record}





\section{V13: multiplet cancellation of mod-two determinant holonomy}
\label{sec:v19v13}

V11 identifies each real fermion sign obstruction with a mod-two
spectral-flow character, while V12 certifies signs on zero-free chambers.
The Standard Model contains a multiplet of chiral species, so the relevant
object is the tensor product of their determinant or Pfaffian lines.
This note proves the exact addition law for the combined obstruction.
Individual lines may be nonorientable even when their product is
orientable.

\subsection{Tensor products add first Stiefel--Whitney classes}

Let \(R_f\to X^\theta\) be the real fixed determinant line of species
\(f\), and let \(m_f\) be its multiplicity.  Define
\[
 R_{\rm tot}:=\bigotimes_f R_f^{\otimes m_f}.                  \tag{16.1}
\]

\begin{theorem}[Multiplet orientation law]
\label{thm:v19v13-w1-sum}
The combined real line satisfies
\[
 w_1(R_{\rm tot})
 =\sum_f m_f\,w_1(R_f)
 \quad\hbox{in }H^1(X^\theta;\mathbb Z/2).                    \tag{16.2}
\]
Equivalently, for every loop \(\gamma\) in \(X^\theta\),
\[
 {\rm Hol}_{R_{\rm tot}}(\gamma)
 =\prod_f{\rm Hol}_{R_f}(\gamma)^{m_f},                        \tag{16.3}
\]
and the combined line is orientable if and only if
\[
 \sum_fm_f
 \left\langle w_1(R_f),[\gamma]\right\rangle=0\pmod2          \tag{16.4}
\]
for every loop.
\end{theorem}

\begin{proof}
Choose local unit frames for every \(R_f\).  Their transition functions
are signs \(\epsilon_{f,ab}\in\{\pm1\}\).  The tensor-product frame has
transition sign
\[
 \epsilon_{{\rm tot},ab}
 =\prod_f\epsilon_{f,ab}^{m_f}.                               \tag{16.5}
\]
Passing from multiplicative signs to additive \(\mathbb Z/2\) cocycles
turns (16.5) into (16.2).  Multiplication around a loop proves (16.3).
A real line is orientable exactly when all of its sign holonomies are one,
which gives (16.4).
\end{proof}

\begin{corollary}[Spectral-flow form]
\label{cor:v19v13-flow-sum}
Suppose the species families are generic along \(\gamma\).  For a real
symmetric determinant family let \(\nu_f(\gamma)\) be the sum of crossing
kernel dimensions, and for a skew Pfaffian family let it be the sum of
half-kernel dimensions.  Then
\[
 \left\langle w_1(R_{\rm tot}),[\gamma]\right\rangle
 =\sum_fm_f\nu_f(\gamma)\pmod2.                               \tag{16.6}
\]
\end{corollary}

\begin{proof}
Apply Theorems \ref{thm:v19v11-det-sign},
\ref{thm:v19v11-pf-sign}, and \ref{thm:v19v11-w1} to each species, then
use Theorem \ref{thm:v19v13-w1-sum}.
\end{proof}

\begin{corollary}[The square of a real line is canonically oriented]
\label{cor:v19v13-square}
For every metrized real line \(R\), the tensor square \(R^{\otimes2}\) is
orientable.  A local unit vector \(e\) defines the global positive tensor
\(e\otimes e\), independent of the replacement \(e\mapsto-e\).
\end{corollary}

\begin{proof}
Equation (16.2) gives \(w_1(R^{\otimes2})=2w_1(R)=0\).
The displayed local construction supplies the stated orientation directly.
\end{proof}

The last corollary explains the sign security of an identical even
multiplicity.  It does not authorize adding a mirror species: changing a
multiplicity changes the chiral theory, just as the vectorlike determinant
square did in V7.

\subsection{Real cancellation is weaker than complex phase cancellation}

Let \(L_f\) be the complex determinant lines before restriction to the
reflection-fixed locus, and put \(L_{\rm tot}=\bigotimes_fL_f^{\otimes
m_f}\).

\begin{proposition}[Two independent descent equations]
\label{prop:v19v13-two-obstructions}
The complex and Real obstruction classes obey
\[
 c_1(L_{\rm tot})=\sum_fm_fc_1(L_f),\qquad
 w_1(R_{\rm tot})=\sum_fm_fw_1(R_f).                           \tag{16.7}
\]
Vanishing of the second sum does not imply vanishing of the first, nor
does it force the \(U(1)\) holonomy of a flat \(L_{\rm tot}\) to be
trivial.
\end{proposition}

\begin{proof}
The first identity follows by multiplying \(U(1)\) transition functions;
the second is Theorem \ref{thm:v19v13-w1-sum}.  A \(U(1)\) phase retains
more information than its possible restriction to a sign.  For example,
two flat factors with the same holonomy \(e^{i\alpha}\) have trivial
doubled sign character but product holonomy \(e^{2i\alpha}\), which need
not be one.
\end{proof}

\begin{theorem}[Conditional positive multiplet chamber]
\label{thm:v19v13-positive-chamber}
Let \(C\) be a connected component of \(X^\theta\).  Assume:
\begin{enumerate}
 \item the complex product line \(L_{\rm tot}\) passes the V8 descent test;
 \item \(R_{\rm tot}|_C\) passes (16.4);
 \item the combined determinant or Pfaffian section
       \(\Delta_{\rm tot}=\prod_f\Delta_f^{m_f}\) has no zero on \(C\).
\end{enumerate}
Then a Real orientation can be chosen so that
\(\Delta_{\rm tot}>0\) everywhere on \(C\).  A V12 gap certificate for the
combined block is sufficient for the third item.
\end{theorem}

\begin{proof}
The first two assumptions give a scalar reflection-real coefficient on
\(C\).  It is a continuous nonzero real function, so Proposition
\ref{prop:v19v9-sign} makes its sign constant.  Choose the Real orientation
to make it positive at one reference point; it is then positive throughout
\(C\).  The last assertion is Theorem \ref{thm:v19v12-chamber}.
\end{proof}

\begin{criterion}[Finite global-sign audit]
\label{crit:v19v13-audit}
On a finite configuration-space complex, choose a cycle basis as in V8.
For each fermion species and each cycle, compute the V11 crossing parity.
Arrange the results as a matrix over \(\mathbb Z/2\).  The multiplicity
vector \(m=(m_f)\) cancels the Real sign obstruction precisely when
\[
 \nu m=0\quad\hbox{over }\mathbb Z/2.                          \tag{16.8}
\]
The \(U(1)\) transition phases must still be audited separately.
\end{criterion}

\begin{proof}
Each row of (16.8) is condition (16.4) for one cycle-basis element.
Linearity over \(\mathbb Z/2\) and the cycle-basis theorem of V8 extend it
to every loop.  Proposition \ref{prop:v19v13-two-obstructions} proves the
last sentence.
\end{proof}

\subsection{Executable certificate}

The verifier
\[
 \texttt{scripts/verify\_sm\_v19\_v13\_multiplet\_sign\_holonomy.py}
                                                                  \tag{16.9}
\]
uses three cycle generators and three nontrivial species characters whose
sum is zero.  It checks multiplication of their sign holonomies, exhibits
the obstruction left by an incomplete multiplet, verifies that an
identical real-line square is orientable, and gives a
\(\mathbb Z/12\mathbb Z\) witness where Real parity cancels but complex
phase holonomy does not.  It emits
\[
 \texttt{artifacts/sm\_v19\_v13\_multiplet\_sign\_holonomy.json}.
                                                                  \tag{16.10}
\]
The hashes are
\[
\begin{array}{c|l}
\hbox{object}&\hbox{SHA--256}\\ \hline
\hbox{verifier}&
\texttt{A6DAC4B3F4AA3B353E00FCC0445C51B7DD7B9D054D3D4D7CA646A88E81B64AF0}\\
\hbox{canonical payload}&
\texttt{3B0AC8054191FBCD5BF77213D8573C530B25DC786E7F3C0622B19712AFBAAA99}.
\end{array}                                                     \tag{16.11}
\]

\begin{firewall}[The cancellation matrix has not been evaluated for the Standard Model]
\label{fire:v19v13}
Equation (16.8) is the exact missing global-sign test, not its Standard
Model evaluation.  The regulator-specific cycle basis, crossing parities,
complex holonomies, determinant zeros, and reflected-Grassmann Gram
factorization remain to be constructed.  Perturbative charge-anomaly
cancellation does not fill those entries.
\end{firewall}

\begin{decision}[V14 acquisition]
\label{dec:v19v13-next}
V14 will return from determinant lines to the local transfer kernel.  It
will prove that direct sums and unitary flavour rotations preserve the V7
positive cross matrix, and will derive a Schur-complement criterion for
off-diagonal flavour or Yukawa cross blocks.  This will identify which
multiplet couplings preserve the reflected-Grassmann cone before the
global phase audit.
\end{decision}

\begin{assessment}[V13 acquisition]
\label{ass:v19v13}
The global sign problem is now correctly posed for the full fermion
multiplet: mod-two crossing characters add, and their sum must vanish on
every loop.  This cancellation can occur without vectorlike doubling, but
it neither cancels the complex determinant phase nor supplies local
Grassmann positivity.
\end{assessment}

\begin{record}[V13 append record]
\label{rec:v19v13}
V13 is appended to the validated V12 whose installed SHA--256 was
\texttt{F40D6AC282D1F06519F25E1EFF4F7A4DD569EF4F0FF9A77EC3669ACF87AE1C40}.
After installation only V13 is the active numeric SM note; frozen archives
and unrelated notes are unchanged.
\end{record}




\section{V14: multiplet cross-block positivity and the Schur gate}
\label{sec:v19v14}

V13 gives the global mod-two cancellation law for the full fermion
multiplet.  Local reflection positivity remains a separate matrix
question: the complete temporal cross matrix must be positive before the
V7 exterior-power Gram theorem applies.  This note proves stability under
multiplet assembly and unitary flavour changes, then gives an exact
Schur-complement test for off-diagonal cross couplings.

\subsection{Direct sums and flavour rotations}

\begin{proposition}[Stable operations on the cross matrix]
\label{prop:v19v14-stable}
Let \(C_f\geq0\) be Hermitian cross matrices on finite one-particle spaces
\(H_f\).  Then
\[
 C_{\rm dir}:=\bigoplus_f C_f\geq0.                            \tag{17.1}
\]
For every unitary \(U\) on \(\bigoplus_fH_f\),
\[
 C_U:=UC_{\rm dir}U^*\geq0.                                   \tag{17.2}
\]
Consequently, the Grassmann cross exponentials associated with
\(C_{\rm dir}\) and \(C_U\) belong to the reflected-square cone of V7.
\end{proposition}

\begin{proof}
For \(v=(v_f)_f\),
\[
 \langle v,C_{\rm dir}v\rangle
 =\sum_f\langle v_f,C_fv_f\rangle\geq0.
\]
For (17.2),
\(\langle v,C_Uv\rangle=\langle U^*v,C_{\rm dir}U^*v\rangle\geq0\).
Theorem \ref{thm:v19v7-gram} then puts both cross exponentials in
\({\cal C}_\Theta\).  The unitary change of generators induces the
corresponding unitary maps on every exterior power, so it does not alter
the Gram conclusion.
\end{proof}

Thus ordinary flavour-basis changes cause no local positivity problem.
An off-diagonal block that changes the positive operator itself must pass
the following test.

\subsection{The exact block criterion}

Let
\[
 C=\begin{pmatrix}A&Y\\Y^*&B\end{pmatrix},                    \tag{17.3}
\]
where \(A=A^*\), \(B=B^*\), and \(Y\) is rectangular.

\begin{theorem}[Strict Schur-complement criterion]
\label{thm:v19v14-strict-schur}
If \(A>0\), then
\[
 C\geq0
 \quad\Longleftrightarrow\quad
 S:=B-Y^*A^{-1}Y\geq0.                                       \tag{17.4}
\]
Indeed,
\[
 C=
 \begin{pmatrix}I&0\\Y^*A^{-1}&I\end{pmatrix}
 \begin{pmatrix}A&0\\0&S\end{pmatrix}
 \begin{pmatrix}I&A^{-1}Y\\0&I\end{pmatrix}.                  \tag{17.5}
\]
\end{theorem}

\begin{proof}
The first and third factors in (17.5) are adjoints, and direct
multiplication gives (17.3), so \(S\geq0\) implies \(C\geq0\).
Conversely, for any \(v\), evaluate the quadratic form of \(C\) at
\((-A^{-1}Yv,v)\).  The result is \(\langle v,Sv\rangle\), which must be
nonnegative if \(C\geq0\).
\end{proof}

Wilson temporal projectors are singular, so the Moore--Penrose form is
needed.

\begin{theorem}[Singular Schur-complement criterion]
\label{thm:v19v14-singular-schur}
Let \(A\geq0\) and let \(A^+\) be its Moore--Penrose inverse.  Then the
block matrix (17.3) is positive if and only if
\[
 \operatorname{Ran}Y\subseteq\operatorname{Ran}A,\qquad
 B-Y^*A^+Y\geq0.                                               \tag{17.6}
\]
When (17.6) holds, putting \(X=A^+Y\) gives \(Y=AX\) and
\[
 C=
 \begin{pmatrix}I&0\\X^*&I\end{pmatrix}
 \begin{pmatrix}A&0\\0&B-Y^*A^+Y\end{pmatrix}
 \begin{pmatrix}I&X\\0&I\end{pmatrix}.                         \tag{17.7}
\]
\end{theorem}

\begin{proof}
Assume first that \(C\geq0\).  Its principal block gives \(A\geq0\).
For \(u\in\ker A\), \(v\) arbitrary, and \(t\in\mathbb C\), positivity at
\((tu,v)\) gives
\[
 2\operatorname{Re}\!\left(\overline t\,\langle u,Yv\rangle\right)
 +\langle v,Bv\rangle\geq0.                                   \tag{17.8}
\]
Varying the magnitude and phase of \(t\) forces
\(\langle u,Yv\rangle=0\).  Therefore
\(\operatorname{Ran}Y\perp\ker A\), which in finite dimension is the first
condition in (17.6).  Evaluation at \((-A^+Yv,v)\) then gives the second
condition.

Conversely, the range condition implies \(AA^+Y=Y\).  Hence, with
\(X=A^+Y\), one has \(AX=Y\) and
\[
 X^*AX=Y^*A^+Y.
\]
Multiplication proves (17.7), whose middle factor is positive by (17.6).
\end{proof}

\begin{corollary}[Quantitative small-coupling gate]
\label{cor:v19v14-small}
If
\[
 A\geq aI,\qquad B\geq bI,\qquad a,b>0,
\]
then \(C\geq0\) whenever
\[
 \|Y\|_{\rm op}^2\leq ab.                                     \tag{17.9}
\]
If the last inequality is strict, then \(C>0\).
\end{corollary}

\begin{proof}
Since \(A^{-1}\leq a^{-1}I\),
\[
 B-Y^*A^{-1}Y
 \geq\left(b-\frac{\|Y\|_{\rm op}^2}{a}\right)I.               \tag{17.10}
\]
Apply Theorem \ref{thm:v19v14-strict-schur}.
\end{proof}

\begin{corollary}[Projector support gate]
\label{cor:v19v14-projector}
Let \(A=\kappa P\), where \(P\) is an orthogonal projector and
\(\kappa>0\).  Positivity of (17.3) requires
\[
 (I-P)Y=0.                                                     \tag{17.11}
\]
On \(\operatorname{Ran}P\), the remaining condition is
\[
 B-\kappa^{-1}Y^*Y\geq0.                                      \tag{17.12}
\]
\end{corollary}

\begin{proof}
Here \(\operatorname{Ran}A=\operatorname{Ran}P\) and
\(A^+=\kappa^{-1}P\).  The two conditions in Theorem
\ref{thm:v19v14-singular-schur} become (17.11)--(17.12).
\end{proof}

\subsection{Placement of Standard Model couplings}

\begin{criterion}[Local multiplet reflection gate]
\label{crit:v19v14-local}
After temporal gauge transport and fermionic reflection phases have been
absorbed into the fibre identification, assemble every term that crosses
the reflection plane into one Hermitian matrix \(C_{\rm tot}\).
If \(C_{\rm tot}\geq0\), then
\[
 \exp\!\left(\sum_{i,j}\Theta(\psi_i)
 (C_{\rm tot})_{ij}\psi_j\right)\in{\cal C}_\Theta.            \tag{17.13}
\]
Reflected mass, spatial, Higgs--Yukawa, and gauge interactions lying
strictly in the two open halves enter as \(\Theta(B)B\) and need not satisfy
a Schur inequality.  Any such interaction placed directly on a crossing
hinge must instead pass (17.4) or (17.6).
\end{criterion}

\begin{proof}
Equation (17.13) is Theorem \ref{thm:v19v7-gram}.  The treatment of
one-side terms is the last assertion of that theorem.  The two block
criteria are necessary and sufficient for a crossing block of the form
(17.3).
\end{proof}

\begin{remark}[Separation from the global phase]
\label{rem:v19v14-global}
The criterion is local and basis invariant.  Even if it holds at every
finite crossing hinge, the determinant lines must still pass the V8
complex-holonomy test and the V13 combined mod-two test.  Conversely, those
global cancellations cannot repair a negative eigenvalue of
\(C_{\rm tot}\).
\end{remark}

\subsection{Executable certificate}

The verifier
\[
 \texttt{scripts/verify\_sm\_v19\_v14\_fermion\_cross\_block\_positivity.py}
                                                                  \tag{17.14}
\]
checks exact rational factorizations for both the strict and singular
Schur complements.  It exhibits a range-violating block with quadratic
form \(-1\), verifies positivity under a rational orthogonal flavour
rotation, and checks the quantitative bound
\(a=2\), \(b=3\), \(\|Y\|=2\), whose Schur lower bound is \(1\).
It emits
\[
 \texttt{artifacts/sm\_v19\_v14\_fermion\_cross\_block\_positivity.json}.
                                                                  \tag{17.15}
\]
The hashes are
\[
\begin{array}{c|l}
\hbox{object}&\hbox{SHA--256}\\ \hline
\hbox{verifier}&
\texttt{7A7D991F1792D43D4390835805FD3DB552538C139AD98F8C9214E53E1B1207B5}\\
\hbox{canonical payload}&
\texttt{81EC8FB705E8809F3EFC9131AC18B141EEE362F5B6D19DF8C6C926C35CC06ADA}.
\end{array}                                                     \tag{17.16}
\]

\begin{firewall}[A local positive cross matrix is conditional]
\label{fire:v19v14}
No explicit chiral Standard Model regulator has yet been shown to produce
a Hermitian \(C_{\rm tot}\) satisfying these conditions on every
gauge--Higgs--metric background.  The Schur gate also does not control
determinant zeros, global phase holonomy, cutoff removal, or the continuum
Hamiltonian spectrum.
\end{firewall}

\begin{decision}[V15 checkpoint]
\label{dec:v19v14-next}
V15 is the next mandatory companion audit.  If it contributes no new
matter-sector input, V16 should construct a finite combined bosonic and
fermionic transfer theorem from the V4, V6, V7, and V14 positive factors,
with all global determinant-line assumptions stated explicitly.
\end{decision}

\begin{assessment}[V14 acquisition]
\label{ass:v19v14}
The local fermion gate is now exact at the multiplet level.  Direct sums
and unitary flavour rotations preserve the reflected-Grassmann cone.
Off-diagonal crossing couplings preserve it precisely when their strict or
singular Schur complement is positive.
\end{assessment}

\begin{record}[V14 append record]
\label{rec:v19v14}
V14 is appended to the validated V13 whose installed SHA--256 was
\texttt{2369598A829DC010B18BC28B4ECDD28460F08FFAEF2FC069BFAA227889B461D6}.
After installation only V14 is the active numeric SM note; frozen archives
and unrelated notes are unchanged.
\end{record}



\section{V15: mandatory companion audit before transfer assembly}
\label{sec:v19v15}

V14 completed the local Schur gate for fermion cross blocks.  The standing
checkpoint rule now requires another read-only companion audit before the
bosonic and fermionic transfer factors are assembled.

\begin{audit}[Companion state]
\label{aud:v19v15-files}
The active Yang--Mills note remains
\[
 \texttt{Renewal\_Yang\_Mills\_Mass\_Gap\_Research\_Notes\_V25.tex},
 \qquad
 \operatorname{SHA256}=
 \texttt{5D6759B911E076304A17D09955E6777C3ACD6C99D3B47CD42BBA75EB0EA96BE6}.
                                                                  \tag{18.1}
\]
The active Navier--Stokes note remains
\[
 \texttt{Renewal\_NS\_Stability\_Research\_Notes\_V26.tex},
 \qquad
 \operatorname{SHA256}=
 \texttt{82C887F8C2C69E4F28FAD7C3DC8DE5B9099CB5A4BF431EBA1FFF3E91935978AD}.
                                                                  \tag{18.2}
\]
These are the same files and hashes recorded at V10.
\end{audit}

\begin{proposition}[Transfer direction is unchanged]
\label{prop:v19v15-direction}
The companion notes supply no new gauge--Higgs cross kernel, fermionic
reflection cone, determinant-line phase, or combined transfer theorem.
V16 should therefore assemble the finite positive sectors already proved
in V4, V6, V7, and V14, while retaining the global chiral conditions of
V8--V13 as explicit hypotheses.
\end{proposition}

\begin{proof}
Yang--Mills V25 still ends at its finite Laurent/Haar certificate gate and
retains the exact-link coarsening rule only for a later continuum stage.
Navier--Stokes V26 still ends with rank-one Oseen-kernel norms and a
programme for improved constants.  Neither conclusion acts on the
boundary Hilbert--Grassmann space used in V14.  The unchanged hashes
exclude a later uninspected appendix.
\end{proof}

\begin{decision}[V16 acquisition]
\label{dec:v19v15-next}
V16 will state and prove a finite gauge--Higgs--fermion reflection theorem.
The proof will tensor the V4 gauge heat operator with the V6 balanced Higgs
operator, insert the V7/V14 Grassmann cone, and compress by simultaneous
gauge averaging.  It will state separately the hypotheses needed to
interpret a chiral determinant globally.
\end{decision}

\begin{firewall}[Audit boundary]
\label{fire:v19v15}
This audit proves no new positivity or continuum estimate.  It only
confirms that transfer assembly is the next nonduplicative acquisition.
\end{firewall}

\begin{record}[V15 append record]
\label{rec:v19v15}
V15 is appended to the validated V14 whose installed SHA--256 was
\texttt{7C6177ACE0EE3CFB0D7A1F90C60D989A79A7C67403438E69391AD76C511AF0D3}.
After installation only V15 is the active numeric SM note.
\end{record}



\section{V16: a finite hybrid gauge--Higgs--fermion transfer square}
\label{sec:v19v16}

V4 gives a positive heat-kernel gauge transfer, V6 gives a positive
balanced or reflected-pair Higgs transfer, and V7 with V14 gives a
reflected-square cone for fermionic cross blocks.  This note assembles
those ingredients without replacing Berezin integration by an ordinary
positive measure.  The common structure is a tensor product of a bosonic
Hilbert quadratic form with a fermionic coefficient Hilbert form.

\subsection{The fermionic cone induces a positive coefficient form}

Let \({\mathfrak A}_+\) and \({\cal L}\) be as in V7, and let
\[
 W_F=\sum_{\alpha}\Theta(C_\alpha)C_\alpha
 \in{\cal C}_\Theta,                                           \tag{19.1}
\]
with each \(C_\alpha\) homogeneous.  For even
\(A,B\in{\mathfrak A}_+\), define
\[
 (A,B)_{W_F}:={\cal L}\!\left(\Theta(A)B\,W_F\right).           \tag{19.2}
\]

\begin{lemma}[Fermionic coefficient Hilbert form]
\label{lem:v19v16-fermion-form}
The form (19.2) is positive semidefinite and has the explicit expression
\[
 (A,B)_{W_F}
 =\sum_\alpha\langle AC_\alpha,BC_\alpha\rangle_{\rm coeff},   \tag{19.3}
\]
where the coefficient inner product is the one in Lemma
\ref{lem:v19v7-coefficient}.  After quotienting its null space and
completing, it defines a finite-dimensional Hilbert space
\({\cal H}_{F,W}\).
\end{lemma}

\begin{proof}
Insert (19.1) into (19.2).  Since \(A\) and \(B\) are even, they pass
through reflected homogeneous factors without a graded sign and commute
with \(C_\alpha\).  The anti-automorphism law gives
\[
 {\cal L}\!\left(\Theta(A)B\Theta(C_\alpha)C_\alpha\right)
 ={\cal L}\!\left(\Theta(AC_\alpha)(BC_\alpha)\right).
\]
Equation (10.2) is exactly the coefficient inner product, proving (19.3).
Its value at \(A=B\) is a sum of squared norms.  The standard null-space
quotient of a positive semidefinite sesquilinear form is a Hilbert space;
finite dimensionality makes completion automatic.
\end{proof}

\subsection{The hybrid tensor-square theorem}

Let \((X,\nu)\) be a bosonic boundary measure space and let \(T_B\geq0\)
be a bounded self-adjoint operator on \({\cal H}_B=L^2(X,\nu)\).
For an even polynomial-valued boundary amplitude
\(\Psi:X\to{\mathfrak A}_+\), regard its coefficient class as an element
of \({\cal H}_B\otimes{\cal H}_{F,W}\).

\begin{theorem}[Hybrid boson--Grassmann positivity]
\label{thm:v19v16-hybrid}
Under the preceding assumptions,
\[
 {\cal Q}_{B,F}(\Psi)
 :=\left\langle\Psi,
 (T_B\otimes I_{{\cal H}_{F,W}})\Psi\right\rangle
 \geq0.                                                        \tag{19.4}
\]
If \(T_B\) has a kernel \(K_B(x,y)\), then (19.4) is the rigorous meaning
of
\[
 \iint K_B(x,y)\,
 {\cal L}\!\left(\Theta(\Psi(x))\Psi(y)W_F\right)
 \,d\nu(x)d\nu(y)\geq0.                                       \tag{19.5}
\]
\end{theorem}

\begin{proof}
Choose an orthonormal basis \(e_1,\ldots,e_r\) of
\({\cal H}_{F,W}\) and write
\(\Psi=\sum_{\ell=1}^r\psi_\ell\otimes e_\ell\).  Then
\[
 {\cal Q}_{B,F}(\Psi)
 =\sum_{\ell=1}^r\langle\psi_\ell,T_B\psi_\ell\rangle_{{\cal H}_B}
 \geq0.                                                        \tag{19.6}
\]
When a kernel representation is available, expanding the left side in
the coefficient basis and using (19.2) gives (19.5).  Boundedness and
square integrability justify the finite coefficient interchange.
\end{proof}

\begin{corollary}[Gauge compression and one-side weights]
\label{cor:v19v16-compression}
Let \({\cal P}\) be any orthogonal projection on the hybrid boundary
Hilbert space and let \(M\) be multiplication by a bounded nonnegative
reflection-plane weight.  Then
\[
 {\cal K}_{\rm hyb}
 :=M^{1/2}{\cal P}(T_B\otimes I){\cal P}M^{1/2}\geq0.           \tag{19.7}
\]
\end{corollary}

\begin{proof}
For every \(\Psi\),
\[
 \langle\Psi,{\cal K}_{\rm hyb}\Psi\rangle
 =\left\|(T_B^{1/2}\otimes I)
 {\cal P}M^{1/2}\Psi\right\|^2\geq0.                           \tag{19.8}
\]
\end{proof}

\subsection{Finite coupled reflection theorem}

\begin{theorem}[Conditional finite gauge--Higgs--fermion reflection
positivity]
\label{thm:v19v16-coupled}
Consider a finite reflected lattice slab and impose the following
conditions.
\begin{enumerate}
 \item Gauge plaquettes have the heat-kernel transfer of V4.
 \item Every crossing Higgs factor is either balanced as in Theorem
       \ref{thm:v19v6-balanced}, or occurs in a reflected \(d,-d\)
       two-step pair as in Theorem \ref{thm:v19v6-paired}.
 \item After temporal transport and reflection phases are absorbed into
       the fibre identification, the full fermionic cross matrix
       \(C_{\rm tot}\) is positive, as tested by V14.
 \item All remaining action factors occur as reflection-matched one-side
       factors, and the bosonic factors are nonnegative and integrable.
\end{enumerate}
Then the finite unintegrated fermion theory is reflection positive on its
even gauge-invariant observable algebra:
\[
 {\cal E}\!\left(\Theta(F)F\right)\geq0.                       \tag{19.9}
\]
The statement remains true after a measurable direct integral over
reflection-fixed parameters for which all four conditions hold almost
everywhere and the resulting form is integrable.
\end{theorem}

\begin{proof}
The gauge and Higgs crossing operators are positive by V4 and V6, so their
tensor product is a positive bosonic operator \(T_B\).  By V7,
\(C_{\rm tot}\geq0\) makes its Grassmann exponential an element \(W_F\) of
\({\cal C}_\Theta\).  Reflected fermionic one-side factors multiply this
density as \(\Theta(B)B\) and remain in the cone.  Lemma
\ref{lem:v19v16-fermion-form} and Theorem \ref{thm:v19v16-hybrid} combine
the sectors.

Simultaneous boundary gauge averaging is an orthogonal projection on the
bosonic and fermionic coefficient spaces.  Reflection-plane bosonic
potentials give the multiplication \(M\).  Corollary
\ref{cor:v19v16-compression} therefore makes the reflected expectation of
the boundary amplitude produced by \(F\) nonnegative, which is (19.9).
For a direct integral, the conditional quadratic form is nonnegative
almost everywhere; integrating that nonnegative function preserves the
inequality.
\end{proof}

\begin{criterion}[Additional hypotheses for a chiral interpretation]
\label{crit:v19v16-chiral}
For Theorem \ref{thm:v19v16-coupled} to define a chiral Standard Model
measure on the gauge-field quotient, rather than a local Grassmann
calculation in chosen bases, the regulator must also prove:
\begin{enumerate}
 \item V8 complex determinant-line descent;
 \item V9 reflection-real descent;
 \item V13 cancellation of the combined mod-two sign character;
 \item compatibility of the local cone decompositions under every
       determinant-line transition.
\end{enumerate}
A V12 gap certificate may replace explicit V11 crossing enumeration only
on a zero-free connected chamber.
\end{criterion}

\begin{proof}
The first three items are respectively necessary for a global scalar
fermion phase, its reflection reality, and its consistent real
orientation.  The fourth ensures that the local positive coefficient
forms in Lemma \ref{lem:v19v16-fermion-form} patch to one global form.
Without it, (19.9) is chart dependent.  The final sentence is Theorem
\ref{thm:v19v12-chamber}.
\end{proof}

\subsection{Executable certificate}

The verifier
\[
 \texttt{scripts/verify\_sm\_v19\_v16\_hybrid\_transfer\_square.py}
                                                                  \tag{19.10}
\]
constructs exact rational bosonic and fermionic Gram matrices, verifies
their tensor-Gram identity, compresses by a nontrivial orthogonal gauge
average, dresses by unequal positive half-weights, and checks three
quadratic forms as explicit squared norms.  It emits
\[
 \texttt{artifacts/sm\_v19\_v16\_hybrid\_transfer\_square.json}.
                                                                  \tag{19.11}
\]
The hashes are
\[
\begin{array}{c|l}
\hbox{object}&\hbox{SHA--256}\\ \hline
\hbox{verifier}&
\texttt{A0E868DBBD1172F278750EFEBABA9C0AF8C84E0D672917CB24FDAE8BF16396F2}\\
\hbox{canonical payload}&
\texttt{84486E2AA00F2F3530495DA2D2B99EE13B31ABE57D4223DF6F5844C94C66D3B9}.
\end{array}                                                     \tag{19.12}
\]

\begin{firewall}[Finite conditional coupling is not the SM--Einstein
continuum]
\label{fire:v19v16}
Theorem \ref{thm:v19v16-coupled} is finite and conditional.  No chiral
regulator has been shown to satisfy the global patching hypotheses, and
no positive, normalizable Einstein metric measure has been constructed.
The result supplies neither cutoff-uniform tightness nor
Osterwalder--Schrader reconstruction, clustering, a physical Hamiltonian,
or a simultaneous continuum limit.
\end{firewall}

\begin{decision}[V17 acquisition]
\label{dec:v19v16-next}
V17 will go upstream to the metric measure.  It will analyze the Euclidean
Einstein--Hilbert conformal mode and prove the precise unbounded-action
obstruction to using \(e^{-S_{\rm EH}}\) as the missing positive metric
weight.  Any later gravitational transfer proposal must first answer that
obstruction without invalidating the reflection structure proved here.
\end{decision}

\begin{assessment}[V16 acquisition]
\label{ass:v19v16}
The finite local Standard Model sectors now possess one rigorous hybrid
reflection-positive boundary form under explicit hypotheses.  The proof
uses a bosonic operator square and a fermionic coefficient square, so it
does not misinterpret Berezin integration as a positive measure.  The
largest upstream obstruction is now the gravitational metric measure.
\end{assessment}

\begin{record}[V16 append record]
\label{rec:v19v16}
V16 is appended to the validated V15BELLO? whose installed SHA--256 was
\texttt{E96BEC?}.
After installation only V16 is the active numeric SM note; frozen archives
and unrelated notes are unchanged.
\end{record}



\section{V17: quadratic-curvature stabilization and its positivity gate}
\label{sec:v19v17}

\begin{record}[Append-only correction to the V16 append record]
\label{rec:v19v17-v16-correction}
The final record in V16 contains corrupted prose and an incomplete hash.
It is superseded by this append-only correction: V16 was appended to the
validated V15 whose installed SHA--256 was
\[
 \texttt{EB4EC553C0CB64A3187EFCCEF08248A772DE9097F993D17A0F238304EA5B0BE9}.
                                                               \tag{20.1}
\]
The installed V16 SHA--256 was
\[
 \texttt{FC6571074E73866E65C6CF32610942E7B1B1E0BC9F031241829403C8A7824282}.
                                                               \tag{20.2}
\]
No earlier bytes are changed.
\end{record}

The V16 decision proposed proving the Euclidean Einstein--Hilbert
conformal instability.  That proof already belongs to the retained
archive: the gravity ledger in
\[
 \texttt{Renewal\_SM\_Einstein\_Continuum\_Research\_Notes\_V100\_f18.tex}
\]
records unbounded action along fixed-volume conformal oscillations.  The
archive SHA--256 is
\[
 \texttt{AE0B9CFE523D98B7C1C5E43C8EFAC4C656A650DC1F432019753CD65D54500D36}.
                                                               \tag{20.3}
\]
We do not repeat that result.  Instead, this note asks what a positive
scalar-curvature-square regulator repairs and whether that repair fits the
positive transfer structure of V16.

\subsection{Exact lower bounds from an \(R^2\) stabilizer}

Let \(M\) be a closed four-manifold, let \(\kappa>0\),
\(\alpha>0\), and \(\Lambda\in\mathbb R\).  Consider
\[
 S_{\alpha,\Lambda}(g)
 :=\int_M\left(\alpha R_g^2-\kappa R_g+2\kappa\Lambda\right)
 \,dV_g.                                                        \tag{20.4}
\]
The second and third terms use the Euclidean Einstein--Hilbert sign
convention relevant to the archived conformal obstruction.

\begin{theorem}[Curvature-square completion]
\label{thm:v19v17-completion}
For every smooth Riemannian metric \(g\),
\[
\begin{split}
 S_{\alpha,\Lambda}(g)
 &=
 \alpha\int_M
 \left(R_g-\frac{\kappa}{2\alpha}\right)^2\,dV_g\\
 &\quad+
 \left(2\kappa\Lambda-\frac{\kappa^2}{4\alpha}\right)
 \operatorname{Vol}_g(M).                                     \tag{20.5}
\end{split}
\]
Consequently, on the fixed-volume slice
\(\operatorname{Vol}_g(M)=V\),
\[
 S_{\alpha,\Lambda}(g)
 \geq V\left(2\kappa\Lambda-\frac{\kappa^2}{4\alpha}\right).   \tag{20.6}
\]
If \(\Lambda\geq\kappa/(8\alpha)\), the action is nonnegative without
fixing the volume.
\end{theorem}

\begin{proof}
The pointwise identity
\[
 \alpha R^2-\kappa R
 =\alpha\left(R-\frac{\kappa}{2\alpha}\right)^2
  -\frac{\kappa^2}{4\alpha}
\]
integrates to (20.5).  Dropping the nonnegative square proves (20.6).
The stated condition makes the remaining volume coefficient nonnegative.
\end{proof}

\begin{theorem}[Scalar-curvature coercive estimate]
\label{thm:v19v17-coercive}
For every \(g\),
\[
 S_{\alpha,\Lambda}(g)
 \geq
 \frac{\alpha}{2}\int_MR_g^2\,dV_g
 +
 \left(2\kappa\Lambda-\frac{\kappa^2}{2\alpha}\right)
 \operatorname{Vol}_g(M).                                     \tag{20.7}
\]
Thus, on a fixed-volume slice, every action sublevel has a uniform
\(L^2\)-scalar-curvature bound.
\end{theorem}

\begin{proof}
Young's identity gives
\[
 \alpha R^2-\kappa R
 -\frac{\alpha}{2}R^2+\frac{\kappa^2}{2\alpha}
 =\frac{\alpha}{2}\left(R-\frac{\kappa}{\alpha}\right)^2
 \geq0.                                                        \tag{20.8}
\]
Integrate (20.8) and add the cosmological term.  If
\(S_{\alpha,\Lambda}(g)\leq C\) and the volume equals \(V\), rearranging
(20.7) bounds \(\int R_g^2\,dV_g\).
\end{proof}

\begin{firewall}[Scalar curvature does not control metric compactness]
\label{fire:v19v17-scalar}
Estimate (20.7) controls only scalar curvature.  It does not bound the Weyl
tensor or traceless Ricci tensor, prevent collapse, fix diffeomorphism
gauge, or make a regulated metric action proper.  It stabilizes one
curvature direction; it does not construct a gravitational measure.
\end{firewall}

\subsection{The removal estimate is singular}

\begin{proposition}[Degeneration as \(\alpha\downarrow0\)]
\label{prop:v19v17-removal}
For fixed \(V>0\), \(\kappa>0\), and \(\Lambda\), the proved lower bound
in (20.6) tends to \(-\infty\) as \(\alpha\downarrow0\):
\[
 V\left(2\kappa\Lambda-\frac{\kappa^2}{4\alpha}\right)
 \longrightarrow-\infty.                                     \tag{20.9}
\]
Therefore Theorem \ref{thm:v19v17-completion} supplies no lower bound
uniform in removal of the \(R^2\) regulator.
\end{proposition}

\begin{proof}
The term \(-V\kappa^2/(4\alpha)\) diverges to \(-\infty\), while
\(2V\kappa\Lambda\) is fixed.
\end{proof}

This concerns the proved estimate.  It does not assert that the exact
infimum attains the bound on every manifold.

\subsection{No positive real Gaussian linearization has the required sign}

\begin{theorem}[Positive auxiliary-field moment obstruction]
\label{thm:v19v17-auxiliary}
Let \(\alpha>0\).  There is no finite positive measure \(\mu\) on
\(\mathbb R\), with exponential moments near zero, such that
\[
 e^{-\alpha r^2}=\int_{\mathbb R}e^{\sigma r}\,d\mu(\sigma)
 \quad\hbox{for all \(r\) near zero}.                          \tag{20.10}
\]
The correct Gaussian linearization is
\[
 e^{-\alpha r^2}
 =
 \frac1{\sqrt{4\pi\alpha}}
 \int_{\mathbb R}
 e^{-\sigma^2/(4\alpha)}e^{i\sigma r}\,d\sigma,                \tag{20.11}
\]
which uses an imaginary coupling.
\end{theorem}

\begin{proof}
If (20.10) held, evaluation at zero would give \(\mu(\mathbb R)=1\).
Differentiation under the integral gives
\[
 \left.\frac{d^2}{dr^2}\right|_{r=0}
 \int e^{\sigma r}\,d\mu(\sigma)
 =\int\sigma^2\,d\mu(\sigma)\geq0.                             \tag{20.12}
\]
The target has second derivative \(-2\alpha<0\), a contradiction.
Equation (20.11) is the elementary Fourier transform of the centered
Gaussian.  Completing the square with a real coupling instead produces
\(e^{+\alpha r^2}\).
\end{proof}

\begin{corollary}[Interface with the V16 positive transfer]
\label{cor:v19v17-transfer}
The standard local Hubbard--Stratonovich representation of
\(e^{-\alpha R^2}\) cannot add a real auxiliary variable with a positive
Gaussian measure and a real linear curvature coupling to the V16 bosonic
Hilbert space.  The Fourier representation (20.11) introduces a complex
phase, so reflection positivity requires a separate proof.
\end{corollary}

\begin{proof}
A positive real Gaussian auxiliary field would give the positive Laplace
representation excluded by Theorem \ref{thm:v19v17-auxiliary}.  The
remaining representation is (20.11), whose integrand is not nonnegative.
\end{proof}

\subsection{Executable certificate}

The verifier
\[
 \texttt{scripts/verify\_sm\_v19\_v17\_quadratic\_curvature\_stabilizer.py}
                                                                  \tag{20.13}
\]
checks the completion identity and coercive bound with exact rational
quadrature data, records the negative second-moment obstruction, and
evaluates the degenerating fixed-volume lower bound at
\(\alpha=1,1/2,1/4,1/8\).  It emits
\[
 \texttt{artifacts/sm\_v19\_v17\_quadratic\_curvature\_stabilizer.json}.
                                                                  \tag{20.14}
\]
The hashes are
\[
\begin{array}{c|l}
\hbox{object}&\hbox{SHA--256}\\ \hline
\hbox{verifier}&
\texttt{C679EC9CE26F725AE1D90ED82E3947D48D84AE3182F695DB1A715518C8BC691F}\\
\hbox{canonical payload}&
\texttt{3E93ECA648270FC9E7168B579F77F59321798575E203FD01D109EC6873962E83}.
\end{array}                                                     \tag{20.15}
\]

\begin{firewall}[Finite stabilization is not the Einstein continuum]
\label{fire:v19v17}
Theorems \ref{thm:v19v17-completion} and
\ref{thm:v19v17-coercive} give finite-\(\alpha\) lower bounds.  No
reflection-positive, normalizable Einstein metric measure follows.
There is no cutoff-uniform tightness, diffeomorphism quotient, or removal
of \(\alpha\), and no simultaneous Standard Model--Einstein limit.
\end{firewall}

\begin{decision}[V18 acquisition]
\label{dec:v19v17-next}
V18 will test the Gaussian higher-derivative propagator directly.  Its
partial fractions have a negative residue.  The next note will construct
a positive-time test vector with a negative reflected quadratic form,
deciding whether this \(R^2\)-type stabilization can enter the V16
positive transfer by the standard Gaussian route.
\end{decision}

\begin{assessment}[V17 acquisition]
\label{ass:v19v17}
The archived conformal instability was not repeated.  A positive \(R^2\)
term gives an exact fixed-volume lower bound and scalar-curvature
coercivity, but the estimate degenerates under removal and the standard
auxiliary-field representation requires a complex coupling.
\end{assessment}

\begin{record}[V17 append record]
\label{rec:v19v17}
V17 is appended to the validated V16 whose installed SHA--256 was
\texttt{FC6571074E73866E65C6CF32610942E7B1B1E0BC9F031241829403C8A7824282}.
After installation only V17 is the active numeric SM note; frozen archives
and unrelated notes are unchanged.
\end{record}




\section{V18: a two-time reflection counterexample for higher derivatives}
\label{sec:v19v18}

V17 proved that an \(R^2\) term repairs a lower bound at fixed regulator
strength but did not establish reflection positivity.  The standard
Gaussian higher-derivative covariance already fails that test.  The
failure is finite and explicit: two positive times give a reflected
covariance matrix with negative determinant.

\subsection{The higher-derivative covariance}

Fix \(\alpha>0\) and a spatial momentum whose ordinary Euclidean frequency
is \(\omega>0\).  Put
\[
 \Omega:=\sqrt{\omega^2+\alpha^{-1}}>\omega.
\]
The scalar covariance for the quadratic operator
\[
 (-\partial_t^2+\omega^2)
 \left[1+\alpha(-\partial_t^2+\omega^2)\right]
\]
has Fourier transform
\[
 \widehat C_\alpha(p_0)
 =
 \frac1{(p_0^2+\omega^2)
 \left[1+\alpha(p_0^2+\omega^2)\right]}
 =
 \frac1{p_0^2+\omega^2}-\frac1{p_0^2+\Omega^2}.                \tag{21.1}
\]
Therefore
\[
 C_\alpha(t)
 =
 \frac{e^{-\omega|t|}}{2\omega}
 -
 \frac{e^{-\Omega|t|}}{2\Omega}.                              \tag{21.2}
\]

\begin{proof}
The algebraic identity
\[
 \frac1{x(1+\alpha x)}
 =\frac1x-\frac1{x+\alpha^{-1}}
\]
gives (21.1).  Fourier inversion of
\((p_0^2+a^2)^{-1}\) gives \(e^{-a|t|}/(2a)\), proving (21.2).
\end{proof}

\subsection{The reflected kernel is indefinite}

For a time-translation-invariant Gaussian, reflection positivity across
\(t=0\) requires
\[
 \sum_{i,j=1}^n\overline c_i c_j C_\alpha(t_i+t_j)\geq0
 \quad\hbox{whenever }t_i>0.                                  \tag{21.3}
\]

\begin{theorem}[Strict two-time obstruction]
\label{thm:v19v18-two-time}
For any two distinct positive times \(t_1,t_2\), the matrix
\[
 K_{ij}:=C_\alpha(t_i+t_j),\qquad 1\leq i,j\leq2,              \tag{21.4}
\]
has strictly negative determinant.  Hence (21.3) fails.
\end{theorem}

\begin{proof}
Set
\[
 a=\frac1{2\omega},\quad b=\frac1{2\Omega},\quad
 u_i=e^{-\omega t_i},\quad v_i=e^{-\Omega t_i}.
\]
Equation (21.2) gives
\[
 K=a\,uu^{\mathsf T}-b\,vv^{\mathsf T}.
\]
A direct \(2\times2\) calculation yields
\[
 \det K=-ab(u_1v_2-u_2v_1)^2.                                 \tag{21.5}
\]
If \(t_1\neq t_2\), then
\[
 \frac{u_1}{u_2}=e^{-\omega(t_1-t_2)}
 \neq e^{-\Omega(t_1-t_2)}=\frac{v_1}{v_2},
\]
so the square in (21.5) is nonzero.  Thus \(\det K<0\), and a real
symmetric matrix with negative determinant has a negative eigenvalue.
\end{proof}

\begin{corollary}[Smooth compactly supported witness]
\label{cor:v19v18-smooth}
There is a real \(f\in C_c^\infty((0,\infty))\) such that
\[
 \int_0^\infty\!\!\int_0^\infty
 f(s)C_\alpha(s+t)f(t)\,ds\,dt<0.                             \tag{21.6}
\]
\end{corollary}

\begin{proof}
Choose a vector \(c\in\mathbb R^2\) with \(c^{\mathsf T}Kc<0\), whose
existence follows from Theorem \ref{thm:v19v18-two-time}.  Replace the
point masses at \(t_1,t_2\) by disjoint unit-mass smooth mollifiers of
radius \(\varepsilon\), weighted by \(c_1,c_2\).  Continuity of
\(C_\alpha(s+t)\) makes the resulting quadratic form converge to
\(c^{\mathsf T}Kc<0\) as \(\varepsilon\downarrow0\).  It is therefore
negative for sufficiently small \(\varepsilon\).
\end{proof}

A particularly transparent coefficient vector cancels the positive slow
mode:
\[
 c=(u_2,-u_1).
\]
Then
\[
 c\cdot u=0,\qquad
 c^{\mathsf T}Kc=-b(c\cdot v)^2<0.                            \tag{21.7}
\]

\subsection{Positive spectral measures cannot give \(p^{-4}\) decay}

The negative pole residue in (21.1) is forced by the desired ultraviolet
improvement.

\begin{theorem}[Stieltjes ultraviolet obstruction]
\label{thm:v19v18-stieltjes}
Let \(\rho\) be a nonzero positive measure on \([0,\infty)\) and suppose
\[
 G(x)=\int_{[0,\infty)}\frac{d\rho(s)}{x+s}<\infty,
 \qquad x>0.                                                   \tag{21.8}
\]
Then
\[
 xG(x)=\int\frac{x}{x+s}\,d\rho(s)
 \uparrow \rho([0,\infty])\in(0,\infty]
 \quad\hbox{as }x\to\infty.                                   \tag{21.9}
\]
In particular, a nonzero covariance with \(G(x)=O(x^{-1-\varepsilon})\)
for some \(\varepsilon>0\) has no representation (21.8) by a positive
spectral measure.
\end{theorem}

\begin{proof}
For each \(s\geq0\), the function \(x/(x+s)\) increases to one.
The monotone convergence theorem gives (21.9).  If
\(G(x)=O(x^{-1-\varepsilon})\), then \(xG(x)\to0\), contradicting
(21.9) for nonzero \(\rho\).
\end{proof}

\begin{corollary}[No positive K\"allén--Lehmann weight]
\label{cor:v19v18-kl}
The covariance (21.1), which is asymptotic to
\(\alpha^{-1}p_0^{-4}\), cannot have a nonzero positive
K\"allén--Lehmann/Stieltjes spectral measure.  Its negative residue is not
a removable partial-fraction convention.
\end{corollary}

\begin{proof}
Apply Theorem \ref{thm:v19v18-stieltjes} with
\(x=p_0^2+\omega^2\).  Alternatively, the second pole in (21.1) has
coefficient \(-1\), whereas every atom of a positive spectral measure has
nonnegative mass.
\end{proof}

\subsection{Executable certificate}

The verifier
\[
 \texttt{scripts/verify\_sm\_v19\_v18\_higher\_derivative\_reflection\_obstruction.py}
                                                                  \tag{21.10}
\]
uses the exact decays \(u_t=2^{-t}\), \(v_t=3^{-t}\) at times one and two.
It obtains
\[
 \det(uu^{\mathsf T}-vv^{\mathsf T})=-\frac1{1296}
\]
and the same negative value from the slow-mode cancellation vector.
It also checks (21.1) at three rational momenta and samples the monotone
quantity \(xG(x)\) for a positive two-atom Stieltjes measure.  It emits
\[
 \texttt{artifacts/sm\_v19\_v18\_higher\_derivative\_reflection\_obstruction.json}.
                                                                  \tag{21.11}
\]
The hashes are
\[
\begin{array}{c|l}
\hbox{object}&\hbox{SHA--256}\\ \hline
\hbox{verifier}&
\texttt{D2AB4CFF02C1599D67DCDABF0D5F52CD2584DF61798440067F2EA2993285B1EE}\\
\hbox{canonical payload}&
\texttt{7D3EE52817BC85D714299D9285B3F33EB51299C7040BAC71F1405FBFB0BEA5DA}.
\end{array}                                                     \tag{21.12}
\]

\begin{firewall}[The scalar counterexample has a precise scope]
\label{fire:v19v18}
Theorem \ref{thm:v19v18-two-time} rules out the standard factorized
higher-derivative Gaussian covariance as an Osterwalder--Schrader-positive
field.  It does not prove that every \(R^2\) gravitational formulation
fails after constraints, gauge reduction, auxiliary changes of variables,
or restriction to physical observables.  It does prove that the naive
positive Gaussian transfer cannot be inserted into V16.
\end{firewall}

\begin{decision}[V19 acquisition]
\label{dec:v19v18-next}
V19 will return to a second-order, reflection-positive lattice metric
reference and ask how much of the Einstein interaction can be included on
a bounded conformal sector.  A local nonnegative cutoff can be split
between the two time halves and preserves the transfer square; the note
will derive the exact relative-form margin and keep removal of that cutoff
as a separate tail problem.
\end{decision}

\begin{assessment}[V18 acquisition]
\label{ass:v19v18}
The straightforward higher-derivative Gaussian repair is incompatible
with reflection positivity: two positive times already give a negative
quadratic form.  More generally, positive spectral mixtures cannot improve
a nonzero scalar covariance from \(p^{-2}\) to \(p^{-4}\).  A viable metric
reference must retain a second-order positive transfer or use a
non-Gaussian construction with a new positivity proof.
\end{assessment}

\begin{record}[V18 append record]
\label{rec:v19v18}
V18 is appended to the validated V17 whose installed SHA--256 was
\texttt{8C51C43B04AE6FDCE07589855437FCE5508F2A012F8C4C22DC4DF042E6BF56DC}.
After installation only V18 is the active numeric SM note; frozen archives
and unrelated notes are unchanged.
\end{record}




\section{V19: a sharp reflection-positive conformal stiffness threshold}
\label{sec:v19v19}

The archived programme already introduced positive Gaussian gravitational
references and a strong-mismatch branch for conformal response.  V18 now
shows why replacing that reference by a \(p^{-4}\) higher-derivative
Gaussian is not reflection positive.  We return to the second-order
reference and determine exactly when it overcomes the Euclidean
Einstein conformal kinetic sign at a single reflected edge.

\subsection{The scale-factor variable linearizes the conformal gradient}

Let \(\bar g\) be a fixed smooth metric on a closed four-manifold and write
\[
 g=e^{2\phi}\bar g=a^2\bar g,\qquad a=e^\phi>0.
\]
The retained conformal transformation identity is
\[
 \int_M R_g\,dV_g
 =
 \int_M\left(\bar R\,a^2+6|\bar\nabla a|^2\right)dV_{\bar g}.
                                                               \tag{22.1}
\]
Thus, with the Euclidean sign used in V17,
\[
 S_{\rm EH}[a]
 =
 -6\kappa\int_M|\bar\nabla a|^2\,dV_{\bar g}
 -\kappa\int_M\bar R\,a^2\,dV_{\bar g}
 +2\kappa\Lambda\int_Ma^4\,dV_{\bar g}.                        \tag{22.2}
\]
Equation (22.1) is used here as retained input rather than claimed as a
new conformal-instability result.

On a finite reflection-symmetric weighted graph, add a positive Gaussian
reference stiffness \(\gamma\geq0\).  The conformal action becomes
\[
 S_{\gamma}(a)
 =
 \frac{\tau}{2}\sum_{e=\{x,y\}}w_e(a_x-a_y)^2
 +\sum_xv_x\left(2\kappa\Lambda a_x^4-\kappa\bar R_xa_x^2\right),
 \qquad
 \tau:=\gamma-12\kappa.                                       \tag{22.3}
\]
Here \(w_e,v_x>0\) and \(a_x\in(0,\infty)\).

\subsection{Necessary and sufficient edge positivity}

\begin{theorem}[Sharp single-edge threshold]
\label{thm:v19v19-edge}
For \(w>0\), the reflected edge kernel
\[
 K_{\tau,w}(s,t)
 :=\exp\!\left[-\frac{\tau w}{2}(s-t)^2\right],
 \qquad s,t>0,                                                 \tag{22.4}
\]
is positive definite if and only if
\[
 \tau\geq0,\qquad\hbox{equivalently}\qquad\gamma\geq12\kappa. \tag{22.5}
\]
\end{theorem}

\begin{proof}
If \(\tau\geq0\), then
\[
 K_{\tau,w}(s,t)
 =
 e^{-\tau ws^2/2}e^{-\tau wt^2/2}e^{\tau wst}.                 \tag{22.6}
\]
The final factor is the Fock Gram kernel of Lemma
\ref{lem:v19v6-fock}; diagonal dressing preserves positivity by
Corollary \ref{cor:v19v6-dressing}.  This includes \(\tau=0\), when the
kernel is constant.

If \(\tau<0\), choose distinct \(s,t>0\).  The evaluation matrix is
\[
 \begin{pmatrix}
 1&\exp\!\left(\frac{|\tau|w}{2}(s-t)^2\right)\\
 \exp\!\left(\frac{|\tau|w}{2}(s-t)^2\right)&1
 \end{pmatrix}.
\]
Its determinant is
\[
 1-\exp\!\left(|\tau|w(s-t)^2\right)<0,                       \tag{22.7}
\]
so the kernel is not positive definite.
\end{proof}

The threshold is sharper than an action lower bound: it is an exact
two-point condition for the reflected transfer kernel.

\subsection{Quartic confinement of the onsite factor}

\begin{lemma}[Onsite lower bound]
\label{lem:v19v19-onsite}
Assume \(\Lambda>0\).  For \(a\geq0\),
\[
 2\kappa\Lambda a^4-\kappa\bar R\,a^2
 \geq-\frac{\kappa(\bar R_+)^2}{8\Lambda},
 \qquad \bar R_+:=\max\{\bar R,0\}.                            \tag{22.8}
\]
Consequently,
\[
 \int_0^\infty
 \exp\!\left[-v(2\kappa\Lambda a^4-\kappa\bar R a^2)\right]da
 <\infty
 \quad(v>0).                                                   \tag{22.9}
\]
\end{lemma}

\begin{proof}
Put \(z=a^2\geq0\).  The quadratic
\(2\kappa\Lambda z^2-\kappa\bar Rz\) has its constrained minimum at
\(z=\bar R_+/(4\Lambda)\), with value (22.8).  Its positive leading
coefficient gives quartic decay in \(a\), proving (22.9).
\end{proof}

\begin{theorem}[Finite conformal-sector reflection positivity]
\label{thm:v19v19-finite}
Let the graph, weights, background curvature, and reflection be invariant
under time reflection.  If
\[
 \Lambda>0,\qquad\gamma\geq12\kappa,                           \tag{22.10}
\]
then the finite measure with density \(e^{-S_\gamma(a)}\prod_xda_x\) is
normalizable and reflection positive.  For every bounded observable \(F\)
supported in the positive half,
\[
 \int \overline{F(\theta a)}F(a)e^{-S_\gamma(a)}
 \prod_xda_x\geq0.                                             \tag{22.11}
\]
\end{theorem}

\begin{proof}
Normalizability follows from Lemma \ref{lem:v19v19-onsite}: the nonnegative
edge action may be discarded when bounding the partition function above,
and the remaining finite product of one-site integrals is finite.

Every edge crossing the reflection plane contributes the positive kernel
(22.4) by Theorem \ref{thm:v19v19-edge}.  Products of these kernels are
positive by tensoring their Fock feature maps.  Edges and onsite terms
strictly in the two open halves give conjugate one-side amplitudes.
Reflection-plane onsite factors are nonnegative and split into equal
square roots.  Integrating either half produces a boundary amplitude, and
the reflected integral is its squared norm in the tensor feature space.
This proves (22.11), first for compactly supported cylinder functions and
then for bounded \(F\) by dominated approximation.
\end{proof}

\begin{remark}[Relation to the hybrid Standard Model transfer]
\label{rem:v19v19-hybrid}
The conformal transfer in Theorem \ref{thm:v19v19-finite} can be tensored
with the conditional gauge--Higgs--fermion transfer of V16.  This produces
a finite reflection-positive conformal-sector model when the chiral
patching assumptions of V16 also hold.  It still fixes the background
conformal class and does not integrate transverse metric, lapse, shift, or
diffeomorphism-orbit variables.
\end{remark}

\subsection{The positive stiffness cannot be removed in this family}

\begin{proposition}[Nonremovable threshold]
\label{prop:v19v19-nonremoval}
Let \(\kappa>0\).  Any real path \(\gamma_\varepsilon\to0\) crosses the
region \(\gamma_\varepsilon<12\kappa\).  Hence the additive Gaussian
stabilizer in (22.3) cannot be removed while preserving single-edge
reflection positivity.
\end{proposition}

\begin{proof}
For sufficiently small \(\varepsilon\),
\(\gamma_\varepsilon<12\kappa\), so
\(\tau_\varepsilon=\gamma_\varepsilon-12\kappa<0\).
Theorem \ref{thm:v19v19-edge} then supplies a negative two-point
determinant.
\end{proof}

This is a no-go result for the specific real additive stiffness family.
It does not exclude a different metric variable, contour, constraint
reduction, or non-Gaussian transfer construction.

\subsection{Executable certificate}

The verifier
\[
 \texttt{scripts/verify\_sm\_v19\_v19\_conformal\_scale\_transfer\_threshold.py}
                                                                  \tag{22.12}
\]
checks the threshold arithmetic at \(\kappa=2\), an exact degree-six Fock
Gram decomposition in the positive region, a negative-stiffness
two-point matrix with determinant \(-3\), the quartic onsite minimum
\(-4/3\), and loss of positivity along a sequence
\(\gamma=30,24,18,12,6,0\).  It emits
\[
 \texttt{artifacts/sm\_v19\_v19\_conformal\_scale\_transfer\_threshold.json}.
                                                                  \tag{22.13}
\]
The hashes are
\[
\begin{array}{c|l}
\hbox{object}&\hbox{SHA--256}\\ \hline
\hbox{verifier}&
\texttt{9B3E0C989DDEFD79A790D81659BA8FD6C03C6235E676A74302278F79CE66E5B7}\\
\hbox{canonical payload}&
\texttt{EB4082B5FCD97F635DAFAAC88692755EFBB969C0F6D00ACCF45B3B0540C27465}.
\end{array}                                                     \tag{22.14}
\]

\begin{firewall}[The strong-stiffness branch is not physical closure]
\label{fire:v19v19}
Theorem \ref{thm:v19v19-finite} proves a finite reflection-positive
conformal scalar sector only for \(\gamma\geq12\kappa\).  The bare
Einstein value \(\gamma=0\) lies in the excluded region.  No transverse
metric measure, constraint algebra, Lorentzian continuation, regulator
removal, or coupled continuum follows.
\end{firewall}

\begin{decision}[V20 checkpoint]
\label{dec:v19v19-next}
V20 is the mandatory companion audit.  After it, V21 should examine
whether exact lapse/shift or Hamiltonian constraints project out the
negative conformal transfer direction.  A constraint may repair the gate
only if its projected boundary kernel is proved positive; formal gauge
language is insufficient.
\end{decision}

\begin{assessment}[V19 acquisition]
\label{ass:v19v19}
The archived strong-mismatch idea now has a sharp reflection-positivity
meaning.  A second-order Gaussian stiffness makes the conformal scale
sector finite and reflection positive exactly when
\(\gamma\geq12\kappa\).  The same exact edge test proves that this
stiffness cannot be removed within the additive real Gaussian family.
\end{assessment}

\begin{record}[V19 append record]
\label{rec:v19v19}
V19 is appended to the validated V18 whose installed SHA--256 was
\texttt{4997AB2009F6D89FF082BF1285E25148A674129E2FFBF54FD5B8EC30C20CDEAD}.
After installation only V19 is the active numeric SM note; frozen archives
and unrelated notes are unchanged.
\end{record}




\section{V20: mandatory companion audit at the metric constraint gate}
\label{sec:v19v20}

The V19 threshold is a finite conformal transfer result.  Before testing
constraint projection, the five-version rule requires another read-only
audit of the active companion notes.

\begin{audit}[Unchanged companion files]
\label{aud:v19v20-files}
The active Yang--Mills note remains
\[
 \texttt{Renewal\_Yang\_Mills\_Mass\_Gap\_Research\_Notes\_V25.tex},
 \qquad
 \operatorname{SHA256}=
 \texttt{5D6759B911E076304A17D09955E6777C3ACD6C99D3B47CD42BBA75EB0EA96BE6}.
                                                                  \tag{23.1}
\]
The active Navier--Stokes note remains
\[
 \texttt{Renewal\_NS\_Stability\_Research\_Notes\_V26.tex},
 \qquad
 \operatorname{SHA256}=
 \texttt{82C887F8C2C69E4F28FAD7C3DC8DE5B9099CB5A4BF431EBA1FFF3E91935978AD}.
                                                                  \tag{23.2}
\]
Both hashes agree with the V15 audit.
\end{audit}

\begin{proposition}[No imported constraint theorem]
\label{prop:v19v20-no-import}
Neither companion note contains a lapse, shift, Hamiltonian constraint,
metric boundary projection, or theorem that removes the V19 negative
conformal edge mode.  V21 should proceed without importing a companion
estimate.
\end{proposition}

\begin{proof}
Yang--Mills V25 remains at a finite Laurent/Haar computation and its later
exact-link routing constraint.  Navier--Stokes V26 remains at pointwise
rank-one Oseen derivative norms.  Their unchanged final programmes contain
no operator on the gravitational boundary space and no negative spectral
projection of the conformal kernel.
\end{proof}

\begin{decision}[V21 acquisition]
\label{dec:v19v20-next}
V21 will prove the exact compression criterion for an invariant constraint
projection.  If the projection commutes with the transfer kernel,
positivity is restored precisely when it annihilates the negative spectral
subspace.  The V19 two-point mode will give a concrete test of what a
proposed gravitational constraint must remove.
\end{decision}

\begin{firewall}[Audit boundary]
\label{fire:v19v20}
The unchanged companion state neither repairs nor worsens the metric
positivity obstruction.  It only fixes the next nonduplicative step.
\end{firewall}

\begin{record}[V20 append record]
\label{rec:v19v20}
V20 is appended to the validated V19 whose installed SHA--256 was
\texttt{D1D46802A9F9FDB6818593EB6F8D7B6B13A39F7CC3C15EF9DD70B159033C8909}.
After installation only V20 is the active numeric SM note.
\end{record}



\section{V21: when a constraint projection removes a negative transfer mode}
\label{sec:v19v21}

V19 exhibited a negative conformal edge mode below the stiffness threshold.
V20 found no companion constraint theorem.  This note gives the exact
finite operator criterion.  For a genuine symmetry constraint whose
orthogonal projector commutes with the transfer operator, positivity is
recovered precisely when every negative spectral vector is removed.

\subsection{The commuting compression criterion}

Let \(K=K^*\) be a bounded operator on a Hilbert space \({\cal H}\), let
\[
 E_-:=\mathbf 1_{(-\infty,0)}(K)                               \tag{24.1}
\]
be its negative spectral projection, and let \(P=P^*=P^2\) be a constraint
projection.

\begin{theorem}[Exact invariant-constraint criterion]
\label{thm:v19v21-compression}
Assume
\[
 [P,K]=0.                                                       \tag{24.2}
\]
Then the compression of \(K\) to the constrained space is positive,
\[
 PKP\geq0,                                                      \tag{24.3}
\]
if and only if
\[
 PE_-=0.                                                        \tag{24.4}
\]
\end{theorem}

\begin{proof}
Commutation with \(K\) implies commutation with every bounded Borel
function of \(K\), in particular with \(E_-\).  If (24.4) holds, then
\(\operatorname{Ran}P\) lies in the zero and positive spectral subspaces of
\(K\), so
\[
 \langle v,Kv\rangle\geq0\qquad(v\in\operatorname{Ran}P),
\]
which is (24.3).

Conversely, if \(PE_-\neq0\), choose \(u\) with
\(v:=PE_-u\neq0\).  Since \(P\) and \(E_-\) commute,
\(v\in\operatorname{Ran}P\cap\operatorname{Ran}E_-\).  The spectral
measure of \(v\) is supported in \((-\infty,0)\), hence
\[
 \langle v,Kv\rangle
 =\int_{(-\infty,0)}\lambda\,d\mu_v(\lambda)<0.                \tag{24.5}
\]
Thus \(PKP\) is not positive.
\end{proof}

\begin{corollary}[Compact group averaging]
\label{cor:v19v21-group}
Let a compact group \(G\) act unitarily by \(U_g\), suppose
\(U_gK=KU_g\), and let
\[
 P_G:=\int_GU_g\,dg                                             \tag{24.6}
\]
be Haar averaging.  The invariant transfer \(P_GKP_G\) is positive if and
only if the negative spectral subspace of \(K\) contains no nonzero
\(G\)-invariant vector.
\end{corollary}

\begin{proof}
Haar averaging is the orthogonal projection onto the invariant subspace
and commutes with \(K\).  Condition \(P_GE_-=0\) says exactly that no
negative spectral vector is invariant.  Apply Theorem
\ref{thm:v19v21-compression}.
\end{proof}

\subsection{The V19 two-point mode}

For \(q>0\), the negative-stiffness edge kernel evaluated at two distinct
field values has the form
\[
 K_q=
 \begin{pmatrix}1&e^q\\e^q&1\end{pmatrix}.                    \tag{24.7}
\]
Let
\[
 P_+=\frac12\begin{pmatrix}1&1\\1&1\end{pmatrix},
 \qquad
 P_-=\frac12\begin{pmatrix}1&-1\\-1&1\end{pmatrix}.             \tag{24.8}
\]

\begin{proposition}[The constraint must kill the difference mode]
\label{prop:v19v21-two-point}
The spectral resolution is
\[
 K_q=(1+e^q)P_+ +(1-e^q)P_-.                                  \tag{24.9}
\]
Thus \(P_+K_qP_+\geq0\), while
\[
 P_-K_qP_-=(1-e^q)P_-<0.                                      \tag{24.10}
\]
Equivalently, for the difference vector \(d=(1,-1)\),
\[
 d^{\mathsf T}K_qd=2(1-e^q)<0.                                \tag{24.11}
\]
\end{proposition}

\begin{proof}
The vectors \((1,1)\) and \((1,-1)\) are eigenvectors with eigenvalues
\(1+e^q\) and \(1-e^q\).  Equations (24.9)--(24.11) follow.
\end{proof}

\begin{corollary}[Fixed configurations defeat a group constraint]
\label{cor:v19v21-fixed}
On a finite evaluation space, suppose a constraint group fixes two
distinct conformal configurations individually.  Then their evaluation
difference is invariant.  For the negative-stiffness kernel, that
invariant vector has the negative form (24.11), so group averaging does
not restore positivity.
\end{corollary}

\begin{proof}
Individual invariance of the two evaluation vectors makes their
difference invariant.  Proposition \ref{prop:v19v21-two-point} makes its
quadratic form negative, and Corollary \ref{cor:v19v21-group} applies.
\end{proof}

If instead a two-element group exchanges the two evaluation vectors, its
average is \(P_+\) and it removes \(P_-\).  This finite example shows what
must be proved about an actual gravitational constraint; it does not
identify that swap with a lapse, shift, or Hamiltonian gauge orbit.

\subsection{Why commutation is a material hypothesis}

\begin{proposition}[Noncommuting projections can mix signs]
\label{prop:v19v21-noncommuting}
Let
\[
 K_0=\begin{pmatrix}-1&0\\0&3\end{pmatrix},
 \qquad
 P=\frac12\begin{pmatrix}1&1\\1&1\end{pmatrix}.
\]
Then \(P\) overlaps the negative eigenspace of \(K_0\), but
\[
 PK_0P=P\geq0,\qquad [P,K_0]\neq0.                             \tag{24.12}
\]
Therefore (24.4) is not a necessary condition without (24.2).
\end{proposition}

\begin{proof}
Direct multiplication gives (24.12).  Also
\(P(1,0)^{\mathsf T}\neq0\), so the range of \(P\) overlaps the negative
eigenline.
\end{proof}

For a transfer construction, such accidental mixing is insufficient by
itself.  One must also prove that the constrained space is preserved by
time evolution; that invariance is exactly the reason condition (24.2)
is natural.

\begin{criterion}[Gravitational constraint gate]
\label{crit:v19v21-gravity}
A proposed finite lapse, shift, diffeomorphism, or Hamiltonian constraint
repairs the V19 conformal transfer only after proving all of the following:
\begin{enumerate}
 \item it defines a bounded orthogonal boundary projection \(P\);
 \item \(P\) commutes with the regulated transfer \(K\);
 \item \(PE_-=0\) for the full negative spectral projection, not only for
       one sampled two-point matrix;
 \item the projected observables retain the desired physical metric
       degrees of freedom.
\end{enumerate}
\end{criterion}

\begin{proof}
The first two items make Theorem \ref{thm:v19v21-compression} applicable.
The third is its necessary and sufficient positivity condition.  The
fourth is necessary because the projection onto constants would trivially
remove every difference mode while also removing conformal dynamics.
\end{proof}

\subsection{Executable certificate}

The verifier
\[
 \texttt{scripts/verify\_sm\_v19\_v21\_constraint\_compression.py}
                                                                  \tag{24.13}
\]
uses exact rational matrices.  It decomposes
\(\left(\begin{smallmatrix}1&2\\2&1\end{smallmatrix}\right)
=3P_+-P_-\), checks the positive and negative compressions, realizes
\(P_+\) as a two-element swap average, obtains \(-2\) on the difference
vector, and verifies the noncommuting example (24.12).  It emits
\[
 \texttt{artifacts/sm\_v19\_v21\_constraint\_compression.json}. \tag{24.14}
\]
The hashes are
\[
\begin{array}{c|l}
\hbox{object}&\hbox{SHA--256}\\ \hline
\hbox{verifier}&
\texttt{C6B641FF43BF8A9B5F2DEBD49BC2DEF5F0343B376A486DDE1D9F63BA137E03D7}\\
\hbox{canonical payload}&
\texttt{0C7F776BAC374DD3FAD3F1649C7EC4D28555053DEC1E524F45FCA56E214B06AB}.
\end{array}                                                     \tag{24.15}
\]

\begin{firewall}[Formal constraints are not a positive projector]
\label{fire:v19v21}
The Wheeler--DeWitt equation or a formal lapse integral does not by itself
supply the bounded \(P\) in Criterion \ref{crit:v19v21-gravity}.
Noncompact group averaging may be distributional, the constraint range
may fail to be closed, and the negative conformal mode may be invariant.
None of these analytic issues has been resolved here.
\end{firewall}

\begin{decision}[V22 acquisition]
\label{dec:v19v21-next}
V22 will test the simplest genuine restriction already present in the
archived conformal argument: fixed volume.  Linearization makes it a
weighted mean-zero constraint.  The note will prove that this removes the
uniform conformal mode but retains every nonconstant graph-Laplacian mode,
on which the Euclidean Einstein kinetic quadratic form is strictly
negative.
\end{decision}

\begin{assessment}[V21 acquisition]
\label{ass:v19v21}
Constraint projection can repair an indefinite transfer, but only through
a concrete spectral mechanism.  For commuting symmetry constraints the
criterion is exact: every negative spectral vector must be removed.  No
lapse, shift, or Hamiltonian projector satisfying that criterion has yet
been constructed.
\end{assessment}

\begin{record}[V21 append record]
\label{rec:v19v21}
V21 is appended to the validated V20 whose installed SHA--256 was
\texttt{3C2DA46FBF98C7970BBF7B80246D82C0CA20092F869840C4F12715F8299AEF0C}.
After installation only V21 is the active numeric SM note; frozen archives
and unrelated notes are unchanged.
\end{record}



\section{V22: fixed volume retains every oscillatory conformal instability}
\label{sec:v19v22}

V21 shows that a constraint repairs transfer positivity only by eliminating
the negative spectral space.  The fixed-volume restriction used in the
archived conformal-instability theorem does the opposite at linear order:
it removes the constant scale mode and leaves the nonconstant graph
Laplacian modes, precisely where the Euclidean Einstein kinetic Hessian is
negative.

\subsection{Exact fixed-volume conformal paths}

Let \(G=(V,E)\) be a finite connected graph with edge weights \(w_e>0\)
and vertex volume weights \(v_x>0\).  Put
\[
 V_0:=\sum_xv_x,\qquad
 {\cal V}(\phi):=\sum_xv_xe^{4\phi_x}.                         \tag{25.1}
\]
The tangent space to \({\cal V}(\phi)=V_0\) at \(\phi=0\) is
\[
 H_0:=\left\{h\in\mathbb R^V:\sum_xv_xh_x=0\right\}.           \tag{25.2}
\]

\begin{lemma}[Every volume tangent has an exact constrained path]
\label{lem:v19v22-path}
For \(h\in H_0\), define
\[
 c_h(\varepsilon)
 :=-\frac14\log\left(
 \frac1{V_0}\sum_xv_xe^{4\varepsilon h_x}\right),
 \qquad
 \phi_x(\varepsilon):=\varepsilon h_x+c_h(\varepsilon).       \tag{25.3}
\]
Then
\[
 {\cal V}(\phi(\varepsilon))=V_0,\qquad
 \phi(0)=0,\qquad \phi'(0)=h.                                 \tag{25.4}
\]
\end{lemma}

\begin{proof}
Substitution of (25.3) gives
\[
 \sum_xv_xe^{4\phi_x(\varepsilon)}
 =e^{4c_h(\varepsilon)}
  \sum_xv_xe^{4\varepsilon h_x}=V_0.
\]
Also \(c_h(0)=0\), and differentiation gives
\[
 c_h'(0)=-\frac{\sum_xv_xh_x}{V_0}=0
\]
by (25.2).  Hence \(\phi'(0)=h\).
\end{proof}

\subsection{The kinetic Hessian is negative on the constraint tangent}

Write \(a_x=e^{\phi_x}\).  The discrete conformal part of the Euclidean
Einstein kinetic action is
\[
 S_{\rm kin}(\phi)
 :=-6\kappa\sum_{e=\{x,y\}}w_e
 \left(e^{\phi_x}-e^{\phi_y}\right)^2,
 \qquad \kappa>0.                                              \tag{25.5}
\]
Let \(L_G\) be the weighted graph Laplacian, so
\[
 h^{\mathsf T}L_Gh
 =\sum_{e=\{x,y\}}w_e(h_x-h_y)^2.                              \tag{25.6}
\]

\begin{theorem}[Fixed-volume negative Hessian]
\label{thm:v19v22-negative}
For every \(h\in H_0\), the exact fixed-volume path (25.3) satisfies
\[
 \left.\frac{d^2}{d\varepsilon^2}\right|_{\varepsilon=0}
 S_{\rm kin}(\phi(\varepsilon))
 =-12\kappa\,h^{\mathsf T}L_Gh.                               \tag{25.7}
\]
If \(h\neq0\), the right side is strictly negative.
\end{theorem}

\begin{proof}
The common correction \(c_h(\varepsilon)\) factors from every edge:
\[
 e^{\phi_x(\varepsilon)}-e^{\phi_y(\varepsilon)}
 =e^{c_h(\varepsilon)}
 \left(e^{\varepsilon h_x}-e^{\varepsilon h_y}\right).         \tag{25.8}
\]
At zero, the parenthesized difference vanishes and its first derivative is
\(h_x-h_y\).  Since \(c_h(0)=c_h'(0)=0\), the second derivative at zero of
the squared difference is \(2(h_x-h_y)^2\).  Summing (25.5) gives (25.7).

On a connected graph, \(h^{\mathsf T}L_Gh=0\) only for constant \(h\).
A constant vector in \(H_0\) is zero because every \(v_x>0\).  Thus every
nonzero fixed-volume tangent has a strictly negative kinetic Hessian.
\end{proof}

\begin{corollary}[Dimension of the surviving negative sector]
\label{cor:v19v22-dimension}
If \(|V|=N\geq2\), then \(H_0\) has dimension \(N-1\), and
\(-L_G\) is negative definite on \(H_0\) in the quadratic-form sense.
Fixed volume removes only the one-dimensional uniform scaling direction.
\end{corollary}

\begin{proof}
Equation (25.2) is one nonzero linear constraint, so
\(\dim H_0=N-1\).  Theorem \ref{thm:v19v22-negative} gives definiteness.
\end{proof}

\begin{corollary}[Uniform-volume projection fails the V21 test]
\label{cor:v19v22-projection}
For equal vertex weights, let
\[
 P_0=I-\frac1N{\bf1}{\bf1}^{\mathsf T}.                        \tag{25.9}
\]
Then \(P_0L_G=L_GP_0=L_G\).  For the linearized kinetic operator
\(K_{\rm lin}=-L_G\), its negative spectral projection is \(P_0\), and
\[
 P_0K_{\rm lin}P_0=-L_G<0\quad\hbox{on }\operatorname{Ran}P_0. \tag{25.10}
\]
Thus the fixed-volume projection retains, rather than annihilates, the
negative spectral subspace.
\end{corollary}

\begin{proof}
The connected graph Laplacian has kernel
\(\operatorname{span}\{{\bf1}\}\) and is positive on its orthogonal
complement.  Hence (25.9) is both the nonconstant-mode projection and the
negative spectral projection of \(-L_G\).  Equation (25.10) follows, and
Theorem \ref{thm:v19v21-compression} shows that the constraint cannot
restore positivity.
\end{proof}

For unequal \(v_x\), the Euclidean orthogonal projection onto (25.2) need
not commute with \(L_G\), but Theorem \ref{thm:v19v22-negative} still gives
strict negativity of the compressed quadratic form.

\subsection{Scope relative to the full Einstein action}

The curvature and cosmological terms in (22.2) contribute additional
onsite second variations.  On one fixed graph they can alter the sign of
low modes.  The present theorem isolates the kinetic fact needed for the
constraint question: fixed volume does not remove any nonconstant
conformal direction.  The archived continuum conformal argument uses
oscillatory modes to make this negative derivative term dominate; that
argument is not repeated here.

\subsection{Executable certificate}

The verifier
\[
 \texttt{scripts/verify\_sm\_v19\_v22\_fixed\_volume\_conformal\_mode.py}
                                                                  \tag{25.11}
\]
constructs a weighted four-vertex path, an exact basis of its weighted
mean-zero tangent space, and the restricted Laplacian.  Its leading
principal minors are
\[
 11,\quad65,\quad216,
\]
so the restriction is positive definite.  It also verifies the uniform
mean-zero projector and a two-site Einstein kinetic second variation
equal to \(-96\) at \(\kappa=2\).  It emits
\[
 \texttt{artifacts/sm\_v19\_v22\_fixed\_volume\_conformal\_mode.json}.
                                                                  \tag{25.12}
\]
The hashes are
\[
\begin{array}{c|l}
\hbox{object}&\hbox{SHA--256}\\ \hline
\hbox{verifier}&
\texttt{820ADF70D9D80936863A6841825CB0BA1663502B01415AA96B63C8D217D60C65}\\
\hbox{canonical payload}&
\texttt{13BF5C2FE034CEC6DD94845A295740EC58C16B248A9D095AEC6395302895E21C}.
\end{array}                                                     \tag{25.13}
\]

\begin{firewall}[A volume constraint is not the Hamiltonian constraint]
\label{fire:v19v22}
Theorem \ref{thm:v19v22-negative} rules out fixed volume as the missing
projection in V21.  It does not analyze the full nonlinear ADM Hamiltonian
and momentum constraints, their measure, or their quantum operator
domains.  It also does not prove that all onsite contributions are
negative on a fixed finite graph.
\end{firewall}

\begin{decision}[V23 acquisition]
\label{dec:v19v22-next}
V23 will go upstream to the linearized ADM constraint at nonzero spatial
momentum.  It will count the scalar, vector, and tensor components,
show exactly whether the Hamiltonian constraint and spatial
diffeomorphism quotient eliminate the pure conformal mode, and test
reflection positivity on the surviving transverse--traceless sector.
\end{decision}

\begin{assessment}[V22 acquisition]
\label{ass:v19v22}
The simplest gravitational restriction now has a decisive answer.  Fixed
volume removes the uniform scale direction but retains an
\((N-1)\)-dimensional negative kinetic sector on every connected
\(N\)-vertex graph.  Any successful repair must use the actual
gravitational constraints rather than volume fixing.
\end{assessment}

\begin{record}[V22 append record]
\label{rec:v19v22}
V22 is appended to the validated V21 whose installed SHA--256 was
\texttt{B44B6C47085AD13C6698BA6D343099CF15BAB2993093FD834CD625ED5FA9380E}.
After installation only V22 is the active numeric SM note; frozen archives
and unrelated notes are unchanged.
\end{record}



\section{V23: nonzero-momentum ADM reduction and positive TT transfer}
\label{sec:v19v23}

V22 proved that fixed volume retains every oscillatory conformal direction.
That restriction is not the linearized Hamiltonian constraint.  This note
therefore goes upstream to the canonical constraint surface of Einstein
gravity about flat spacetime.  At every nonzero spatial momentum, the
constraint quotient removes the scalar and vector components exactly and
leaves two transverse--traceless (TT) canonical pairs.  Their free
Euclidean covariance is then reflection positive.

\subsection{The configuration constraint quotient}

Fix a nonzero spatial covector (k\in\mathbb R^3), put

\[
 \omega:=|k|>0,
 \qquad Q_k:=\operatorname{Sym}^2(\mathbb R^3),                 \tag{26.1}
\]

and use the Euclidean tensor pairing
(\langle p,h\rangle=p^{ij}h_{ij}).  The linearized Hamiltonian
constraint and the spatial-diffeomorphism gauge map at (k) are

\[
 C_k(h):=k_ik_jh_{ij}-\omega^2 h_i{}^i,
 \qquad
 (D_k\xi)_{ij}:=k_i\xi_j+k_j\xi_i.                            \tag{26.2}
\]

The harmless Fourier factor (i) has been suppressed; equivalently one
may work separately with the real sine and cosine modes.

\begin{theorem}[Exact nonzero-momentum configuration quotient]
\label{thm:v19v23-config}
For every (k\neq0),

\[
 C_kD_k=0,
 \qquad
 \dim\bigl(\ker C_k/\operatorname{Ran}D_k\bigr)=2.             \tag{26.3}
\]

Every class has a unique TT representative satisfying

\[
 k^ih_{ij}=0,
 \qquad h_i{}^i=0.                                            \tag{26.4}
\]

After a rotation taking (k) to ((0,0,\omega)), that representative is

\[
 h=q_+e^++q_\times e^\times,
 \quad
 e^+=\begin{pmatrix}1&0&0\\0&-1&0\\0&0&0\end{pmatrix},
 \quad
 e^\times=\begin{pmatrix}0&1&0\\1&0&0\\0&0&0\end{pmatrix}. \tag{26.5}
\]
\end{theorem}

\begin{proof}
Direct contraction gives

\[
 C_k(D_k\xi)
 =2\omega^2(k\cdot\xi)-\omega^2,2(k\cdot\xi)=0.             \tag{26.6}
\]

Thus the quotient in (26.3) is defined.  Rotate coordinates so that
(k=(0,0,\omega)).  Then

\[
 C_k(h)=-\omega^2(h_{11}+h_{22}),                              \tag{26.7}
\]

whereas \(D_k\xi\) changes \(h_{13},h_{23},h_{33}\) by
\(\omega\xi_1,\omega\xi_2,2\omega\xi_3\), respectively, and changes
none of \(h_{11},h_{22},h_{12}\).  Since \(\omega>0\), every gauge class
has a unique representative with
\(h_{13}=h_{23}=h_{33}=0\).  Equation (26.7) then says
\(h_{22}=-h_{11}\), leaving precisely the two coefficients in (26.5).
Those matrices obey (26.4).  This also proves the dimension statement.
\end{proof}

\begin{corollary}[The pure conformal mode is absent]
\label{cor:v19v23-conformal}
For a pure trace perturbation (h_{ij}=2\varphi\delta_{ij}),

\[
 C_k(h)=-4\omega^2\varphi.                                    \tag{26.8}
\]

Consequently a nonzero-momentum pure conformal perturbation satisfies the
linearized Hamiltonian constraint only when \(\varphi=0\).
\end{corollary}

\begin{proof}
Substitution into (26.2) gives
\(2\varphi\omega^2-6\varphi\omega^2=-4\varphi\omega^2\).
The conclusion follows from \(\omega>0\).
\end{proof}

\subsection{Momentum constraint, normal gauge, and the symplectic quotient}

For (p\in Q_k), define the linearized momentum constraint and the gauge
direction generated by the Hamiltonian constraint by

\[
 (M_kp)_i:=k_jp_{ij},
 \qquad
 (N_kn)_{ij}:=n(k_ik_j-\omega^2\delta_{ij}).                   \tag{26.9}
\]

An overall nonzero normalization or sign in (N_k) has no effect on its
range.

\begin{theorem}[Reduced linearized ADM phase space]
\label{thm:v19v23-phase}
For every (k\neq0),

\[
 N_k\mathbb R\subset\ker M_k,
 \qquad
 \dim\bigl(\ker M_k/N_k\mathbb R\bigr)=2.                    \tag{26.10}
\]

The canonical tensor pairing descends to a nondegenerate pairing

\[
 \frac{\ker M_k}{N_k\mathbb R}
 \ \times\ 
 \frac{\ker C_k}{\operatorname{Ran}D_k},                     \tag{26.11}
\]

and both factors have the TT representatives (26.5).  Hence the reduced
phase space at (k) has dimension four: two configuration amplitudes and
their two conjugate momenta.
\end{theorem}

\begin{proof}
First,

\[
 k_j(N_kn)_{ij}
 =n(k_i\omega^2-\omega^2k_i)=0.                              \tag{26.12}
\]

For (k=(0,0,\omega)), the three independent equations (M_kp=0)
are (p_{13}=p_{23}=p_{33}=0).  On this three-dimensional kernel,
(N_kn) is proportional to
(\operatorname{diag}(1,1,0)).  Quotienting by this transverse-trace
direction leaves (p_{22}=-p_{11}) and (p_{12}), so (26.10) and the TT
representatives follow.

It remains to check that (26.11) is intrinsic.  If (p\in\ker M_k), then

\[
 \langle p,D_k\xi\rangle=2\xi_i k_jp_{ij}=0.                  \tag{26.13}
\]

If (h\in\ker C_k), then

\[
 \langle N_kn,h\rangle=nC_k(h)=0.                            \tag{26.14}
\]

Thus changing either representative does not change the pairing.  In the
TT representatives it is the ordinary nondegenerate pairing of the
((+,\times)) coefficients.  This proves nondegeneracy and the phase-space
dimension.
\end{proof}

This is also the canonical count
(12-2\cdot4=4): the linearized Hamiltonian and three momentum
constraints are first class at the flat background, and each constraint
removes one constraint coordinate and one gauge coordinate.

\subsection{Reflection positivity of the reduced free covariance}

Let (h_{0\mu}=0) after solving the lapse and shift constraints.  With
the normalization (e^A_{ij}e^B_{ij}=2\delta^{AB}), substitution of the
TT field (h_{ij}=\sum_Aq_Ae^A_{ij}) into the Euclidean Fierz--Pauli
quadratic functional gives, for the spatial pair (k,-k),

\[
 S^{(2)}_{\rm TT}
 =\frac{c_g}{2}\sum_{A\in\{+,\times\}}
   \int_{\mathbb R}
   \left(|\dot q_A(t)|^2+\omega^2|q_A(t)|^2\right),dt,
 \qquad c_g>0.                                                \tag{26.15}
\]

Indeed every divergence and trace term vanishes by (26.4); the remaining
derivative square is positive and the polarization normalization gives
(26.15).  Reality is imposed by (q_A(-k)=\overline{q_A(k)}), or by using
the two real sine/cosine copies.

\begin{theorem}[TT Osterwalder--Schrader positivity]
\label{thm:v19v23-os}
The centered Gaussian defined by (26.15) has covariance

\[
 C_{AB}(t-s;k)
 =\delta_{AB}\frac{e^{-\omega|t-s|}}{2c_g\omega}.             \tag{26.16}
\]

For every pair of test functions \(f_A\) supported in \(t>0\), reflection
\(\theta t=-t\) gives

\[
 \begin{split}
 \langle f,f\rangle_\theta
 &:={\sum_A}\int_0^\infty\!\!\int_0^\infty
 \overline{f_A(t)},C_{AA}(-t-s;k),f_A(s)\,dt\,ds\\
 &=\frac1{2c_g\omega}\sum_A
 \left|\int_0^\infty e^{-\omega t}f_A(t)\,dt\right|^2
 \geq0.                                                       \tag{26.17}
 \end{split}
\]
\end{theorem}

\begin{proof}
The inverse of (c_g(-\partial_t^2+\omega^2)) on the real line is
(26.16).  For (t,s>0),
(|-t-s|=t+s), so its reflected kernel factorizes as

\[
 C_{AA}(-t-s;k)
 =\frac{e^{-\omega t}e^{-\omega s}}{2c_g\omega}.              \tag{26.18}
\]

Substitution produces the squares in (26.17).
\end{proof}

\begin{corollary}[Positive free physical transfer]
\label{cor:v19v23-transfer}
At each nonzero spatial momentum, the constraint-reduced free Einstein
sector admits the standard positive oscillator transfer semigroup with
one-particle energy \(\omega=|k|\), independently for the \(+\) and
(\times) polarizations.
\end{corollary}

\begin{proof}
Theorem \ref{thm:v19v23-os} supplies the reflection-positive Gaussian
covariance.  Its time-translation factor in the reflected kernel is
(e^{-\omega t}), hence reconstruction gives the nonnegative oscillator
Hamiltonian with energy \(\omega\) on each polarization.
\end{proof}

\subsection{Relation to the conformal instability}

There is no contradiction between V22 and Theorem
\ref{thm:v19v23-os}.  Fixed volume is a codimension-one restriction on
off-shell conformal fields and retains nonconstant scalar directions.  The
linearized ADM construction instead imposes all four gravitational
constraints and quotients all four gauge directions before forming the
physical Gaussian.  At (k\neq0), its scalar direction is absent by
Corollary \ref{cor:v19v23-conformal}.  This identifies a viable free
physical-sector route; it does not make the unreduced Euclidean
Einstein--Hilbert integral convergent.

\subsection{Executable certificate}

The exact verifier

\[
 \texttt{scripts/verify\_sm\_v19\_v23\_linearized\_adm\_tt.py} \tag{26.19}
\]

uses symmetric-tensor coordinates
((11,22,33,12,13,23)) at (k=(0,0,2)).  It checks the ranks and
annihilations in Theorems \ref{thm:v19v23-config} and
\ref{thm:v19v23-phase}, obtains the pure-trace constraint value (-8),
and constructs the reflected covariance at decay (1/2),

\[
 \begin{pmatrix}
  1/4&1/8&1/16\\
  1/8&1/16&1/32\\
  1/16&1/32&1/64
 \end{pmatrix}
 =
 \begin{pmatrix}1/2\\1/4\\1/8\end{pmatrix}
 \begin{pmatrix}1/2&1/4&1/8\end{pmatrix}.                    \tag{26.20}
\]

It emits

\[
 \texttt{artifacts/sm\_v19\_v23\_linearized\_adm\_tt.json}.  \tag{26.21}
\]

The hashes are

\[
\begin{array}{c|l}
\hbox{object}&\hbox{SHA--256}\\ \hline
\hbox{verifier}&
\texttt{F7E29DDEBD82358E851614476AA95E8F558878F4EC0B7FB43E7424AEBAE04C67}\\
\hbox{canonical payload}&
\texttt{D61C544901BED9CA3555EB25B4A598115597CC62A1BD83ED35EDC751934B7946}.
\end{array}                                                     \tag{26.22}
\]

\begin{firewall}[Free nonzero modes are not nonlinear quantum gravity]
\label{fire:v19v23}
Theorems \ref{thm:v19v23-phase} and \ref{thm:v19v23-os} concern the
linearized theory about a fixed flat background and (k\neq0).  They do
not construct a global nonlinear constraint surface or gauge slice, a
Faddeev--Popov or reduced measure, a positive interacting transfer
operator, or a background-independent continuum limit.  The spatial zero
mode is excluded: (26.7) loses rank there and (26.16) becomes infrared
singular.
\end{firewall}

\begin{decision}[V24 acquisition]
\label{dec:v19v23-next}
V24 will isolate the excluded (k=0) sector.  It will prove the precise
failure of a stationary massless Gaussian covariance, then give explicit
finite-volume boundary or mean-zero conditions under which the inverse
kinetic operator and reflection form are well defined.  This separates an
infrared normalization problem from the nonzero-mode constraint result.
\end{decision}

\begin{assessment}[V23 acquisition]
\label{ass:v19v23}
The actual linearized ADM constraints succeed where fixed volume fails:
at every nonzero spatial momentum they remove the conformal scalar and
all vector components, leaving exactly two TT canonical pairs.  The free
TT covariance has an explicit positive reflected-square factorization.
The next obstruction is the spatial zero mode and, beyond it, nonlinear
constraint reduction.
\end{assessment}

\begin{record}[V23 append record]
\label{rec:v19v23}
V23 is appended to the validated V22 whose installed SHA--256 was
\texttt{068FB4064D9DDA29D83E859E2BF5F327B611E01CC884E669BE246AC01092BE16}.
After installation only V23 is the active numeric SM note; frozen archives
and unrelated notes are unchanged.
\end{record}



\section{V24: the homogeneous gravitational mode and infrared prescriptions}
\label{sec:v19v24}

V23 deliberately assumed nonzero spatial momentum.  This note identifies
what fails at \(k=0\).  The symbols of all four linearized constraints and
all four infinitesimal gauge directions vanish there, while the TT
oscillator frequency tends to zero.  The resulting one-dimensional
kinetic action is a free particle: it has a positive heat semigroup but no
normalizable translation-invariant Gaussian vacuum on the whole time
line.  Two explicit infrared prescriptions recover a well-defined
reflection form without pretending that the zero mode has been solved
dynamically.

\subsection{Rank loss of the ADM symbol}

Retain the maps \(C_k,D_k,M_k,N_k\) from (26.2) and (26.9).

\begin{proposition}[The zero-momentum quotient is singular]
\label{prop:v19v24-rank}
At \(k=0\),

\[
 C_0=0,\qquad D_0=0,\qquad M_0=0,\qquad N_0=0.                \tag{27.1}
\]

Thus the configuration and momentum constraint maps both have rank zero,
and neither infinitesimal gauge map removes a homogeneous tensor
component.  In particular, the dimension-four quotient of Theorem
\ref{thm:v19v23-phase} does not extend continuously to \(k=0\).
\end{proposition}

\begin{proof}
Every term in (26.2) and (26.9) contains at least one factor of \(k\);
the scalar symbols contain two.  Substitution of \(k=0\) gives (27.1).
For \(k\neq0\), V23 found constraint ranks \(1\) and \(3\), and gauge
ranks \(3\) and \(1\), respectively.  All four ranks therefore jump at
zero.
\end{proof}

On a spatial torus this is a genuine finite-dimensional homogeneous
sector.  Transversality is vacuous at \(k=0\), the linear Hamiltonian
constraint does not impose tracelessness, and constant spatial
diffeomorphisms have zero derivative.  Some homogeneous perturbations
describe torus metric moduli.  They cannot be counted as two ordinary
nonzero-momentum graviton polarizations.

\subsection{No stationary Gaussian inverse for the free-particle mode}

For one real homogeneous amplitude, the quadratic action is

\[
 S_0(q)=\frac c2\int_{\mathbb R}|\dot q(t)|^2\,dt,
 \qquad c>0.                                                   \tag{27.2}
\]

\begin{theorem}[Massless zero-mode vacuum obstruction]
\label{thm:v19v24-no-vacuum}
There is no centered, time-translation-invariant Gaussian process on
\(\mathbb R\), with finite point variance, whose covariance \(K(t-s)\)
is an inverse of the kinetic operator in (27.2):

\[
 -cK''=\delta_0                                                     \tag{27.3}
\]

in the sense of distributions.
\end{theorem}

\begin{proof}
Reality and stationarity make \(K(u)\) even.  Since
\((|u|)''=2\delta_0\), every even distributional solution of (27.3)
which is a function has the form

\[
 K(u)=a-\frac{|u|}{2c}.                                        \tag{27.4}
\]

Here \(a=K(0)\geq0\).  Positivity of the covariance matrix at the two
times \(0,u\) requires

\[
 \begin{pmatrix}a&K(u)\\K(u)&a\end{pmatrix}\succeq0,
 \qquad\hbox{hence}\qquad |K(u)|\leq a.                         \tag{27.5}
\]

For sufficiently large \(|u|\), (27.4) has
\(|K(u)|>a\), contradicting (27.5).  If \(a=0\), every \(u\neq0\)
already gives the contradiction.
\end{proof}

The obstruction is normalization rather than negativity of the
free-particle Hamiltonian.  The heat semigroup
\(e^{-tP^2/(2c)}\) is positive, but its zero-energy generalized
eigenfunction is constant in \(q\) and is not square integrable.  Hence it
does not define the stationary probability measure assumed in the
Osterwalder--Schrader Gaussian construction.

\subsection{Prescription I: remove the spatial mean}

Let the spatial manifold be the flat torus
\(\mathbb T_L^3=(\mathbb R/L\mathbb Z)^3\).  Impose on every retained
linearized tensor component

\[
 \int_{\mathbb T_L^3}h_{ij}(t,x)\,dx=0.                        \tag{27.6}
\]

\begin{theorem}[Mean-zero finite-volume TT covariance]
\label{thm:v19v24-mean-zero}
Condition (27.6) removes exactly the Fourier coefficient \(k=0\).
All remaining momenta are

\[
 k_n=\frac{2\pi}{L}n,\qquad
 n\in\mathbb Z^3\setminus\{0\},\qquad
 |k_n|\geq\frac{2\pi}{L}.                                     \tag{27.7}
\]

For smooth positive-time TT test tensors \(f_A(t,n)\), the covariance
obtained by summing (26.16) over \(n\neq0\) obeys

\[
 \begin{split}
 \langle f,f\rangle_\theta
 =\sum_{n\neq0}\frac1{2c_g|k_n|}
 \sum_{A\in\{+,\times\}}
 \left|
 \int_0^\infty e^{-|k_n|t}f_A(t,n)\,dt
 \right|^2
 \geq0.                                                       \tag{27.8}
 \end{split}
\]

At a finite spatial Fourier cutoff the sum is finite; without that cutoff
it defines a positive distribution on spatially smooth test tensors.
\end{theorem}

\begin{proof}
Fourier integration in (27.6) selects the coefficient \(n=0\), proving
its removal and (27.7).  Apply Theorem \ref{thm:v19v23-os} to each
remaining momentum.  Every summand is a nonnegative square.  Smoothness
gives rapid decay of the spatial Fourier coefficients, so the
distributional sum converges on the stated test space.
\end{proof}

This prescription fixes the global spatial average.  It is appropriate
when that average is boundary data or a chosen torus modulus; it is not a
consequence of the local linearized constraints.

\subsection{Prescription II: reflection-symmetric Dirichlet time slab}

The homogeneous mode can instead be retained on the finite interval
\([-T,T]\), with \(T>0\), by imposing

\[
 q(-T)=q(T)=0.                                                 \tag{27.9}
\]

\begin{theorem}[Dirichlet zero-mode reflection positivity]
\label{thm:v19v24-dirichlet}
The operator \(c(-\partial_t^2)\) with (27.9) is strictly positive and
has Green function

\[
 G_T(t,s)=
 \frac{(T+\min\{t,s\})(T-\max\{t,s\})}{2cT}.                  \tag{27.10}
\]

It is invariant under \(t\mapsto-t\).  For every test function supported
in \(0<t<T\),

\[
 \begin{split}
 \langle f,f\rangle_{\theta,T}
 &:=\int_0^T\!\!\int_0^T
 \overline{f(t)}G_T(-t,s)f(s)\,dt\,ds\\
 &=\frac1{2cT}
 \left|\int_0^T(T-t)f(t)\,dt\right|^2\geq0.                  \tag{27.11}
 \end{split}
\]
\end{theorem}

\begin{proof}
The Poincaré inequality for \(H_0^1([-T,T])\) makes the Dirichlet
quadratic form \(c\int|\dot q|^2\) strictly positive.  The function
(27.10) vanishes at both endpoints, is continuous at \(t=s\), is affine
away from \(t=s\), and its first derivative has jump \(-1/c\).
Consequently \(c(-\partial_t^2)G_T(\,\cdot\,,s)=\delta_s\).
The formula is unchanged after simultaneously replacing \(t,s\) by
\(-t,-s\).  If \(t,s>0\), then \(-t<s\), and (27.10) gives

\[
 G_T(-t,s)=\frac{(T-t)(T-s)}{2cT}.                            \tag{27.12}
\]

Substitution yields the square in (27.11).
\end{proof}

There is no stationary infinite-slab limit.  For fixed \(t,s>0\),

\[
 G_T(-t,s)
 =\frac{T}{2c}-\frac{t+s}{2c}+\frac{ts}{2cT},                 \tag{27.13}
\]

whose first term diverges as \(T\to\infty\).  Removing that divergence by
hand changes the covariance and returns the obstruction of Theorem
\ref{thm:v19v24-no-vacuum}.

\subsection{Executable certificate}

The verifier

\[
 \texttt{scripts/verify\_sm\_v19\_v24\_zero\_mode\_ir.py}      \tag{27.14}
\]

checks the four zero ranks in (27.1).  For \(c=a=1\) and time separation
six, the only stationary inverse candidate has the two-point covariance
matrix

\[
 \begin{pmatrix}1&-2\\-2&1\end{pmatrix},
 \qquad \det=-3,                                               \tag{27.15}
\]

which certifies Theorem \ref{thm:v19v24-no-vacuum}.  For the Dirichlet
choice \(T=5,c=2\) and positive times \(1,2,3\), it checks the exact
reflected matrix

\[
 \begin{pmatrix}
 4/5&3/5&2/5\\
 3/5&9/20&3/10\\
 2/5&3/10&1/5
 \end{pmatrix}
 =\frac1{20}
 \begin{pmatrix}4\\3\\2\end{pmatrix}
 \begin{pmatrix}4&3&2\end{pmatrix}\succeq0.                  \tag{27.16}
\]

It emits

\[
 \texttt{artifacts/sm\_v19\_v24\_zero\_mode\_ir.json}.         \tag{27.17}
\]

The hashes are

\[
\begin{array}{c|l}
\hbox{object}&\hbox{SHA--256}\\ \hline
\hbox{verifier}&
\texttt{2247068391A49C15D5C1DF4BBB59363B1AFDDE8DC20A9491838F9869053E4CC0}\\
\hbox{canonical payload}&
\texttt{5B6CAB3E97265CE0CD62B4B51BC01284BB6FF2B6D26DF0A4F24F2A19442E0E6A}.
\end{array}                                                     \tag{27.18}
\]

\begin{firewall}[An infrared prescription is extra global data]
\label{fire:v19v24}
Neither (27.6) nor (27.9) follows from the local ADM constraint symbols.
The first fixes homogeneous spatial data; the second chooses a finite-time
boundary state and breaks time-translation invariance.  V24 establishes
well-defined positive free reflection forms under those choices.  It does
not select the physically correct gravitational zero-mode state, solve
the nonlinear integrated constraints, or construct an infinite-volume
interacting vacuum.
\end{firewall}

\begin{decision}[V25 companion audit and V26 acquisition]
\label{dec:v19v24-next}
V25 will perform the required Yang--Mills and Navier--Stokes companion
audit.  V26 will then examine the first nonlinear obstruction to extending
V23: the failure of a naive implicit-function argument for the compact
flat-background constraint map when its adjoint has homogeneous zero
modes.  The aim is to identify the exact quadratic compatibility
condition rather than assume that every linear TT seed integrates to an
exact constrained solution.
\end{decision}

\begin{assessment}[V24 acquisition]
\label{ass:v19v24}
The nonzero-momentum TT result now has a precise infrared boundary.  At
\(k=0\) the linearized constraint and gauge symbols all vanish, and the
free-particle kinetic operator has no stationary Gaussian inverse with
finite variance.  Spatial mean-zero data remove that sector and preserve
the V23 sum of squares; symmetric Dirichlet time boundaries retain it and
give a different explicit reflected square.
\end{assessment}

\begin{record}[V24 append record]
\label{rec:v19v24}
V24 is appended to the validated V23 whose installed SHA--256 was
\texttt{510246BD7984D23ABF49FF6A58F420CDD69C3A3416759167C1AD463DC502BC8B}.
After installation only V24 is the active numeric SM note; frozen archives
and unrelated notes are unchanged.
\end{record}



\section{V25: compulsory companion audit after the infrared split}
\label{sec:v19v25}

This is the required five-version checkpoint following V20.  The newest
active Yang--Mills and Navier--Stokes notes in the shared
\texttt{latex\_notes} folder were read without modification.  The audit
tests whether either programme changes the next Einstein acquisition
after the nonzero-mode ADM reduction of V23 and the zero-mode analysis of
V24.

\subsection{Files inspected}

\begin{record}[V25 companion snapshot]
\label{rec:v19v25-snapshot}
The files inspected were

\[
\begin{array}{c|r|r}
\hbox{file and SHA--256}&\hbox{bytes}&\hbox{lines}\\ \hline
\texttt{Renewal\_Yang\_Mills\_Mass\_Gap\_Research\_Notes\_V27.tex}
&385758&9119\\
\multicolumn{3}{l}{
\texttt{3D80B53EDC2C565509BED36604346980C6A86F4BCC42A3A560C1368C7B685EE5}}
\\
\texttt{Renewal\_NS\_Stability\_Research\_Notes\_V26.tex}
&278283&4761\\
\multicolumn{3}{l}{
\texttt{82C887F8C2C69E4F28FAD7C3DC8DE5B9099CB5A4BF431EBA1FFF3E91935978AD}}.
\end{array}                                                     \tag{28.1}
\]

At V20 the active Yang--Mills note was V25 and the Navier--Stokes note
was V26.  Thus Yang--Mills has advanced by two versions, while the
Navier--Stokes source is unchanged.
\end{record}

\subsection{Yang--Mills reading}

Yang--Mills V26 completed a signed spectral grid over the first finite
field prime with restart-safe provenance.  V27 performed the exact
separable inverse transform, proved no aliasing on its enumerated Laurent
support, found no residues outside that support, and checked a complete
forward round trip.  The result is still a one-prime coefficient residue
array; the remaining-prime CRT lift, reflection-paired inequality,
interacting measure, regulator removal, Osterwalder--Schrader
reconstruction, and mass gap remain open.

Its Proposition V27 on finite-field conjugation is methodologically
relevant: algebraic conjugation of the transform values does not reverse
momentum unless a separate coefficient-reflection identity is proved.
No finite-field coefficient or transform identity is a gravitational ADM
input.  The lesson retained here is only that V23's positivity must remain
tied to the actual time-reflection kernel (26.18), not to an unrelated
formal involution.

\subsection{Navier--Stokes reading}

Navier--Stokes V26 computes closed pointwise rank-one norms for the first
and second derivatives of its Oseen kernel.  The first-derivative
refinement yields no gain; the second yields a limited improvement over a
specified radial range.  Those estimates concern a far-past
Navier--Stokes stability ledger.  They neither alter the Einstein
constraint symbol nor supply the nonlinear compatibility condition
announced at V24.

\begin{decision}[V25 audit decision]
\label{dec:v19v25-audit}
No theorem, constant, or executable artifact is imported from either
companion.  The V26 direction is unchanged: analyze the nonlinear ADM
constraint map near a compact flat background and expose the cokernel
compatibility condition before invoking any implicit-function argument.
In carrying out later reflection tests, the physical time reflection and
its action on all fields must be proved explicitly.
\end{decision}

\begin{proof}
The Yang--Mills acquisitions are finite-field reconstruction statements
for a gauge response polynomial, and their own firewall leaves physical
reflection positivity open.  The Navier--Stokes acquisitions are
pointwise bounds for derivatives of a heat kernel.  Neither statement has
the domain, variables, or conclusion needed for an ADM constraint
theorem.  V24 already isolates the next unresolved Einstein question:
whether nonzero linear TT seeds satisfy the nonlinear integrated
constraints.  Hence no change of mathematical target is justified.
\end{proof}

\begin{assessment}[V25 acquisition]
\label{ass:v19v25}
The companion audit is current and produces no cross-program theorem.
It does confirm the present proof discipline: reflection positivity is a
property of a specified reflected kernel and field involution, while
constraint solvability requires a separate range and cokernel analysis.
\end{assessment}

\begin{record}[V25 append record]
\label{rec:v19v25}
V25 is appended to the validated V24 whose installed SHA--256 was
\texttt{F4FBEBB8ED72E8AA5A0D320FE627CF579A01A3A0ACFB21B2FC9950FEFF5E57DA}.
After installation only V25 is the active numeric SM note; the two
companions, frozen archives, and all unrelated notes are unchanged.  The
next compulsory companion audit is V30.
\end{record}



\section{V26: quadratic ADM compatibility on the compact flat background}
\label{sec:v19v26}

V23 identified two positive free TT modes at every nonzero spatial
momentum, and V24 showed that the homogeneous sector cannot simply be
ignored when one asks for a stationary vacuum.  The first nonlinear test
now joins those conclusions.  On a compact flat Cauchy surface, the
linearized scalar constraint has a constant cokernel.  Pairing the
second-order constraint with that constant produces a positive
compatibility charge.  Consequently a generic nonzero TT solution of the
linearized constraints is not tangent to a curve of exact vacuum data
through the static flat torus.

\subsection{The missing surjectivity}

Let \((\mathbb T^3,\delta)\) be a flat three-torus.  For a metric \(g\)
and second fundamental form \(K\), use the vacuum constraint convention

\[
 {\cal H}(g,K):=R(g)+(\operatorname{tr}_gK)^2-|K|_g^2=0,
 \qquad
 {\cal M}_i(g,K):=\nabla^j
 \bigl(K_{ij}-(\operatorname{tr}_gK)g_{ij}\bigr)=0.            \tag{29.1}
\]

At the time-symmetric flat datum \((\delta,0)\), the scalar part of the
linearization is

\[
 Lh:=D{\cal H}_{(\delta,0)}(h,A)
 =\partial_i\partial_jh_{ij}
  -\Delta(\operatorname{tr}_\delta h).                         \tag{29.2}
\]

\begin{proposition}[Constant-lapse cokernel]
\label{prop:v19v26-cokernel}
For every sufficiently regular periodic symmetric tensor \(h\),

\[
 \int_{\mathbb T^3}Lh\,dx=0.                                  \tag{29.3}
\]

Hence the constant function belongs to \(\ker L^*\), \(L\) is not
surjective onto any function space containing constants, and an ordinary
implicit-function theorem cannot solve the full scalar constraint near
\((\delta,0)\) without first imposing a compatibility condition or adding
a complementary variable.
\end{proposition}

\begin{proof}
Both terms in (29.2) are total second derivatives.  Their integrals vanish
by periodicity.  Equation (29.3) is precisely
\(\langle1,Lh\rangle=0\), which puts \(1\) in the formal-adjoint kernel.
Surjectivity would require the range to contain a nonzero constant, which
is impossible by (29.3).
\end{proof}

The same background also has constant-shift cokernel directions for the
momentum constraint.  V26 treats the constant-lapse compatibility charge;
the three momentum charges remain separate.

\subsection{Second variation of total scalar curvature}

Write

\[
 {\cal S}(g):=\int_{\mathbb T^3}R(g)\,dV_g.                    \tag{29.4}
\]

\begin{lemma}[Flat TT Hessian]
\label{lem:v19v26-hessian}
If \(h\) is transverse and traceless with respect to \(\delta\), then

\[
 D^2{\cal S}_\delta(h,h)
 =-\frac12\int_{\mathbb T^3}|\nabla h|^2\,dx.                 \tag{29.5}
\]

For any \(C^2\) metric curve \(g(\varepsilon)\) with
\(g(0)=\delta\) and \(g'(0)=h\), the second derivative of
\({\cal S}(g(\varepsilon))\) is still (29.5), independently of
\(g''(0)\).
\end{lemma}

\begin{proof}
On a closed manifold the first variation with respect to a covariant
metric perturbation is

\[
 D{\cal S}_g(h)=-\int\langle
 \operatorname{Ric}_g-\tfrac12R_gg,h\rangle_g\,dV_g.           \tag{29.6}
\]

It vanishes at the flat metric.  The flat linearized Ricci tensor is

\[
 (D\operatorname{Ric}_\delta h)_{ij}
 =\frac12\left(
 \partial_k\partial_i h_{jk}
 +\partial_k\partial_j h_{ik}
 -\Delta h_{ij}
 -\partial_i\partial_j\operatorname{tr}h\right).              \tag{29.7}
\]

For transverse traceless \(h\), this is
\(-\tfrac12\Delta h_{ij}\), while \(DR_\delta(h)=0\).  Differentiating
(29.6) at \(\delta\) therefore gives

\[
 D^2{\cal S}_\delta(h,h)
 =\frac12\int h_{ij}\Delta h_{ij}\,dx
 =-\frac12\int|\nabla h|^2\,dx.                               \tag{29.8}
\]

Along a curved path the chain rule adds
\(D{\cal S}_\delta(g''(0))\), which is zero because the first variation
vanishes at \(\delta\).
\end{proof}

\subsection{The constant-lapse compatibility charge}

\begin{theorem}[Vacuum quadratic compatibility identity]
\label{thm:v19v26-taub}
Let \((g(\varepsilon),K(\varepsilon))\) be a \(C^2\) family of exact
vacuum solutions of (29.1), with

\[
 g(0)=\delta,\quad K(0)=0,\quad
 h:=g'(0),\quad A:=K'(0).                                     \tag{29.9}
\]

Assume \(h\) is transverse traceless.  Write

\[
 \tau:=\operatorname{tr}_\delta A,\qquad
 A_0:=A-\frac{\tau}{3}\delta.                                 \tag{29.10}
\]

Then the necessary second-order identity is

\[
 \boxed{\quad
 \frac14\int_{\mathbb T^3}|\nabla h|^2\,dx
 +\int_{\mathbb T^3}|A_0|^2\,dx
 =\frac23\int_{\mathbb T^3}\tau^2\,dx .
 \quad}                                                       \tag{29.11}
\]
\end{theorem}

\begin{proof}
Multiply the pointwise Hamiltonian constraint by \(dV_g\) and integrate:

\[
 {\cal S}(g(\varepsilon))
 +\int_{\mathbb T^3}
 \left((\operatorname{tr}_{g(\varepsilon)}K(\varepsilon))^2
       -|K(\varepsilon)|_{g(\varepsilon)}^2\right)
 dV_{g(\varepsilon)}=0.                                      \tag{29.12}
\]

At \(\varepsilon=0\), \(K=0\).  Therefore derivatives of the metric,
inverse metric, and volume form do not contribute to the second
derivative of the terms quadratic in \(K\); only \(A=K'(0)\) remains.
Twice differentiating (29.12) and applying Lemma
\ref{lem:v19v26-hessian} gives

\[
 -\frac12\int|\nabla h|^2\,dx
 +2\int\left(\tau^2-|A|^2\right)\,dx=0.                       \tag{29.13}
\]

The orthogonal decomposition (29.10) has
\(|A|^2=|A_0|^2+\tau^2/3\).  Substitution in (29.13), followed by division
by two and rearrangement, gives (29.11).  No unknown second derivative of
the family appears, because the first variation of (29.12) at
\((\delta,0)\) vanishes.
\end{proof}

\begin{corollary}[Compact TT linearization instability]
\label{cor:v19v26-instability}
Under the hypotheses of Theorem \ref{thm:v19v26-taub}, if \(A\) is also
traceless, then

\[
 \frac14\int|\nabla h|^2\,dx+\int|A|^2\,dx=0.                 \tag{29.14}
\]

Consequently \(A=0\) and \(h\) is spatially parallel.  If \(h\) also has
zero spatial mean, as in (27.6), then \(h=0\).  Thus no nonzero-momentum TT
seed from V23, with tracefree first-order extrinsic curvature, is tangent
to such an exact vacuum family through \((\delta,0)\).
\end{corollary}

\begin{proof}
Put \(\tau=0\) in (29.11).  Both terms on the left are nonnegative, so
both vanish.  Vanishing of \(\nabla h\) makes \(h\) constant on the flat
torus, and its mean-zero condition then makes it zero.
\end{proof}

Constant \(h\) with \(A=0\) is not excluded: it is tangent to the
finite-dimensional family of flat torus metrics.  The obstruction is to
oscillatory TT gravitons about a compact spacetime with a static
time-translation symmetry.

\subsection{The homogeneous compensation channel}

\begin{corollary}[Required mean-curvature backreaction]
\label{cor:v19v26-backreaction}
Suppose

\[
 A=A_0+\frac{\tau}{3}\delta,\qquad
 \partial^j(A_0)_{ij}=0,\qquad
 \operatorname{tr}A_0=0,\qquad \tau=\hbox{constant}.          \tag{29.15}
\]

Then the linearized momentum constraint holds, and the necessary
constant-lapse condition is

\[
 \frac23\,\tau^2\operatorname{Vol}(\mathbb T^3)
 =\frac14\int|\nabla h|^2\,dx+\int|A_0|^2\,dx.                \tag{29.16}
\]

Thus a nonzero TT seed requires a nonzero homogeneous expansion or
contraction at first order if no other sector contributes.
\end{corollary}

\begin{proof}
The linearized momentum constraint obtained from (29.1) is

\[
 \partial^j(A_{ij}-\tau\delta_{ij})=0.                        \tag{29.17}
\]

Conditions (29.15) make both terms divergence free.  Equation (29.16) is
(29.11) with constant \(\tau\).
\end{proof}

Equation (29.16) is necessary, not sufficient.  It identifies the
homogeneous variable that a nonlinear construction must solve together
with the TT amplitudes; higher orders and the remaining three
compatibility charges may impose further conditions.

\subsection{Positive matter energy strengthens the static obstruction}

For Einstein--matter data use

\[
 R(g)+(\operatorname{tr}_gK)^2-|K|_g^2
 =16\pi G\,\rho.                                               \tag{29.18}
\]

\begin{corollary}[Matter-coupled constant-lapse identity]
\label{cor:v19v26-matter}
Assume the background and first matter variation have zero energy density,
\(\rho(0)=\rho'(0)=0\), and retain the hypotheses of Theorem
\ref{thm:v19v26-taub}.  Then

\[
 \frac14\int|\nabla h|^2
 +\int|A_0|^2
 +8\pi G\int\rho''(0)
 =\frac23\int\tau^2.                                         \tag{29.19}
\]

If \(\int\rho''(0)\geq0\), matter cannot repair the
\(\tau=0\) obstruction; it increases the homogeneous compensation
required by (29.19).
\end{corollary}

\begin{proof}
Twice differentiate the integrated form of (29.18).  The assumptions on
\(\rho\) eliminate derivatives of the volume form multiplying
\(\rho(0)\) or \(\rho'(0)\).  The right side is therefore
\(16\pi G\int\rho''(0)\).  Repeating the rearrangement from (29.13) gives
(29.19).
\end{proof}

The sign assumption in this corollary is an explicit hypothesis.  It
holds for the integrated quadratic energy of the usual positive-energy
free bosonic sectors, but V26 does not prove it for the full constrained,
gauge-fixed, chiral Standard Model measure.

\subsection{Executable certificate}

The verifier

\[
 \texttt{scripts/verify\_sm\_v19\_v26\_taub\_compatibility.py} \tag{29.20}
\]

uses normalized torus average and the exact TT seeds

\[
 h=\cos z\,\operatorname{diag}(1,-1,0),\qquad
 A_0=\cos z\,\operatorname{diag}(1,-1,0).                     \tag{29.21}
\]

It checks

\[
 \left\langle|\nabla h|^2\right\rangle=1,\qquad
 \left\langle|A_0|^2\right\rangle=1,\qquad
 Q_{\rm TT}=\frac14+1=\frac54.                                \tag{29.22}
\]

The compensating homogeneous trace must satisfy
\(\tau^2=15/8\), for which
\(Q_{\rm TT}-(2/3)\tau^2=0\).  A positive matter test also increases the
charge.  The emitted artifact is

\[
 \texttt{artifacts/sm\_v19\_v26\_taub\_compatibility.json}.    \tag{29.23}
\]

The hashes are

\[
\begin{array}{c|l}
\hbox{object}&\hbox{SHA--256}\\ \hline
\hbox{verifier}&
\texttt{C39C8ABCD8417929AC9F64F48B5D02B00C1AB7E62C014F4BBFD4EDDFDBCF1E6F}\\
\hbox{canonical payload}&
\texttt{C9C9B0C795C7857BDF5391CF9FB1BF972D0CA59F56C66561FDCEB017B019BD51}.
\end{array}                                                     \tag{29.24}
\]

\begin{firewall}[One Taub charge is not nonlinear existence]
\label{fire:v19v26}
V26 derives one necessary constant-lapse compatibility identity on a
compact flat torus with \(\Lambda=0\).  It does not prove sufficiency,
solve the nonlinear constraints, analyze the three constant-shift
charges, handle a background without Killing fields, or construct a
positive gravitational measure.  Equation (29.19) is conditional on the
stated matter-energy sign and is not yet a Standard Model theorem.
\end{firewall}

\begin{decision}[V27 acquisition]
\label{dec:v19v26-next}
V27 will derive the three constant-shift quadratic compatibility charges
by differentiating the integrated momentum constraint against the
translation Killing fields.  It will determine whether TT Fourier pairs
carry a nonzero obstruction, how opposite momenta cancel, and how the
matter momentum density enters.  Only after all four compact-background
charges are explicit will the programme choose a nonlinear constraint
solver.
\end{decision}

\begin{assessment}[V26 acquisition]
\label{ass:v19v26}
The first nonlinear extension test fails for isolated compact TT data.
The constant lapse lies in the scalar-constraint cokernel, and its
second-order compatibility charge is a sum of TT metric-gradient and
TT momentum squares.  A homogeneous mean-curvature mode supplies the
only compensation identified here; nonnegative matter energy raises its
required magnitude.  This is the precise nonlinear role of the V24 zero
mode.
\end{assessment}

\begin{record}[V26 append record]
\label{rec:v19v26}
V26 is appended to the validated V25 whose installed SHA--256 was
\texttt{8D5FEE4DA0280E8ECF72AA7AB43592DA19EFECB09A8623CD63F29D841C319828}.
After installation only V26 is the active numeric SM note; frozen
archives and unrelated notes are unchanged.
\end{record}



\section{V27: the three translation compatibility charges}
\label{sec:v19v27}

V26 derived the constant-lapse charge on the compact flat background.
The remaining cokernel directions are the three constant shifts.  Their
quadratic charges are the total spatial momentum of the linearized TT
field.  This note derives them directly from the exact canonical momentum
map, so no coordinate expansion of the second-order Christoffel symbols
is needed.

\subsection{Exact momentum-map identity}

Let

\[
 \pi^{ij}:=\sqrt{\det g}\,
 \bigl(K^{ij}-(\operatorname{tr}_gK)g^{ij}\bigr)               \tag{30.1}
\]

be the ADM momentum density, with an inessential overall gravitational
constant suppressed.  For a vector field \(X\), the gravitational
momentum constraint smeared against \(X\) is

\[
 {\cal C}_g(X)
 :=\int_{\mathbb T^3}X^i{\cal C}_{g,i}\,dx
 =\int_{\mathbb T^3}\pi^{ij}({\cal L}_Xg)_{ij}\,dx.            \tag{30.2}
\]

The second equality follows by integrating the covariant divergence in
the local momentum constraint; it also fixes the sign convention used
below.

\begin{theorem}[Constant-shift quadratic compatibility]
\label{thm:v19v27-shift}
Let \((g(\varepsilon),\pi(\varepsilon))\) be a \(C^2\) family satisfying
the exact vacuum momentum constraint on \(\mathbb T^3\), with

\[
 g(0)=\delta,\qquad \pi(0)=0,\qquad
 h=g'(0),\qquad p=\pi'(0).                                    \tag{30.3}
\]

For each constant translation \(X=X^a\partial_a\),

\[
 Q_X(h,p):=\int_{\mathbb T^3}p^{ij}({\cal L}_Xh)_{ij}\,dx
 =X^a\int_{\mathbb T^3}p^{ij}\partial_a h_{ij}\,dx=0.         \tag{30.4}
\]
\end{theorem}

\begin{proof}
Because \(X\) is a Killing field of \(\delta\),
\({\cal L}_X\delta=0\).  Also \(\pi(0)=0\).  Therefore the value and first
derivative of (30.2) vanish at \(\varepsilon=0\), while its second
derivative is

\[
 \left.\frac{d^2}{d\varepsilon^2}\right|_0{\cal C}_g(X)
 =2\int_{\mathbb T^3}p^{ij}({\cal L}_Xh)_{ij}\,dx.             \tag{30.5}
\]

Terms containing \(\pi''(0)\) multiply
\({\cal L}_X\delta=0\), and terms containing \(g''(0)\) multiply
\(\pi(0)=0\).  The exact constraint makes (30.5) zero, proving (30.4).
\end{proof}

Thus the unknown second-order correction again drops out.  The four
background Killing fields give four necessary quadratic conditions:
(29.11) from time translation and the three components of (30.4) from
spatial translation.

\subsection{Fourier form and the phase test}

Use normalized torus measure and expand real TT tensors as

\[
 h_{ij}(x)=\sum_k h_{ij}(k)e^{ik\cdot x},\qquad
 p_{ij}(x)=\sum_k p_{ij}(k)e^{ik\cdot x},\qquad
 h(-k)=\overline{h(k)},\quad p(-k)=\overline{p(k)}.            \tag{30.6}
\]

Let \({\cal K}_+\) contain one representative of every pair
\(\{k,-k\}\), excluding zero.

\begin{proposition}[Translation charge in TT Fourier variables]
\label{prop:v19v27-fourier}
The three charges in (30.4) are

\[
 Q_a(h,p)
 =2\sum_{k\in{\cal K}_+}k_a\,
 \operatorname{Im}\bigl(
 \overline{h_{ij}(k)}p_{ij}(k)\bigr).                          \tag{30.7}
\]

Reality pairing of \(k\) with \(-k\) does not force (30.7) to vanish.
It makes the expression real.  Vanishing is an additional phase and
momentum-balance condition.
\end{proposition}

\begin{proof}
Parseval's identity gives

\[
 Q_a
 =\sum_k i k_a\,p^{ij}(-k)h_{ij}(k)
 =\sum_k i k_a\,\overline{p^{ij}(k)}h_{ij}(k).                 \tag{30.8}
\]

Pair the term at \(k\) with that at \(-k\).  If
\(z=\overline{p^{ij}(k)}h_{ij}(k)\), their sum is
\(ik_a(z-\overline z)=-2k_a\operatorname{Im}z\), which is
(30.7).
\end{proof}

\begin{corollary}[Standing and traveling TT modes]
\label{cor:v19v27-waves}
Let \(e^+=\operatorname{diag}(1,-1,0)\), so
\(e^+:e^+=2\), and use normalized measure on a torus with a \(z\)-circle
of period \(2\pi\).

\[
\begin{array}{c|c|c}
h&p&(Q_x,Q_y,Q_z)\\ \hline
\cos z\,e^+&\cos z\,e^+&(0,0,0)\\
\cos z\,e^+&\sin z\,e^+&(0,0,-1).
\end{array}                                                     \tag{30.9}
\]

A second nonzero mode with the opposite phase relation has charge
\((0,0,1)\), so a counterpropagating pair can satisfy all three
translation conditions without either member vanishing.
\end{corollary}

\begin{proof}
Only the \(z\)-derivative is nonzero and
\(\partial_z(\cos z\,e^+)=-\sin z\,e^+\).  Hence

\[
 Q_z=-2\langle\cos z\sin z\rangle=0
\]

for the first row, while

\[
 Q_z=-2\langle\sin^2z\rangle=-1
\]

for the second.  Reversing the relative sign of \(p\) reverses the
charge.  The sum of charges is additive.
\end{proof}

For example, the second member may use the orthogonal
\(\times\)-polarization, so cancellation need not cancel the fields
themselves.  Equation (30.7) is the familiar distinction between a
standing real wave and initial data carrying directed wave momentum.

\subsection{Independence from the homogeneous lapse compensator}

\begin{proposition}[Mean curvature does not cancel translation charge]
\label{prop:v19v27-independent}
Let \(h\) be traceless and let the first extrinsic-curvature variation be

\[
 A=A_0+\frac{\tau}{3}\delta,\qquad \tau=\hbox{constant}.       \tag{30.10}
\]

At the flat background, (30.1) gives

\[
 p=A-\tau\delta=A_0-\frac{2\tau}{3}\delta.                    \tag{30.11}
\]

The homogeneous trace term contributes zero to every \(Q_a\).  Therefore
choosing \(\tau\) to satisfy the constant-lapse identity (29.16) does not
alter any translation charge.
\end{proposition}

\begin{proof}
The contraction of the last term in (30.11) with
\(\partial_a h\) is

\[
 -\frac{2\tau}{3}\int\delta^{ij}\partial_a h_{ij}\,dx
 =-\frac{2\tau}{3}\int\partial_a(\operatorname{tr}h)\,dx=0.   \tag{30.12}
\]

Thus \(Q_a\) depends only on the tracefree part \(A_0\).
\end{proof}

V26 and V27 consequently impose independent requirements.  Homogeneous
expansion can balance the positive energy charge, while the traveling
wave momenta must balance among modes or against matter.

\subsection{Matter momentum}

Write the exact coupled momentum constraint in smeared canonical form as

\[
 {\cal C}_g(X)+{\cal C}_m(X)=0.                               \tag{30.13}
\]

\begin{corollary}[Total-momentum compatibility]
\label{cor:v19v27-matter}
Assume the matter momentum functional and its first variation vanish at
the background.  Define

\[
 P_X^{\,m}:=\frac12
 \left.\frac{d^2}{d\varepsilon^2}\right|_0{\cal C}_m(X).
                                                                    \tag{30.14}
\]

Then every exact Einstein--matter family satisfies

\[
 Q_X(h,p)+P_X^{\,m}=0                                         \tag{30.15}
\]

for the three constant shifts.
\end{corollary}

\begin{proof}
Twice differentiate (30.13), use (30.5), and divide by two.
\end{proof}

Equation (30.15) permits a matter excitation with opposite total momentum
to compensate one traveling graviton mode.  It does not assert that such
initial data solve the local constraints or that their quantum measure is
positive.

\subsection{Executable certificate}

The verifier

\[
 \texttt{scripts/verify\_sm\_v19\_v27\_momentum\_taub\_charges.py}
                                                                    \tag{30.16}
\]

evaluates (30.9) using the exact identities
\(e^+:e^+=2\),
\(\langle\sin^2z\rangle=1/2\), and
\(\langle\sin z\cos z\rangle=0\).  It checks the nonzero traveling
charge, a counterpropagating cancellation, and an equal opposite matter
charge.  It emits

\[
 \texttt{artifacts/sm\_v19\_v27\_momentum\_taub\_charges.json}. \tag{30.17}
\]

The hashes are

\[
\begin{array}{c|l}
\hbox{object}&\hbox{SHA--256}\\ \hline
\hbox{verifier}&
\texttt{0C2655C288D04B2F985150B2B8B0452B336C8A6272BEA0FE6BB39F857DC86978}\\
\hbox{canonical payload}&
\texttt{10C64844F6CB97D92F1B69E8DD304608FFFE2FAC28E2A4812BAC0DE637009EB1}.
\end{array}                                                     \tag{30.18}
\]

\begin{firewall}[Compatibility is still not a constraint solution]
\label{fire:v19v27}
Equations (29.11) and (30.4) are four necessary conditions obtained by
pairing with the Killing cokernel of the static flat torus.  They are not
sufficient for the pointwise Hamiltonian and momentum constraints.
The calculation does not cover noncompact boundary fluxes, nonconstant
shifts, higher-order obstructions, or quantum interactions.  Matter
compensation in (30.15) is only a total-momentum condition.
\end{firewall}

\begin{decision}[V28 acquisition]
\label{dec:v19v27-next}
V28 will move from necessary conditions to an exact solvable nonlinear
branch.  It will use the constant-mean-curvature conformal method on the
flat torus, first with a constant transverse traceless seed, solve the
Lichnerowicz equation in closed form, and verify both ADM constraints.
This will test whether the homogeneous mean-curvature channel identified
in V26 gives genuine initial data rather than only a quadratic
compensator.
\end{decision}

\begin{assessment}[V27 acquisition]
\label{ass:v19v27}
All four Killing-cokernel charges of the compact static background are now
explicit at quadratic order.  The lapse charge is energy-like and needs a
homogeneous mean-curvature contribution; the three shift charges are the
TT field's total spatial momentum and require counterpropagating or matter
balance.  These conditions expose exactly why the unconstrained product
of free V23 modes cannot yet be promoted to a nonlinear compact
Einstein--Standard Model measure.
\end{assessment}

\begin{record}[V27 append record]
\label{rec:v19v27}
V27 is appended to the validated V26 whose installed SHA--256 was
\texttt{BEC55BBACA751A7381ADCCD8128ACB15A5BAFB444743661D140557CE83BD6ACD}.
After installation only V27 is the active numeric SM note; frozen
archives and unrelated notes are unchanged.
\end{record}



\section{V28: an exact nonlinear CMC branch on the flat torus}
\label{sec:v19v28}

V26 found that the constant-lapse charge of TT data can be balanced by a
homogeneous mean curvature, while V27 showed that this balance is
independent of total spatial momentum.  The next question is sufficiency.
For constant transverse-traceless data, the constant-mean-curvature
(CMC) conformal method answers it exactly: the Lichnerowicz equation
becomes algebraic, and the compensating relation from V26 produces a
one-parameter family satisfying both nonlinear ADM constraints.

\subsection{Conformal reduction of the constraints}

Let \(\delta\) be a flat metric on \(\mathbb T^3\).  Choose a symmetric
tensor \(\sigma\) and a number \(\tau\) such that

\[
 \operatorname{tr}_\delta\sigma=0,\qquad
 \partial_j\sigma^{ij}=0,\qquad \partial_i\tau=0.              \tag{31.1}
\]

For a positive function \(\psi\), define

\[
 g_{ij}=\psi^4\delta_{ij},\qquad
 K_{ij}=\psi^{-2}\sigma_{ij}
       +\frac{\tau}{3}\psi^4\delta_{ij}.                       \tag{31.2}
\]

\begin{lemma}[CMC constraint reduction]
\label{lem:v19v28-reduction}
The data (31.2) satisfy the vacuum momentum constraint for every positive
\(\psi\).  Their vacuum Hamiltonian constraint is equivalent to

\[
 -8\Delta_\delta\psi
 -|\sigma|_\delta^2\psi^{-7}
 +\frac23\tau^2\psi^5=0.                                     \tag{31.3}
\]
\end{lemma}

\begin{proof}
Put \(A_{ij}:=K_{ij}-(\tau/3)g_{ij}=\psi^{-2}\sigma_{ij}\).
Raising indices with \(g\) gives

\[
 A^{ij}=\psi^{-10}\sigma^{ij}.                                \tag{31.4}
\]

For \(g=e^{2u}\delta\), \(u=2\log\psi\), its Christoffel symbols are

\[
 \Gamma^i{}_{jk}
 =\delta^i_k\partial_ju+\delta^i_j\partial_ku
  -\delta_{jk}\partial^iu.                                   \tag{31.5}
\]

In \(\nabla_jA^{ij}\), differentiating \(\psi^{-10}\) contributes
\(-10\psi^{-11}(\partial_j\psi)\sigma^{ij}\).  The two Christoffel
terms contribute
\((4+6)\psi^{-11}(\partial_j\psi)\sigma^{ij}\); the trace term in
(31.5) vanishes by (31.1).  Hence

\[
 \nabla_jA^{ij}=\psi^{-10}\partial_j\sigma^{ij}=0.             \tag{31.6}
\]

Since \(\tau\) is constant, this is the momentum constraint.

The conformal scalar-curvature and norm identities in three dimensions
are

\[
 R(g)=-8\psi^{-5}\Delta_\delta\psi,\qquad
 |A|_g^2=\psi^{-12}|\sigma|_\delta^2,\qquad
 |K|_g^2=|A|_g^2+\frac{\tau^2}{3}.                            \tag{31.7}
\]

Substitution into
\(R(g)+\tau^2-|K|_g^2=0\), followed by multiplication by \(\psi^5\),
gives (31.3).
\end{proof}

\subsection{Closed-form constant solution}

\begin{theorem}[Constant TT--CMC solution]
\label{thm:v19v28-constant}
Suppose \(\sigma\) has constant coefficients,
\[
 s^2:=|\sigma|_\delta^2>0,\qquad \tau\neq0.                   \tag{31.8}
\]
Then (31.3) has exactly one positive constant solution:

\[
 \psi_*^{12}=\frac{3s^2}{2\tau^2},
 \qquad
 \psi_*=\left(\frac{3s^2}{2\tau^2}\right)^{1/12}.             \tag{31.9}
\]

The resulting \(g\) and \(K\) solve both nonlinear vacuum ADM
constraints.
\end{theorem}

\begin{proof}
For constant \(\psi\), the Laplacian term in (31.3) vanishes, and the
remaining equation is

\[
 -s^2\psi^{-7}+\frac23\tau^2\psi^5=0.                         \tag{31.10}
\]

Multiplication by \(\psi^7>0\) gives
\((2/3)\tau^2\psi^{12}=s^2\), which has the unique positive solution
(31.9).  Lemma \ref{lem:v19v28-reduction} then verifies both constraints.
\end{proof}

If \(s=0\) and \(\tau\neq0\), or if \(s>0\) and \(\tau=0\), no positive
constant solution exists.  When both vanish, every positive constant
\(\psi\) gives time-symmetric flat data.  The nonzero TT and CMC pieces
must therefore balance.

\subsection{An exact family through the static flat datum}

\begin{corollary}[Nonlinear realization of the V26 compensator]
\label{cor:v19v28-family}
Let \(\sigma_1\) be a nonzero constant TT tensor and let
\(\tau_1\neq0\) satisfy

\[
 |\sigma_1|_\delta^2=\frac23\tau_1^2.                         \tag{31.11}
\]

For every real \(\lambda\), set

\[
 g(\lambda)=\delta,\qquad
 K(\lambda)=\lambda\left(\sigma_1+\frac{\tau_1}{3}\delta\right).
                                                                    \tag{31.12}
\]

Then \((g(\lambda),K(\lambda))\) is an exact vacuum solution of both ADM
constraints, passes through \((\delta,0)\), and has

\[
 g'(0)=0,\qquad K'(0)=\sigma_1+\frac{\tau_1}{3}\delta.         \tag{31.13}
\]
\end{corollary}

\begin{proof}
Use the conformal data
\(\sigma=\lambda\sigma_1\), \(\tau=\lambda\tau_1\), and \(\psi=1\).
Equation (31.11) makes (31.10) vanish for every \(\lambda\).  All
coefficients are spatially constant, so the momentum constraint
vanishes.  This gives (31.12) and (31.13).
\end{proof}

For this family, the V26 identity with \(h=0\) is exactly

\[
 \int|\sigma_1|^2
 =\frac23\tau_1^2\operatorname{Vol}(\mathbb T^3).             \tag{31.14}
\]

Thus the quadratic condition was not merely formal: in the homogeneous
CMC sector it is the exact all-orders Hamiltonian constraint.  The V27
translation charges also vanish because the data are spatially constant.

\subsection{Explicit rational branch}

Take

\[
 \sigma_1=\operatorname{diag}(2,2,-4),\qquad \tau_1=6.         \tag{31.15}
\]

Then

\[
 \operatorname{tr}\sigma_1=0,\qquad
 |\sigma_1|^2=24=\frac23(6)^2,\qquad \psi_*=1,                \tag{31.16}
\]

and

\[
 K(\lambda)=\lambda\,\operatorname{diag}(4,4,-2),\qquad
 (\operatorname{tr}K)^2=|K|^2=36\lambda^2.                   \tag{31.17}
\]

Since \(R(\delta)=0\), (31.17) verifies the Hamiltonian constraint
directly, and constancy verifies the momentum constraint.

\subsection{Executable certificate}

The verifier

\[
 \texttt{scripts/verify\_sm\_v19\_v28\_cmc\_constant\_branch.py}
                                                                    \tag{31.18}
\]

checks (31.15)--(31.17), the algebraic Lichnerowicz residual, and all
three momentum residuals exactly over the rationals.  It emits

\[
 \texttt{artifacts/sm\_v19\_v28\_cmc\_constant\_branch.json}.  \tag{31.19}
\]

The hashes are

\[
\begin{array}{c|l}
\hbox{object}&\hbox{SHA--256}\\ \hline
\hbox{verifier}&
\texttt{9981C4400AA211427015DA0FD14CD22AC0B541C8B11F53334B6B0B1EAA236A6A}\\
\hbox{canonical payload}&
\texttt{59760A8514ABC0ED8EED2F833C410A609F0CDFDF8DC022E3D674F01D7288EB86}.
\end{array}                                                     \tag{31.20}
\]

\begin{firewall}[The exact branch is homogeneous]
\label{fire:v19v28}
The family (31.12) proves nonlinear sufficiency only for spatially
constant CMC data.  It contains anisotropic homogeneous expansion but no
oscillatory graviton, no matter field, and no Euclidean probability
measure.  It does not show that the nonzero TT modes of V23 integrate to
exact data, nor that the resulting Lorentzian development is global.
\end{firewall}

\begin{decision}[V29 acquisition]
\label{dec:v19v28-next}
V29 will allow a spatially varying TT conformal seed
\(\sigma(x)\not\equiv0\) with constant \(\tau\neq0\).  It will prove
existence and uniqueness of a positive solution to (31.3) on the flat
torus, using explicit constant sub- and supersolutions where possible and
a maximum-principle uniqueness argument.  Degeneracy at zeros of
\(\sigma\) will be treated explicitly rather than hidden in a positive
lower-bound assumption.
\end{decision}

\begin{assessment}[V28 acquisition]
\label{ass:v19v28}
The homogeneous mean-curvature channel is a genuine nonlinear solution
mechanism.  The CMC conformal ansatz reduces both constraints to one
scalar equation, and constant TT data give a unique positive algebraic
solution.  Under the V26 balance, this produces an exact one-parameter
vacuum family through the static flat torus.  Spatially varying TT waves
remain the next test.
\end{assessment}

\begin{record}[V28 append record]
\label{rec:v19v28}
V28 is appended to the validated V27 whose installed SHA--256 was
\texttt{E0A18C4343466DC07B828BAE4FB78B16A1D7A880901E3997C75117A6E19A5D89}.
After installation only V28 is the active numeric SM note; frozen
archives and unrelated notes are unchanged.
\end{record}



\section{V29: variable TT seeds in the CMC conformal constraint}
\label{sec:v19v29}

V28 solved the CMC conformal constraint when the TT seed has constant
norm.  The same route in fact works for every smooth, nonzero TT seed on
the flat torus.  A constant lower barrier fails if the seed vanishes
somewhere.  The proof below removes that bottleneck by solving one
positive linear elliptic equation and scaling its solution to obtain a
strictly positive lower barrier.

\subsection{Statement and the degenerate cases}

Let

\[
 a(x):=|\sigma(x)|_\delta^2\geq0,\qquad
 b:=\frac23\tau^2.                                            \tag{32.1}
\]

The Lichnerowicz equation (31.3) is

\[
 {\cal L}(\psi):=-8\Delta\psi-a(x)\psi^{-7}+b\psi^5=0,
 \qquad \psi>0.                                                \tag{32.2}
\]

\begin{theorem}[Existence and uniqueness for a variable TT seed]
\label{thm:v19v29-existence}
Let \(\sigma\) be a smooth TT tensor on the flat three-torus,
\(\sigma\not\equiv0\), and let \(\tau\neq0\) be constant.  Then (32.2)
has a unique positive smooth solution.  Consequently (31.2) gives a
smooth exact solution of both vacuum ADM constraints.
\end{theorem}

The theorem permits \(a\) to vanish on sets of arbitrary shape.  It does
not assume a positive lower bound for \(|\sigma|\).

\begin{proposition}[Sharp elementary degeneracies]
\label{prop:v19v29-degenerate}
On the flat torus:
\begin{enumerate}
\item if \(a\not\equiv0\) and \(b=0\), (32.2) has no positive solution;
\item if \(a\equiv0\) and \(b>0\), (32.2) has no positive solution;
\item if \(a\equiv0=b\), the positive solutions are precisely the
constant functions.
\end{enumerate}
\end{proposition}

\begin{proof}
Integrate (32.2).  In the first case it gives
\(-\int a\psi^{-7}=0\), impossible.  In the second it gives
\(\int b\psi^5=0\), also impossible.  In the third case
\(\Delta\psi=0\), and every harmonic function on a compact connected
torus is constant.
\end{proof}

\subsection{A positive lower barrier through zeros of the seed}

Assume henceforth \(a\not\equiv0\) and \(b>0\).  Let \(v\) be the unique
solution of

\[
 (-8\Delta+b)v=a.                                             \tag{32.3}
\]

\begin{lemma}[Auxiliary positivity]
\label{lem:v19v29-auxiliary}
The solution \(v\) of (32.3) is smooth and strictly positive.
\end{lemma}

\begin{proof}
The operator \(-8\Delta+b\) is strictly positive and invertible on the
torus because \(b>0\).  Elliptic regularity gives a smooth solution.
The weak maximum principle gives \(v\geq0\), and the strong maximum
principle, together with \(a\geq0\) and \(a\not\equiv0\), gives
\(v>0\) everywhere.
\end{proof}

\begin{lemma}[Ordered sub- and supersolutions]
\label{lem:v19v29-barriers}
There exist constants \(\varepsilon>0\) and \(M>0\) such that

\[
 \underline\psi:=\varepsilon v,\qquad
 \overline\psi:=M                                             \tag{32.4}
\]

satisfy

\[
 0<\underline\psi\leq\overline\psi,\qquad
 {\cal L}(\underline\psi)\leq0,\qquad
 {\cal L}(\overline\psi)\geq0.                                \tag{32.5}
\]
\end{lemma}

\begin{proof}
Using (32.3),

\[
 \begin{split}
 {\cal L}(\varepsilon v)
 &=a\left(\varepsilon-\varepsilon^{-7}v^{-7}\right)
 +b\left(-\varepsilon v+\varepsilon^5v^5\right).              \tag{32.6}
 \end{split}
\]

Choose \(\varepsilon\) so small that, pointwise,

\[
 \varepsilon^8v^7\leq1,\qquad
 \varepsilon^4v^4\leq1.                                      \tag{32.7}
\]

Both brackets in (32.6) are then nonpositive.  This remains true where
\(a=0\), because the second bracket supplies the inequality there.
Next choose \(M\) large enough that

\[
 bM^{12}\geq\max_{\mathbb T^3}a,\qquad
 M\geq\varepsilon\max_{\mathbb T^3}v.                         \tag{32.8}
\]

Since
\({\cal L}(M)=M^{-7}(bM^{12}-a)\), equations (32.5) follow.
\end{proof}

\subsection{Existence and uniqueness}

\begin{proof}[Proof of Theorem \ref{thm:v19v29-existence}]
On the compact interval
\([\min\underline\psi,\max\overline\psi]\subset(0,\infty)\), put

\[
 F(x,s):=-a(x)s^{-7}+bs^5,\qquad
 \partial_sF=7a(x)s^{-8}+5bs^4>0.                             \tag{32.9}
\]

Choose \(\Lambda>\max\partial_sF\) on this interval and define
\[
 H(x,s):=\Lambda s-F(x,s).                                   \tag{32.10}
\]
Then \(H\) is nondecreasing in \(s\).  Starting from
\(\psi_0=\underline\psi\), solve successively

\[
 (-8\Delta+\Lambda)\psi_{n+1}=H(x,\psi_n).                    \tag{32.11}
\]

The lower inequality in (32.5), the upper inequality in (32.5), and the
maximum principle show inductively that

\[
 \underline\psi\leq\psi_n\leq\psi_{n+1}\leq\overline\psi.
                                                                    \tag{32.12}
\]

Uniform elliptic estimates and compactness give a limit \(\psi\) solving
(32.2); the positive lower barrier keeps the singular power uniformly
regular.  Bootstrapping gives smoothness.

For uniqueness, let \(\psi_1,\psi_2>0\) solve (32.2) and put
\(w=\psi_1-\psi_2\).  The mean-value theorem gives

\[
 -8\Delta w+c(x)w=0,\qquad
 c(x)=\int_0^1
 \partial_sF\bigl(x,\psi_2+t(\psi_1-\psi_2)\bigr)\,dt>0.
                                                                    \tag{32.13}
\]

At a positive maximum of \(w\), the two terms on the left of (32.13)
cannot sum to zero.  The same argument applied to \(-w\) rules out a
negative minimum.  Hence \(w=0\).

Finally Lemma \ref{lem:v19v28-reduction} converts this unique \(\psi\)
into exact vacuum constraint data.
\end{proof}

\subsection{What the theorem supplies to the nonlinear programme}

For every smooth TT momentum seed \(\sigma\not\equiv0\) and every
nonzero constant mean curvature, V29 constructs

\[
 g=\psi^4\delta,\qquad
 K=\psi^{-2}\sigma+\frac{\tau}{3}\psi^4\delta                 \tag{32.14}
\]

with no residual Hamiltonian or momentum constraint.  The conformal
metric response is determined rather than freely prescribed.  Therefore
this theorem does not contradict the V26 obstruction for a family with
an independently chosen first-order TT metric perturbation.

In particular, scaling
\(\sigma_\lambda=\lambda\sigma_0\) and
\(\tau_\lambda=\lambda\tau_0\) produces exact data for every
\(\lambda\neq0\).  Whether the unique \(\psi_\lambda\) extends smoothly
to \(\lambda=0\) with a prescribed flat limit is a singular
Lyapunov--Schmidt question, because at \(\lambda=0\) equation (32.2)
reduces to \(\Delta\psi=0\) and regains its constant kernel.

\subsection{Executable barrier certificate}

The continuum proof is analytic.  The verifier

\[
 \texttt{scripts/verify\_sm\_v19\_v29\_variable\_cmc\_barriers.py}
                                                                    \tag{32.15}
\]

provides an exact finite periodic analogue of the delicate barrier step.
On a four-site cycle it takes

\[
 a=(0,1,4,1),\qquad b=2                                      \tag{32.16}
\]

and solves \((-8\Delta_{\rm cyc}+b)v=a\) exactly:

\[
 v=\left(\frac{100}{153},\frac{25}{34},
          \frac{134}{153},\frac{25}{34}\right)>0.             \tag{32.17}
\]

With \(\varepsilon=1/100\) and \(M=2\), it verifies at every site that
\({\cal L}(\varepsilon v)\leq0\),
\({\cal L}(M)\geq0\), and
\(\varepsilon v\leq M\), including the site where \(a=0\).
It emits

\[
 \texttt{artifacts/sm\_v19\_v29\_variable\_cmc\_barriers.json}. \tag{32.18}
\]

The hashes are

\[
\begin{array}{c|l}
\hbox{object}&\hbox{SHA--256}\\ \hline
\hbox{verifier}&
\texttt{20D097B5B3EEB48BFA52133645A36DD86231C81B168B9F0FA9315704B3F2065A}\\
\hbox{canonical payload}&
\texttt{5F693EB5B41148CAFEF18A35AEB4EAC13F0647E5019DA10A04F19FCF189837E3}.
\end{array}                                                     \tag{32.19}
\]

\begin{firewall}[Constraint existence is not continuum quantization]
\label{fire:v19v29}
Theorem \ref{thm:v19v29-existence} is a classical CMC vacuum
initial-data theorem on a compact flat torus.  It prescribes conformally
flat \(g\) and does not allow an independent TT metric seed, nonconstant
mean curvature, Standard Model sources, or a quantum measure.  It gives
neither global Lorentzian evolution nor Euclidean reflection positivity.
The discrete artifact illustrates the barrier algebra and is not a
substitute for the continuum elliptic proof.
\end{firewall}

\begin{decision}[V30 audit and V31 acquisition]
\label{dec:v19v29-next}
V30 will perform the compulsory companion audit.  V31 will then resolve
the singular small-amplitude limit
\(\sigma_\lambda=\lambda\sigma_0\),
\(\tau_\lambda=\lambda\tau_0\) by separating the constant mode from the
mean-zero conformal factor.  The target is a rigorous branch through
\(\psi=1\), with the exact average balance on
\((\sigma_0,\tau_0)\) and the leading \(O(\lambda^2)\) conformal
backreaction computed.
\end{decision}

\begin{assessment}[V29 acquisition]
\label{ass:v19v29}
The CMC route now handles every smooth nonzero TT conformal momentum seed,
even when its norm vanishes locally.  A positive linear auxiliary
solution supplies the missing lower barrier, monotone iteration gives
existence, and the strict scalar nonlinearity gives uniqueness.  This is
an exact nonlinear constraint sector; connecting it smoothly to the
static flat datum is the next kernel-splitting problem.
\end{assessment}

\begin{record}[V29 append record]
\label{rec:v19v29}
V29 is appended to the validated V28 whose installed SHA--256 was
\texttt{1AA3A21B7C45C0BFB405259053997C52D3D36FDD941C2796F9D1628CBDD1F7C0}.
After installation only V29 is the active numeric SM note; frozen
archives and unrelated notes are unchanged.
\end{record}



\section{V30: compulsory companion audit after nonlinear CMC existence}
\label{sec:v19v30}

This is the required audit following V25.  The newest active companion
notes in the shared folder were inspected read-only before beginning the
small-amplitude kernel split announced in V29.

\begin{record}[V30 companion snapshot]
\label{rec:v19v30-snapshot}
The inspected files were

\[
\begin{array}{c|r|r}
\hbox{file and SHA--256}&\hbox{bytes}&\hbox{lines}\\ \hline
\texttt{Renewal\_Yang\_Mills\_Mass\_Gap\_Research\_Notes\_V29.tex}
&401357&9446\\
\multicolumn{3}{l}{
\texttt{F10A2A9A585B2C5DFDD33522B4722DC3D36BD70A5078391D42A2CC90A3EC595B}}
\\
\texttt{Renewal\_NS\_Stability\_Research\_Notes\_V26.tex}
&278283&4761\\
\multicolumn{3}{l}{
\texttt{82C887F8C2C69E4F28FAD7C3DC8DE5B9099CB5A4BF431EBA1FFF3E91935978AD}}.
\end{array}                                                     \tag{33.1}
\]

Since V25, Yang--Mills has advanced from V27 to V29 and
Navier--Stokes has not changed.
\end{record}

\subsection{Companion findings}

Yang--Mills V28 proved that changing its exact row-batch partition leaves
the finite-field result invariant, but measured no host acceleration.
V29 therefore moved upstream and proved that two disjoint prime jobs have
schedule-independent exact outputs; an executed gate authorized
concurrency two.  Its own firewall still leaves the remaining primes,
CRT reconstruction, physical reflection pairing, interacting measure,
regulator removal, and mass gap open.

Navier--Stokes V26 remains the rank-one Oseen-kernel derivative
refinement already read at V25.  It supplies no new statement for this
audit.

\begin{decision}[V30 audit decision]
\label{dec:v19v30-audit}
No companion theorem, numerical constant, or artifact is imported.
Yang--Mills process-level concurrency does not resolve the constant kernel
of the CMC Lichnerowicz operator, and the Navier--Stokes kernel norms are
analytically unrelated.  V31 will proceed with the mean/mean-zero
Lyapunov--Schmidt split announced in V29.  Reflection positivity will
remain a later, separate test on the resulting constrained variables.
\end{decision}

\begin{proof}
The V31 problem is the loss of invertibility of \(-8\Delta\) on constants
when \((\sigma,\tau)\to0\).  Neither companion constructs an inverse on
the mean-zero scalar sector or supplies the scalar average compatibility
equation.  Their claim boundaries also prevent importing a continuum or
positivity conclusion.  The announced direction therefore remains the
first unresolved dependency.
\end{proof}

\begin{assessment}[V30 acquisition]
\label{ass:v19v30}
The audit is current and leaves the nonlinear Einstein direction
unchanged.  The next task is to turn the family of nonzero-amplitude CMC
solutions from V29 into a smooth exact branch through the flat datum.
\end{assessment}

\begin{record}[V30 append record]
\label{rec:v19v30}
V30 is appended to the validated V29 whose installed SHA--256 was
\texttt{02D53AB3BD8EBA36768EC83D1DA6740207D565714F1B9CA22D101628882C9870}.
After installation only V30 is the active numeric SM note; the companions,
frozen archives, and all unrelated notes are unchanged.  The next
compulsory companion audit is V35.
\end{record}



\section{V31: a smooth small-amplitude CMC branch through flat data}
\label{sec:v19v31}

V29 constructs a unique CMC conformal solution whenever
\((\sigma,\tau)\neq(0,0)\), but at zero amplitude the Laplacian regains a
constant kernel.  This note resolves that singular limit by separating
the conformal factor into its spatial mean and mean-zero part.  The
constant-lapse balance from V26 becomes exactly the scalar
Lyapunov--Schmidt equation, and the leading conformal backreaction can be
computed.

\subsection{Scaled equation and necessary flat-limit balance}

Fix a smooth nonzero TT tensor \(\sigma_0\) and a nonzero constant
\(\tau_0\).  Put

\[
 a(x):=|\sigma_0(x)|^2,\qquad
 b:=\frac23\tau_0^2,\qquad
 \sigma_\lambda=\lambda\sigma_0,\qquad
 \tau_\lambda=\lambda\tau_0.                                 \tag{34.1}
\]

For \(\mu=\lambda^2\), equation (31.3) becomes

\[
 -8\Delta\psi+\mu f(x,\psi)=0,\qquad
 f(x,s):=-a(x)s^{-7}+bs^5.                                   \tag{34.2}
\]

Let \(\langle u\rangle\) denote the normalized spatial average.

\begin{proposition}[Necessary balance for the unit flat limit]
\label{prop:v19v31-necessary}
If positive solutions of (34.2) exist for \(\lambda\neq0\) and
\(\psi_\lambda\to1\) uniformly as \(\lambda\to0\), then

\[
 \boxed{\quad b=\langle a\rangle,\quad\hbox{equivalently}\quad
 \frac23\tau_0^2=\left\langle|\sigma_0|^2\right\rangle.\quad}  \tag{34.3}
\]
\end{proposition}

\begin{proof}
Integrate (34.2), divide by \(\mu>0\), and obtain

\[
 \left\langle-a\psi_\lambda^{-7}
 +b\psi_\lambda^5\right\rangle=0.                            \tag{34.4}
\]

Uniform convergence to one permits passage to the limit, giving
\(-\langle a\rangle+b=0\).
\end{proof}

If the limiting constant is allowed to be \(c_*\neq1\), the same argument
gives

\[
 c_*^{12}=\frac{\langle a\rangle}{b}.                         \tag{34.5}
\]

Thus every nonzero pair has a natural flat scale, but a branch through
the chosen metric \(\delta\) requires (34.3).

\subsection{Mean/mean-zero Lyapunov--Schmidt split}

Fix \(0<\alpha<1\).  Let

\[
 P u:=\langle u\rangle,\qquad Q:=I-P,\qquad
 C^{k,\alpha}_0:=\{u\in C^{k,\alpha}:\langle u\rangle=0\}.     \tag{34.6}
\]

Write

\[
 \psi=c+u,\qquad c\in\mathbb R,\qquad u\in C^{2,\alpha}_0.    \tag{34.7}
\]

\begin{theorem}[Analytic branch through \(\psi=1\)]
\label{thm:v19v31-branch}
Assume the exact balance (34.3).  There are \(\mu_0>0\) and unique
real-analytic functions

\[
 c:[0,\mu_0)\to\mathbb R,\qquad
 u:[0,\mu_0)\to C^{2,\alpha}_0                                \tag{34.8}
\]

with \(c(0)=1\), \(u(0)=0\), and \(c(\mu)+u(\mu)>0\), such that
(34.2) holds.  For \(|\lambda|<\sqrt{\mu_0}\),

\[
 \psi_\lambda:=c(\lambda^2)+u(\lambda^2)                     \tag{34.9}
\]

is therefore an even real-analytic function of \(\lambda\).  For
\(\lambda\neq0\), it is the unique positive solution from Theorem
\ref{thm:v19v29-existence}.
\end{theorem}

\begin{proof}
First solve the mean-zero equation

\[
 {\cal F}(\mu,c,u):=-8\Delta u+\mu Qf(x,c+u)=0.                \tag{34.10}
\]

At \((\mu,c,u)=(0,1,0)\), its derivative in \(u\) is

\[
 D_u{\cal F}(0,1,0)=-8\Delta:
 C^{2,\alpha}_0\longrightarrow C^{0,\alpha}_0,                \tag{34.11}
\]

an isomorphism.  The analytic implicit-function theorem gives a unique
analytic \(u=U(\mu,c)\) near \((0,1)\), with \(U(0,c)=0\).

The remaining scalar equation is the continuously extended average
constraint

\[
 {\cal G}(\mu,c):=
 P f\bigl(x,c+U(\mu,c)\bigr)=0.                               \tag{34.12}
\]

At \(\mu=0\),

\[
 {\cal G}(0,c)=-\langle a\rangle c^{-7}+bc^5.                 \tag{34.13}
\]

Balance (34.3) gives
\({\cal G}(0,1)=0\), and

\[
 \partial_c{\cal G}(0,1)
 =7\langle a\rangle+5b=12b>0.                                \tag{34.14}
\]

A second analytic implicit-function theorem therefore gives a unique
\(c=c(\mu)\) near one.  Equations (34.10) and (34.12) are respectively
the mean-zero and mean parts of (34.2), so their solution satisfies the
full equation.  Positivity follows by shrinking \(\mu_0\).  Analyticity
in \(\mu=\lambda^2\) proves evenness.  For nonzero \(\lambda\), global
positive uniqueness is Theorem \ref{thm:v19v29-existence}.
\end{proof}

\subsection{Leading conformal backreaction}

\begin{theorem}[Second-order coefficient]
\label{thm:v19v31-leading}
Under (34.3),

\[
 \psi_\lambda=1+\lambda^2w+O(\lambda^4)
 \quad\hbox{in }C^{2,\alpha}.                                \tag{34.15}
\]

Let \(w_0\in C^{2,\alpha}_0\) be the unique solution

\[
 -8\Delta w_0=a-b,\qquad \langle w_0\rangle=0.                \tag{34.16}
\]

Then

\[
 w=w_0+\overline w,\qquad
 \overline w
 =-\frac{7}{12b}\langle a w_0\rangle
 =-\frac{14}{3b}\left\langle|\nabla w_0|^2\right\rangle
 \leq0.                                                       \tag{34.17}
\]

Equality holds precisely when \(a\) is constant.
\end{theorem}

\begin{proof}
Insert \(\psi=1+\mu w+O(\mu^2)\) into (34.2).  The coefficient of
\(\mu\) gives

\[
 -8\Delta w-a+b=0.                                           \tag{34.18}
\]

Its mean-zero part is (34.16); the constant part of \(w\) is not fixed at
this order.  Expand the average equation (34.4) one order further:

\[
 0=\left\langle(7a+5b)w\right\rangle.                         \tag{34.19}
\]

Write \(w=w_0+\overline w\), use
\(\langle w_0\rangle=0\) and \(\langle a\rangle=b\), and solve:

\[
 12b\,\overline w+7\langle a w_0\rangle=0.                   \tag{34.20}
\]

Multiplying (34.16) by \(w_0\), integrating by parts, and using its zero
mean gives

\[
 \langle a w_0\rangle
 =8\left\langle|\nabla w_0|^2\right\rangle.                   \tag{34.21}
\]

Equations (34.20)--(34.21) prove (34.17).  The gradient term vanishes
exactly when \(w_0\) is constant, hence zero; (34.16) then says \(a=b\).
\end{proof}

\subsection{Exact constrained family and its tangent}

\begin{corollary}[Smooth exact initial data through \((\delta,0)\)]
\label{cor:v19v31-data}
Define

\[
 g_\lambda=\psi_\lambda^4\delta,\qquad
 K_\lambda=\lambda\psi_\lambda^{-2}\sigma_0
 +\frac{\lambda\tau_0}{3}\psi_\lambda^4\delta.                \tag{34.22}
\]

Then, for sufficiently small \(|\lambda|\), these are smooth exact vacuum
solutions of both ADM constraints and

\[
 g_0=\delta,\quad K_0=0,\quad
 g'_0=0,\quad
 K'_0=\sigma_0+\frac{\tau_0}{3}\delta,\quad
 g''_0=8w\,\delta.                                            \tag{34.23}
\]
\end{corollary}

\begin{proof}
Lemma \ref{lem:v19v28-reduction} and Theorem
\ref{thm:v19v31-branch} give the exact constraints.  Evenness of
\(\psi_\lambda\) gives \(\psi'_0=0\), and differentiating (34.22) gives
(34.23).  The expansion
\(\psi_\lambda^4=1+4\lambda^2w+O(\lambda^4)\) yields
\(g''_0=8w\delta\).
\end{proof}

The tangent in (34.23) satisfies all four quadratic compatibility
conditions: (34.3) is the V26 lapse balance with \(h=0\), while the V27
translation charges vanish because \(g'_0=0\).  The spatially varying
conformal response first appears at second order.

\subsection{Explicit Fourier example}

Take

\[
 \sigma_0=\cos z\,\operatorname{diag}(1,-1,0).                \tag{34.24}
\]

Then

\[
 a=1+\cos2z,\qquad \langle a\rangle=b=1,\qquad
 \tau_0^2=\frac32.                                            \tag{34.25}
\]

Equations (34.16)--(34.17) give

\[
 w_0=\frac1{32}\cos2z,\qquad
 \langle aw_0\rangle=\frac1{64},\qquad
 \overline w=-\frac7{768},                                   \tag{34.26}
\]

and hence

\[
 \psi_\lambda
 =1+\lambda^2\left(\frac1{32}\cos2z-\frac7{768}\right)
 +O(\lambda^4).                                               \tag{34.27}
\]

\subsection{Executable certificate}

The verifier

\[
 \texttt{scripts/verify\_sm\_v19\_v31\_small\_cmc\_branch.py} \tag{34.28}
\]

checks every rational coefficient in (34.25)--(34.27), including the
Poisson residual, the averaged next-order residual, and
\(8\langle|\nabla w_0|^2\rangle=\langle aw_0\rangle\).  It emits

\[
 \texttt{artifacts/sm\_v19\_v31\_small\_cmc\_branch.json}.     \tag{34.29}
\]

The hashes are

\[
\begin{array}{c|l}
\hbox{object}&\hbox{SHA--256}\\ \hline
\hbox{verifier}&
\texttt{7C8024AC8A5BDBE467E73DB1013ABA8727E1F904AB7DE45F97CEAA0B655EEEDB}\\
\hbox{canonical payload}&
\texttt{D9EAAA39F80168716C9D02B8D0A2D54C01D4EB78BEC03EDD1FF404DD498D040B}.
\end{array}                                                     \tag{34.30}
\]

\begin{firewall}[The branch solves initial constraints, not the quantum theory]
\label{fire:v19v31}
Corollary \ref{cor:v19v31-data} is an exact local branch of classical
vacuum initial data with conformally flat metric tangent.  It does not
permit an independent first-order TT metric displacement, include
Standard Model sources, construct the Lorentzian development, or define
a Euclidean measure.  Analyticity of the elliptic constraint solver does
not imply reflection positivity of a pushforward measure.
\end{firewall}

\begin{decision}[V32 acquisition]
\label{dec:v19v31-next}
V32 will establish the exact measure-theoretic condition under which a
nonlinear spatial constraint solver preserves reflection positivity:
the map must commute with time reflection and send positive-time output
observables to positive-time input observables.  The theorem will then be
tested on the time-slice CMC map, with the time parity of
\(\sigma,\tau,g,K\) stated explicitly.  This is the next bridge from the
classical constraint branch to the continuum reconstruction programme.
\end{decision}

\begin{assessment}[V31 acquisition]
\label{ass:v19v31}
The singular zero-amplitude CMC limit is resolved.  The V26 average
energy balance is necessary and sufficient for a unique analytic
constraint branch through the chosen flat scale, and V31 computes its
leading spatial and homogeneous conformal backreaction.  This provides
exact nonlinear initial data with spatially varying TT momentum; quantum
reflection positivity remains a separate measure-level problem.
\end{assessment}

\begin{record}[V31 append record]
\label{rec:v19v31}
V31 is appended to the validated V30 whose installed SHA--256 was
\texttt{7F5C01B84875D6EC15DBFC60F34C60B434FCD54B9FB7E7B3807449925536DA0E}.
After installation only V31 is the active numeric SM note; frozen
archives and unrelated notes are unchanged.
\end{record}


\section{V32: reflection positivity through a nonlinear constraint map}
\label{sec:v19v32}

V31 produced an exact analytic branch of constrained initial data, and its
decision asked for the measure-theoretic condition under which a nonlinear
spatial constraint solver preserves reflection positivity.  This note
proves that condition, shows by exact examples that each of its two
hypotheses is indispensable, and then tests the time-slice CMC map.  The
test has two parts.  The solution map is real-analytic and even in the
seed, so it is automatically equivariant for the reflection under which
the extrinsic curvature is odd.  That odd parity is forced, not chosen:
the reflected form of a real Euclidean time derivative is a positive
square only with the odd sign.  But an odd field that is
\(L^2\)-continuous across the reflection plane must vanish there.  Hence
a sharp-time constraint solver receives the zero seed in continuous time,
and the nonvacuous version of Decision \ref{dec:v19v31-next} lives on a
Euclidean time lattice.

\subsection{Reflection-positive triples and their pushforward}

\begin{definition}[Reflection-positive triple]
\label{def:v19v32-triple}
A reflection-positive triple \((\mu,\Theta,{\cal H}_+)\) consists of a
probability space \((\Omega,{\cal F},\mu)\), a measurable involution
\(\Theta:\Omega\to\Omega\) with \(\Theta_*\mu=\mu\), and a closed
subspace \({\cal H}_+\subset L^2(\mu)\) such that

\[
 \langle F,F\rangle_\Theta
 :=\int_\Omega\overline{F(\Theta\omega)}\,F(\omega)\,d\mu(\omega)
 \geq0\qquad(F\in{\cal H}_+).                                  \tag{35.1}
\]

Since \(\Theta\) preserves \(\mu\), the map \(F\mapsto F\circ\Theta\) is
a unitary involution of \(L^2(\mu)\), so (35.1) is a bounded Hermitian
form and it suffices to verify it on a dense subset of \({\cal H}_+\).
\end{definition}

\begin{theorem}[Pushforward criterion]
\label{thm:v19v32-pushforward}
Let \((\mu,\Theta,{\cal H}_+)\) be reflection positive.  Let
\((\Omega',{\cal F}')\) be a measurable space with a measurable
involution \(\Theta'\), let \(\Phi:\Omega\to\Omega'\) be measurable, and
put \(\nu:=\Phi_*\mu\).  Let \({\cal H}'_+\subset L^2(\nu)\) be a closed
subspace.  Assume

\[
 \Phi\circ\Theta=\Theta'\circ\Phi\quad\mu\hbox{-a.e.},
 \qquad\hbox{and}\qquad
 G\circ\Phi\in{\cal H}_+\quad(G\in{\cal H}'_+).                \tag{35.2}
\]

Then \(\Theta'_*\nu=\nu\), and \((\nu,\Theta',{\cal H}'_+)\) is
reflection positive.  More precisely,

\[
 \langle G,G\rangle_{\Theta'}
 =\langle G\circ\Phi,G\circ\Phi\rangle_\Theta\geq0
 \qquad(G\in{\cal H}'_+).                                       \tag{35.3}
\]
\end{theorem}

\begin{proof}
Invariance: \(\Theta'_*\nu=(\Theta'\circ\Phi)_*\mu
=(\Phi\circ\Theta)_*\mu=\Phi_*(\Theta_*\mu)=\Phi_*\mu=\nu\).
The map \(G\mapsto G\circ\Phi\) is an isometry
\(L^2(\nu)\to L^2(\mu)\) because \(\nu\) is the image measure.  For
\(G\in{\cal H}'_+\), the change of variables and the first equation in
(35.2) give

\[
 \int\overline{G(\Theta'\omega')}G(\omega')\,d\nu
 =\int\overline{G(\Theta'\Phi\omega)}G(\Phi\omega)\,d\mu
 =\int\overline{(G\circ\Phi)(\Theta\omega)}(G\circ\Phi)(\omega)\,d\mu.
\]

This is (35.3); the right side is nonnegative by the second condition
in (35.2) and (35.1).
\end{proof}

The first condition says that \(\Phi\) is reflection equivariant; the
second says that positive-time output observables depend only on
positive-time input.  Both are necessary in the sense that either one may
fail alone and destroy positivity.

\subsection{Both hypotheses are necessary}

\begin{lemma}[A two-time Gaussian pair is reflection positive]
\label{lem:v19v32-pair}
Let \(0<r<1\), let \(\mu\) be the centred Gaussian law on
\(\Omega=\mathbb R^2\) with coordinates \((x_{-1},x_1)\), unit variances,
and covariance \(r\), and let \(\Theta(x_{-1},x_1)=(x_1,x_{-1})\).  Let
\({\cal H}_+\) be the \(L^2\)-closure of the polynomials in \(x_1\).
Then \((\mu,\Theta,{\cal H}_+)\) is reflection positive.
\end{lemma}

\begin{proof}
Let \(z,\xi_+,\xi_-\) be independent standard normals and set
\(x_{\pm1}=\sqrt r\,z+\sqrt{1-r}\,\xi_\pm\).  This realizes \(\mu\),
since the variances are \(r+(1-r)=1\) and the covariance is \(r\).
Conditionally on \(z\), the two coordinates are independent and
identically distributed, so for a polynomial \(p\),

\[
 \langle p(x_1),p(x_1)\rangle_\Theta
 =E\bigl[\overline{p(x_{-1})}\,p(x_1)\bigr]
 =E_z\bigl[\,|m_p(z)|^2\bigr]\geq0,
 \qquad m_p(z):=E[p(x_1)\mid z].                               \tag{35.4}
\]

Density of polynomials in \({\cal H}_+\) and boundedness of the form
finish the proof.
\end{proof}

\begin{proposition}[Two exact failures]
\label{prop:v19v32-failures}
In the setting of Lemma \ref{lem:v19v32-pair}, take
\(\Omega'=\mathbb R^2\), \(\Theta'=\Theta\), and let \({\cal H}'_+\) be
the closure of the polynomials in the second coordinate \(y_1\).
\begin{enumerate}
\item The slice-local map \(\Phi_a(x_{-1},x_1)=(-x_{-1},x_1)\) satisfies
      the locality condition in (35.2) but not equivariance, and

\[
 \langle y_1,y_1\rangle_{\Theta'}=-r<0
 \qquad\hbox{under }(\Phi_a)_*\mu.                              \tag{35.5}
\]

\item The map \(\Phi_b(x_{-1},x_1)=(x_{-1}-x_1,\;x_1-x_{-1})\) is
      equivariant but violates locality, and

\[
 \langle y_1,y_1\rangle_{\Theta'}=-2(1-r)<0
 \qquad\hbox{under }(\Phi_b)_*\mu.                              \tag{35.6}
\]
\end{enumerate}
In both cases the reflected Gram matrix of \(\{1,y_1,y_1^2\}\) has a
negative principal minor, so the pushforward is not reflection positive.
\end{proposition}

\begin{proof}
For \(\Phi_a\): \(\Phi_a\Theta(x)=(-x_1,x_{-1})\) while
\(\Theta'\Phi_a(x)=(x_1,-x_{-1})\), so equivariance fails; the output
\(y_1=x_1\) is a positive-time input, so locality holds.  The reflected
form of \(y_1\) is \(E[y_{-1}y_1]=E[(-x_{-1})x_1]=-r\).

For \(\Phi_b\): with \(y=\Phi_b(x)\) one has \(y_{-1}=-y_1\), hence
\(\Phi_b(\Theta x)=(x_1-x_{-1},x_{-1}-x_1)=\Theta'\Phi_b(x)\), so
equivariance holds; but \(y_1\) depends on \(x_{-1}\), so locality
fails.  The reflected form is
\(E[y_{-1}y_1]=-E[(x_1-x_{-1})^2]=-(2-2r)\).

The Gram statement follows because \(\langle y_1,y_1\rangle_{\Theta'}\)
is itself a one-by-one principal minor.
\end{proof}

\subsection{The induced parity of a Euclidean time derivative}

Let \(q\) be the real stationary Gaussian process on \(\mathbb R\) with
covariance

\[
 C(t,s)=c\,e^{-\omega|t-s|},\qquad c,\omega>0,                 \tag{35.7}
\]

which is the V23 covariance (26.16) of one TT polarization, with
\(c=1/(2c_g\omega)\).  Realize it on path space with the path reflection
\((\Theta\omega)(t)=\omega(-t)\), which preserves the law because
(35.7) is even.  For \(f\in C_c^\infty(\mathbb R)\) write
\(q(f)=\int f(t)q(t)\,dt\), and define the derivative field by

\[
 \dot q(f):=-q(f'),\qquad (Rf)(t):=f(-t).                       \tag{35.8}
\]

Let \({\cal H}_+\) be the closure of polynomials in the variables
\(q(f)\) with \(\operatorname{supp}f\subset(0,\infty)\).  Reflection
positivity of \((\mu,\Theta,{\cal H}_+)\) is the content of Theorem
\ref{thm:v19v23-os}.

\begin{proposition}[The derivative is odd, and only the odd sign is positive]
\label{prop:v19v32-derivative}
For every test function \(f\),

\[
 \Theta\,\dot q(f)=-\dot q(Rf).                                 \tag{35.9}
\]

If \(\operatorname{supp}f\subset(0,\infty)\), then

\[
 \langle\dot q(f),\dot q(f)\rangle_\Theta
 =c\,\omega^2\left(\int_0^\infty f(t)e^{-\omega t}\,dt\right)^2\geq0,
                                                                    \tag{35.10}
\]

whereas the even assignment
\(\widetilde\Theta\,\dot q(f):=\dot q(Rf)\) gives the negative of the
right side of (35.10), which is strictly negative for every nonzero
\(f\geq0\).
\end{proposition}

\begin{proof}
On observables, \(\Theta q(f)=q(Rf)\).  Since \((Rf)'=-R(f')\),

\[
 \Theta\,\dot q(f)=-q(R(f'))=q((Rf)')=-\dot q(Rf),
\]

which is (35.9).  Next, with \(\operatorname{supp}f\subset(0,\infty)\)
and \(\operatorname{supp}Rf\subset(-\infty,0)\),

\[
 \langle\dot q(f),\dot q(f)\rangle_\Theta
 =-E[\dot q(Rf)\,\dot q(f)]
 =-\iint(Rf)'(u)\,f'(s)\,C(u,s)\,du\,ds.
\]

On the support of the integrand \(u<0<s\), where
\(C(u,s)=c\,e^{-\omega(s-u)}\) is smooth, so two integrations by parts
give

\[
 -\iint(Rf)(u)f(s)\,\partial_u\partial_sC(u,s)\,du\,ds
 =c\,\omega^2\iint f(-u)f(s)e^{-\omega(s-u)}\,du\,ds,
\]

because \(\partial_u\partial_s e^{-\omega(s-u)}=-\omega^2e^{-\omega(s-u)}\).
Substituting \(t=-u\) gives (35.10).  The even assignment differs by an
overall sign.
\end{proof}

Thus if the Euclidean extrinsic curvature is a real time derivative of
the spatial metric, its reflection parity is odd by (35.9), and any even
assignment contradicts (35.10).  The lattice form of this fact is exact:
for the Gaussian chain \(E[q_nq_m]=r^{|n-m|}\) with link variables
\(d_{n+1/2}=q_{n+1}-q_n\), the site reflection sends \(d_{n+1/2}\) to
\(-d_{-n-1/2}\), and for \(n,m\geq0\)

\[
 \langle d_{n+1/2},d_{m+1/2}\rangle_\Theta
 =(1-r)^2r^{\,n+m},                                             \tag{35.11}
\]

a rank-one positive matrix, while the even assignment gives its
negative.  Both statements are checked exactly in the certificate below.

\subsection{Sharp-time odd fields vanish}

\begin{theorem}[An \(L^2\)-continuous odd field vanishes on the reflection plane]
\label{thm:v19v32-sharp}
Let \((\mu,\Theta,{\cal H}_+)\) be reflection positive and let
\(t\mapsto\sigma(t)\in L^2(\mu)\) be real valued with

\[
 \sigma(t)\circ\Theta=-\sigma(-t)\quad(t\in\mathbb R),
 \qquad
 \sigma(t)\in{\cal H}_+\quad(t>0).                              \tag{35.12}
\]

Then

\[
 E[\sigma(-t)\sigma(t)]\leq0\qquad(t>0).                        \tag{35.13}
\]

If \(t\mapsto\sigma(t)\) is \(L^2\)-continuous at \(t=0\), then
\(\sigma(0)=0\) almost surely.  If in addition \(E[\sigma(t)^2]\) does
not depend on \(t\), then \(\sigma(t)=0\) almost surely for every
\(t\).
\end{theorem}

\begin{proof}
For \(t>0\), (35.1) and (35.12) give
\(0\leq\langle\sigma(t),\sigma(t)\rangle_\Theta
=E[(\sigma(t)\circ\Theta)\sigma(t)]=-E[\sigma(-t)\sigma(t)]\), which is
(35.13).  If \(\sigma(\pm t)\to\sigma(0)\) in \(L^2\) as \(t\downarrow0\),
then the inner products converge and
\(E[\sigma(0)^2]=\lim E[\sigma(-t)\sigma(t)]\leq0\), so
\(\sigma(0)=0\).  Constant variance then forces every
\(\sigma(t)\) to vanish.
\end{proof}

\begin{corollary}[No sharp-time extrinsic-curvature seed in continuous time]
\label{cor:v19v32-no-seed}
Let a reflection-positive Euclidean gravitational process carry an
extrinsic-curvature seed \(\sigma(t)\) that is odd in the sense of
(35.12) and stationary in variance.  If \(\sigma\) has an
\(L^2\)-continuous restriction to sharp times, then \(\sigma\equiv0\), and
the time-slice CMC map of V29 returns time-symmetric flat data at every
time.  Consequently a nontrivial odd seed has no \(L^2\)-continuous
sharp-time restriction; its coincident-time covariance carries a singular
part, as the term \(2\omega\delta(t-s)\) in
\(\partial_t\partial_sC\) for (35.7) illustrates.
\end{corollary}

\begin{proof}
Apply Theorem \ref{thm:v19v32-sharp} to each spatial component of
\(\sigma\) paired with a fixed smooth TT test tensor, and to
\(\tau\).  With zero seed, (32.2) has only constant solutions, which give
\(g=\hbox{const}\cdot\delta\) and \(K=0\).
\end{proof}

Physically, (35.13) is the Euclidean shadow of
\(\langle\Omega,pe^{-2tH}p\,\Omega\rangle\geq0\): the real Euclidean
derivative equals \(i\) times the Euclidean evolution of the
self-adjoint momentum, and the resulting sign forbids a real random
variable with the required negative variance at the plane.

\subsection{The CMC solution map is analytic and even in the seed}

Fix \(0<\alpha<1\) and let \({\cal T}^{0,\alpha}\) denote the
\(C^{0,\alpha}\) transverse-traceless tensors on the flat torus.  Put

\[
 {\cal U}:=\{(\sigma,\tau)\in{\cal T}^{0,\alpha}\times\mathbb R:
 \sigma\not\equiv0,\ \tau\neq0\},                                \tag{35.14}
\]

an open set.

\begin{lemma}[V29 for H\"older seeds]
\label{lem:v19v32-hoelder}
For \((\sigma,\tau)\in{\cal U}\), equation (32.2) has a unique positive
solution \(\psi\in C^{2,\alpha}\).
\end{lemma}

\begin{proof}
The proof of Theorem \ref{thm:v19v29-existence} uses only Schauder
theory and the maximum principle.  With \(a=|\sigma|^2\in C^{0,\alpha}\),
the auxiliary solution of (32.3) lies in \(C^{2,\alpha}\) and is
strictly positive by the strong maximum principle.  Each monotone
iterate (32.11) has a \(C^{0,\alpha}\) right side and therefore lies in
\(C^{2,\alpha}\), with bounds uniform in \(n\) by the barriers.
Arzel\`a--Ascoli gives a \(C^2\) limit solving (32.2), and one more
Schauder step gives \(C^{2,\alpha}\).  Uniqueness is unchanged.
\end{proof}

\begin{theorem}[Analytic even solution map]
\label{thm:v19v32-analytic}
The solution map

\[
 \Psi:{\cal U}\to C^{2,\alpha},\qquad
 (\sigma,\tau)\mapsto\psi                                         \tag{35.15}
\]

is real-analytic, and

\[
 \Psi(-\sigma,-\tau)=\Psi(\sigma,\tau).                          \tag{35.16}
\]

Consequently the constrained data (31.2) satisfy

\[
 g(-\sigma,-\tau)=g(\sigma,\tau),\qquad
 K(-\sigma,-\tau)=-K(\sigma,\tau).                                \tag{35.17}
\]
\end{theorem}

\begin{proof}
Define, on the open subset of
\({\cal T}^{0,\alpha}\times\mathbb R\times C^{2,\alpha}\) where
\(\psi>0\),

\[
 {\cal N}(\sigma,\tau,\psi)
 :=-8\Delta\psi-|\sigma|^2\psi^{-7}+\tfrac23\tau^2\psi^5
 \in C^{0,\alpha}.                                                \tag{35.18}
\]

The map is polynomial in \(\sigma\) and \(\tau\), and the Nemytskii maps
\(\psi\mapsto\psi^{-7},\psi^5\) are real-analytic
\(C^{2,\alpha}\to C^{2,\alpha}\) on \(\{\psi>0\}\), so \({\cal N}\) is
real-analytic.  At a solution with \(\tau\neq0\),

\[
 D_\psi{\cal N}=-8\Delta+c(x),\qquad
 c=7|\sigma|^2\psi^{-8}+\tfrac{10}3\tau^2\psi^4
 \geq\tfrac{10}3\tau^2\min\psi^4>0.                              \tag{35.19}
\]

The operator \(-8\Delta+c\) with \(c>0\) is injective on \(C^{2,\alpha}\)
by the maximum principle and is Fredholm of index zero from
\(C^{2,\alpha}\) to \(C^{0,\alpha}\), hence an isomorphism.  The analytic
implicit-function theorem yields, near each point of \({\cal U}\), a
unique analytic branch of positive solutions.  By Lemma
\ref{lem:v19v32-hoelder} that branch coincides with \(\Psi\), which is
therefore analytic on \({\cal U}\).  Since \({\cal N}\) depends on
\((\sigma,\tau)\) only through \(|\sigma|^2\) and \(\tau^2\), uniqueness
gives (35.16).  Finally (31.2) is linear in \((\sigma,\tau)\) at fixed
\(\psi\), which gives (35.17).
\end{proof}

Theorem \ref{thm:v19v31-branch} is the special case of (35.16) along the
ray \((\lambda\sigma_0,\lambda\tau_0)\): evenness in \(\lambda\) is the
parity statement.

\begin{corollary}[Slice-local equivariance]
\label{cor:v19v32-slicewise}
Let \(\Omega\) be a space of seed paths
\(\omega=(\sigma(\cdot),\tau(\cdot))\) with values in \({\cal U}\), and
let \(\Theta\) act by
\((\Theta\omega)(t)=(-\sigma(-t),-\tau(-t))\).  Define the slice-wise
map

\[
 \Phi(\omega)(t):=\bigl(g,K\bigr)\bigl(\sigma(t),\tau(t)\bigr),
 \qquad
 (\Theta'\omega')(t):=\bigl(g(-t),-K(-t)\bigr).                   \tag{35.20}
\]

Then \(\Phi\) is measurable, equivariant, and positive-time local, so
Theorem \ref{thm:v19v32-pushforward} applies to any reflection-positive
seed law on \(\Omega\) whose positive-time subspace is generated by
sharp-time observables.
\end{corollary}

\begin{proof}
Measurability follows from continuity of \(\Psi\).  Equivariance is
(35.17) evaluated at \(-t\).  Locality holds because the output at time
\(t\) is a function of the input at the same time.
\end{proof}

\subsection{What the test decides}

Corollary \ref{cor:v19v32-slicewise} is the positive half of the V31
decision: nothing in the nonlinear constraint solver obstructs
reflection positivity.  Theorem \ref{thm:v19v32-sharp} is the negative
half: the sharp-time seed hypothesis is empty in continuous time, because
a stationary odd seed with sharp-time values that vary continuously in
\(L^2\) vanishes.  The two halves are compatible: Corollary
\ref{cor:v19v32-slicewise} is not vacuous on a Euclidean time lattice,
where odd variables sit on time links, the link reflection is exactly
the sign rule (35.11), and every link variable is a genuine random
variable.  Moreover the pushforward (35.3) is only a change of variables:
(31.2) is inverted by
\(\psi=(\det g/\det\delta)^{1/12}\), \(\tau=\operatorname{tr}_gK\),
\(\sigma=\psi^2(K-\tfrac\tau3g)\), so the output algebra equals the seed
algebra and the reconstructed semigroup is the seed's own.  The
constraint solver adds kinematics, not dynamics.

\begin{warning}[Signature of the slice constraint]
\label{warn:v19v32-signature}
Equations (29.1) and (31.3) are the Lorentzian Gauss--Codazzi
constraints.  For a hypersurface of a Ricci-flat Riemannian
four-manifold the Hamiltonian constraint reads
\(R-(\operatorname{tr}K)^2+|K|^2=0\), whose CMC reduction is
\(-8\Delta\psi+|\sigma|^2\psi^{-7}-\tfrac23\tau^2\psi^5=0\).  Its
nonlinearity has the opposite monotonicity, so the uniqueness and
monotone-iteration arguments of V29 do not transfer.  Parity (35.17) is
unaffected because both equations depend on the seed only through
\(|\sigma|^2\) and \(\tau^2\).  Which slice constraint a reduced
Euclidean measure must satisfy is a separate question, recorded here and
not decided.
\end{warning}

\subsection{Executable certificate}

The verifier

\[
 \texttt{scripts/verify\_sm\_v19\_v32\_pushforward\_reflection\_parity.py}
                                                                    \tag{35.21}
\]

works over the rationals with \(r=1/2\).  It checks that the reflected
Gram matrix of \(\{1,y_1,y_1^2\}\) has all principal minors nonnegative
for the Gaussian pair, evaluates (35.5) and (35.6) as \(-1/2\) and
\(-1\) with a negative principal minor in each case, evaluates the
equivariant local cubic image \(E[x_{-1}^3x_1^3]=9r+6r^3=21/4>0\),
verifies (35.11) for \(n,m\in\{0,1,2\}\) together with the failure of the
even assignment, records the coincident-plane odd correlation
\(E[d_{-1/2}d_{1/2}]=-(1-r)^2=-1/4\), and checks (35.16)--(35.17) on the
exact constant branch (31.15).  It emits

\[
 \texttt{artifacts/sm\_v19\_v32\_pushforward\_reflection\_parity.json}.
                                                                    \tag{35.22}
\]

The hashes are

\[
\begin{array}{c|l}
\hbox{object}&\hbox{SHA--256}\\ \hline
\hbox{verifier}&
\texttt{0FF4531F50C31D736BAB6AE91B837B7DCFA9DB86828DA606B3421F56BD91BD8B}\\
\hbox{canonical payload}&
\texttt{08194607F131F81068E1C5223E622868A1FBC70B3BB4984FA05384E8537B0CE0}.
\end{array}                                                     \tag{35.23}
\]

\begin{firewall}[Positivity transport is not a gravitational measure]
\label{fire:v19v32}
Theorem \ref{thm:v19v32-pushforward} transports reflection positivity
from a seed law that must already possess it; no interacting or
gravitational seed law with Einstein dynamics has been constructed.
Corollary \ref{cor:v19v32-slicewise} is a change of variables and does
not produce a positive transfer operator with the reduced Einstein
Hamiltonian.  Theorem \ref{thm:v19v32-sharp} concerns \(L^2\)-continuous
sharp-time restrictions only; distribution-valued odd fields are not
excluded.  The Euclidean-signature constraint of Warning
\ref{warn:v19v32-signature} has not been analyzed.
\end{firewall}

\begin{decision}[V33 acquisition]
\label{dec:v19v32-next}
V33 will place the constraint solver on a Euclidean time lattice.  The
seed will be a reflection-positive Gaussian Markov chain on sites, with
transverse-traceless link differences as the odd seed and a link scalar
as the mean curvature; the site reflection acts on links by the sign
rule (35.11).  The note will prove that the link-wise CMC map is
equivariant and positive-time local for the half-shifted link algebra,
transport reflection positivity by Theorem
\ref{thm:v19v32-pushforward}, compute the exact reflected link Gram, and
state which lattice-time limit of the link variables recovers the
singular coincident covariance of Corollary \ref{cor:v19v32-no-seed}.
\end{decision}

\begin{assessment}[V32 acquisition]
\label{ass:v19v32}
The measure-theoretic condition of Decision \ref{dec:v19v31-next} is
exact and sharp: equivariance plus positive-time locality transports
reflection positivity, and each hypothesis alone is insufficient.  The
CMC solver is analytic and even in its seed, hence equivariant with an
odd extrinsic curvature, and that parity is forced by the reflected form
of a real Euclidean time derivative.  The same parity makes a
continuous-time sharp-time seed vanish, so the nonvacuous test moves to
lattice time.  The programme also records that the slice constraint of a
Euclidean four-geometry has the opposite \(K^2\) sign.
\end{assessment}

\begin{record}[V32 append record]
\label{rec:v19v32}
V32 is appended to the validated V31 whose installed SHA--256 was
\texttt{2700B0136C14853505125A4F0C9E34A8893785326287B8FD2380407495D78205}.
The preamble gains the previously missing declaration
\(\backslash\)\texttt{newtheorem\{criterion\}}, used by V21; no proof text
is altered.  After installation only V32 is the active numeric SM note;
frozen archives and unrelated notes are unchanged.
\end{record}


\section{V33: link seeds on a Euclidean time lattice and transported positivity}
\label{sec:v19v33}

V32 showed that a sharp-time odd seed vanishes in continuous time, and
that the nonvacuous version of the constraint-solver test lives on a
Euclidean time lattice.  This note carries it out.  A stationary
reversible Markov chain on time sites is reflection positive; its link
differences are automatically odd; the link-wise CMC map is equivariant
and local; and Theorem \ref{thm:v19v32-pushforward} transports positivity
to the constrained data.  The exact link covariances are computed, their
continuum limit reproduces the singular coincident term of Corollary
\ref{cor:v19v32-no-seed}, and a barrier argument shows that the
constrained metric acquires a heavy tail from links of small mean
curvature.

\subsection{Reversible chains are site-reflection positive}

Let \(E\) be a Polish space, \(\pi\) a probability measure on \(E\), and
\(P\) a Markov kernel on \(E\) that is reversible with respect to
\(\pi\):

\[
 \pi(dx)P(x,dy)=\pi(dy)P(y,dx).                                 \tag{36.1}
\]

Let \(\mu\) be the law of the stationary two-sided chain
\((q_n)_{n\in\mathbb Z}\) with marginal \(\pi\) and transition \(P\),
on \(\Omega=E^{\mathbb Z}\).  Let \(\theta\) be the path reflection
\((\theta q)_n=q_{-n}\), and let \({\cal H}_+\) be the \(L^2(\mu)\)-closure
of the bounded functions of \((q_n)_{n\geq0}\).

\begin{theorem}[Site-reflection positivity of reversible chains]
\label{thm:v19v33-chain}
\(\theta_*\mu=\mu\), and \((\mu,\theta,{\cal H}_+)\) is a
reflection-positive triple.  For every bounded \(F\) measurable with
respect to \((q_n)_{n\geq0}\),

\[
 \langle F,F\rangle_\theta
 =E\Bigl[\,\bigl|E[F\mid q_0]\bigr|^2\Bigr]\geq0.                \tag{36.2}
\]
\end{theorem}

\begin{proof}
For \(n_1<\dots<n_k\), the joint law of \((q_{n_1},\dots,q_{n_k})\) is
\(\pi(dx_1)P^{n_2-n_1}(x_1,dx_2)\cdots P^{n_k-n_{k-1}}(x_{k-1},dx_k)\).
Iterating (36.1) reverses the order of the factors, so this law is
symmetric under \((x_1,\dots,x_k)\mapsto(x_k,\dots,x_1)\).  Applied to
the times \(-n_k<\dots<-n_1\), this gives \(\theta_*\mu=\mu\).

By the two-sided Markov property, conditionally on \(q_0\) the future
\((q_n)_{n\geq0}\) and the past \((q_{-n})_{n\geq0}\) are independent, and
by reflection invariance they have the same conditional law.  Since
\(\theta\) fixes \(q_0\), \(E[F\circ\theta\mid q_0]=E[F\mid q_0]\).
Hence

\[
 \langle F,F\rangle_\theta
 =E\bigl[\overline{F\circ\theta}\,F\bigr]
 =E\Bigl[\overline{E[F\circ\theta\mid q_0]}\,E[F\mid q_0]\Bigr]
 =E\Bigl[\bigl|E[F\mid q_0]\bigr|^2\Bigr].
\]

Boundedness of the form extends (36.2) to \({\cal H}_+\).
\end{proof}

\begin{corollary}[Gaussian product chains]
\label{cor:v19v33-gaussian}
Let \(E=\mathbb R^d\) and let the coordinates of \(q_n\) be independent
stationary Gaussian chains with

\[
 E[q^{(j)}_nq^{(j)}_m]=c_j\,r_j^{|n-m|},\qquad c_j>0,\ 0<r_j<1.
                                                                    \tag{36.3}
\]

Then \(\mu\) satisfies Theorem \ref{thm:v19v33-chain}, and for
\(n,m\geq0\) the reflected site Gram is the rank-one matrix
\(\langle q^{(j)}_n,q^{(j)}_m\rangle_\theta=c_jr_j^{\,n+m}\).
\end{corollary}

\begin{proof}
Each coordinate is the stationary chain of the Gaussian kernel with
mean \(r_jx\) and variance \(c_j(1-r_j^2)\), which is reversible for its
Gaussian invariant law; products of reversible kernels are reversible for
the product law.  The Gram entry is
\(E[q^{(j)}_{-n}q^{(j)}_m]=c_jr_j^{\,n+m}\).
\end{proof}

With the V23 rates, one takes for each retained TT Fourier mode
\(r_k=e^{-\varepsilon|k|}\) and \(c_k=1/(2c_g|k|)\); the homogeneous
mean-curvature seed needs its own stationary chain with some
\(r_0<1\), which is exactly the infrared prescription discussed in V24.

\subsection{Link seeds and the transported positivity}

Fix a lattice spacing \(\varepsilon>0\) and define the link variables

\[
 d_{n+1/2}:=\frac{q_{n+1}-q_n}{\varepsilon}\qquad(n\in\mathbb Z).
                                                                    \tag{36.4}
\]

\begin{lemma}[Induced link parity and locality]
\label{lem:v19v33-links}
Under the site reflection,

\[
 d_{n+1/2}\circ\theta=-d_{-n-1/2}.                                \tag{36.5}
\]

The link algebra \({\cal A}^{\rm link}_+\) of bounded functions of
\((d_{n+1/2})_{n\geq0}\) is contained in the site algebra of
\((q_n)_{n\geq0}\), so its closure \({\cal H}^{\rm link}_+\) lies in
\({\cal H}_+\).
\end{lemma}

\begin{proof}
\(d_{n+1/2}\circ\theta=(q_{-n-1}-q_{-n})/\varepsilon
=-(q_{-n}-q_{-n-1})/\varepsilon=-d_{-n-1/2}\).  A link with
\(n\geq0\) has both endpoints in \(\{0,1,2,\dots\}\).
\end{proof}

Now let \(E=V\oplus\mathbb R\), where \(V\) is a finite-dimensional space
of smooth transverse-traceless tensors on the flat torus spanned by
finitely many real Fourier modes, and let the chain be that of Corollary
\ref{cor:v19v33-gaussian} in a basis of \(V\oplus\mathbb R\).  Write
\(d_{n+1/2}=(\sigma_{n+1/2},\tau_{n+1/2})\).

\begin{theorem}[Transported positivity of link-wise constrained data]
\label{thm:v19v33-transport}
Almost surely every link seed lies in the set \({\cal U}\) of (35.14).
Define, link by link,

\[
 (g_{n+1/2},K_{n+1/2}):=\bigl(g,K\bigr)\bigl(\sigma_{n+1/2},\tau_{n+1/2}\bigr)
                                                                    \tag{36.6}
\]

by (31.2) with \(\psi=\Psi(\sigma_{n+1/2},\tau_{n+1/2})\), and let
\(\Theta'\) act on output paths by
\((g,K)_{n+1/2}\mapsto(g_{-n-1/2},-K_{-n-1/2})\).  Then the law
\(\nu\) of \((g_{n+1/2},K_{n+1/2})_{n\in\mathbb Z}\) is
\(\Theta'\)-invariant and reflection positive for the closure of the
bounded functions of \((g_{n+1/2},K_{n+1/2})_{n\geq0}\).  Every output
link is an exact solution of both vacuum constraints (29.1).
\end{theorem}

\begin{proof}
The link seed is a nondegenerate Gaussian vector in \(V\oplus\mathbb R\)
(its covariance is \(2c_j(1-r_j)/\varepsilon^2>0\) in each coordinate),
so \(\sigma_{n+1/2}\neq0\) and \(\tau_{n+1/2}\neq0\) almost surely for
each of countably many links.  On \({\cal U}\), \(\Psi\) is analytic by
Theorem \ref{thm:v19v32-analytic}, so (36.6) is measurable.  Equivariance
of the link-wise map is (35.17) combined with (36.5): the reflected seed
at link \(n+1/2\) is \(-(\sigma,\tau)_{-n-1/2}\), whose image is
\((g_{-n-1/2},-K_{-n-1/2})\).  Locality is Lemma
\ref{lem:v19v33-links}.  Theorem \ref{thm:v19v32-pushforward} applies with
the reflection-positive triple of Theorem \ref{thm:v19v33-chain}.  The
constraint statement is Lemma \ref{lem:v19v28-reduction} with Lemma
\ref{lem:v19v32-hoelder}.
\end{proof}

\subsection{Exact link covariances and the singular continuum limit}

\begin{proposition}[Link covariances of one chain component]
\label{prop:v19v33-covariance}
For a component with parameters \((c,r)\) in (36.3):
\begin{enumerate}
\item same-side covariance
\[
 E[d_{n+1/2}d_{m+1/2}]
 =\varepsilon^{-2}\times
 \begin{cases}
  2c(1-r), & n=m,\\
  -c(1-r)^2r^{|n-m|-1}, & n\neq m;
 \end{cases}                                                   \tag{36.7}
\]
\item reflected Gram, for \(n,m\geq0\),
\[
 \langle d_{n+1/2},d_{m+1/2}\rangle_\theta
 =\varepsilon^{-2}c(1-r)^2r^{\,n+m};                            \tag{36.8}
\]
\item window identity
\[
 \sum_{m=-M}^{M}E[d_{1/2}d_{m+1/2}]
 =2c\,\varepsilon^{-2}r^{M}(1-r)\longrightarrow0
 \qquad(M\to\infty).                                            \tag{36.9}
\]
\end{enumerate}
\end{proposition}

\begin{proof}
Expand \(E[(q_{n+1}-q_n)(q_{m+1}-q_m)]\) with (36.3).  For \(n=m\) it is
\(c(2-2r)\).  For \(n<m\) with \(j=m-n\geq1\) it is
\(c(r^{j}-r^{j-1}-r^{j+1}+r^{j})=-c\,r^{j-1}(1-r)^2\).  For (36.8) use
(36.5): the reflected Gram is
\(-E[d_{-n-1/2}d_{m+1/2}]\), and
\(E[(q_{-n}-q_{-n-1})(q_{m+1}-q_m)]
=c(r^{n+m+1}-r^{n+m}-r^{n+m+2}+r^{n+m+1})=-c\,r^{n+m}(1-r)^2\).
For (36.9), the window sum telescopes to
\(E[(q_1-q_0)(q_{M+1}-q_{-M})]/\varepsilon^2
=c(r^{M}-r^{M+1}-r^{M+1}+r^{M})/\varepsilon^2\).
\end{proof}

\begin{corollary}[The lattice limit reproduces the singular coincident covariance]
\label{cor:v19v33-limit}
Put \(r=e^{-\omega\varepsilon}\) and let \(n\varepsilon\to t\),
\(m\varepsilon\to s\) as \(\varepsilon\to0\).  Then
\begin{enumerate}
\item for \(t\neq s\), \(E[d_{n+1/2}d_{m+1/2}]\to-c\,\omega^2e^{-\omega|t-s|}\);
\item for \(t,s>0\),
      \(\langle d_{n+1/2},d_{m+1/2}\rangle_\theta\to c\,\omega^2e^{-\omega(t+s)}\),
      the square in (35.10);
\item the same-link variance is \(2c\omega/\varepsilon+O(1)\), and
      \(\varepsilon\,E[d_{n+1/2}^2]\to2c\omega\), the coefficient of
      \(\delta(t-s)\) in \(\partial_t\partial_sC\) for (35.7).
\end{enumerate}
\end{corollary}

\begin{proof}
Insert \(1-r=\omega\varepsilon+O(\varepsilon^2)\) and
\(r^{j}=e^{-\omega j\varepsilon}\) into (36.7)--(36.8).
\end{proof}

Identity (36.9) is the lattice statement that the derivative field has
zero spectral weight at zero frequency: the positive same-link atom
exactly balances the negative tail.  Theorem \ref{thm:v19v32-sharp}
says that no continuous covariance can do this; on the lattice the atom
is a legitimate random variable, which is why the link seed exists.

\subsection{Barriers and the heavy tail of the constrained metric}

\begin{proposition}[Two-sided constant barriers]
\label{prop:v19v33-barriers}
Let \((\sigma,\tau)\in{\cal U}\), \(a=|\sigma|^2\), \(b=\tfrac23\tau^2\),
and let \(\psi=\Psi(\sigma,\tau)\).  Then

\[
 \psi\leq\Bigl(\frac{\max a}{b}\Bigr)^{1/12}
 \quad\hbox{everywhere, and}\quad
 \psi\geq\Bigl(\frac{\min a}{b}\Bigr)^{1/12}
 \quad\hbox{if }\min a>0.                                        \tag{36.10}
\]
\end{proposition}

\begin{proof}
With \(M^{12}=\max a/b\), the constant \(M\) satisfies
\({\cal L}(M)=M^{-7}(bM^{12}-a)\geq0\).  Put \(w=\psi-M\).  The
mean-value form of \({\cal L}(\psi)-{\cal L}(M)\leq0\) is
\(-8\Delta w+c(x)w\leq0\) with \(c>0\) as in (32.13).  At a positive
maximum of \(w\) one has \(-8\Delta w\geq0\) and \(cw>0\), a
contradiction; hence \(w\leq0\).  The lower barrier
\(m^{12}=\min a/b\) has \({\cal L}(m)\leq0\), and the same argument
applied to \(m-\psi\) gives \(\psi\geq m\).
\end{proof}

\begin{theorem}[Moment threshold of the link metric]
\label{thm:v19v33-moments}
In the setting of Theorem \ref{thm:v19v33-transport}, suppose
\(\tau_{n+1/2}\) is independent of \(\sigma_{n+1/2}\) with a density
that is bounded and, near zero, bounded below by a positive constant.
Write \(\operatorname{tr}_\delta g=3\psi^4\).
\begin{enumerate}
\item For \(0<p<\tfrac32\),
      \(E\bigl[(\max_x\operatorname{tr}_\delta g_{n+1/2})^p\bigr]<\infty\).
\item If \(\Pr[\min_x|\sigma_{n+1/2}|^2>0]>0\), then for every
      \(x\) and every \(p\geq\tfrac32\),
      \(E\bigl[(\operatorname{tr}_\delta g_{n+1/2})(x)^p\bigr]=\infty\).
\end{enumerate}
\end{theorem}

\begin{proof}
By (36.10), \(\psi^4\leq(3\max|\sigma|^2/(2\tau^2))^{1/3}\).  The
Gaussian coefficients give \(E[(\max|\sigma|^2)^{p/3}]<\infty\), and
independence reduces the first claim to
\(E[|\tau|^{-2p/3}]<\infty\), which holds for a bounded density exactly
when \(2p/3<1\).  For the second claim, on the event
\(\min|\sigma|^2>0\) the lower barrier gives
\(\psi^4\geq(3\min|\sigma|^2/(2\tau^2))^{1/3}\), and
\(E[|\tau|^{-2p/3}]=\infty\) for \(2p/3\geq1\) when the density is bounded
below near zero.  Independence and positivity of the event finish the
proof.
\end{proof}

The hypothesis of the second claim holds, for instance, when \(V\)
contains the constant transverse-traceless tensors with positive
variance.  The mechanism is the nonlinear form of the V26 balance: a link
whose mean curvature is small but whose TT momentum is not must carry a
large conformal volume.  A Gaussian seed therefore produces a metric
whose distribution has a power tail of exponent \(3/2\).

\subsection{The output algebra is the seed algebra}

\begin{proposition}[No added dynamics]
\label{prop:v19v33-no-dynamics}
The link-wise map (36.6) is injective on \({\cal U}\) with the
continuous inverse
\(\psi=(\det g/\det\delta)^{1/12}\), \(\tau=\operatorname{tr}_gK\),
\(\sigma=\psi^{2}(K-\tfrac\tau3g)\).  Hence the positive-time output
algebra coincides with \({\cal A}^{\rm link}_+\), the reflected forms
agree by (35.3), and the Hilbert space and contraction semigroup
reconstructed from the constrained data are unitarily those
reconstructed from the link seed algebra.
\end{proposition}

\begin{proof}
The formulas invert (31.2) because \(\det(\psi^4\delta)=\psi^{12}\),
\(\operatorname{tr}_g(\psi^{-2}\sigma)=\psi^{-6}\operatorname{tr}_\delta\sigma=0\),
and \(\operatorname{tr}_g(\tfrac\tau3g)=\tau\).  Generated
\(\sigma\)-algebras of a bijection with measurable inverse coincide, and
the rest is Theorem \ref{thm:v19v32-pushforward}.
\end{proof}

Note also that \({\cal A}^{\rm link}_+\) is a proper subalgebra of the
site algebra: it contains no information on the absolute level of
\(q_0\).  The reconstructed space of the constrained data is therefore
the difference sector of the seed chain.

\subsection{Executable certificate}

The verifier

\[
 \texttt{scripts/verify\_sm\_v19\_v33\_lattice\_link\_pushforward.py}
                                                                    \tag{36.11}
\]

works over the rationals with \(\varepsilon=1\), modes
\((c,r)=(1,\tfrac12)\) and \((2,\tfrac14)\).  It checks the rank-one
site Gram of Corollary \ref{cor:v19v33-gaussian}, the closed forms
(36.7) for \(n,m\leq3\), the window identity (36.9) with \(M=5\)
(value \(1/32\)), the block rank-one reflected link Gram (36.8) for two
modes and \(n,m\in\{0,1\}\) with all principal minors nonnegative, the
parity rule (36.5) on a sample path, and the constant barriers (36.10)
with \(b=1\) and \(a\in\{1,4096\}\), giving \(M=2\) and \(m=1\).  It
emits

\[
 \texttt{artifacts/sm\_v19\_v33\_lattice\_link\_pushforward.json}.
                                                                    \tag{36.12}
\]

The hashes are

\[
\begin{array}{c|l}
\hbox{object}&\hbox{SHA--256}\\ \hline
\hbox{verifier}&
\texttt{3EC69FE37013B0784FBE579A0F1943E135EF58EB209B07D803036A15B3172FD4}\\
\hbox{canonical payload}&
\texttt{6AE0810EFD92B58A17AC4ED636178655A5CD7123B5F4223F5FBC3A36C9518B8E}.
\end{array}                                                     \tag{36.13}
\]

\begin{firewall}[A constrained free chain is not quantum gravity]
\label{fire:v19v33}
The seed of Theorem \ref{thm:v19v33-transport} is a Gaussian chain of
free TT modes with an ad hoc homogeneous chain; its reconstructed
semigroup is the free one, and Proposition
\ref{prop:v19v33-no-dynamics} shows the constraint solver does not
change it.  No interacting reduced Hamiltonian, no Standard Model
source, no lattice-spacing limit of the constrained data, and no
Lorentzian development are constructed.  The heavy tail of Theorem
\ref{thm:v19v33-moments} is a property of this seed, not a theorem
about a physical vacuum.  The slice constraint used is Lorentzian, as
recorded in Warning \ref{warn:v19v32-signature}.
\end{firewall}

\begin{decision}[V34 acquisition]
\label{dec:v19v33-next}
V34 will analyze the Euclidean-signature slice constraint of Warning
\ref{warn:v19v32-signature},
\(-8\Delta\psi+|\sigma|^2\psi^{-7}-\tfrac23\tau^2\psi^5=0\).  It will
prove the integrated balance, characterize the constant solutions, locate
the exact bifurcation values \(\tfrac23\tau^2\psi_*^4=\tfrac23|k|^2\) at
which the linearization about a constant solution acquires a kernel, and
prove by a Crandall--Rabinowitz argument in the even one-dimensional
sector that nonconstant positive solutions bifurcate.  The purpose is to
decide whether a Euclidean-signature constraint solver can be
single-valued, which Theorem \ref{thm:v19v32-pushforward} requires.
\end{decision}

\begin{assessment}[V33 acquisition]
\label{ass:v19v33}
The lattice-time test of the V31 decision is complete and positive:
reversible chains are site-reflection positive, link differences are
automatically odd, and the link-wise CMC solver transports positivity
exactly, with an explicit rank-one reflected link Gram whose continuum
limit is the singular coincident covariance.  The solver is a change of
variables, so it adds constraints but no dynamics, and it turns a
Gaussian seed into a metric with a power tail of exponent \(3/2\).  The
next question is the signature of the slice constraint itself.
\end{assessment}

\begin{record}[V33 append record]
\label{rec:v19v33}
V33 is appended to the validated V32 whose installed SHA--256 was
\texttt{D22BD4C238838126B0C67A29E70B1A1A08B9AA3D0E856818D55E5F4BC9117426}.
After installation only V33 is the active numeric SM note; frozen
archives and unrelated notes are unchanged.
\end{record}


\section{V34: the Euclidean-signature slice constraint is multivalued}
\label{sec:v19v34}

Warning \ref{warn:v19v32-signature} recorded that a hypersurface of a
Ricci-flat Riemannian four-manifold obeys the Gauss--Codazzi relation
with the opposite sign of the extrinsic terms.  Theorem
\ref{thm:v19v32-pushforward} needs a single-valued constraint solver.
This note shows that the Euclidean-signature CMC equation does not
provide one: the constant solution is nondegenerate below an explicit
threshold, but a nonconstant branch bifurcates from it, and the exact
third-order computation shows the branch turns back into the region
where the constant solution is nondegenerate.  Hence the same seed admits
several exact Euclidean constrained slices.  The Lorentzian equation of
V29 has none of this, because its linearization is strictly positive.

\subsection{The Euclidean CMC equation}

On the flat torus with normalized measure, and with
\(a=|\sigma|^2\geq0\), \(b=\tfrac23\tau^2>0\), the Euclidean slice
constraint \(R-(\operatorname{tr}K)^2+|K|^2=0\) becomes, through the
conformal identities (31.7),

\[
 {\cal L}_E(\psi):=-8\Delta\psi+a\,\psi^{-7}-b\,\psi^{5}=0,
 \qquad\psi>0.                                                   \tag{37.1}
\]

\begin{proposition}[Balance, degeneracies, and constant solutions]
\label{prop:v19v34-basic}
Every positive solution satisfies
\(\langle a\psi^{-7}\rangle=b\langle\psi^5\rangle\).  If exactly one of
\(a\not\equiv0\), \(b>0\) holds, there is no positive solution.  A
constant \(\psi\) solves (37.1) if and only if \(a\) is constant and
\(a=b\psi^{12}\); in that case the constant solution is unique.
\end{proposition}

\begin{proof}
Integrate (37.1).  The two degenerate cases give
\(\langle a\psi^{-7}\rangle=0\) or \(b\langle\psi^5\rangle=0\), both
impossible.  For constant \(\psi\), (37.1) reads
\(a=b\psi^{12}\), which forces \(a\) constant and determines
\(\psi\).
\end{proof}

\subsection{Linearization and the exact bifurcation values}

\begin{proposition}[Kernel of the linearization at a constant solution]
\label{prop:v19v34-kernel}
Let \(a\equiv s^2>0\) and \(\psi_*=(s^2/b)^{1/12}\).  Then

\[
 D{\cal L}_E(\psi_*)=-8\Delta-12\,b\,\psi_*^{4},                   \tag{37.2}
\]

whose kernel on the torus is nontrivial exactly when

\[
 12\,b\,\psi_*^{4}=8|k|^2,\quad\hbox{that is}\quad
 \tfrac23\tau^2\psi_*^4=\tfrac23|k|^2,                             \tag{37.3}
\]

for some nonzero dual-lattice vector \(k\).  For the Lorentzian
equation (32.2) the linearization at the same point is
\(-8\Delta+12b\psi_*^4\), which is strictly positive for every seed.
\end{proposition}

\begin{proof}
Differentiate (37.1): \(D{\cal L}_E(\psi)=-8\Delta-7a\psi^{-8}-5b\psi^4\).
At \(\psi_*\), \(a\psi_*^{-8}=b\psi_*^{4}\), which gives (37.2).  On
\(e^{ik\cdot x}\) the operator (37.2) acts by \(8|k|^2-12b\psi_*^4\),
whence (37.3).  The Lorentzian sign follows from (35.19).
\end{proof}

Below the first threshold, \(12b\psi_*^4<8|k_1|^2\) with \(k_1\) a
shortest nonzero dual vector, the constant solution is nondegenerate,
so it is locally unique by the implicit-function theorem.  This does
not exclude other solutions far from it, and the next theorem produces
them nearby in the parameter.

\subsection{A bifurcating nonconstant branch}

Take a torus whose \(z\)-period is \(2\pi\), and consider the
one-dimensional even sector: solutions depending on \(z\) only and even
in \(z\).  Normalize the constant seed so that \(\psi\equiv1\) solves
(37.1) for every \(b\): put \(a\equiv b\), that is
\(|\sigma|^2=b\), \(\tau^2=\tfrac32b\).  Then (37.1) reads

\[
 F(b,u):=-8u''+b\bigl[(1+u)^{-7}-(1+u)^5\bigr]=0,
 \qquad\psi=1+u,                                                  \tag{37.4}
\]

with the expansion

\[
 (1+u)^{-7}-(1+u)^5=-12u+18u^2-94u^3+O(u^4).                     \tag{37.5}
\]

\begin{theorem}[Crandall--Rabinowitz bifurcation at \(b_1=2/3\)]
\label{thm:v19v34-bifurcation}
Let \(X=C^{2,\alpha}_{\rm even}(\mathbb R/2\pi\mathbb Z)\) and
\(Y=C^{0,\alpha}_{\rm even}(\mathbb R/2\pi\mathbb Z)\).  There exist
\(\delta>0\) and real-analytic functions \(b:(-\delta,\delta)\to(0,\infty)\)
and \(v:(-\delta,\delta)\to X\) with \(b(0)=2/3\), \(v(0)=0\),
\(\langle v(s),\cos z\rangle=0\), such that

\[
 \psi_s:=1+s\cos z+s\,v(s)                                         \tag{37.6}
\]

is a positive nonconstant solution of (37.1) with data
\(a\equiv b(s)\), \(\tau^2=\tfrac32b(s)\), for every \(0<|s|<\delta\).
Near \((2/3,0)\) these, the constant solution, and their translates by
\(\pi\) in \(z\) are the only solutions in \(X\).  Moreover

\[
 b(s)=\frac23-\frac83s^2+O(s^4),\qquad
 \psi_s=1+s\cos z+s^2\Bigl(\frac34-\frac{\cos2z}4\Bigr)
 +s^3\,\frac{7}{24}\cos3z+O(s^4).                                  \tag{37.7}
\]
\end{theorem}

\begin{proof}
\(F:(0,\infty)\times\{u\in X:1+u>0\}\to Y\) is real-analytic, and
\(F(b,0)=0\).  Its linearization
\(D_uF(b,0)=-8\partial_z^2-12b\) acts on \(\cos kz\) by
\(8k^2-12b\).  At \(b_1=2/3\) the eigenvalue at \(k=1\) vanishes, the
eigenvalue at \(k=0\) is \(-8\), and those at \(k\geq2\) are
\(8k^2-8>0\); in the even sector \(\cos z\) spans a one-dimensional
kernel, and the range is the closed hyperplane orthogonal to
\(\cos z\), of codimension one, because the operator is self-adjoint
Fredholm.  The mixed derivative is
\(D_bD_uF(b_1,0)\cos z=-12\cos z\), which is not in the range.  The
Crandall--Rabinowitz theorem at a simple eigenvalue gives the curve
(37.6), its local uniqueness statement, and analyticity because \(F\)
is analytic.  Positivity holds for small \(|s|\).  Nonconstancy holds
because the \(\cos z\) coefficient of \(u\) is \(s\neq0\).  The
translates by \(\pi\) correspond to \(s\mapsto-s\).

For (37.7), write \(u=s\cos z+s^2u_2+s^3u_3+O(s^4)\) and
\(b=b_1+sb'+s^2b''+O(s^3)\) and insert (37.5) into (37.4).  At order
\(s\) the equation is \((-8\partial_z^2-12b_1)\cos z=0\).  At order
\(s^2\),

\[
 (-8\partial_z^2-12b_1)u_2-12b'\cos z+18b_1\cos^2z=0.              \tag{37.8}
\]

Projecting on \(\cos z\) gives \(-6b'+18b_1\langle\cos^3z\rangle=0\),
so \(b'=0\); solving on the complement with
\(\cos^2z=\tfrac12+\tfrac12\cos2z\) gives
\(u_2=\tfrac34-\tfrac14\cos2z\).  At order \(s^3\),

\[
 (-8\partial_z^2-12b_1)u_3-12b''\cos z+36b_1\cos z\,u_2
 -94b_1\cos^3z=0.                                                  \tag{37.9}
\]

Using \(\cos z\,u_2=\tfrac58\cos z-\tfrac18\cos3z\) and
\(\cos^3z=\tfrac34\cos z+\tfrac14\cos3z\), the source terms equal
\(-32\cos z-\tfrac{56}3\cos3z\).  The \(\cos z\) component forces
\(b''=-8/3\), and the \(\cos3z\) component gives
\(64\,u_3=\tfrac{56}3\cos3z\), that is \(u_3=\tfrac7{24}\cos3z\).
Evenness of \(b\) in \(s\) follows from the \(\pi\)-translation
symmetry, so the error in \(b\) is \(O(s^4)\).
\end{proof}

\begin{corollary}[No single-valued Euclidean solver]
\label{cor:v19v34-multivalued}
For every \(b\) in an interval \((\tfrac23-\eta,\tfrac23)\), the
Euclidean CMC equation with constant seed \(|\sigma|^2=b\),
\(\tau^2=\tfrac32b\) has at least three distinct positive solutions on the
torus with \(z\)-period \(2\pi\): the constant \(1\) and the two
nonconstant translates \(\psi_{\pm s}\) with \(b(s)=b\).  On this
interval the constant solution is nondegenerate.  Consequently the
Euclidean equation does not determine the slice as a function of the
seed: any single-valued Euclidean solver is a selection rule not
implied by (37.1), and Theorem \ref{thm:v19v32-pushforward} can be
applied only after such a rule is fixed and proved measurable and
equivariant.  A selection rule covariant under \(z\)-translation must
choose the constant slice for constant seeds, because \(\psi_s\) and
\(\psi_{-s}\) are translates of each other.
\end{corollary}

\begin{proof}
Since \(b''<0\), \(b(s)<2/3\) for small \(s\neq0\) and \(b\) is even
with \(b(s)\to2/3\) as \(s\to0\); the intermediate value theorem provides
\(s\) for each \(b\) in the interval, and \(\psi_s\neq\psi_{-s}\)
because their \(\cos z\) coefficients differ.  Nondegeneracy below
\(2/3\) is Proposition \ref{prop:v19v34-kernel}.  The remaining
statements are immediate: the pair \(\psi_{\pm s}\) is exchanged by
\(z\mapsto z+\pi\), so a translation-covariant single-valued choice at
a translation-invariant seed must be translation invariant, hence
constant.
\end{proof}

Contrast this with the Lorentzian equation, where uniqueness makes the
solver canonical and translation covariance, parity, and analyticity
are automatic consequences (Theorem \ref{thm:v19v32-analytic}).

\subsection{The one-dimensional first integral}

In the \(z\)-only sector with constant \(a=s^2\), (37.1) is the
Newtonian equation \(8\psi''=s^2\psi^{-7}-b\psi^5=-W'(\psi)\) with

\[
 W(\psi)=\frac{s^2}{6}\psi^{-6}+\frac b6\psi^{6},
 \qquad 4(\psi')^2+W(\psi)=E.                                      \tag{37.10}
\]

\(W\) is strictly convex on \((0,\infty)\), tends to infinity at both
ends, and has its unique minimum at \(\psi_*\) with
\(W''(\psi_*)=12b\psi_*^4\).  Every \(z\)-only solution therefore lies
on a closed level curve of (37.10), is periodic, and is bounded above and
below by the two roots of \(W=E\).  Small oscillations have frequency
\(\sqrt{12b\psi_*^4/8}\), which is (37.3) with \(|k|\) replaced by the
frequency.  The Lorentzian equation has \(W_L=-\tfrac{s^2}6\psi^{-6}-\tfrac b6\psi^6\),
concave, with a unique unstable equilibrium and no closed orbits, which
is the mechanical form of V29 uniqueness.

\subsection{Executable certificate}

The verifier

\[
 \texttt{scripts/verify\_sm\_v19\_v34\_euclidean\_slice\_bifurcation.py}
                                                                    \tag{37.11}
\]

implements exact cosine-series arithmetic over the rationals.  It checks
the Taylor coefficients (37.5) from binomial series, the eigenvalues
\(8k^2\mp12b_1\) for \(k\leq3\) in both signatures, the values
\(b_k=2k^2/3\), the vanishing of \(b'\), the series \(u_2\) and \(u_3\),
the value \(b''=-8/3\), the exact vanishing of the order-two and
order-three residuals (37.8)--(37.9) as cosine polynomials, and the
integrated balance at order two.  It emits

\[
 \texttt{artifacts/sm\_v19\_v34\_euclidean\_slice\_bifurcation.json}.
                                                                    \tag{37.12}
\]

The hashes are

\[
\begin{array}{c|l}
\hbox{object}&\hbox{SHA--256}\\ \hline
\hbox{verifier}&
\texttt{CAA7A238601EC8DB4A8A7E6C3B55C71C84113B7AC82EB3DF2608C830B8F9728A}\\
\hbox{canonical payload}&
\texttt{A967C5A1A722851E5B635F25D7EA81416E12799CAFA8733999766DAD1C90D5E5}.
\end{array}                                                     \tag{37.13}
\]

\begin{firewall}[Scope of the Euclidean multivaluedness]
\label{fire:v19v34}
Theorem \ref{thm:v19v34-bifurcation} is a local bifurcation result for
constant seeds in the even one-dimensional sector of one torus.  It does
not classify all Euclidean solutions, prove or disprove uniqueness for
small seeds far below the threshold, or decide which slice constraint a
reduced Euclidean gravitational measure must satisfy.  It does not
construct any measure.
\end{firewall}

\begin{decision}[V35 audit and V36 acquisition]
\label{dec:v19v34-next}
V35 is the compulsory companion audit.  V36 will then return to the
Lorentzian reduced picture, which V32--V34 single out as the only
single-valued option, and ask what dynamics the seed must carry: it will
compute the exact second-order expansion of the constrained slice volume
\(\langle\psi^6\rangle\) along the V31 branch and along the lattice
link seeds, and compare the linearized reduced evolution of the seed
with the V23 oscillator to state precisely where the free chain and the
Einstein evolution first differ.
\end{decision}

\begin{assessment}[V34 acquisition]
\label{ass:v19v34}
The signature question of V32 is decided at the level of solvability.
The Lorentzian CMC constraint has a unique analytic solution for every
admissible seed; the Euclidean one has an explicit subcritical pitchfork
at \(\tfrac23\tau^2\psi_*^4=\tfrac23|k|^2\), so a Euclidean slice solver
is multivalued precisely where its constant solution is nondegenerate.
A reduced-variable Euclidean measure for gravity must therefore solve
the Lorentzian constraint pathwise or carry a selection rule, and the
first option is the one already in use since V28.
\end{assessment}

\begin{record}[V34 append record]
\label{rec:v19v34}
V34 is appended to the validated V33 whose installed SHA--256 was
\texttt{9B0FD9315293ED1B2E85A0836AB7B0DA4B598AC950D01E589531FC42D6CEE597}.
After installation only V34 is the active numeric SM note; frozen
archives and unrelated notes are unchanged.
\end{record}


\section{V35: compulsory companion audit after the signature decision}
\label{sec:v19v35}

This is the required audit following V30.  Every active companion note in
the shared folder, including two series that did not exist at V30, was
inspected read-only before the seed-dynamics question announced in V34 is
taken up.

\begin{record}[V35 companion snapshot]
\label{rec:v19v35-snapshot}
The inspected files were

\[
\begin{array}{c|r|r}
\hbox{file and SHA--256}&\hbox{bytes}&\hbox{lines}\\ \hline
\texttt{Renewal\_Yang\_Mills\_Mass\_Gap\_Research\_Notes\_V30.tex}
&407339&10535\\
\multicolumn{3}{l}{
\texttt{C777BD838A65CB56FFE00EAF0283A1C00044D6FC83E47DBA65E49E8C2CC51790}}
\\
\texttt{Renewal\_NS\_Stability\_Research\_Notes\_V26.tex}
&278283&5319\\
\multicolumn{3}{l}{
\texttt{82C887F8C2C69E4F28FAD7C3DC8DE5B9099CB5A4BF431EBA1FFF3E91935978AD}}
\\
\texttt{Renewal\_GRH\_Cofinal\_Visibility\_Attack\_V50.tex}
&463785&14578\\
\multicolumn{3}{l}{
\texttt{E6EE171D8AB266216F348DF9FCD1FD533349DBF13CA85E4C19FA4DDF3212051A}}
\\
\texttt{Renewal\_YM\_Parallel\_Research\_Notes\_V5.tex}
&54174&1351\\
\multicolumn{3}{l}{
\texttt{306F1BB84CFA8A8D47EA569EC17E9545F9867C8D32B548EC9D9C444B4F3C0512}}.
\end{array}                                                     \tag{38.1}
\]

Since V30, Yang--Mills has advanced from V29 to V30, Navier--Stokes is
unchanged, and the GRH and parallel Yang--Mills series are new to this
audit.  During the sweep the GRH file was renamed at least once by its
own programme; the snapshot records the file present when the hash was
taken.
\end{record}

\subsection{Companion findings}

Yang--Mills V30 is itself an audit.  It reads SM V29 correctly, retains
the flat-torus CMC theorem as a possible later gravitational constraint
interface, imports nothing, and authorizes concurrent prime production
under its V29 gate.  Its firewall leaves the finite coefficient
reconstruction, interacting measure, regulator removal, reconstruction,
clustering, and mass gap open.  It also records that SM V25 listed
\(9119\) lines for the Yang--Mills V27 source that actually had
\(10020\); the hash there was correct, and no SM theorem depends on the
line count.

Navier--Stokes V26 is unchanged since V25 and V30: closed rank-one norms
of the first and second Oseen-kernel derivatives, with no stability or
regularity claim.

GRH V50 is a companion sweep and an exact Bernstein-coefficient ledger on
a height-twelve slab.  It imports only two proof disciplines from
Yang--Mills V20 and an older SM V12 finite-net theorem, and makes no
analytic claim relevant to constraint solving or reflection positivity.

Yang--Mills Parallel V5 closes its Programme I with a success stop on a
weak-coupling response margin and opens Programme II on sector tightness.
Its imported list cites companion SM V74--V78 of an earlier cycle for
abstract energy--entropy implications only.  It supplies no reflection
structure, transfer kernel, or constraint statement.

\begin{decision}[V35 audit decision]
\label{dec:v19v35-audit}
No companion theorem, numerical constant, or executable artifact is
imported.  The V34 direction stands: the seed of the constrained lattice
chain must be examined for dynamics.  V36 will prove exactly where the
constrained free chain of V33 and an Einstein development first
disagree, and will compute the second-order slice volume along the V31
branch.
\end{decision}

\begin{proof}
The V32--V34 results concern parity, pushforward positivity, and the
solvability of the two slice constraints.  None of the four companions
treats a gravitational constraint, a Euclidean time reflection acting on
canonical pairs, or a Markov seed law; the Yang--Mills V30 reading of SM
V29 confirms that the interface is deferred on their side as well.
Hence there is no cross-programme dependency to resolve, and the first
unresolved item on the SM side is the mismatch between the seed
dynamics and the Einstein evolution equations, which Proposition
\ref{prop:v19v33-no-dynamics} made explicit.
\end{proof}

\begin{assessment}[V35 acquisition]
\label{ass:v19v35}
The audit is current, covers two new series, and leaves the direction
unchanged.  The programme's own record since V30 is: an exact
positivity-transport criterion with sharp counterexamples, forced odd
parity of the extrinsic curvature with a sharp-time no-go, a lattice-time
realization with exact link Grams and a power-tail metric, and a proof
that the Euclidean-signature slice constraint is multivalued while the
Lorentzian one is canonical.
\end{assessment}

\begin{record}[V35 append record]
\label{rec:v19v35}
V35 is appended to the validated V34 whose installed SHA--256 was
\texttt{E82227C2505443D8B83D19CB81A66415C1DB640E052E67A4E9E982F53F6A9E5C}.
After installation only V35 is the active numeric SM note; the
companions, frozen archives, and all unrelated notes are unchanged.  The
next compulsory companion audit is V40.
\end{record}


\section{V36: stationarity against compact Einstein evolution}
\label{sec:v19v36}

V33 built a reflection-positive chain of exactly constrained slices and
V35 asked where that chain and an Einstein development first disagree.
The answer is: at the first step, and for a structural reason.  Two
exact identities of any Einstein development are violated in expectation
by every stationary, sign-symmetric chain of constrained slices: the
metric--curvature step pairing is strictly negative for Einstein
evolution but vanishes for the chain, and the mean curvature of CMC
slices strictly increases for Einstein evolution but is stationary for
the chain.  The second identity shows more: on compact slices, no
time-translation-invariant process can carry nontrivial Einstein
evolution at all.  The note closes with the second-order slice volume
along the V31 branch.

\subsection{Two exact identities of Einstein evolution}

Let a four-dimensional Ricci-flat manifold of either signature be
foliated by compact slices with induced metric \(g(t)\), extrinsic
curvature \(K(t)\), lapse \(N>0\), and shift \(X\), so that

\[
 \partial_tg_{ij}=-2NK_{ij}+({\cal L}_Xg)_{ij},
 \qquad
 \nabla_i\bigl(K^{ij}-\tau g^{ij}\bigr)=0,
 \qquad \tau:=\operatorname{tr}_gK.                              \tag{39.1}
\]

The first relation defines \(K\); the second is the Codazzi constraint,
identical in both signatures.

\begin{proposition}[Step pairing]
\label{prop:v19v36-pairing}
On every slice on which \(\tau\) is spatially constant,

\[
 \int\partial_tg_{ij}\,K^{ij}\,dV_g=-2\int N|K|_g^2\,dV_g\leq0,
                                                                    \tag{39.2}
\]

with equality if and only if \(K\equiv0\) on the slice.
\end{proposition}

\begin{proof}
Insert (39.1).  The shift term is
\(2\int\nabla_iX_jK^{ij}dV_g=-2\int X_j\nabla_iK^{ij}dV_g
=-2\int X_j\nabla^j\tau\,dV_g=0\), by the Codazzi constraint and
constancy of \(\tau\).  The lapse term gives (39.2).
\end{proof}

\begin{proposition}[Mean-curvature monotonicity on CMC slices]
\label{prop:v19v36-monotone}
For a Lorentzian vacuum development with the ADM evolution
\(\partial_tK_{ij}=N(R_{ij}-2K_{ik}K^k{}_j+\tau K_{ij})-\nabla_i\nabla_jN
+({\cal L}_XK)_{ij}\) and the Hamiltonian constraint (29.1),

\[
 \partial_t\tau=N|K|_g^2-\Delta_gN+X(\tau).                        \tag{39.3}
\]

If the slices have constant mean curvature, then

\[
 \operatorname{Vol}_g\,\frac{d\tau}{dt}=\int N|K|_g^2\,dV_g\geq0,   \tag{39.4}
\]

with equality if and only if \(K\equiv0\).  In particular \(\tau\) is
strictly increasing along every CMC vacuum development with \(K\not\equiv0\),
and no such development is time periodic.
\end{proposition}

\begin{proof}
Differentiate \(\tau=g^{ij}K_{ij}\):
\(\partial_t\tau=g^{ij}\partial_tK_{ij}-g^{ik}g^{jl}\partial_tg_{kl}K_{ij}\).
The first term is \(N(R-2|K|^2+\tau^2)-\Delta_gN+g^{ij}({\cal L}_XK)_{ij}\);
the second, by (39.1), is \(2N|K|^2-({\cal L}_Xg)^{ij}K_{ij}\).  Since
\(\tau\) is a scalar, \(X(\tau)=g^{ij}({\cal L}_XK)_{ij}-({\cal L}_Xg)^{ij}K_{ij}\).
Hence \(\partial_t\tau=N(R+\tau^2)-\Delta_gN+X(\tau)\), and the
Hamiltonian constraint \(R+\tau^2=|K|^2\) gives (39.3).  On a CMC slice
\(X(\tau)=0\), and integrating (39.3) over the closed slice removes
\(\Delta_gN\), which gives (39.4).  Strict positivity for \(K\not\equiv0\)
follows from \(N>0\); periodicity would force \(\tau\) to return to its
initial value.
\end{proof}

The Milne slicing of Minkowski space, \(g=t^2g_{\mathbb H^3}\),
\(N=1\), \(K=-g/t\), \(\tau=-3/t\), \(|K|^2=3/t^2\), confirms the sign
convention in (39.3)--(39.4) and is checked exactly in the certificate.

\subsection{Two exact defects of the stationary constrained chain}

Return to the constrained chain of Theorem \ref{thm:v19v33-transport}.
Its joint seed law is a centred Gaussian, hence invariant under the
global sign flip \(S\) that reverses every link seed simultaneously.
Under \(S\), every \(g_{n+1/2}\) is even and every \(K_{n+1/2}\) is odd,
by (35.17).

\begin{theorem}[Vanishing step pairing]
\label{thm:v19v36-step}
Let the joint law of the link seeds be invariant under \(S\).  For all
\(m,n\), with all contractions taken in the metric \(g_{n+1/2}\),

\[
 E\Bigl[\int(g_{m+1/2})_{ij}\,K_{n+1/2}^{ij}\,dV_{g_{n+1/2}}\Bigr]=0,
                                                                    \tag{39.5}
\]

whenever the integrand is integrable, which holds for the Gaussian
chain.  Consequently, for every bounded positive lapse
\(N=N(g_{n+1/2})\) and every shift \(X\),

\[
 D_n:=E\Bigl[\int\Bigl(\frac{g_{n+3/2}-g_{n+1/2}}{\varepsilon}
 +2NK_{n+1/2}-{\cal L}_Xg_{n+1/2}\Bigr)_{ij}K^{ij}_{n+1/2}\,dV\Bigr]
 =2E\Bigl[\int N|K_{n+1/2}|^2dV\Bigr]>0,                          \tag{39.6}
\]

whereas (39.2) makes the corresponding integrand vanish pathwise for
any Einstein development.
\end{theorem}

\begin{proof}
The integrand in (39.5) is the product of an \(S\)-even functional of
the seeds and an \(S\)-odd one, so it changes sign under \(S\), while
the law is \(S\)-invariant; its expectation is therefore zero.  For
integrability in the Gaussian chain, note that the trace-free part of
\(K\) is \(\delta\)-orthogonal to \(g_{m+1/2}\propto\delta\), so the
integrand equals \(\tau_{n+1/2}\int\psi_{m+1/2}^4\psi_{n+1/2}^{2}dx\);
the barriers (36.10) bound it by
\(C|\tau_{n+1/2}|^{2/3}|\tau_{m+1/2}|^{-2/3}\) times polynomial moments of
the TT coefficients, and H\"older's inequality with exponents \(5\) and
\(5/4\) reduces integrability to \(E|\tau|^{-5/6}<\infty\), which holds
for a Gaussian.  In (39.6) the difference term vanishes by (39.5), the
shift term vanishes pathwise by the Codazzi constraint exactly as in the
proof of Proposition \ref{prop:v19v36-pairing}, and the lapse term is
positive because \(K_{n+1/2}\neq0\) almost surely.
\end{proof}

\begin{theorem}[Vanishing mean-curvature increment]
\label{thm:v19v36-tau}
For any stationary chain of exactly constrained CMC slices,

\[
 E[\tau_{n+3/2}-\tau_{n+1/2}]=0,                                   \tag{39.7}
\]

whereas (39.4) gives, for every Einstein development with positive
lapse and \(K\not\equiv0\), a strictly positive increment on every time
interval.  In particular, no stationary process of exactly constrained
CMC slices with \(\Pr[K\not\equiv0]>0\) is an Einstein development in
expectation, for any choice of lapse and shift.
\end{theorem}

\begin{proof}
Stationarity gives equal marginals at the two links, hence (39.7).  The
Einstein increment is the time integral of (39.4), which is positive
when \(K\not\equiv0\) on a set of times of positive measure; for a
smooth development this holds as soon as \(K\not\equiv0\) initially.
\end{proof}

\begin{corollary}[Stationary or periodic compact CMC developments are flat]
\label{cor:v19v36-no-stationary}
Let a Lorentzian vacuum CMC development on compact slices have a mean
curvature that is constant in time, or periodic in time.  Then
\(K\equiv0\) on every slice, the slices are scalar flat, and on
\(\mathbb T^3\) they are flat.
\end{corollary}

\begin{proof}
By (39.4), \(\tau\) is nondecreasing, so constancy or periodicity forces
\(\int N|K|^2dV=0\) at every time, hence \(K\equiv0\).  The Hamiltonian
constraint (29.1) then gives \(R(g)=0\), and a scalar-flat metric on
\(\mathbb T^3\) is flat.
\end{proof}

\begin{remark}[Consequence for time-translation-invariant processes]
\label{rem:v19v36-stationary}
If a reflection-positive, time-translation-invariant process is to be
reconstructed into a theory whose classical configurations are
Lorentzian vacuum CMC developments on compact slices, those
configurations inherit stationarity of the mean curvature, and Corollary
\ref{cor:v19v36-no-stationary} restricts them to the flat, time-symmetric
datum.  Every configuration with nonzero graviton content is
non-stationary.  This is a statement about the classical limit; it is
not a theorem about the reconstructed quantum theory.
\end{remark}

This is the nonlinear meaning of the V26 obstruction.  There the
compact TT linearization was unstable because the constant lapse charge
\(\tfrac23\int\tau^2\) must balance the TT energy; here (39.4) says the
balancing mean curvature cannot stay constant in time.  The reduced
compact Einstein dynamics is intrinsically non-autonomous, and a
stationary seed can represent it only at zero graviton energy density.

\subsection{Second-order slice volume along the V31 branch}

\begin{proposition}[Slice volume and the kinetic identity]
\label{prop:v19v36-volume}
For every exact CMC solution of V29 with normalized measure,

\[
 \int|K|_g^2\,dV_g=\tau^2\langle\psi^6\rangle+8\langle|\nabla\psi|^2\rangle.
                                                                    \tag{39.8}
\]

Along the V31 branch under the balance (34.3),

\[
 \langle\psi_\lambda^6\rangle
 =1-\frac{28}{b}\,\lambda^2\bigl\langle|\nabla w_0|^2\bigr\rangle+O(\lambda^4),
                                                                    \tag{39.9}
\]

which for the Fourier example (34.24) is \(1-\tfrac7{128}\lambda^2+O(\lambda^4)\).
\end{proposition}

\begin{proof}
The Hamiltonian constraint gives
\(|K|_g^2=\tau^2+R(g)=\tau^2-8\psi^{-5}\Delta\psi\), and
\(dV_g=\psi^6dx\); integrating by parts yields (39.8).  For (39.9),
\(\psi_\lambda^6=1+6\lambda^2w+O(\lambda^4)\) with
\(\langle w\rangle=\overline w\) from (34.17), and
\(6\overline w=-\tfrac{28}b\langle|\nabla w_0|^2\rangle\); the Fourier
value is \(6\cdot(-7/768)=-7/128\).
\end{proof}

Thus spatially varying TT momentum lowers the volume of the balanced CMC
slice at second order, by an amount proportional to the Dirichlet energy
of the conformal response.

\subsection{Executable certificate}

The verifier

\[
 \texttt{scripts/verify\_sm\_v19\_v36\_evolution\_defect.py}       \tag{39.10}
\]

works over the rationals.  It checks (39.9) for the Fourier example,
constructs two exact constant constrained links with \(\psi=1\) and
\(\psi=2\) from \(\sigma=\operatorname{diag}(2,2,-4)\) and
\(64\sigma\) with \(\tau=6\), verifies both constraints on each, forms a
sign-symmetric correlated two-link seed law with same-sign probability
\(3/4\), and evaluates the step pairing (39.5) in both the \(\delta\)
and the \(g\) contraction as exactly zero while the Einstein side
\(-2E\int|K|^2dV\) equals \(-2340\).  It checks the gauge average for
\(X=\sin z\,\partial_z\) and the Milne identity
\(d\tau/dt=N|K|^2\) at three rational times.  It emits

\[
 \texttt{artifacts/sm\_v19\_v36\_evolution\_defect.json}.          \tag{39.11}
\]

The hashes are

\[
\begin{array}{c|l}
\hbox{object}&\hbox{SHA--256}\\ \hline
\hbox{verifier}&
\texttt{71A4AF560BBD5EA4F4EEB34EF070BC12A8CD8AC2422B9E22B8179015C54BABD6}\\
\hbox{canonical payload}&
\texttt{07737C26B8AD059E64E97ACBCFE0D0D945829CC938F83B593ED3578F8234FA5A}.
\end{array}                                                     \tag{39.12}
\]

\begin{firewall}[What the defects do and do not show]
\label{fire:v19v36}
Theorems \ref{thm:v19v36-step} and \ref{thm:v19v36-tau} show that the
V33 chain, and more generally any stationary sign-symmetric chain of
constrained slices, is not an Einstein development in expectation.  They
do not construct a non-stationary process, a reduced Hamiltonian, or an
infinite-volume limit.  Corollary \ref{cor:v19v36-no-stationary}
concerns the Lorentzian vacuum classical limit on compact slices; it
does not treat noncompact slices, matter sources, a cosmological
constant, or Euclidean-signature developments, for which (39.3) has
different signs.
\end{firewall}

\begin{decision}[V37 acquisition]
\label{dec:v19v36-next}
The stationary compact target is now excluded beyond the linearized
level, so V37 goes upstream to the volume dependence.  It will prove
that along the balanced branch the required mean curvature is fixed by
the TT energy density, \(\tfrac23\tau^2=\langle|\sigma|^2\rangle\), so
that at fixed total TT energy on a torus of side \(L\) the mean
curvature and the constrained conformal response vanish as
\(L\to\infty\), with explicit rates; and it will state the exact sense in
which the free TT chain becomes Einstein-consistent in that
zero-density limit.  The purpose is to decide whether the
reflection-positive stationary programme should move to infinite volume
or to non-stationary cosmological seeds.
\end{decision}

\begin{assessment}[V36 acquisition]
\label{ass:v19v36}
The seed-dynamics question is answered exactly.  A stationary
sign-symmetric chain of constrained slices violates the Einstein step
pairing and the CMC monotonicity identity in expectation at the very
first step, and the latter identity excludes every time-translation-invariant
process from representing nontrivial compact Einstein evolution.  Free
TT gravitons on a torus are therefore intrinsically cosmological beyond
linear order; stationarity survives only at zero energy density.
\end{assessment}

\begin{record}[V36 append record]
\label{rec:v19v36}
V36 is appended to the validated V35 whose installed SHA--256 was
\texttt{3043E0842C1DDD850C3CA154B8A6EC894593033D0492C5AB96BA195F8E755CB8}.
After installation only V36 is the active numeric SM note; frozen
archives and unrelated notes are unchanged.
\end{record}


\section{V37: volume dependence of the balanced branch and the two regimes}
\label{sec:v19v37}

V36 showed that stationarity is compatible with compact Einstein
evolution only at zero graviton energy density.  This note makes the
volume dependence exact.  On a torus of side \(L\), the balanced branch
of V31 behaves in two opposite ways: at fixed total TT energy the mean
curvature vanishes like \(L^{-3/2}\) and the conformal response converges
to the asymptotically flat maximal-slice response; at fixed energy
density the homogeneous conformal response grows like \(L^2\), so the
small-amplitude branch through flat data shrinks to zero radius.  The
first regime is where a stationary reflection-positive seed is
consistent; the second is cosmological.  A balanced odd mean curvature is
then added to the lattice chain without losing reflection positivity.

\subsection{Exact scaling on the torus of side \(L\)}

Let \(\mathbb T^3_L=(\mathbb R/L\mathbb Z)^3\), and now use unnormalized
integrals.  For a smooth TT tensor \(\sigma_0\), write

\[
 a=|\sigma_0|^2,\qquad E_\sigma:=\int_{\mathbb T^3_L}a\,dx,\qquad
 b=\langle a\rangle=\frac{E_\sigma}{L^3},                          \tag{40.1}
\]

and let \(w_0^L\) be the mean-zero solution of
\(-8\Delta w_0^L=a-\langle a\rangle\).

\begin{proposition}[Balanced branch in unnormalized form]
\label{prop:v19v37-scaling}
On \(\mathbb T^3_L\), the balance (34.3), the mean shift (34.17), and
the second-order volume (39.9) read

\[
 \tau_0^2=\frac{3E_\sigma}{2L^3},\qquad
 \overline w_L=-\frac{7}{12E_\sigma}\int a\,w_0^L\,dx,\qquad
 \langle\psi_\lambda^6\rangle
 =1-\frac{7\lambda^2}{2E_\sigma}\int a\,w_0^L\,dx+O(\lambda^4),      \tag{40.2}
\]

and \(\int a\,w_0^L\,dx=8\int|\nabla w_0^L|^2\,dx\).
\end{proposition}

\begin{proof}
Insert \(b=E_\sigma/L^3\) and \(\langle aw_0\rangle=L^{-3}\int aw_0^L\)
into (34.3), (34.17), and (39.9); the last identity is (34.21) in
unnormalized form.
\end{proof}

\subsection{Fixed total energy: convergence to the asymptotically flat response}

Let \(G_1\) be the mean-zero Green function of \(-\Delta\) on the unit
torus, \(-\Delta G_1=\delta-1\), and \(G_L(x):=L^{-1}G_1(x/L)\), so that
\(-\Delta G_L=\delta-L^{-3}\) on \(\mathbb T^3_L\).

\begin{lemma}[Structure of the periodic Green function]
\label{lem:v19v37-green}
\(H_1(x):=G_1(x)-\dfrac1{4\pi|x|}\) extends to a smooth function on
\(\{|x|<\tfrac12\}\).  Consequently, for \(|x|\leq L/4\),

\[
 G_L(x)=\frac1{4\pi|x|}+H_L(x),\qquad
 |H_L(x)|\leq\frac{C_H}{L},\qquad
 C_H:=\sup_{|y|\leq1/4}|H_1(y)|<\infty.                             \tag{40.3}
\]
\end{lemma}

\begin{proof}
On \(\{|x|<\tfrac12\}\), \(G_1\) has its only singularity at \(0\), and
\(-\Delta H_1=(\delta-1)-\delta=-1\) in the sense of distributions.  A
distributional solution of \(\Delta H_1=1\) is smooth by elliptic
regularity.  The scaling \(G_L(x)=L^{-1}G_1(x/L)\) gives
\(H_L(x)=L^{-1}H_1(x/L)\), and \(|x|\leq L/4\) means \(|x/L|\leq1/4\).
\end{proof}

\begin{theorem}[Infinite-volume limit of the conformal response]
\label{thm:v19v37-limit}
Let \(a\geq0\) be smooth with support in a ball of radius \(L/8\),
and set

\[
 U_\infty:=\frac1{32\pi}\iint\frac{a(x)a(y)}{|x-y|}\,dx\,dy.          \tag{40.4}
\]

Then

\[
 \Bigl|\int a\,w_0^L\,dx-U_\infty\Bigr|\leq\frac{C_HE_\sigma^2}{8L}.
                                                                    \tag{40.5}
\]

Hence, as \(L\to\infty\) at fixed \(a\),

\[
 \tau_0^2=\frac{3E_\sigma}{2L^3}\to0,\qquad
 \overline w_L\to-\frac{7U_\infty}{12E_\sigma},\qquad
 \langle\psi_\lambda^6\rangle-1\to-\frac{7\lambda^2U_\infty}{2E_\sigma}
 \ \hbox{at order }\lambda^2,                                       \tag{40.6}
\]

and \(w_0^L\to w_0^\infty:=\tfrac1{32\pi}\int a(y)|x-y|^{-1}dy\)
uniformly on compact sets, which is the second-order response of the
asymptotically flat maximal-slice Lichnerowicz equation
\(-8\Delta\psi=|\sigma|^2\psi^{-7}\), \(\psi\to1\) at infinity.
\end{theorem}

\begin{proof}
Since \(G_L\) has mean zero, \(w_0^L(x)=\tfrac18\int G_L(x-y)a(y)\,dy\)
is the mean-zero solution.  Therefore
\(\int aw_0^L=\tfrac18\iint a(x)G_L(x-y)a(y)\).  For \(x,y\) in the
support, \(|x-y|\leq L/4\), so (40.3) applies and the difference from
\(U_\infty\) is \(\tfrac18\iint a(x)H_L(x-y)a(y)\), bounded by
\(C_HE_\sigma^2/(8L)\).  Equations (40.6) follow from (40.2).  The local
uniform convergence of \(w_0^L\) is (40.3) applied inside the integral,
and \(w_0^\infty\) solves \(-8\Delta w_0^\infty=a\) on \(\mathbb R^3\)
with decay at infinity, which is the order-\(\lambda^2\) equation of the
asymptotically flat problem.
\end{proof}

The limiting torus branch is therefore the asymptotically flat
maximal-slice solution up to the constant conformal factor
\((1+\lambda^2\overline w_\infty)^4\), which is a global rescaling of the
flat background and carries no local information.

\begin{corollary}[Stationarity violation on the branch]
\label{cor:v19v37-rate}
For the Lorentzian development of the V31 data with unit lapse, the
initial mean-curvature rate (39.4) is

\[
 \frac{d\tau}{dt}\Big|_{t=0}
 =\lambda^2\tau_0^2+O(\lambda^4)
 =\frac{3\lambda^2E_\sigma}{2L^3}+O(\lambda^4).                      \tag{40.7}
\]

The relative change of the mean curvature over a fixed time \(T\) at
this rate is \(\lambda T\tau_0+O(\lambda^3)\), which vanishes like
\(L^{-3/2}\) at fixed total energy.
\end{corollary}

\begin{proof}
By (39.8), \(\int|K_\lambda|^2dV=\lambda^2\tau_0^2\int\psi_\lambda^6+8\int|\nabla\psi_\lambda|^2\),
and \(\nabla\psi_\lambda=O(\lambda^2)\), \(\int\psi_\lambda^6=L^3(1+O(\lambda^2))\).
Divide by the volume and use (40.2).  The initial mean curvature is
\(\lambda\tau_0\).
\end{proof}

\subsection{Fixed energy density: the cosmological regime}

\begin{proposition}[Growth of the homogeneous response at fixed density]
\label{prop:v19v37-density}
Let \(a=1+\cos(kz)\) with \(k=4\pi/L\) on a torus of \(z\)-period
\(L\), so that \(b=\langle a\rangle=1\) for every \(L\).  Then

\[
 w_0^L=\frac{\cos kz}{8k^2},\qquad
 \langle aw_0^L\rangle=\frac1{16k^2},\qquad
 \overline w_L=-\frac{7}{192k^2}=-\frac{7L^2}{3072\pi^2}.             \tag{40.8}
\]

Hence the second-order volume coefficient \(6\overline w_L\) grows like
\(L^2\): the second-order term of (34.15) is of order one as soon as
\(\lambda^2\) is of order \(k^2\propto L^{-2}\).  The expansion is
therefore not uniform in \(L\), and at fixed density the small-amplitude
description through flat data covers only amplitudes
\(\lambda=O(1/L)\).
\end{proposition}

\begin{proof}
Direct substitution into (34.16)--(34.17) with \(b=1\): the Poisson
energy identity gives \(\langle aw_0\rangle=\langle\cos^2kz\rangle/(8k^2)\),
and \(\lambda^2\overline w_L\) is of order one exactly when
\(\lambda^2k^{-2}\) is.
\end{proof}

At fixed density the required mean curvature does not vanish
(\(\tau_0^2=\tfrac32\) here), the CMC monotonicity rate (39.4) stays of
order one, and the small-amplitude balanced branch about flat data
disappears.  This is the regime of Corollary
\ref{cor:v19v36-no-stationary}: a finite-density graviton gas on a
torus is a cosmology, not a perturbation of a static vacuum.

\subsection{A balanced odd mean curvature for the lattice chain}

\begin{proposition}[Balanced chain remains reflection positive]
\label{prop:v19v37-balanced}
In the setting of Theorem \ref{thm:v19v33-transport}, let \(\phi\) be an
additional independent scalar chain of Corollary
\ref{cor:v19v33-gaussian} and define, link by link,

\[
 \tau_{n+1/2}:=\operatorname{sgn}(\phi_{n+1}-\phi_n)
 \sqrt{\tfrac32\langle|\sigma_{n+1/2}|^2\rangle}.                    \tag{40.9}
\]

Then \(\tau_{n+1/2}\) is an almost surely nonzero, measurable, odd
function of the link seeds, the link-wise constrained data built from
\((\sigma_{n+1/2},\tau_{n+1/2})\) satisfy the conclusions of Theorem
\ref{thm:v19v33-transport}, and for a spatially constant seed the
conformal factor is identically one.
\end{proposition}

\begin{proof}
The sign factor is odd under the induced link parity (36.5) and the
square root is even, so (40.9) is odd; it vanishes only on the null
event \(\phi_{n+1}=\phi_n\).  The map from link seeds to
\((\sigma,\tau)\) is link-local and equivariant, so composing it with the
CMC map preserves the hypotheses of Theorem
\ref{thm:v19v32-pushforward}.  For constant \(\sigma\), (31.9) gives
\(\psi_*^{12}=3|\sigma|^2/(2\tau^2)=1\).
\end{proof}

For nonconstant links the amplitude of \(\sigma_{n+1/2}\) is of order
\(\varepsilon^{-1/2}\) by (36.7), outside the small-amplitude branch;
the balanced chain is therefore exactly constrained and reflection
positive, but its conformal response is not controlled by (40.2).

\subsection{Executable certificate}

The verifier

\[
 \texttt{scripts/verify\_sm\_v19\_v37\_volume\_scaling.py}         \tag{40.10}
\]

checks (40.8) exactly for \(k^2\in\{1,\tfrac14,\tfrac1{16}\}\),
including the Poisson energy identity and the \(L^2\) growth of
\(\overline w_L\); the fixed-energy relations \(b=E_\sigma/L^3\),
\(\tau_0^2=3E_\sigma/(2L^3)\), and the \(L\)-independence of the
coefficient of \(\int aw_0^L\) in \(\overline w_L\) for \(E_\sigma=6\)
and \(L\in\{2,4,8\}\); and the \(1/L\) scaling of the bound (40.5).  It
emits

\[
 \texttt{artifacts/sm\_v19\_v37\_volume\_scaling.json}.            \tag{40.11}
\]

The hashes are

\[
\begin{array}{c|l}
\hbox{object}&\hbox{SHA--256}\\ \hline
\hbox{verifier}&
\texttt{9B0D8687D802AE2DDF57FA26C8BF0A4C6D3A75AAFBC46E72FFD2777F5AB7752D}\\
\hbox{canonical payload}&
\texttt{73733B7CA558A595FF6ABD09883BB018D3BDCE47C3DE8BA7B7E30DA0EEA037C7}.
\end{array}                                                     \tag{40.12}
\]

\begin{firewall}[Two regimes, no measure]
\label{fire:v19v37}
Theorem \ref{thm:v19v37-limit} is a classical second-order statement
about the balanced branch; it does not construct an infinite-volume
measure, control the large-amplitude chain, or treat the asymptotically
flat problem beyond its order-\(\lambda^2\) equation.  Proposition
\ref{prop:v19v37-density} shows only that the small-amplitude branch
fails at fixed density; it does not describe the cosmological solutions.
No Standard Model coupling appears.
\end{firewall}

\begin{decision}[V38 acquisition]
\label{dec:v19v37-next}
The reflection structure itself points to the non-stationary regime:
an odd mean curvature vanishes on the reflection plane, which is exactly
a moment of time symmetry, and the Einstein monotonicity (39.4) makes
\(\tau\) increase through zero there.  V38 will prove that a
time-inhomogeneous Markov chain whose law is invariant under the site
reflection is still reflection positive, so that stationarity can be
dropped from Theorem \ref{thm:v19v33-chain}; it will show that the
odd-parity link mean curvature of such a chain is a bounce, negative
before and positive after the plane; and it will construct the
reconstructed Hilbert space at the plane with its non-autonomous family
of positive contractions.  This replaces the excluded stationary
compact target by a reflection-positive compact cosmology.
\end{decision}

\begin{assessment}[V37 acquisition]
\label{ass:v19v37}
The volume dependence is exact and separates two regimes.  At fixed
total energy the balanced branch converges to the asymptotically flat
maximal-slice response with an \(O(E_\sigma^2/L)\) error, the required
mean curvature vanishes like \(L^{-3/2}\), and stationarity is restored
in the limit; at fixed density the homogeneous conformal response grows
like \(L^2\) and the branch through flat data disappears.  The
reflection-positive stationary programme is thus an infinite-volume,
zero-density statement, while compact finite-density gravity must be
treated as a reflection-symmetric bounce.
\end{assessment}

\begin{record}[V37 append record]
\label{rec:v19v37}
V37 is appended to the validated V36 whose installed SHA--256 was
\texttt{6399AC240E7ECE901485500754222667F9C4378C1F23B4E46D4A89C12E926AEC}.
After installation only V37 is the active numeric SM note; frozen
archives and unrelated notes are unchanged.
\end{record}


\section{V38: reflection-symmetric bounce chains and reconstruction at the plane}
\label{sec:v19v38}

V36 excluded stationary processes as carriers of nontrivial compact
Einstein evolution, and V37 identified the fixed-density regime as
cosmological.  This note drops stationarity.  A time-inhomogeneous
Markov chain whose law is invariant under the site reflection is still
reflection positive, with the same conditional-expectation formula as
V33.  Its odd link mean curvature is a bounce: it has opposite means on
the two sides of the plane and its expected increments are even about
the plane, which is exactly the pattern the Einstein monotonicity
(39.4) demands.  Both V36 obstructions are consequences of stationarity
or global sign symmetry and disappear.  The reconstructed Hilbert space
at the plane is identified exactly as \(L^2\) of the time-zero marginal,
and the non-autonomous evolution is a family of positivity-preserving
contractions between the time-\(n\) spaces.

\subsection{Reflection-symmetric chains}

Let \(E\) be Polish, \(\pi_0\) a probability measure on \(E\), and
\((P_n)_{n\geq0}\) Markov kernels on \(E\).  Construct a law \(\mu\) on
\(E^{\mathbb Z}\) as follows: draw \(q_0\sim\pi_0\); conditionally on
\(q_0\), let \((q_n)_{n\geq0}\) be the Markov chain started at \(q_0\)
with transitions \(P_0,P_1,\dots\), and let \((q_{-n})_{n\geq0}\) be an
independent copy of the same chain started at the same \(q_0\).  Call
\(\mu\) the reflection-symmetric chain of \((\pi_0,(P_n))\).  Write
\(\pi_n\) for the marginal of \(q_n\) and \(\theta\) for the site
reflection.

\begin{theorem}[Reflection positivity without stationarity]
\label{thm:v19v38-chain}
\(\theta_*\mu=\mu\), and for every bounded \(F\) measurable with respect
to \((q_n)_{n\geq0}\),

\[
 \langle F,F\rangle_\theta=E\Bigl[\bigl|E[F\mid q_0]\bigr|^2\Bigr]\geq0.
                                                                    \tag{41.1}
\]

Thus \((\mu,\theta,{\cal H}_+)\) is a reflection-positive triple.  The
process is a two-sided Markov chain.
\end{theorem}

\begin{proof}
Reflection invariance holds by construction: \(\theta\) exchanges the two
conditionally independent copies, which have the same conditional law.
Given \(q_0\), past and future are independent and \(\theta\) fixes
\(q_0\), so \(E[F\circ\theta\mid q_0]=E[F\mid q_0]\), and the computation
in the proof of Theorem \ref{thm:v19v33-chain} gives (41.1).  The
forward Markov property at times \(n\geq0\) is built in; at times
\(n<0\) it holds because the time reversal of a Markov chain is Markov
and the future beyond \(0\) depends on the past only through
\(q_0\).
\end{proof}

Stationarity and reversibility, used in Theorem \ref{thm:v19v33-chain}
only to make the law reflection invariant, are thus not needed.

\begin{corollary}[Gaussian bounce chain]
\label{cor:v19v38-gaussian}
Let \(E=\mathbb R^d\), \(\pi_0=N(0,\Sigma_0)\), and
\(P_n(q,\cdot)=N(A_nq+b_n,C_n)\).  Then Theorem
\ref{thm:v19v38-chain} applies.  With link variables (36.4), the
induced parity (36.5) holds, and for \(n,m\geq0\)

\[
 \langle d_{n+1/2},d_{m+1/2}\rangle_\theta
 =E\bigl[m_n(q_0)\,m_m(q_0)\bigr],
 \qquad m_n(q_0):=E[d_{n+1/2}\mid q_0],                              \tag{41.2}
\]

a Gram matrix.  In one dimension with \(A_n=r\), \(\Sigma_0=v_0\),

\[
 m_n(q_0)=\frac{r^n(r-1)q_0+t_n}{\varepsilon},\qquad
 t_n:=E[d_{n+1/2}]=\frac{c_{n+1}-c_n}{\varepsilon},\qquad
 c_{n+1}=rc_n+b_n,\ c_0=0,                                          \tag{41.3}
\]

so that

\[
 \langle d_{n+1/2},d_{m+1/2}\rangle_\theta
 =\frac{r^{n+m}(1-r)^2v_0+\varepsilon^2t_nt_m}{\varepsilon^2}.        \tag{41.4}
\]
\end{corollary}

\begin{proof}
Equation (41.2) is (41.1) polarized, with \(F=d_{n+1/2}\) and the mirror
identity \(E[d_{-n-1/2}\mid q_0]=-m_n(q_0)\).  For (41.3),
\(E[q_n\mid q_0]=r^nq_0+c_n\).  Then (41.4) is
\(E[m_nm_m]\) with \(E[q_0]=0\), \(E[q_0^2]=v_0\).
\end{proof}

\subsection{The odd mean curvature is a bounce}

Let the link mean curvature \(\tau_{n+1/2}\) be any odd link-local
function of the seeds, for instance \(\phi_{n+1}-\phi_n\) for a scalar
component or the balanced choice (40.9).

\begin{proposition}[Bounce structure]
\label{prop:v19v38-bounce}
For a reflection-symmetric chain,

\[
 E[\tau_{-n-1/2}]=-E[\tau_{n+1/2}],
 \qquad
 E[\tau_{n+3/2}-\tau_{n+1/2}]=E[\tau_{-n-1/2}-\tau_{-n-3/2}]
 \qquad(n\geq0).                                                    \tag{41.5}
\]

Hence the expected mean curvature is odd about the plane and its
expected increments are even about it.  If the increments are positive
on the future side, the expected mean curvature is negative before the
plane, zero on it in the sense of the limit of (35.13), and positive
after it, increasing throughout: the pattern required by the Einstein
identity (39.4).  In the Gaussian bounce chain with
\(\tau_{n+1/2}=d_{n+1/2}\), \(r=\tfrac12\), \(b_n=n+1\), the future link
means are \(1,\tfrac32,\tfrac74,\tfrac{15}8,\tfrac{31}{16}\) with
increments \(\tfrac12,\tfrac14,\tfrac18,\tfrac1{16}\).
\end{proposition}

\begin{proof}
Both identities are \(\theta\)-invariance of \(\mu\) combined with the
oddness \(\tau_{n+1/2}\circ\theta=-\tau_{-n-1/2}\).  The numbers follow
from (41.3).
\end{proof}

\begin{proposition}[The two V36 obstructions are absent]
\label{prop:v19v38-lifted}
For a reflection-symmetric chain of constrained slices:
\begin{enumerate}
\item the mean-curvature increment identity (39.7) is not implied; a
      chain with positive expected increments on the future side
      exists by Proposition \ref{prop:v19v38-bounce};
\item the step-pairing identity (39.5) is not implied, because a drift
      \(b_n\neq0\) breaks the global sign symmetry \(S\) while preserving
      reflection symmetry.
\end{enumerate}
Reflection positivity is retained in both cases by Theorem
\ref{thm:v19v38-chain} and Theorem \ref{thm:v19v32-pushforward}.
\end{proposition}

\begin{proof}
Theorem \ref{thm:v19v36-tau} used stationarity and Theorem
\ref{thm:v19v36-step} used \(S\)-invariance; neither hypothesis holds
for the drifted bounce chain, whose law is nevertheless
\(\theta\)-invariant by construction.  The link-wise CMC map remains
equivariant and local, so the pushforward criterion applies.
\end{proof}

This removes the parity and stationarity obstructions only.  Whether a
reflection-symmetric chain can be chosen so that its expected increments
equal the Einstein rates (39.4) and (39.2) is a dynamical question not
addressed here.

\subsection{Reconstruction at the plane}

\begin{theorem}[Hilbert space and non-autonomous contractions]
\label{thm:v19v38-reconstruction}
Let \(\mu\) be a reflection-symmetric chain and let \({\cal H}_+\) be the
closure of the bounded functions of \((q_n)_{n\geq0}\).
\begin{enumerate}
\item The null space of \(\langle\cdot,\cdot\rangle_\theta\) on
      \({\cal H}_+\) is \(\{F:E[F\mid q_0]=0\}\), and
      \(F\mapsto E[F\mid q_0]\) induces a unitary from the completed
      quotient \({\cal H}\) onto \(L^2(E,\pi_0)\).
\item For \(n\geq0\) and \(f\in L^2(\pi_n)\), the class of \(f(q_n)\) in
      \({\cal H}\) is \(P_0P_1\cdots P_{n-1}f\), where each
      \(P_j:L^2(\pi_{j+1})\to L^2(\pi_j)\) is a positivity-preserving
      contraction with \(P_j1=1\).
\item For \(f\in L^2(\pi_n)\), \(g\in L^2(\pi_m)\),
\[
 \langle f(q_n),g(q_m)\rangle_\theta
 =\bigl\langle P_0\cdots P_{n-1}f,\;P_0\cdots P_{m-1}g\bigr\rangle_{L^2(\pi_0)}.
                                                                    \tag{41.6}
\]
\end{enumerate}
In the stationary reversible case all \(P_j\) equal one self-adjoint
contraction \(P\) on \(L^2(\pi)\), and (41.6) reads
\(\langle f,P^{n+m}g\rangle\).
\end{theorem}

\begin{proof}
Polarizing (41.1) gives
\(\langle F,G\rangle_\theta=\langle E[F\mid q_0],E[G\mid q_0]\rangle_{L^2(\pi_0)}\),
which identifies the null space and makes the induced map isometric;
it is onto because \(f(q_0)\mapsto f\).  The tower property gives
\(E[f(q_n)\mid q_0]=(P_0\cdots P_{n-1}f)(q_0)\); each \(P_j\) is a
conditional expectation restricted to functions of \(q_{j+1}\), hence a
positivity-preserving unital contraction in \(L^2\) by conditional
Jensen.  Equation (41.6) is the polarized form once more.  In the
stationary reversible case \(P\) is self-adjoint on \(L^2(\pi)\), and
\(\langle Pf,Pg\rangle\) with \(n=m=1\) equals \(\langle f,P^2g\rangle\).
\end{proof}

The physical state space of a reflection-positive compact cosmology is
therefore \(L^2\) of the law of the time-symmetric slice, and its
evolution is the non-autonomous family \((P_n)\).  Self-adjointness and
positivity of a single transfer operator, which encode a
time-independent Hamiltonian, are exactly the properties that
Corollary \ref{cor:v19v36-no-stationary} forbids for nontrivial compact
vacuum gravity; the bounce structure replaces them by contractions
between distinct time-\(n\) spaces.

\subsection{Executable certificate}

The verifier

\[
 \texttt{scripts/verify\_sm\_v19\_v38\_bounce\_chain.py}           \tag{41.7}
\]

builds the Gaussian bounce chain with \(r=\tfrac12\), \(v_0=1\),
noise variance \(\tfrac14\), and drifts \(b_n=n+1\).  It computes the
conditional offsets \(c_n\), the future link means and increments, the
mirrored link means, and evaluates the reflected link Gram both from the
explicit cross covariances \(E[q_{-a}q_b]=r^{a+b}v_0+c_ac_b\) and from
the closed form (41.4), checking equality and all principal minors for
\(n,m\leq2\); it checks the site Gram, the oddness of the mirrored means,
and the evenness (41.5) of the increments.  It emits

\[
 \texttt{artifacts/sm\_v19\_v38\_bounce\_chain.json}.              \tag{41.8}
\]

The hashes are

\[
\begin{array}{c|l}
\hbox{object}&\hbox{SHA--256}\\ \hline
\hbox{verifier}&
\texttt{305B0DC3A4C2EE53714FD5EE07EAE6D105DD04B14655D03A0C4A82F5C0277C3E}\\
\hbox{canonical payload}&
\texttt{AB4B185A7C582717125E6E2ADA0F9A795556AB978F8E352882FD107008E09C38}.
\end{array}                                                     \tag{41.9}
\]

\begin{firewall}[A bounce chain is a kinematic frame, not a solution]
\label{fire:v19v38}
Theorem \ref{thm:v19v38-chain} and Theorem
\ref{thm:v19v38-reconstruction} are general facts about
reflection-symmetric Markov laws.  No chain has been shown to reproduce
the Einstein rates, no Euclidean gravitational action selects the
kernels \(P_n\), no lattice-spacing or infinite-volume limit is taken,
and no Standard Model sector is coupled.  The bounce example is a
Gaussian illustration.
\end{firewall}

\begin{decision}[V39 acquisition]
\label{dec:v19v38-next}
V39 will ask what the Einstein identities require of the kernels
\(P_n\) at the linearized level about the plane.  It will take the
reflection-symmetric TT chain with a time-dependent conditional mean
and covariance, impose that the expected mean-curvature increments equal
the constant-lapse rate (39.4) computed from the V26 quadratic charge
at second order, and determine whether that fixes the homogeneous
component of the chain uniquely.  This is the first place where the
Einstein dynamics, rather than the constraints alone, enters the seed
law.
\end{decision}

\begin{assessment}[V38 acquisition]
\label{ass:v19v38}
Stationarity is removed from the reflection-positivity theorem, the odd
mean curvature of a reflection-symmetric chain is a bounce with the
Einstein sign pattern, both V36 obstructions are shown to be artefacts
of stationarity and sign symmetry, and the reconstructed state space at
the time-symmetric slice is identified exactly with its non-autonomous
contractions.  The compact programme now has a consistent
reflection-positive frame in which the Einstein rates can be imposed.
\end{assessment}

\begin{record}[V38 append record]
\label{rec:v19v38}
V38 is appended to the validated V37 whose installed SHA--256 was
\texttt{47429A3DD9BCC74D237E0A49AE0764EA5115A64FDB5EFF225DADD0BA671A71B5}.
After installation only V38 is the active numeric SM note; frozen
archives and unrelated notes are unchanged.  The next compulsory
companion audit is V40.
\end{record}


\section{V39: toroidal rigidity at the reflection plane and the negative cosmological constant}
\label{sec:v19v39}

V38 removed stationarity and found that an odd mean curvature is a
sign-changing profile about the reflection plane.  Before imposing
Einstein rates on such a chain, this note asks what the reflection plane
is classically.  A reflection-symmetric development has \(K=0\) on the
plane, hence a maximal slice; but on \(\mathbb T^3\) the classical
rigidity theorem for nonnegative scalar curvature makes every
time-symmetric vacuum or dominant-energy slice flat.  The vacuum
reflection-symmetric programme on the torus is therefore classically
empty, and its sign pattern, where nonempty, is a volume maximum rather
than a bounce.  A negative cosmological constant changes both
conclusions: the plane becomes a strict volume minimum, and
time-symmetric slices carrying gravitational waves exist provided
\(\Lambda\) balances the wave energy at second order.  The exact V28
branch is identified with Kasner data as a consistency check.

\subsection{Maximal slices and the volume}

For a CMC development with lapse \(N>0\) and shift \(X\), the volume
evolves by \(\partial_t\,dV_g=(-N\tau+\operatorname{div}X)\,dV_g\), so

\[
 \frac{d}{dt}\operatorname{Vol}_g=-\tau\int N\,dV_g.                 \tag{42.1}
\]

\begin{proposition}[Volume at a maximal slice]
\label{prop:v19v39-maximal}
Let a CMC vacuum development pass through a slice with \(\tau=0\) at
\(t=0\).  Then \(\dot{\operatorname{Vol}}(0)=0\) and

\[
 \ddot{\operatorname{Vol}}(0)
 =-\frac{\int N|K|^2dV}{\operatorname{Vol}}\int N\,dV\leq0,           \tag{42.2}
\]

with strict inequality if \(K\not\equiv0\) on the slice.  If moreover
the development is reflection symmetric, \(g(-t)=g(t)\),
\(K(-t)=-K(t)\), then \(K(0)=0\) and the volume is stationary to third
order at the plane.
\end{proposition}

\begin{proof}
Differentiate (42.1) at \(\tau=0\): only \(-\dot\tau\int N\,dV\)
survives, and (39.4) gives \(\dot\tau\).  Reflection symmetry makes
\(K(0)=-K(0)\), so \(|K|^2=O(t^2)\), \(\dot\tau=O(t^2)\),
\(\tau=O(t^3)\), and \(\dot{\operatorname{Vol}}=O(t^3)\).
\end{proof}

Thus in vacuum the Einstein sign pattern of V38, an odd increasing
mean curvature, is a maximal-volume turnaround, expanding before the
plane and contracting after it; the word bounce in V38 should be read as
turnaround in the vacuum case.

\subsection{Toroidal rigidity}

\begin{theorem}[Time-symmetric slices on the torus are flat]
\label{thm:v19v39-rigidity}
Let \((g,K,\rho)\) be initial data on \(\mathbb T^3\) for Einstein--matter
with the Hamiltonian constraint (29.18) and energy density
\(\rho\geq0\).  If \(K\equiv0\), then \(g\) is flat and \(\rho\equiv0\).
The same holds under the weaker hypothesis \(\tau\equiv0\).
\end{theorem}

\begin{proof}
With \(K=0\), or with \(\tau=0\), (29.18) gives
\(R(g)=|K|_g^2+16\pi G\rho\geq0\).  By the theorem of Schoen--Yau and
Gromov--Lawson, a metric on \(\mathbb T^3\) with nonnegative scalar
curvature is flat, and then \(R=0\) forces \(|K|^2+16\pi G\rho=0\).
\end{proof}

\begin{corollary}[The vacuum reflection-symmetric programme on the torus is empty]
\label{cor:v19v39-empty}
Every reflection-symmetric Einstein--matter development on
\(\mathbb T^3\) with \(\rho\geq0\) is the static flat torus with
\(\rho\equiv0\).  Every non-flat vacuum CMC development on
\(\mathbb T^3\) has mean curvature of one strict sign for all time, so
it never passes through a maximal slice.
\end{corollary}

\begin{proof}
Reflection symmetry gives \(K(0)=0\), and Theorem
\ref{thm:v19v39-rigidity} gives flat time-symmetric data, whose vacuum
development is static flat.  For the second statement, a maximal slice
would be time symmetric or have \(\tau=0\) with \(K\neq0\); the theorem
excludes both unless flat, and (39.4) makes \(\tau\) monotone, so it
cannot cross zero.
\end{proof}

This is the sharpest form of the V26 obstruction and closes the
vacuum torus route for the reflection-symmetric chain of V38: the
classical limit at the plane is flat.  Combined with Theorem
\ref{thm:v19v32-sharp}, an \(L^2\)-continuous odd extrinsic curvature
vanishes on the plane in the quantum process as well, so the quantum
process cannot evade the rigidity through a nonzero classical \(K\)
there.

\subsection{The exact V28 branch is Kasner data}

\begin{proposition}[Kasner identification]
\label{prop:v19v39-kasner}
The V28 data \(g=\delta\), \(K=\lambda\operatorname{diag}(4,4,-2)\),
\(\tau=6\lambda\), are the \(t=-1/(6\lambda)\) slice of the Kasner
vacuum metric

\[
 -dt^2+t^{4/3}dx^2+t^{4/3}dy^2+t^{-2/3}dz^2,
 \qquad (p_1,p_2,p_3)=\bigl(\tfrac23,\tfrac23,-\tfrac13\bigr),         \tag{42.3}
\]

up to constant rescalings of the coordinates.  Along it,
\(\tau=-1/t\), \(|K|^2=1/t^2\), \(\operatorname{Vol}\propto t\), and
(39.4) and (42.1) hold exactly with \(N=1\).
\end{proposition}

\begin{proof}
For \(g=\operatorname{diag}(t^{2p_i})\) and \(N=1\), (39.1) gives
\(K^i{}_j=-p_i\delta^i_j/t\), \(\tau=-\sum p_i/t=-1/t\), and
\(|K|^2=\sum p_i^2/t^2=1/t^2\); the Kasner relations
\(\sum p_i=\sum p_i^2=1\) are the vacuum equations.  At the V28 slice,
\(p_i=-K^i{}_it=K^i{}_i/(6\lambda)=(4,4,-2)/6\), which satisfies both
relations.  Then \(\dot\tau=1/t^2=|K|^2\) and
\(\dot{\operatorname{Vol}}=\operatorname{Vol}/t=-\tau\operatorname{Vol}\).
\end{proof}

The homogeneous compensator of V26 and V28 is therefore an anisotropic
Kasner expansion, monotone in \(\tau\), never time symmetric: exactly
what Corollary \ref{cor:v19v39-empty} predicts.

\subsection{A negative cosmological constant}

With a cosmological constant the Hamiltonian constraint and the
mean-curvature evolution read

\[
 R+\tau^2-|K|^2=2\Lambda,\qquad
 \partial_t\tau=N\bigl(|K|^2+2\Lambda\bigr)-\Delta_gN+X(\tau),         \tag{42.4}
\]

the second by the proof of Proposition \ref{prop:v19v36-monotone}
with the modified constraint.  The CMC conformal equation (31.3)
becomes

\[
 -8\Delta\psi-|\sigma|^2\psi^{-7}+\Bigl(\tfrac23\tau^2-2\Lambda\Bigr)\psi^5=0.
                                                                    \tag{42.5}
\]

\begin{proposition}[Bounce for \(\Lambda<0\)]
\label{prop:v19v39-bounce}
Let \(\Lambda<0\) and let a CMC development be reflection symmetric
about a slice at \(t=0\).  Then, with the CMC lapse,

\[
 \dot\tau(0)=2\Lambda\langle N\rangle<0,\qquad
 \ddot{\operatorname{Vol}}(0)=-2\Lambda\langle N\rangle\int N\,dV>0,   \tag{42.6}
\]

so the plane is a strict local minimum of the volume: a bounce.  The
mean curvature is positive before the plane and negative after it.
\end{proposition}

\begin{proof}
At the plane \(K=0\), and the CMC lapse condition, that
\(\partial_t\tau\) be spatially constant, integrates (42.4) to
\(\dot\tau(0)\operatorname{Vol}=2\Lambda\int N\,dV\).  Then (42.1) gives
(42.6), and the sign of \(\tau\) near \(t=0\) follows from
\(\tau(0)=0\), \(\dot\tau(0)<0\).
\end{proof}

\begin{theorem}[Second-order balance for time-symmetric waves]
\label{thm:v19v39-balance}
Let \(g(\lambda)\) be a \(C^2\) family of metrics on \(\mathbb T^3\)
with \(g(0)=\delta\), \(g'(0)=h\) transverse traceless, and let
\(\Lambda(\lambda)=\lambda^2\Lambda_2+o(\lambda^2)\).  If each
\(g(\lambda)\) is time-symmetric vacuum data with cosmological constant
\(\Lambda(\lambda)\), that is \(R(g(\lambda))=2\Lambda(\lambda)\), then

\[
 \Lambda_2=-\frac18\bigl\langle|\nabla h|^2\bigr\rangle\leq0,          \tag{42.7}
\]

with equality only for \(h=0\) when \(h\) has zero mean.  For
\(h=\cos z\,e^+\), \(\Lambda_2=-1/8\).
\end{theorem}

\begin{proof}
Integrate \(R(g(\lambda))=2\Lambda(\lambda)\) against \(dV_{g(\lambda)}\):
\({\cal S}(g(\lambda))=2\Lambda(\lambda)\operatorname{Vol}(g(\lambda))\).
By Lemma \ref{lem:v19v26-hessian}, the left side is
\(-\tfrac{\lambda^2}4\int|\nabla h|^2+o(\lambda^2)\), independently of
\(g''(0)\); the right side is
\(2\lambda^2\Lambda_2\operatorname{Vol}(\delta)+o(\lambda^2)\).  Compare
coefficients.  The Fourier value uses \(\langle|\nabla h|^2\rangle=1\)
from (29.22).
\end{proof}

Thus on the torus a time-symmetric slice can carry gravitational waves
only if a negative cosmological constant of second order in the wave
amplitude balances their Dirichlet energy, in the same way that in V26
a homogeneous mean curvature balanced it and in V28 that balance was
Kasner expansion.  With \(\Lambda<0\) the coefficient
\(\tfrac23\tau^2-2\Lambda\) in (42.5) is positive even at \(\tau=0\), so
the V29 existence and uniqueness proof applies verbatim on every slice,
including the plane, once the background conformal class is allowed to
carry the even wave \(h\); the conformally flat class alone cannot,
because integrating (42.5) with \(|\sigma|=0\) gives a contradiction.

\subsection{Executable certificate}

The verifier

\[
 \texttt{scripts/verify\_sm\_v19\_v39\_toroidal\_rigidity.py}        \tag{42.8}
\]

checks over the rationals that the V28 data at \(\lambda=1/3\) satisfy
the Hamiltonian constraint and correspond to Kasner exponents
\((2/3,2/3,-1/3)\) at \(t=-1/2\), with \(\sum p_i=\sum p_i^2=1\); that
along Kasner \(\dot\tau=|K|^2\) and
\(\dot{\operatorname{Vol}}=-\tau\operatorname{Vol}\) at three rational
times; that (42.7) gives \(-1/8\) for the Fourier wave, consistent with
the V26 Hessian; and that (42.6) is positive for \(\Lambda=-1/8\).  It
emits

\[
 \texttt{artifacts/sm\_v19\_v39\_toroidal\_rigidity.json}.           \tag{42.9}
\]

The hashes are

\[
\begin{array}{c|l}
\hbox{object}&\hbox{SHA--256}\\ \hline
\hbox{verifier}&
\texttt{5ED269D32BAA11BE61A262788996822EB76ADFCD8671E7106D3B1325CECF779D}\\
\hbox{canonical payload}&
\texttt{95850A96DF987677DC5ABE01B7F8E78A5A4C0465D58161E0795064850C13692F}.
\end{array}                                                     \tag{42.10}
\]

\begin{firewall}[External input and scope]
\label{fire:v19v39}
Theorem \ref{thm:v19v39-rigidity} imports the Schoen--Yau and
Gromov--Lawson rigidity theorem for nonnegative scalar curvature on the
three-torus; it is the only external theorem used since the restart
beyond standard elliptic theory and bifurcation theory.  The
\(\Lambda<0\) statements are classical: no measure, chain, or
reconstruction with a cosmological constant has been built, and Theorem
\ref{thm:v19v39-balance} is a necessary condition whose sufficiency is
deferred.
\end{firewall}

\begin{decision}[V40 audit and V41 acquisition]
\label{dec:v19v39-next}
V40 is the compulsory companion audit.  V41 will prove sufficiency of
(42.7): an analytic branch of time-symmetric vacuum data
\(g_\lambda=\psi_\lambda^4(\delta+\lambda h)\) with
\(R(g_\lambda)=2\Lambda(\lambda)\), \(\Lambda(\lambda)\) even and
analytic with leading coefficient (42.7), by the mean/mean-zero
Lyapunov--Schmidt split of V31 applied to the conformal Laplacian of
\(\delta+\lambda h\).  This supplies the classical content of the
reflection plane for the \(\Lambda<0\) torus programme, after which the
reflection-symmetric chain of V38 can be seeded at a bounce.
\end{decision}

\begin{assessment}[V39 acquisition]
\label{ass:v19v39}
The reflection plane of a reflection-symmetric development is a
time-symmetric slice, and on the torus with nonnegative energy density
such a slice is flat by toroidal rigidity.  The vacuum
reflection-symmetric compact programme is therefore classically empty,
which explains the monotone mean curvature of the exact Kasner branch
behind V28.  A negative cosmological constant of second order in the
wave amplitude restores time-symmetric slices with gravitational waves
and turns the plane into a strict volume minimum, a genuine bounce.
\end{assessment}

\begin{record}[V39 append record]
\label{rec:v19v39}
V39 is appended to the validated V38 whose installed SHA--256 was
\texttt{6D7A997AA887027695008436455EDC66BFA2361F4D66B6E956241B6DE96C931A}.
After installation only V39 is the active numeric SM note; frozen
archives and unrelated notes are unchanged.
\end{record}


\section{V40: compulsory companion audit and ten-version sweep after toroidal rigidity}
\label{sec:v19v40}

This is the required audit following V35 and the ten-version sweep of
notes generated by the parallel programmes.  Every active companion
file in the shared folder was inspected read-only before the
\(\Lambda<0\) time-symmetric branch announced in V39 is constructed.

\begin{record}[V40 companion snapshot]
\label{rec:v19v40-snapshot}
The inspected files were

\[
\begin{array}{c|r|r}
\hbox{file and SHA--256}&\hbox{bytes}&\hbox{lines}\\ \hline
\texttt{Renewal\_Yang\_Mills\_Mass\_Gap\_Research\_Notes\_V31.tex}
&432885&11118\\
\multicolumn{3}{l}{
\texttt{73712A413552FD14065BA67CA448B57F3742DED3CC0EFF1A74512BBCF2B83DDB}}
\\
\texttt{Renewal\_NS\_Stability\_Research\_Notes\_V26.tex}
&278283&5319\\
\multicolumn{3}{l}{
\texttt{82C887F8C2C69E4F28FAD7C3DC8DE5B9099CB5A4BF431EBA1FFF3E91935978AD}}
\\
\texttt{Renewal\_GRH\_Cofinal\_Visibility\_Attack\_V51.tex}
&489091&15170\\
\multicolumn{3}{l}{
\texttt{61927F9477BD72FAD280D5EF4042919932C7687ED895DE4571A71F2D81450134}}
\\
\texttt{Renewal\_YM\_Parallel\_Research\_Notes\_V5.tex}
&54174&1351\\
\multicolumn{3}{l}{
\texttt{306F1BB84CFA8A8D47EA569EC17E9545F9867C8D32B548EC9D9C444B4F3C0512}}.
\end{array}                                                     \tag{43.1}
\]

Since V35, Yang--Mills has advanced from V30 to V31, GRH from V50 to
V51, and the Navier--Stokes and parallel Yang--Mills files are
unchanged.
\end{record}

\subsection{Companion findings}

Yang--Mills V31 withdraws its prime-production route: it proves that
the completed coefficient artifact certifies a single anchor box, that a
box-uniform route would need more than \(6.69\times10^{12}\) boxes, and
that a certified nonzero constant class would refute rather than close
its self-similar comparison.  It redirects to a nested Gaussian tower
and records that its own ten-iteration sweep is due at its V40.  Its
firewall leaves the Gaussian anomaly, interacting measure, regulator
removal, reconstruction, clustering, and mass gap open.

GRH V51 proves a sharp total-height budget \(t_1+t_2+t_3\leq15\) for
its three-pair heat sextic and a complete classification above
threshold, closing the finite model that occupied its V10--V50.  It is a
theorem about one sextic and claims nothing about GRH.

Navier--Stokes V26 and Yang--Mills Parallel V5 are unchanged since V35.

\begin{decision}[V40 audit decision]
\label{dec:v19v40-audit}
No companion theorem, numerical constant, or executable artifact is
imported.  The V39 direction stands: V41 constructs the analytic
\(\Lambda<0\) time-symmetric branch on the torus.
\end{decision}

\begin{proof}
The Yang--Mills withdrawal concerns a finite-field sign instrument
and a self-similar Wilson comparison; the GRH result concerns
real-rootedness of a heat-flowed sextic.  Neither contains a scalar
curvature rigidity statement, a conformal constraint equation, a
reflection-symmetric Markov law, or a cosmological constant.  The one
methodological parallel, that both companions moved upstream after
proving their finite instruments could not decide the continuum
question, is the same move made here in V36 and V39; it is a shared
discipline, not an imported result.
\end{proof}

\begin{assessment}[V40 acquisition]
\label{ass:v19v40}
The audit is current.  Since V35 the programme has proved that
stationary constrained chains fail the Einstein step and monotonicity
identities, separated the fixed-energy and fixed-density regimes,
removed stationarity from the reflection-positivity theorem, and shown
by toroidal rigidity that the vacuum reflection plane on
\(\mathbb T^3\) is flat, so that a negative cosmological constant of
second order in the wave amplitude is required for a nontrivial
compact bounce.
\end{assessment}

\begin{record}[V40 append record]
\label{rec:v19v40}
V40 is appended to the validated V39 whose installed SHA--256 was
\texttt{7A1DC7903DB2B572FBCD889153F7B402EF338120AD68D2EE83397D6E1D3659A8}.
After installation only V40 is the active numeric SM note; the
companions, frozen archives, and all unrelated notes are unchanged.  The
next compulsory companion audit is V45.
\end{record}


\section{V41: the analytic time-symmetric branch with a negative cosmological constant}
\label{sec:v19v41}

V39 showed that a time-symmetric slice of the torus can carry a
gravitational wave only if a cosmological constant of second order in
the wave amplitude balances its Dirichlet energy.  This note proves
sufficiency.  The cosmological constant is itself the compatibility
variable that removes the constant cokernel of V26, so a single analytic
implicit-function argument produces the branch, with no Lyapunov--Schmidt
reduction.  The second-order conformal response is computed in closed
form for the Fourier wave, and the resulting data are the classical
bounce slice of the \(\Lambda<0\) torus programme.

\subsection{The conformal equation with a wave background}

Fix a smooth transverse-traceless \(h\) on the flat torus, and put
\(\gamma_\lambda:=\delta+\lambda h\), a Riemannian metric for
\(|\lambda|<\lambda_h:=1/\max\|h\|_{\rm op}\).  For \(\psi>0\) the
three-dimensional conformal identity gives

\[
 R(\psi^4\gamma)=\psi^{-5}\bigl(-8\Delta_\gamma\psi+R(\gamma)\psi\bigr),
                                                                    \tag{44.1}
\]

so time-symmetric vacuum data \((\psi^4\gamma_\lambda,0)\) with
cosmological constant \(\Lambda\) are exactly the positive solutions of

\[
 {\cal E}(\lambda,\psi,\Lambda):=
 -8\Delta_{\gamma_\lambda}\psi+R(\gamma_\lambda)\psi-2\Lambda\psi^5=0.
                                                                    \tag{44.2}
\]

If \((\psi,\Lambda)\) solves (44.2) then so does
\((k\psi,k^{-4}\Lambda)\) for every \(k>0\); this scaling freedom is
fixed by \(\langle\psi\rangle=1\), the flat average.

\begin{theorem}[Analytic time-symmetric branch]
\label{thm:v19v41-branch}
Fix \(0<\alpha<1\).  There are \(\lambda_0>0\) and unique real-analytic
maps

\[
 u:(-\lambda_0,\lambda_0)\to C^{2,\alpha}_0,\qquad
 \Lambda:(-\lambda_0,\lambda_0)\to\mathbb R,                          \tag{44.3}
\]

with \(u(0)=0\), \(\Lambda(0)=0\), such that
\(\psi_\lambda:=1+u(\lambda)>0\) and
\({\cal E}(\lambda,\psi_\lambda,\Lambda(\lambda))=0\).  Moreover
\(u'(0)=0\), \(\Lambda'(0)=0\), and

\[
 \Lambda(\lambda)=-\frac{\lambda^2}8\bigl\langle|\nabla h|^2\bigr\rangle
 +O(\lambda^3),
 \qquad
 -8\Delta u_2=\langle R_2\rangle-R_2,\quad u=\lambda^2u_2+O(\lambda^3),
                                                                    \tag{44.4}
\]

where \(R(\gamma_\lambda)=\lambda^2R_2+O(\lambda^3)\) and
\(\langle R_2\rangle=-\tfrac14\langle|\nabla h|^2\rangle\).
Consequently \(g_\lambda:=\psi_\lambda^4\gamma_\lambda\) with
\(K_\lambda=0\) is an exact solution of the constraints (42.4) with
\(\Lambda=\Lambda(\lambda)<0\) for \(0<|\lambda|<\lambda_0\) whenever
\(h\) is not constant.
\end{theorem}

\begin{proof}
Consider
\({\cal E}:(-\lambda_h,\lambda_h)\times\{u\in C^{2,\alpha}_0:1+u>0\}\times\mathbb R
\to C^{0,\alpha}\), \((\lambda,u,\Lambda)\mapsto{\cal E}(\lambda,1+u,\Lambda)\).
The coefficients of \(\Delta_{\gamma_\lambda}\) and the function
\(R(\gamma_\lambda)\) are rational in the entries of \(\gamma_\lambda\)
and its first two derivatives, hence real-analytic in \(\lambda\) with
values in \(C^{k,\alpha}\) for every \(k\); the map is polynomial in
\(u\) and \(\Lambda\).  At \((0,0,0)\), \(\gamma_0=\delta\), \(R=0\), and
the derivative with respect to \((u,\Lambda)\) is

\[
 (v,M)\longmapsto-8\Delta v-2M:\quad
 C^{2,\alpha}_0\times\mathbb R\to C^{0,\alpha}=C^{0,\alpha}_0\oplus\mathbb R.
                                                                    \tag{44.5}
\]

The first summand is the isomorphism (34.11) onto \(C^{0,\alpha}_0\),
and \(M\mapsto-2M\) is onto the constants, so (44.5) is an
isomorphism.  The analytic implicit-function theorem gives (44.3).
Positivity holds for small \(\lambda_0\).

Differentiate \({\cal E}=0\) at \(\lambda=0\):
\(-8\Delta u'(0)+DR_\delta(h)-2\Lambda'(0)=0\).  For transverse
traceless \(h\), (29.7) gives \(DR_\delta(h)=\partial_i\partial_jh_{ij}-\Delta\operatorname{tr}h=0\);
averaging gives \(\Lambda'(0)=0\), and then \(\Delta u'(0)=0\) with zero
mean gives \(u'(0)=0\).  Hence \(g_\lambda=\delta+\lambda h+O(\lambda^2)\)
has \(g'(0)=h\), and Theorem \ref{thm:v19v39-balance} applied to the
family \(g_\lambda\) with \(R(g_\lambda)=2\Lambda(\lambda)\) yields the
first relation in (44.4).  At order \(\lambda^2\) the equation reads
\(-8\Delta u_2+R_2-2\Lambda_2=0\), because \(\Delta_{\gamma_\lambda}-\Delta=O(\lambda)\)
acts on \(u=O(\lambda^2)\); averaging gives \(\Lambda_2=\langle R_2\rangle/2\)
and the mean-zero part gives the second relation.  The identity
\(\langle R_2\rangle=-\tfrac14\langle|\nabla h|^2\rangle\) is Lemma
\ref{lem:v19v26-hessian}, since
\({\cal S}(\gamma_\lambda)=\int R(\gamma_\lambda)dV_{\gamma_\lambda}
=\lambda^2\int R_2\,dx+O(\lambda^3)\) for traceless \(h\).  Finally,
\(\Lambda(\lambda)<0\) for small \(\lambda\neq0\) because
\(\langle|\nabla h|^2\rangle>0\) unless \(h\) is constant.
\end{proof}

The cosmological constant plays, for time-symmetric data, exactly the
role that the homogeneous mean curvature played in V26--V31: it is the
one-dimensional complement of the range of the linearized constraint.

\subsection{Closed form for the Fourier wave}

\begin{proposition}[Explicit second-order branch]
\label{prop:v19v41-fourier}
For \(h=\cos z\,e^+\), the metric
\(\gamma_\lambda=dz^2+(1+\lambda\cos z)dx^2+(1-\lambda\cos z)dy^2\) has

\[
 R(\gamma_\lambda)=-\lambda^2\Bigl(\frac14+\frac74\cos2z\Bigr)+O(\lambda^3),
                                                                    \tag{44.6}
\]

and therefore

\[
 \Lambda(\lambda)=-\frac{\lambda^2}8+O(\lambda^3),\qquad
 \psi_\lambda=1+\frac{7\lambda^2}{128}\cos2z+O(\lambda^3).            \tag{44.7}
\]
\end{proposition}

\begin{proof}
For \(dz^2+a(z)^2dx^2+b(z)^2dy^2\), the scalar curvature is
\(R=-2(a''/a+b''/b+a'b'/(ab))\).  With
\(a=(1+\lambda\cos z)^{1/2}\), \(b=(1-\lambda\cos z)^{1/2}\), expansion
to second order gives
\(a''/a+b''/b=\lambda^2(\tfrac14+\tfrac34\cos2z)+O(\lambda^3)\) and
\(a'b'/(ab)=-\tfrac{\lambda^2}4\sin^2z+O(\lambda^3)
=\lambda^2(-\tfrac18+\tfrac18\cos2z)+O(\lambda^3)\); the sum is
\(\lambda^2(\tfrac18+\tfrac78\cos2z)\), which gives (44.6).  Then
\(\langle R_2\rangle=-\tfrac14\), in agreement with
\(\langle|\nabla h|^2\rangle=1\) from (29.22), so \(\Lambda_2=-1/8\);
and \(-8\Delta u_2=\tfrac74\cos2z\) gives \(u_2=\tfrac7{128}\cos2z\).
\end{proof}

The bounce slice of the \(\Lambda<0\) torus programme for this wave is
therefore, to second order,

\[
 g_\lambda=\Bigl(1+\frac{7\lambda^2}{128}\cos2z\Bigr)^4
 \bigl(dz^2+(1+\lambda\cos z)dx^2+(1-\lambda\cos z)dy^2\bigr),
 \qquad K_\lambda=0,\qquad \Lambda=-\frac{\lambda^2}8.                 \tag{44.8}
\]

By Proposition \ref{prop:v19v39-bounce}, its development has a strict
volume minimum at the plane with \(\dot\tau(0)=2\Lambda\langle N\rangle<0\).

\subsection{Executable certificate}

The verifier

\[
 \texttt{scripts/verify\_sm\_v19\_v41\_lambda\_time\_symmetric\_branch.py}
                                                                    \tag{44.9}
\]

implements exact series in \(\lambda\) with cosine-polynomial
coefficients, including square roots, inverses, and derivatives, and
evaluates \(R(\gamma_\lambda)\) through order \(\lambda^2\).  It checks
that the orders \(\lambda^0\) and \(\lambda^1\) vanish, that
\(R_2=-\tfrac14-\tfrac74\cos2z\), that \(\Lambda_2=-1/8\), that
\(u_2=\tfrac7{128}\cos2z\) with zero order-two residual, and the V26
consistency \(\langle R_2\rangle=-\tfrac14\langle|\nabla h|^2\rangle\).
It emits

\[
 \texttt{artifacts/sm\_v19\_v41\_lambda\_time\_symmetric\_branch.json}.
                                                                    \tag{44.10}
\]

The hashes are

\[
\begin{array}{c|l}
\hbox{object}&\hbox{SHA--256}\\ \hline
\hbox{verifier}&
\texttt{B52EE06EAD3935EC7118CB121653813E63AAAB500EA1C33A3F1B62952F05449C}\\
\hbox{canonical payload}&
\texttt{0C1C395A501F1B33440803817FDDECB76F9FF130D4D5D6966F8D84BB3C25FDE2}.
\end{array}                                                     \tag{44.11}
\]

\begin{firewall}[A classical slice with a tuned constant]
\label{fire:v19v41}
Theorem \ref{thm:v19v41-branch} is a classical existence theorem for
time-symmetric data in which \(\Lambda\) is determined by the wave.  A
physical theory has one fixed \(\Lambda\); here each amplitude selects
its own, and the inverse question, which waves are admissible at a
given \(\Lambda<0\), is not addressed.  No Lorentzian development beyond
the initial rates, no reflection-positive chain with \(\Lambda\), no
Standard Model source, and no measure are constructed.
\end{firewall}

\begin{decision}[V42 acquisition]
\label{dec:v19v41-next}
V42 will invert the roles: fix \(\Lambda<0\) and describe the set of
time-symmetric torus data near the branch, proving that for each fixed
\(\Lambda<0\) small the admissible wave amplitudes form a sphere
\(\langle|\nabla h|^2\rangle=-8\Lambda+O(\Lambda^{3/2})\) in the TT
mode space, so that the classical bounce slice at fixed \(\Lambda\) is
a compact family.  It will then define the \(\Lambda<0\)
reflection-symmetric chain of V38 seeded on that family and verify the
pushforward hypotheses.
\end{decision}

\begin{assessment}[V41 acquisition]
\label{ass:v19v41}
The necessary condition of V39 is sufficient: a negative cosmological
constant of second order in the amplitude is exactly what makes a
time-symmetric torus slice with a gravitational wave solvable, and the
cosmological constant is the compatibility variable that V26 found
missing.  The Fourier wave gives an explicit second-order bounce slice.
The compact \(\Lambda<0\) torus programme now has its classical plane.
\end{assessment}

\begin{record}[V41 append record]
\label{rec:v19v41}
V41 is appended to the validated V40 whose installed SHA--256 was
\texttt{FEC9891645B58B207F3C8F864DD5401FE6B1D19A98D23097E00DFA5064E92773}.
After installation only V41 is the active numeric SM note; frozen
archives and unrelated notes are unchanged.
\end{record}


\section{V42: the admissible wave sphere at fixed cosmological constant and the seeded chain}
\label{sec:v19v42}

V41 let the wave choose its cosmological constant.  A physical theory
has one \(\Lambda\), so this note inverts the roles.  For a fixed small
\(\Lambda<0\), the time-symmetric torus slices near flat form, in each
finite TT mode space, an analytic sphere whose Dirichlet radius is
\(-8\Lambda\) to leading order.  Near any point of that sphere the
\(\Lambda\)-Lichnerowicz equation has a unique analytic solution for
nearby backgrounds, TT momenta, and mean curvatures, because the
linearization at the plane has lowest eigenvalue \(-8\Lambda+O(\Lambda^2)\).
This gives a single-valued solver on a neighbourhood of the plane
datum, and the reflection-symmetric chain of V38 can be seeded there
with bounded noise and pushed forward with Theorem
\ref{thm:v19v32-pushforward}.

\subsection{Analytic dependence on the wave}

Let \(V_N\) be a finite-dimensional space of smooth mean-zero
transverse-traceless tensors spanned by finitely many real Fourier
modes, with the Dirichlet form \(Q(h):=\langle|\nabla h|^2\rangle\),
positive definite on \(V_N\).

\begin{proposition}[Analytic branch over \(V_N\)]
\label{prop:v19v42-analytic}
There are a neighbourhood \(U\) of \(0\) in \(V_N\) and unique
real-analytic maps \(u:U\to C^{2,\alpha}_0\), \(\Lambda:U\to\mathbb R\)
with \(u(0)=0\), \(\Lambda(0)=0\), such that
\({\cal E}(h,1+u(h),\Lambda(h))=0\) in the notation of (44.2) with
\(\gamma=\delta+h\).  Moreover

\[
 \Lambda(h)=-\frac18Q(h)+O(\|h\|^3),\qquad u(h)=O(\|h\|^2).          \tag{45.1}
\]
\end{proposition}

\begin{proof}
Replace the scalar parameter \(\lambda\) in the proof of Theorem
\ref{thm:v19v41-branch} by \(h\in V_N\); the map is analytic in \(h\)
because \(R(\delta+h)\) and \(\Delta_{\delta+h}\) are rational in
\(h\) and its derivatives, and the derivative (44.5) is unchanged.
Restricting to the line \(\lambda h\) gives (44.4), hence (45.1), with
the cubic remainder uniform on \(U\) by analyticity.
\end{proof}

\subsection{The admissible sphere at fixed \(\Lambda\)}

\begin{theorem}[Admissible waves at fixed \(\Lambda<0\)]
\label{thm:v19v42-sphere}
Let \(S:=\{\omega\in V_N:Q(\omega)=1\}\).  There is \(\Lambda_0>0\) such
that for every \(\Lambda_*\in(-\Lambda_0,0)\) the set

\[
 {\cal A}(\Lambda_*):=\{h\in U:\Lambda(h)=\Lambda_*\}                  \tag{45.2}
\]

is a compact real-analytic hypersurface of \(V_N\), the graph
\(\{r(\omega)\omega:\omega\in S\}\) of an analytic function
\(r:S\to(0,\infty)\) with

\[
 r(\omega)^2=-8\Lambda_*\bigl(1+O(|\Lambda_*|^{1/2})\bigr),
 \qquad\hbox{that is}\qquad
 Q(h)=-8\Lambda_*+O(|\Lambda_*|^{3/2})\ \hbox{on }{\cal A}(\Lambda_*).  \tag{45.3}
\]

Each \(h\in{\cal A}(\Lambda_*)\) gives exact time-symmetric data
\((\psi^4(\delta+h),0)\) with \(\psi=1+u(h)\) for the constraints
(42.4) at the fixed constant \(\Lambda_*\).
\end{theorem}

\begin{proof}
Write \(h=r\omega\) with \(\omega\in S\), \(r\geq0\).  By (45.1),
\(\Lambda(r\omega)=-r^2/8+r^3\rho(r,\omega)\) with \(\rho\) analytic
and bounded.  Put \(s:=\sqrt{-8\Lambda_*}\) and seek \(r=s(1+\eta)\).
The equation \(\Lambda(r\omega)=\Lambda_*\) becomes
\((1+\eta)^2-1=8s(1+\eta)^3\rho\), that is
\(\eta(2+\eta)=8s(1+\eta)^3\rho(s(1+\eta),\omega)\).  At \(s=0\) the
solution is \(\eta=0\) and the derivative in \(\eta\) is \(2\neq0\); the
analytic implicit-function theorem gives a unique small analytic
\(\eta(s,\omega)=O(s)\), uniformly on the compact \(S\).  Hence
\(r(\omega)=s(1+O(s))\), which is (45.3).  Compactness and analyticity of
the graph follow from those of \(S\) and \(r\).  The data statement is
Proposition \ref{prop:v19v42-analytic}.
\end{proof}

For the two-mode space spanned by \(\cos z\,e^+\) and
\(\cos z\,e^\times\), \(Q(x,y)=x^2+y^2\) and \({\cal A}(-1/8)\) is, to
leading order, the unit circle.

\subsection{Local solvability of the \(\Lambda\)-Lichnerowicz equation}

On a background \(\gamma\) with \(\gamma\)-transverse-traceless
\(\sigma\) and constant \(\tau\), the constraints (42.4) for
\(g=\psi^4\gamma\), \(K=\psi^{-2}\sigma+\tfrac\tau3\psi^4\gamma\) reduce,
as in Lemma \ref{lem:v19v28-reduction}, to

\[
 {\cal L}_\Lambda(\gamma,\sigma,\tau;\psi):=
 -8\Delta_\gamma\psi+R(\gamma)\psi-|\sigma|_\gamma^2\psi^{-7}
 +\Bigl(\tfrac23\tau^2-2\Lambda\Bigr)\psi^5=0.                          \tag{45.4}
\]

\begin{theorem}[Single-valued solver near the plane datum]
\label{thm:v19v42-solver}
Fix \(\Lambda_*\in(-\Lambda_0,0)\) and \(h_*\in{\cal A}(\Lambda_*)\),
with plane solution \(\psi_*=1+u(h_*)\).  There is a neighbourhood
\({\cal W}\) of \((\delta+h_*,0,0)\) in
\(C^{2,\alpha}\hbox{-metrics}\times C^{0,\alpha}\hbox{-tensors}\times\mathbb R\)
on which (45.4) with \(\Lambda=\Lambda_*\) has a unique solution
\(\psi=\Psi_\Lambda(\gamma,\sigma,\tau)\) near \(\psi_*\), depending
analytically on \((\gamma,\sigma,\tau)\), and even in
\((\sigma,\tau)\mapsto(-\sigma,-\tau)\).
\end{theorem}

\begin{proof}
At the plane datum, \(\sigma=0\), \(\tau=0\), and the
\(\psi\)-derivative of (45.4) is

\[
 D_\psi{\cal L}_\Lambda=-8\Delta_{\gamma_*}+c,\qquad
 c:=R(\gamma_*)-10\Lambda_*\psi_*^4.                                   \tag{45.5}
\]

With \(\|h_*\|^2\) of order \(|\Lambda_*|\) by (45.3),
\(R(\gamma_*)=O(|\Lambda_*|)\), \(\psi_*=1+O(|\Lambda_*|)\), so
\(c=O(|\Lambda_*|)\) and \(\Delta_{\gamma_*}-\Delta=O(|\Lambda_*|^{1/2})\)
as operators \(C^{2,\alpha}\to C^{0,\alpha}\).  The operator (45.5) is
self-adjoint on \(L^2(dV_{\gamma_*})\) and is a small perturbation of
\(-8\Delta\), whose spectrum is \(\{0\}\cup[8|k_1|^2,\infty)\) with
simple eigenvalue \(0\).  Its lowest eigenvalue is therefore
\(\langle c\rangle_{\gamma_*}+O(|\Lambda_*|^2)\), and integrating
(44.2) at the plane, \(\int R(\gamma_*)\psi_*dV=2\Lambda_*\int\psi_*^5dV\),
gives \(\langle c\rangle=2\Lambda_*-10\Lambda_*+O(|\Lambda_*|^2)=-8\Lambda_*+O(|\Lambda_*|^2)>0\).
The remaining spectrum stays above \(8|k_1|^2-O(|\Lambda_*|^{1/2})>0\).
Hence (45.5) is an isomorphism \(C^{2,\alpha}\to C^{0,\alpha}\), and the
analytic implicit-function theorem gives \(\Psi_\Lambda\).  Evenness
follows from uniqueness, since (45.4) depends on \((\sigma,\tau)\) only
through \(|\sigma|^2\) and \(\tau^2\).
\end{proof}

For the Fourier wave, \(c=\lambda^2(1-\tfrac74\cos2z)\) with
\(\Lambda_*=-\lambda^2/8\); the lowest eigenvalue of
\(-8\partial_z^2+c\) on even functions is
\(\lambda^2-\tfrac{49}{1024}\lambda^4+O(\lambda^6)\), positive for
\(\lambda^2<1024/49\) at this order.

\subsection{The seeded \(\Lambda<0\) chain}

\begin{proposition}[Seeded reflection-symmetric constrained chain]
\label{prop:v19v42-chain}
Fix \(\Lambda_*\) and \(h_*\) as above, and a neighbourhood
\({\cal W}\) from Theorem \ref{thm:v19v42-solver}.  Let \(q_n=(h_n,\phi_n)\)
be the reflection-symmetric chain of Theorem \ref{thm:v19v38-chain}
with \(q_0=(h_0,\phi_0)\), \(h_0\) distributed on \({\cal A}(\Lambda_*)\)
near \(h_*\), and forward kernels \(P_n\) of bounded support, so small
that for \(|n|\leq n_0\) the link data

\[
 \gamma_{n+1/2}:=\delta+\tfrac12(h_n+h_{n+1}),\qquad
 \sigma_{n+1/2}:=\Pi_{\gamma}\frac{h_{n+1}-h_n}{\varepsilon},\qquad
 \tau_{n+1/2}:=\frac{\phi_{n+1}-\phi_n}{\varepsilon}                    \tag{45.6}
\]

lie in \({\cal W}\), where \(\Pi_\gamma\) is the \(\gamma\)-TT
projection.  Then the link-wise constrained data
\((g,K)_{n+1/2}\) obtained from \(\Psi_\Lambda\) are exact solutions of
(42.4) at the fixed \(\Lambda_*\) for \(|n|<n_0\), the map from seeds to
data is equivariant and positive-time local, and the pushforward law is
reflection positive on the algebra of links \(0\leq n<n_0\).
\end{proposition}

\begin{proof}
Bounded noise keeps finitely many steps inside \({\cal W}\).  Under the
site reflection, \(\gamma_{n+1/2}\) is even and both quantities in
(45.6) built from differences are odd; \(\Pi_\gamma\) is linear and
depends only on the even \(\gamma\).  Evenness of \(\Psi_\Lambda\) in
\((\sigma,\tau)\) makes \(g\) even and \(K\) odd, so the link-wise map
is equivariant; it is local because it uses sites \(n,n+1\).  Theorem
\ref{thm:v19v38-chain} gives reflection positivity of the seed and
Theorem \ref{thm:v19v32-pushforward} transports it; the restriction to
\(n<n_0\) is a closed subspace of \({\cal H}_+\).
\end{proof}

At the plane the seed is time symmetric in law, the classical datum is
a bounce slice of V41 type, and the mean curvature changes sign across
the plane by Proposition \ref{prop:v19v39-bounce}.  This is the
reflection-positive compact frame with the correct classical plane; its
dynamics remain those of the seed.

\subsection{Executable certificate}

The verifier

\[
 \texttt{scripts/verify\_sm\_v19\_v42\_fixed\_lambda\_sphere.py}      \tag{45.7}
\]

checks over the rationals that three rational points lie on the
leading-order admissible circle \(x^2+y^2=1\) for \(\Lambda_*=-1/8\) in
the two-mode space, that the radial derivative is nonzero, that the
plane coefficient \(c\) has mean \(-8\Lambda_*\), that the second-order
lowest eigenvalue is \(\lambda^2-\tfrac{49}{1024}\lambda^4\) and
positive at \(\lambda^2=1/4\), and that the Lichnerowicz coefficient
at the plane is \(-2\Lambda_*>0\).  It emits

\[
 \texttt{artifacts/sm\_v19\_v42\_fixed\_lambda\_sphere.json}.        \tag{45.8}
\]

The hashes are

\[
\begin{array}{c|l}
\hbox{object}&\hbox{SHA--256}\\ \hline
\hbox{verifier}&
\texttt{A9CDEEB0B52A6134CAC785D3A4F7961915CE77F4D93CF324D0E0E2DB4BEADDCF}\\
\hbox{canonical payload}&
\texttt{8ADAA2F6BE1F5FAFC8A9E2F0E79DFE0962AE2C680B35171053F86326010985EC}.
\end{array}                                                     \tag{45.9}
\]

\begin{firewall}[Finite modes, finite steps, seed dynamics]
\label{fire:v19v42}
Theorem \ref{thm:v19v42-sphere} is proved in a finite-dimensional
mode space; the full TT space has no compact sphere.  Theorem
\ref{thm:v19v42-solver} is local near one plane datum, and Proposition
\ref{prop:v19v42-chain} controls only finitely many steps with bounded
noise.  The kernels \(P_n\) are unspecified, so the chain carries no
Einstein dynamics, and no lattice-spacing or mode-number limit is
taken.  No Standard Model field is present.
\end{firewall}

\begin{decision}[V43 acquisition]
\label{dec:v19v42-next}
V43 will impose the first Einstein rate on the seeded chain: it will
require the expected mean-curvature increment across the first link to
equal \(\varepsilon\,\dot\tau(0)=2\varepsilon\Lambda_*\langle N\rangle\)
from (42.6) with unit lapse, and the expected metric step to equal
\(-2\varepsilon E[K]\), and determine the mean of the forward kernel
\(P_0\) that achieves both at second order in the wave amplitude.  This
is the first entry of Einstein dynamics into the seed law.
\end{decision}

\begin{assessment}[V42 acquisition]
\label{ass:v19v42}
At fixed \(\Lambda<0\) the classical bounce slices form an analytic
sphere of Dirichlet radius \(-8\Lambda\) in each finite mode space, the
\(\Lambda\)-Lichnerowicz equation is uniquely and analytically solvable
near every such slice because its plane linearization has lowest
eigenvalue \(-8\Lambda\), and the reflection-symmetric chain can be
seeded on the sphere and pushed forward with exact constraints and
reflection positivity.  The compact \(\Lambda<0\) frame is now
complete at the kinematic level.
\end{assessment}

\begin{record}[V42 append record]
\label{rec:v19v42}
V42 is appended to the validated V41 whose installed SHA--256 was
\texttt{91BC631BBCB450ACF4E6C5F9D7938C3E87EFA0AACE1C19C11970B21417DCCA76}.
After installation only V42 is the active numeric SM note; frozen
archives and unrelated notes are unchanged.
\end{record}


\section{V43: corrected mean-curvature rate with \(\Lambda\), the turnaround, and the first Einstein step}
\label{sec:v19v43}

While preparing the first Einstein rate for the seeded chain, the
cosmological term in the mean-curvature evolution (42.4) was rederived
and found to carry the wrong coefficient.  This note records the
correction, which reverses the sign conclusion of Proposition
\ref{prop:v19v39-bounce}: with \(\Lambda<0\) the reflection plane of
the torus programme is a strict volume maximum, a turnaround, as it is
in vacuum, and the negative constant strengthens focusing rather than
producing a bounce.  Nothing else in V39--V42 depends on that sign.
The note then computes the exact Taylor coefficients of the Einstein
development at the plane and shows that, at linear order in the wave,
the first step of a reflection-symmetric chain consistent with them is
the discrete free oscillator \(h_1=\cos(\varepsilon\omega)h_0\) on each
mode, together with a homogeneous mean-curvature drift
\(-\varepsilon\Lambda/2\).

\subsection{Erratum: the cosmological term in the mean-curvature evolution}

With cosmological constant the vacuum ADM evolution of the extrinsic
curvature is

\[
 \partial_tK_{ij}=N\bigl(R_{ij}-2K_{ik}K^k{}_j+\tau K_{ij}-\Lambda g_{ij}\bigr)
 -\nabla_i\nabla_jN+({\cal L}_XK)_{ij},                              \tag{46.1}
\]

because the effective stress of \(\Lambda\) has \(\rho=\Lambda/8\pi G\),
\(S_{ij}=-(\Lambda/8\pi G)g_{ij}\), hence
\(S_{ij}-\tfrac12g_{ij}(S-\rho)=(\Lambda/8\pi G)g_{ij}\).

\begin{proposition}[Corrected rate]
\label{prop:v19v43-rate}
Under (46.1) and the constraint \(R+\tau^2-|K|^2=2\Lambda\),

\[
 \partial_t\tau=N\bigl(|K|_g^2-\Lambda\bigr)-\Delta_gN+X(\tau).        \tag{46.2}
\]

The second formula in (42.4) and the statements (42.6) are to be
replaced by (46.2) and by

\[
 \dot\tau(0)=-\Lambda\langle N\rangle,\qquad
 \ddot{\operatorname{Vol}}(0)=\Lambda\langle N\rangle\int N\,dV,        \tag{46.3}
\]

at a reflection plane \(K(0)=0\) with the CMC lapse.  For
\(\Lambda<0\) the plane is a strict local maximum of the volume, the
mean curvature is negative before and positive after it, and the
development expands to the plane and recollapses.
\end{proposition}

\begin{proof}
Repeat the proof of Proposition \ref{prop:v19v36-monotone} with
(46.1): the trace of the new term is \(-3N\Lambda\), so
\(\partial_t\tau=N(R+\tau^2-3\Lambda)-\Delta_gN+X(\tau)\), and the
constraint gives (46.2).  At the plane, integrating (46.2) with
\(\partial_t\tau\) spatially constant gives
\(\dot\tau(0)\operatorname{Vol}=-\Lambda\int N\,dV\), and (42.1) gives
\(\ddot{\operatorname{Vol}}(0)=-\dot\tau(0)\int N\,dV\).  The sign
statements follow from \(\Lambda<0\).
\end{proof}

The anti-de Sitter metric \(-dt^2+\cos^2(t/\ell)\,g_{\mathbb H^3(\ell)}\)
with \(\Lambda=-3/\ell^2\) confirms (46.2) exactly: \(K(0)=0\),
\(R=-6/\ell^2=2\Lambda\), \(\tau=3\ell^{-1}\tan(t/\ell)\),
\(|K|^2=3\ell^{-2}\tan^2(t/\ell)\), and
\(\dot\tau=3\ell^{-2}\sec^2(t/\ell)=|K|^2-\Lambda\), with volume
\(\propto\cos^3(t/\ell)\) maximal at \(t=0\).  The certificate below
checks this at rational values of \(\tan(t/\ell)\).

\begin{corollary}[The torus reflection plane is always a turnaround]
\label{cor:v19v43-turnaround}
For time-symmetric torus data with any \(\Lambda\): if \(\Lambda>0\)
no non-flat data exist, since \(R=2\Lambda>0\) contradicts toroidal
rigidity; if \(\Lambda=0\) the data are flat and static; if
\(\Lambda<0\) the data of V41--V42 exist and their development has a
strict volume maximum at the plane.  Thus the word bounce in V38, V39,
V41, and V42 is to be read as maximal-volume turnaround throughout.
\end{corollary}

\begin{proof}
Theorem \ref{thm:v19v39-rigidity} covers \(\Lambda\geq0\), and
Proposition \ref{prop:v19v43-rate} covers \(\Lambda<0\).
\end{proof}

\subsection{Taylor coefficients of the development at the plane}

\begin{proposition}[Second-order development from time-symmetric data]
\label{prop:v19v43-taylor}
Let \((g_*,0)\) be time-symmetric data with constant \(\Lambda\), and
let the development use \(N=1\), \(X=0\).  Then

\[
 g(t)=g_*-t^2\bigl(\operatorname{Ric}(g_*)-\Lambda g_*\bigr)+O(t^3),
 \qquad
 \tau(t)=-\Lambda t+O(t^2),\qquad
 K(t)=t\bigl(\operatorname{Ric}(g_*)-\Lambda g_*\bigr)+O(t^2).            \tag{46.4}
\]

If \(g_*=\psi_*^4(\delta+\lambda h)\) is the V41 datum with \(h\) a TT
Fourier mode of frequency \(|k|=\omega\), then to first order in
\(\lambda\)

\[
 \operatorname{Ric}(g_*)-\Lambda g_*=\frac{\lambda\omega^2}{2}h+O(\lambda^2),
 \qquad
 g(t)=\delta+\lambda\Bigl(1-\frac{\omega^2t^2}2\Bigr)h+O(\lambda^2)+O(t^3).
                                                                    \tag{46.5}
\]
\end{proposition}

\begin{proof}
With \(N=1\), \(X=0\): \(\dot g=-2K\), and (46.1) at \(K=0\) gives
\(\dot K(0)=\operatorname{Ric}(g_*)-\Lambda g_*\); hence
\(\ddot g(0)=-2\dot K(0)\), which is (46.4), and
\(\dot\tau(0)=R-3\Lambda=-\Lambda\).  For (46.5), (29.7) gives
\(\operatorname{Ric}(\delta+\lambda h)=-\tfrac\lambda2\Delta h+O(\lambda^2)
=\tfrac{\lambda\omega^2}2h+O(\lambda^2)\) for a TT mode, while
\(\Lambda g_*=O(\lambda^2)\) and \(\psi_*=1+O(\lambda^2)\).
\end{proof}

Equation (46.5) is the wave equation \(\ddot h=\Delta h\) at the plane,
that is, the V23 oscillator, with the time-symmetric initial condition
\(\dot h(0)=0\).

\subsection{The first Einstein step of a reflection-symmetric chain}

\begin{theorem}[First-step consistency]
\label{thm:v19v43-step}
Let \((q_n)=(h_n,\phi_n)\) be a reflection-symmetric chain of
Corollary \ref{cor:v19v38-gaussian} type on the finite TT mode space
\(V_N\oplus\mathbb R\), seeded at the plane on the admissible set of
Theorem \ref{thm:v19v42-sphere}, with lattice spacing \(\varepsilon\).
Require that the conditional mean of the first step match (46.4) to
order \(\varepsilon^2\) at linear order in the wave amplitude, and that
the expected first-link mean curvature match \(\tau(\varepsilon/2)\).
Then, mode by mode with frequency \(\omega_k\),

\[
 E[h_1\mid h_0]=\Bigl(1-\frac{\varepsilon^2\omega_k^2}2\Bigr)h_0+O(\varepsilon^3)
 =\cos(\varepsilon\omega_k)h_0+O(\varepsilon^3),
 \qquad
 E[\phi_1-\phi_0]=-\frac{\varepsilon^2\Lambda}{2}+O(\varepsilon^3).     \tag{46.6}
\]

Conversely the symmetric Verlet recursion
\(h_{n+1}=2\cos(\varepsilon\omega_k)h_n-h_{n-1}\) with the reflection
condition \(h_{-1}=h_1\) has exactly \(h_1=\cos(\varepsilon\omega_k)h_0\),
so the first step (46.6) is the one-step restriction of the exact
discrete free oscillator, whose transfer matrix is that of the V23
covariance.
\end{theorem}

\begin{proof}
The site \(1\) sits at \(t=\varepsilon\); (46.5) with
\(t=\varepsilon\) gives the first relation.  The link \(1/2\) sits at
\(t=\varepsilon/2\), and (46.4) gives
\(\tau(\varepsilon/2)=-\Lambda\varepsilon/2+O(\varepsilon^2)\); since
\(\tau_{1/2}=(\phi_1-\phi_0)/\varepsilon\), the second relation follows.
For the converse, put \(n=0\) in the Verlet recursion and use
\(h_{-1}=h_1\).
\end{proof}

Thus the Einstein dynamics enters the seed law for the first time, and
at linear order it selects precisely the free TT oscillator of V23 with
time-symmetric phase, while the cosmological constant supplies a
deterministic positive drift of the homogeneous mean curvature.  The
noise covariance of the first step is not fixed by these mean
conditions; it is the quantum datum.

\subsection{Executable certificate}

The verifier

\[
 \texttt{scripts/verify\_sm\_v19\_v43\_first\_einstein\_step.py}    \tag{46.7}
\]

checks (46.2) on the anti-de Sitter slicing for three rational
\(\ell\) and three rational values of \(\tan(t/\ell)\), including that
the uncorrected rate fails; the plane constraint \(R=2\Lambda\); the
negativity of (46.3) for \(\Lambda=-1/8\); the order-\(\varepsilon^2\)
agreement of the first-step coefficient with \(\cos(\varepsilon\omega)\)
and the symmetric Verlet identity; and the homogeneous drift
\(E[\phi_1-\phi_0]=-\varepsilon^2\Lambda/2\).  It emits

\[
 \texttt{artifacts/sm\_v19\_v43\_first\_einstein\_step.json}.        \tag{46.8}
\]

The hashes are

\[
\begin{array}{c|l}
\hbox{object}&\hbox{SHA--256}\\ \hline
\hbox{verifier}&
\texttt{7247EC0D8EDCFA309167354241F3836375FB04F662B30BA163FE29A59B6FF6BA}\\
\hbox{canonical payload}&
\texttt{55F7AC2677E9E0D270B71DCF9B4264CA61341E9BAB8EBCD35937710159C0D37A}.
\end{array}                                                     \tag{46.9}
\]

\begin{firewall}[Scope of the correction and of the first step]
\label{fire:v19v43}
The erratum changes only the sign discussion around (42.6); the
existence theorems of V41--V42 and all reflection-positivity results are
unaffected.  Theorem \ref{thm:v19v43-step} fixes conditional means at
one step and at linear order in the wave; it does not fix the noise,
higher steps, second-order terms, or a continuum limit, and it does not
construct a measure with Einstein dynamics.
\end{firewall}

\begin{decision}[V44 acquisition]
\label{dec:v19v43-next}
V44 will extend the first-step consistency to all steps at linear order
by proving that a reflection-symmetric Gaussian chain whose conditional
means follow the symmetric Verlet recursion and whose noise is
stationary in the oscillator's action--angle variables has the V23
reflected covariance in the limit \(\varepsilon\to0\), and it will
determine whether the second-order homogeneous drift
\(E[\phi_{n+1}-\phi_n]\) can be made to follow (46.2) in expectation
with a covariance that stays reflection positive.  The V45 audit
follows.
\end{decision}

\begin{assessment}[V43 acquisition]
\label{ass:v19v43}
The cosmological term in the mean-curvature rate is corrected, which
makes every torus reflection plane a maximal-volume turnaround: for
\(\Lambda>0\) no data, for \(\Lambda=0\) flat, for \(\Lambda<0\) the
V41--V42 slices with recollapse.  The first Einstein step at the plane
is exact algebra and selects the discrete free oscillator on TT modes
plus a deterministic homogeneous drift, so the Einstein dynamics has now
entered the seed law at the level of conditional means.
\end{assessment}

\begin{record}[V43 append record]
\label{rec:v19v43}
V43 is appended to the validated V42 whose installed SHA--256 was
\texttt{BB6316FC479130EF070A6159D3AD9318E7A9F6528FF6091F7CD18C989BFEEC4E}.
After installation only V43 is the active numeric SM note; frozen
archives and unrelated notes are unchanged.
\end{record}


\section{V44: Euclidean throat, Lorentzian turnaround, and the exact spectral consistency of the linear seed}
\label{sec:v19v44}

V43 imposed the Lorentzian Taylor coefficients of the development on
the conditional mean of the Euclidean chain.  That is the wrong object:
the conditional mean of a Euclidean Markov chain is the decaying
Euclidean classical solution, and the Lorentzian development is
recovered from the Euclidean process by analytic continuation of its
correlation functions, not from its conditional means.  This note
records the correct relation.  At the reflection plane the Euclidean
and Lorentzian developments from the same time-symmetric data have
opposite second derivatives: the plane is a minimal slice, a throat, of
the Euclidean geometry and a maximal slice of the Lorentzian one.  The
correct linear-order consistency is spectral, and the reflection-symmetric
Ornstein--Uhlenbeck chain satisfies it exactly at every lattice
spacing: its reconstructed Hamiltonian is the oscillator with the
Einstein frequency on each TT mode.

\subsection{Euclidean development from time-symmetric data}

For a Riemannian four-manifold with \(\operatorname{Ric}=\Lambda g\),
foliated with unit normal and the convention (39.1), the Gauss--Codazzi
constraint and the extrinsic-curvature evolution read

\[
 R-\tau^2+|K|^2=2\Lambda,\qquad
 \partial_tK_{ij}=N\bigl(-R_{ij}+\Lambda g_{ij}+2K_{ik}K^k{}_j-\tau K_{ij}\bigr)
 +\nabla_i\nabla_jN+({\cal L}_XK)_{ij}.                               \tag{47.1}
\]

\begin{proposition}[Throat and turnaround]
\label{prop:v19v44-throat}
Let \((g_*,0)\) be time-symmetric data with \(R(g_*)=2\Lambda\), and
let \(g_E(t)\), \(g_L(t)\) be the Euclidean and Lorentzian developments
with \(N=1\), \(X=0\).  Then

\[
 \ddot g_E(0)=2\bigl(\operatorname{Ric}(g_*)-\Lambda g_*\bigr)=-\ddot g_L(0),
 \qquad
 \dot\tau_E(0)=\Lambda=-\dot\tau_L(0),
 \qquad
 \ddot{\operatorname{Vol}}_E(0)=-\Lambda\operatorname{Vol}=-\ddot{\operatorname{Vol}}_L(0).
                                                                    \tag{47.2}
\]

For \(\Lambda<0\) the plane is a strict volume minimum of the Euclidean
geometry and a strict volume maximum of the Lorentzian one.
\end{proposition}

\begin{proof}
From (47.1) at \(K=0\), \(\dot K_E=-\operatorname{Ric}+\Lambda g_*\)
and \(\ddot g_E=-2\dot K_E\); compare with (46.4).  The trace gives
\(\dot\tau_E=-R+3\Lambda=\Lambda\) by the constraint, and (42.1) gives
the volume.
\end{proof}

The anti-de Sitter and hyperbolic slicings
\(-dt^2+\cos^2(t/\ell)g_{\mathbb H^3}\) and
\(dt^2+\cosh^2(t/\ell)g_{\mathbb H^3}\), both with \(\Lambda=-3/\ell^2\),
realize (47.2) exactly.  For the TT wave (46.5), the Euclidean
development is \(h_E(t)=\cosh(\omega t)h_0+O(\lambda^2)\), the Wick
rotation of \(\cos(\omega t)h_0\).

\subsection{What a Euclidean chain's conditional mean is}

\begin{proposition}[Conditional means are decaying Euclidean solutions]
\label{prop:v19v44-mean}
For one TT mode of frequency \(\omega\), let the reflection-symmetric
Gaussian chain of Corollary \ref{cor:v19v38-gaussian} have forward
kernel \(N(rq,v(1-r^2))\) with \(r=e^{-\varepsilon\omega}\) and
\(v=1/(2c_g\omega)\).  Then \(E[q_n\mid q_0]=e^{-n\varepsilon\omega}q_0\)
for \(n\geq0\), which is the Euclidean classical solution
\(\ddot q=\omega^2q\) that decays away from the plane in both
directions; it is neither \(\cos(n\varepsilon\omega)q_0\) nor
\(\cosh(n\varepsilon\omega)q_0\).  Condition (46.6) is therefore not a
consistency requirement for a Euclidean seed, and Theorem
\ref{thm:v19v43-step} is to be read only as the identification of the
Lorentzian symmetric Verlet step with (46.5).
\end{proposition}

\begin{proof}
The forward kernel gives \(E[q_{n+1}\mid q_n]=rq_n\) and the tower
property gives \(r^nq_0\).  The mirror construction gives the same
decay for \(n<0\).  The symmetric solutions of the Euclidean and
Lorentzian mode equations through \(q_0\) with zero velocity are
\(\cosh\) and \(\cos\); the decaying solution is the one selected by
integrating out the future.
\end{proof}

Likewise the homogeneous drift of a Euclidean chain, if it is to follow
the Euclidean classical solution in mean, is
\(E[\phi_1-\phi_0]=\varepsilon^2\Lambda/2+O(\varepsilon^3)<0\) for
\(\Lambda<0\), the opposite sign to (46.6), by (47.2).

\subsection{Exact spectral consistency of the Ornstein--Uhlenbeck seed}

\begin{theorem}[Reconstructed Hamiltonian of the linear seed]
\label{thm:v19v44-spectrum}
For the chain of Proposition \ref{prop:v19v44-mean}, the reconstruction
of Theorem \ref{thm:v19v38-reconstruction} gives
\({\cal H}=L^2(\mathbb R,N(0,v))\) and one transfer contraction
\(P\) with

\[
 P\,\mathrm{He}_n\bigl(q/\sqrt v\bigr)=e^{-n\varepsilon\omega}\,\mathrm{He}_n\bigl(q/\sqrt v\bigr)
 \qquad(n=0,1,2,\dots),                                              \tag{47.3}
\]

where \(\mathrm{He}_n\) are the probabilists' Hermite polynomials.
Hence \(P=e^{-\varepsilon H}\) with \(H=\omega\,\widehat N\), the
oscillator Hamiltonian with spectrum \(\{0,\omega,2\omega,\dots\}\), at
every lattice spacing \(\varepsilon\).  For the finite TT mode space
\(V_N\) the reconstructed Hamiltonian is \(\sum_k\omega_k\widehat N_k\)
with \(\omega_k=|k|\), which is the Fock Hamiltonian of the linearized
Einstein sector of Corollary \ref{cor:v19v23-transfer}.
\end{theorem}

\begin{proof}
The Ornstein--Uhlenbeck kernel acts on a polynomial \(f\) by
\((Pf)(q)=E[f(rq+\sqrt{v(1-r^2)}\,\xi)]\) with \(\xi\) standard normal.
The generating function of the Hermite polynomials,
\(\sum_n\mathrm{He}_n(x)s^n/n!=e^{sx-s^2/2}\), gives
\(E[e^{s(rx+\sqrt{1-r^2}\xi)-s^2/2}]=e^{srx-s^2r^2/2}\), whose
coefficient of \(s^n/n!\) is \(r^n\mathrm{He}_n(x)\); this is (47.3)
after scaling by \(\sqrt v\).  Hermite polynomials form an orthogonal
basis of \(L^2(N(0,v))\), so \(P\) is diagonal with eigenvalues
\(r^n=e^{-n\varepsilon\omega}\).  Independent modes give the tensor
product.
\end{proof}

\begin{corollary}[Linear Einstein consistency of the seed is exact and spectral]
\label{cor:v19v44-consistency}
A reflection-symmetric Gaussian chain on \(V_N\) reconstructs the
linearized Einstein Hamiltonian if and only if its forward kernel on
each mode has \(r_k=e^{-\varepsilon|k|}\); the stationary variance
\(v_k\) fixes only the normalization of the mode.  No condition on
conditional means beyond this is required or permitted.
\end{corollary}

\begin{proof}
By (47.3) the spectrum is \(\{n\log(1/r_k)/\varepsilon\}\), which
equals \(\{n|k|\}\) exactly when \(r_k=e^{-\varepsilon|k|}\); the
variance enters (47.3) only through the argument scaling.
\end{proof}

Thus the seed already carries the Einstein dynamics at linear order,
exactly, and the analytic continuation of its two-point function
\(v\,e^{-\omega|t|}\) to \(v\,e^{-i\omega t}\) is where the Lorentzian
\(\cos(\omega t)\) of (46.5) lives.  What the seed does not carry is the
nonlinear content: the constraint sphere of V42 at the plane, and the
second-order rates.

\subsection{Executable certificate}

The verifier

\[
 \texttt{scripts/verify\_sm\_v19\_v44\_ou\_spectrum\_throat.py}     \tag{47.4}
\]

implements exact Gaussian-moment polynomial arithmetic and checks
(47.3) for \(n\leq5\) at \(r=1/2\); checks on the anti-de Sitter and
hyperbolic slicings with \(\ell=2\) that the second derivatives of the
scale factor are opposite, that \(\dot\tau_L(0)=-\Lambda\) and
\(\dot\tau_E(0)=\Lambda\), and that the Euclidean volume is a minimum
while the Lorentzian one is a maximum; and records the negative
Euclidean homogeneous drift.  It emits

\[
 \texttt{artifacts/sm\_v19\_v44\_ou\_spectrum\_throat.json}.        \tag{47.5}
\]

The hashes are

\[
\begin{array}{c|l}
\hbox{object}&\hbox{SHA--256}\\ \hline
\hbox{verifier}&
\texttt{20C85C6DACC42841E7A2BFC177B978201E2D12DA7EBC87B17A573BD1F8BB49D8}\\
\hbox{canonical payload}&
\texttt{84F1DAF89F706146184039B5615C118998D13F014FE508D461CEE3EE5BA51E73}.
\end{array}                                                     \tag{47.6}
\]

\begin{firewall}[Linear only]
\label{fire:v19v44}
Theorem \ref{thm:v19v44-spectrum} is the free Fock sector; it contains
no interaction, no constraint sphere, no \(\Lambda\), and no
second-order rate.  Proposition \ref{prop:v19v44-throat} is classical.
The reinterpretation of V43 removes a condition; it adds no
construction.
\end{firewall}

\begin{decision}[V45 audit and V46 acquisition]
\label{dec:v19v44-next}
V45 is the compulsory companion audit.  V46 will address the first
nonlinear content of the seed at the plane: the plane law must be
supported on the admissible sphere \({\cal A}(\Lambda_*)\) of V42, which
no Gaussian is.  It will prove that conditioning a reflection-symmetric
Gaussian chain on \(\{h_0\in{\cal A}(\Lambda_*)\}\), a \(\theta\)-invariant
event, preserves reflection positivity and the Markov structure, and
will compute the resulting plane law and its first two moments in the
two-mode space.
\end{decision}

\begin{assessment}[V44 acquisition]
\label{ass:v19v44}
The relation between the Euclidean seed and the Lorentzian development
is now exact: opposite second derivatives at the plane, a Euclidean
throat continuing to a Lorentzian turnaround, and a spectral rather
than mean-value consistency condition.  The reflection-symmetric
Ornstein--Uhlenbeck seed reconstructs the linearized Einstein
Hamiltonian exactly at every lattice spacing.  V43's first-step
condition is withdrawn as a requirement and retained as the Lorentzian
Verlet identity.
\end{assessment}

\begin{record}[V44 append record]
\label{rec:v19v44}
V44 is appended to the validated V43 whose installed SHA--256 was
\texttt{B09BC5B4AC12895EB719114F296AD9950F8C56AEB062AEAB5180645D767FD89F}.
After installation only V44 is the active numeric SM note; frozen
archives and unrelated notes are unchanged.
\end{record}


\section{V45: compulsory companion audit after the spectral consistency}
\label{sec:v19v45}

This is the required audit following V40.  Every active companion note
in the shared folder was inspected read-only before the constrained
plane law announced in V44 is constructed.

\begin{record}[V45 companion snapshot]
\label{rec:v19v45-snapshot}
The inspected files were

\[
\begin{array}{c|r|r}
\hbox{file and SHA--256}&\hbox{bytes}&\hbox{lines}\\ \hline
\texttt{Renewal\_Yang\_Mills\_Mass\_Gap\_Research\_Notes\_V32.tex}
&452376&11540\\
\multicolumn{3}{l}{
\texttt{F1DE1DC4AA157BC4DF069FEDB4F34C3032565655348CB9AE35AC95F5695744D3}}
\\
\texttt{Renewal\_NS\_Stability\_Research\_Notes\_V26.tex}
&278283&5319\\
\multicolumn{3}{l}{
\texttt{82C887F8C2C69E4F28FAD7C3DC8DE5B9099CB5A4BF431EBA1FFF3E91935978AD}}
\\
\texttt{Renewal\_GRH\_Cofinal\_Visibility\_Attack\_V51.tex}
&489091&15170\\
\multicolumn{3}{l}{
\texttt{61927F9477BD72FAD280D5EF4042919932C7687ED895DE4571A71F2D81450134}}
\\
\texttt{Renewal\_YM\_Parallel\_Research\_Notes\_V5.tex}
&54174&1351\\
\multicolumn{3}{l}{
\texttt{306F1BB84CFA8A8D47EA569EC17E9545F9867C8D32B548EC9D9C444B4F3C0512}}.
\end{array}                                                     \tag{48.1}
\]

Since V40, Yang--Mills has advanced from V31 to V32; the other three
files are unchanged.
\end{record}

\subsection{Companion findings}

Yang--Mills V32 turns its archived router Schur matrix into an exact
Gaussian renormalization recursion, proves that the recursion preserves
transversality, evaluates it exactly at level one and at level-two
corners, and proves that the blocked kernel keeps its long-wavelength
Maxwell coefficient while its corner values are rescaled by factors
between \(1/4\) and \(15/28\).  It thereby refutes its self-similar
tower hypothesis at tree level and turns to a smeared dyadic blocking.
Its firewall leaves the one-loop shell coefficient, any fixed point, the
Haar sign, and all continuum statements open.

Navier--Stokes V26, GRH V51, and Yang--Mills Parallel V5 are unchanged
since V40.

\begin{decision}[V45 audit decision]
\label{dec:v19v45-audit}
No companion theorem, numerical constant, or executable artifact is
imported.  The V44 direction stands: V46 constructs the constrained
plane law by conditioning the reflection-symmetric chain on the
admissible sphere.
\end{decision}

\begin{proof}
The Yang--Mills recursion concerns transverse Gaussian kernels on a
four-dimensional lattice under dyadic blocking; the present programme's
Gaussian object is the reflection-symmetric Ornstein--Uhlenbeck chain,
whose exact spectrum was established in Theorem
\ref{thm:v19v44-spectrum} without any blocking step.  The Yang--Mills
finding that blocking preserves the long-wavelength coefficient while
altering short-distance corners has no counterpart in a time-lattice
chain, which has no spatial blocking.  There is no shared operator or
hypothesis, and the V46 task, conditioning on a \(\theta\)-invariant
event, is a measure-theoretic step independent of every companion.
\end{proof}

\begin{assessment}[V45 acquisition]
\label{ass:v19v45}
The audit is current.  Since V40 the programme has constructed the
analytic \(\Lambda<0\) time-symmetric branch and its fixed-\(\Lambda\)
admissible sphere, corrected the cosmological mean-curvature rate,
identified the Euclidean throat with the Lorentzian turnaround, and
proved that the reflection-symmetric Ornstein--Uhlenbeck seed
reconstructs the linearized Einstein Hamiltonian exactly.
\end{assessment}

\begin{record}[V45 append record]
\label{rec:v19v45}
V45 is appended to the validated V44 whose installed SHA--256 was
\texttt{449282A4A5A0E5524DE1BBE4DF4C0D7F11FE426067F9D557EF6BF4B8C1C52361}.
After installation only V45 is the active numeric SM note; the
companions, frozen archives, and all unrelated notes are unchanged.  The
next compulsory companion audit is V50.
\end{record}


\section{V46: the constrained plane law by disintegration on the admissible sphere}
\label{sec:v19v46}

V44 established that the linear seed carries the Einstein Hamiltonian
exactly and that the missing content is nonlinear: at the plane the
Hamiltonian constraint with fixed \(\Lambda_*<0\) forces the even wave
\(h_0\) onto the admissible sphere \({\cal A}(\Lambda_*)\) of V42, which
carries no Gaussian mass.  This note conditions the reflection-symmetric
chain on that sphere by the coarea disintegration of the plane marginal
along the analytic function \(\Lambda(h)\).  Because the event depends
only on the plane variable, the conditioned law is again a
reflection-symmetric chain and remains reflection positive with the
same conditional-expectation formula.  In the two-mode space the plane
law is uniform on the circle to leading order, its moments are explicit,
and the constraint holds as an exact Ward identity in expectation.

\subsection{Conditioning on a plane event}

Let \(\mu\) be the reflection-symmetric chain of \((\pi_0,(P_n))\) as
in V38.

\begin{proposition}[Plane conditioning preserves the structure]
\label{prop:v19v46-conditioning}
Let \(\pi_0'\) be any probability measure on \(E\).  The
reflection-symmetric chain \(\mu'\) of \((\pi_0',(P_n))\) satisfies
Theorem \ref{thm:v19v38-chain}.  If \(\pi_0'=\pi_0(\cdot\mid B)\) for a
measurable \(B\) with \(\pi_0(B)>0\), then \(\mu'=\mu(\cdot\mid q_0\in B)\).
More generally, if \(\pi_0=\int\pi_0^s\,\rho(ds)\) is a disintegration
of \(\pi_0\) along a measurable map \(F:E\to\mathbb R\), so that
\(\pi_0^s\) is carried by \(\{F=s\}\), then \(\mu=\int\mu^s\rho(ds)\)
with \(\mu^s\) the reflection-symmetric chain of \((\pi_0^s,(P_n))\),
and each \(\mu^s\) is reflection positive.
\end{proposition}

\begin{proof}
The construction in V38 uses \(\pi_0\) only to draw \(q_0\); the
kernels and the mirror are unchanged, so Theorem
\ref{thm:v19v38-chain} applies to any plane law.  Conditioning
\(\mu\) on \(\{q_0\in B\}\) conditions only the draw of \(q_0\), by the
conditional independence of past and future given \(q_0\), which gives
\(\mu'\).  The disintegration statement is the same with
\(\pi_0(\cdot\mid B)\) replaced by the family \(\pi_0^s\), since
\(\mu=\int(\hbox{law of the chain given }q_0=x)\,\pi_0(dx)\) and the
inner law depends on \(x\) only.
\end{proof}

\subsection{The coarea plane law on the admissible sphere}

Take \(E=V_N\oplus\mathbb R\), \(\pi_0=N(0,\Sigma)\otimes N(0,v_0)\) with
\(\Sigma\) positive definite on \(V_N\), and the analytic map
\(\Lambda:U\to\mathbb R\) of Proposition \ref{prop:v19v42-analytic},
extended measurably outside \(U\).

\begin{theorem}[Constrained plane law]
\label{thm:v19v46-plane}
For \(\Lambda_*\in(-\Lambda_0,0)\), the disintegration of \(\pi_0\) along
\(\Lambda\) at the value \(\Lambda_*\) is

\[
 \pi_0^{\Lambda_*}(dh\,d\phi)
 =\frac1{Z(\Lambda_*)}\,
 \frac{e^{-\frac12\langle h,\Sigma^{-1}h\rangle}}{|\nabla\Lambda(h)|}\,
 \sigma_{{\cal A}(\Lambda_*)}(dh)\otimes N(0,v_0)(d\phi),               \tag{49.1}
\]

with \(\sigma_{{\cal A}}\) the surface measure of the analytic
hypersurface \({\cal A}(\Lambda_*)\) and \(Z\) the normalization.  The
reflection-symmetric chain \(\mu^{\Lambda_*}\) of
\((\pi_0^{\Lambda_*},(P_n))\) is reflection positive, and its plane
variable satisfies the Hamiltonian constraint at \(\Lambda_*\) almost
surely.  Its first moments satisfy the constraint Ward identity

\[
 E^{\Lambda_*}\bigl[Q(h_0)\bigr]=-8\Lambda_*+O(|\Lambda_*|^{3/2}).       \tag{49.2}
\]
\end{theorem}

\begin{proof}
On \(U\), \(\nabla\Lambda\neq0\) along \({\cal A}(\Lambda_*)\) by the
radial derivative in the proof of Theorem \ref{thm:v19v42-sphere}, so
the coarea formula disintegrates the Gaussian density along the level
sets of \(\Lambda\) with the weight \(|\nabla\Lambda|^{-1}\), which
gives (49.1); the \(\phi\) factor is untouched because \(\Lambda\) does
not depend on \(\phi\).  Reflection positivity is Proposition
\ref{prop:v19v46-conditioning}.  The constraint holds because
\({\cal A}(\Lambda_*)\) consists of exact time-symmetric data by
Theorem \ref{thm:v19v42-sphere}, and (49.2) is (45.3) integrated
against (49.1).
\end{proof}

\begin{corollary}[Two-mode space]
\label{cor:v19v46-two-mode}
For the two modes \(\cos z\,e^+\), \(\cos z\,e^\times\) with equal
variance, \(Q=x^2+y^2\) and \(\Sigma\propto I\).  To leading order in
\(\Lambda_*\), \({\cal A}(\Lambda_*)\) is the circle \(x^2+y^2=-8\Lambda_*\),
\(|\nabla\Lambda|=\sqrt{-8\Lambda_*}/4\) is constant on it, and
(49.1) is the uniform law on the circle.  Hence

\[
 E[x_0^2]=E[y_0^2]=-4\Lambda_*,\qquad
 E[x_0^4]=24\Lambda_*^2,\qquad
 E[x_0^2y_0^2]=8\Lambda_*^2,                                         \tag{49.3}
\]

all odd moments vanish, and \(E[Q(h_0)]=-8\Lambda_*\) exactly at this
order.  With the Ornstein--Uhlenbeck kernel of ratio \(r\) and
stationary variance \(v\) on each mode, the reflected Gram of
\(\{x_1,x_1^2\}\) is

\[
 \begin{pmatrix}
  -4r^2\Lambda_*&0\\
  0&24r^4\Lambda_*^2-8r^2v(1-r^2)\Lambda_*+v^2(1-r^2)^2
 \end{pmatrix}\geq0.                                                 \tag{49.4}
\]
\end{corollary}

\begin{proof}
Rotation invariance of the Gaussian and of the leading circle make the
weight constant.  The moments are the circle averages
\(\langle\cos^2\rangle=\tfrac12\), \(\langle\cos^4\rangle=\tfrac38\),
\(\langle\cos^2\sin^2\rangle=\tfrac18\) times powers of
\(r^2=-8\Lambda_*\).  For (49.4), \(Px=rx\) and
\(Px^2=r^2x^2+v(1-r^2)\); apply (41.1) with the uniform plane law.
\end{proof}

The conditioned chain is no longer Gaussian: its site laws are the
plane law transported by the Gaussian kernels, so \(q_n\) is a
Gaussian mixture over the sphere.  Its reconstructed space is
\(L^2(\pi_0^{\Lambda_*})\), which for the two-mode case is
\(L^2\) of the circle times \(L^2(N(0,v_0))\): the plane Hilbert space
of a constrained system with one angular degree of freedom per mode
pair.

\subsection{What conditioning does and does not achieve}

Conditioning makes the Hamiltonian constraint at fixed \(\Lambda_*\)
hold on the plane; it does not make it hold on the other links, where
the link-wise solver of Proposition \ref{prop:v19v42-chain} is still
needed, and it does not touch the momentum constraint, which the TT
seeds satisfy by construction.  Conditioning on the plane also does
not alter the reconstructed transfer contractions \(P_n\); it alters the
state.  The spectral consistency of V44 therefore survives, while the
plane state now satisfies the constraint Ward identity (49.2).

\subsection{Executable certificate}

The verifier

\[
 \texttt{scripts/verify\_sm\_v19\_v46\_constrained\_plane\_law.py}   \tag{49.5}
\]

checks (49.3) over the rationals for \(\Lambda_*=-1/8\), the constancy
of \(|\nabla\Lambda|=1/4\) on the unit circle, the Ward identity
\(E[Q]=-8\Lambda_*\), and the positivity of (49.4) with \(r=1/2\),
\(v=2\), together with the fact that the conditioned second moment
differs from the Gaussian one.  It emits

\[
 \texttt{artifacts/sm\_v19\_v46\_constrained\_plane\_law.json}.      \tag{49.6}
\]

The hashes are

\[
\begin{array}{c|l}
\hbox{object}&\hbox{SHA--256}\\ \hline
\hbox{verifier}&
\texttt{8CB1878A8465E529A9BC8A858E3A8994E9477D67F5DDBBF7105EDCAC2FDFA8DA}\\
\hbox{canonical payload}&
\texttt{4216B41627FA7166E89CD975124D3F7547BB6D88F5B6CF0742E829EB318D89BE}.
\end{array}                                                     \tag{49.7}
\]

\begin{firewall}[A constrained state on a free chain]
\label{fire:v19v46}
The kernels remain the free Ornstein--Uhlenbeck kernels; only the plane
state is constrained.  Off the plane the links are constrained only
after the solver of V42 is applied, and the dynamics of the chain is
still free.  The two-mode computation is leading order in
\(\Lambda_*\).  No continuum, infinite-mode, or Standard Model
statement is made.
\end{firewall}

\begin{decision}[V47 acquisition]
\label{dec:v19v46-next}
V47 will ask whether the constraint can be carried off the plane by
the chain itself rather than by the solver: it will compute the
expected constraint defect
\(E^{\Lambda_*}[R(\delta+h_n)+\tau_n^2-|K_n|^2-2\Lambda_*]\) at the first
link of the free chain at second order, identify the terms that the
free kernels cannot cancel, and state the minimal modification of the
forward kernel, a drift in the homogeneous mode, that cancels the
expected defect.  This is the first genuinely nonlinear dynamical
condition on the seed.
\end{decision}

\begin{assessment}[V46 acquisition]
\label{ass:v19v46}
The plane state of the compact \(\Lambda<0\) programme is now a
constrained, non-Gaussian, reflection-positive law obtained by
disintegrating the free plane marginal on the admissible sphere.  It
satisfies the Hamiltonian constraint exactly on the plane and the
constraint Ward identity in expectation, with explicit moments in the
two-mode space, while the linear spectral consistency of V44 is
retained.
\end{assessment}

\begin{record}[V46 append record]
\label{rec:v19v46}
V46 is appended to the validated V45 whose installed SHA--256 was
\texttt{E9C322B0E476200513927757EE7858F79CD36682CF1F1FE984D9175226118104}.
After installation only V46 is the active numeric SM note; frozen
archives and unrelated notes are unchanged.
\end{record}


\section{V47: the expected constraint defect at the first link and the zero-point balance}
\label{sec:v19v47}

V46 constrained the plane.  This note asks how far the free kernels
carry the Hamiltonian constraint to the first link without the solver.
The answer is exact: the spatially averaged expected defect at the
first link consists of a term of order \(\varepsilon^{-1}\), the
zero-point energy of the link momenta, and a term of order one, the
kink energy of the mirror-symmetric mean path together with the
cosmological constant.  Both can be cancelled by the homogeneous mode
alone: its stationary variance absorbs the zero-point term through the
quantum form of the V26 balance \(\tfrac23E[\tau^2]=E\langle|\sigma|^2\rangle\),
and its rate or a small drift absorbs the order-one residue.  This is
the first nonlinear dynamical condition on the seed and the first
appearance of the cosmological-constant problem in the programme: the
TT zero-point sum grows with the number of retained modes.

\subsection{Setting}

Let \(V_N\) be spanned by real TT modes \(\cos(k\cdot x)e^A\) with
frequencies \(\omega_k=|k|\), normalized so that for
\(h=\sum x_ke_k\), \(\langle|h|^2\rangle=\sum x_k^2\) and
\(Q(h)=\sum\omega_k^2x_k^2\).  Take the reflection-symmetric chain with
Ornstein--Uhlenbeck forward kernels of ratio \(r_k\) and stationary
variance \(v_k\) on each TT mode, ratio \(r_0\) and variance \(v_0\) on
the homogeneous scalar \(\phi\), and the constrained plane law of
Theorem \ref{thm:v19v46-plane}, whose second moments are
\(E[x_{k,0}^2]=:m_k\) with \(\sum\omega_k^2m_k=-8\Lambda_*\) at leading
order.  Define at the first link

\[
 \bar h=\tfrac12(h_0+h_1),\qquad
 \sigma=\frac{h_1-h_0}{\varepsilon},\qquad
 \tau=\frac{\phi_1-\phi_0}{\varepsilon},                              \tag{50.1}
\]

and the spatially averaged Hamiltonian-constraint defect, to second
order in the fields,

\[
 {\cal D}:=\bigl\langle R(\delta+\bar h)\bigr\rangle+\tfrac23\tau^2
 -\langle|\sigma|^2\rangle-2\Lambda_*.                                 \tag{50.2}
\]

\subsection{Exact expected defect}

\begin{theorem}[Expected defect at the first link]
\label{thm:v19v47-defect}
For every lattice spacing and every choice of the kernel parameters,

\[
 E[{\cal D}]
 =-\frac14\sum_k\omega_k^2\Bigl[\Bigl(\frac{1+r_k}2\Bigr)^2m_k+\frac{v_k(1-r_k^2)}4\Bigr]
 +\frac{4v_0(1-r_0)}{3\varepsilon^2}
 -\sum_k\frac{(1-r_k)^2m_k+v_k(1-r_k^2)}{\varepsilon^2}
 -2\Lambda_*.                                                          \tag{50.3}
\]

It is affine in \(v_0\) with positive slope, so for every \(r_0<1\)
there is exactly one \(v_0^*\) with \(E[{\cal D}]=0\), and
\(E[{\cal D}]<0\) when the homogeneous mode is absent.
\end{theorem}

\begin{proof}
Write \(h_1=r_kh_0+\xi\) mode by mode with \(\xi\) centred, independent
of \(h_0\), of variance \(v_k(1-r_k^2)\).  Then
\(\bar h_k=\tfrac{1+r_k}2x_{k,0}+\tfrac12\xi_k\) and
\(\sigma_k=\tfrac{r_k-1}\varepsilon x_{k,0}+\tfrac1\varepsilon\xi_k\),
whose second moments give the first and third sums, using
\(\langle R(\delta+\bar h)\rangle=-\tfrac14Q(\bar h)+O(|\bar h|^3)\)
from Lemma \ref{lem:v19v26-hessian} and the normalization of \(Q\).
Similarly \(E[\tau^2]=\bigl[(1-r_0)^2v_0+v_0(1-r_0^2)\bigr]/\varepsilon^2
=2v_0(1-r_0)/\varepsilon^2\).  All other terms in (50.2) are cubic or
vanish by symmetry of the plane law.  Affinity in \(v_0\) and the sign
statements are read off from (50.3) together with
\(-\tfrac14\sum\omega_k^2m_k=2\Lambda_*\).
\end{proof}

\subsection{Zero-point balance and the order-one residue}

\begin{corollary}[Small-spacing structure]
\label{cor:v19v47-structure}
Put \(r_k=e^{-\varepsilon\omega_k}\), \(r_0=e^{-\varepsilon\omega_0}\).
Then

\[
 E[{\cal D}]
 =\frac1\varepsilon\Bigl[\frac43v_0\omega_0-2\sum_kv_k\omega_k\Bigr]
 +\Bigl[-\frac23v_0\omega_0^2+2\sum_kv_k\omega_k^2+8\Lambda_*\Bigr]
 +O(\varepsilon).                                                       \tag{50.4}
\]

The singular term vanishes if and only if

\[
 \frac23\,v_0\omega_0=\sum_kv_k\omega_k,
 \qquad\hbox{equivalently}\qquad
 \frac23E[\tau^2]=E\langle|\sigma|^2\rangle+O(1),                        \tag{50.5}
\]

the quantum form of the V26 balance.  Under (50.5) the order-one term
is \(\sum_kv_k\omega_k(2\omega_k-\omega_0)+8\Lambda_*\), which vanishes
for the homogeneous rate

\[
 \omega_0=\frac{2\sum_kv_k\omega_k^2+8\Lambda_*}{\sum_kv_k\omega_k}
                                                                    \tag{50.6}
\]

whenever the right side is positive; otherwise a drift
\(E[\phi_1-\phi_0]=\varepsilon\beta\) with
\(\tfrac23\beta^2=-\bigl[\sum_kv_k\omega_k(2\omega_k-\omega_0)+8\Lambda_*\bigr]\)
cancels it, which is possible exactly when that bracket is nonpositive.
\end{corollary}

\begin{proof}
Expand \(1-r^2=2\varepsilon\omega-2\varepsilon^2\omega^2+O(\varepsilon^3)\),
\((1-r)^2=\varepsilon^2\omega^2+O(\varepsilon^3)\), and
\(((1+r)/2)^2=1+O(\varepsilon)\) in (50.3), and use
\(-\tfrac14\sum\omega_k^2m_k=2\Lambda_*\) and
\(\sum\omega_k^2m_k=-8\Lambda_*\).  A drift adds \(\beta^2\) to
\(E[\tau^2]\), hence \(\tfrac23\beta^2\) to \(E[{\cal D}]\).
\end{proof}

Three remarks fix the meaning of (50.4)--(50.6).  First, the
\(\varepsilon^{-1}\) term is the coincident-time singularity of
Corollary \ref{cor:v19v33-limit}: the link momenta are white at
coincidence, and a quadratic constraint evaluated on them has a
divergent expectation that only another quadratic white variable, the
homogeneous mean curvature, can cancel.  Second, with the V23
normalization \(v_k=1/(2c_g\omega_k)\) the sum \(\sum v_k\omega_k\) is
proportional to the number of retained modes, so the balancing
homogeneous variance grows without bound as modes are added: this is
the cosmological-constant problem in its exact lattice form, confined
here to a finite mode space.  Third, the order-one term contains
\(8\Lambda_*=-\sum\omega_k^2m_k\), the kink energy of the mirror-symmetric
mean path, whose slope at the plane is \(-\omega_kx_{k,0}\) rather than
zero; this is the Euclidean cusp of Proposition
\ref{prop:v19v44-mean}, not a Lorentzian velocity.

\subsection{Executable certificate}

The verifier

\[
 \texttt{scripts/verify\_sm\_v19\_v47\_first\_link\_defect.py}       \tag{50.7}
\]

evaluates (50.3) exactly for the two-mode chain with
\(\omega_k=1\), \(r_k=1/2\), \(v_k=2\), \(\varepsilon=1/4\),
\(\Lambda_*=-1/8\), the uniform plane law, and \(r_0=1/3\): it finds
\(E\langle|\sigma|^2\rangle=52\), \(E[Q(\bar h)]=21/16\), the balancing
variance \(v_0^*=29997/8192\), zero defect there, a negative defect
without the homogeneous mode, and checks the small-spacing expansion
of \((1-r^2)/\varepsilon^2\) on a rational surrogate.  It emits

\[
 \texttt{artifacts/sm\_v19\_v47\_first\_link\_defect.json}.          \tag{50.8}
\]

The hashes are

\[
\begin{array}{c|l}
\hbox{object}&\hbox{SHA--256}\\ \hline
\hbox{verifier}&
\texttt{1E11433AD19FAB1A7E7319812ADBAE60A45EE8A0C5F194FEAF3B2297CAB9F932}\\
\hbox{canonical payload}&
\texttt{01AA8AE6AAB18E9805AD00C73976015A6F20DEB1912DB45C3D3A3AFF6BC13D91}.
\end{array}                                                     \tag{50.9}
\]

\begin{firewall}[Expectation, second order, first link, finite modes]
\label{fire:v19v47}
Theorem \ref{thm:v19v47-defect} controls the expectation of the
averaged defect, not its variance, its pointwise value, or the momentum
constraint; it is second order in the fields and concerns one link.
The balance (50.5) is a finite-mode statement whose mode-number
divergence is exhibited, not resolved.  No measure with a continuum
limit is constructed.
\end{firewall}

\begin{decision}[V48 acquisition]
\label{dec:v19v47-next}
V48 will treat the variance of the defect: it will compute
\(\operatorname{Var}[{\cal D}]\) at the first link for the balanced chain,
show that it is of order \(\varepsilon^{-2}\) with a coefficient that
cannot be cancelled by any homogeneous Gaussian mode because the
fourth cumulant of a sum of independent white variables is positive,
and thereby prove that the free chain cannot satisfy the constraint
pathwise at the first link even in the balanced case.  This will fix
the precise sense in which the link solver of V42 is indispensable.
\end{decision}

\begin{assessment}[V47 acquisition]
\label{ass:v19v47}
The free kernels carry the constraint to the first link only in
expectation and only with a tuned homogeneous mode.  The expected
defect is exact, its singular part is the zero-point energy of the
white link momenta and is cancelled by the quantum V26 balance, and
its finite part is cancelled by the homogeneous rate or a drift.  The
mode-number growth of the balancing variance is the cosmological
constant problem in exact form.
\end{assessment}

\begin{record}[V47 append record]
\label{rec:v19v47}
V47 is appended to the validated V46 whose installed SHA--256 was
\texttt{587B3FA080707DBDEE7A821CD001DD5A42F6C6DFE1B9943E58B1153DD0A58002}.
After installation only V47 is the active numeric SM note; frozen
archives and unrelated notes are unchanged.
\end{record}


\section{V48: the defect variance and why the constraint cannot be a Ward identity}
\label{sec:v19v48}

V47 cancelled the expected constraint defect at the first link with a
tuned homogeneous mode.  This note shows that the cancellation cannot be
made pathwise by any homogeneous mode independent of the TT modes: the
variance of the defect has a lower bound of order \(\varepsilon^{-2}\)
coming from the TT noise alone, and variances of independent
contributions add.  Consequently a free chain with an independent
homogeneous seed violates the Hamiltonian constraint on almost every
path at the first link, however its parameters are tuned; only a
homogeneous variable that is a function of the TT link seeds, as in
the balanced choice (40.9) or the full solver of V42, can enforce it.
This fixes the precise sense in which the constraint is a pathwise
identity rather than a Ward identity in expectation.

\subsection{Variance lower bound}

In the setting of Theorem \ref{thm:v19v47-defect}, write the defect as

\[
 {\cal D}=A+B-2\Lambda_*,\qquad
 A:=-\tfrac14Q(\bar h)-\langle|\sigma|^2\rangle,\qquad
 B:=\tfrac23\tau^2,                                                    \tag{51.1}
\]

where \(A\) depends only on the TT variables and \(B\) only on the
homogeneous variables.

\begin{theorem}[Defect variance]
\label{thm:v19v48-variance}
Suppose the homogeneous variables are independent of the TT variables,
with any law.  Then

\[
 \operatorname{Var}[{\cal D}]\geq\operatorname{Var}[A]
 \geq2\sum_k\Bigl(\frac1{\varepsilon^2}+\frac{\omega_k^2}{16}\Bigr)^2v_k^2(1-r_k^2)^2
 \geq\frac{2}{\varepsilon^4}\sum_kv_k^2(1-r_k^2)^2.                      \tag{51.2}
\]

With \(r_k=e^{-\varepsilon\omega_k}\) the last bound is
\(8\varepsilon^{-2}\sum_kv_k^2\omega_k^2+O(\varepsilon^{-1})\).  In
particular the defect is not almost surely zero, and
\(\Pr[|{\cal D}|\geq c/\varepsilon]>0\) for some \(c>0\) independent of
\(\varepsilon\).
\end{theorem}

\begin{proof}
Independence gives \(\operatorname{Var}[A+B]=\operatorname{Var}A+\operatorname{Var}B\geq\operatorname{Var}A\).
For \(\operatorname{Var}A\), condition on the plane variables
\(x=(x_{k,0})\).  Mode by mode, with \(\xi_k\) the centred Gaussian
noise of variance \(s_k^2=v_k(1-r_k^2)\),

\[
 A=\sum_k\bigl(\alpha_kx_{k,0}^2+\beta_kx_{k,0}\xi_k+\gamma_k\xi_k^2\bigr),
 \qquad
 \gamma_k=-\frac1{\varepsilon^2}-\frac{\omega_k^2}{16},                  \tag{51.3}
\]

for suitable \(\alpha_k,\beta_k\).  Given \(x\), the modes are
independent, and \(\operatorname{Cov}(\xi_k,\xi_k^2)=0\) for a centred
Gaussian, so
\(\operatorname{Var}(A\mid x)=\sum_k(\beta_k^2x_{k,0}^2s_k^2+2\gamma_k^2s_k^4)\).
The law of total variance gives
\(\operatorname{Var}A\geq E[\operatorname{Var}(A\mid x)]\geq2\sum_k\gamma_k^2s_k^4\),
which is (51.2).  The expansion uses \(1-r_k^2=2\varepsilon\omega_k+O(\varepsilon^2)\).
Finally, if \(|{\cal D}|<c/\varepsilon\) almost surely then
\(E[{\cal D}^2]\leq c^2/\varepsilon^2\), which contradicts
\(E[{\cal D}^2]\geq\operatorname{Var}[{\cal D}]\geq8\varepsilon^{-2}\sum v_k^2\omega_k^2\)
for small \(\varepsilon\) once \(c^2<8\sum_kv_k^2\omega_k^2\); hence the
last claim.
\end{proof}

\begin{corollary}[No independent homogeneous seed enforces the constraint]
\label{cor:v19v48-no-independent}
For every reflection-symmetric chain whose homogeneous seed is
independent of the TT seeds, the Hamiltonian constraint fails at the
first link on a set of positive probability, with a violation of order
\(\varepsilon^{-1}\).  The V47 tuning removes the mean, not the
fluctuation.
\end{corollary}

\subsection{Dependent homogeneous modes}

The lower bound uses only the independence of \(B\) from \(A\).  If
instead \(\tau\) is a function of the TT link seeds, the two terms can
cancel pathwise.  The balanced choice (40.9),
\(\tfrac23\tau^2=\langle|\sigma|^2\rangle\), removes exactly the
\(\langle|\sigma|^2\rangle\) term of \(A\), leaving
\(-\tfrac14Q(\bar h)-2\Lambda_*\), whose fluctuation is of order one and
whose expectation vanishes at the plane; it does not remove that
order-one term.  The full solver of Theorem \ref{thm:v19v42-solver}
removes everything by construction: it replaces the naive link metric
\(\delta+\bar h\) by \(\psi^4(\delta+\bar h)\) with \(\psi\) chosen
pathwise so that the constraint holds exactly.  Theorem
\ref{thm:v19v48-variance} therefore says that the solver is not a
convenience but a necessity: the constraint surface has measure zero
for every product law, and the seed must be pushed onto it by a
deterministic, seed-dependent map.

\subsection{Exact two-mode value}

For the balanced two-mode chain of V47 (\(\omega_k=1\), \(r_k=\tfrac12\),
\(v_k=2\), \(\varepsilon=\tfrac14\), \(r_0=\tfrac13\),
\(v_0=29997/8192\)), an exact fourth-moment computation gives

\[
 E[{\cal D}]=0,\qquad
 \operatorname{Var}[{\cal D}]=\frac{16632525}{2048}\approx8121,
 \qquad
 \operatorname{Var}[A]\geq\frac{594441}{256}\approx2322,
 \qquad
 \operatorname{Var}[B]=\frac{11108889}{2048}\approx5424,               \tag{51.4}
\]

so the homogeneous tuning that cancels the mean adds more variance than
the TT noise itself.

\subsection{Executable certificate}

The verifier

\[
 \texttt{scripts/verify\_sm\_v19\_v48\_defect\_variance.py}          \tag{51.5}
\]

represents the defect as a polynomial in the plane variables, the two
TT noises, and the homogeneous increment, and evaluates its first and
second moments exactly with uniform-circle and Gaussian moment rules.
It confirms (51.4), the vanishing mean at the balanced variance, and
the inequalities \(\operatorname{Var}[{\cal D}]\geq\) TT lower bound and
\(\geq\operatorname{Var}[B]\).  It emits

\[
 \texttt{artifacts/sm\_v19\_v48\_defect\_variance.json}.             \tag{51.6}
\]

The hashes are

\[
\begin{array}{c|l}
\hbox{object}&\hbox{SHA--256}\\ \hline
\hbox{verifier}&
\texttt{F907DF857FD45741C4B63D1E497586A99A3A9A27DA295108D0DC492C8F83D36E}\\
\hbox{canonical payload}&
\texttt{CC1F208BD1EE2641DC5A7BFF43069869F5C9A8A4C3C39F42DD44979EC254C7FD}.
\end{array}                                                     \tag{51.7}
\]

\begin{firewall}[First link, second order]
\label{fire:v19v48}
The theorem concerns the second-order averaged defect at the first
link.  It does not compute the variance after the solver is applied,
which is zero by construction, nor the variance of the momentum
constraint, nor anything about the continuum limit.  The order-one
fluctuation of \(-\tfrac14Q(\bar h)-2\Lambda_*\) under the balanced
choice is stated, not bounded.
\end{firewall}

\begin{decision}[V49 acquisition]
\label{dec:v19v48-next}
V49 will quantify the solver's action on the free chain at the first
link: it will compute the conformal factor \(\psi_{1/2}=\Psi_\Lambda(\gamma,\sigma,\tau)\)
of Theorem \ref{thm:v19v42-solver} to first order in the defect, namely
\(\psi_{1/2}-\psi_*=(-8\Delta_{\gamma_*}+c)^{-1}(\hbox{defect})+\dots\),
show that the \(\varepsilon^{-1}\) part of the defect is spatially
constant to leading order and is therefore absorbed by the homogeneous
component of \(\psi\) with coefficient \(1/\langle c\rangle=-1/(8\Lambda_*)\),
and determine the size of the resulting random conformal rescaling of
the first-link metric, \(\psi_{1/2}^4-1=O(1/(\varepsilon\Lambda_*))\).
The V50 audit follows.
\end{decision}

\begin{assessment}[V48 acquisition]
\label{ass:v19v48}
The Hamiltonian constraint cannot be enforced in law by an independent
homogeneous seed: the first-link defect has variance of order
\(\varepsilon^{-2}\) from the TT noise alone, variances of independent
parts add, and the V47 tuning even increases the fluctuation.  The
constraint is a pathwise identity, which is exactly what the
seed-dependent solver of V42 provides.
\end{assessment}

\begin{record}[V48 append record]
\label{rec:v19v48}
V48 is appended to the validated V47 whose installed SHA--256 was
\texttt{D0788E29004D2C760E0E621A731E74A9F64012414A557078D50E8F576FB51266}.
After installation only V48 is the active numeric SM note; frozen
archives and unrelated notes are unchanged.
\end{record}


\section{V49: the solver on white link momenta, conformal divergence, and time smearing}
\label{sec:v19v49}

V48 showed that the constraint is a pathwise identity that only the
seed-dependent solver can enforce.  This note computes what the solver
does to the free seed at a link.  Because the link momenta are white at
coincidence, their squared norm is of order \(\varepsilon^{-1}\), and the
integrated Lichnerowicz balance then forces the conformal factor to be
large: the maximal value of \(\psi^{12}\) on the link is at least the
ratio of the mean squared momentum to the effective mean-curvature
coefficient, and its expectation grows like \(\varepsilon^{-1}\).  The
constrained link metrics of the free chain therefore have no limit as
the time spacing vanishes.  The remedy is dictated by V32: constraints
must be imposed on momenta smeared over a fixed physical time, whose
variance stays bounded, and the perturbative regime of the V42 solver
becomes a condition on the smearing time.

\subsection{An integrated lower bound on the conformal factor}

\begin{proposition}[Maximal conformal factor from the integrated balance]
\label{prop:v19v49-bound}
Let \(\psi>0\) solve the \(\Lambda\)-Lichnerowicz equation (45.4) on a
background \(\gamma\) with \(\gamma\)-TT tensor \(\sigma\), constant
\(\tau\), and \(b:=\tfrac23\tau^2-2\Lambda>0\).  Let \(M:=\max\psi\) and
\(\rho:=\max|R(\gamma)|\).  Then

\[
 M^{12}\Bigl(b+\rho\,M^{-4}\Bigr)\geq\langle|\sigma|_\gamma^2\rangle_\gamma,
 \qquad\hbox{hence}\qquad
 M^{12}\geq\frac{\langle|\sigma|^2_\gamma\rangle_\gamma}{b+\rho}
 \ \hbox{ whenever }M\geq1,                                            \tag{52.1}
\]

with \(\langle\cdot\rangle_\gamma\) the \(dV_\gamma\)-average.  For a
constant seed on the flat class the bound is an equality:
\(M^{12}=3|\sigma|^2/(2\tau^2-6\Lambda)\).
\end{proposition}

\begin{proof}
Integrate (45.4) against \(dV_\gamma\); the Laplacian term vanishes,
so \(b\langle\psi^5\rangle=\langle|\sigma|^2\psi^{-7}\rangle-\langle R\psi\rangle\).
Bound the left side by \(bM^5\), the first term on the right from below
by \(\langle|\sigma|^2\rangle M^{-7}\), and the second by \(\rho M\).
Multiply by \(M^7\).  For the constant seed, \(\psi\) is the constant
solution (31.9) with \(b\) replaced by \(\tfrac23\tau^2-2\Lambda\).
\end{proof}

\subsection{Divergence of the constrained link metric of the free chain}

\begin{theorem}[No small-spacing limit for constrained white links]
\label{thm:v19v49-divergence}
Consider the reflection-symmetric chain with Ornstein--Uhlenbeck TT
kernels \(r_k=e^{-\varepsilon\omega_k}\), variances \(v_k\), on a
finite mode space, any plane law with bounded second moments, any
homogeneous seed with \(b\geq-2\Lambda_*>0\), and the link data (45.6).
Let \(\psi_{n+1/2}\) be any positive solution of (45.4) at link
\(n+1/2\) with \(\rho\leq\rho_0\).  Then

\[
 E\Bigl[\max_x\psi_{n+1/2}^{12}\Bigr]
 \geq\frac{1}{b_{\max}+\rho_0}\Bigl(\frac2\varepsilon\sum_kv_k\omega_k-C\Bigr)
 \longrightarrow\infty\qquad(\varepsilon\to0),                          \tag{52.2}
\]

where \(b_{\max}\) bounds \(b\) and \(C\) is independent of
\(\varepsilon\).  In particular the constrained link metrics
\(\psi_{n+1/2}^4\gamma_{n+1/2}\) do not converge, in any sense that
controls the maximal conformal factor, as the time spacing vanishes.
\end{theorem}

\begin{proof}
Apply (52.1) on the event \(M\geq1\), and note \(M^{12}\geq1\) on its
complement, so
\(E[M^{12}]\geq E[\langle|\sigma|^2\rangle]/(b_{\max}+\rho_0)-1\).  By
(36.7), \(E\langle|\sigma|^2\rangle=\sum_k[(1-r_k)^2m_k+v_k(1-r_k^2)]/\varepsilon^2
=\tfrac2\varepsilon\sum_kv_k\omega_k+O(1)\), and the \(\gamma\)-average
differs from the flat one by a factor \(1+O(\|h\|)\).  This gives
(52.2).
\end{proof}

With the V23 normalization the divergent coefficient is
\(\sum_kv_k\omega_k=N_{\rm modes}/(2c_g)\): the zero-point energy of the
retained modes, the same quantity that appeared in the V47 balance,
now acting pathwise.  Its effect is a random global conformal
enlargement of every constrained link, of order
\(\psi^4\gtrsim(\varepsilon\lambda^2)^{-1/3}\) for wave amplitude
\(\lambda\), and the V42 solver leaves its perturbative neighbourhood
as soon as \(\varepsilon^{-1}\sum v_k\omega_k\) exceeds \(|\Lambda_*|\).

\subsection{Time-smeared momenta}

\begin{proposition}[Smeared link momenta have bounded variance]
\label{prop:v19v49-smearing}
For \(T=m\varepsilon\) fixed, define the smeared link momentum
\(\bar\sigma_{n,T}:=(h_{n+m}-h_n)/T\).  For an Ornstein--Uhlenbeck mode,

\[
 \operatorname{Var}\bigl[\bar\sigma_{n,T}\bigr]=\frac{2v\,(1-r^m)}{T^2}
 \longrightarrow\frac{2v(1-e^{-\omega T})}{T^2}\qquad(\varepsilon\to0),
                                                                    \tag{52.3}
\]

bounded uniformly in \(\varepsilon\), and the reflected Gram of smeared
momenta is the Gram (36.8) with \(r\) replaced by \(r^m\) and
\(\varepsilon\) by \(T\), hence positive.  The solver (45.4) applied
with \(\bar\sigma_{n,T}\) in place of \(\sigma_{n+1/2}\) stays in the
perturbative neighbourhood of Theorem \ref{thm:v19v42-solver} whenever

\[
 \frac{2}{T}\sum_kv_k\omega_k\ll|\Lambda_*|,                            \tag{52.4}
\]

that is, when the smearing time exceeds the zero-point energy of the
retained modes divided by the wave energy density.
\end{proposition}

\begin{proof}
\(h_{n+m}-h_n\) has variance \(2v(1-r^m)\) by (36.3), and
\(r^m=e^{-\omega T}\).  The smeared variables are link variables of the
decimated chain with spacing \(T\), which is again a reflection-symmetric
Ornstein--Uhlenbeck chain with ratio \(r^m\) when \(m\) divides the
site index; (36.8) applies.  Condition (52.4) is the requirement that
the expected squared smeared momentum, \(\tfrac2T\sum v_k\omega_k+O(1)\),
be small compared with the plane coefficient \(-8\Lambda_*\) of
(45.5).
\end{proof}

This is the constructive content of Theorem \ref{thm:v19v32-sharp}:
odd fields exist only smeared in time, so the constraints, which are
quadratic in them, can only be imposed on time-smeared data.  The
smearing time is a new parameter of the constrained chain, and (52.4)
shows it cannot be sent to zero at fixed mode number and amplitude.

\subsection{Executable certificate}

The verifier

\[
 \texttt{scripts/verify\_sm\_v19\_v49\_link\_conformal\_divergence.py}
                                                                    \tag{52.5}
\]

checks that the constant seed saturates (52.1), that a negative
\(\Lambda\) lowers the bound, that the unsmeared link variance with
the rational surrogate \(r=1-\varepsilon\omega\) equals
\(2v\omega/\varepsilon-v\omega^2\) and grows through
\(\varepsilon\in\{1/4,\dots,1/32\}\), and that the smeared variance
(52.3) at \(T=1\) stays within \([2.4,2.8]\) for the same spacings.
It emits

\[
 \texttt{artifacts/sm\_v19\_v49\_link\_conformal\_divergence.json}.  \tag{52.6}
\]

The hashes are

\[
\begin{array}{c|l}
\hbox{object}&\hbox{SHA--256}\\ \hline
\hbox{verifier}&
\texttt{26E9A378AD1AB3E2E93E95BB699D155F0428BEA0A18D213CEB29C65D420ED4A1}\\
\hbox{canonical payload}&
\texttt{9D49A9AF55B41B7E73E4C0A75282848CE2A8BD339F1B9C45BECA92793445B795}.
\end{array}                                                     \tag{52.7}
\]

\begin{firewall}[Lower bounds and a prescription]
\label{fire:v19v49}
Theorem \ref{thm:v19v49-divergence} is a lower bound in expectation on
the maximal conformal factor; it does not describe the constrained
metric otherwise.  Proposition \ref{prop:v19v49-smearing} introduces
the smearing time as a prescription with a stated regime; it does not
prove that the smeared constrained chain has a limit as the mode number
grows, nor relate \(T\) to any physical scale.  No continuum measure is
constructed.
\end{firewall}

\begin{decision}[V50 audit and V51 acquisition]
\label{dec:v19v49-next}
V50 is the compulsory companion audit.  V51 will define the smeared
constrained chain precisely: decimate the reflection-symmetric chain to
spacing \(T\), impose the plane law of V46, apply the solver of V42 to
the decimated links under (52.4), and prove reflection positivity,
exact constraints on every decimated link, and the exact reflected
Gram of the decimated seed.  It will then examine whether the mode-number
growth of \(\sum v_k\omega_k\) can be absorbed into \(\Lambda_*\), which
is the lattice form of the cosmological-constant renormalization.
\end{decision}

\begin{assessment}[V49 acquisition]
\label{ass:v19v49}
The solver applied to white link momenta produces conformal factors
whose expected maximum grows like \(\varepsilon^{-1}\), so the naive
constrained chain has no small-spacing limit.  The obstruction is the
zero-point energy of the retained modes acting pathwise, and the
remedy is to impose constraints on momenta smeared over a fixed
physical time, in the perturbative regime (52.4).  This converts the
V32 sharp-time theorem into a concrete rule for the constrained
construction.
\end{assessment}

\begin{record}[V49 append record]
\label{rec:v19v49}
V49 is appended to the validated V48 whose installed SHA--256 was
\texttt{9274615E658451491B31548FA5511A1910E96980B50ED50049E715B8DA64F442}.
After installation only V49 is the active numeric SM note; frozen
archives and unrelated notes are unchanged.
\end{record}


\section{V50: compulsory companion audit and ten-version sweep after the smearing rule}
\label{sec:v19v50}

This is the required audit following V45 and the ten-version sweep
after V40.  Every active companion note in the shared folder was
inspected read-only before the smeared constrained chain announced in
V49 is defined.

\begin{record}[V50 companion snapshot]
\label{rec:v19v50-snapshot}
The inspected files were

\[
\begin{array}{c|r|r}
\hbox{file and SHA--256}&\hbox{bytes}&\hbox{lines}\\ \hline
\texttt{Renewal\_Yang\_Mills\_Mass\_Gap\_Research\_Notes\_V32.tex}
&452376&11540\\
\multicolumn{3}{l}{
\texttt{F1DE1DC4AA157BC4DF069FEDB4F34C3032565655348CB9AE35AC95F5695744D3}}
\\
\texttt{Renewal\_NS\_Stability\_Research\_Notes\_V26.tex}
&278283&5319\\
\multicolumn{3}{l}{
\texttt{82C887F8C2C69E4F28FAD7C3DC8DE5B9099CB5A4BF431EBA1FFF3E91935978AD}}
\\
\texttt{Renewal\_GRH\_Cofinal\_Visibility\_Attack\_V51.tex}
&489091&15170\\
\multicolumn{3}{l}{
\texttt{61927F9477BD72FAD280D5EF4042919932C7687ED895DE4571A71F2D81450134}}
\\
\texttt{Renewal\_YM\_Parallel\_Research\_Notes\_V5.tex}
&54174&1351\\
\multicolumn{3}{l}{
\texttt{306F1BB84CFA8A8D47EA569EC17E9545F9867C8D32B548EC9D9C444B4F3C0512}}.
\end{array}                                                     \tag{53.1}
\]

All four files are byte-identical to the V45 snapshot; no companion
has advanced during V46--V49.
\end{record}

\subsection{Sweep of the ten versions V40--V49}

Since the V40 sweep the programme has: proved sufficiency of the
\(\Lambda<0\) balance and the analytic time-symmetric branch (V41);
described the admissible wave sphere at fixed \(\Lambda\), the local
single-valued solver, and the seeded chain (V42); corrected the
cosmological term in the mean-curvature rate and identified every torus
reflection plane as a turnaround (V43); related the Euclidean throat to
the Lorentzian turnaround and proved the exact spectral consistency of
the Ornstein--Uhlenbeck seed (V44); constrained the plane law by
disintegration (V46); computed the expected first-link defect and the
zero-point balance (V47); proved the defect-variance lower bound
(V48); and derived the conformal divergence of white constrained links
together with the time-smearing rule (V49).

Two internal corrections were made during the sweep, V43 on the sign
of the cosmological rate and V44 on the meaning of the first-step
condition.  Both are recorded in place; neither affects an existence
theorem or a positivity theorem.

\begin{decision}[V50 audit decision]
\label{dec:v19v50-audit}
No companion theorem, numerical constant, or executable artifact is
imported.  The V49 direction stands: V51 defines the smeared
constrained chain and examines the absorption of the zero-point sum
into \(\Lambda_*\).
\end{decision}

\begin{proof}
The companions are unchanged since V45, whose audit already found no
shared operator or hypothesis with the reflection-symmetric chain.  The
V51 task depends only on decimation of a Markov chain, the V42 solver,
and the V47 balance, all internal.
\end{proof}

\begin{assessment}[V50 acquisition]
\label{ass:v19v50}
The audit is current and the ten-version sweep finds the direction
consistent: from the excluded stationary compact target through
toroidal rigidity, the \(\Lambda<0\) branch, the constrained plane, and
the white-momentum obstruction, to the smearing rule.  The programme's
open core is unchanged in kind: a seed with interacting Einstein
dynamics and a limit in mode number and lattice spacing, with the
Standard Model sectors of V6--V16 still uncoupled.
\end{assessment}

\begin{record}[V50 append record]
\label{rec:v19v50}
V50 is appended to the validated V49 whose installed SHA--256 was
\texttt{3B3AC2A3B830D0AA25A265E35E098C733D0B96A9997B8914F95CF811BB3F8770}.
After installation only V50 is the active numeric SM note; the
companions, frozen archives, and all unrelated notes are unchanged.  The
next compulsory companion audit is V55.
\end{record}


\section{V51: the smeared constrained chain and the zero-point dichotomy}
\label{sec:v19v51}

V49 replaced sharp link momenta by momenta smeared over a fixed
physical time.  This note defines the resulting chain exactly: the
reflection-symmetric chain decimated to spacing \(T\), with the
constrained plane law of V46 and the V42 solver on the decimated links.
Reflection positivity, exact constraints on every decimated link, and
the reflected Gram are established.  The zero-point energy of the
retained modes then reappears as a finite but mode-number-divergent sum,
and the note proves a dichotomy: either the homogeneous seed variance
that balances it diverges quadratically in the mode cutoff, or the
constraint is normal ordered, in which case the momentum square loses
its sign and the monotone structure that made the solver single-valued
is lost.  This is the exact lattice form of the cosmological-constant
problem inside the programme.

\subsection{Decimation}

\begin{proposition}[Decimated chain]
\label{prop:v19v51-decimation}
Let \(\mu\) be the reflection-symmetric chain of \((\pi_0,(P_n))\) on
sites \(\varepsilon\mathbb Z\), and let \(m\geq1\), \(T=m\varepsilon\).
The process \((q_{mj})_{j\in\mathbb Z}\) is the reflection-symmetric
chain of \((\pi_0,(P_{mj}P_{mj+1}\cdots P_{mj+m-1})_{j\geq0})\).  For
stationary Ornstein--Uhlenbeck kernels of ratio \(r\) it is the chain
of ratio \(r^m=e^{-\omega T}\), the smeared momenta
\(\bar\sigma_{j+1/2}=(q_{m(j+1)}-q_{mj})/T\) are odd, and for
\(j,j'\geq0\)

\[
 \langle\bar\sigma_{j+1/2},\bar\sigma_{j'+1/2}\rangle_\theta
 =\frac{v(1-r^m)^2r^{m(j+j')}}{T^2}.                                   \tag{54.1}
\]
\end{proposition}

\begin{proof}
The mirror construction commutes with decimation because
\(m\mathbb Z\) is symmetric and contains \(0\); the composed kernels are
the \(m\)-step kernels.  For Ornstein--Uhlenbeck kernels the \(m\)-step
kernel has ratio \(r^m\), and (54.1) is (36.8) with \(r\to r^m\),
\(\varepsilon\to T\).
\end{proof}

\subsection{The smeared constrained chain}

\begin{theorem}[Smeared constrained chain]
\label{thm:v19v51-chain}
Fix \(\Lambda_*\in(-\Lambda_0,0)\), a finite mode space \(V_N\), a
smearing time \(T\), and kernels of bounded support such that the
decimated link data
\((\gamma_{j+1/2},\bar\sigma_{j+1/2},\bar\tau_{j+1/2})\), formed as in
(45.6) from the decimated sites, lie in the neighbourhood \({\cal W}\)
of Theorem \ref{thm:v19v42-solver} for \(|j|<j_0\).  Seed the plane
with the constrained law (49.1).  Then the decimated link data pushed
through \(\Psi_\Lambda\) are exact solutions of the constraints (42.4)
at \(\Lambda_*\) for \(|j|<j_0\), the plane datum is an exact
time-symmetric solution, and the law of the constrained decimated data
is reflection positive on the algebra of links \(0\leq j<j_0\).
\end{theorem}

\begin{proof}
Combine Proposition \ref{prop:v19v51-decimation}, Proposition
\ref{prop:v19v46-conditioning} for the plane law, Theorem
\ref{thm:v19v42-solver} on each decimated link, and Theorem
\ref{thm:v19v32-pushforward} with the locality of decimated links in
the decimated positive-time algebra, exactly as in Proposition
\ref{prop:v19v42-chain}.
\end{proof}

\subsection{The zero-point sum at fixed smearing time}

Let the TT modes be the dual-lattice vectors \(0<|k|\leq K\) of
\(\mathbb Z^3\), with the Einstein normalization
\(v_k\omega_k=1/(2c_g)\) from (26.16).  The expected smeared momentum
square at the first decimated link is, by (36.7) with \(r\to r^m\),

\[
 E\langle|\bar\sigma|^2\rangle
 =\sum_k\frac{(1-e^{-\omega_kT})^2m_k+v_k(1-e^{-2\omega_kT})}{T^2}
 =:{\cal K}_N(T)+Z_N(T),                                              \tag{54.2}
\]

with \({\cal K}_N\) the kink term and \(Z_N\) the zero-point term.

\begin{proposition}[Growth of the zero-point term]
\label{prop:v19v51-growth}
With the Einstein normalization,

\[
 Z_N(T)=\frac1{2c_gT^2}\sum_{0<|k|\leq K}\frac{1-e^{-2|k|T}}{|k|}
 \geq\frac{(1-e^{-2T})\,N(K)}{2c_gT^2K},
 \qquad N(K):=\#\{k\in\mathbb Z^3:0<|k|\leq K\}\geq\frac{K^3}6\ (K\geq2),
                                                                    \tag{54.3}
\]

so \(Z_N(T)\) grows at least like \(K^2\) as the mode cutoff grows,
for every fixed \(T>0\).  The exact mode counts are
\(N(2)=32\), \(N(3)=122\), \(N(4)=256\), \(N(5)=514\), \(N(6)=924\),
all exceeding \(K^3\).
\end{proposition}

\begin{proof}
Each summand is at least \((1-e^{-2T})/K\) for \(|k|\leq K\).  The
ball of radius \(K\) contains the box
\(\{-\lfloor K/\sqrt3\rfloor,\dots,\lfloor K/\sqrt3\rfloor\}^3\) minus
the origin, so \(N(K)\geq(2\lfloor K/\sqrt3\rfloor+1)^3-1\); for
\(K=2,3\) this is \(26\), and for \(K\geq4\) it is at least
\((2K/\sqrt3-1)^3-1\geq K^3/6\).  The listed counts are exact
enumerations.
\end{proof}

\begin{theorem}[Zero-point dichotomy]
\label{thm:v19v51-dichotomy}
For the smeared constrained chain with the Einstein normalization, one
of the following must hold as the mode cutoff \(K\) grows at fixed
\(T\) and \(\Lambda_*\).
\begin{enumerate}
\item The expected smeared defect at the first link is cancelled by the
      homogeneous seed, in which case its variance satisfies
      \(\tfrac43v_0(1-e^{-\omega_0T})/T^2\geq Z_N(T)+O(1)\), hence
      \(v_0\) grows at least like \(K^2\), and by the argument of
      Theorem \ref{thm:v19v49-divergence} the balancing homogeneous
      momentum \(\bar\tau^2\) drives the conformal factor of every
      constrained link to infinity through \(b=\tfrac23\bar\tau^2-2\Lambda_*\).
\item The constraint is normal ordered, \(|\bar\sigma|^2\) being
      replaced by \(|\bar\sigma|^2-E|\bar\sigma|^2\), in which case the
      coefficient \(a\) of \(\psi^{-7}\) in (45.4) takes negative values
      on a set of positive probability, the derivative
      \(7a\psi^{-8}+5b\psi^4\) is no longer positive, and the
      comparison, uniqueness, and barrier arguments of V29, V32, and
      V33 no longer apply.
\end{enumerate}
\end{theorem}

\begin{proof}
The expected defect at the first decimated link is (50.3) with the
decimated parameters; its \(Z_N\) term is deterministic and negative,
and only the homogeneous term \(\tfrac23E[\bar\tau^2]=\tfrac43v_0(1-e^{-\omega_0T})/T^2\)
is positive and adjustable, which gives the inequality in the first
alternative; the growth is Proposition \ref{prop:v19v51-growth}.  A
large \(b\) does not enlarge \(\psi\) by itself, but with
\(\bar\tau^2\) of order \(v_0/T^2\) random and independent of the TT
seeds, Theorem \ref{thm:v19v48-variance} shows the link defect has
variance of order \(v_0^2\), so on a set of positive probability
\(|\bar\sigma|^2\) exceeds \(b\) by a factor growing with \(K\), and
(52.1) makes \(\max\psi^{12}\) grow accordingly.  In the second
alternative, \(|\bar\sigma|^2-E|\bar\sigma|^2\) is negative whenever
the smeared momentum is small, an event of positive probability for a
Gaussian seed; the sign of \(\partial_sF\) in (32.9) is then
indefinite.
\end{proof}

Neither alternative is acceptable as it stands.  The first is the
cosmological-constant problem: the homogeneous sector must absorb a
vacuum energy that diverges with the cutoff.  The second removes the
divergence at the price of the positivity structure that has carried
the constraint solver since V28, and V34 showed what indefinite
coefficients do to uniqueness.  A third route, changing the seed
variances \(v_k\) so that \(\sum_kv_k(1-e^{-2\omega_kT})\) converges,
abandons the Einstein normalization of the linearized sector and hence
the exact spectral consistency of V44.

\subsection{Executable certificate}

The verifier

\[
 \texttt{scripts/verify\_sm\_v19\_v51\_smeared\_chain\_zero\_point.py}
                                                                    \tag{54.4}
\]

checks (54.1) for \(v=2\), \(r=1/2\), \(m=3\), \(\varepsilon=1/4\)
and \(j,j'\leq2\) with all principal minors nonnegative, that
smearing reduces the link variance from \(32\) to \(56/9\), the mode
counts \(N(K)\) for \(K\leq6\) and the monotone growth of
\(N(K)/(2K)\), the positivity of the balancing variance at finite
cutoff, and that the normal-ordered momentum square is negative at zero
momentum.  It emits

\[
 \texttt{artifacts/sm\_v19\_v51\_smeared\_chain\_zero\_point.json}.  \tag{54.5}
\]

The hashes are

\[
\begin{array}{c|l}
\hbox{object}&\hbox{SHA--256}\\ \hline
\hbox{verifier}&
\texttt{AB981389EE79CFA3125E9C13499E78281C1767941BD754876C6F1336A0B1378C}\\
\hbox{canonical payload}&
\texttt{E4917050993F0B43620B06F7B737B4C4D085623BBACE1CD4CB871F366EFF80E7}.
\end{array}                                                     \tag{54.6}
\]

\begin{firewall}[A dichotomy, not a resolution]
\label{fire:v19v51}
Theorem \ref{thm:v19v51-dichotomy} states two unacceptable
alternatives at fixed \(T\) and finite steps; it does not prove that
no other construction exists, does not analyze the normal-ordered
equation beyond the loss of monotonicity, and does not treat the
momentum constraint or the Standard Model sectors.  The smeared
constrained chain of Theorem \ref{thm:v19v51-chain} is exact only for
finitely many links inside the solver's neighbourhood.
\end{firewall}

\begin{decision}[V52 acquisition]
\label{dec:v19v51-next}
V52 will go upstream of the dichotomy to the object that generates it,
the linearized constraint as a quadratic form in the free Fock space.
It will show that the spatially averaged Hamiltonian constraint of the
linearized theory is, in the reconstructed Fock space of V44, the
operator \(\tfrac23\hat\tau^2-\sum_k\omega_k\widehat N_k-\hbox{(zero point)}\)
up to the kink term, identify its zero-point part with
\(Z_N\), and determine whether the physical requirement is the
vanishing of the operator on states or of its normal-ordered form,
which is the standard resolution and fixes which alternative of
Theorem \ref{thm:v19v51-dichotomy} the reconstructed theory selects.
\end{decision}

\begin{assessment}[V51 acquisition]
\label{ass:v19v51}
The smeared constrained chain is exact and reflection positive on
finitely many links, and it exposes the zero-point energy of the
retained modes as a term that grows at least quadratically with the
mode cutoff.  Balancing it with the homogeneous seed inflates every
constrained link; normal ordering it destroys the positivity that
makes the solver single-valued.  The programme has reached the
cosmological-constant problem in exact form and must decide it at the
operator level.
\end{assessment}

\begin{record}[V51 append record]
\label{rec:v19v51}
V51 is appended to the validated V50 whose installed SHA--256 was
\texttt{BE9C0187CCE621189E896587575211AF0069F19AF96368078CB63727D4C2FDE0}.
After installation only V51 is the active numeric SM note; frozen
archives and unrelated notes are unchanged.
\end{record}


\section{V52: the averaged linear constraint as an operator and the induced physical space}
\label{sec:v19v52}

V51 found that imposing the constraint pathwise on the Euclidean seed
leads to a dichotomy between a divergent homogeneous variance and a
loss of positivity.  This note goes upstream to the object that
generates both: the spatially averaged linearized Hamiltonian
constraint, viewed as an operator on the reconstructed Fock space of
V44.  Its transverse-traceless part is the free graviton Hamiltonian
density, whose zero-point part is the divergent sum of V47 and V51; its
homogeneous part is the square of the mean curvature.  Imposing the
constraint on states rather than on paths makes normal ordering
legitimate, because the normal-ordered graviton energy is a positive
operator and no pathwise sign structure is needed.  The induced
physical inner product then identifies the physical states with
graviton number states, each carrying its own mean curvature, and
exhibits the vacuum as a singular double root: the zero-mode problem
of V24 in operator form.

\subsection{The constraint density on one mode}

For a TT mode with amplitude \(x\), Euclidean action
\(\tfrac{c_g}2\int(\dot x^2+\omega^2x^2)dt\), momentum \(p=c_g\dot x\),
and Hamiltonian \(H=p^2/(2c_g)+c_g\omega^2x^2/2\), the linearized
constraint density of (29.11) with \(\sigma=-\dot x/2\) per mode gives

\[
 \langle|\sigma|^2\rangle+\tfrac14Q
 =\frac{p^2}{4c_g^2}+\frac{\omega^2x^2}4
 =\frac{H}{2c_g}.                                                     \tag{55.1}
\]

\begin{proposition}[Averaged linear constraint operator]
\label{prop:v19v52-operator}
In the reconstructed space \({\cal F}\otimes L^2(\mathbb R,d\tau)\) of
V44, with \({\cal F}\) the Fock space of the modes \(k\) and the
homogeneous sector represented by the mean curvature \(\tau\), the
spatially averaged linearized Hamiltonian constraint with cosmological
constant is the operator

\[
 \widehat{\cal C}
 =\tfrac23\hat\tau^2-\frac1{2c_g}\sum_k\omega_k\Bigl(\widehat N_k+\tfrac12\Bigr)-2\Lambda,
 \qquad
 Z:=\frac1{4c_g}\sum_k\omega_k,                                       \tag{55.2}
\]

whose zero-point part \(Z\) is the sharp-time value of the smeared
sum \(Z_N(T)\) of (54.2), and the normal-ordered constraint is
\(\widehat{\cal C}_{\rm no}=\widehat{\cal C}+Z\).
\end{proposition}

\begin{proof}
Sum (55.1) over modes and use \(H_k=\omega_k(\widehat N_k+\tfrac12)\)
from Theorem \ref{thm:v19v44-spectrum}.  The identification of \(Z\)
with the \(T\to0\) limit of the coefficient of the divergent term
follows from (54.2), whose zero-point term is the Euclidean expectation
of \(\langle|\bar\sigma|^2\rangle\) in the reconstructed vacuum, that
is \(\sum_kp_k^2/(4c_g^2)\) smeared; as \(T\to0\) it approaches the
sharp expectation \(\sum_k\omega_k/(4c_g)\).
\end{proof}

\subsection{Normal ordering on states versus on paths}

\begin{proposition}[Positivity is retained at the operator level]
\label{prop:v19v52-positivity}
\(\widehat{\cal C}_{\rm no}+2\Lambda=\tfrac23\hat\tau^2-\tfrac1{2c_g}\sum\omega_k\widehat N_k\)
is a difference of two commuting nonnegative operators, and the
graviton energy \(\sum\omega_k\widehat N_k\) is nonnegative on every
state.  In contrast, the pathwise normal-ordered momentum square
\(|\bar\sigma|^2-E|\bar\sigma|^2\) of Theorem \ref{thm:v19v51-dichotomy}
takes negative values on paths.  The loss of positivity in the second
alternative of that theorem is therefore a property of the pathwise
Euclidean imposition, not of the constraint.
\end{proposition}

\begin{proof}
Both statements are read off: \(\widehat N_k\geq0\) and \(\hat\tau^2\geq0\)
commute; a random variable with mean zero and positive variance is
negative with positive probability.
\end{proof}

This resolves the V51 dichotomy in favour of the second alternative,
provided the constraint is imposed on the reconstructed states.  The
pathwise solver of V42 remains the classical, saddle-point form of the
same condition, which is why it needed a positive \(|\sigma|^2\).

\subsection{The induced physical space}

Let \(|n\rangle\), \(n=(n_k)\), be the Fock number states with
graviton energies \(E_n=\sum_k\omega_kn_k\).  The normal-ordered
constraint acts on \(|n\rangle\otimes\phi(\tau)\) as multiplication by
\(f_n(\tau):=\tfrac23\tau^2-E_n/(2c_g)-2\Lambda\).

\begin{theorem}[Induced inner product and physical states]
\label{thm:v19v52-induced}
Let \(\Lambda=0\).  For every \(n\) with \(E_n>0\), the induced
(group-averaged) inner product
\(\langle\Psi,\delta(\widehat{\cal C}_{\rm no})\Phi\rangle\) restricted to
the sector \(|n\rangle\otimes L^2(d\tau)\) is

\[
 \langle\phi,\delta(f_n)\psi\rangle
 =\sum_{\pm}\frac{\overline{\phi(\pm\tau_n)}\,\psi(\pm\tau_n)}{\tfrac43\tau_n},
 \qquad
 \tau_n=\sqrt{\frac{3E_n}{4c_g}},                                      \tag{55.3}
\]

so the physical space of the excited sectors is
\(\bigoplus_{E_n>0}\mathbb C^2\), one two-dimensional space per number
state, spanned by the graviton state with mean curvature \(\pm\tau_n\):
the quantum form of the V26 balance \(\tfrac23\tau^2=\)~TT energy
density.  For the vacuum \(n=0\), \(f_0(\tau)=\tfrac23\tau^2\) has a
double root at \(\tau=0\) and \(\delta(f_0)\) is not a well-defined
distribution, so the vacuum sector has no induced inner product by
this procedure.  With the un-normal-ordered constraint the vacuum
acquires the root \(\tau_0^2=\tfrac32Z\), growing with the cutoff.
\end{theorem}

\begin{proof}
For \(E_n>0\), \(f_n\) has the two simple roots \(\pm\tau_n\) with
\(|f_n'(\pm\tau_n)|=\tfrac43\tau_n\), and \(\delta(f_n)=\sum_\pm\delta(\tau\mp\tau_n)/|f_n'|\)
by the standard composition rule for the delta distribution with a
function having simple zeros; this gives (55.3), whose Gram matrix on
the two point evaluations is positive definite.  For \(n=0\), \(f_0\)
vanishes to second order at \(0\), and \(\delta(f_0)\) would require
\(\int\phi(\tau)\delta(\tfrac23\tau^2)d\tau\), which diverges for
\(\phi(0)\neq0\).  For the un-normal-ordered operator replace \(E_n\)
by \(E_n+2c_gZ\).
\end{proof}

The vacuum singularity is the operator form of the V24 zero-mode
problem: the flat static torus is the only state with \(\tau=0\), and
its homogeneous sector has no normalizable induced state.  A negative
cosmological constant shifts every root by \(-2\Lambda\) and removes the
double root, \(f_0(\tau)=\tfrac23\tau^2-2\Lambda\) having simple roots
\(\pm\sqrt{3|\Lambda|}\) at the plane; but by V43 those roots
describe a turnaround with \(\tau\neq0\), which the V41 time-symmetric
slice has only through its wave content.  The precise reconciliation,
which requires the second-order constraint operator, is deferred.

\subsection{Executable certificate}

The verifier

\[
 \texttt{scripts/verify\_sm\_v19\_v52\_constraint\_operator.py}      \tag{55.4}
\]

checks (55.1) at three rational points with \(c_g=3/4\), the roots and
weights (55.3) for \(E_n\in\{1,4,9/4\}\), the vanishing derivative at
the vacuum double root, and the growth of the un-normal-ordered
vacuum root with a stand-in cutoff.  It emits

\[
 \texttt{artifacts/sm\_v19\_v52\_constraint\_operator.json}.         \tag{55.5}
\]

The hashes are

\[
\begin{array}{c|l}
\hbox{object}&\hbox{SHA--256}\\ \hline
\hbox{verifier}&
\texttt{5AF3084E43F605A3D2089B090D0A9933862F5EC5B8D6A0FE9FBB35364B31C926}\\
\hbox{canonical payload}&
\texttt{0D86845D38368F16054185E5B87417DF35A838EEE1516D112810757C30807450}.
\end{array}                                                     \tag{55.6}
\]

\begin{firewall}[Averaged, linear, and formal in the vacuum sector]
\label{fire:v19v52}
The constraint treated here is spatially averaged and linearized; the
local constraint, the momentum constraint, and the second-order terms
that couple \(\tau\) to the waves dynamically are not treated.  The
induced inner product is a standard formal device made precise here
only sector by sector.  Nothing is said about interacting states, the
mode-number limit, or the Standard Model sectors.
\end{firewall}

\begin{decision}[V53 acquisition]
\label{dec:v19v52-next}
V53 will reconcile the operator picture with the Euclidean chain: it
will define the physical Euclidean process as the chain of V44
conditioned, at the reconstructed level, on the normal-ordered averaged
constraint, show that this conditioning is equivalent to restricting
the plane law to the graviton-number sectors with their induced
mean-curvature values, and compute the resulting expected mean
curvature and its parity across the plane.  The aim is to replace the
pathwise solver by the state-level constraint for the averaged part
while keeping the pathwise solver for the local, mean-zero part of the
constraint, whose positivity structure is unaffected by normal
ordering.
\end{decision}

\begin{assessment}[V52 acquisition]
\label{ass:v19v52}
The averaged linear constraint is the free graviton energy density
minus the homogeneous term, its zero-point part is exactly the
divergent sum of V47 and V51, and normal ordering is legitimate on
states because the graviton energy is a positive operator.  The
induced physical space pairs every graviton number state with a mean
curvature fixed by the quantum V26 balance, while the vacuum is a
singular double root, the operator form of the V24 zero mode.  The
pathwise dichotomy of V51 is thereby resolved in favour of state-level
imposition of the averaged constraint.
\end{assessment}

\begin{record}[V52 append record]
\label{rec:v19v52}
V52 is appended to the validated V51 whose installed SHA--256 was
\texttt{F50AE38D46A989055D1899400991485E66E61BB23BFF33956E86B11039CE9759}.
After installation only V52 is the active numeric SM note; frozen
archives and unrelated notes are unchanged.
\end{record}


\section{V53: graviton-number plane states as reflection-symmetric chains}
\label{sec:v19v53}

V52 imposed the averaged linear constraint on states and found that the
physical states of the linearized compact theory are graviton number
states, each paired with a mean curvature fixed by the quantum V26
balance.  This note returns those states to the Euclidean side.  A
graviton number state at the plane is exactly a reflection-symmetric
chain whose plane law is the Gaussian multiplied by the squared Hermite
polynomial: no new construction is needed, the chain of V38 with a
different plane law.  Its reflected two-point function carries the Bose
factor \(2n+1\), its constraint expectation is
\((n+\tfrac12)\omega/(2c_g)\), and its mean curvature has vanishing
odd mean across the plane, because the induced physical state is the
symmetric superposition of the expanding and contracting branches.

\subsection{Number states are plane laws}

Fix one TT mode with Ornstein--Uhlenbeck kernel of ratio \(r\) and
stationary variance \(v\), and the reflection-symmetric chain of
Corollary \ref{cor:v19v38-gaussian} with plane law \(\pi_0=N(0,v)\).
Let \(\mathrm{He}_n\) be the probabilists' Hermite polynomials and put

\[
 \pi_0^{(n)}(dx):=\frac{\mathrm{He}_n(x/\sqrt v)^2}{n!}\,N(0,v)(dx).      \tag{56.1}
\]

\begin{theorem}[Number-state chain]
\label{thm:v19v53-chain}
\(\pi_0^{(n)}\) is a probability measure, and the reflection-symmetric
chain \(\mu^{(n)}\) of \((\pi_0^{(n)},P)\) is reflection positive.  Under
the unitary \(L^2(\pi_0^{(n)})\to L^2(\pi_0)\),
\(F\mapsto F\cdot\mathrm{He}_n(\cdot/\sqrt v)/\sqrt{n!}\), the
reconstructed plane vector \(1\) of \(\mu^{(n)}\) is the number state
\(\mathrm{He}_n/\sqrt{n!}\) of Theorem \ref{thm:v19v44-spectrum}, and
the transfer contraction is unchanged.  For \(j,j'\geq0\),

\[
 \langle q_j,q_{j'}\rangle_\theta^{(n)}=E^{(n)}\bigl[q_{-j}q_{j'}\bigr]
 =(2n+1)\,v\,r^{\,j+j'},                                              \tag{56.2}
\]

a rank-one positive matrix with the Bose factor \(2n+1\).
\end{theorem}

\begin{proof}
\(E[\mathrm{He}_n^2]=n!\) under the standard normal gives the
normalization.  Reflection positivity is Proposition
\ref{prop:v19v46-conditioning} with the plane law (56.1).  The
reconstructed space of \(\mu^{(n)}\) is \(L^2(\pi_0^{(n)})\) by Theorem
\ref{thm:v19v38-reconstruction}, and multiplication by
\(\mathrm{He}_n/\sqrt{n!}\) is unitary onto \(L^2(\pi_0)\) by (56.1),
sending \(1\) to the normalized number state; the kernel acts on
functions of \(q_0\) in the same way in both representations.  For
(56.2), \(E^{(n)}[q_{-j}q_{j'}]=E^{(n)}[(P^jq)(q_0)(P^{j'}q)(q_0)]
=r^{j+j'}E^{(n)}[q_0^2]\) by (41.2), and
\(E^{(n)}[q_0^2]=v\,E[x^2\mathrm{He}_n(x)^2]/n!=(2n+1)v\) by the
recursion \(x\mathrm{He}_n=\mathrm{He}_{n+1}+n\mathrm{He}_{n-1}\) and
orthogonality.
\end{proof}

Multi-mode number states \(|n\rangle=\otimes_k|n_k\rangle\) are the
product plane laws \(\otimes_k\pi_0^{(n_k)}\); the chain is the product
chain.

\subsection{Constraint expectation and the induced mean curvature}

\begin{proposition}[Constraint expectation in a number state]
\label{prop:v19v53-constraint}
With the Einstein normalization \(v_k=1/(2c_g\omega_k)\), the TT part of
the averaged linear constraint (55.1) has, in the number-state chain,
plane expectation

\[
 E^{(n)}\Bigl[\tfrac14Q(h_0)\Bigr]+E^{(n)}\bigl[\langle|\sigma|^2\rangle\bigr]_{\rm sharp}
 =\sum_k\Bigl(n_k+\tfrac12\Bigr)\frac{\omega_k}{2c_g},                 \tag{56.3}
\]

where the sharp momentum expectation is the equipartition partner of
the configuration term.  The induced mean curvature of V52 is
\(\tau_n=\sqrt{3E_n/(4c_g)}\) with \(E_n=\sum\omega_kn_k\) after normal
ordering, and \(\sqrt{3(E_n+2c_gZ)/(4c_g)}\) without it.
\end{proposition}

\begin{proof}
\(E^{(n)}[\tfrac14\omega_k^2x_k^2]=\tfrac14\omega_k^2(2n_k+1)v_k
=(2n_k+1)\omega_k/(8c_g)\), and the momentum term contributes the same
by the virial identity of the oscillator, giving (56.3); the mean
curvature values are Theorem \ref{thm:v19v52-induced}.
\end{proof}

\subsection{Parity of the mean curvature across the plane}

The induced physical state of Theorem \ref{thm:v19v52-induced} in the
sector \(n\) is a vector in \(\mathbb C^2\), the two branches
\(\tau=\pm\tau_n\).  A reflection-symmetric Euclidean chain can carry
only the symmetric combination.

\begin{proposition}[No odd mean in a reflection-symmetric state]
\label{prop:v19v53-parity}
Let the homogeneous plane law be even in \(\phi_0\) and the homogeneous
kernel be a Gaussian kernel with zero drift.  Then for every link
\(E[\tau_{j+1/2}]=0\), while \(E[\tau_{j+1/2}^2]>0\).  In particular the
branches \(\pm\tau_n\) enter only through the second moment, and any
state with a nonzero odd mean requires either a drift, as in
Corollary \ref{cor:v19v38-gaussian}, or a plane law that is not even.
\end{proposition}

\begin{proof}
With an even plane law and a driftless symmetric kernel, the joint law
of the homogeneous path is invariant under \(\phi\mapsto-\phi\), which
reverses every link increment, so odd moments vanish; the second moment
is \(2v_0(1-r_0)/\varepsilon^2>0\) as in (50.3).
\end{proof}

This is the Euclidean face of the induced \(\mathbb C^2\): the
Hartle--Hawking-type symmetric superposition of the expanding and
contracting branches is the only one representable by a real,
reflection-symmetric plane law without drift.  Selecting a branch
requires a drift, and V38 showed drifts are compatible with reflection
positivity.

\subsection{Executable certificate}

The verifier

\[
 \texttt{scripts/verify\_sm\_v19\_v53\_number\_state\_chain.py}      \tag{56.4}
\]

computes \(E[\mathrm{He}_n^2]=n!\), \(E^{(n)}[x^2]=2n+1\), the
vanishing odd moment, and the rank-one Gram (56.2) exactly for
\(n\leq5\) with \(r=1/2\), \(v=2\), and checks (56.3) for \(n\leq3\)
with \(c_g=1/4\), \(\omega=2\).  It emits

\[
 \texttt{artifacts/sm\_v19\_v53\_number\_state\_chain.json}.         \tag{56.5}
\]

The hashes are

\[
\begin{array}{c|l}
\hbox{object}&\hbox{SHA--256}\\ \hline
\hbox{verifier}&
\texttt{F9C10688EE3239E817F381872709C9EBEE1109447F3A0BCA8E206FE042A8B348}\\
\hbox{canonical payload}&
\texttt{DF1EA60068BF10E59793152766567E78D8EEC4163FB0E86B6803B7389D9C45D8}.
\end{array}                                                     \tag{56.6}
\]

\begin{firewall}[Free number states only]
\label{fire:v19v53}
The number-state chains are exact but free: their kernels are the
Ornstein--Uhlenbeck kernels, the constraint is the linearized averaged
one, and the induced mean curvature is a label, not a dynamical
variable.  The local constraint, the second-order coupling of
\(\tau\) to the waves, the mode-number limit, and the Standard Model
sectors are untouched.
\end{firewall}

\begin{decision}[V54 acquisition]
\label{dec:v19v53-next}
V54 will combine the state-level averaged constraint of V52--V53 with
the pathwise local constraint: it will split the Hamiltonian constraint
into its spatial average and its mean-zero part, impose the average on
the plane state by the number-state construction, and apply the V42
solver only to the mean-zero part, whose coefficient enters (45.4)
through \(|\sigma|^2-\langle|\sigma|^2\rangle\) with the zero point
removed by the average.  It will determine whether the mean-zero
Lichnerowicz problem retains the positivity needed for a single-valued
solver when \(|\sigma|^2\) is replaced by its mean-zero part plus the
state-level average, which is nonnegative.  The V55 audit follows.
\end{decision}

\begin{assessment}[V53 acquisition]
\label{ass:v19v53}
The physical graviton number states of the linearized compact theory
are realized exactly as reflection-symmetric Euclidean chains with
Hermite-weighted plane laws, with Bose-enhanced reflected correlations
and the expected constraint value \((n+\tfrac12)\omega/(2c_g)\) per
mode.  The induced mean curvature is a two-branch label whose odd mean
cannot appear in a driftless reflection-symmetric chain; the symmetric
superposition is the Euclidean default and a drift selects a branch.
\end{assessment}

\begin{record}[V53 append record]
\label{rec:v19v53}
V53 is appended to the validated V52 whose installed SHA--256 was
\texttt{8EF037ACF279F78DA0AD3B49EE86384202221EDE4B33F7C060E9D691CB56FF07}.
After installation only V53 is the active numeric SM note; frozen
archives and unrelated notes are unchanged.
\end{record}


\section{V54: the split constraint and the incompatibility of normal ordering with pathwise positivity}
\label{sec:v19v54}

V52--V53 imposed the spatially averaged constraint on states and V53's
decision proposed to keep the pathwise solver for the mean-zero part
only.  This note tests that proposal exactly.  Writing the squared
momentum as a state-level normal-ordered average plus a pathwise
mean-zero fluctuation, the coefficient of \(\psi^{-7}\) in the
Lichnerowicz equation becomes \(A_{\rm phys}(n)+a_0(x)\), and the V29
structure requires it to be nonnegative.  For a family of modes with a
single spatial frequency the supremum of the fluctuation equals the
full spatial average, whose expectation exceeds the normal-ordered
value by exactly the vacuum value, so the coefficient is negative
somewhere with positive probability in every graviton number state.
Normal ordering at the state level and positivity of the pathwise
solver are therefore incompatible in general.  The pathwise solver
survives only as the classical and semiclassical device it was in
V41--V42; the quantum constraint must be imposed on states in full.

\subsection{The split coefficient}

At a decimated link with smeared momentum \(\bar\sigma\), write

\[
 |\bar\sigma|^2=\langle|\bar\sigma|^2\rangle+a_0(x),\qquad
 \langle a_0\rangle=0,                                                 \tag{57.1}
\]

and define the state-level normal-ordered average in the number state
\(n\) by

\[
 A_{\rm phys}(n):=E^{(n)}\bigl[\langle|\bar\sigma|^2\rangle\bigr]
 -E^{(0)}\bigl[\langle|\bar\sigma|^2\rangle\bigr]
 =\sum_k\frac{2n_kv_k(1-e^{-\omega_kT})^2}{T^2}\geq0.                  \tag{57.2}
\]

The split Lichnerowicz equation replaces \(|\bar\sigma|^2\) in (45.4) by
\(A_{\rm phys}(n)+a_0(x)\).

\begin{proposition}[Formula (57.2)]
\label{prop:v19v54-aphys}
In the number-state chain of Theorem \ref{thm:v19v53-chain} decimated
to spacing \(T\), \(E^{(n)}\langle|\bar\sigma|^2\rangle
=\sum_k[(1-e^{-\omega_kT})^2(2n_k+1)v_k+v_k(1-e^{-2\omega_kT})]/T^2\),
and subtracting the vacuum value gives (57.2).
\end{proposition}

\begin{proof}
Apply (36.7) with \(r\to e^{-\omega_kT}\), \(\varepsilon\to T\), and
\(E^{(n)}[x_{k,0}^2]=(2n_k+1)v_k\) from (56.2).
\end{proof}

\subsection{Positivity failure for single-frequency families}

\begin{theorem}[The split coefficient is negative with positive probability]
\label{thm:v19v54-failure}
Let the retained modes be \(\cos(k\cdot x)e^A\) for one fixed \(k\) and
both polarizations.  Then \(|\bar\sigma|^2=S\,(1+\cos2k\cdot x)\) with
\(S:=\langle|\bar\sigma|^2\rangle\), so \(a_0=S\cos2k\cdot x\) and
\(\min_x\bigl(A_{\rm phys}(n)+a_0\bigr)=A_{\rm phys}(n)-S\).  In every
number state,

\[
 E^{(n)}[S]-A_{\rm phys}(n)=E^{(0)}\bigl[\langle|\bar\sigma|^2\rangle\bigr]>0,
                                                                    \tag{57.3}
\]

hence \(\Pr^{(n)}[S>A_{\rm phys}(n)]>0\), and on that event the split
coefficient is negative near the zeros of \(1+\cos2k\cdot x\).  The
comparison, uniqueness, and barrier arguments of V29 and V33 do not
apply there, and by V34 uniqueness can genuinely fail.
\end{theorem}

\begin{proof}
For \(\bar\sigma=s_+\cos(k\cdot x)e^++s_\times\cos(k\cdot x)e^\times\),
\(|\bar\sigma|^2=2(s_+^2+s_\times^2)\cos^2(k\cdot x)=(s_+^2+s_\times^2)(1+\cos2k\cdot x)\)
and \(S=s_+^2+s_\times^2\).  Equation (57.3) is the definition (57.2).
A random variable whose mean exceeds a constant exceeds it with positive
probability.
\end{proof}

\begin{corollary}[Vacuum links have no positive coefficient]
\label{cor:v19v54-vacuum}
In the vacuum sector \(A_{\rm phys}(0)=0\), so the split coefficient
is \(a_0\), which is negative on a set of positive spatial measure for
every path with \(\bar\sigma\not\equiv0\).  The split solver is never
single-valued on vacuum links.
\end{corollary}

For families with several spatial frequencies the supremum of
\(a_0\) is smaller relative to its mean, and the failure probability
decreases with the number of frequencies and with the graviton number,
since \(A_{\rm phys}(n)\) grows with \(n\) while the fluctuation scale
of \(a_0\) does not; but no family makes it zero, because \(a_0\) is a
centred variable with positive variance and \(A_{\rm phys}\) is finite.

\subsection{Consequence}

The two ways of handling the zero-point energy identified in V51 and
V52 cannot be combined pathwise: subtracting it at the state level
removes exactly the positive floor that made the pathwise coefficient
nonnegative.  The consistent position is the one V52 reached for the
average and now extends to the whole constraint: impose the quantum
constraint on states, where the graviton energy is a positive operator
and no pathwise sign is needed, and use the pathwise Lichnerowicz
solver only for classical data, as in V41--V42, or for semiclassical
states in which \(A_{\rm phys}(n)\) dominates the fluctuation of
\(a_0\) with high probability.

\subsection{Executable certificate}

The verifier

\[
 \texttt{scripts/verify\_sm\_v19\_v54\_split\_constraint\_positivity.py}
                                                                    \tag{57.4}
\]

checks (57.2) and (57.3) exactly for the two-mode family with the
rational stand-in \(e^{-\omega T}=1/4\), \(v=2\), \(T=1\), \(n\leq3\),
the single-frequency norm identity, and a negative coefficient when
\(S>A_{\rm phys}\).  It emits

\[
 \texttt{artifacts/sm\_v19\_v54\_split\_constraint\_positivity.json}.
                                                                    \tag{57.5}
\]

The hashes are

\[
\begin{array}{c|l}
\hbox{object}&\hbox{SHA--256}\\ \hline
\hbox{verifier}&
\texttt{ADE9F847C1E1B8FE3F269D9AE5E9A6226DB3CA28D509A168A7769349E47363E3}\\
\hbox{canonical payload}&
\texttt{2B91DFE27AF74852B05EDE67368AF5D2C0C8CD90AB899E04D0F9019B4B75C73B}.
\end{array}                                                     \tag{57.6}
\]

\begin{firewall}[A negative result about one construction]
\label{fire:v19v54}
Theorem \ref{thm:v19v54-failure} refutes the split pathwise solver as
a universal rule; it does not refute state-level quantization, does
not bound the failure probability for multi-frequency families, and
does not construct the state-level local constraint.
\end{firewall}

\begin{decision}[V55 audit and V56 acquisition]
\label{dec:v19v54-next}
V55 is the compulsory companion audit.  V56 will begin the
state-level treatment of the local constraint: it will expand the
Hamiltonian constraint to second order about the flat torus in the TT
Fock variables, identify its mean-zero Fourier components at
linear order with the V23 constraint that the TT reduction already
solves, and write the second-order mean-zero components as explicit
quadratic forms in creation and annihilation operators, determining
which of them fail to annihilate the number states and hence which
states of V53 are only first-order physical.
\end{decision}

\begin{assessment}[V54 acquisition]
\label{ass:v19v54}
The proposal to normal order at the state level while solving the
mean-zero constraint pathwise fails: for single-frequency mode families
the split coefficient is negative with positive probability in every
number state, and identically so in the vacuum, so the single-valued
pathwise solver is lost exactly where the zero-point energy is
removed.  The programme therefore moves the entire quantum constraint
to the state level and retains the Lichnerowicz solver for classical
and semiclassical data.
\end{assessment}

\begin{record}[V54 append record]
\label{rec:v19v54}
V54 is appended to the validated V53 whose installed SHA--256 was
\texttt{7C1C61812B5A7306A7568048FFE5DB107712A68441558D89AC94F2803651A899}.
After installation only V54 is the active numeric SM note; frozen
archives and unrelated notes are unchanged.
\end{record}


\section{V55: compulsory companion audit after the state-level decision}
\label{sec:v19v55}

This is the required audit following V50.  Every active companion note
in the shared folder was inspected read-only before the state-level
treatment of the local constraint announced in V54 begins.

\begin{record}[V55 companion snapshot]
\label{rec:v19v55-snapshot}
The inspected files were

\[
\begin{array}{c|r|r}
\hbox{file and SHA--256}&\hbox{bytes}&\hbox{lines}\\ \hline
\texttt{Renewal\_Yang\_Mills\_Mass\_Gap\_Research\_Notes\_V32.tex}
&452376&11540\\
\multicolumn{3}{l}{
\texttt{F1DE1DC4AA157BC4DF069FEDB4F34C3032565655348CB9AE35AC95F5695744D3}}
\\
\texttt{Renewal\_NS\_Stability\_Research\_Notes\_V26.tex}
&278283&5319\\
\multicolumn{3}{l}{
\texttt{82C887F8C2C69E4F28FAD7C3DC8DE5B9099CB5A4BF431EBA1FFF3E91935978AD}}
\\
\texttt{Renewal\_GRH\_Cofinal\_Visibility\_Attack\_V52.tex}
&504761&15514\\
\multicolumn{3}{l}{
\texttt{995305F804828109C77F2E105081AB0B3AE87F103ECAF34F67E1364E06BE92E0}}
\\
\texttt{Renewal\_YM\_Parallel\_Research\_Notes\_V5.tex}
&54174&1351\\
\multicolumn{3}{l}{
\texttt{306F1BB84CFA8A8D47EA569EC17E9545F9867C8D32B548EC9D9C444B4F3C0512}}.
\end{array}                                                     \tag{58.1}
\]

Since V50, GRH has advanced from V51 to V52; the other three files are
unchanged.
\end{record}

\subsection{Companion findings}

GRH V52 proves a general-degree total-height budget for its heat-flowed
polynomials, a ladder-form barrier for critical blocks, and a two-stage
coalescence theorem with two exact no-go witnesses closing its naive
coalescence routes.  Its firewall states that nothing there proves its
target for more real factors than pairs, and that GRH remains open.
Its assumption ledger lists Rolle's theorem, Walsh's coincidence
theorem, and exact rational arithmetic.

Yang--Mills V32, Navier--Stokes V26, and Yang--Mills Parallel V5 are
unchanged since V50.

\begin{decision}[V55 audit decision]
\label{dec:v19v55-audit}
No companion theorem, numerical constant, or executable artifact is
imported.  The V54 direction stands: V56 expands the local Hamiltonian
constraint to second order in the TT Fock variables.
\end{decision}

\begin{proof}
The GRH results concern real-rootedness under a backward heat flow of
polynomials in one complex variable; the SM task concerns quadratic
forms in creation and annihilation operators on a Fock space over TT
modes.  No operator, hypothesis, or constant is shared, and the V56
task is an internal second-order expansion whose linear part is the V23
reduction.
\end{proof}

\begin{assessment}[V55 acquisition]
\label{ass:v19v55}
The audit is current.  Since V50 the programme has defined the smeared
constrained chain, proved the zero-point dichotomy, moved the averaged
constraint to the operator level where normal ordering is legitimate,
realized graviton number states as reflection-symmetric chains, and
shown that a split pathwise solver is incompatible with state-level
normal ordering.  The quantum constraint is now to be imposed on
states; the pathwise solver is retained for classical data.
\end{assessment}

\begin{record}[V55 append record]
\label{rec:v19v55}
V55 is appended to the validated V54 whose installed SHA--256 was
\texttt{0504BD958DF562A18501B5B2FC05CDEFF2DCAA302CB5A71FE55EB7503F594F24}.
After installation only V55 is the active numeric SM note; the
companions, frozen archives, and all unrelated notes are unchanged.  The
next compulsory companion audit is V60.
\end{record}


\section{V56: the second-order constraint of plane-symmetric waves and its operator form}
\label{sec:v19v56}

V54 moved the quantum constraint to the state level.  This note begins
its second-order treatment in the class where everything is exact: TT
waves depending on one coordinate, with both polarizations and any
finite set of frequencies.  The exact scalar curvature of such metrics
is a Riccati expression in the \(z\)-slices' extrinsic curvature, its
second-order part is an explicit quadratic form, and its Fourier
components separate into the spatial average, which is the V52 operator
and restricts states, and the mean-zero components, which are solved
by a second-order scalar response linear in the source and therefore
impose nothing on the states.  This is the quantum form of the V26--V31
split: the constant mode is a compatibility condition, the oscillatory
modes are absorbed by the conformal factor.  The second-order
conformal-factor operator is a finite quadratic form in creation and
annihilation operators and needs no sign condition.

\subsection{Exact curvature of plane-symmetric metrics}

\begin{proposition}[Riccati form]
\label{prop:v19v56-riccati}
For \(g=dz^2+G_{ab}(z)dx^adx^b\) on the torus, with
\(K:=\tfrac12G^{-1}G'\) the mixed extrinsic curvature of the flat
\(z\)-slices,

\[
 R(g)=-2\,(\operatorname{tr}K)'-(\operatorname{tr}K)^2-\operatorname{tr}(K^2).
                                                                    \tag{59.1}
\]

For \(G=I+H\) with \(H\) symmetric traceless, the second-order part is

\[
 R_2=\tfrac12\bigl(\operatorname{tr}H^2\bigr)''-\tfrac14\operatorname{tr}\bigl(H'^2\bigr).
                                                                    \tag{59.2}
\]
\end{proposition}

\begin{proof}
The slices \(z=\hbox{const}\) are flat, so the Gauss--Codazzi
contraction for the unit normal \(\partial_z\) gives (59.1); for
diagonal \(G=\operatorname{diag}(a^2,b^2)\) it reduces to the formula
used in Proposition \ref{prop:v19v41-fourier}.  Expanding
\(G^{-1}=I-H+H^2+\dots\), \(K=\tfrac12H'-\tfrac12HH'+O(3)\), so
\(\operatorname{tr}K=-\tfrac14(\operatorname{tr}H^2)'+O(3)\) since
\(\operatorname{tr}H'=0\), \((\operatorname{tr}K)^2=O(4)\), and
\(\operatorname{tr}K^2=\tfrac14\operatorname{tr}H'^2+O(3)\); substitution
gives (59.2).
\end{proof}

\subsection{Fourier components of the second-order constraint}

Let \(H=\sum_kc_k\cos kz\), \(c_k=x_ke^++y_ke^\times\), with
\(\operatorname{tr}(c_kc_{k'})=2(x_kx_{k'}+y_ky_{k'})\), and let the TT
momentum be \(A=\sum_kd_k\cos kz\), \(d_k=p_ke^++\bar p_ke^\times\).

\begin{theorem}[Second-order constraint components]
\label{thm:v19v56-components}
The second-order Hamiltonian constraint density
\(R_2-|A|^2+\tfrac23\tau^2\) has cosine coefficients: at \(q=0\),

\[
 {\cal H}_0=\tfrac23\tau^2-\sum_k\Bigl[\tfrac14k^2(x_k^2+y_k^2)+(p_k^2+\bar p_k^2)\Bigr],
                                                                    \tag{59.3}
\]

and at \(q\neq0\),

\[
 {\cal H}_q=\sum_{k+k'=q}\Bigl[(x_kx_{k'}+y_ky_{k'})\bigl(-\tfrac12q^2+\tfrac14kk'\bigr)
 -(p_kp_{k'}+\bar p_k\bar p_{k'})\Bigr]
 +\sum_{|k-k'|=q}\Bigl[(x_kx_{k'}+y_ky_{k'})\bigl(-\tfrac12q^2-\tfrac14kk'\bigr)
 -(p_kp_{k'}+\bar p_k\bar p_{k'})\Bigr],                              \tag{59.4}
\]

sums over ordered pairs.  Equation (59.3) is the V26 balance
(29.11) and, in the normalization of (55.1), the V52 operator (55.2)
without cosmological constant.
\end{theorem}

\begin{proof}
Insert the Fourier series into (59.2) using
\(\cos kz\cos k'z=\tfrac12[\cos(k-k')z+\cos(k+k')z]\) and
\(\sin kz\sin k'z=\tfrac12[\cos(k-k')z-\cos(k+k')z]\); the term
\(k=k'\) of \(\operatorname{tr}H'^2\) gives the constant
\(-\tfrac14k^2(x_k^2+y_k^2)\) and the term of \(|A|^2\) gives
\(p_k^2+\bar p_k^2\).  The certificate checks (59.4) against the direct
cosine algebra for two frequencies.
\end{proof}

\subsection{Average restricts states, mean-zero parts are solved}

\begin{theorem}[Second-order scalar response]
\label{thm:v19v56-response}
Let \(g=\psi^4(dz^2+G)\) with \(\psi=1+w\) and \(w\) of second order.
The second-order Hamiltonian constraint is satisfied for all
\(q\neq0\) by the unique mean-zero solution of
\(-8\Delta w=\langle{\cal H}\rangle-{\cal H}\), namely

\[
 w_q=-\frac{{\cal H}_q}{8q^2}\qquad(q\neq0),                          \tag{59.5}
\]

where \({\cal H}_q\) is (59.4) with \(\tau\) omitted, and it is
satisfied at \(q=0\) if and only if (59.3) vanishes.  In the Fock
representation of V44, with \(x_k,p_k\) the mode coordinates and
momenta, (59.5) defines the second-order conformal-factor operator
\(\hat w_q\) as a finite quadratic form in creation and annihilation
operators on every finite mode space; it is bounded on each number
sector and requires no sign condition.  The only state condition at
second order is \(\widehat{\cal H}_0\Psi=0\), the normal-ordered form
of which is Theorem \ref{thm:v19v52-induced}.
\end{theorem}

\begin{proof}
By (44.1), \(R(\psi^4\gamma)=\psi^{-5}(-8\Delta_\gamma\psi+R(\gamma)\psi)\),
and at second order with \(\gamma=\delta+H\), \(-8\Delta_\gamma w=-8\Delta w+O(3)\);
the constraint density becomes \(-8\Delta w+{\cal H}+O(3)\).  Its
mean-zero part is solved by (59.5) because \(-8\Delta\cos qz=8q^2\cos qz\),
and its average is (59.3), which \(w\) cannot affect since
\(\langle\Delta w\rangle=0\).  In the Fock representation, (59.4) is a
finite sum of products of two mode operators, hence a finite quadratic
form; on a number sector all such operators are bounded.
\end{proof}

Thus at second order the state-level quantization is consistent: the
mean-zero constraints define the dependent conformal degree of freedom
as an operator, exactly as V31 defined \(w_0\) classically, and only
the average, the constant cokernel of V26, restricts the physical
states.  The positivity that the pathwise solver needed in V29 and lost
in V54 is not needed here because the second-order response is linear.

\subsection{Executable certificate}

The verifier

\[
 \texttt{scripts/verify\_sm\_v19\_v56\_plane\_wave\_second\_order.py}
                                                                    \tag{59.6}
\]

computes \(R_2\) for the two-frequency wave with amplitudes
\((x_1,y_1)=(1,\tfrac12)\), \((x_2,y_2)=(-\tfrac13,2)\) both by the
direct cosine algebra of (59.2) and by the pair formula (59.4),
obtaining the identical series
\(-\tfrac{637}{144}-\tfrac43\cos z-\tfrac{35}{16}\cos2z-\tfrac{16}3\cos3z-\tfrac{259}9\cos4z\);
checks the average \(-Q/4\) with \(Q=637/36\); recovers the V41
single-mode series; and verifies that (59.5) cancels the mean-zero
part exactly.  It emits

\[
 \texttt{artifacts/sm\_v19\_v56\_plane\_wave\_second\_order.json}.   \tag{59.7}
\]

The hashes are

\[
\begin{array}{c|l}
\hbox{object}&\hbox{SHA--256}\\ \hline
\hbox{verifier}&
\texttt{AEA55BB6C16225D221A52E079071DD2A462D13BAD8327C098459A37672FC6825}\\
\hbox{canonical payload}&
\texttt{85A3CE8C302C395839F3D75C87C7E8570E5A38A434A90253EE75E7FF4A409BAB}.
\end{array}                                                     \tag{59.8}
\]

\begin{firewall}[Plane symmetry, second order, no momentum constraint]
\label{fire:v19v56}
The class is plane symmetric and the expansion stops at second order;
the momentum constraint, the vector sector, the third-order terms that
first couple \(\tau\) dynamically to the waves, operator ordering
ambiguities in (59.4), and the mode-number limit are not treated.
The exact Riccati form (59.1) is classical.
\end{firewall}

\begin{decision}[V57 acquisition]
\label{dec:v19v56-next}
V57 will treat the momentum constraint in the same class: for
plane-symmetric data the momentum constraint reduces to ordinary
differential identities in \(z\), and the note will prove that the
second-order vector response is determined by the TT quadratic sources
exactly as the scalar one, identify the resulting shift-vector operator,
and confirm that no further state condition arises.  It will then
state the complete second-order physical-state condition of the
plane-symmetric compact theory and its induced inner product,
including the cosmological constant.
\end{decision}

\begin{assessment}[V56 acquisition]
\label{ass:v19v56}
In the plane-symmetric class the second-order Hamiltonian constraint
is exact and explicit.  Its spatial average is the V52 operator and
restricts states; its mean-zero components are solved by a linear
scalar response that becomes a bounded quadratic-form operator on
every finite-mode number sector.  State-level quantization is
therefore consistent at second order without any pathwise sign
condition, resolving the tension recorded in V51 and V54.
\end{assessment}

\begin{record}[V56 append record]
\label{rec:v19v56}
V56 is appended to the validated V55 whose installed SHA--256 was
\texttt{75F4395FAD8B4E1954A08A292552D5A748E03D6C0DC97E3FFE30EB3BBD773894}.
After installation only V56 is the active numeric SM note; frozen
archives and unrelated notes are unchanged.
\end{record}


\section{V57: the plane-symmetric momentum constraint, its exact vector solution, and the translation charge}
\label{sec:v19v57}

V56 handled the Hamiltonian constraint of plane-symmetric waves at
second order.  This note completes the second-order constraint system
with the momentum constraint.  For metrics \(dz^2+G(z)\) it reduces
exactly to two ordinary differential identities in \(z\): the
transverse components form a homogeneous linear equation whose full
solution space is a two-parameter family of homogeneous transverse
momenta, and the longitudinal component is a linear equation with a
quadratic source whose periodic solvability is exactly the vanishing of
the \(z\)-translation charge of V27.  For standing waves the charge
vanishes identically and the longitudinal momentum is a bounded
quadratic-form operator; for traveling components the charge is the
graviton momentum operator and restricts the states.  Together with
V52 and V56 this gives the complete second-order physical-state
condition of the plane-symmetric compact theory.

\subsection{Exact reduction}

For \(g=dz^2+G_{ab}(z)dx^adx^b\) the nonzero Christoffel symbols are
\(\Gamma^z_{ab}=-\tfrac12G'_{ab}\) and
\(\Gamma^a_{zb}=\tfrac12(G^{-1}G')^a{}_b\).  Put \(M:=G^{-1}G'\).

\begin{proposition}[Momentum constraint as ordinary differential identities]
\label{prop:v19v57-ode}
For \(z\)-dependent \(K^{ij}\) and \(\tau=\operatorname{tr}K\), the
momentum constraint \(\nabla_j(K^{ij}-\tau g^{ij})=0\) is equivalent to

\[
 \partial_zK^{az}+M^a{}_bK^{bz}+\tfrac12(\operatorname{tr}M)K^{az}=0,
 \qquad
 \partial_zK^{zz}+\tfrac12(\operatorname{tr}M)K^{zz}
 =\tfrac12G'_{ab}K^{ab}+\partial_z\tau.                                \tag{60.1}
\]
\end{proposition}

\begin{proof}
\(\nabla_jK^{ij}=\partial_jK^{ij}+\Gamma^i_{jl}K^{lj}+\Gamma^j_{jl}K^{il}\)
with only \(z\)-derivatives nonzero; insert the Christoffel symbols,
using \(\Gamma^b_{bz}=\tfrac12\operatorname{tr}M\), and
\(\nabla_j(\tau g^{ij})=g^{iz}\partial_z\tau\).
\end{proof}

\begin{theorem}[Exact transverse solution and longitudinal solvability]
\label{thm:v19v57-solution}
\begin{enumerate}
\item The periodic solutions of the transverse equation in (60.1) are
      exactly

\[
 K^{az}=(\det G)^{-1/2}\,(G^{-1})^{ab}C_b,\qquad C\in\mathbb R^2,      \tag{60.2}
\]

      the homogeneous transverse momenta.
\item The longitudinal equation has a periodic solution if and only if

\[
 \int_0^L(\det G)^{1/2}\Bigl(\tfrac12G'_{ab}K^{ab}+\partial_z\tau\Bigr)dz=0,
                                                                    \tag{60.3}
\]

      and then \(K^{zz}\) is determined up to the homogeneous term
      \(c\,(\det G)^{-1/2}\).  For constant \(\tau\), (60.3) is the
      vanishing of the exact \(z\)-translation charge (30.2) with
      \(X=\partial_z\).
\end{enumerate}
\end{theorem}

\begin{proof}
For (60.2), differentiate: \(\partial_z[(\det G)^{-1/2}G^{-1}]
=(\det G)^{-1/2}[-\tfrac12(\operatorname{tr}M)G^{-1}-G^{-1}G'G^{-1}]\),
and \(M^a{}_b(G^{-1}C)^b=(G^{-1}G'G^{-1}C)^a\), so the three terms cancel;
the equation is first order linear with a two-dimensional solution
space, so these are all solutions, and they are periodic because
\(G\) is.  For the longitudinal equation, \((\det G)^{1/2}\) is an
integrating factor:
\(\partial_z[(\det G)^{1/2}K^{zz}]=(\det G)^{1/2}(\tfrac12G'_{ab}K^{ab}+\partial_z\tau)\),
whose periodic solvability is (60.3).  With \(\pi^{ij}=\sqrt g(K^{ij}-\tau g^{ij})\)
and \(({\cal L}_{\partial_z}g)_{ab}=G'_{ab}\), the charge (30.2) is
\(\int\sqrt g\,K^{ab}G'_{ab}-\tau\int\sqrt g\,G^{ab}G'_{ab}\), and the
second integral is \(2\tau\int(\sqrt g)'=0\).
\end{proof}

\subsection{Second order for TT waves}

With \(G=I+H\), \(H=\sum_kc_k\cos kz\), and TT momentum
\(A=\sum_kd_k\cos kz\) in the \(ab\) block, the longitudinal source is
\(\tfrac12\operatorname{tr}(H'A)+O(3)\).

\begin{proposition}[Longitudinal momentum of standing waves]
\label{prop:v19v57-standing}
For cosine modes, \(\operatorname{tr}(H'A)\) is a sine series, its
average vanishes, the charge condition (60.3) holds identically at
second order, and

\[
 K^{zz}_{(2)}(z)=\tfrac12\int^z\operatorname{tr}(H'A)\,dz'
 =\sum_{k,k'}\frac{k\,(x_kp_{k'}+y_k\bar p_{k'})}{2}
 \Bigl[\frac{\cos(k+k')z}{k+k'}+\frac{\cos(k-k')z}{k-k'}\Bigr]+c,
                                                                    \tag{60.4}
\]

where the second fraction is omitted when \(k=k'\).  The transverse second-order momentum
is a homogeneous constant.  In the Fock representation,
\(\widehat K^{zz}_{(2)}\) is a finite quadratic form in the mode
operators, bounded on every number sector, and imposes no state
condition.
\end{proposition}

\begin{proof}
\(\operatorname{tr}(H'A)=-\sum_{k,k'}2k(x_kp_{k'}+y_k\bar p_{k'})\sin kz\cos k'z\),
and \(2\sin kz\cos k'z=\sin(k+k')z+\sin(k-k')z\), a sine series with
zero average; integrate term by term.  The transverse equation has no
second-order source because every term in it contains \(K^{az}\), which
vanishes at first order.
\end{proof}

\begin{proposition}[Traveling components carry charge]
\label{prop:v19v57-traveling}
If sine modes are admitted, \(H=x\cos z\,e^+\), \(A=p\sin z\,e^+\)
give \(\langle\operatorname{tr}(H'A)\rangle=-xp\neq0\), so (60.3)
fails at second order unless the total charge of all components
vanishes.  In the Fock representation with traveling-wave modes the
charge is the graviton momentum operator
\(\widehat P_z=\sum_kk(\widehat N_{k,+}-\widehat N_{k,-})\), whose
kernel is the zero-momentum subspace; this is a second state condition,
with discrete spectrum and no induced-inner-product subtlety.
\end{proposition}

\begin{proof}
\(\operatorname{tr}(H'A)=-2xp\sin^2z\) averages to \(-xp\); this is
the traveling charge of Corollary \ref{cor:v19v27-waves}.  In a basis
of right- and left-moving modes the quadratic charge is diagonal with
eigenvalues \(k(n_{k,+}-n_{k,-})\).
\end{proof}

\subsection{Complete second-order physical-state condition}

\begin{theorem}[Second-order physical states of the plane-symmetric compact theory]
\label{thm:v19v57-physical}
In the plane-symmetric class with a finite mode space, the constraints
through second order, with cosmological constant \(\Lambda\), are
equivalent to: the mean-zero Hamiltonian components define the
conformal-factor operator (59.5); the longitudinal momentum component
defines \(\widehat K^{zz}_{(2)}\) by (60.4); the transverse momenta are
homogeneous constants; and the states satisfy

\[
 \Bigl(\tfrac23\hat\tau^2-\frac1{2c_g}\sum_k\omega_k\widehat N_k-2\Lambda\Bigr)\Psi=0
 \quad\hbox{(normal ordered)},\qquad
 \widehat P_z\Psi=0.                                                   \tag{60.5}
\]

The induced inner product of the first condition is (55.3) shifted by
\(-2\Lambda\), with the vacuum double root removed when \(\Lambda<0\);
the second condition is a projection onto zero total momentum.
\end{theorem}

\begin{proof}
Combine Theorem \ref{thm:v19v56-response}, Proposition
\ref{prop:v19v57-standing}, Proposition \ref{prop:v19v57-traveling},
and Theorem \ref{thm:v19v52-induced} with \(f_n(\tau)=\tfrac23\tau^2-E_n/(2c_g)-2\Lambda\),
whose roots \(\pm\sqrt{3(E_n/(2c_g)+2\Lambda)/2}\) are simple for
\(E_n/(2c_g)+2\Lambda>0\), in particular for the vacuum when
\(\Lambda<0\).
\end{proof}

For \(\Lambda<0\) the vacuum acquires \(\tau_0^2=-3\Lambda\), which is
the mean curvature of the turnaround rate \(\dot\tau=-\Lambda\) of V43
integrated over unit time; the exact matching of this label with the
classical time-symmetric slice of V41, whose \(\tau\) vanishes at the
plane and grows away from it, is a third-order question.

\subsection{Executable certificate}

The verifier

\[
 \texttt{scripts/verify\_sm\_v19\_v57\_plane\_momentum\_constraint.py}
                                                                    \tag{60.6}
\]

checks (60.2) exactly on \(G=\operatorname{diag}((1+z)^2,(1+2z)^2)\) at
\(z=1/2\) with \(C=(1,-3)\), computes the sine series of
\(\operatorname{tr}(H'A)\) and the cosine series (60.4) for two
cosine modes, verifies the derivative identity, and records the
traveling charge \(-3/2\) for \(x=2\), \(p=3/4\).  It emits

\[
 \texttt{artifacts/sm\_v19\_v57\_plane\_momentum\_constraint.json}.  \tag{60.7}
\]

The hashes are

\[
\begin{array}{c|l}
\hbox{object}&\hbox{SHA--256}\\ \hline
\hbox{verifier}&
\texttt{B1F9C80500AE0F152414FE3509B728887B387F276FD7045C90A1D0A976DCF379}\\
\hbox{canonical payload}&
\texttt{7D62917D8C27133D48C680BCAE892E32FAD4243CCDA06804D7DEE62C1B653DE4}.
\end{array}                                                     \tag{60.8}
\]

\begin{firewall}[Plane symmetry and second order]
\label{fire:v19v57}
Everything is confined to the plane-symmetric class and to second
order; the general three-dimensional momentum constraint, third-order
couplings, operator-ordering choices in (60.4), and the mode-number
limit are not treated.  Theorem \ref{thm:v19v57-physical} is a
statement about a finite-mode model of the linearized-plus-quadratic
compact theory, not about full quantum gravity.
\end{firewall}

\begin{decision}[V58 acquisition]
\label{dec:v19v57-next}
V58 will return to the Euclidean side with the complete second-order
conditions: it will construct the reflection-symmetric chain whose
plane law is the product of number-state laws restricted to zero total
momentum and to the induced mean-curvature values, show that the
second-order conformal-factor and longitudinal-momentum operators
(59.5) and (60.4) have Euclidean expectations given by exact
Hermite moments, and compute the resulting second-order corrections to
the reflected two-point functions.  This is the first Euclidean object
in the programme that satisfies all constraints of its order at the
state level.
\end{decision}

\begin{assessment}[V57 acquisition]
\label{ass:v19v57}
The plane-symmetric momentum constraint is exactly reduced and solved:
transverse momenta are a two-parameter homogeneous family, the
longitudinal momentum is fixed by an integrating factor subject to the
translation charge, and at second order the charge is the graviton
momentum operator.  The complete second-order physical-state condition
of the plane-symmetric compact theory is the pair (60.5), with all
other constraint components solved for dependent operators.
\end{assessment}

\begin{record}[V57 append record]
\label{rec:v19v57}
V57 is appended to the validated V56 whose installed SHA--256 was
\texttt{EC5D0144D9D4BDA6564997E23BBC650F81496AFA2F13501484160EA736CA99F3}.
After installation only V57 is the active numeric SM note; frozen
archives and unrelated notes are unchanged.
\end{record}


\section{V58: second-order expectations in graviton number states and the vacuum conformal response}
\label{sec:v19v58}

V56--V57 solved the second-order mean-zero constraints of the
plane-symmetric class for dependent operators.  This note evaluates
those operators in the graviton number states of V53.  The
second-order conformal factor at wavenumber \(2k\) has, in the state
with \(n_k\) gravitons in mode \(k\), an expectation that is affine in
\(n_k\) with a strictly positive vacuum value: the zero-point
fluctuations of every mode curve the conformal factor at twice the
mode's wavenumber.  Normal ordering removes the vacuum value and leaves
the graviton contribution, which for one graviton and the Fourier mode
of V41 is the classical coefficient \(7/128\) times the quantum
amplitude.  The second-order longitudinal momentum has zero expectation
in every number state by parity.  The Euclidean chain reproduces the
configuration parts of these expectations through Hermite moments and
the momentum parts only through the reconstruction, not through sharp
link variances.

\subsection{Diagonal pairs and the conformal operator}

Let the retained modes be \(\cos kz\) in both polarizations, with the
Fock representation of V44: \(\langle n|x_k^2|n\rangle=(2n_k+1)v_k\),
\(\langle n|p_k^2|n\rangle=(2n_k+1)c_g\omega_k/2\), \(v_k=1/(2c_g\omega_k)\),
so that \(\langle x_k^2\rangle\langle p_k^2\rangle=(n_k+\tfrac12)^2\).

\begin{theorem}[Conformal response in a number state]
\label{thm:v19v58-conformal}
For \(q=2k\), the only pair in (59.4) with nonzero number-state
expectation is the diagonal pair \(k=k'\), with coefficient
\(-\tfrac12q^2+\tfrac14k^2=-\tfrac7{16}q^2\).  Hence

\[
 \langle n|\widehat{\cal H}_{2k}|n\rangle
 =-2(2n_k+1)\Bigl[\frac{7}{16}(2k)^2v_k+\frac{c_g\omega_k}2\Bigr],
 \qquad
 \langle n|\hat w_{2k}|n\rangle
 =\frac{(2n_k+1)}{16k^2}\Bigl[\frac74k^2v_k+\frac{c_g\omega_k}2\Bigr]>0,
                                                                    \tag{61.1}
\]

the factor \(2\) counting the two polarizations, and
\(\langle n|\hat w_q|n\rangle=0\) for \(q\) not twice a retained
wavenumber.  In particular the vacuum has a nonzero second-order
conformal response at every doubled wavenumber, and the normal-ordered
operator \(:\hat w_{2k}:=\hat w_{2k}-\langle0|\hat w_{2k}|0\rangle\) has
expectation proportional to \(n_k\).
\end{theorem}

\begin{proof}
Off-diagonal pairs \(k\neq k'\) contribute products \(x_kx_{k'}\) or
\(p_kp_{k'}\) of independent modes with zero mean.  For \(k=k'\),
(59.4) gives \((x_k^2+y_k^2)(-\tfrac12(2k)^2+\tfrac14k^2)-(p_k^2+\bar p_k^2)\),
and the coefficient is \(-2k^2+\tfrac14k^2=-\tfrac74k^2=-\tfrac7{16}(2k)^2\).
Insert the oscillator expectations and (59.5).
\end{proof}

\begin{corollary}[Classical limit and one graviton]
\label{cor:v19v58-classical}
For a classical standing wave \(x_k=\lambda\), \(p_k=0\), one
polarization, (59.4)--(59.5) give \(w_{2k}=\tfrac7{128}\lambda^2\) for
every \(k\), which for \(k=1\) is the V41 coefficient \(7/128\).  In the
normal-ordered quantum expression the one-graviton state has
\(\langle1|:\hat w_{2k}:|1\rangle=\tfrac1{8k^2}[\tfrac74k^2v_k+\tfrac{c_g\omega_k}2]\),
the classical value with \(\lambda^2\) replaced by the quantum
amplitude \(2v_k\) plus the momentum share.
\end{corollary}

\begin{proof}
Set \(n_k=1\) in (61.1) and subtract the vacuum value; the classical
statement is the diagonal term of (59.4) with \(p=0\).
\end{proof}

\subsection{Longitudinal momentum}

\begin{proposition}[Vanishing expectation]
\label{prop:v19v58-longitudinal}
\(\langle n|\widehat K^{zz}_{(2)}|n\rangle=0\) for every number state,
with any symmetric ordering of the products \(x_kp_{k'}\).
\end{proposition}

\begin{proof}
(60.4) is bilinear in \(x\) and \(p\); for \(k\neq k'\) the factors are
independent with zero mean, and for \(k=k'\) the symmetrized product
\(x_kp_k+p_kx_k\) has zero expectation in every number state.
\end{proof}

\subsection{What the Euclidean chain reproduces}

\begin{proposition}[Euclidean side]
\label{prop:v19v58-euclidean}
In the number-state chain of Theorem \ref{thm:v19v53-chain}, the plane
expectation of every configuration monomial \(x_kx_{k'}\) equals its
Fock expectation, by the Hermite moment
\(E^{(n)}[x_k^2]=(2n_k+1)v_k\).  The momentum monomials \(p_kp_{k'}\)
are not plane observables of the chain: the sharp link variance
\(E[\sigma_{1/2}^2]\) is of order \(\varepsilon^{-1}\) by (36.7), while
the Fock value \((2n_k+1)c_g\omega_k/2\) is recovered only through the
reconstruction of Theorem \ref{thm:v19v38-reconstruction}, where the
momentum is the generator conjugate to \(x_k\) in \(L^2(\pi_0^{(n)})\).
Consequently the expectation (61.1) is a statement about the
reconstructed state, and its Euclidean evaluation requires smeared
momenta with the reflected form (35.10), not link variances.
\end{proposition}

\begin{proof}
The first statement is (56.2) at \(j=j'=0\).  The second is Theorem
\ref{thm:v19v32-sharp} together with Corollary \ref{cor:v19v33-limit}:
the sharp link momentum has no finite coincident variance, whereas the
reconstructed momentum operator is bounded on each number sector.
\end{proof}

\subsection{Executable certificate}

The verifier

\[
 \texttt{scripts/verify\_sm\_v19\_v58\_second\_order\_expectations.py}
                                                                    \tag{61.2}
\]

computes (61.1) exactly for \(c_g=3/4\), \(k\in\{1,2\}\), \(n\leq3\),
including the uncertainty product \((n+\tfrac12)^2\), the coefficient
\(-\tfrac7{16}q^2\), the positive vacuum value, and the normal-ordered
values; it reproduces the V41 classical coefficient \(7/128\) and
records the vanishing longitudinal expectation.  For \(k=1\), the
vacuum response is \(37/384\) and the one-graviton normal-ordered
response is \(37/192\).  It emits

\[
 \texttt{artifacts/sm\_v19\_v58\_second\_order\_expectations.json}.  \tag{61.3}
\]

The hashes are

\[
\begin{array}{c|l}
\hbox{object}&\hbox{SHA--256}\\ \hline
\hbox{verifier}&
\texttt{0ACFB92235B3906054E35977B1BD930CCAA458A2F0FCAF40566E387187B18EAA}\\
\hbox{canonical payload}&
\texttt{848CED3F6AC41B2DF8CE003EC0EDE2DE6C2365959F6D1AC15BBB29EC7A33C962}.
\end{array}                                                     \tag{61.4}
\]

\begin{firewall}[Expectations of dependent operators]
\label{fire:v19v58}
The results are expectations of second-order dependent operators in
free number states of a finite-mode plane-symmetric model; they are
not a construction of interacting states, do not treat ordering of
the \(x\)-\(p\) products beyond symmetrization, and do not take a
mode-number limit, where the vacuum response summed over modes is a
further zero-point quantity.
\end{firewall}

\begin{decision}[V59 acquisition]
\label{dec:v19v58-next}
V59 will sum the vacuum conformal response over modes: it will show
that \(\sum_k\langle0|\hat w_{2k}|0\rangle\cos2kz\) is a divergent
Fourier series as the mode cutoff grows, identify its divergence with
the zero-point structure of V47 and V52, and determine the exact effect
of normal ordering the constraint on the conformal operator, namely
whether \(:\hat w_q:\) coincides with the response to the normal-ordered
constraint.  The V60 audit follows.
\end{decision}

\begin{assessment}[V58 acquisition]
\label{ass:v19v58}
The second-order dependent operators of the plane-symmetric compact
theory have exact expectations in graviton number states: the conformal
factor at doubled wavenumber is affine in the occupation number with a
positive zero-point value, the longitudinal momentum averages to zero,
and the classical V41 coefficient is recovered as the one-graviton
normal-ordered response.  The Euclidean chain reproduces configuration
moments exactly and momentum moments only through reconstruction.
\end{assessment}

\begin{record}[V58 append record]
\label{rec:v19v58}
V58 is appended to the validated V57 whose installed SHA--256 was
\texttt{C1AC9FA08C82DCECA4F8F627E3E7BD1FB7585BE2F6C01F2E5529E9E304CC73E0}.
After installation only V58 is the active numeric SM note; frozen
archives and unrelated notes are unchanged.
\end{record}


\section{V59: the summed vacuum response, translation invariance, and normal ordering}
\label{sec:v19v59}

V58 found a positive vacuum conformal response at every doubled
wavenumber in the standing-wave model.  This note sums it over modes
and finds a logarithmic caustic, then shows that the response is an
artifact of the standing-wave truncation: in the translation-invariant
model, with sine and cosine modes of equal weight, every mean-zero
second-order operator has zero expectation in the vacuum and in every
momentum eigenstate, so the entire zero-point content of the
second-order constraint sits in the spatial average, exactly where V52
placed it.  Mean-zero responses appear only in states that break
translation symmetry, such as standing-wave number states and the
coherent states that describe the classical branch of V41.  Normal
ordering commutes with the linear response.

\subsection{The standing-wave vacuum response}

\begin{proposition}[Logarithmic caustic]
\label{prop:v19v59-log}
With the Einstein normalization, (61.1) gives
\(\langle0|\hat w_{2k}|0\rangle=C_0/k\) with
\(C_0=\tfrac1{16}\bigl(\tfrac7{8c_g}+\tfrac{c_g}2\bigr)\), independent of
the polarization count beyond the factor already included.  The mode
sum of the vacuum second-order conformal factor over all cosine modes
\(k\leq K\) converges as \(K\to\infty\), in \(L^2\) and pointwise for
\(z\notin\pi\mathbb Z\), to

\[
 w_{\rm vac}(z)=C_0\sum_{k\geq1}\frac{\cos2kz}k=-C_0\log|2\sin z|,
                                                                    \tag{62.1}
\]

which is unbounded at \(z\in\pi\mathbb Z\).
\end{proposition}

\begin{proof}
Insert \(v_k=1/(2c_gk)\), \(\omega_k=k\) into (61.1).  The series
\(\sum\cos(2kz)/k=-\log|2\sin z|\) is the classical Fourier expansion,
convergent for \(z\notin\pi\mathbb Z\) and in \(L^2\) since
\(\sum k^{-2}<\infty\); at \(z=\pi/2\) it is
\(\sum(-1)^k/k=-\log2\), which the certificate brackets by alternating
partial sums.
\end{proof}

The caustic is where all standing modes have a common antinode; it is
a property of the cosine truncation, whose vacuum is not translation
invariant.

\subsection{Translation invariance removes mean-zero vacuum responses}

Let the mode space contain, for each \(k\), both \(\cos kz\) and
\(\sin kz\) in each polarization, with equal variances, so that
\(z\)-translations act unitarily on the Fock space and the vacuum is
invariant.

\begin{theorem}[Momentum eigenstates have no mean-zero response]
\label{thm:v19v59-invariance}
Let \(\hat O_q\) be any operator with definite wavenumber \(q\neq0\)
under \(z\)-translations, \(U_a\hat O_qU_a^{-1}=e^{iqa}\hat O_q\), for
instance \(\hat w_q\), \(\widehat{\cal H}_q\), or \(\widehat K^{zz}_{(2),q}\).
Then \(\langle\Psi|\hat O_q|\Psi\rangle=0\) for every eigenstate
\(\Psi\) of \(U_a\), in particular for the vacuum and for every
momentum eigenstate.  Hence in the translation-invariant model the
vacuum second-order conformal factor is spatially constant, and the
zero-point content of the second-order constraint is entirely in
\(\widehat{\cal H}_0\).
\end{theorem}

\begin{proof}
\(\langle\Psi|\hat O_q|\Psi\rangle=\langle U_a\Psi|U_a\hat O_qU_a^{-1}|U_a\Psi\rangle
=e^{iqa}\langle\Psi|\hat O_q|\Psi\rangle\) for every \(a\), which forces
zero for \(q\neq0\).  For a single mode pair,
\((x_c\cos kz+x_s\sin kz)^2=\tfrac12(x_c^2+x_s^2)+\tfrac12(x_c^2-x_s^2)\cos2kz+x_cx_s\sin2kz\),
whose mean-zero coefficients have zero vacuum expectation when the
variances are equal; this is the explicit form of the general
argument.
\end{proof}

\begin{corollary}[Where the responses live]
\label{cor:v19v59-where}
Mean-zero second-order responses are nonzero only in states that break
translation symmetry: standing-wave number states of the truncated
model, and coherent states \(|\alpha\rangle\) with
\(\langle\alpha|x_k|\alpha\rangle=\lambda\), for which
\(\langle\alpha|\hat w_{2k}|\alpha\rangle\) equals the classical V41
value \(\tfrac7{128}\lambda^2\) plus the translation-invariant vacuum
value zero.  The classical branch of V41 is thus the coherent-state
sector of the quantum model, not its number-state sector.
\end{corollary}

\begin{proof}
For a coherent state the expectation of a quadratic form is the
classical value at the coherent amplitude plus the vacuum expectation,
and the latter vanishes by Theorem \ref{thm:v19v59-invariance}.
\end{proof}

\subsection{Normal ordering and the response}

\begin{proposition}[Normal ordering commutes with the linear response]
\label{prop:v19v59-no}
Since \(\hat w_q=-\widehat{\cal H}_q/(8q^2)\) is linear in the constraint
operator, \(:\hat w_q:=-:\widehat{\cal H}_q:/(8q^2)\); normal ordering
the constraint and normal ordering its response are the same
operation, and in the translation-invariant model both are
unnecessary for \(q\neq0\) because the vacuum expectation already
vanishes.  Normal ordering is needed only for \(\widehat{\cal H}_0\).
\end{proposition}

\begin{proof}
Normal ordering is linear, and (59.5) is linear.
\end{proof}

This completes the operator-level account of the zero-point energy in
the plane-symmetric model: it is a c-number in \(\widehat{\cal H}_0\), to
be absorbed by normal ordering or, equivalently, by the cosmological
constant, and nowhere else.

\subsection{Executable certificate}

The verifier

\[
 \texttt{scripts/verify\_sm\_v19\_v59\_vacuum\_response\_sum.py}     \tag{62.2}
\]

checks \(\langle0|\hat w_{2k}|0\rangle=C_0/k\) for \(k\leq3\) with
\(c_g=3/4\), the alternation and contraction of the partial sums of
\(\sum(-1)^k/k\) through twelve terms with \(-\log2\) bracketed, the
single-mode-pair identity at two points, and the vanishing mean-zero
vacuum response for equal variances.  It emits

\[
 \texttt{artifacts/sm\_v19\_v59\_vacuum\_response\_sum.json}.        \tag{62.3}
\]

The hashes are

\[
\begin{array}{c|l}
\hbox{object}&\hbox{SHA--256}\\ \hline
\hbox{verifier}&
\texttt{8E95A7BF1E8D158C1366A474E25A846FD4216A901D48F4CD1B1637AF2DD004E7}\\
\hbox{canonical payload}&
\texttt{2A6E3730DAC099A005875FB4019B3A2128B988CB8BABF2E74261A1B6BECB3980}.
\end{array}                                                     \tag{62.4}
\]

\begin{firewall}[Second order, plane symmetric, finite modes]
\label{fire:v19v59}
The statements are exact within the plane-symmetric second-order
model.  The logarithmic series is a property of the cosine truncation;
the translation-invariant statement is general but says nothing about
third order, about the mode-number behaviour of \(\widehat{\cal H}_0\),
which remains the divergent zero-point sum, or about interacting
states.
\end{firewall}

\begin{decision}[V60 audit and V61 acquisition]
\label{dec:v19v59-next}
V60 is the compulsory companion audit and ten-version sweep.  V61 will
address the coherent-state sector identified in Corollary
\ref{cor:v19v59-where}: it will define the Euclidean reflection-symmetric
chain whose reconstructed plane state is a coherent state, namely the
Gaussian chain with a nonzero reflection-symmetric mean path, show that
its second-order responses reproduce the classical V41 branch plus the
translation-invariant vacuum, and compute its induced mean-curvature
distribution from the averaged constraint, which for a coherent state
is a superposition over graviton numbers and hence over mean-curvature
labels.
\end{decision}

\begin{assessment}[V59 acquisition]
\label{ass:v19v59}
The vacuum conformal response of V58 sums to a logarithmic caustic in
the standing-wave truncation and vanishes identically in the
translation-invariant model, where every mean-zero second-order
response is zero in momentum eigenstates.  The zero-point energy is
confined to the averaged constraint, normal ordering acts there alone,
and the classical branch of V41 is identified as the coherent-state
sector of the quantum model.
\end{assessment}

\begin{record}[V59 append record]
\label{rec:v19v59}
V59 is appended to the validated V58 whose installed SHA--256 was
\texttt{0CBC796312B66EAB186305756B6638A9F8C1BF9A27A07DEE0E0948952CA8C5F5}.
After installation only V59 is the active numeric SM note; frozen
archives and unrelated notes are unchanged.
\end{record}


\section{V60: compulsory companion audit and ten-version sweep after the operator-level constraint}
\label{sec:v19v60}

This is the required audit following V55 and the ten-version sweep
after V50.  Every active companion note in the shared folder was
inspected read-only before the coherent-state chain announced in V59
is constructed.

\begin{record}[V60 companion snapshot]
\label{rec:v19v60-snapshot}
The inspected files were

\[
\begin{array}{c|r|r}
\hbox{file and SHA--256}&\hbox{bytes}&\hbox{lines}\\ \hline
\texttt{Renewal\_Yang\_Mills\_Mass\_Gap\_Research\_Notes\_V32.tex}
&452376&11540\\
\multicolumn{3}{l}{
\texttt{F1DE1DC4AA157BC4DF069FEDB4F34C3032565655348CB9AE35AC95F5695744D3}}
\\
\texttt{Renewal\_NS\_Stability\_Research\_Notes\_V26.tex}
&278283&5319\\
\multicolumn{3}{l}{
\texttt{82C887F8C2C69E4F28FAD7C3DC8DE5B9099CB5A4BF431EBA1FFF3E91935978AD}}
\\
\texttt{Renewal\_GRH\_Cofinal\_Visibility\_Attack\_V52.tex}
&504761&15514\\
\multicolumn{3}{l}{
\texttt{995305F804828109C77F2E105081AB0B3AE87F103ECAF34F67E1364E06BE92E0}}
\\
\texttt{Renewal\_YM\_Parallel\_Research\_Notes\_V5.tex}
&54174&1351\\
\multicolumn{3}{l}{
\texttt{306F1BB84CFA8A8D47EA569EC17E9545F9867C8D32B548EC9D9C444B4F3C0512}}.
\end{array}                                                     \tag{63.1}
\]

All four files are byte-identical to the V55 snapshot, and no new
note series has appeared in the folder.
\end{record}

\subsection{Sweep of the ten versions V50--V59}

Since the V50 sweep the programme has: defined the smeared constrained
chain and proved the zero-point dichotomy (V51); moved the averaged
constraint to the operator level, where normal ordering is legitimate,
and derived the induced physical space with its vacuum double root
(V52); realized graviton number states as reflection-symmetric chains
(V53); shown that a split pathwise solver is incompatible with
state-level normal ordering (V54); reduced the plane-symmetric
Hamiltonian constraint at second order to explicit Fourier components
whose mean-zero parts are solved by a linear scalar response (V56);
done the same for the momentum constraint, with the translation charge
as the only further state condition (V57); evaluated the dependent
operators in number states (V58); and shown that mean-zero vacuum
responses vanish in the translation-invariant model, confining the
zero-point energy to the averaged constraint (V59).

The direction has shifted once during the sweep, at V54, from
pathwise to state-level imposition of the quantum constraint; the
pathwise solver is retained for classical and coherent-state data.
No erratum was needed in V50--V59.

\begin{decision}[V60 audit decision]
\label{dec:v19v60-audit}
No companion theorem, numerical constant, or executable artifact is
imported.  The V59 direction stands: V61 constructs the coherent-state
chain and its induced mean-curvature distribution.
\end{decision}

\begin{proof}
The companions are unchanged since V55, whose audit found no shared
operator or hypothesis with the reflection-symmetric chain and its
Fock reconstruction.  The V61 task uses only the shifted Gaussian plane
law, Theorem \ref{thm:v19v38-chain}, and the induced inner product of
Theorem \ref{thm:v19v52-induced}.
\end{proof}

\begin{assessment}[V60 acquisition]
\label{ass:v19v60}
The audit is current and the sweep finds the direction coherent: the
compact \(\Lambda<0\) programme now has a state-level second-order
constraint theory in the plane-symmetric class, with the zero-point
energy isolated in one c-number, the classical branch identified as
the coherent-state sector, and a Euclidean realization of every number
state.  The open core is unchanged in kind: third order and beyond,
the general three-dimensional class, the mode-number limit, and the
coupling of the Standard Model sectors.
\end{assessment}

\begin{record}[V60 append record]
\label{rec:v19v60}
V60 is appended to the validated V59 whose installed SHA--256 was
\texttt{D76C80D009D97AA57D887EB8EDBF302DC7222080A30280A4CC7441071DD2C6D9}.
After installation only V60 is the active numeric SM note; the
companions, frozen archives, and all unrelated notes are unchanged.  The
next compulsory companion audit is V65.
\end{record}


\section{V61: the coherent-state chain and the singular sector of the classical branch}
\label{sec:v19v61}

V59 identified the classical branch of V41 with the coherent-state
sector.  This note constructs the Euclidean coherent-state chain: a
reflection-symmetric Ornstein--Uhlenbeck chain whose plane law is a
shifted Gaussian.  Its reconstructed plane vector is the coherent state
with the classical amplitude, its reflected correlations are those of
the vacuum plus the classical mean, and its second-order responses are
the classical V41 values plus the vacuum values of V58--V59.  The
averaged constraint then distributes the coherent state over graviton
number sectors with Poisson weights, and each sector carries its own
induced mean curvature.  When the cosmological constant is tuned to the
classical balance of V41, the sector at the mean occupation number is
exactly the double root of V52, and all sectors below it have no real
root at all: the classical time-symmetric slice sits on the singular
locus of the induced inner product, and a Poisson fraction of the
coherent state is unphysical.

\subsection{Shifted plane laws are coherent states}

Fix one mode with ratio \(r\), vacuum variance \(v=1/(2c_g\omega)\), and
let \(\pi_0^{(\lambda)}:=N(\lambda,v)\).

\begin{theorem}[Coherent-state chain]
\label{thm:v19v61-chain}
The reflection-symmetric chain of \((\pi_0^{(\lambda)},P)\) is
reflection positive; its reconstructed plane vector corresponds, under
the unitary \(L^2(N(\lambda,v))\to L^2(N(0,v))\),
\(F\mapsto F\,e^{\lambda x/v-\lambda^2/(2v)}\), to the normalized
coherent state with \(\langle x\rangle=\lambda\), that is
\(|\alpha\rangle\) with \(|\alpha|^2=\lambda^2c_g\omega/2\); the
transfer contraction is unchanged; and for \(j,j'\geq0\)

\[
 E^{(\lambda)}[q_{-j}q_{j'}]=r^{\,j+j'}\bigl(v+\lambda^2\bigr),
 \qquad
 E^{(\lambda)}[q_n]=\lambda r^{|n|}.                                    \tag{64.1}
\]
\end{theorem}

\begin{proof}
Reflection positivity is Proposition \ref{prop:v19v46-conditioning}.
The density of \(N(\lambda,v)\) relative to \(N(0,v)\) is
\(e^{\lambda x/v-\lambda^2/(2v)}\), whose square root is the stated
multiplier, and this function is the exponential vector of the Fock
representation, the coherent state with mean \(\lambda\); the
identification of \(|\alpha|^2\) uses \(x=(a+a^\dagger)/\sqrt{2c_g\omega}\).
The moments follow from (41.2) with \(E^{(\lambda)}[q_0^2]=v+\lambda^2\)
and the decaying conditional mean of Proposition
\ref{prop:v19v44-mean}.
\end{proof}

The mean path \(\lambda r^{|n|}\) is the Euclidean decaying solution
with a cusp at the plane, as in V44; the classical Lorentzian
\(\lambda\cos\omega t\) is its analytic continuation.

\subsection{Second-order responses}

\begin{proposition}[Coherent responses]
\label{prop:v19v61-response}
In the coherent state of Theorem \ref{thm:v19v61-chain},
\(\langle\alpha|\hat w_{2k}|\alpha\rangle=\tfrac7{128}\lambda^2+\langle0|\hat w_{2k}|0\rangle\),
which is \(\tfrac7{128}\lambda^2+C_0/k\) in the cosine model and
\(\tfrac7{128}\lambda^2\) in the translation-invariant model, and
\(\langle\alpha|\widehat K^{zz}_{(2)}|\alpha\rangle=0\).
\end{proposition}

\begin{proof}
A quadratic form in a coherent state equals its classical value at the
mean plus its vacuum expectation; apply Corollary
\ref{cor:v19v58-classical}, Proposition \ref{prop:v19v59-log}, Theorem
\ref{thm:v19v59-invariance}, and Proposition
\ref{prop:v19v58-longitudinal} with \(\langle p\rangle=0\) at the plane.
\end{proof}

\subsection{Induced mean-curvature sectors of the coherent state}

The coherent state has Poisson occupation weights
\(e^{-|\alpha|^2}|\alpha|^{2n}/n!\) and normal-ordered energy
\(E_{\rm cl}=\omega|\alpha|^2=c_g\omega^2\lambda^2/2\), the classical
energy of the standing wave.  The classical balance of V41 at the
time-symmetric plane is \(2\Lambda=-E_{\rm cl}/(2c_g)=-\omega^2\lambda^2/4\),
that is \(\Lambda=-\omega^2\lambda^2/8=-\tfrac18\langle|\nabla h|^2\rangle\).

\begin{theorem}[Sectors of the coherent state at the classical balance]
\label{thm:v19v61-sectors}
With \(\Lambda\) fixed at the classical balance, the induced inner
product of Theorem \ref{thm:v19v52-induced} in the sector with \(n\)
gravitons has roots

\[
 \tau_n^2=\frac{3\omega}{4c_g}\bigl(n-|\alpha|^2\bigr).                 \tag{64.2}
\]

Hence: sectors with \(n<|\alpha|^2\) have no real root and carry no
physical state; the sector \(n=|\alpha|^2\), when \(|\alpha|^2\) is an
integer, is the double root and has no induced inner product, exactly
as the vacuum in V52; sectors with \(n>|\alpha|^2\) are physical with
two branches \(\pm\tau_n\).  The Poisson weight of the unphysical
sectors is \(e^{-|\alpha|^2}\sum_{n<|\alpha|^2}|\alpha|^{2n}/n!\), which
is \(e^{-|\alpha|^2}\) for \(|\alpha|^2\leq1\).
\end{theorem}

\begin{proof}
Insert \(2\Lambda=-\omega|\alpha|^2/(2c_g)\) into
\(f_n(\tau)=\tfrac23\tau^2-n\omega/(2c_g)-2\Lambda\) and solve.  The
root structure follows from the sign of \(n-|\alpha|^2\), and the
double root from \(f_n'(0)=0\) when the constant term vanishes.  The
weights are the Poisson law of the coherent state.
\end{proof}

The classical time-symmetric slice, \(\tau=0\) with the wave energy
balanced by \(\Lambda\), is therefore the sector \(n=|\alpha|^2\) of the
coherent state: the one sector where the induced inner product is
singular.  The neighbouring physical sectors describe the same wave
with a small nonzero mean curvature of either sign, precisely the
turnaround branches of V43 displaced from the plane by one graviton's
worth of energy.  This is the quantum reason the classical plane is
special: it is the boundary between the unphysical and the physical
Poisson sectors of the coherent state.

\subsection{Executable certificate}

The verifier

\[
 \texttt{scripts/verify\_sm\_v19\_v61\_coherent\_state\_chain.py}    \tag{64.3}
\]

checks (64.1) for \(r=1/2\), \(v=1/3\), \(\lambda=1\); the coherent
energy and the V41 balance \(\Lambda=-\omega^2\lambda^2/8=-1/2\) for
\(c_g=3/4\), \(\omega=2\); the sector values (64.2) with
\(|\alpha|^2=3/4\), negative for \(n=0\) and positive for \(n\geq1\);
the truncated Poisson weight of the unphysical sector, approximately
\(0.472\); and the coherent responses \(79/768\) and \(7/128\) in the
two models.  It emits

\[
 \texttt{artifacts/sm\_v19\_v61\_coherent\_state\_chain.json}.       \tag{64.4}
\]

The hashes are

\[
\begin{array}{c|l}
\hbox{object}&\hbox{SHA--256}\\ \hline
\hbox{verifier}&
\texttt{388FD4C0D12C08C55B103340F976AA52F69F430B9165BD77165F4E606F423E4F}\\
\hbox{canonical payload}&
\texttt{9D7B9345DC4EA5985E17FA540FDF4BADF61D5C00E28327E0F0FCF8358124797B}.
\end{array}                                                     \tag{64.5}
\]

\begin{firewall}[Free coherent states, averaged constraint]
\label{fire:v19v61}
The coherent-state chain is free, and the sector analysis uses the
averaged linear constraint only.  The unphysical fraction is a
statement about the induced inner product of that constraint; it does
not decide how the third-order coupling of \(\tau\) to the waves
modifies the sector structure, and it does not concern the general
three-dimensional class.
\end{firewall}

\begin{decision}[V62 acquisition]
\label{dec:v19v61-next}
V62 will examine the singular sector directly: it will compute the
induced inner product of the averaged constraint with the third-order
term that couples \(\hat\tau\) to the graviton number through the
mean-curvature rate (46.2), namely the term \(\propto\hat\tau\,\widehat N\)
that lifts the double root, and determine whether the classical sector
becomes physical with a finite induced norm once the lowest-order
dynamics of \(\tau\) is included.  This is the operator form of the
turnaround rate \(\dot\tau=|K|^2-\Lambda\) at the plane.
\end{decision}

\begin{assessment}[V61 acquisition]
\label{ass:v19v61}
The classical branch of V41 is realized as a coherent-state chain with
exact moments and responses, and the averaged constraint distributes it
over graviton sectors with Poisson weights and induced mean curvatures
\(\tau_n^2\propto n-|\alpha|^2\).  At the classical balance the
time-symmetric slice is the singular sector of the induced inner
product, flanked by physical turnaround sectors, and a Poisson
fraction of the coherent state is unphysical.
\end{assessment}

\begin{record}[V61 append record]
\label{rec:v19v61}
V61 is appended to the validated V60 whose installed SHA--256 was
\texttt{A540EC3BF9BFC3B46114336961EC6A49D55CE5D5B1844D582A553C4DC39F48A6}.
After installation only V61 is the active numeric SM note; frozen
archives and unrelated notes are unchanged.
\end{record}


\section{V62: the exact mean-curvature dependence and the two-branch physical space}
\label{sec:v19v62}

V61 asked whether a higher-order term of the constraint lifts the
double root at the classical sector.  The answer is no, for an exact
reason: for constant-mean-curvature data the Hamiltonian constraint
depends on \(\tau\) through \(\tfrac23\tau^2\) exactly, at every order
in the wave, because \(|K|_g^2=|A|_g^2+\tau^2/3\) is an algebraic
identity.  The singular locus of the induced inner product is
therefore structural: it is the set of states whose \(\tau\)-independent
constraint eigenvalue vanishes, and no interaction term moves it.  The
consequence is the two-branch structure familiar from constraints
quadratic in an internal time: the physical space is a direct sum of a
contracting and an expanding branch, the Euclidean reflection swaps
them, the reconstructed Euclidean plane vector is their symmetric
combination, and the classical time-symmetric slice is the measure-zero
junction.  The physical norm of a coherent state is finite and is
bounded explicitly.

\subsection{Exactness of the \(\tau\)-dependence}

\begin{proposition}[Algebraic identity]
\label{prop:v19v62-exact}
For any Riemannian \(g\), any \(g\)-traceless symmetric \(A\), and
\(K=A+\tfrac\tau3g\): \(\operatorname{tr}_gK=\tau\) and
\(|K|_g^2=|A|_g^2+\tau^2/3\).  Hence for CMC data the Hamiltonian
constraint reads exactly

\[
 {\cal H}=R(g)-|A|_g^2+\tfrac23\tau^2-2\Lambda,                        \tag{65.1}
\]

with \(\tau\) entering only through \(\tfrac23\tau^2\), at every order of
any expansion of \(g\) and \(A\).
\end{proposition}

\begin{proof}
\(|K|^2=|A|^2+\tfrac{2\tau}3\operatorname{tr}_gA+\tfrac{\tau^2}9|g|_g^2
=|A|^2+\tau^2/3\) since \(\operatorname{tr}_gA=0\) and \(|g|_g^2=3\).
The certificate checks this on a rational non-diagonal metric.
\end{proof}

\begin{theorem}[The singular locus is structural]
\label{thm:v19v62-locus}
Let \(\widehat{\cal C}=\tfrac23\hat\tau^2-\widehat E-2\Lambda\) be the
averaged constraint operator on \({\cal F}\otimes L^2(d\tau)\), where
\(\widehat E\) is the spatial average of \(|A|_g^2-R(g)\) to any
finite order, an operator on the Fock space commuting with
\(\hat\tau\).  On every spectral subspace of \(\widehat E\) with
eigenvalue \(e\), the induced inner product of Theorem
\ref{thm:v19v52-induced} is well defined with two simple roots
\(\pm\sqrt{3(e+2\Lambda)/2}\) if \(e+2\Lambda>0\), empty if
\(e+2\Lambda<0\), and singular exactly when \(e+2\Lambda=0\).  No term
of \(\widehat E\), at any order, changes this trichotomy; it only moves
the eigenvalues \(e\).
\end{theorem}

\begin{proof}
By Proposition \ref{prop:v19v62-exact} the \(\tau\)-dependence is
exactly quadratic, so the constraint is diagonal in the joint spectral
representation of \(\widehat E\) and \(\hat\tau\), with symbol
\(f_e(\tau)=\tfrac23\tau^2-e-2\Lambda\); the root structure of a
quadratic polynomial is the stated trichotomy.
\end{proof}

The classical time-symmetric slice of V41, whose wave energy is
balanced by \(\Lambda\), is precisely \(e+2\Lambda=0\) in the classical
limit: it is not a physical state but the junction of the two branches.

\subsection{Two branches and the reflection}

\begin{theorem}[Branch decomposition and reflection]
\label{thm:v19v62-branches}
Let \({\cal H}_{\rm phys}=\bigoplus_{e+2\Lambda>0}\mathbb C^2_e\) be the
induced physical space, with the branch weights
\(1/|f_e'(\pm\tau_e)|=3/(4\tau_e)\) per branch.  Write
\({\cal H}_{\rm phys}={\cal H}_+\oplus{\cal H}_-\) according to the sign
of \(\tau\).  Then:
\begin{enumerate}
\item the time reflection acts as \(\hat\tau\mapsto-\hat\tau\), hence as
      the unitary swap \({\cal H}_+\leftrightarrow{\cal H}_-\);
\item the reconstructed plane vector of every reflection-symmetric
      driftless Euclidean chain, projected to \({\cal H}_{\rm phys}\),
      lies in the symmetric subspace \(\{\psi_+=\psi_-\}\);
\item reflection positivity of the seed is the statement that the
      induced form is nonnegative on the symmetric subspace, which
      holds because both branch weights are positive.
\end{enumerate}
\end{theorem}

\begin{proof}
The first item is the oddness of \(\tau\) (35.12).  The second is
Proposition \ref{prop:v19v53-parity}: a driftless reflection-symmetric
chain has a \(\tau\)-even plane state, whose components on the two
roots are equal.  The third is (55.3) with equal components:
\(\sum_\pm3|\psi(\pm\tau_e)|^2/(4\tau_e)=\tfrac32|\psi(\tau_e)|^2/\tau_e\geq0\).
\end{proof}

Selecting a branch, a cosmological arrow, requires a drift in the
homogeneous chain (Corollary \ref{cor:v19v38-gaussian}); the Euclidean
default is the symmetric combination.

\subsection{Finite physical norm of coherent states}

\begin{proposition}[Coherent physical norm]
\label{prop:v19v62-norm}
In the setting of Theorem \ref{thm:v19v61-sectors}, let the homogeneous
plane wavefunction be \(\phi\in L^2(d\tau)\), bounded.  The physical
norm of the projected coherent state is

\[
 \|P_{\rm phys}\alpha\|^2=\sum_{n>|\alpha|^2}e^{-|\alpha|^2}\frac{|\alpha|^{2n}}{n!}
 \sum_\pm\frac{3|\phi(\pm\tau_n)|^2}{4\tau_n}
 \leq\frac{3\|\phi\|_\infty^2}{2\tau_{n_1}}\bigl(1-e^{-|\alpha|^2}\bigr),
                                                                    \tag{65.2}
\]

with \(n_1\) the smallest integer exceeding \(|\alpha|^2\) and
\(\tau_n\) from (64.2); it is finite and strictly positive.  With
\(3\omega/(4c_g)=4\) and \(|\alpha|^2=3/4\), the sectors \(n=1,3,7\)
have \(\tau_n=1,3,5\) and branch weights \(3/4,1/4,3/20\).
\end{proposition}

\begin{proof}
Insert (55.3) into the Poisson expansion; \(\tau_n\) increases with
\(n\), so \(1/\tau_n\leq1/\tau_{n_1}\), and the Poisson weights of
\(n\geq1\) sum to \(1-e^{-|\alpha|^2}\) when \(|\alpha|^2\leq1\).  The
double-root sector, if any, has Poisson weight but no induced norm and
is excluded from the projection as a set of measure zero in
\(\tau\).
\end{proof}

\subsection{Executable certificate}

The verifier

\[
 \texttt{scripts/verify\_sm\_v19\_v62\_exact\_tau\_dependence.py}    \tag{65.3}
\]

checks Proposition \ref{prop:v19v62-exact} on a rational non-diagonal
metric with a \(g\)-traceless \(A\) and \(\tau=7/3\), the perfect-square
sectors and branch weights of Proposition \ref{prop:v19v62-norm}, and
the consistency of the partial lower bound and the upper bound in
(65.2) for the truncated Poisson tail.  It emits

\[
 \texttt{artifacts/sm\_v19\_v62\_exact\_tau\_dependence.json}.       \tag{65.4}
\]

The hashes are

\[
\begin{array}{c|l}
\hbox{object}&\hbox{SHA--256}\\ \hline
\hbox{verifier}&
\texttt{15CA614B372532FA71DE7A54B7F879A2BA1773DF6E6A82D26B6BF981EAC323F1}\\
\hbox{canonical payload}&
\texttt{62BEFFFB10C08AC1BBF9D25355F4441CD767698639F651F9008A989BB757CC39}.
\end{array}                                                     \tag{65.5}
\]

\begin{firewall}[Averaged constraint, CMC, finite modes]
\label{fire:v19v62}
The exactness statement is for CMC data and the averaged constraint;
the operator \(\widehat E\) is treated abstractly, and its eigenvalues
beyond second order are not computed.  The branch structure is that
of a constraint quadratic in one variable; interacting modifications
of the mean-zero constraints, the general class, and the mode-number
limit remain open.
\end{firewall}

\begin{decision}[V63 acquisition]
\label{dec:v19v62-next}
V63 will use the branch structure to define the physical Euclidean
process of one branch: the reflection-symmetric chain with a
homogeneous drift chosen so that its reconstructed plane vector,
projected to \({\cal H}_{\rm phys}\), lies in \({\cal H}_-\), the
expanding branch after the plane in the V43 convention.  It will show
that such a drift preserves reflection positivity (V38), compute the
expected mean curvature on the first links, and compare it with the
turnaround rate \(\dot\tau=|K|^2-\Lambda\) evaluated in the coherent
state, which is the first dynamical check of the branch chain against
the corrected Einstein rate.
\end{decision}

\begin{assessment}[V62 acquisition]
\label{ass:v19v62}
The mean-curvature dependence of the CMC Hamiltonian constraint is
exactly quadratic, so the singular locus of the induced inner product
is structural and the physical space of the averaged constraint is a
direct sum of two branches swapped by the Euclidean reflection.
Reflection-symmetric chains represent the symmetric combination, a
drift selects a branch, the classical time-symmetric slice is the
junction, and coherent states have finite physical norm.
\end{assessment}

\begin{record}[V62 append record]
\label{rec:v19v62}
V62 is appended to the validated V61 whose installed SHA--256 was
\texttt{85850FEEF86DDC6B744486F77B959783433CA15E3B2D2A1998AB95C0D1FCE1A9}.
After installation only V62 is the active numeric SM note; frozen
archives and unrelated notes are unchanged.
\end{record}


\section{V63: real Euclidean plane states are branch symmetric; drifts are Euclidean velocities}
\label{sec:v19v63}

V62 left the choice of a branch to a drift in the homogeneous chain.
This note shows that no real reflection-positive Euclidean process can
make that choice.  The reconstructed plane state of any real chain is a
real function of the plane variables, and the Fourier transform of a
real function has equal modulus at \(\pm\tau\): the two branches of
V62 always carry equal weight.  A real drift does not touch the plane
state at all; it modifies the transfer contraction, which then ceases
to satisfy detailed balance, and it gives the homogeneous mode a
Euclidean mean velocity, the throat rate \(\Lambda\) of V44, whose
Lorentzian counterpart is imaginary.  The remark in V38 and V53 that a
drift selects a branch is therefore withdrawn: a cosmological arrow at
the plane requires a complex Euclidean weight, outside the positive
measures of this programme.  The first-link Euclidean mean curvature of
the drifted chain is computed and matched to the Euclidean rate.

\subsection{Branch symmetry of real plane states}

\begin{theorem}[Real plane states weigh both branches equally]
\label{thm:v19v63-symmetry}
Let \(\mu\) be any reflection-symmetric chain with real plane law
\(\pi_0\) on \(V\oplus\mathbb R\), and let \(\Psi\in L^2(\pi_0)\) be a
real function, in particular the reconstructed plane vector \(1\) or
any real observable.  Write \(\hat\Psi(h,\tau)\) for its partial Fourier
transform in the homogeneous variable \(\phi\), after transporting
\(\pi_0\) to Lebesgue measure by the square root of its density.  Then
\(|\hat\Psi(h,\tau)|=|\hat\Psi(h,-\tau)|\) for almost every
\((h,\tau)\).  Consequently the projection of \(\Psi\) onto the
physical space of Theorem \ref{thm:v19v62-branches} has components of
equal norm in \({\cal H}_+\) and \({\cal H}_-\), whatever the kernels,
drifts, or plane law.
\end{theorem}

\begin{proof}
The transported function \(\Psi\sqrt{d\pi_0/d\phi}\) is real; the
Fourier transform of a real function satisfies
\(\hat F(-\tau)=\overline{\hat F(\tau)}\), which gives the equality of
moduli.  The branch components are the restrictions to \(\pm\tau_e\),
weighted by the common factor \(3/(4\tau_e)\).
\end{proof}

\begin{corollary}[No arrow from a positive Euclidean measure]
\label{cor:v19v63-no-arrow}
A physical state with unequal branch weights, in particular a
one-branch state, is not the reconstructed plane vector of any
reflection-positive process with a real, positive measure.  It
requires a complex plane weight, for instance \(e^{i\tau_0\phi}\), and
its reflection positivity would have to be proved for a modified
reflection.  The statement in V38 and V53 that a drift selects a branch
is withdrawn.
\end{corollary}

\begin{proof}
Theorem \ref{thm:v19v63-symmetry} applies to every real plane state,
drift or not.
\end{proof}

\subsection{What a real drift does}

\begin{proposition}[Drifts act on the transfer, not on the plane state]
\label{prop:v19v63-drift}
For the drifted Gaussian chain of Corollary \ref{cor:v19v38-gaussian}
with kernels \(N(r_0\phi+b_n,s^2)\) on the homogeneous mode:
\begin{enumerate}
\item the plane law and hence the reconstructed plane vector are
      unchanged;
\item the transfer contraction satisfies detailed balance with respect
      to the stationary Gaussian if and only if \(b_n=0\); for
      \(b_n\neq0\) it is the composition of a translation with the
      Ornstein--Uhlenbeck semigroup and is not self-adjoint, so it is
      not \(e^{-\varepsilon H}\) for any self-adjoint \(H\);
\item the expected link mean curvatures are
      \(E[\tau_{j+1/2}]=(c_{j+1}-c_j)/\varepsilon\) with
      \(c_{j+1}=r_0c_j+b_j\), \(c_0=0\), the Euclidean classical
      velocity of the mean path.
\end{enumerate}
\end{proposition}

\begin{proof}
The first item is the mirror construction of V38.  For the second,
the detailed-balance ratio
\(\pi(x)P(x,y)/(\pi(y)P(y,x))=\exp\{[(x-ry-b)^2-(y-rx-b)^2]/(2s^2)+(y^2-x^2)/(2v)\}\)
equals one identically exactly when \(b=0\); a translation is unitary
and does not commute with the Ornstein--Uhlenbeck semigroup, so the
composition is not self-adjoint.  The third item is (41.3).
\end{proof}

\begin{proposition}[Euclidean rate on the first link]
\label{prop:v19v63-rate}
Choosing \(b_0=\varepsilon^2\Lambda/2\) and \(b_j=0\) for \(j\geq1\)
gives \(E[\tau_{1/2}]=\varepsilon\Lambda/2\), the Euclidean turnaround
rate \(\dot\tau_E(0)=\Lambda\) of (47.2) integrated over half a link,
and \(E[\tau_{j+1/2}]=(r_0-1)r_0^{j-1}\varepsilon\Lambda/2\) for
\(j\geq1\), decaying into the bulk.  Its analytic continuation is the
Lorentzian rate \(-\Lambda\) of (46.3); the sign reversal is the
Euclidean--Lorentzian relation of Proposition \ref{prop:v19v44-throat},
not a branch choice.
\end{proposition}

\begin{proof}
Insert the drift into (41.3).
\end{proof}

\subsection{Consequence for the programme}

The Euclidean side of the compact programme can represent the
time-symmetric superposition of the two turnaround branches, and it can
carry Euclidean mean velocities, but it cannot carry a cosmological
arrow at the reflection plane.  This is consistent with, and sharpens,
the parity theorem of V32 and the branch structure of V62: real
reflection-positive processes describe the junction, and the arrow is
a property of the reconstructed state that must be imposed after
reconstruction, by projecting onto \({\cal H}_+\) or \({\cal H}_-\),
which is legitimate because the two branches are orthogonal there.
The Euclidean measure itself remains symmetric.

\subsection{Executable certificate}

The verifier

\[
 \texttt{scripts/verify\_sm\_v19\_v63\_branch\_symmetry.py}          \tag{66.1}
\]

checks the Hermitian symmetry \(X_{-k}=\overline{X_k}\) of the exact
length-four discrete Fourier transform of a real rational sequence in
Gaussian rationals, the link means of the drifted chain with
\(\varepsilon=1/5\), \(\Lambda=-1/2\), \(r_0=1/3\), giving
\(E[\tau_{1/2}]=-1/20\), and the detailed-balance ratio of the
Gaussian kernel, zero log-ratio without drift and nonzero with drift
\(1/4\).  It emits

\[
 \texttt{artifacts/sm\_v19\_v63\_branch\_symmetry.json}.             \tag{66.2}
\]

The hashes are

\[
\begin{array}{c|l}
\hbox{object}&\hbox{SHA--256}\\ \hline
\hbox{verifier}&
\texttt{A792FFF97A3472594AA4D9EB2C7F9BAA6E1FFDC6742CF3138C28F0552E94672D}\\
\hbox{canonical payload}&
\texttt{3502E3BC8DAC54A3EEFC88F9E8FEA84A3862FB789E547A9F6DE4026C9E5656D5}.
\end{array}                                                     \tag{66.3}
\]

\begin{firewall}[Real measures only]
\label{fire:v19v63}
The no-arrow statement concerns real positive Euclidean measures with
the site reflection; complex weights with a modified reflection are
not analyzed.  The Euclidean rate matching is at the level of
conditional means of a Gaussian chain, as in V44, and is not a
dynamical theorem.
\end{firewall}

\begin{decision}[V64 acquisition]
\label{dec:v19v63-next}
V64 will consolidate the operator-level programme into a single
theorem for the finite-mode plane-symmetric compact model: reflection
positivity of the seed, exact spectral consistency, the induced
two-branch physical space with its constraint conditions (60.5), the
branch-symmetric plane states, the second-order dependent operators,
and the classical coherent sector, stated with all hypotheses and
firewalls in one place, as the reference point for the third-order
and three-dimensional extensions.  The V65 audit follows.
\end{decision}

\begin{assessment}[V63 acquisition]
\label{ass:v19v63}
Real reflection-positive Euclidean processes are branch symmetric at
the plane: the Fourier transform of a real plane state has equal
moduli at opposite mean curvatures, so no drift or plane law selects a
cosmological arrow.  Drifts act on the transfer, break detailed
balance, and encode Euclidean mean velocities, whose first-link value
reproduces the throat rate \(\Lambda\).  The arrow is a property of
the reconstructed state, imposed by projection onto one branch.
\end{assessment}

\begin{record}[V63 append record]
\label{rec:v19v63}
V63 is appended to the validated V62 whose installed SHA--256 was
\texttt{328874681C5883CCBAF41B44826952C3CE6B8EC8B988993DC58D5A4FA8866004}.
After installation only V63 is the active numeric SM note; frozen
archives and unrelated notes are unchanged.
\end{record}


\section{V64: consolidated theorem for the finite-mode plane-symmetric compact model}
\label{sec:v19v64}

V32--V63 built, piece by piece, an operator-level account of the
compact \(\Lambda<0\) torus programme in the plane-symmetric,
finite-mode, second-order setting.  This note states the result as one
theorem with all hypotheses, so that later extensions have a fixed
reference, and certifies mechanically that every executable
certificate quoted since V32 still reproduces its recorded hashes.
Nothing new is proved; everything is cross-referenced.

\subsection{Hypotheses}

\begin{definition}[The finite-mode plane-symmetric compact model]
\label{def:v19v64-model}
Fix: a flat three-torus with \(z\)-period \(2\pi\); a finite set
\({\cal K}\) of positive integers; the TT modes
\(\cos kz\,e^A\) and \(\sin kz\,e^A\), \(k\in{\cal K}\),
\(A\in\{+,\times\}\), with frequencies \(\omega_k=k\) and the Einstein
normalization \(v_k=1/(2c_g\omega_k)\); a homogeneous scalar mode
\(\phi\) with a stationary Gaussian law of variance \(v_0\); a
cosmological constant \(\Lambda<0\); a lattice spacing \(\varepsilon\)
and a smearing time \(T=m\varepsilon\).  The Euclidean seed is the
reflection-symmetric chain of Theorem \ref{thm:v19v38-chain} with
Ornstein--Uhlenbeck kernels of ratios \(e^{-\varepsilon\omega_k}\) on
the TT modes and any bounded-support or Gaussian kernel on \(\phi\),
with a plane law that is any probability measure on the plane
variables.
\end{definition}

\subsection{Statement}

\begin{theorem}[Consolidated statement]
\label{thm:v19v64-consolidated}
In the model of Definition \ref{def:v19v64-model}:
\begin{enumerate}
\item \textup{(Positivity.)}  The seed is reflection positive for every
      plane law, with reflected form
      \(\langle F,F\rangle_\theta=E[|E[F\mid q_0]|^2]\)
      (Theorem \ref{thm:v19v38-chain}, Proposition
      \ref{prop:v19v46-conditioning}); link differences are odd under
      the reflection (Lemma \ref{lem:v19v33-links}); any equivariant,
      positive-time-local map of the seed transports positivity
      (Theorem \ref{thm:v19v32-pushforward}); and no \(L^2\)-continuous
      odd field has a nonzero sharp-time restriction on the plane
      (Theorem \ref{thm:v19v32-sharp}).
\item \textup{(Spectral consistency.)}  The reconstructed plane space is
      \(L^2\) of the plane law with the contractions of Theorem
      \ref{thm:v19v38-reconstruction}; with the Gaussian plane law the
      transfer is \(e^{-\varepsilon\sum_k\omega_k\widehat N_k}\) exactly,
      the linearized Einstein Fock Hamiltonian (Theorem
      \ref{thm:v19v44-spectrum}).
\item \textup{(Classical plane.)}  Time-symmetric torus data with a
      wave exist exactly when \(\Lambda\) balances the Dirichlet energy
      at second order (Theorems \ref{thm:v19v39-balance},
      \ref{thm:v19v41-branch}); at fixed \(\Lambda\) they form the
      analytic sphere of Theorem \ref{thm:v19v42-sphere}; the plane is
      a Lorentzian volume maximum and a Euclidean volume minimum
      (Propositions \ref{prop:v19v43-rate}, \ref{prop:v19v44-throat}),
      and on the torus it is flat unless \(\Lambda<0\) (Theorem
      \ref{thm:v19v39-rigidity}).
\item \textup{(Second-order constraints.)}  In the plane-symmetric class
      the Hamiltonian constraint through second order has the exact
      Fourier components (59.3)--(59.4); the mean-zero components are
      solved by the linear scalar response (59.5), the longitudinal
      momentum by (60.4), the transverse momenta are homogeneous
      constants (Theorem \ref{thm:v19v57-solution}), and the
      \(\tau\)-dependence is exactly \(\tfrac23\tau^2\) (Proposition
      \ref{prop:v19v62-exact}).
\item \textup{(Physical states.)}  The states satisfy (60.5): the
      normal-ordered averaged constraint and zero total momentum.  The
      induced physical space is the two-branch sum of Theorem
      \ref{thm:v19v62-branches}, with sectors labelled by graviton
      number and mean curvature \(\pm\tau_e\), \(\tfrac23\tau_e^2=e+2\Lambda\);
      the sector \(e+2\Lambda=0\) is singular and contains the classical
      time-symmetric slice (Theorems \ref{thm:v19v52-induced},
      \ref{thm:v19v61-sectors}).  Every real plane state weighs the two
      branches equally (Theorem \ref{thm:v19v63-symmetry}).
\item \textup{(Zero point.)}  The zero-point energy of the retained modes
      sits entirely in the averaged constraint, as the c-number
      \(Z\) of (55.2), and is removed by normal ordering, which is
      legitimate on states; mean-zero vacuum responses vanish in the
      translation-invariant model (Theorem \ref{thm:v19v59-invariance}).
      Imposing the constraint pathwise on the Euclidean seed instead
      leads to the divergent variance, conformal divergence, and
      positivity loss of Theorems \ref{thm:v19v48-variance},
      \ref{thm:v19v49-divergence}, \ref{thm:v19v51-dichotomy},
      \ref{thm:v19v54-failure}.
\item \textup{(States realized.)}  Graviton number states and coherent
      states are the chains with Hermite-weighted and shifted-Gaussian
      plane laws (Theorems \ref{thm:v19v53-chain},
      \ref{thm:v19v61-chain}); the classical branch of item 3 is the
      coherent-state sector, with second-order responses given by the
      classical values plus vacuum values (Propositions
      \ref{prop:v19v58-longitudinal}, \ref{prop:v19v61-response}).
\end{enumerate}
\end{theorem}

\begin{proof}
Each item is the cited result, whose hypotheses are those of the
definition or weaker.
\end{proof}

\subsection{What is not in the statement}

The theorem does not contain: third-order terms, in particular the
dynamical coupling of \(\tau\) to the waves that makes the reduced
dynamics non-autonomous (V36); the general three-dimensional class,
where the momentum constraint is not an ordinary differential
identity; any limit in the mode number, where \(Z\) diverges (V51) and
the balancing structures of V47 fail; any limit in the lattice spacing
beyond the exact-in-\(\varepsilon\) spectral statement; interacting
kernels; a cosmological arrow, which needs a complex weight (V63);
and any Standard Model sector, whose finite reflection-positive forms
of V6--V16 remain uncoupled to the gravitational chain.

\subsection{Consolidation certificate}

The verifier

\[
 \texttt{scripts/verify\_sm\_v19\_v64\_consolidation\_hashes.py}    \tag{67.1}
\]

parses the predecessor note, the installed V63 whose SHA--256 is
recorded below, extracts every verifier name and its two quoted
hashes from V32 onward, recomputes the script hash, re-executes the
script, and compares the emitted payload hash.  All twenty-six
certificates from V32 to V63 reproduce their quoted hashes.  It emits

\[
 \texttt{artifacts/sm\_v19\_v64\_consolidation\_hashes.json}.       \tag{67.2}
\]

The hashes are

\[
\begin{array}{c|l}
\hbox{object}&\hbox{SHA--256}\\ \hline
\hbox{verifier}&
\texttt{E65EE16F5026661A0DFDE44A86F001363B62EC4C4FEB7768FFDDCFA5371D7D21}\\
\hbox{canonical payload}&
\texttt{2E05FCC579D4AB71D8BC806B303B368AE85CDEC7FB0E0E2E5907899F8D2CB556}.
\end{array}                                                     \tag{67.3}
\]

\begin{firewall}[A consolidation, not a continuum]
\label{fire:v19v64}
Theorem \ref{thm:v19v64-consolidated} concerns a finite-mode,
plane-symmetric, second-order model with free kernels.  It is not a
construction of quantum gravity on the torus, not a continuum limit,
and not a Standard Model--Einstein theory.  Its value is that every
hypothesis and every gap is now explicit in one place.
\end{firewall}

\begin{decision}[V65 audit and V66 acquisition]
\label{dec:v19v64-next}
V65 is the compulsory companion audit.  V66 will begin the third
order in the plane-symmetric class, where the mean curvature first
couples dynamically: it will expand the Riccati form (59.1) to third
order, extract the cubic constraint components, and determine the
first correction to the graviton energies \(e\) in (60.5), hence the
first shift of the branch labels \(\tau_e\), which is where the
turnaround dynamics of V43 enters the state conditions.
\end{decision}

\begin{assessment}[V64 acquisition]
\label{ass:v19v64}
The operator-level programme of V32--V63 is consolidated into one
theorem for the finite-mode plane-symmetric compact model, with
positivity, exact spectral consistency, the classical plane, the
second-order constraints, the two-branch physical space, the
zero-point accounting, and the realized states, and all twenty-six
certificates since V32 are mechanically re-verified.  The listed
absences define the remaining programme.
\end{assessment}

\begin{record}[V64 append record]
\label{rec:v19v64}
V64 is appended to the validated V63 whose installed SHA--256 was
\texttt{411201E03820F4DCC05241F27FA2F8BBFEFBBAC84F8286F67431BD9EDC5807E1}.
After installation only V64 is the active numeric SM note; frozen
archives and unrelated notes are unchanged.
\end{record}


\section{V65: compulsory companion audit after the consolidation}
\label{sec:v19v65}

This is the required audit following V60.  Every active companion note
in the shared folder was inspected read-only before the third-order
expansion announced in V64 begins.

\begin{record}[V65 companion snapshot]
\label{rec:v19v65-snapshot}
The inspected files were

\[
\begin{array}{c|r|r}
\hbox{file and SHA--256}&\hbox{bytes}&\hbox{lines}\\ \hline
\texttt{Renewal\_Yang\_Mills\_Mass\_Gap\_Research\_Notes\_V32.tex}
&452376&11540\\
\multicolumn{3}{l}{
\texttt{F1DE1DC4AA157BC4DF069FEDB4F34C3032565655348CB9AE35AC95F5695744D3}}
\\
\texttt{Renewal\_NS\_Stability\_Research\_Notes\_V26.tex}
&278283&5319\\
\multicolumn{3}{l}{
\texttt{82C887F8C2C69E4F28FAD7C3DC8DE5B9099CB5A4BF431EBA1FFF3E91935978AD}}
\\
\texttt{Renewal\_GRH\_Cofinal\_Visibility\_Attack\_V52.tex}
&504761&15514\\
\multicolumn{3}{l}{
\texttt{995305F804828109C77F2E105081AB0B3AE87F103ECAF34F67E1364E06BE92E0}}
\\
\texttt{Renewal\_YM\_Parallel\_Research\_Notes\_V5.tex}
&54174&1351\\
\multicolumn{3}{l}{
\texttt{306F1BB84CFA8A8D47EA569EC17E9545F9867C8D32B548EC9D9C444B4F3C0512}}.
\end{array}                                                     \tag{68.1}
\]

All four files are byte-identical to the V60 snapshot, and no new
note series has appeared.
\end{record}

\begin{decision}[V65 audit decision]
\label{dec:v19v65-audit}
No companion theorem, numerical constant, or executable artifact is
imported.  The V64 direction stands: V66 begins the third order of the
plane-symmetric class.
\end{decision}

\begin{proof}
The companions are unchanged since V55; the V66 task is an internal
expansion of the exact Riccati form (59.1), whose second-order part
was certified in V56 without any companion input.
\end{proof}

\begin{assessment}[V65 acquisition]
\label{ass:v19v65}
The audit is current.  Since V60 the programme has realized coherent
states as chains, located the classical branch on the singular locus
of the induced inner product, proved that locus structural through
the exact \(\tau\)-dependence, shown that real Euclidean plane states
are branch symmetric, and consolidated V32--V63 into one theorem with
a mechanical re-verification of all its certificates.
\end{assessment}

\begin{record}[V65 append record]
\label{rec:v19v65}
V65 is appended to the validated V64 whose installed SHA--256 was
\texttt{9FCF076E9A2B8639910DEE09A6A0F43087B6119FF195B9E39A811A4E48C44C0C}.
After installation only V65 is the active numeric SM note; the
companions, frozen archives, and all unrelated notes are unchanged.  The
next compulsory companion audit is V70.
\end{record}


\section{V66: the third order vanishes and the fourth-order balance of the classical plane}
\label{sec:v19v66}

V64 announced the third order of the plane-symmetric class as the
place where the mean curvature first couples dynamically.  The third
order turns out to be empty: for two-by-two traceless matrices every
cubic invariant that enters the Riccati form vanishes, and the
conformal and momentum terms contribute no cubic piece either.  The
first correction beyond second order is therefore quartic.  For the
time-symmetric classical plane with the Fourier wave the fourth-order
balance is computed exactly, giving the cosmological constant to
fourth order in the amplitude, \(\Lambda=-\lambda^2/8-145\lambda^4/2048+O(\lambda^6)\).
In the state picture this is the first interaction correction to the
sector labels of (60.5): the graviton energies acquire terms quadratic
in the occupation numbers.

\subsection{The third order vanishes}

\begin{theorem}[No cubic terms in the plane-symmetric class]
\label{thm:v19v66-cubic}
Let \(H\) and \(H'\) be two-by-two symmetric traceless matrices.  Then
\(\operatorname{tr}H^3=0\), \(\operatorname{tr}(HH'^2)=0\), and in the
expansion of (59.1) with \(G=I+H\) the third-order part of the scalar
curvature is

\[
 R_3=-\tfrac13\bigl(\operatorname{tr}H^3\bigr)''+\tfrac12\operatorname{tr}(HH'^2)=0.
                                                                    \tag{69.1}
\]

Likewise \(|A|_\gamma^2\) has no third-order term for a
\(\gamma\)-traceless \(A\), the conformal Laplacian correction
\((\Delta_\gamma-\Delta)w_2\) is fourth order, and the exact
\(\tau\)-dependence is quadratic.  Hence for the time-symmetric plane
\(w_3=0\) and \(\Lambda_3=0\), and for general data the third-order
constraint is identically satisfied by the second-order dependent
operators.
\end{theorem}

\begin{proof}
A traceless symmetric two-by-two matrix satisfies
\(M^2=\tfrac12\operatorname{tr}(M^2)\,I\), so
\(H^3=\tfrac12\operatorname{tr}(H^2)H\) has zero trace and
\(HH'^2=\tfrac12\operatorname{tr}(H'^2)H\) has zero trace.  In the
expansion \(K=\tfrac12(I-H+H^2-\dots)H'\), the third-order part of
\(\operatorname{tr}K\) is \(\tfrac12\operatorname{tr}(H^2H')=\tfrac16(\operatorname{tr}H^3)'\)
and that of \(\operatorname{tr}K^2\) is \(-\tfrac12\operatorname{tr}(HH'^2)\),
while \((\operatorname{tr}K)^2\) starts at fourth order; this gives
(69.1).  For \(|A|^2_\gamma=\operatorname{tr}((G^{-1}A)^2)\), the cubic
part is \(-2\operatorname{tr}(HA_0^2)+2\operatorname{tr}(A_0A_2)\) with
\(A_2\) the second-order trace correction \(\tfrac12\operatorname{tr}(HA_0)I\)
plus a traceless part; both terms vanish by the same identity and by
\(\operatorname{tr}A_0=0\).  The Laplacian on functions of \(z\) is
\(f''+\tfrac12(\ln\det G)'f'\) with \(\ln\det G=-\tfrac12\operatorname{tr}H^2+O(4)\),
so its correction on \(w_2\) is fourth order.  The remaining claims
follow from Proposition \ref{prop:v19v62-exact} and the expansion of
(44.1).  The certificate confirms \(R_3=0\) for the Fourier wave.
\end{proof}

\subsection{Fourth-order balance of the time-symmetric plane}

\begin{theorem}[Fourth-order balance]
\label{thm:v19v66-fourth}
For the time-symmetric plane \(g=\psi^4\gamma_\lambda\) with
\(\psi=1+\lambda^2w_2+\lambda^4w_4+O(\lambda^6)\), \(\langle w\rangle=0\),
and \(R(g)=2\Lambda(\lambda)\), the averaged fourth-order equation is

\[
 2\Lambda_4=\langle R_4\rangle-4\bigl\langle(\ln\det G)'_2\,w_2'\bigr\rangle
 +\langle R_2w_2\rangle,                                               \tag{69.2}
\]

where \(R(\gamma_\lambda)=\lambda^2R_2+\lambda^4R_4+O(\lambda^6)\) and
\((\ln\det G)'_2\) is the second-order part of \((\ln\det G)'\).  For
the Fourier wave \(h=\cos z\,e^+\),

\[
 R_4=-\frac5{16}-\cos2z-\frac{11}{16}\cos4z,\qquad
 \langle R_2w_2\rangle=-\frac{49}{1024},\qquad
 \bigl\langle(\ln\det G)'_2w_2'\bigr\rangle=-\frac7{128},              \tag{69.3}
\]

so that

\[
 \Lambda(\lambda)=-\frac{\lambda^2}8-\frac{145}{2048}\lambda^4+O(\lambda^6).
                                                                    \tag{69.4}
\]
\end{theorem}

\begin{proof}
Expand \(R(\psi^4\gamma)=\psi^{-5}(-8\Delta_\gamma\psi+R(\gamma)\psi)\)
to fourth order, using \(\Delta_\gamma f=f''+\tfrac12(\ln\det G)'f'\):
the fourth-order part is
\(-8w_4''-4(\ln\det G)'_2w_2'+R_4+R_2w_2-5w_2(-8w_2''+R_2)\), and
\(-8w_2''+R_2=2\Lambda_2\) by the second-order equation, so the last
term averages to \(-10\Lambda_2\langle w_2\rangle=0\); averaging the
rest gives (69.2).  For the Fourier wave, the exact series of
\(R(\gamma_\lambda)\) to fourth order, obtained from the diagonal
formula of Proposition \ref{prop:v19v41-fourier}, gives \(R_4\) as
stated; \(w_2=\tfrac7{128}\cos2z\) and \(R_2=-\tfrac14-\tfrac74\cos2z\)
give \(\langle R_2w_2\rangle\); and
\(\ln\det G=\ln(1-\lambda^2\cos^2z)\) gives
\((\ln\det G)'_2=\sin2z\), so \(\langle(\ln\det G)'_2w_2'\rangle
=\langle\sin2z\cdot(-\tfrac7{64}\sin2z)\rangle=-\tfrac7{128}\).  Then
\(2\Lambda_4=-\tfrac5{16}+\tfrac7{32}-\tfrac{49}{1024}=-\tfrac{145}{1024}\).
\end{proof}

The fourth-order coefficient is negative: the self-interaction of the
standing wave requires a slightly more negative cosmological constant
than the quadratic balance, by \(145/2048\) of the fourth power of the
amplitude.  In the sector labels of (60.5), where \(2\Lambda\) equals
minus the normal-ordered graviton energy density at the classical
plane, this is the first term quadratic in the occupation number: the
graviton energy \(e_n\) of the standing mode acquires a correction
proportional to \(n^2\), the plane-symmetric analogue of a
graviton--graviton interaction energy, with the coherent-state
amplitude \(\lambda^2\leftrightarrow2v_kn\) of Corollary
\ref{cor:v19v58-classical}.

\subsection{Executable certificate}

The verifier

\[
 \texttt{scripts/verify\_sm\_v19\_v66\_fourth\_order\_balance.py}    \tag{69.5}
\]

extends the exact series machinery of V41 to fourth order, confirms
that the orders \(0\), \(1\), and \(3\) of \(R(\gamma_\lambda)\) vanish
and that the second order is the V41 series, computes \(R_4\), the
averages in (69.3), and \(\Lambda_4=-145/2048\).  It emits

\[
 \texttt{artifacts/sm\_v19\_v66\_fourth\_order\_balance.json}.       \tag{69.6}
\]

The hashes are

\[
\begin{array}{c|l}
\hbox{object}&\hbox{SHA--256}\\ \hline
\hbox{verifier}&
\texttt{528EE728B081BCD6972514C371B713688F0FDC1909FF2660EF3681442D7ABE80}\\
\hbox{canonical payload}&
\texttt{BD3052BB5D5CD99D492EC177BCC575866F1F26C55E2FC9BEE08CD43C57F16CEE}.
\end{array}                                                     \tag{69.7}
\]

\begin{firewall}[One wave, time-symmetric plane, fourth order]
\label{fire:v19v66}
Theorem \ref{thm:v19v66-fourth} is computed for one Fourier wave at
the time-symmetric plane; the general fourth-order constraint
operator with its mode couplings, the corresponding shifts of every
sector label, and the ordering of the quartic operator products are
not derived.  The vanishing third order is exact for the two-by-two
class only.
\end{firewall}

\begin{decision}[V67 acquisition]
\label{dec:v19v66-next}
V67 will derive the general fourth-order averaged constraint of the
plane-symmetric class as a quartic form in the mode amplitudes and
momenta, using the two-by-two identities of Theorem
\ref{thm:v19v66-cubic} to reduce every quartic invariant to products
of traces, and will extract the number-conserving part of the
corresponding Fock operator, which is the first interaction correction
to the sector energies \(e_n\) of the physical-state condition (60.5).
\end{decision}

\begin{assessment}[V66 acquisition]
\label{ass:v19v66}
The third order of the plane-symmetric class vanishes identically by
two-by-two trace identities, so the second-order theory of V56--V62 is
exact through third order, and the first correction is quartic.  For
the classical Fourier plane the fourth-order balance is exact,
\(\Lambda_4=-145/2048\), which is the first interaction correction to
the sector labels of the physical states.
\end{assessment}

\begin{record}[V66 append record]
\label{rec:v19v66}
V66 is appended to the validated V65 whose installed SHA--256 was
\texttt{480395E143C28E5D7A310B2048DCCA02A85E6E4BCB8B07CE00670DDA9A0FC287}.
After installation only V66 is the active numeric SM note; frozen
archives and unrelated notes are unchanged.
\end{record}


\section{V67: the quartic averaged constraint of one standing mode and its Fock diagonal}
\label{sec:v19v67}

V66 showed that the first correction beyond second order is quartic
and computed it for the classical plane.  This note derives the full
quartic averaged constraint of one standing mode with its momentum, as
an exact quartic form in the amplitude and momentum, and extracts its
number-conserving part in the Fock representation.  The result is the
first interaction correction to the sector energies of the
physical-state condition (60.5): a shift affine in the occupation
number plus a term quadratic in it, the plane-symmetric
graviton self-energy and graviton--graviton interaction of the mode.
The ordering ambiguity of the mixed term is a constant per mode and is
recorded.

\subsection{The quartic form}

Take one mode \(h=x\cos kz\,e^+\), TT momentum \(A_0=p\cos kz\,e^+\),
with the \(\gamma\)-traceless completion of \(A\), and the second-order
conformal response \(w_2=(\tfrac7{128}x^2+\tfrac{p^2}{32k^2})\cos2kz\)
from (59.4)--(59.5).

\begin{theorem}[Quartic averaged constraint of one mode]
\label{thm:v19v67-quartic}
The fourth-order part of the spatially averaged Hamiltonian constraint
\(\langle R(\psi^4\gamma)-|A|_g^2\rangle-2\Lambda\) is

\[
 F_4=-\frac{145}{1024}k^2x^4+\frac{109}{256}x^2p^2+\frac{3}{16k^2}p^4-2\Lambda_4,
                                                                    \tag{70.1}
\]

obtained from
\(F_4=\langle R_4\rangle-4\langle(\ln\det G)'_2w_2'\rangle+\langle R_2w_2\rangle
+12\langle w_2|A_0|^2\rangle-2\Lambda_4\),
where \(\langle R_4\rangle=-\tfrac5{16}k^2x^4\),
\(\langle R_2w_2\rangle=-\tfrac{49}{1024}k^2x^4-\tfrac7{256}x^2p^2\),
\(-4\langle(\ln\det G)'_2w_2'\rangle=\tfrac7{32}k^2x^4+\tfrac18x^2p^2\),
\(12\langle w_2|A_0|^2\rangle=\tfrac{21}{64}x^2p^2+\tfrac3{16k^2}p^4\),
and \(|A|^2_\gamma=|A_0|^2_\delta\) exactly for the diagonal
\(\gamma\)-traceless completion.  At \(p=0\) it reproduces (69.4).
\end{theorem}

\begin{proof}
The fourth-order expansion of \(\psi^{-5}(-8\Delta_\gamma\psi+R\psi)\)
is as in the proof of Theorem \ref{thm:v19v66-fourth}, now with the
momentum part of \(w_2\); \(|A|_g^2=\psi^{-12}|A|_\gamma^2\) contributes
\(-12w_2|A_0|^2\) at fourth order; and for \(G=\operatorname{diag}(1+xc,1-xc)\),
\(c=\cos kz\), the \(\gamma\)-traceless completion
\(A=\operatorname{diag}(pc(1+xc),-pc(1-xc))\) has
\(G^{-1}A=pc\,e^+\) exactly, so \(|A|_\gamma^2=2p^2c^2=|A_0|_\delta^2\)
with no higher-order terms.  The averages use
\(\langle\cos^22kz\rangle=\langle\sin^22kz\rangle=\tfrac12\),
\((\ln\det G)'_2=kx^2\sin2kz\), and the exact series of
\(R(\gamma)\) to fourth order, whose \(k\)-dependence is the overall
\(k^2\) of a two-derivative expression.
\end{proof}

\subsection{Fock diagonal and the interaction energy}

In the Fock representation with \(\langle n|x^2|n\rangle=(2n+1)v\),
\(v=1/(2c_gk)\), the Weyl-symmetrized diagonal moments are

\[
 \langle n|x^4|n\rangle=3v^2(2n^2+2n+1),\qquad
 \langle n|x^2p^2|n\rangle_W=\frac{2n^2+2n-1}4,\qquad
 \langle n|p^4|n\rangle=\frac{3(2n^2+2n+1)}{16v^2}.                   \tag{70.2}
\]

\begin{theorem}[Fourth-order sector shift]
\label{thm:v19v67-shift}
The number-conserving part of (70.1) shifts the sector energy of
(60.5) by

\[
 e_4(n)=-\frac{145}{1024}k^2\cdot3v^2(2n^2+2n+1)
 +\frac{109}{256}\cdot\frac{2n^2+2n-1}4
 +\frac3{16k^2}\cdot\frac{3(2n^2+2n+1)}{16v^2},                        \tag{70.3}
\]

whose coefficient of \(n^2\) is

\[
 \kappa_2=-\frac{435}{512}k^2v^2+\frac{109}{512}+\frac{9}{128k^2v^2}
 =-\frac{435}{2048c_g^2}+\frac{109}{512}+\frac{9c_g^2}{32},             \tag{70.4}
\]

the plane-symmetric graviton--graviton interaction energy of the mode,
and whose constant part is the fourth-order vacuum energy, to be
absorbed by normal ordering together with \(Z\).  Any other symmetric
ordering of \(x^2p^2\) changes \(e_4(n)\) by an \(n\)-independent
constant per mode.  For \(c_g=3/4\), \(k=1\):
\(e_4(0)=-\tfrac{83}{384}\), \(e_4(1)=-\tfrac{11}{48}\),
\(e_4(2)=-\tfrac{49}{192}\), \(e_4(3)=-\tfrac{113}{384}\),
\(\kappa_2=-\tfrac5{768}\).
\end{theorem}

\begin{proof}
Insert (70.2) into (70.1) and collect powers of \(n\); the
identities (70.2) follow from \(x=\sqrt v(a+a^\dagger)\),
\(p=(a-a^\dagger)/(2i\sqrt v)\), and the diagonal matrix elements of
\((a\pm a^\dagger)^4\) and of \((a+a^\dagger)^2(a-a^\dagger)^2\), which are
\(6n^2+6n+3\) and \(-(2n^2+2n-1)\).  Orderings of \(x^2p^2\) differ by
multiples of the identity by the canonical commutation relation.
\end{proof}

In the physical-state condition, the sector label becomes
\(\tfrac23\tau_n^2=e_n^{(2)}+e_4(n)+2\Lambda\) with
\(e_n^{(2)}=n\omega_k/(2c_g)\): the branch mean curvatures of a
many-graviton sector are shifted by the interaction, negative for the
parameters above.  The classical coherent-state correspondence of
V61 identifies \(\lambda^2\leftrightarrow2vn\) and reproduces
\(2\Lambda_4=-\tfrac{145}{1024}k^2\lambda^4\) from the \(x^4\) term
alone, as it must, since a coherent state has \(\langle p\rangle=0\) at
the plane but \(\langle p^2\rangle\neq0\), and the \(p\)-dependent terms
of (70.1) then contribute the vacuum values.

\subsection{Executable certificate}

The verifier

\[
 \texttt{scripts/verify\_sm\_v19\_v67\_quartic\_constraint.py}       \tag{70.5}
\]

computes \(R_2\) and \(R_4\) with the exact series engine at \(k=1\),
assembles (70.1) with coefficients \(-145/1024\), \(109/256\), \(3/16\),
checks the classical limit against V66, and evaluates (70.3)--(70.4)
for \(c_g=3/4\).  It emits

\[
 \texttt{artifacts/sm\_v19\_v67\_quartic\_constraint.json}.          \tag{70.6}
\]

The hashes are

\[
\begin{array}{c|l}
\hbox{object}&\hbox{SHA--256}\\ \hline
\hbox{verifier}&
\texttt{5C1948B9730B89E16EA9628A931A1790745F8B9C666B0E310CB4BB8747CC0F8B}\\
\hbox{canonical payload}&
\texttt{16A167E46400091218A483F7224AF8F06111CD0C5CA909E4E040492EEA29E145}.
\end{array}                                                     \tag{70.7}
\]

\begin{firewall}[One mode, one polarization, averaged constraint]
\label{fire:v19v67}
The quartic form is derived for a single cosine mode with one
polarization; mode couplings, the second polarization, the mean-zero
quartic components, and the momentum constraint at fourth order are
not treated.  The sector shift is the number-conserving part only;
number-changing quartic terms mix sectors at higher order in
perturbation theory and are not analyzed.
\end{firewall}

\begin{decision}[V68 acquisition]
\label{dec:v19v67-next}
V68 will extend Theorem \ref{thm:v19v67-quartic} to two modes
\(k\neq k'\) with both polarizations, deriving the quartic cross terms
\(x_k^2x_{k'}^2\), \(x_k^2p_{k'}^2\), and the polarization-mixing
terms by the two-by-two trace identities, and extracting the
number-conserving part \(\kappa_{kk'}n_kn_{k'}\), the mode--mode
interaction energy; it will state whether the interaction matrix is
positive or negative definite on the two-mode space.
\end{decision}

\begin{assessment}[V67 acquisition]
\label{ass:v19v67}
The quartic averaged constraint of one standing mode is an exact
quartic form in amplitude and momentum, and its Fock diagonal gives
the first interaction correction to the sector energies of the
physical states: a graviton self-energy and a graviton--graviton
interaction \(\kappa_2n^2\), computed in closed form, with the
classical fourth-order balance of V66 as its coherent-state limit.
\end{assessment}

\begin{record}[V67 append record]
\label{rec:v19v67}
V67 is appended to the validated V66 whose installed SHA--256 was
\texttt{6C97546F2A5CDC49E22C085BD93E950D0650E86E8EC0D28B42D7279B1F2BE034}.
After installation only V67 is the active numeric SM note; frozen
archives and unrelated notes are unchanged.
\end{record}


\section{V68: the two-mode quartic constraint and the mode--mode interaction matrix}
\label{sec:v19v68}

V67 gave the quartic averaged constraint of one standing mode.  This
note extends it to two standing modes of frequencies \(1\) and \(2\) in
one polarization, reconstructing the quartic form in the four
variables \((x_1,x_2,p_1,p_2)\) exactly by interpolation on rational
points, with the exact series engine supplying \(\langle R_4\rangle\)
and the second-order responses supplying the remaining terms.  The
form has nine monomials.  Its Fock diagonal gives the sector energies
to fourth order as a quadratic polynomial in the occupation numbers,
whose off-diagonal coefficient is the mode--mode interaction energy.
For the reference coupling the self-interactions are weakly negative
and equal, the cross interaction is large and positive, and the
interaction matrix is indefinite.

\subsection{The quartic form}

Let \(h=(x_1\cos z+x_2\cos2z)e^+\), \(A_0=(p_1\cos z+p_2\cos2z)e^+\), and
let \(w_2=-\sum_{q\neq0}{\cal H}_q\cos qz/(8q^2)\) with the components
\({\cal H}_q\) of (59.4) including the momentum terms.

\begin{theorem}[Two-mode quartic averaged constraint]
\label{thm:v19v68-form}
The fourth-order part of the averaged constraint,
\(F_4=\langle R_4\rangle-4\langle(\ln\det G)'_2w_2'\rangle+\langle R_2w_2\rangle
+12\langle w_2|A_0|^2\rangle-2\Lambda_4\), is the quartic form

\[
\begin{aligned}
 F_4+2\Lambda_4
 &=-\frac{145}{1024}x_1^4-\frac{145}{256}x_2^4-\frac{95}{72}x_1^2x_2^2
 +\frac{109}{256}\bigl(x_1^2p_1^2+x_2^2p_2^2\bigr)\\
 &\quad+\frac{179}{36}x_1x_2p_1p_2
 +\frac{3}{16}p_1^4+\frac{3}{64}p_2^4+\frac{10}{3}p_1^2p_2^2 .
\end{aligned}                                                          \tag{71.1}
\]

It is even under \((x_1,p_1)\mapsto(-x_1,-p_1)\), the translation
\(z\mapsto z+\pi\); the single-mode terms agree with (70.1) with
\(k=1\) and \(k=2\); and no monomials \(x_1^2p_2^2\), \(x_2^2p_1^2\)
occur.
\end{theorem}

\begin{proof}
The parity restricts the quartic monomials to nineteen candidates.
The value of \(F_4+2\Lambda_4\) at a rational point is computed
exactly: \(\langle R_4\rangle\) by the series of Proposition
\ref{prop:v19v41-fourier} with the cosine series \(x_1\cos z+x_2\cos2z\)
in place of \(\cos z\); \(R_2\), \(w_2\), \((\ln\det G)'_2=-2x(z)x'(z)\),
and \(|A_0|^2=2(p_1\cos z+p_2\cos2z)^2\) by trigonometric-series
algebra.  Twenty-seven rational points give an overdetermined linear
system for the nineteen coefficients, solved exactly; consistency on
all points proves that \(F_4\) is a quartic form on the candidate
basis, and (71.1) lists the nonzero coefficients.  The single-mode
limits are Theorem \ref{thm:v19v67-quartic} with \(k^2=1,4\).
\end{proof}

\subsection{Fock diagonal and the interaction matrix}

\begin{theorem}[Sector energies to fourth order for two modes]
\label{thm:v19v68-energies}
With \(v_k=1/(2c_gk)\) and the Weyl moments (70.2) mode by mode, the
number-conserving part of (71.1) is a quadratic polynomial
\(e_4(n_1,n_2)\) with

\[
 \kappa_{11}=\kappa_{22}=-\frac{435}{2048c_g^2}+\frac{109}{512}+\frac{9c_g^2}{32},
 \qquad
 \kappa_{12}=-\frac{95}{18}v_1v_2+\frac{5}{6v_1v_2}
 =-\frac{95}{144c_g^2}+\frac{20c_g^2}{3},
                                                                    \tag{71.2}
\]

using \(v_1v_2=1/(8c_g^2)\).  For \(c_g=3/4\),
\(\kappa_{11}=\kappa_{22}=-5/768\), \(\kappa_{12}=835/324\), the
interaction matrix \(\bigl(\begin{smallmatrix}\kappa_{11}&\kappa_{12}/2\\ \kappa_{12}/2&\kappa_{22}\end{smallmatrix}\bigr)\)
has trace \(-5/384\) and determinant \(-713940175/429981696<0\), and is
indefinite.  The monomial \(x_1x_2p_1p_2\) does not contribute to the
diagonal because \(\langle n|x_kp_k|n\rangle_W=0\).
\end{theorem}

\begin{proof}
The diagonal expectation of a product monomial factorizes over the two
independent oscillators, and the \(n_1n_2\) coefficient collects the
monomials \(x_1^2x_2^2\) and \(p_1^2p_2^2\) with weights
\(4v_1v_2\) and \(4/(16v_1v_2)\); the diagonal coefficients are those of
Theorem \ref{thm:v19v67-shift}, which are \(k\)-independent because
\(k^2v_k^2=1/(4c_g^2)\).
\end{proof}

Thus at fourth order the plane-symmetric compact model has a weakly
attractive graviton self-interaction and a strongly repulsive
interaction between the two standing modes, for the reference coupling;
the sign pattern depends on \(c_g\) through (71.2), and the branch
mean curvatures of mixed sectors are shifted accordingly.

\subsection{Executable certificate}

The verifier

\[
 \texttt{scripts/verify\_sm\_v19\_v68\_two\_mode\_quartic.py}        \tag{71.3}
\]

reconstructs (71.1) exactly by interpolation, checks the single-mode
limits against V67, and computes \(\kappa_{11},\kappa_{22},\kappa_{12}\),
the trace and the determinant of the interaction matrix for
\(c_g=3/4\).  It emits

\[
 \texttt{artifacts/sm\_v19\_v68\_two\_mode\_quartic.json}.           \tag{71.4}
\]

The hashes are

\[
\begin{array}{c|l}
\hbox{object}&\hbox{SHA--256}\\ \hline
\hbox{verifier}&
\texttt{4229B4B1CD64DC02569AE7E6CE0A4C619A44591AB348749F5C91E3EF6552381F}\\
\hbox{canonical payload}&
\texttt{E504C18BDE7EF0FCBAEA9226E0216FD682164715D35E350619FA8CBF505890A5}.
\end{array}                                                     \tag{71.5}
\]

\begin{firewall}[Two modes, one polarization]
\label{fire:v19v68}
The second polarization, non-adjacent frequencies, the mean-zero
quartic components, and number-changing terms are not treated.  The
sign statements are for one value of \(c_g\).
\end{firewall}

\begin{decision}[V69 acquisition]
\label{dec:v19v68-next}
V69 will add the second polarization to the single-mode quartic
constraint, using a two-by-two matrix series for the non-diagonal
metric \(I+x\cos kz\,e^++y\cos kz\,e^\times\), and determine the
polarization-mixing quartic terms and their Fock diagonal, which
decides whether the \(+\) and \(\times\) gravitons of one mode interact.
The V70 audit follows.
\end{decision}

\begin{assessment}[V68 acquisition]
\label{ass:v19v68}
The two-mode quartic averaged constraint is an exact nine-term form,
and its Fock diagonal gives fourth-order sector energies with a
computed mode--mode interaction: for the reference coupling, equal
weakly negative self-terms and a large positive cross term, an
indefinite interaction matrix.  The first interacting structure of the
compact model is now explicit.
\end{assessment}

\begin{record}[V68 append record]
\label{rec:v19v68}
V68 is appended to the validated V67 whose installed SHA--256 was
\texttt{AA57ADE88449C5DDCB53E6E603DBB587056F5DF5DF001B64D80ED97071834265}.
After installation only V68 is the active numeric SM note; frozen
archives and unrelated notes are unchanged.
\end{record}


\section{V69: polarization mixing of one mode and the attraction of its two gravitons}
\label{sec:v19v69}

V67--V68 treated one polarization.  This note adds the second.  The
metric is then non-diagonal, and the exact scalar curvature is computed
with a two-by-two matrix series in the trigonometric algebra.  The
outcome is governed by rotation invariance in the polarization plane:
every quartic term is a function of \(|x|^2=x^2+y^2\),
\(|p|^2=p_x^2+p_y^2\), \(x\cdot p\), and the wedge \(x\wedge p=xp_y-yp_x\).
The curvature terms are those of V67 with \(x^2\) replaced by
\(|x|^2\); the only new structure comes from the momentum norm on the
perturbed background, whose fourth-order part is exactly
\(2\cos^4kz\,(x\wedge p)^2\), a term absent for parallel polarizations.
Its Fock diagonal is \(-\tfrac32n_+n_\times\) plus lower terms,
independent of the coupling: the \(+\) and \(\times\) gravitons of one
standing mode attract.

\subsection{Rotation invariance and the momentum norm}

\begin{proposition}[Two-by-two identities]
\label{prop:v19v69-identities}
For \(H=c(z)(xe^++ye^\times)\) and \(A_0=c(z)(p_xe^++p_ye^\times)\):
\(e^{+2}=e^{\times2}=I\), \(e^+e^\times=-e^\times e^+=J\) with
\(J^2=-I\), \(H^2=c^2|x|^2I\), \(\operatorname{tr}(HA_0)=2c^2\,x\cdot p\),
\(HA_0=c^2\bigl((x\cdot p)I+(x\wedge p)J\bigr)\), and
\(\det(I+H)=1-c^2|x|^2\).
\end{proposition}

\begin{proof}
Direct multiplication of the Pauli-type matrices.
\end{proof}

\begin{theorem}[Curvature and momentum norm with two polarizations]
\label{thm:v19v69-norm}
Let \(G=I+\lambda H\) and let \(A=\lambda A_0+\lambda^2A_2\) with
\(A_2=c^2(x\cdot p)I\) the \(\gamma\)-traceless completion through
second order.  Then
\begin{enumerate}
\item \(R(\gamma_\lambda)\) has vanishing orders \(0,1,3\), and its
      averages are rotation invariant:
      \(\langle R_2\rangle=-\tfrac14k^2|x|^2\),
      \(\langle R_4\rangle=-\tfrac5{16}k^2|x|^4\), with
      \(R_2=-\tfrac14k^2|x|^2-\tfrac74k^2|x|^2\cos2kz\);
\item \(|A|_\gamma^2=\operatorname{tr}((G^{-1}A)^2)\) has second-order
      part \(2c^2|p|^2\), no third-order part, and fourth-order part
      \(2c^4(x\wedge p)^2\), whose average is \(\tfrac34(x\wedge p)^2\).
\end{enumerate}
\end{theorem}

\begin{proof}
For the first item, a rotation of the polarization plane by angle
\(\theta\) about the \(z\)-axis maps \(e^+\to\cos2\theta\,e^++\sin2\theta\,e^\times\)
and is an isometry of the flat background, so every scalar built from
\(G\) is invariant under the rotation of \((x,y)\) and \((p_x,p_y)\);
a rotation-invariant quartic in \((x,y)\) alone is a multiple of
\(|x|^4\), and the multiple is fixed by the diagonal case of V66.  The
certificate confirms the averages by the matrix series at
\((x,y)=(\tfrac35,\tfrac45)\) and \((\tfrac12,\tfrac13)\).  For the second
item, with \(G^{-1}=I-\lambda H+\lambda^2H^2-\dots\), write
\(G^{-1}A=\lambda A_0+\lambda^2B_2+\lambda^3B_3+\dots\) with
\(B_2=A_2-HA_0=-c^2(x\wedge p)J\) and
\(B_3=-HA_2+H^2A_0=c^3\bigl(|x|^2A_0/c-(x\cdot p)H/c\bigr)\) by
Proposition \ref{prop:v19v69-identities}.  Then the fourth-order part of
\(\operatorname{tr}((G^{-1}A)^2)\) is
\(\operatorname{tr}(B_2^2)+2\operatorname{tr}(A_0B_3)
=-2c^4(x\wedge p)^2+2c^4\bigl(2|x|^2|p|^2-2(x\cdot p)^2\bigr)=2c^4(x\wedge p)^2\),
using \(\operatorname{tr}(A_0^2)=2c^2|p|^2\),
\(\operatorname{tr}(A_0H)=2c^2x\cdot p\), and Lagrange's identity.  The
third-order part \(2\operatorname{tr}(A_0B_2)=0\) since \(J\) is
traceless against symmetric matrices.  The average uses
\(\langle\cos^4kz\rangle=\tfrac38\).
\end{proof}

\subsection{The quartic form with two polarizations}

\begin{theorem}[Single-mode quartic form, both polarizations]
\label{thm:v19v69-form}
The fourth-order averaged constraint of one standing mode with both
polarizations is

\[
 F_4+2\Lambda_4=-\frac{145}{1024}k^2|x|^4+\frac{109}{256}|x|^2|p|^2
 +\frac{3}{16k^2}|p|^4-\frac34\,(x\wedge p)^2.                          \tag{72.1}
\]
\end{theorem}

\begin{proof}
The terms \(\langle R_4\rangle\), \(\langle R_2w_2\rangle\),
\(-4\langle(\ln\det G)'_2w_2'\rangle\), and \(12\langle w_2|A_0|^2\rangle\)
are those of Theorem \ref{thm:v19v67-quartic} with \(x^2\to|x|^2\) and
\(p^2\to|p|^2\), since \(R_2\), \(w_2\), \(\det G\), and \(|A_0|^2_\delta\)
are rotation invariant and \(w_2=(\tfrac7{128}|x|^2+\tfrac{|p|^2}{32k^2})\cos2kz\)
by (59.4) with both polarizations.  The only new term is
\(-\langle(|A|_\gamma^2)_4\rangle=-\tfrac34(x\wedge p)^2\) from Theorem
\ref{thm:v19v69-norm}.
\end{proof}

\subsection{Fock diagonal: the two gravitons of one mode attract}

\begin{theorem}[Polarization interaction]
\label{thm:v19v69-interaction}
With independent oscillators for the two polarizations, the
number-conserving part of (72.1) contains the cross term
\(\kappa_{+\times}n_+n_\times\) with

\[
 \kappa_{+\times}=-\frac{145}{512c_g^2}+\frac{109}{128}+\frac{3c_g^2}{8}-\frac32,
                                                                    \tag{72.2}
\]

where the last term is the diagonal of \(-\tfrac34(x\wedge p)^2\):
\(\langle(x\wedge p)^2\rangle_W=\tfrac12(2n_++1)(2n_\times+1)\).  For
\(c_g=3/4\), \(\kappa_{+\times}=-271/288\).  The wedge contribution
\(-\tfrac32\) is independent of the coupling and of \(k\).
\end{theorem}

\begin{proof}
\((x\wedge p)^2=x^2p_y^2+y^2p_x^2-2xp_y\,yp_x\); the mixed term has
zero Weyl diagonal, and
\(\langle x^2\rangle\langle p_y^2\rangle=(2n_++1)v\cdot(2n_\times+1)/(4v)\)
and symmetrically, giving the stated diagonal.  The other cross terms
come from \(|x|^4=x^4+2x^2y^2+y^4\), \(|x|^2|p|^2\), and \(|p|^4\)
with the moments (70.2), exactly as in Theorem \ref{thm:v19v68-energies}
with \(v_1=v_2=v\).
\end{proof}

The polarization interaction is thus attractive for the reference
coupling and remains negative for every \(c_g\) in the interval where
the curvature and momentum terms are small compared with \(3/2\).

\subsection{Executable certificate}

The verifier

\[
 \texttt{scripts/verify\_sm\_v19\_v69\_polarization\_mixing.py}      \tag{72.3}
\]

implements the two-by-two trigonometric-series engine, computes
\(R(\gamma_\lambda)\) and \(|A|^2_\gamma\) to fourth order at four
rational points including mixed polarizations, checks the vanishing
orders, the rotation-invariant averages, the exact fourth-order
momentum norm \(\tfrac34(x\wedge p)^2\), and evaluates (72.2).  It emits

\[
 \texttt{artifacts/sm\_v19\_v69\_polarization\_mixing.json}.         \tag{72.4}
\]

The hashes are

\[
\begin{array}{c|l}
\hbox{object}&\hbox{SHA--256}\\ \hline
\hbox{verifier}&
\texttt{BAFE6B54354E5C368FFE948FAB65C682DB5797AD8CD29D740016E3F7703236BA}\\
\hbox{canonical payload}&
\texttt{CDFCD0C6B06D277EA9E0C21A5A393EBBDF6E96E27247C260EB4BAD22A4D67D51}.
\end{array}                                                     \tag{72.5}
\]

\begin{firewall}[One mode, averaged constraint, Weyl ordering]
\label{fire:v19v69}
The statement is for one standing mode; the traceless completion
\(A_2\) was chosen without a second-order TT freedom; the ordering of
\(xp_yyp_x\) is Weyl.  Mean-zero quartic components and the momentum
constraint at fourth order are not treated.
\end{firewall}

\begin{decision}[V70 audit and V71 acquisition]
\label{dec:v19v69-next}
V70 is the compulsory companion audit and ten-version sweep.  V71 will
assemble the complete fourth-order sector energy of the finite-mode
plane-symmetric model with both polarizations and all pairs of modes,
using the rotation-invariant structure of this note and the two-mode
cross terms of V68, and will state the resulting quadratic form in the
occupation numbers together with its definiteness on the full mode
space, the first complete interacting spectrum of the programme.
\end{decision}

\begin{assessment}[V69 acquisition]
\label{ass:v19v69}
Polarization mixing is fixed by rotation invariance up to one new
term, the fourth-order momentum norm \(\tfrac34(x\wedge p)^2\), computed
exactly with a matrix series.  Its Fock diagonal makes the two
gravitons of one standing mode attract with a coupling-independent
strength \(3/2\), the first exact interaction constant of the compact
model.
\end{assessment}

\begin{record}[V69 append record]
\label{rec:v19v69}
V69 is appended to the validated V68 whose installed SHA--256 was
\texttt{0A6C9536879E0E38D3925D60CAF2B237C4E3E03BD4D445A0106EDC14A98436DA}.
After installation only V69 is the active numeric SM note; frozen
archives and unrelated notes are unchanged.
\end{record}


\section{V70: compulsory companion audit and ten-version sweep after the quartic sector}
\label{sec:v19v70}

This is the required audit following V65 and the ten-version sweep
after V60.  Every active companion note in the shared folder was
inspected read-only before the full fourth-order spectrum announced in
V69 is assembled.

\begin{record}[V70 companion snapshot]
\label{rec:v19v70-snapshot}
The inspected files were

\[
\begin{array}{c|r|r}
\hbox{file and SHA--256}&\hbox{bytes}&\hbox{lines}\\ \hline
\texttt{Renewal\_Yang\_Mills\_Mass\_Gap\_Research\_Notes\_V33.tex}
&472637&11954\\
\multicolumn{3}{l}{
\texttt{8DC2C41F18C64C2A87D4EA904040DD201DE109AEA75F0EB0346DFC6524F5A697}}
\\
\texttt{Renewal\_NS\_Stability\_Research\_Notes\_V26.tex}
&278283&5319\\
\multicolumn{3}{l}{
\texttt{82C887F8C2C69E4F28FAD7C3DC8DE5B9099CB5A4BF431EBA1FFF3E91935978AD}}
\\
\texttt{Renewal\_GRH\_Cofinal\_Visibility\_Attack\_V53.tex}
&517708&15814\\
\multicolumn{3}{l}{
\texttt{09E1840B82BA8433F407444217C9289EEFA00E245F9F952F8D36A791E1092C57}}
\\
\texttt{Renewal\_YM\_Parallel\_Research\_Notes\_V5.tex}
&54174&1351\\
\multicolumn{3}{l}{
\texttt{306F1BB84CFA8A8D47EA569EC17E9545F9867C8D32B548EC9D9C444B4F3C0512}}.
\end{array}                                                     \tag{73.1}
\]

Since V65, Yang--Mills has advanced from V32 to V33, with its V32 file
still present at the time of the snapshot, and GRH from V52 to V53;
the other two files are unchanged.
\end{record}

\subsection{Companion findings}

Yang--Mills V33 defines a transversely averaged dyadic blocking, proves
its Ward property and constrained level map, and shows that on the
coordinate axes the Gaussian kernel recursion has an exactly solvable
local fixed point reached geometrically at rate \(1/4\), replacing the
negative diagnosis of its V32 by a constructive Gaussian tower.  Its
firewall leaves the off-axis fixed point, one-loop coefficients, the
interacting contraction, the Haar sign, and all continuum statements
open.

GRH V53 proves a reals-only coalescence reduction and a sufficient
condition for a real-rooted diagonal in terms of a strip of zeros; its
firewall states that GRH remains open.

Navier--Stokes V26 and Yang--Mills Parallel V5 are unchanged since V55.

\subsection{Sweep of the ten versions V60--V69}

Since the V60 sweep the programme has: realized coherent states as
chains and located the classical plane on the singular sector (V61);
proved the exact quadratic \(\tau\)-dependence and the two-branch
physical space (V62); shown real plane states are branch symmetric and
withdrawn the drift-selects-branch remark (V63); consolidated V32--V63
with a mechanical hash re-verification (V64); proved that the third
order vanishes and computed the fourth-order balance
\(\Lambda_4=-145/2048\) (V66); derived the quartic averaged constraint
of one mode and its Fock diagonal (V67); the two-mode form with its
mode--mode interaction (V68); and the polarization mixing with the
coupling-independent attraction \(-3/2\) between the two gravitons of
a mode (V69).  One erratum was recorded, in V63.

\begin{decision}[V70 audit decision]
\label{dec:v19v70-audit}
No companion theorem, numerical constant, or executable artifact is
imported.  The V69 direction stands: V71 assembles the complete
fourth-order sector energy of the finite-mode plane-symmetric model.
\end{decision}

\begin{proof}
The Yang--Mills Gaussian fixed point concerns dyadic blocking of a
four-dimensional gauge kernel; the SM quartic forms concern the
averaged Hamiltonian constraint of plane-symmetric gravitational waves,
computed by exact trigonometric series with no blocking step.  The
GRH results concern real-rootedness of heat-flowed polynomials.  No
operator or hypothesis is shared; the V71 task combines V67--V69 only.
\end{proof}

\begin{assessment}[V70 acquisition]
\label{ass:v19v70}
The audit is current and the sweep finds the direction coherent.  The
compact \(\Lambda<0\) plane-symmetric programme has advanced from the
consolidated second-order theory to the first interacting structures
at fourth order, with exact interaction constants.  The open core is
unchanged in kind: the mode-number limit, the general
three-dimensional class, number-changing terms beyond the diagonal,
and the coupling of the Standard Model sectors.
\end{assessment}

\begin{record}[V70 append record]
\label{rec:v19v70}
V70 is appended to the validated V69 whose installed SHA--256 was
\texttt{78BBA1739A4E8324B2B82801E5CF38513654C7D7EAC8124DB878EED25ABA4363}.
After installation only V70 is the active numeric SM note; the
companions, frozen archives, and all unrelated notes are unchanged.  The
next compulsory companion audit is V75.
\end{record}


\section{V71: pair interactions of standing modes and the fourth-order sector energy}
\label{sec:v19v71}

V68 computed the quartic cross terms of the pair of frequencies
\((1,2)\).  This note computes them for the pairs \((1,3)\), \((2,3)\),
and \((1,4)\) by the same exact interpolation, exhibits the resonant
odd monomials that appear when one frequency is three times the other,
and assembles the fourth-order sector energy of the finite-mode
plane-symmetric model in one polarization as an explicit quadratic
form in the occupation numbers.  Its off-diagonal entries have no
uniform sign: for the reference coupling the pairs \((1,2)\), \((1,3)\),
\((2,3)\) repel and the pair \((1,4)\) attracts.  With the polarization
structure of V69 this gives the complete number-conserving quartic
Hamiltonian of the model, and shows that the interaction matrix is
indefinite.

\subsection{Pair forms}

\begin{theorem}[Quartic cross terms of a pair]
\label{thm:v19v71-pairs}
For modes \(k<l\) in one polarization, the quartic averaged constraint
contains the single-mode terms of (70.1) for each mode and the cross
terms listed below; all other monomials vanish except, for \(l=3k\),
the resonant monomials \(x_k^3x_l\), \(x_kx_lp_k^2\), \(x_k^2p_kp_l\),
\(p_k^3p_l\), whose Weyl diagonal vanishes.

\[
\begin{array}{c|c|c|c}
(k,l)&x_k^2x_l^2&p_k^2p_l^2&x_kx_lp_kp_l\\ \hline
(1,2)&-95/72&10/3&179/36\\
(1,3)&-2605/1024&15/16&1059/256\\
(2,3)&-767/200&78/25&771/100\\
(1,4)&-3859/900&34/75&3551/900
\end{array}                                                            \tag{74.1}
\]

No monomials \(x_k^2p_l^2\) or \(x_l^2p_k^2\) occur.  For \((1,3)\) the
resonant coefficients are
\(x_1^3x_3:-109/256\), \(x_1x_3p_1^2:153/128\), \(x_1^2p_1p_3:109/128\),
\(p_1^3p_3:3/4\).
\end{theorem}

\begin{proof}
For each pair the value of \(F_4+2\Lambda_4\) at a rational point is
computed as in Theorem \ref{thm:v19v68-form}, now on the full
thirty-five-monomial quartic basis without a parity restriction, since
\(z\mapsto z+\pi\) is not a symmetry when one frequency is odd and the
other is not a multiple; an overdetermined exact system fixes the
coefficients and is verified on all points.  The resonant monomials
arise from \(\langle\cos^3kz\cos lz\rangle\neq0\) exactly when
\(l=3k\).
\end{proof}

\subsection{Sector energy}

\begin{theorem}[Fourth-order sector energy, one polarization]
\label{thm:v19v71-energy}
For a finite set of standing modes in one polarization, the
number-conserving part of the fourth-order averaged constraint is

\[
 e_4(n)=\sum_k\bigl(\kappa_2n_k^2+\kappa_1n_k+\kappa_0\bigr)
 +\sum_{k<l}\kappa_{kl}\,n_kn_l+\sum_{k<l}\bigl(\hbox{terms linear in }n_k,n_l\bigr),
                                                                    \tag{74.2}
\]

with \(\kappa_2\) from (70.4) and

\[
 \kappa_{kl}=c^{xx}_{kl}\cdot4v_kv_l+c^{pp}_{kl}\cdot\frac{1}{4v_kv_l},
 \qquad v_k=\frac1{2c_gk},                                              \tag{74.3}
\]

where \(c^{xx}_{kl}\), \(c^{pp}_{kl}\) are the first two columns of
(74.1); the \(x_kx_lp_kp_l\) and resonant monomials do not contribute to
the diagonal.  For \(c_g=3/4\):
\(\kappa_{12}=835/324\), \(\kappa_{13}=515/6912\),
\(\kappa_{23}=25363/2700\), \(\kappa_{14}=-3587/4050\); the interaction
matrix on the modes \(1,2,3,4\) is indefinite.
\end{theorem}

\begin{proof}
Products over independent oscillators factorize; the \(n_kn_l\)
coefficient of \(\langle x_k^2\rangle\langle x_l^2\rangle=(2n_k+1)(2n_l+1)v_kv_l\)
is \(4v_kv_l\), and of \(\langle p_k^2\rangle\langle p_l^2\rangle\) is
\(4/(16v_kv_l)\); the mixed monomial has zero Weyl diagonal since
\(\langle n|x_kp_k|n\rangle_W=0\), and the resonant monomials contain an
odd power of a single mode's variables.  Indefiniteness follows from
the mixed signs of the off-diagonal entries together with the small
diagonal \(\kappa_2=-5/768\).
\end{proof}

Combined with Theorem \ref{thm:v19v69-interaction}, the two-polarization
model has, in addition, the coupling-independent attraction \(-\tfrac32\)
between the two gravitons of each mode and the rotation-invariant
completion of the cross terms; the resulting quadratic form in all
occupation numbers is the complete number-conserving quartic
Hamiltonian of the finite-mode plane-symmetric compact model, and its
sector labels \(\tfrac23\tau_n^2=e^{(2)}_n+e_4(n)+2\Lambda\) are the
fourth-order branch mean curvatures of the physical states.

\subsection{Executable certificate}

The verifier

\[
 \texttt{scripts/verify\_sm\_v19\_v71\_pair\_interactions.py}        \tag{74.4}
\]

reconstructs the quartic forms of the four pairs exactly, checks every
single-mode coefficient against (70.1) with the appropriate \(k^2\),
confirms the \((1,2)\) form of V68, lists all nonzero monomials
including the resonant ones for \((1,3)\), and computes (74.3) for
\(c_g=3/4\).  It emits

\[
 \texttt{artifacts/sm\_v19\_v71\_pair\_interactions.json}.           \tag{74.5}
\]

The hashes are

\[
\begin{array}{c|l}
\hbox{object}&\hbox{SHA--256}\\ \hline
\hbox{verifier}&
\texttt{73F86CFB8E48B377382568295ED656B87F11F74C51F57823168F15241719D580}\\
\hbox{canonical payload}&
\texttt{BAA8EF206A8B5097CECDAB076020ABCE4648AF5B1C51AF87A2728CD4DBC238CF}.
\end{array}                                                     \tag{74.6}
\]

\begin{firewall}[Finite pairs, diagonal part, no closed formula]
\label{fire:v19v71}
The cross coefficients are computed for four pairs and no closed
formula in \((k,l)\) is claimed; the resonant terms and all
number-changing terms are outside the diagonal statement; the
two-polarization cross terms between different modes are inferred from
rotation invariance in each polarization plane jointly and not
recomputed.  Nothing is said about the mode-number limit, where the
sums in (74.2) grow.
\end{firewall}

\begin{decision}[V72 acquisition]
\label{dec:v19v71-next}
V72 will examine the number-changing part of the quartic constraint:
the resonant monomials of the pair \((1,3)\) connect sectors with
\((n_1,n_3)\to(n_1\mp3,n_3\pm1)\), so the sector labels of (60.5) are
not exact eigenlabels beyond the diagonal approximation.  It will set
up the constraint as a matrix on the finite-mode Fock space truncated
at fixed total energy, compute its exact eigenvalues on the lowest
resonant block by rational linear algebra, and quantify the mixing of
the branch mean curvatures.
\end{decision}

\begin{assessment}[V71 acquisition]
\label{ass:v19v71}
The pair interactions of standing modes are exact, resonant odd terms
appear when one frequency is three times another, and the fourth-order
sector energy of the finite-mode plane-symmetric model in one
polarization is an explicit indefinite quadratic form in the
occupation numbers.  Together with V69 this is the complete
number-conserving quartic Hamiltonian of the compact model.
\end{assessment}

\begin{record}[V71 append record]
\label{rec:v19v71}
V71 is appended to the validated V70 whose installed SHA--256 was
\texttt{0EA5FD3E39E47CECA971B8D0B2D8B5DA3F55D7A3681BA65F42203906A843D066}.
After installation only V71 is the active numeric SM note; frozen
archives and unrelated notes are unchanged.
\end{record}


\section{V72: resonant mixing of degenerate sectors and the loss of number labels}
\label{sec:v19v72}

V71 found that the pair of frequencies \((1,3)\) carries resonant
quartic monomials that change occupation numbers.  This note evaluates
their effect on the lowest resonant block, the two free sectors with
three gravitons of frequency \(1\) and one graviton of frequency \(3\),
which are exactly degenerate at second order.  The number-changing
part connects them with an off-diagonal element whose square is
rational, the fourth-order diagonal splitting is tiny, and the exact
eigenvectors are nearly equal superpositions: the mixing satisfies
\(\sin^22\theta=10506528/10510753\).  The number labels of V52--V53
are therefore not eigenlabels of the constraint in resonant blocks; the
physical sectors, and their branch mean curvatures, are those of the
diagonalized quartic constraint.

\subsection{The resonant block}

With the Einstein normalization \(v_1=1/(2c_g)\), \(v_3=v_1/3\), the
free sectors \(|3,0\rangle\) and \(|0,1\rangle\) both have
\(e^{(2)}=3/(2c_g)\).  The quartic form of the pair \((1,3)\) has the
resonant monomials
\(c_1x_1^3x_3+c_2x_1x_3p_1^2+c_3x_1^2p_1p_3+c_4p_1^3p_3\) with
\((c_1,c_2,c_3,c_4)=(-\tfrac{109}{256},\tfrac{153}{128},\tfrac{109}{128},\tfrac34)\)
from Theorem \ref{thm:v19v71-pairs}.

\begin{theorem}[Exact two-level mixing]
\label{thm:v19v72-mixing}
On the block spanned by \(|3,0\rangle\) and \(|0,1\rangle\), the
averaged constraint operator through fourth order is the symmetric
matrix \(\bigl(\begin{smallmatrix}d_1&m\\m&d_2\end{smallmatrix}\bigr)\)
with \(d_i=e^{(2)}+e_4\) the diagonal energies of Theorem
\ref{thm:v19v71-energy} and

\[
 m=\sqrt2\Bigl[c_1v_1^2-\frac{c_2}4+\frac{3c_3}4-\frac{3c_4}{16v_1^2}\Bigr],
 \qquad
 m^2=2\Bigl[c_1v_1^2-\frac{c_2}4+\frac{3c_3}4-\frac{3c_4}{16v_1^2}\Bigr]^2.
                                                                    \tag{75.1}
\]

Its eigenvalues are \(\tfrac12(d_1+d_2)\pm\sqrt{\tfrac14(d_1-d_2)^2+m^2}\)
and the mixing angle satisfies
\(\sin^22\theta=m^2/\bigl(\tfrac14(d_1-d_2)^2+m^2\bigr)\).  For
\(c_g=3/4\): \(d_1=44789/27648\), \(d_2=14843/9216\),
\(d_1-d_2=65/6912\), \(m^2=36481/663552\), discriminant
\(10510753/191102976\), and \(\sin^22\theta=10506528/10510753\).
\end{theorem}

\begin{proof}
Only the terms lowering \(n_1\) by three and raising \(n_3\) by one
connect the two states, and in each resonant monomial that term is the
product of the pure annihilation part \(a_1^3\), with
\(\langle0|a_1^3|3\rangle=\sqrt6\), and the pure creation part
\(a_3^\dagger\); the ordering of the operators does not affect this
coefficient.  With \(x_k=\sqrt{v_k}(a_k+a_k^\dagger)\) and
\(p_k=(a_k-a_k^\dagger)/(2i\sqrt{v_k})\), the four monomials contribute
\(\sqrt6\,[c_1v_1^{3/2}v_3^{1/2}-\tfrac{c_2}4(v_3/v_1)^{1/2}+\tfrac{c_3}4(v_1/v_3)^{1/2}-\tfrac{c_4}{16}(v_1^{3/2}v_3^{1/2})^{-1}]\),
which with \(v_3=v_1/3\) is (75.1).  The diagonal entries are the
Weyl diagonals of the full pair form, in which the resonant monomials
vanish.  The certificate evaluates every quantity exactly.
\end{proof}

\begin{corollary}[Number labels are not eigenlabels in resonant blocks]
\label{cor:v19v72-labels}
The physical sectors of Theorem \ref{thm:v19v57-physical} in the
resonant block are the two eigenvectors of Theorem
\ref{thm:v19v72-mixing}, nearly equal superpositions of three
frequency-one gravitons and one frequency-three graviton, with branch
mean curvatures \(\tfrac23\tau_\pm^2=\lambda_\pm+2\Lambda\) split by
\(2\sqrt{\tfrac14(d_1-d_2)^2+m^2}\approx0.469\) for the reference
coupling.  The induced inner product of Theorem
\ref{thm:v19v52-induced} must be taken in this eigenbasis; the
number-state chains of V53 represent superpositions of the two
physical sectors.
\end{corollary}

\begin{proof}
The averaged constraint commutes with \(\hat\tau\) and is diagonal in
the eigenbasis of its \(\tau\)-independent part by Proposition
\ref{prop:v19v62-exact}; the sector labels are its eigenvalues.
\end{proof}

The near-maximal mixing is a consequence of the exact degeneracy
\(3\omega_1=\omega_3\) of standing modes on the torus, which the
quartic constraint resolves at first order in degenerate perturbation
theory.  Every triple \((k,3k)\) has such blocks, and higher resonances
\((k,l)\) with \(l/k\) odd appear at higher orders.

\subsection{Executable certificate}

The verifier

\[
 \texttt{scripts/verify\_sm\_v19\_v72\_resonant\_mixing.py}          \tag{75.2}
\]

imports the unchanged V71 pair reconstruction, checks the resonant
coefficients, computes the diagonal energies, the bracket
\(-191/1152\), \(m^2\), the splitting, the discriminant, and
\(\sin^22\theta\), all exactly for \(c_g=3/4\).  It emits

\[
 \texttt{artifacts/sm\_v19\_v72\_resonant\_mixing.json}.             \tag{75.3}
\]

The hashes are

\[
\begin{array}{c|l}
\hbox{object}&\hbox{SHA--256}\\ \hline
\hbox{verifier}&
\texttt{3B2729D4C595B02881A1F9C01E67150038598F18B2E7DA443CCA835E370BBB74}\\
\hbox{canonical payload}&
\texttt{6F9800B496F8695E4330E135B58DC90F0C0BADF1D77FC2BF037420130FF6B68F}.
\end{array}                                                     \tag{75.4}
\]

\begin{firewall}[One block, fourth order, Weyl diagonals]
\label{fire:v19v72}
Only the lowest resonant block of one pair is diagonalized, and only
through fourth order; the diagonal entries use Weyl ordering, whose
ambiguity shifts both entries by constants and leaves the mixing
essentially unchanged; higher blocks, other pairs, both polarizations,
and the mode-number limit are not treated.
\end{firewall}

\begin{decision}[V73 acquisition]
\label{dec:v19v72-next}
V73 will go upstream of the resonance: on the torus every frequency
ratio is rational, so resonant blocks are dense in the free spectrum,
and the diagonal approximation of V67--V71 is never uniformly valid.
It will state the resonance structure of the finite-mode model exactly,
prove that the quartic constraint restricted to any fixed free-energy
shell is a finite symmetric matrix whose eigenvalues are the sector
labels, and bound the size of that matrix as a function of the shell
energy and the mode cutoff, which decides whether a mode-number limit of
the sector spectrum can be posed at all.
\end{decision}

\begin{assessment}[V72 acquisition]
\label{ass:v19v72}
Number-changing quartic terms mix degenerate free sectors almost
maximally in the lowest resonant block, with an exactly computed
splitting.  The physical sectors and their branch mean curvatures are
those of the diagonalized constraint, not of the graviton numbers,
which corrects the labelling used since V52 in resonant blocks and
points upstream to the resonance structure of the torus spectrum.
\end{assessment}

\begin{record}[V72 append record]
\label{rec:v19v72}
V72 is appended to the validated V71 whose installed SHA--256 was
\texttt{EC21B7ED47C2EF0F37BE933891C5F650564994C5B6E550D46BFCE881AAEF87B4}.
After installation only V72 is the active numeric SM note; frozen
archives and unrelated notes are unchanged.
\end{record}


\section{V73: free-energy shells, parity blocks, and what a mode-number limit of the sector spectrum requires}
\label{sec:v19v73}

V72 showed that resonant quartic terms mix degenerate free sectors.
This note states the structure exactly.  On the torus every frequency
ratio is rational, so the free spectrum is organized in integer
shells, and the quartic constraint preserves each shell: on a shell it
is a finite symmetric matrix whose eigenvalues are the sector labels.
The only exact selection rule is the parity of the total graviton
number, which splits every shell into two blocks; within a block the
energy-preserving quartic moves connect every state once the shell
energy exceeds three.  The shell basis is independent of the mode
cutoff once the cutoff exceeds the shell energy, so the sector spectrum
at fixed shell energy has a trivially attained mode-number limit for
its shell-preserving part; what does not saturate are the c-number
zero point, removed by normal ordering, and the linear self-energy
sums over all higher modes, whose convergence is not established.

\subsection{Shells and the shell matrix}

Let the modes be \(k=1,\dots,K\) with two polarizations, and for an
occupation vector \(n=(n_{k,A})\) let \(E(n)=\sum_kk(n_{k,+}+n_{k,\times})\).

\begin{theorem}[Shell structure of the quartic constraint]
\label{thm:v19v73-shell}
\begin{enumerate}
\item The number-changing part of the quartic averaged constraint
      consists of monomials of the types
      \(a^\dagger_{k_1}a^\dagger_{k_2}a_{k_3}a_{k_4}\) with
      \(k_1+k_2=k_3+k_4\), and
      \(a^\dagger_{k_1}a^\dagger_{k_2}a^\dagger_{k_3}a_{k_4}\) with
      \(k_1+k_2+k_3=k_4\), together with their adjoints, up to
      polarization labels; every such monomial preserves \(E\).
\item The shell \({\cal S}_E\) of occupation vectors with \(E(n)=E\) is
      finite, of dimension equal to the number of two-coloured
      partitions of \(E\) with parts at most \(K\), and independent of
      \(K\) for \(K\geq E\).
\item The averaged constraint through fourth order is, on
      \({\cal S}_E\), a finite real symmetric matrix commuting with
      \(\hat\tau\); its eigenvalues are the sector labels of the
      physical-state condition (60.5), and the total graviton number
      modulo two is conserved, so the matrix is block diagonal with at
      least two blocks.
\end{enumerate}
\end{theorem}

\begin{proof}
The quartic form is a sum of products of four mode operators; in the
Fock representation a monomial changes \(E\) by \(\sum(\hbox{created})-\sum(\hbox{annihilated})\)
frequencies, and the shell-preserving ones are those listed; energy
conservation of the resonant monomials of Theorem
\ref{thm:v19v71-pairs} is the case \((k,k,k)\to(3k)\).  Terms with an
odd number of creation plus annihilation operators do not occur in a
quartic form, and the listed types change the total number by \(0\) or
\(\pm2\), so number parity is conserved.  The shell count is the
generating function \(\prod_{k\leq K}(1-q^k)^{-2}\), and parts larger
than \(E\) cannot occur.  Symmetry of the matrix follows from the
Hermiticity of the Weyl-ordered form.
\end{proof}

\subsection{Connectivity and growth}

\begin{proposition}[Connectivity within parity classes]
\label{prop:v19v73-connectivity}
Let the shell graph have an edge between two occupation vectors when
some listed move, with any polarization assignment, maps one to the
other.  Then for \(K\geq3\) and \(E\geq4\) the graph has exactly two
connected components, the two parity classes.  The dimensions and
largest components are

\[
\begin{array}{c|c|c|c|c|c}
K&E=2&E=3&E=4&E=6&E=8\\ \hline
2&5/1&8/1&14/1&30/1&55/1\\
3&5/1&10/6&18/12&49/27&110/62\\
4&5/1&10/6&20/12&59/31&149/77\\
6&5/1&10/6&20/12&65/35&179/93
\end{array}                                                            \tag{76.1}
\]

with entries dimension/largest component.  For \(K=2\) no move
exists and every state is its own block.
\end{proposition}

\begin{proof}
Exact enumeration by the certificate.  With \(K=2\) the only candidate
move \((1,1)\to(2)\) does not preserve energy and \((1,1,1)\to(3)\) is
unavailable.
\end{proof}

Since the true quartic form restricts the polarization assignments,
the true shell graph is a subgraph of the one counted; the parity
blocks are therefore a lower bound on the block decomposition, and the
counted components are an upper bound on connectivity.  The dimension
grows with \(E\) like the two-coloured partition function, which is
subexponential in \(E\) but unbounded, and saturates in \(K\) at
\(K=E\).

\subsection{What a mode-number limit requires}

\begin{theorem}[Mode-number limit at fixed shell energy]
\label{thm:v19v73-limit}
Fix a shell energy \(E\) and let the cutoff \(K\to\infty\).  The
shell-preserving part of the quartic constraint on \({\cal S}_E\) is
independent of \(K\) for \(K\geq E\).  The diagonal part depends on
\(K\) through two structures:
\begin{enumerate}
\item the c-number zero point, the sum over all modes of the vacuum
      values of the single-mode quartic terms together with
      \(Z\) of (55.2), which diverges with \(K\) and is removed by
      normal ordering;
\item the self-energy of each occupied mode \(k\), linear in \(n_k\),
      \(\sigma_k(K)=\sum_{l\leq K,\,l\neq k}\bigl(2c^{xx}_{kl}v_kv_l+c^{pp}_{kl}/(8v_kv_l)\bigr)\),
      arising from the vacuum fluctuations of every other mode through
      the pair cross terms of Theorem \ref{thm:v19v71-pairs}.
\end{enumerate}
The sector spectrum on \({\cal S}_E\) has a mode-number limit if and
only if the self-energies \(\sigma_k(K)\), \(k\leq E\), converge or are
renormalized by a \(k\)-dependent counterterm; this is not decided
here, since \(c^{xx}_{kl}\), \(c^{pp}_{kl}\) are known for four pairs
only.
\end{theorem}

\begin{proof}
Shell preservation and the \(K\)-independence of the basis are Theorem
\ref{thm:v19v73-shell}.  The diagonal of a pair form on a state with
\(n_l=0\) contributes \(c^{xx}_{kl}(2n_k+1)v_kv_l+c^{pp}_{kl}(2n_k+1)/(16v_kv_l)\),
whose \(n_k\)-linear part is the stated summand; the constant parts
are the zero point.  Nothing else in the fourth-order diagonal depends
on modes outside the shell.
\end{proof}

The structure is that of a mass renormalization: the zero point is
the vacuum energy, \(\sigma_k(K)\) is the one-loop self-energy of a
graviton mode from all other modes, and the finite shell matrix is the
interaction.  The programme's mode-number problem is thereby reduced,
for the averaged constraint of the plane-symmetric model, to the
asymptotics of the pair coefficients \(c_{kl}\) in \(l\), which V74 can
attack with the interpolation engine at larger \(l\).

\subsection{Executable certificate}

The verifier

\[
 \texttt{scripts/verify\_sm\_v19\_v73\_shell\_resonance.py}          \tag{76.2}
\]

enumerates the shells, the energy-preserving moves for \(K=3\), namely
\((1,1,1)\leftrightarrow(3)\) and \((1,3)\leftrightarrow(2,2)\), and the
connected components with all polarization assignments, producing
(76.1) and checking the growth in \(E\) and \(K\) and the presence of
the V72 block.  It emits

\[
 \texttt{artifacts/sm\_v19\_v73\_shell\_resonance.json}.             \tag{76.3}
\]

The hashes are

\[
\begin{array}{c|l}
\hbox{object}&\hbox{SHA--256}\\ \hline
\hbox{verifier}&
\texttt{508ED38B1D5C22136E60944FD9E15248F6678819C89514BD42708E2AB8C454C1}\\
\hbox{canonical payload}&
\texttt{703B9D30E756F73684C4B3D097AC175626EC564CDC9C48E7718624D6C02674AD}.
\end{array}                                                     \tag{76.4}
\]

\begin{firewall}[Counting, not diagonalizing]
\label{fire:v19v73}
The shell matrices are not computed beyond the V72 block, the
connectivity count uses unrestricted polarization assignments, and the
convergence of the self-energies is left open.  Fourth order, plane
symmetry, and the averaged constraint remain the setting.
\end{firewall}

\begin{decision}[V74 acquisition]
\label{dec:v19v73-next}
V74 will compute the pair coefficients \(c^{xx}_{1l}\), \(c^{pp}_{1l}\)
for \(l\) up to eight with the interpolation engine, determine their
growth in \(l\), and hence the behaviour of the self-energy sum
\(\sigma_1(K)\) with the Einstein normalization; if it diverges, state
the counterterm structure needed for a mode-number limit of the
sector spectrum.  The V75 audit follows.
\end{decision}

\begin{assessment}[V73 acquisition]
\label{ass:v19v73}
The quartic constraint of the finite-mode plane-symmetric model is a
finite symmetric matrix on each integer free-energy shell, block
diagonal by total-number parity, with a shell basis that saturates at
cutoff equal to the shell energy.  Its mode-number limit reduces to the
zero point, removed by normal ordering, and to the graviton
self-energy sums over all higher modes, which are the next object.
\end{assessment}

\begin{record}[V73 append record]
\label{rec:v19v73}
V73 is appended to the validated V72 whose installed SHA--256 was
\texttt{51F923E223543B009E4502E17ACC9F17D29A9794B51022BBE892578C5CE18B3F}.
After installation only V73 is the active numeric SM note; frozen
archives and unrelated notes are unchanged.
\end{record}


\section{V74: growth of the pair coefficients and the divergent graviton self-energy}
\label{sec:v19v74}

V73 reduced the mode-number limit of the sector spectrum to the
self-energy sums \(\sigma_k(K)\).  This note computes the pair
coefficients \(c^{xx}_{1l}\), \(c^{pp}_{1l}\) exactly for \(l\) up to
eight and finds that \(c^{xx}_{1l}/l^2\) approaches \(-1/4\) while
\(c^{pp}_{1l}l^2\) approaches about six.  With the Einstein
normalization the self-energy summand of mode \(1\) from mode \(l\)
therefore grows linearly in \(l\), and \(\sigma_1(K)\) diverges
quadratically.  Since a term linear in \(n_k\) is a shift of the mode
frequency, the divergence is a frequency renormalization: with
\(k\)-dependent frequency counterterms the finite shell matrices of
V73 have a cutoff-independent limit, while without them the sector
labels of every occupied mode run off with the cutoff.  This is the
exact form, in the plane-symmetric model, of the ultraviolet
renormalization that the continuum programme must perform.

\subsection{Pair coefficients}

\begin{proposition}[Exact pair coefficients \((1,l)\)]
\label{prop:v19v74-pairs}
The coefficients of \(x_1^2x_l^2\) and \(p_1^2p_l^2\) in the quartic
averaged constraint of the pair \((1,l)\) are

\[
\begin{array}{c|c|c|c|c}
l&c^{xx}_{1l}&c^{xx}_{1l}/l^2&c^{pp}_{1l}&c^{pp}_{1l}\,l^2\\ \hline
2&-95/72&-95/288&10/3&40/3\\
3&-2605/1024&-2605/9216&15/16&135/16\\
4&-3859/900&-3859/14400&34/75&544/75\\
5&-60229/9216&-60229/230400&13/48&325/48\\
6&-90983/9800&-90983/352800&222/1225&7992/1225\\
7&-462025/36864&-462025/1806336&25/192&1225/192\\
8&-258505/15876&-258505/1016064&130/1323&8320/1323
\end{array}                                                            \tag{77.1}
\]

The ratios \(c^{xx}_{1l}/l^2\) decrease in modulus toward
\(-0.254\) at \(l=8\), and \(c^{pp}_{1l}l^2\) decreases toward
\(6.29\).
\end{proposition}

\begin{proof}
Exact interpolation as in Theorem \ref{thm:v19v71-pairs} for each
pair; the values are certified.
\end{proof}

The observed behaviour \(c^{xx}_{1l}\approx-l^2/4\) and
\(c^{pp}_{1l}\approx6/l^2\) is not proved for all \(l\); it is
recorded as computed through \(l=8\).

\subsection{The self-energy sum}

\begin{theorem}[Self-energy of mode \(1\)]
\label{thm:v19v74-self-energy}
With \(v_l=1/(2c_gl)\), the linear-in-\(n_1\) contribution of the
vacuum of mode \(l\) to the sector energy is
\(\sigma_{1l}=2\bigl(c^{xx}_{1l}v_1v_l+c^{pp}_{1l}/(16v_1v_l)\bigr)\),
and for \(c_g=3/4\) the partial sums are

\[
 \sigma_1(K)=\sum_{l=2}^K\sigma_{1l}:\quad
 \tfrac{835}{648},\ \tfrac{54985}{41472},\ \tfrac{33907}{38400},\ \tfrac{13223}{129600},\
 -\tfrac{1228877}{1270080},\ -\tfrac{1155419}{501760},\ -\tfrac{203314693}{52254720}
                                                                    \tag{77.2}
\]

for \(K=2,\dots,8\).  If \(c^{xx}_{1l}=-l^2/4+O(l)\) and
\(c^{pp}_{1l}=O(l^{-2})\), then
\(\sigma_{1l}=-l/(8c_g^2)+O(1)\) and \(\sigma_1(K)=-K^2/(16c_g^2)+O(K)\).
\end{theorem}

\begin{proof}
The summand is the \(n_1\)-linear part of
\(c^{xx}_{1l}(2n_1+1)v_1v_l+c^{pp}_{1l}(2n_1+1)/(16v_1v_l)\) at
\(n_l=0\), by Theorem \ref{thm:v19v73-limit}.  With
\(v_1v_l=1/(4c_g^2l)\) the first term is \(2c^{xx}_{1l}/(4c_g^2l)\), which
is \(-l/(8c_g^2)+O(1)\) under the stated asymptotics, and the second is
\(c^{pp}_{1l}\cdot4c_g^2l/8=O(l^{-1})\); summing gives the quadratic
divergence.  The partial sums are certified.
\end{proof}

\subsection{Frequency renormalization}

\begin{proposition}[Linear terms are frequency shifts]
\label{prop:v19v74-frequency}
A term \(\sigma_kn_k\) in the sector energy is equivalent to replacing
the free frequency in the second-order label
\(e^{(2)}_n=\sum_kn_k\omega_k/(2c_g)\) by \(\omega_k+2c_g\sigma_k\).
Consequently the shell matrices of Theorem \ref{thm:v19v73-shell} have
a cutoff-independent limit if and only if the frequencies are
renormalized by counterterms \(\delta\omega_k(K)=-2c_g\sigma_k(K)+\hbox{finite}\),
one per mode, in addition to the normal ordering of the zero point.
The shell-preserving part, which fixes the mixing of V72, needs no
counterterm.
\end{proposition}

\begin{proof}
Immediate from the form of \(e^{(2)}\) and Theorem
\ref{thm:v19v73-limit}.
\end{proof}

Thus the mode-number limit of the plane-symmetric sector spectrum
requires exactly two renormalizations at fourth order, a vacuum energy
and a frequency shift per mode, both of the ultraviolet type familiar
from perturbative field theory, and no others at this order.  Whether
the renormalized frequencies can be chosen consistently with the
Einstein normalization of the linearized sector, on which the exact
spectral consistency of V44 rests, is the question that the
counterterms raise and that this note does not answer.

\subsection{Executable certificate}

The verifier

\[
 \texttt{scripts/verify\_sm\_v19\_v74\_self\_energy\_growth.py}      \tag{77.3}
\]

imports the unchanged V71 reconstruction and computes (77.1) and
(77.2) exactly, checking the pair \((1,2)\) against V68.  It emits

\[
 \texttt{artifacts/sm\_v19\_v74\_self\_energy\_growth.json}.         \tag{77.4}
\]

The hashes are

\[
\begin{array}{c|l}
\hbox{object}&\hbox{SHA--256}\\ \hline
\hbox{verifier}&
\texttt{FCD87855183FE22BE88567A117C657266D808EBCA4C9FA64B4B50FA878C75932}\\
\hbox{canonical payload}&
\texttt{8A7A57A79DAA9529147EEBFB22D666C5A329FCC94D159BC8350B00C26D5D8196}.
\end{array}                                                     \tag{77.5}
\]

\begin{firewall}[Asymptotics observed, not proved]
\label{fire:v19v74}
The growth laws of the pair coefficients are observed through
\(l=8\) and not proved; the self-energy divergence is conditional on
them.  The counterterm statement is structural and does not construct
the renormalized model or relate the counterterms to the Einstein
normalization.
\end{firewall}

\begin{decision}[V75 audit and V76 acquisition]
\label{dec:v19v74-next}
V75 is the compulsory companion audit.  V76 will attempt to prove the
asymptotics of \(c^{xx}_{1l}\) by deriving the pair form analytically
from (59.2) and the second-order responses: the cross term arises from
\(\langle R_4\rangle\), whose two-mode part can be reduced by the
two-by-two identities to averages of products of four cosines with
explicit frequencies, so that \(c^{xx}_{1l}\) is a closed rational
function of \(l\); this would make Theorem \ref{thm:v19v74-self-energy}
unconditional.
\end{decision}

\begin{assessment}[V74 acquisition]
\label{ass:v19v74}
The pair coefficients grow like \(l^2\) in the configuration sector and
decay like \(l^{-2}\) in the momentum sector, so the graviton
self-energy from higher modes diverges quadratically with the cutoff,
conditionally on the observed asymptotics.  The divergence is a
frequency renormalization, one counterterm per mode, which together
with normal ordering is all that the fourth-order sector spectrum
needs for a mode-number limit; the compatibility of those counterterms
with the Einstein normalization is the next upstream question.
\end{assessment}

\begin{record}[V74 append record]
\label{rec:v19v74}
V74 is appended to the validated V73 whose installed SHA--256 was
\texttt{38435EBC73FC32ABFD6A9CB289E3CE636FD8A8D53D4D4CB3DA03C9BAFAFB61F2}.
After installation only V74 is the active numeric SM note; frozen
archives and unrelated notes are unchanged.
\end{record}


\section{V75: compulsory companion audit after the self-energy divergence}
\label{sec:v19v75}

This is the required audit following V70.  Every active companion note
in the shared folder was inspected read-only before the analytic pair
asymptotics announced in V74 are attempted.

\begin{record}[V75 companion snapshot]
\label{rec:v19v75-snapshot}
The inspected files were

\[
\begin{array}{c|r|r}
\hbox{file and SHA--256}&\hbox{bytes}&\hbox{lines}\\ \hline
\texttt{Renewal\_Yang\_Mills\_Mass\_Gap\_Research\_Notes\_V33.tex}
&473284&11963\\
\multicolumn{3}{l}{
\texttt{B417E58E5CF6BFBE9F3FB0DD21106355674D525AE9BF8464EFD7AA580D0134F9}}
\\
\texttt{Renewal\_NS\_Stability\_Research\_Notes\_V26.tex}
&278283&5319\\
\multicolumn{3}{l}{
\texttt{82C887F8C2C69E4F28FAD7C3DC8DE5B9099CB5A4BF431EBA1FFF3E91935978AD}}
\\
\texttt{Renewal\_GRH\_Cofinal\_Visibility\_Attack\_V53.tex}
&517708&15814\\
\multicolumn{3}{l}{
\texttt{09E1840B82BA8433F407444217C9289EEFA00E245F9F952F8D36A791E1092C57}}
\\
\texttt{Renewal\_YM\_Parallel\_Research\_Notes\_V5.tex}
&54174&1351\\
\multicolumn{3}{l}{
\texttt{306F1BB84CFA8A8D47EA569EC17E9545F9867C8D32B548EC9D9C444B4F3C0512}}.
\end{array}                                                     \tag{78.1}
\]

Since V70, the Yang--Mills V33 file has been edited in place by nine
lines and \(647\) bytes with a new hash, its V32 predecessor has been
removed, and the other three files are unchanged.
\end{record}

\begin{decision}[V75 audit decision]
\label{dec:v19v75-audit}
No companion theorem, numerical constant, or executable artifact is
imported.  The V74 direction stands: V76 derives the pair asymptotics
analytically.
\end{decision}

\begin{proof}
The in-place edit of Yang--Mills V33 does not change its section list
or its firewall, which concerns Gaussian dyadic blocking; the SM task is
the analytic reduction of a two-mode average of products of cosines,
internal to V56--V74.
\end{proof}

\begin{assessment}[V75 acquisition]
\label{ass:v19v75}
The audit is current.  Since V70 the programme has assembled the pair
interactions and the fourth-order sector energy, diagonalized the
lowest resonant block exactly, established the shell and parity
structure of the quartic constraint with its cutoff saturation, and
identified the quadratically divergent graviton self-energy as a
frequency renormalization, one counterterm per mode.
\end{assessment}

\begin{record}[V75 append record]
\label{rec:v19v75}
V75 is appended to the validated V74 whose installed SHA--256 was
\texttt{B0F585B26FF82A202F63FDC7DA4B6C74C1513005902B7FCEE97571F3EF9C5F30}.
After installation only V75 is the active numeric SM note; the
companions, frozen archives, and all unrelated notes are unchanged.  The
next compulsory companion audit is V80.
\end{record}


\section{V76: closed forms of the pair coefficients and the unconditional self-energy divergence}
\label{sec:v19v76}

V74 observed the growth of the pair coefficients through \(l=8\).  This
note derives them in closed form.  For the diagonal plane-symmetric
metric the scalar curvature is an explicit rational expression in the
wave profile and its first two derivatives, so its second- and
fourth-order parts are polynomial functionals, and the two-mode
averages reduce to products of four trigonometric functions.  The
configuration coefficient is a rational function of \(l\) with leading
term \(-l^2/4\), the momentum coefficient is
\(3[(l-1)^{-2}+(l+1)^{-2}]\), and the quadratic divergence of the
graviton self-energy of V74 becomes a theorem, with a logarithmic
subleading term from the momentum sector.

\subsection{Closed form of the scalar curvature}

\begin{proposition}[Curvature of the diagonal profile]
\label{prop:v19v76-curvature}
For \(g=dz^2+(1+f)dx^2+(1-f)dy^2\) with \(|f|<1\),

\[
 R=\frac{2ff''}{1-f^2}+\frac{f'^2(3+f^2)}{2(1-f^2)^2},
 \qquad
 R_2=2ff''+\tfrac32f'^2,\qquad
 R_4=2f^3f''+\tfrac72f^2f'^2,                                          \tag{79.1}
\]

and \(\langle R_4\rangle=-\tfrac52\langle f^2f'^2\rangle\) on the circle.
\end{proposition}

\begin{proof}
With \(a=\sqrt{1+f}\), \(b=\sqrt{1-f}\):
\(a''/a=f''/(2(1+f))-f'^2/(4(1+f)^2)\),
\(b''/b=-f''/(2(1-f))-f'^2/(4(1-f)^2)\),
\(a'b'/(ab)=-f'^2/(4(1-f^2))\); summing and using
\(\tfrac1{(1+f)^2}+\tfrac1{(1-f)^2}+\tfrac1{1-f^2}=\tfrac{3+f^2}{(1-f^2)^2}\)
in \(R=-2(a''/a+b''/b+a'b'/(ab))\) gives the first formula, whose
expansion gives \(R_2\) and \(R_4\).  Integrating \(f^3f''\) by parts
gives \(-3\langle f^2f'^2\rangle\).  The certificate checks (79.1)
against the series engine for a two-mode profile.
\end{proof}

\subsection{Closed forms of the pair coefficients}

For \(f=x_1\cos z+x_l\cos lz\), \(l\geq2\), put

\[
 A_\mp:=\frac{2l^2\mp3l+2}{16(l\mp1)^2}.                               \tag{79.2}
\]

\begin{theorem}[Pair coefficients in closed form]
\label{thm:v19v76-closed}
The coefficients of \(x_1^2x_l^2\) and \(p_1^2p_l^2\) in the quartic
averaged constraint of the pair \((1,l)\) are

\[
 c^{xx}_{1l}=-\frac58(l^2+1)+12\bigl[(l-1)^2A_-^2+(l+1)^2A_+^2\bigr]-l(A_--A_+),
 \qquad
 c^{pp}_{1l}=3\Bigl[\frac1{(l-1)^2}+\frac1{(l+1)^2}\Bigr],               \tag{79.3}
\]

and consequently

\[
 c^{xx}_{1l}=-\frac{l^2}4+O(l),\qquad
 c^{pp}_{1l}=\frac{6(l^2+1)}{(l^2-1)^2}=\frac6{l^2}+O(l^{-4}).           \tag{79.4}
\]
\end{theorem}

\begin{proof}
At \(p=0\), (70.1) reads
\(F_4+2\Lambda_4=\langle R_4\rangle+\langle R_2w_2\rangle-4\langle(\ln\det G)'_2w_2'\rangle\)
with \(w_2''=(R_2-\langle R_2\rangle)/8\) and \((\ln\det G)_2=-f^2\).
Integrating by parts, \(\langle R_2w_2\rangle=-8\langle w_2'^2\rangle\)
and \(-4\langle(\ln\det G)'_2w_2'\rangle=-8\langle(ff')'w_2\rangle\);
since \((ff')'=\tfrac12R_2+\tfrac14f'^2\), the total is
\(-\tfrac52\langle f^2f'^2\rangle+24\langle w_2'^2\rangle-2\langle f'^2w_2\rangle\).
The \(x_1^2x_l^2\) coefficient of \(\langle f^2f'^2\rangle\) is
\((l^2+1)/4\).  The cross part of \(R_2\) is
\(x_1x_l[(\tfrac{3l}2-l^2-1)\cos(l-1)z-(l^2+1+\tfrac{3l}2)\cos(l+1)z]\),
so the cross part of \(w_2\) is \(x_1x_l[A_-\cos(l-1)z+A_+\cos(l+1)z]\),
which gives the \(x_1^2x_l^2\) parts \(12[(l-1)^2A_-^2+(l+1)^2A_+^2]\)
of \(24\langle w_2'^2\rangle\) and \(-l(A_--A_+)\) of
\(-2\langle f'^2w_2\rangle\), using
\(f'^2\ni lx_1x_l[\cos(l-1)z-\cos(l+1)z]\).  For \(c^{pp}\), the only
source is \(12\langle w_2|A_0|^2\rangle\) with the momentum part of
\(w_2\), namely \(2p_1p_l[\cos(l-1)z/(8(l-1)^2)+\cos(l+1)z/(8(l+1)^2)]\),
against \(2p_1p_l[\cos(l-1)z+\cos(l+1)z]\) in \(|A_0|^2\), giving
\(12\cdot\tfrac14[(l-1)^{-2}+(l+1)^{-2}]\).  The leading terms follow
from \(A_\mp=\tfrac18+O(l^{-1})\); the certificate confirms (79.3)
against the interpolated values for \(2\leq l\leq8\).
\end{proof}

\subsection{The self-energy divergence is unconditional}

\begin{theorem}[Graviton self-energy of mode \(1\)]
\label{thm:v19v76-self-energy}
With the Einstein normalization,

\[
 \sigma_1(K)=\sum_{l=2}^K\Bigl[\frac{c^{xx}_{1l}}{2c_g^2l}+\frac{c^{pp}_{1l}c_g^2l}2\Bigr]
 =-\frac{K^2}{16c_g^2}+O(K)+3c_g^2\log K+O(1),                            \tag{79.5}
\]

so the fourth-order sector label of every state containing a
frequency-one graviton diverges quadratically with the mode cutoff,
and the frequency counterterm of Proposition
\ref{prop:v19v74-frequency} is required, with a subleading
logarithmic part from the momentum sector.  For \(c_g=3/4\) the
partial sums divided by \(K^2\) approach \(-1/9\).
\end{theorem}

\begin{proof}
Insert (79.4) into the summand of Theorem
\ref{thm:v19v74-self-energy} and sum; the certificate evaluates the
exact partial sums through \(K=200\).
\end{proof}

The result is exact and independent of any numerical observation.
It states, within the plane-symmetric fourth-order model, that the
compact torus gravitons acquire a cutoff-dependent frequency shift
quadratic in the cutoff, the lattice form of the ultraviolet
self-energy of perturbative gravity; the Einstein normalization of the
linear sector, which fixes \(v_l=1/(2c_gl)\), is what turns the
\(l^2\) growth of \(c^{xx}_{1l}\) into a linear summand.

\subsection{Executable certificate}

The verifier

\[
 \texttt{scripts/verify\_sm\_v19\_v76\_pair\_closed\_forms.py}       \tag{79.6}
\]

checks (79.1) against the series engine for a two-mode profile,
including the integration-by-parts identity, checks (79.3) against
the interpolated coefficients for \(2\leq l\leq8\), verifies the
asymptotics (79.4) at \(l\) up to \(200\), and evaluates the exact
partial sums (79.5) through \(K=200\) with \(\sigma_1(K)/K^2\)
approaching \(-1/9\).  It emits

\[
 \texttt{artifacts/sm\_v19\_v76\_pair\_closed\_forms.json}.          \tag{79.7}
\]

The hashes are

\[
\begin{array}{c|l}
\hbox{object}&\hbox{SHA--256}\\ \hline
\hbox{verifier}&
\texttt{E27392E31CC97ABAB533864BB81323E748404DCC6100CB2D1CC37DF49799436D}\\
\hbox{canonical payload}&
\texttt{98CAF72F9A3CC1918E615DB3865A2F30B7B77A5001B7E09567C6EB6E8F5F4591}.
\end{array}                                                     \tag{79.8}
\]

\begin{firewall}[Mode one, one polarization, fourth order]
\label{fire:v19v76}
The closed forms are for pairs \((1,l)\); general pairs \((k,l)\) follow
by the same method but are not written out, and the second
polarization adds the rotation-invariant completion of V69.  The
divergence is a statement about the averaged constraint's diagonal in
the finite-mode plane-symmetric model; whether the counterterms can
be chosen compatibly with the Einstein normalization remains the open
upstream question.
\end{firewall}

\begin{decision}[V77 acquisition]
\label{dec:v19v76-next}
V77 will confront the counterterm with the normalization: the
frequency shift \(\delta\omega_1(K)=-2c_g\sigma_1(K)\) changes the
stationary variance of the seed chain, hence the reflected two-point
function of V44, unless the seed normalization \(v_1\) is rescaled in
step; it will determine whether a \(K\)-dependent rescaling of the
seed variances keeps both the linear spectral consistency and the
fourth-order sector spectrum finite, which is the lattice form of
wave-function renormalization, and record the resulting running of
the Einstein coupling \(c_g\) with the cutoff.
\end{decision}

\begin{assessment}[V76 acquisition]
\label{ass:v19v76}
The pair coefficients of the plane-symmetric quartic constraint are
explicit rational functions, the scalar curvature of the diagonal
profile is a closed rational expression, and the quadratic divergence
of the graviton self-energy with the mode cutoff is now a theorem,
with an exact \(-K^2/(16c_g^2)\) leading term and a logarithmic
momentum-sector correction.  The programme has reached the ultraviolet
renormalization of the compact model in exact form.
\end{assessment}

\begin{record}[V76 append record]
\label{rec:v19v76}
V76 is appended to the validated V75 whose installed SHA--256 was
\texttt{68D3AB89EE7436F5BB16DFCF8AE778B453E1A5E82DB52FE201D071F946116BCA}.
After installation only V76 is the active numeric SM note; frozen
archives and unrelated notes are unchanged.
\end{record}


\section{V77: the counterterm against the Einstein normalization, and the intrinsic cutoff of the bare model}
\label{sec:v19v77}

V76 made the graviton self-energy divergence exact.  This note asks
what it does to the seed.  The fourth-order sector label of a
one-graviton state is \(n_1E_1(K)\) with
\(E_1(K)=1/(2c_g)+\sigma_1(K)\), and consistency of the seed with the
constraint, in which the Hamiltonian is the constraint, requires the
seed kernel ratio \(e^{-\varepsilon\omega^{\rm eff}}\) with
\(\omega^{\rm eff}=2c_gE_1(K)>0\).  The exact partial sums show that
\(E_1(K)\) first rises, driven by the positive momentum-sector
contribution, and then falls quadratically, changing sign at a finite
cutoff \(K_*(c_g)\): six for the reference coupling, twelve for
\(c_g=1\), fifty-nine for \(c_g=2\), and beyond four hundred for
\(c_g=5\).  Beyond \(K_*\) no stationary reflection-positive seed is
consistent with the bare fourth-order constraint.  The on-shell
counterterm \(\delta\omega_1(K)=-2c_g\sigma_1(K)\) removes the running
entirely, keeps the seed at its Einstein normalization for every
cutoff, and moves all cutoff dependence into the counterterm, which
grows like \(K^2/(8c_g)\).  The bare plane-symmetric model therefore has
an intrinsic cutoff \(K_*(c_g)\), growing roughly like
\(c_g^2\sqrt{48\log K_*}\), at which its effective coupling diverges.

\subsection{Consistency of the seed with the constraint}

\begin{proposition}[Effective frequency]
\label{prop:v19v77-effective}
Let the constraint be imposed on states as in (60.5) with the
fourth-order diagonal of Theorem \ref{thm:v19v73-limit}.  The
one-graviton sector of mode \(1\) has label \(\tfrac23\tau^2=E_1(K)+2\Lambda\)
with \(E_1(K)=1/(2c_g)+\sigma_1(K)\).  Identifying the constraint energy
with the generator of the seed, the seed must have ratio
\(e^{-\varepsilon\omega^{\rm eff}_1}\), \(\omega^{\rm eff}_1=2c_gE_1(K)\),
and a stationary reflection-positive Ornstein--Uhlenbeck seed exists
if and only if \(E_1(K)>0\).
\end{proposition}

\begin{proof}
The label is (60.5) with \(e=e^{(2)}+e_4\), and \(e_4\) linear in
\(n_1\) with coefficient \(\sigma_1(K)\); a Gaussian kernel of ratio
\(r\) has a stationary law exactly when \(0<r<1\), that is
\(\omega^{\rm eff}>0\) (Corollary \ref{cor:v19v33-gaussian}).
\end{proof}

\subsection{The sign change}

\begin{theorem}[Intrinsic cutoff of the bare model]
\label{thm:v19v77-cutoff}
With the closed forms (79.3), the exact partial sums give a first
negative value of \(E_1(K)\) at

\[
 K_*(\tfrac34)=6,\qquad K_*(1)=12,\qquad K_*(2)=59,                     \tag{80.1}
\]

while for \(c_g=5\), \(E_1(K)\) is positive through \(K=400\), rising
to about \(354\) near \(K=100\) and falling to about \(82\) at
\(K=400\).  Asymptotically \(E_1(K)=\tfrac1{2c_g}-\tfrac{K^2}{16c_g^2}+3c_g^2\log K+O(K)\),
so \(K_*\) is the root of \(K^2\approx48c_g^4\log K+8c_g\), which for
the three computed couplings gives \(5.3\), \(11\), \(56\), consistent
with (80.1).  Below \(K_*\) the effective coupling
\(c_g^{\rm eff}(K):=1/(2E_1(K))\) is finite and diverges as
\(K\uparrow K_*\); for the reference coupling
\(E_1(2),\dots,E_1(5)\approx1.96,1.99,1.55,0.77\) and
\(E_1(6)<0\).
\end{theorem}

\begin{proof}
Direct exact summation, certified through \(K=400\) for four
couplings; the asymptotic form is Theorem
\ref{thm:v19v76-self-energy}.
\end{proof}

\subsection{The on-shell counterterm}

\begin{theorem}[Counterterm compatible with the Einstein normalization]
\label{thm:v19v77-counterterm}
Add to the averaged constraint the mode-diagonal counterterm
\(-\sum_k\sigma_k(K)\widehat N_k\).  Then for every cutoff the
one-graviton labels are the tree values \(\omega_k/(2c_g)\), the seed
keeps its ratios \(e^{-\varepsilon\omega_k}\) and variances
\(v_k=1/(2c_g\omega_k)\), the exact spectral consistency of Theorem
\ref{thm:v19v44-spectrum} is untouched, and the remaining cutoff
dependence of the fourth-order sector spectrum is confined to the
finite shell matrices, which saturate at \(K=E\), and to the c-number
zero point.  The counterterm itself grows like
\(\sigma_1(K)\sim-K^2/(16c_g^2)\).
\end{theorem}

\begin{proof}
The counterterm cancels exactly the \(n_k\)-linear cutoff-dependent
part identified in Theorem \ref{thm:v19v73-limit}; nothing else in
the seed or in the shell-preserving part depends on \(K\) beyond
\(K=E\).
\end{proof}

Thus the frequency counterterm is compatible with the Einstein
normalization in the on-shell scheme: the physical one-graviton energy
is held at its tree value and the bare frequency runs.  The
alternative, keeping the bare model, is a theory with an intrinsic
cutoff \(K_*(c_g)\) beyond which no reflection-positive seed exists;
in the units of the torus of period \(2\pi\), with \(c_g\) the
coefficient of the linearized Einstein action, \(K_*\) is the mode
number at which the graviton self-energy equals its tree energy, the
model's analogue of the Planck scale.

\subsection{Executable certificate}

The verifier

\[
 \texttt{scripts/verify\_sm\_v19\_v77\_running\_coupling.py}         \tag{80.2}
\]

imports the closed forms of V76 and computes \(E_1(K)\) exactly for
\(K\leq400\) and \(c_g\in\{\tfrac34,1,2,5\}\), the first negative
values (80.1), the absence of a sign change through \(K=400\) for
\(c_g=5\), the effective couplings, and the counterterm at
\(K=400\).  It emits

\[
 \texttt{artifacts/sm\_v19\_v77\_running\_coupling.json}.            \tag{80.3}
\]

The hashes are

\[
\begin{array}{c|l}
\hbox{object}&\hbox{SHA--256}\\ \hline
\hbox{verifier}&
\texttt{F19AECD73E021D8EE95B73A5DABB36FC61ADAC8F929E360BD23DBF4B1DB5272C}\\
\hbox{canonical payload}&
\texttt{030611C2434DACAC1C052079553C62D5BF3DB9026CF084795DB2E6B4A5F7190C}.
\end{array}                                                     \tag{80.4}
\]

\begin{firewall}[Mode one, one polarization, on-shell scheme]
\label{fire:v19v77}
The statements concern mode one in one polarization at fourth order;
other modes have their own \(\sigma_k\), the second polarization adds
the V69 completion, and the identification of the constraint energy
with the seed generator is the programme's consistency requirement, not
a theorem about full quantum gravity.  The scheme is on-shell; other
schemes differ by finite parts.
\end{firewall}

\begin{decision}[V78 acquisition]
\label{dec:v19v77-next}
V78 will compute \(\sigma_k(K)\) for \(k=2,3\) from the general pair
closed forms, obtained by the method of V76 with \(x_k\cos kz\) in
place of \(\cos z\), and determine whether the counterterms are a
single function of \(k\), the lattice analogue of a universal
wave-function renormalization, or mode dependent in a way that would
require a nonlocal counterterm; this decides whether the renormalized
plane-symmetric model has a local counterterm structure at fourth
order.
\end{decision}

\begin{assessment}[V77 acquisition]
\label{ass:v19v77}
The self-energy divergence is confronted with the seed: without a
counterterm the bare model loses its reflection-positive seed at an
exactly computed cutoff \(K_*(c_g)\), the mode number where the
effective coupling diverges; with the on-shell counterterm the seed
stays at its Einstein normalization for every cutoff and all remaining
cutoff dependence is finite and confined to shell matrices and the
zero point.  The renormalization of the compact model is thereby fixed
at fourth order for mode one.
\end{assessment}

\begin{record}[V77 append record]
\label{rec:v19v77}
V77 is appended to the validated V76 whose installed SHA--256 was
\texttt{377FAE6CE1D25C76E4A4E0607BB8D1AAB5185A37C3E5EA6DCC92587B833F9DB8}.
After installation only V77 is the active numeric SM note; frozen
archives and unrelated notes are unchanged.
\end{record}


\section{V78: general pair closed forms and the mass-plus-wave-function structure of the counterterm}
\label{sec:v19v78}

V77 fixed the counterterm for mode one.  This note derives the pair
coefficients for every pair of frequencies in closed form, by the same
reduction as V76 with two arbitrary frequencies, and computes the
self-energies of modes two and three.  The counterterms are local in
the field-theoretic sense: the quadratically divergent part is a
graviton mass term, \(\delta m^2=K^2/(4c_g)\), the same for every
mode, and the logarithmically divergent part is proportional to the
squared frequency, a wave-function renormalization.  A graviton mass
counterterm is forbidden in the diffeomorphism-invariant theory, so its
appearance identifies precisely what the plane-symmetric truncation
with fixed gauge lacks: the vector and gauge sector that would cancel
it.  This is the upstream target the note records.

\subsection{General pair closed forms}

For \(f=x_k\cos kz+x_l\cos lz\), \(k\neq l\), put

\[
 A_\mp:=\frac{2k^2+2l^2\mp3kl}{16(l\mp k)^2}.                             \tag{81.1}
\]

\begin{theorem}[Pair coefficients for arbitrary frequencies]
\label{thm:v19v78-closed}

\[
 c^{xx}_{kl}=-\frac58(k^2+l^2)+12\bigl[(l-k)^2A_-^2+(l+k)^2A_+^2\bigr]-kl(A_--A_+),
 \qquad
 c^{pp}_{kl}=3\Bigl[\frac1{(l-k)^2}+\frac1{(l+k)^2}\Bigr].                \tag{81.2}
\]

For fixed \(k\) and \(l\to\infty\), \(c^{xx}_{kl}=-l^2/4+O(l)\) and
\(c^{pp}_{kl}=6/l^2+O(l^{-4})\).
\end{theorem}

\begin{proof}
Repeat the proof of Theorem \ref{thm:v19v76-closed} with general
\(k\): the \(x_k^2x_l^2\) coefficient of \(\langle f^2f'^2\rangle\) is
\((k^2+l^2)/4\), the cross part of \(R_2\) is
\(x_kx_l[(\tfrac{3kl}2-k^2-l^2)\cos(l-k)z-(k^2+l^2+\tfrac{3kl}2)\cos(l+k)z]\),
the cross part of \(f'^2\) is \(klx_kx_l[\cos(l-k)z-\cos(l+k)z]\), and
the momentum cross part of \(|A_0|^2\) is
\(2p_kp_l[\cos(l-k)z+\cos(l+k)z]\); the resonant coincidences
\(l-k=2k\) affect only odd monomials.  The certificate checks (81.2)
against the interpolated forms for the pairs \((1,2)\), \((1,3)\),
\((2,3)\), \((1,4)\), \((2,4)\), \((3,4)\), \((2,5)\).
\end{proof}

\subsection{Self-energies and the counterterm structure}

\begin{theorem}[Mass and wave-function counterterms]
\label{thm:v19v78-counterterm}
With \(\sigma_k(K)=\sum_{l\leq K,\,l\neq k}\bigl[c^{xx}_{kl}/(2c_g^2kl)+c^{pp}_{kl}c_g^2kl/2\bigr]\)
and the on-shell counterterm \(\delta\omega_k(K)=-2c_g\sigma_k(K)\),

\[
 2k\,\delta\omega_k(K)=\frac{K^2}{4c_g}-12c_g^3k^2\log K+O(K),           \tag{81.3}
\]

so that, writing \(\omega_k^2\to\omega_k^2+2\omega_k\delta\omega_k\),
the counterterm is a mode-independent mass term
\(\delta m^2(K)=K^2/(4c_g)\) plus a wave-function term
\(-12c_g^3\log K\) multiplying \(\omega_k^2\).  For \(c_g=3/4\) the
exact values of \(2k\delta\omega_k/K^2\) at \(K=300\) are
\(0.3341\), \(0.3330\), \(0.3306\) for \(k=1,2,3\), approaching
\(1/(4c_g)=1/3\).
\end{theorem}

\begin{proof}
Insert the asymptotics of Theorem \ref{thm:v19v78-closed}:
\(c^{xx}_{kl}/(2c_g^2kl)=-l/(8c_g^2k)+O(1/k)\) sums to
\(-K^2/(16c_g^2k)\), and \(c^{pp}_{kl}c_g^2kl/2=3c_g^2k/l+O(l^{-3})\)
sums to \(3c_g^2k\log K\); multiply by \(-2c_g\) and by \(2k\).  The
exact values are certified.
\end{proof}

\subsection{What the mass counterterm means}

\begin{proposition}[The truncation breaks the massless protection]
\label{prop:v19v78-gauge}
In the full diffeomorphism-invariant theory the graviton is protected
from a mass counterterm; the plane-symmetric model has instead fixed
the shift to zero, the transverse momenta to homogeneous constants
(Theorem \ref{thm:v19v57-solution}), and the gauge freedom of the
\(ab\)-block to the TT parametrization, so its fourth-order averaged
constraint is not the gauge-invariant effective action.  The
mode-independent \(K^2/(4c_g)\) is therefore a property of the
truncation, and the vector and gauge sector omitted since V56 is
exactly what must supply the cancelling contribution.
\end{proposition}

\begin{proof}
The statement about the full theory is the standard Ward-identity
protection of the graviton pole; the list of what the model fixes is
V56--V57; the conclusion is that a gauge-noninvariant truncation can
and here does generate a term that gauge invariance forbids.
\end{proof}

\subsection{Executable certificate}

The verifier

\[
 \texttt{scripts/verify\_sm\_v19\_v78\_general\_pair\_counterterm.py}
                                                                    \tag{81.4}
\]

checks (81.2) against exact interpolation for seven pairs, computes
\(\sigma_k(K)\) and \(\delta\omega_k(K)\) for \(k\leq3\) through
\(K=300\) with \(c_g=3/4\), and verifies that \(2k\delta\omega_k/K^2\)
approaches \(1/3\) within \(0.1\) for every \(k\).  It emits

\[
 \texttt{artifacts/sm\_v19\_v78\_general\_pair\_counterterm.json}.   \tag{81.5}
\]

The hashes are

\[
\begin{array}{c|l}
\hbox{object}&\hbox{SHA--256}\\ \hline
\hbox{verifier}&
\texttt{05DE8C2F6324ED98E0043AAD8AA08C7DDF7246C77DAE39F2DB843A420EED0495}\\
\hbox{canonical payload}&
\texttt{EDF0E7524668AAE432C2125B3B9F53B06586BD80421D34770D8C6730E3D0AAFA}.
\end{array}                                                     \tag{81.6}
\]

\begin{firewall}[Truncated model]
\label{fire:v19v78}
The counterterm structure is that of the plane-symmetric, fixed-gauge,
fourth-order averaged constraint in one polarization; the assertion
that the omitted vector and gauge sector cancels the mass term is a
target, not a theorem.
\end{firewall}

\begin{decision}[V79 acquisition]
\label{dec:v19v78-next}
V79 will reinstate the vector sector of the plane-symmetric class: a
nonzero shift \(X^a(z)\) and the homogeneous transverse momenta
\(K^{az}\) of (60.2), and compute their fourth-order contribution to
the averaged constraint and to \(\sigma_k\).  If the mass term
\(K^2/(4c_g)\) cancels against it, the truncation is repaired; if not,
the remaining term must be attributed to the fixed \(ab\)-gauge and
the note will say so.  The V80 audit and ten-version sweep follow.
\end{decision}

\begin{assessment}[V78 acquisition]
\label{ass:v19v78}
The pair coefficients are now closed rational functions of both
frequencies, the self-energies of the first three modes are exact,
and the counterterm splits into a universal quadratic graviton mass
term and a logarithmic wave-function term proportional to the squared
frequency.  Both are local; the mass term is forbidden by
diffeomorphism invariance and locates the defect of the truncation in
its omitted gauge sector.
\end{assessment}

\begin{record}[V78 append record]
\label{rec:v19v78}
V78 is appended to the validated V77 whose installed SHA--256 was
\texttt{49422116567C1EA43B1249213B133F7F67047BE8C5690E4DC2D012CDD3CDB613}.
After installation only V78 is the active numeric SM note; frozen
archives and unrelated notes are unchanged.
\end{record}


\section{V79: gauge completeness of the plane-symmetric class and the origin of the mass term}
\label{sec:v19v79}

V78 asked whether a vector or gauge sector inside the plane-symmetric
class could cancel the graviton mass counterterm.  It cannot, for an
exact reason: within plane symmetry the metric components
\(g_{zz}\) and \(g_{za}\) are pure gauge up to three constants, a
length modulus and two transverse zero modes, so the parametrization
\(dz^2+G_{ab}(z)dx^adx^b\) used since V56 is a complete gauge fixing of
the class and contains no further sector.  The mass term is therefore
a property of the truncation of the mode sum to collinear modes, and
its cancellation, if it occurs, must come from modes whose momenta are
not parallel to the wave, which lie outside the plane class.  This
turns the three-dimensional problem, so far listed only as an absence,
into a specific computation.

\subsection{Gauge fixing within plane symmetry}

\begin{theorem}[Complete gauge fixing of the plane-symmetric class]
\label{thm:v19v79-gauge}
Let \(g=g_{zz}(z)dz^2+2g_{za}(z)dz\,dx^a+G_{ab}(z)dx^adx^b\) be a
plane-symmetric Riemannian metric on the torus.  Then there is a
plane-symmetric diffeomorphism \(x^a\mapsto x^a+\xi^a(z)\),
\(z\mapsto\tilde z(z)\), after which
\(g=d\tilde z^2+2c_a\,d\tilde z\,dx^a+\widetilde G_{ab}(\tilde z)dx^adx^b\)
with constants \(c_a\); the \(\tilde z\)-period is the invariant
length \(L'=\int_0^L\sqrt{g^{\rm new}_{zz}}\,dz\), where
\(g^{\rm new}_{zz}\) is the \(zz\)-component after the transverse
step.  The gauge-invariant data of the class are the length modulus
\(L'\), the two constants \(c_a\), and the block
\(\widetilde G_{ab}\).  The constants \(c_a\) are conjugate to the
homogeneous transverse momenta of (60.2), and they and \(L'\) are the
only degrees of freedom of the class outside the \(ab\)-block.
\end{theorem}

\begin{proof}
Under \(x^a\mapsto x^a+\xi^a(z)\), \(dx^a\mapsto dx^a+\xi^{a\prime}dz\),
so \(g_{za}\mapsto g_{za}+G_{ab}\xi^{b\prime}\) and
\(g_{zz}\mapsto g_{zz}+2g_{za}\xi^{a\prime}+G_{ab}\xi^{a\prime}\xi^{b\prime}\),
while \(G_{ab}\) is unchanged.  Choose
\(\xi^{a\prime}=-(G^{-1})^{ab}g_{zb}+(G^{-1})^{ab}c_b\) with the constants
\(c_b\) fixed by periodicity, \(\int\xi^{a\prime}dz=0\); then the new
\(g_{za}\) equals \(c_a\), constant.  The new \(g_{zz}\) is a positive
periodic function, and \(\tilde z=\int\sqrt{g_{zz}^{\rm new}}\,dz\)
makes it one, with period \(L'\); this reparametrization preserves the
form of the other components up to the substitution
\(z\to\tilde z\).  The certificate performs the transverse step exactly
for \(G=I\) and oscillating \(g_{za}\).
\end{proof}

\subsection{Consequence for the mass counterterm}

\begin{corollary}[No plane-symmetric cancellation]
\label{cor:v19v79-no-cancel}
The fourth-order averaged constraint of the plane-symmetric class with
homogeneous transverse momenta and constant \(c_a\) has no further
dynamical sector; the counterterm structure of Theorem
\ref{thm:v19v78-counterterm} is the complete fourth-order counterterm
of the class.  In particular the graviton mass term \(K^2/(4c_g)\) is
generated by the sum over collinear modes alone and is not cancelled
within plane symmetry.
\end{corollary}

\begin{proof}
By Theorem \ref{thm:v19v79-gauge} every plane-symmetric configuration
is gauge equivalent to one in the V56 form up to three constants,
which are zero modes without \(z\)-dependence and do not enter the
mode sums.
\end{proof}

\begin{proposition}[Where the cancellation must come from]
\label{prop:v19v79-where}
In the full theory the fourth-order self-energy of a graviton of
momentum \(k\hat z\) receives contributions from vacuum fluctuations of
all modes \(l\in\mathbb Z^3\), and the plane class retains only
\(l\parallel\hat z\).  The isotropic sum over \(l\) differs from the
collinear one already in its second-order coupling structure, since
the pair form of non-collinear modes involves the angle between
\(k\) and \(l\) through the TT projections.  Cancellation of the mass
term, required by diffeomorphism invariance of the full theory, is
therefore a statement about the angular sum, and is the first
three-dimensional computation the programme must perform.
\end{proposition}

\begin{proof}
The plane class is the set of configurations invariant under
translations in \(x\) and \(y\); modes with \(l\not\parallel\hat z\) break
that invariance and are excluded by construction.  The remaining
statements describe the structure of the isotropic sum and are not
computed here.
\end{proof}

\subsection{Executable certificate}

The verifier

\[
 \texttt{scripts/verify\_sm\_v19\_v79\_plane\_gauge\_completeness.py}
                                                                    \tag{82.1}
\]

takes \(G=I\) and oscillating \(g_{za}\) with nonzero means, applies
\(\xi^{a\prime}=-(g_{za}-\langle g_{za}\rangle)\), verifies exactly that
the transformed \(g_{za}\) are the constants \(\langle g_{za}\rangle\),
computes the transformed \(g_{zz}\) as a cosine series, and samples its
positivity.  It emits

\[
 \texttt{artifacts/sm\_v19\_v79\_plane\_gauge\_completeness.json}.   \tag{82.2}
\]

The hashes are

\[
\begin{array}{c|l}
\hbox{object}&\hbox{SHA--256}\\ \hline
\hbox{verifier}&
\texttt{ED6664750189B32CDA992593885C65349ADDA82BAC3523262DFED1327667B1F3}\\
\hbox{canonical payload}&
\texttt{DA1E7FAFBD0920D9E722E5E1CBBA825118AA8F3927F61E0756CC1E9F6A6CD1C0}.
\end{array}                                                     \tag{82.3}
\]

\begin{firewall}[Gauge statement only]
\label{fire:v19v79}
The theorem concerns spatial plane-symmetric diffeomorphisms; the
claim that the isotropic sum cancels the mass term is the expectation
from diffeomorphism invariance and is not proved.  Nothing
three-dimensional is computed.
\end{firewall}

\begin{decision}[V80 audit and V81 acquisition]
\label{dec:v19v79-next}
V80 is the compulsory companion audit and ten-version sweep.  V81 will
begin the three-dimensional computation: the second-order Hamiltonian
constraint density for two TT plane waves with non-parallel momenta
\(k\) and \(l\) on the torus, from the general second variation of the
scalar curvature, reduced to the invariants \(|k|,|l|,k\cdot l\) and the
polarization tensors, which is the input needed for the isotropic
self-energy sum and the test of the mass-term cancellation.
\end{decision}

\begin{assessment}[V79 acquisition]
\label{ass:v19v79}
The plane-symmetric class is completely gauge fixed by the
parametrization in use up to a length modulus and two transverse zero
modes, so the graviton mass counterterm of V78 is intrinsic to the
collinear truncation and cannot be cancelled inside it.  The
three-dimensional computation is now specified: the angular sum over
non-collinear modes.
\end{assessment}

\begin{record}[V79 append record]
\label{rec:v19v79}
V79 is appended to the validated V78 whose installed SHA--256 was
\texttt{D9DFB5038D501894AEECAA688AE5CE11AB34C6D93AA08177D917C106BCD93338}.
After installation only V79 is the active numeric SM note; frozen
archives and unrelated notes are unchanged.
\end{record}


\section{V80: compulsory companion audit and ten-version sweep before the three-dimensional computation}
\label{sec:v19v80}

This is the required audit following V75 and the ten-version sweep
after V70.  Every active companion note in the shared folder was
inspected read-only before the non-collinear second-order constraint
announced in V79 is computed.

\begin{record}[V80 companion snapshot]
\label{rec:v19v80-snapshot}
The inspected files were

\[
\begin{array}{c|r|r}
\hbox{file and SHA--256}&\hbox{bytes}&\hbox{lines}\\ \hline
\texttt{Renewal\_Yang\_Mills\_Mass\_Gap\_Research\_Notes\_V33.tex}
&473284&11963\\
\multicolumn{3}{l}{
\texttt{B417E58E5CF6BFBE9F3FB0DD21106355674D525AE9BF8464EFD7AA580D0134F9}}
\\
\texttt{Renewal\_NS\_Stability\_Research\_Notes\_V26.tex}
&278283&5319\\
\multicolumn{3}{l}{
\texttt{82C887F8C2C69E4F28FAD7C3DC8DE5B9099CB5A4BF431EBA1FFF3E91935978AD}}
\\
\texttt{Renewal\_GRH\_Cofinal\_Visibility\_Attack\_V55.tex}
&537152&16260\\
\multicolumn{3}{l}{
\texttt{D9A771A252139954C008CC1CFDBFC9B87FE4A98B6F053835190C2B8AE955BEA9}}
\\
\texttt{Renewal\_YM\_Parallel\_Research\_Notes\_V5.tex}
&54174&1351\\
\multicolumn{3}{l}{
\texttt{306F1BB84CFA8A8D47EA569EC17E9545F9867C8D32B548EC9D9C444B4F3C0512}}.
\end{array}                                                     \tag{83.1}
\]

Since V75, GRH has advanced from V53 to V55; the other three files are
unchanged.
\end{record}

\subsection{Companion findings}

GRH V55 proves a reduction and a leading-order lemma for a balanced
two-pair Pick expression; its firewall states that its target problem
is open, the two-pair band is not closed, three or more pairs are
untouched, and GRH remains open.  Yang--Mills V33, Navier--Stokes V26,
and Yang--Mills Parallel V5 are unchanged since V75.

\subsection{Sweep of the ten versions V70--V79}

Since the V70 sweep the programme has: assembled the pair interactions
and the fourth-order sector energy (V71); diagonalized the lowest
resonant block exactly (V72); established the shell and parity
structure with cutoff saturation (V73); computed and then proved the
growth of the pair coefficients and the quadratic self-energy
divergence (V74, V76); shown the bare model loses its
reflection-positive seed at an exact cutoff \(K_*(c_g)\) and that the
on-shell counterterm restores the Einstein normalization (V77); derived
the general pair closed forms and identified the counterterm as a
universal graviton mass term plus a wave-function term (V78); and
proved the plane class completely gauge fixed, so that the mass term
must be tested in the isotropic three-dimensional mode sum (V79).  No
erratum was needed in V70--V79.

\begin{decision}[V80 audit decision]
\label{dec:v19v80-audit}
No companion theorem, numerical constant, or executable artifact is
imported.  The V79 direction stands: V81 computes the second-order
constraint density for non-collinear TT plane waves.
\end{decision}

\begin{proof}
The GRH advance concerns Pick expressions of heat-flowed polynomials;
the V81 task is the second variation of the scalar curvature for two
plane waves on the flat torus, an internal Riemannian computation
whose collinear case is (59.2).  Nothing is shared.
\end{proof}

\begin{assessment}[V80 acquisition]
\label{ass:v19v80}
The audit is current.  The plane-symmetric fourth-order programme is
complete at the level of the averaged constraint: exact interaction
constants, exact resonant mixing, exact renormalization structure, and
an exact identification of its defect.  The remaining programme is
three-dimensional: the non-collinear mode sum and the test of the
graviton mass cancellation, then the general constraint system, the
mode-number limit, and the Standard Model coupling.
\end{assessment}

\begin{record}[V80 append record]
\label{rec:v19v80}
V80 is appended to the validated V79 whose installed SHA--256 was
\texttt{14F8AC4202C9CD2D10476BD0A2B918F1A4F4B5A03C64A6B8D32DBBCD201E0C87}.
After installation only V80 is the active numeric SM note; the
companions, frozen archives, and all unrelated notes are unchanged.  The
next compulsory companion audit is V85.
\end{record}


\section{V81: the second-order constraint of two non-collinear waves by an exact three-dimensional engine}
\label{sec:v19v81}

V79 located the test of the graviton mass term in the sum over
non-collinear modes.  This note starts that computation with an exact
three-dimensional engine: metrics are expanded in a formal parameter
with coefficients that are trigonometric polynomials on the torus,
represented by their Fourier coefficients in Gaussian rationals, and
the scalar curvature is computed from the Christoffel symbols and the
Ricci tensor without any symbolic-algebra package.  The engine
reproduces the collinear second-order density of V41 exactly and gives,
for two orthogonal unit-frequency waves with polarizations
\(\operatorname{diag}(1,-1,0)\) and \(\operatorname{diag}(0,1,-1)\), the
second-order scalar curvature in closed form.  Its cross term is
\(2ab\cos x\cos z\): it carries momenta \(k\pm l\), has zero average, and
has no derivative enhancement, in contrast to the collinear cross terms
of V56 whose coefficients grow with the frequencies.

\subsection{The engine}

\begin{proposition}[Exact second-order curvature on the torus]
\label{prop:v19v81-engine}
For \(g=\delta+\varepsilon h\) with \(h\) a trigonometric polynomial
tensor, the expansion \(g^{-1}=\delta-\varepsilon h+\varepsilon^2h^2+O(\varepsilon^3)\),
the Christoffel symbols
\(\Gamma^k_{ij}=\tfrac12g^{kl}(\partial_ig_{lj}+\partial_jg_{li}-\partial_lg_{ij})\),
and \(R_{ij}=\partial_k\Gamma^k_{ij}-\partial_j\Gamma^k_{ik}+\Gamma^k_{kl}\Gamma^l_{ij}-\Gamma^k_{jl}\Gamma^l_{ik}\)
give \(R=g^{ij}R_{ij}\) through order \(\varepsilon^2\) as an exact
finite computation on Fourier coefficients: products are convolutions
of finitely supported coefficient maps and derivatives are
multiplications by \(ik_j\).  For a single TT wave along \(z\) the
engine returns \(R_1=0\) and \(R_2=-\tfrac14-\tfrac74\cos2z\), the
V41 density.
\end{proposition}

\begin{proof}
Each step is a finite sum of products of Fourier coefficients; the
collinear check is the certificate's first assertion.
\end{proof}

\subsection{Two orthogonal waves}

\begin{theorem}[Second-order curvature of an orthogonal pair]
\label{thm:v19v81-pair}
Let \(h=a\cos z\,\operatorname{diag}(1,-1,0)+b\cos x\,\operatorname{diag}(0,1,-1)\),
which is transverse and traceless for each wave.  Then \(R_1=0\) and

\[
 R_2=a^2\Bigl(-\frac14-\frac74\cos2z\Bigr)+b^2\Bigl(-\frac14-\frac74\cos2x\Bigr)
 +2ab\cos x\cos z .                                                       \tag{84.1}
\]

The cross term equals \(ab[\cos(x+z)+\cos(x-z)]\), carries the momenta
\(k\pm l\) with \(|k\pm l|^2=2\), and averages to zero.  The
corresponding second-order conformal response of Theorem
\ref{thm:v19v56-response} is \(w_2^{\rm cross}=-\tfrac{ab}8\cos x\cos z\),
and the averaged constraint receives no cross contribution at second
order.
\end{theorem}

\begin{proof}
The engine evaluates \(R_2\) at \((a,b)=(1,1),(1,-1),(2,1)\); the
bilinear cross part is half the difference of the first two, with
Fourier coefficients \(\tfrac12\) at \((\pm1,0,\pm1)\), and the
quadratic parts are the collinear densities in each direction.  The
response follows from (59.5) with \(|q|^2=2\).
\end{proof}

The absence of any frequency factor in the cross term is the
structural difference from the collinear case, where the pair
components of (59.4) grow like \(k^2\), \(l^2\), and \(kl\): for waves
with orthogonal momenta and these polarizations the curvature
couples the two waves only through the undifferentiated product
\(h_{yy}\)-type overlap.  Since the collinear \(l^2\) growth is what
produced the quadratic self-energy divergence and the mass term of
V76--V78, the isotropic mode sum weights collinear and non-collinear
neighbours very differently, which is the mechanism by which a
cancellation could occur.

\subsection{Executable certificate}

The verifier

\[
 \texttt{scripts/verify\_sm\_v19\_v81\_noncollinear\_second\_order.py}
                                                                    \tag{84.2}
\]

implements the three-dimensional Fourier engine, checks the collinear
density of V41, evaluates \(R_2\) for the orthogonal pair at three
amplitude points, extracts the cross term (84.1), and verifies its
zero average.  It emits

\[
 \texttt{artifacts/sm\_v19\_v81\_noncollinear\_second\_order.json}.  \tag{84.3}
\]

The hashes are

\[
\begin{array}{c|l}
\hbox{object}&\hbox{SHA--256}\\ \hline
\hbox{verifier}&
\texttt{35096391F3D2CCADF5E548D12F1798AF7730D2F39E852817FCA531DFA577E0C0}\\
\hbox{canonical payload}&
\texttt{31F0C68ED550A404749616F49822E3BCEC310809966838BF7DE964B969263816}.
\end{array}                                                     \tag{84.4}
\]

\begin{firewall}[One orthogonal pair at second order]
\label{fire:v19v81}
One pair of momenta and one polarization choice are computed; the
general dependence on \(k\cdot l\) and on the polarization tensors,
the fourth-order terms that enter the self-energy, and the isotropic
sum are not.  The mass-cancellation question remains open.
\end{firewall}

\begin{decision}[V82 acquisition]
\label{dec:v19v81-next}
V82 will extend the engine to fourth order in the metric perturbation
and compute, for the orthogonal pair of Theorem
\ref{thm:v19v81-pair}, the quartic cross coefficients
\(c^{xx}_{kl}\), \(c^{pp}_{kl}\) of the averaged constraint by the same
interpolation as V68, including the conformal response and the
momentum norm on the perturbed background; comparing them with the
collinear values of (81.2) for \(|k|=|l|=1\) will give the first
non-collinear self-energy summand and hence the first test of whether
the isotropic sum can cancel the graviton mass term.
\end{decision}

\begin{assessment}[V81 acquisition]
\label{ass:v19v81}
The programme now has an exact three-dimensional curvature engine and
its first non-collinear result: for orthogonal unit-frequency TT waves
the second-order cross term is \(2ab\cos x\cos z\), without the
frequency enhancement that drives the collinear divergence.  The
mechanism for a cancellation of the graviton mass term in the isotropic
sum is thereby identified, though not yet computed.
\end{assessment}

\begin{record}[V81 append record]
\label{rec:v19v81}
V81 is appended to the validated V80 whose installed SHA--256 was
\texttt{906BC2D0C3DB5CAD3AF9D82266EBF3B1E07ED7CED6FE3757814E65D56C30F789}.
After installation only V81 is the active numeric SM note; frozen
archives and unrelated notes are unchanged.
\end{record}


\section{V82: fourth-order balance of non-collinear pairs and the persistence of the quadratic growth}
\label{sec:v19v82}

V81 identified the absence of frequency enhancement in the
non-collinear second-order cross term as the mechanism by which an
isotropic mode sum might cancel the graviton mass counterterm.  This
note carries the three-dimensional engine to fourth order for the
time-symmetric balance, with the complete order-by-order scheme in
which the conformal factor responds at every order and the
Laplace--Beltrami corrections are computed exactly.  The engine
reproduces the collinear values \(\Lambda_2=-1/8\), \(\Lambda_3=0\),
\(\Lambda_4=-145/2048\) of V41 and V66.  For a unit wave along \(z\)
paired with a wave of frequency \(l\) along \(x\), the fourth-order
cross coefficient is \(-\tfrac5{64}(l^2+1)\) for \(l\leq4\): it still
grows like \(l^2\), five eighths of the collinear rate.  The hoped-for
cancellation does not occur in the time-symmetric configuration
sector; the isotropic sum of the coherent-state self-energy diverges
like the fourth power of the cutoff, and the protection of the massless
graviton must reside in the momentum sector and the dynamics, which
are recorded as the next target.

\subsection{Order-by-order scheme in three dimensions}

Let \(\gamma=\delta+\varepsilon h\) with \(h\) transverse traceless,
\(\psi=1+\sum_{n\geq2}\varepsilon^nw_n\) with flat-mean-zero \(w_n\),
\(\Lambda=\sum_{n\geq2}\varepsilon^n\Lambda_n\), and let
\(\Delta_\gamma=\Delta+\sum_{n\geq1}\varepsilon^nD_n\).  The exact
time-symmetric constraint \(-8\Delta_\gamma\psi+R_\gamma\psi-2\Lambda\psi^5=0\)
gives

\[
 \Lambda_2=\tfrac12\langle R_2\rangle,\qquad
 \Lambda_3=\tfrac12\langle R_3-8D_1w_2\rangle,\qquad
 \Lambda_4=\tfrac12\langle R_4+R_2w_2-8D_1w_3-8D_2w_2\rangle,           \tag{85.1}
\]

with \(w_2\), \(w_3\) the mean-zero solutions of
\(-8\Delta w_n=-(S_n-\langle S_n\rangle)\), \(S_2=R_2\),
\(S_3=R_3-8D_1w_2\).

\begin{proposition}[Exactness and the collinear check]
\label{prop:v19v82-scheme}
The scheme (85.1) is exact through fourth order, the flat average is
the correct average because the equation is written in the
non-densitized form, and for a single TT wave along \(z\) the engine
returns \(\Lambda_2=-\tfrac18\), \(\Lambda_3=0\), \(\Lambda_4=-\tfrac{145}{2048}\),
in agreement with Theorems \ref{thm:v19v41-branch} and
\ref{thm:v19v66-fourth}.  The conformal--TT parametrization
\(g=\psi^4(\delta+\varepsilon h)\) is complete for the spatial metric
up to spatial diffeomorphisms, so the coefficients below are gauge
invariant.
\end{proposition}

\begin{proof}
Expand the constraint and collect orders; \(D_1w_2\) vanishes in the
collinear case because \(h_{zz}=0\) and \(\partial_ih_{ij}=\operatorname{tr}h=0\),
which is why V66 needed only \(D_2\).  Completeness of the
conformal--TT form is the York decomposition applied order by order.
\end{proof}

\subsection{Non-collinear pairs}

\begin{theorem}[Fourth-order cross coefficients]
\label{thm:v19v82-cross}
For \(h=a\cos z\operatorname{diag}(1,-1,0)+b\cos(lx)\operatorname{diag}(0,1,-1)\),
\(\Lambda_3=0\) and \(\Lambda_4=\alpha a^4+\beta_lb^2a^2+\gamma_lb^4\)
with \(\alpha=-\tfrac{145}{2048}\), \(\gamma_l=l^2\alpha\), and

\[
 \beta_1=-\frac5{32},\quad\beta_2=-\frac{25}{64},\quad\beta_3=-\frac{25}{32},\quad\beta_4=-\frac{85}{64},
 \qquad\hbox{that is}\qquad \beta_l=-\frac5{64}(l^2+1)\ (l\leq4).            \tag{85.2}
\]

For the oblique pair \(l=(1,0,1)\) with polarization
\(\tfrac12\bigl(\begin{smallmatrix}-1&0&1\\0&2&0\\1&0&-1\end{smallmatrix}\bigr)\),
\(\beta=-\tfrac{1503}{5120}\).  The collinear comparison value is
\(c^{xx}_{1l}/2=-\tfrac{l^2}8+O(l)\).
\end{theorem}

\begin{proof}
The engine evaluates (85.1) at \((a,b)=(1,0),(0,1),(1,1),(2,1)\); the
quartic form is even in \(a\) and in \(b\) by the translations
\(z\to z+\pi\), \(x\to x+\pi/l\), and the fourth point verifies the
form exactly.  Values are certified.
\end{proof}

\subsection{No cancellation in the configuration sector}

\begin{corollary}[Isotropic sum of the coherent-state self-energy]
\label{cor:v19v82-no-cancel}
Under the coherent-state dictionary of V61, the fourth-order balance
coefficient \(\beta_l\) is the configuration part of the self-energy
of the \(z\)-graviton due to the mode \(l\).  With \(\beta_l\) growing
like \(|l|^2\) for orthogonal as well as collinear momenta, the
isotropic sum over \(l\in\mathbb Z^3\), \(|l|\leq K\), of the summand
\(\beta_lv_1v_l\propto|l|^2/|l|\) diverges like \(K^4\).  The angular
average therefore does not cancel the graviton mass term of V78; it
enlarges its degree of divergence from quadratic to quartic.
\end{corollary}

\begin{proof}
\(\sum_{|l|\leq K}|l|\sim K^4\) in three dimensions; the sign of
\(\beta_l\) is negative in every computed case.
\end{proof}

This is the expected divergence of a graviton mass term with a
three-dimensional momentum cutoff, and it confirms that the
time-symmetric configuration sector, which is all the classical
balance sees, cannot exhibit the cancellation that diffeomorphism
invariance guarantees for the full theory.  The cancellation must
come from the momentum sector, which enters the fourth-order
constraint through \(|A|^2_\gamma\) and the momentum part of the
conformal response, and from the dynamics that relate configuration
and momentum fluctuations in the vacuum; both are absent from a
time-symmetric slice.

\subsection{Executable certificate}

The verifier

\[
 \texttt{scripts/verify\_sm\_v19\_v82\_noncollinear\_fourth\_order.py}
                                                                    \tag{85.3}
\]

implements the fourth-order three-dimensional engine with the exact
Laplace--Beltrami corrections, checks the collinear values, and
computes (85.2) and the oblique value, verifying the quartic-form
structure at a fourth amplitude point.  It emits

\[
 \texttt{artifacts/sm\_v19\_v82\_noncollinear\_fourth\_order.json}.  \tag{85.4}
\]

The hashes are

\[
\begin{array}{c|l}
\hbox{object}&\hbox{SHA--256}\\ \hline
\hbox{verifier}&
\texttt{59A2447B3B6CA6249D5DD23E60931982DD7DFD9703523CE7C31135CB9E6235FC}\\
\hbox{canonical payload}&
\texttt{49EF8AAF1F1D6C9EF1773F44FD75189418AA5CCEC135EF59B52A7A0F9C99E7F6}.
\end{array}                                                     \tag{85.5}
\]

\begin{firewall}[Configuration sector only]
\label{fire:v19v82}
The coefficients are for the time-symmetric balance; the momentum
sector of the fourth-order constraint on a non-collinear background,
the closed form of \(\beta_l\) for all \(l\), and the polarization
dependence are not derived.  The non-cancellation is a statement about
this sector, not a contradiction with the full theory.
\end{firewall}

\begin{decision}[V83 acquisition]
\label{dec:v19v82-next}
V83 will add the momentum sector in three dimensions: the
\(\gamma\)-transverse-traceless completion of a TT momentum on the
perturbed background by the York projection, computed with the engine,
and the resulting \(p\)-dependent fourth-order terms for the
orthogonal pair, so that the full configuration-plus-momentum
self-energy summand can be compared with the requirement of a
massless graviton.
\end{decision}

\begin{assessment}[V82 acquisition]
\label{ass:v19v82}
The three-dimensional fourth-order balance is exact and confirms the
plane-symmetric results as its collinear case.  Non-collinear pairs
have cross coefficients that still grow like the square of the
frequency, so the isotropic configuration self-energy diverges
quartically and the mass counterterm is not cancelled by angular
averaging.  The protection of the massless graviton is thereby placed,
by exclusion, in the momentum sector and the dynamics.
\end{assessment}

\begin{record}[V82 append record]
\label{rec:v19v82}
V82 is appended to the validated V81 whose installed SHA--256 was
\texttt{89F390BE705543A5AB4564DD22AD41CA9E3EC2797D0C8684C9ED4274AA2A118A}.
After installation only V82 is the active numeric SM note; frozen
archives and unrelated notes are unchanged.
\end{record}


\section{V83: the momentum sector in three dimensions, and an erratum for the plane-class momentum completion}
\label{sec:v19v83}

V82 placed the protection of the massless graviton in the momentum
sector.  This note adds that sector to the three-dimensional engine by
the York projection: the TT momentum seed is completed order by order
so that it is traceless and divergence free with respect to the
perturbed metric, with the trace and divergence conditions solved
exactly in Fourier space.  The first result is an erratum.  The
plane-class completion used in V67--V69, which corrected only the
trace of the momentum inside the transverse block, violates the
second-order momentum constraint: it omits the longitudinal component
\(K^{zz}_{(2)}\) that V57 itself derived.  The configuration-sector
results of V66--V82 are untouched, but every momentum-dependent
coefficient since V67 must be replaced by its York value; the corrected
single-mode form is given here, and the affected statements are
listed.  The second result is the momentum-dependent fourth-order
balance of the orthogonal pair, with a nonvanishing coupling between
the amplitude of one wave and the momentum of the other, absent for
collinear pairs.

\subsection{York projection order by order}

\begin{proposition}[Exact York completion]
\label{prop:v19v83-york}
Let \(\gamma=\delta+\varepsilon h\) and let \(A_0\) be a flat TT tensor.
Define \(A=\varepsilon A_0+\varepsilon^2A_2+\varepsilon^3A_3\) by requiring
\(\gamma^{ij}A_{ij}=0\) and \(\nabla^\gamma_jA^{ij}=0\) at orders two and
three, with \(A_n=\tfrac13t_n\delta+L W_n\), \(t_n\) the required flat
trace and \(W_n\) the solution of
\(\Delta W_i+\tfrac13\partial_i\partial\cdot W=S_i\) in Fourier space
(the matrix \(-|q|^2I-\tfrac13qq^{\mathsf T}\) is invertible for
\(q\neq0\)), the flat TT part of \(A_n\) being set to zero.  The
construction succeeds whenever the mean of the source vanishes, which
is the translation-charge condition of V27 and V57, and it holds for
the standing waves used here.
\end{proposition}

\begin{proof}
The trace condition at order \(n\) fixes \(t_n\) from lower orders;
the divergence condition at order \(n\) reads
\(\partial_jA_n^{ij}+(\hbox{known})=0\), and
\(\partial_j(LW)_{ij}=\Delta W_i+\tfrac13\partial_i\partial\cdot W\);
the engine verifies both conditions exactly after the solve.
\end{proof}

\subsection{Erratum: the plane-class completion of V67--V69}

\begin{proposition}[Constraint violation of the diagonal completion]
\label{prop:v19v83-erratum}
The completion \(A=\operatorname{diag}(pc(1+xc),-pc(1-xc))\),
\(A^{zz}=A^{za}=0\), of Theorem \ref{thm:v19v67-quartic}, and its
two-polarization version in Theorem \ref{thm:v19v69-norm}, satisfy
the trace condition but not the momentum constraint: by (60.1) with
\(K^{zz}=0\) the constraint requires \(G'_{ab}K^{ab}=0\), whereas
\(G'_{ab}K^{ab}=2xp\,cc'/(1-x^2c^2)\neq0\).  The missing component is
exactly \(K^{zz}_{(2)}=\tfrac12\int\operatorname{tr}(H'A)\) of
Proposition \ref{prop:v19v57-standing}.  Consequently the
momentum-dependent coefficients of (70.1), (71.1), (72.1), (74.1),
(79.3)--(79.5), (80.1), (81.2)--(81.3), and the momentum parts of the
sector energies, mixing (75.1), and self-energies built from them,
are not values of the constrained model.  The configuration-sector
statements, namely all coefficients of pure amplitude monomials,
\(\Lambda_4\), \(c^{xx}_{kl}\), the leading quadratic divergence, the
mass counterterm \(K^2/(4c_g)\), V79, and V82, are unaffected.
\end{proposition}

\begin{proof}
Direct evaluation of (60.1) for the stated completion.  The
configuration statements have \(p=0\), where the completion is
trivial.
\end{proof}

\begin{theorem}[Corrected single-mode momentum form]
\label{thm:v19v83-single}
With the York completion, the quartic averaged constraint of one
standing mode is

\[
 Q_1=-\frac{145}{1024}k^2x^4-\frac{49}{128}x^2p^2+\frac{7}{64k^2}p^4,
                                                                    \tag{86.1}
\]

in place of (70.1), and the corrected self-interaction coefficient is
\(\kappa_2=-\tfrac{435}{2048c_g^2}-\tfrac{49}{256}+\tfrac{21c_g^2}{128}\),
in place of (70.4).
\end{theorem}

\begin{proof}
The engine, with the balance scheme of V82 extended by the momentum
terms \(-|A|^2_\gamma\psi^{-7}\), evaluates the single-mode quartic
form at \((x,p)=(1,0),(0,1),(1,1),(1,-1),(2,1)\) and confirms the even
quartic structure; \(k\)-dependence follows from scaling as in
Theorem \ref{thm:v19v67-quartic}.  The Fock diagonal uses (70.2).
\end{proof}

\subsection{The orthogonal pair with momenta}

\begin{theorem}[Momentum-dependent balance of the orthogonal pair]
\label{thm:v19v83-pair}
For amplitudes \((a,b)\) and momenta \((p,q)\) of the unit-frequency
waves along \(z\) and \(x\) of Theorem \ref{thm:v19v81-pair}, the
York-completed quartic balance form is

\[
\begin{aligned}
 Q&=-\frac{145}{1024}(a^4+b^4)-\frac5{16}a^2b^2
 -\frac{49}{128}(a^2p^2+b^2q^2)-\frac{63}{128}(a^2q^2+b^2p^2)\\
 &\quad+\frac7{64}(p^4+q^4)+\frac7{16}p^2q^2+\frac{17}{64}\,abpq .
\end{aligned}                                                          \tag{86.2}
\]

Unlike collinear pairs, the amplitude of one wave couples to the
momentum of the other, with coefficient \(-63/128\).
\end{theorem}

\begin{proof}
Differences of engine evaluations at the amplitude and momentum points
listed in the certificate isolate each monomial; the configuration
cross term reproduces \(2\beta_1=-5/16\) of V82.
\end{proof}

\subsection{Consequences}

The York completion is the constrained one, and its Fock diagonal is
invariant under near-identity canonical redefinitions of the momentum
seed, so the corrected sector energies are meaningful; the earlier
values were not, because their completion left the constraint surface.
The leading divergence structure of V76--V78 survives, since it comes
from the configuration coefficients \(c^{xx}_{kl}\); the logarithmic
momentum-sector term, the values of \(K_*\) in V77, and the
wave-function counterterm coefficient must be recomputed from
York-completed pair forms.

\subsection{Executable certificate}

The verifier

\[
 \texttt{scripts/verify\_sm\_v19\_v83\_momentum\_sector\_3d.py}      \tag{86.3}
\]

implements the York completion on the three-dimensional engine,
verifies the trace and divergence conditions at orders two and three,
checks the configuration limits against V82, evaluates (86.1) and
(86.2), and records the superseded V67 values beside the corrected
ones.  It emits

\[
 \texttt{artifacts/sm\_v19\_v83\_momentum\_sector\_3d.json}.         \tag{86.4}
\]

The hashes are

\[
\begin{array}{c|l}
\hbox{object}&\hbox{SHA--256}\\ \hline
\hbox{verifier}&
\texttt{1B0BB15496B88FDE790A4A20AE1048D90417C6A67AD524FFDE54296CBC79253F}\\
\hbox{canonical payload}&
\texttt{6BEA17252F9000BD8CDA7293E78116333A590CF1A2EEAD7B52ACBA703CC19E54}.
\end{array}                                                     \tag{86.5}
\]

\begin{firewall}[Scope of the erratum and of the new values]
\label{fire:v19v83}
The erratum invalidates momentum-dependent numbers only; it does not
change any theorem of V32--V65 or the configuration results.  The
corrected values are for one collinear mode and one orthogonal pair at
unit frequency; the York-completed pair forms for general
frequencies, the corrected self-energies, and the momentum-sector
contribution to the mass cancellation are not yet computed.
\end{firewall}

\begin{decision}[V84 acquisition]
\label{dec:v19v83-next}
V84 will recompute, with the York-completed engine, the collinear pair
forms for \((1,l)\), \(l\leq6\), and the corresponding momentum
coefficients \(c^{pp}_{1l}\) and \(c^{xp}_{1l}\), replacing (79.3) and
(81.2), and will restate the self-energy asymptotics, the cutoffs
\(K_*\), and the counterterm structure of V76--V78 on the corrected
basis.  The V85 audit follows.
\end{decision}

\begin{assessment}[V83 acquisition]
\label{ass:v19v83}
The three-dimensional momentum sector is now constructed exactly by
York projection, and it exposes that the plane-class momentum
completion of V67--V69 violated the momentum constraint; the
configuration results stand and the momentum-dependent numbers are
superseded, with the corrected single-mode form and the orthogonal
pair form given.  Self-correction by an independent exact engine is
the intended function of the certificate discipline.
\end{assessment}

\begin{record}[V83 append record]
\label{rec:v19v83}
V83 is appended to the validated V82 whose installed SHA--256 was
\texttt{F43609ED24720E82A8FCBE22D21482FC629E4182F24CFD3426E3E32A5C3F7A4A}.
After installation only V83 is the active numeric SM note; frozen
archives and unrelated notes are unchanged.
\end{record}


\section{V84: York-completed collinear pair forms and the corrected self-energy}
\label{sec:v19v84}

V83 superseded every momentum-dependent coefficient since V67.  This
note recomputes the collinear pair forms for the frequencies
\((1,l)\), \(l\leq6\), with the York-completed engine, separating the
resonant odd monomials at \(l=3\) by symmetric combinations, and
restates the self-energy of mode one on the corrected basis.  The
configuration coefficients reproduce the closed form (81.2) exactly,
which confirms that the erratum touches the momentum sector only.  The
corrected momentum cross coefficients are smaller than the superseded
ones and a new coupling between the amplitude of the high mode and the
momentum of the low mode appears, of order one.  With these values the
one-graviton sector coefficient changes sign at cutoff four instead of
six for the reference coupling; the quadratic divergence and the
mass-term structure of V76--V78 are unchanged because they come from
the configuration coefficients.

\subsection{Corrected pair forms}

\begin{theorem}[York-completed pair coefficients \((1,l)\)]
\label{thm:v19v84-pairs}
With \(a,p\) the amplitude and momentum of the unit mode and \(b,q\)
those of the mode \(l\), the even part of the York-completed quartic
balance form has the single-mode coefficients
\(-\tfrac{145}{1024}\), \(-\tfrac{49}{128}\), \(\tfrac7{64}\) for
\(a^4,a^2p^2,p^4\) and their \(l\)-scaled counterparts for
\(b^4,b^2q^2,q^4\), the configuration cross coefficient
\(c^{xx}_{1l}\) of (81.2), and the momentum cross coefficients

\[
\begin{array}{c|c|c|c|c}
l&p^2q^2&b^2p^2&a^2q^2&\hbox{superseded }c^{pp}_{1l}\\ \hline
2&35/18&-3/2&-1/24&10/3\\
3&35/64&-279/256&1/256&15/16\\
4&119/450&-74/75&1/200&34/75\\
5&91/576&-725/768&1/256&13/48\\
6&37/350&-2259/2450&29/9800&222/1225
\end{array}                                                            \tag{87.1}
\]

The odd monomials vanish except at \(l=3\), where \(a^3b\) and
\(p^3q\) carry \(-\tfrac{109}{256}\) and \(\tfrac7{16}\).
\end{theorem}

\begin{proof}
For each pair the engine of V83 evaluates the balance at the pure
points and at \((\pm)\) combinations of each pair of variables; the
even and odd parts separate as the half sum and half difference.  The
configuration cross coefficients are asserted equal to (81.2) in the
certificate.
\end{proof}

The superseded \(c^{pp}_{1l}=3[(l-1)^{-2}+(l+1)^{-2}]\) is replaced by
values decaying faster, and the term \(b^2p^2\), absent in the
superseded forms, is now of order one and slowly varying; its
contribution to the self-energy of the low mode from the vacuum of the
high mode is \(2c^{bp}_{1l}v_l/(4v_1)=c^{bp}_{1l}/(2l)\), a logarithmic
sum.

\subsection{Corrected self-energy and cutoff}

\begin{theorem}[Self-energy of mode one, corrected]
\label{thm:v19v84-self-energy}
The \(n_1\)-linear coefficient from the vacuum of mode \(l\) is

\[
 \sigma_{1l}=2c^{xx}_{1l}v_1v_l+\frac{c^{pq}_{1l}}{8v_1v_l}
 +\frac{c^{aq}_{1l}v_1}{2v_l}+\frac{c^{bp}_{1l}v_l}{2v_1},              \tag{87.2}
\]

and for \(c_g=3/4\) the partial sums \(\sigma_1(K)\) for \(K=2,\dots,6\)
are \(\tfrac{235}{2592}\), \(-\tfrac{62615}{165888}\), \(-\tfrac{35209}{30720}\),
\(-\tfrac{225029}{103680}\), \(-\tfrac{17452877}{5080320}\); the
one-graviton coefficient \(E_1(K)=\tfrac1{2c_g}+\sigma_1(K)\) is
negative from \(K=4\) on, so the intrinsic cutoff of Theorem
\ref{thm:v19v77-cutoff} is \(K_*(\tfrac34)=4\) on the corrected basis.
The leading term \(-K^2/(16c_g^2)\) of (79.5) is unchanged, since it
comes from \(c^{xx}_{1l}\); the subleading terms are now the sum of a
decaying momentum cross term and the logarithmic \(b^2p^2\) sum.
\end{theorem}

\begin{proof}
The linear parts of the factorized Weyl diagonals as in Theorem
\ref{thm:v19v73-limit}, now including the two mixed monomials; values
certified.
\end{proof}

The statements of V77 and V78 hold with these replacements: the
counterterm still contains a universal mass term from
\(c^{xx}_{kl}\), and the wave-function-type logarithm now arises from
\(c^{bp}\) rather than from \(c^{pp}\); its coefficient for general
\((k,l)\) is not derived here.

\subsection{Executable certificate}

The verifier

\[
 \texttt{scripts/verify\_sm\_v19\_v84\_york\_pair\_forms.py}         \tag{87.3}
\]

runs the York-completed engine on the five collinear pairs, checks
the configuration coefficients against (81.2) and the single-mode
momentum coefficients against (86.1), tabulates (87.1), and computes
(87.2) with the sign change of \(E_1\).  It emits

\[
 \texttt{artifacts/sm\_v19\_v84\_york\_pair\_forms.json}.            \tag{87.4}
\]

The hashes are

\[
\begin{array}{c|l}
\hbox{object}&\hbox{SHA--256}\\ \hline
\hbox{verifier}&
\texttt{4661B2F77F5AC4EE528D111924D13F208BF628593A5533C9E29E5C3681E4C029}\\
\hbox{canonical payload}&
\texttt{8C68544DA229066CB1F4DA5495106F6BEF59B27796D09376F9B8569591441578}.
\end{array}                                                     \tag{87.5}
\]

\begin{firewall}[Five pairs, one polarization]
\label{fire:v19v84}
Closed forms in \(l\) for the corrected momentum coefficients are not
derived; the two-polarization wedge term of V69, the mode--mode
momentum terms of V68 and V71, and the resonant mixing element of V72
remain to be recomputed on the York basis.
\end{firewall}

\begin{decision}[V85 audit and V86 acquisition]
\label{dec:v19v84-next}
V85 is the compulsory companion audit.  V86 will recompute the
two-polarization single-mode form with the York completion on the
three-dimensional engine, replacing the wedge coefficient of V69, and
the \((1,3)\) resonant mixing element of V72, so that the fourth-order
sector spectrum of the plane-symmetric model is restored on a
constraint-consistent basis.
\end{decision}

\begin{assessment}[V84 acquisition]
\label{ass:v19v84}
The collinear pair forms are rebuilt on the constrained basis: the
configuration coefficients are confirmed, the momentum cross
coefficients are corrected and a new amplitude--momentum coupling is
found, the self-energy is restated with its leading divergence intact,
and the intrinsic cutoff of the bare model moves from six to four.
\end{assessment}

\begin{record}[V84 append record]
\label{rec:v19v84}
V84 is appended to the validated V83 whose installed SHA--256 was
\texttt{1E87DA6123B34FF174DBCCC5D37B883495783AC45643703C11894E74BE16ED8F}.
After installation only V84 is the active numeric SM note; frozen
archives and unrelated notes are unchanged.
\end{record}


\section{V85: compulsory companion audit after the momentum-sector erratum}
\label{sec:v19v85}

This is the required audit following V80.  Every active companion note
present in the shared folder was inspected read-only before the
York-completed recomputation of the two-polarization form and the
resonant mixing announced in V84 begins.

\begin{record}[V85 companion snapshot]
\label{rec:v19v85-snapshot}
The files present were

\[
\begin{array}{c|r|r}
\hbox{file and SHA--256}&\hbox{bytes}&\hbox{lines}\\ \hline
\texttt{Renewal\_NS\_Stability\_Research\_Notes\_V26.tex}
&278283&5319\\
\multicolumn{3}{l}{
\texttt{82C887F8C2C69E4F28FAD7C3DC8DE5B9099CB5A4BF431EBA1FFF3E91935978AD}}
\\
\texttt{Renewal\_GRH\_Cofinal\_Visibility\_Attack\_V55.tex}
&537152&16260\\
\multicolumn{3}{l}{
\texttt{D9A771A252139954C008CC1CFDBFC9B87FE4A98B6F053835190C2B8AE955BEA9}}
\\
\texttt{Renewal\_YM\_Parallel\_Research\_Notes\_V5.tex}
&54174&1351\\
\multicolumn{3}{l}{
\texttt{306F1BB84CFA8A8D47EA569EC17E9545F9867C8D32B548EC9D9C444B4F3C0512}}.
\end{array}                                                     \tag{88.1}
\]

No active Yang--Mills mass-gap note was present at the time of the
snapshot: its programme was between installations, its V33 having
been read at V80 with hash
\texttt{B417E58E5CF6BFBE9F3FB0DD21106355674D525AE9BF8464EFD7AA580D0134F9}.
The other three files are unchanged since V80.
\end{record}

\begin{decision}[V85 audit decision]
\label{dec:v19v85-audit}
No companion theorem, numerical constant, or executable artifact is
imported.  The V84 direction stands: V86 recomputes the
two-polarization form and the resonant mixing on the York basis.
\end{decision}

\begin{proof}
The three files present are unchanged since V80, where no shared
operator or hypothesis was found; the momentary absence of the
Yang--Mills file cannot supply an import.  The V86 task is internal
to the York-completed engine of V83.
\end{proof}

\begin{assessment}[V85 acquisition]
\label{ass:v19v85}
The audit is current.  Since V80 the programme has computed the
non-collinear second- and fourth-order balances, found no angular
cancellation of the graviton mass term in the configuration sector,
constructed the momentum sector by exact York projection, discovered
and recorded the erratum in the plane-class momentum completion of
V67--V69, and rebuilt the collinear pair forms on the constrained
basis with the cutoff of the bare model moving from six to four.
\end{assessment}

\begin{record}[V85 append record]
\label{rec:v19v85}
V85 is appended to the validated V84 whose installed SHA--256 was
\texttt{6597541708E241A33F53C1C907A58F584B2EF4C5DB1039C2C02A6A4C256CCA45}.
After installation only V85 is the active numeric SM note; the
companions, frozen archives, and all unrelated notes are unchanged.  The
next compulsory companion audit is V90.
\end{record}


\section{V86: York-completed polarization form and the corrected resonant mixing}
\label{sec:v19v86}

V84 rebuilt the collinear pair forms.  This note rebuilds the two
remaining momentum-dependent results superseded by the V83 erratum:
the two-polarization single-mode form of V69 and the resonant mixing
of V72.  The two-polarization form is again rotation invariant and is
fixed by four coefficients; the wedge coefficient is \(-5/64\) instead
of \(-3/4\), and the polarization interaction remains attractive at
the reference coupling.  The resonant block \(|3,0\rangle,|0,1\rangle\)
of the pair \((1,3)\) is now mixed only weakly, because the corrected
fourth-order diagonal energies of the two sectors differ by an amount
much larger than the off-diagonal element; the near-maximal mixing of
V72 was an artifact of the constraint-violating completion.

\subsection{Two-polarization single-mode form}

\begin{theorem}[York-completed rotation-invariant form]
\label{thm:v19v86-form}
For \(h=\cos z\,(xe^++ye^\times)\), \(A_0=\cos z\,(pe^++qe^\times)\) with
the York completion, the quartic balance form is

\[
 Q=-\frac{145}{1024}|x|^4-\frac{59}{128}|x|^2|p|^2+\frac{7}{64}|p|^4+\frac{5}{64}(x\cdot p)^2
 =-\frac{145}{1024}|x|^4-\frac{49}{128}|x|^2|p|^2+\frac{7}{64}|p|^4-\frac{5}{64}(x\wedge p)^2,
                                                                    \tag{89.1}
\]

replacing (72.1); the certificate checks rotation invariance at pure
\(\times\) points and at the mixed point \((1,1,1,-1)\).
\end{theorem}

\begin{proof}
The four coefficients are fixed by the pure points and by
\((x,p)=(1,0,1,0)\), \((1,0,0,1)\); the mixed point verifies the
invariant structure.
\end{proof}

\begin{corollary}[Polarization interaction, corrected]
\label{cor:v19v86-kappa}
\(\kappa_{+\times}=-\tfrac{145}{512c_g^2}+2B+2\gamma c_g^2\) with
\(B=-\tfrac{59}{128}\), \(\gamma=\tfrac7{64}\); for \(c_g=3/4\),
\(\kappa_{+\times}=-6001/4608\), replacing \(-271/288\) of V69.  The
attraction between the two gravitons of a mode persists, but the
coupling-independent wedge contribution is now \(2\cdot(-\tfrac5{64})=-\tfrac5{32}\)
in place of \(-\tfrac32\).
\end{corollary}

\begin{proof}
As in Theorem \ref{thm:v19v69-interaction}: the \(n_+n_\times\)
coefficients are \(8v^2\alpha\), \(2(B+D)\) from \(|x|^2|p|^2\),
\(2\gamma c_g^2\) from \(|p|^4\), and \(-2D\) from the wedge; their sum
is \(8v^2\alpha+2B+2\gamma c_g^2\).
\end{proof}

\subsection{Corrected resonant mixing}

\begin{theorem}[York-completed mixing of the \((1,3)\) block]
\label{thm:v19v86-mixing}
The resonant coefficients of the York-completed \((1,3)\) form are
\((c_1,c_2,c_3,c_4)=(-\tfrac{109}{256},-\tfrac1{64},\tfrac{43}{64},\tfrac7{16})\)
for \(a^3b\), \(abp^2\), \(a^2pq\), \(p^3q\), replacing
\((-\tfrac{109}{256},\tfrac{153}{128},\tfrac{109}{128},\tfrac34)\).  With
(75.1), \(m^2=1525225/42467328\); the corrected diagonal energies are
\(d_1=-603061/110592\), \(d_2=9233/36864\), and

\[
 \sin^22\theta=\frac{6498}{1477867}\approx0.0044,                        \tag{89.2}
\]

replacing \(10506528/10510753\).  The two physical sectors are, to
better than one percent in weight, the number states themselves,
split by \(d_1-d_2\approx-5.7\).
\end{theorem}

\begin{proof}
The odd coefficients are half differences of the balance at
\((1,\pm1,1,0)\) and \((1,0,1,\pm1)\) after removing \(c_1\) and
\(c_4\); the diagonals are the Weyl diagonals of the York-completed
even form of V84 at \(l=3\); the formula for \(m\) is Theorem
\ref{thm:v19v72-mixing}, unchanged because it depends only on the
pure-annihilation coefficients.
\end{proof}

The reversal from near-maximal to weak mixing shows how far the
constraint-violating completion had distorted the fourth-order
diagonal: the corrected sector energy of three unit gravitons is
strongly negative, \(e_4(3,0)\approx-7.5\), so the degeneracy of the
free theory is lifted at fourth order by the diagonal terms before the
resonant terms can mix the sectors.

\subsection{Executable certificate}

The verifier

\[
 \texttt{scripts/verify\_sm\_v19\_v86\_york\_polarization\_mixing.py}
                                                                    \tag{89.3}
\]

evaluates (89.1) and its invariance checks, computes
\(\kappa_{+\times}\), extracts the four resonant coefficients of the
\((1,3)\) form, and computes (89.2), recording the superseded values
beside the corrected ones.  It emits

\[
 \texttt{artifacts/sm\_v19\_v86\_york\_polarization\_mixing.json}.   \tag{89.4}
\]

The hashes are

\[
\begin{array}{c|l}
\hbox{object}&\hbox{SHA--256}\\ \hline
\hbox{verifier}&
\texttt{979D86C780A2B2852B0E8DB6D5DE341C2384D20C08F0B4AD215089C0E3249887}\\
\hbox{canonical payload}&
\texttt{01D45B757D983F2F211246AE4027B8AC26DE098A139530E7C0BB0B7268CE9F98}.
\end{array}                                                     \tag{89.5}
\]

\begin{firewall}[Repair complete for the computed cases only]
\label{fire:v19v86}
With V84 and this note every momentum-dependent number quoted since
V67 has a York-completed replacement for the cases computed: one
mode with both polarizations, the collinear pairs \((1,l)\), \(l\leq6\),
and the orthogonal unit pair.  The general pair closed forms in the
momentum sector, the two-mode two-polarization cross terms, and the
mode-number asymptotics of the corrected momentum terms are not
derived.
\end{firewall}

\begin{decision}[V87 acquisition]
\label{dec:v19v86-next}
V87 will return to the three-dimensional question of V82 with the
momentum sector now available: it will compute the York-completed
balance for the orthogonal pair at frequencies \((1,l)\), \(l\leq4\),
including momenta, and compare the growth of the momentum cross
coefficients with the collinear ones, which decides whether the
momentum sector can supply the cancellation of the quartic divergence
found in the configuration sector.
\end{decision}

\begin{assessment}[V86 acquisition]
\label{ass:v19v86}
The repair of the momentum sector is complete for all previously
computed cases: the two-polarization form has wedge coefficient
\(-5/64\), the polarization interaction remains attractive, and the
resonant mixing of the \((1,3)\) block is weak rather than
near-maximal, its earlier value having been an artifact of the
constraint-violating completion.
\end{assessment}

\begin{record}[V86 append record]
\label{rec:v19v86}
V86 is appended to the validated V85 whose installed SHA--256 was
\texttt{35F9666FEBCF3200B6E5A667E78F6D8A095F9ADA8C5AB4936D0DD89B81BC845E}.
After installation only V86 is the active numeric SM note; frozen
archives and unrelated notes are unchanged.
\end{record}


\section{V87: the momentum sector of orthogonal pairs does not cancel the divergence}
\label{sec:v19v87}

V82 found that the configuration self-energy of a graviton from
non-collinear modes grows like the frequency and produces a quartic
divergence in the isotropic sum; V86 completed the momentum sector on
the constrained basis.  This note combines them: the York-completed
fourth-order balance of the orthogonal pairs \((\hat z,l\hat x)\),
\(l\leq4\), with momenta.  The configuration cross coefficient is
exactly \(-\tfrac5{32}(l^2+1)\), and the momentum-sector terms add to
it with the same sign; the mixed amplitude--momentum coupling grows
linearly in \(l\) and is the largest momentum contribution.  The
self-energy summand from an orthogonal mode is therefore negative and
linearly growing, and the isotropic sum diverges like the fourth power
of the cutoff in both sectors.  The averaged fourth-order Hamiltonian
constraint, whatever its sector, does not exhibit the cancellation
that protects the massless graviton in the full theory; the
conclusion is that its fourth-order diagonal is not the physical
graviton energy, and that the protection can only be found once the
dynamics, the reduced Hamiltonian in an internal time, is constructed.

\subsection{Orthogonal pairs with momenta}

\begin{theorem}[York-completed orthogonal pair coefficients]
\label{thm:v19v87-coefficients}
For amplitudes \((a,b)\) and momenta \((p,q)\) of the unit wave along
\(z\) and the wave of frequency \(l\) along \(x\), with polarizations
\(\operatorname{diag}(1,-1,0)\) and \(\operatorname{diag}(0,1,-1)\), the
even cross coefficients of the York-completed balance form are

\[
\begin{array}{c|c|c|c|c}
l&a^2b^2&p^2q^2&a^2q^2&b^2p^2\\ \hline
1&-5/16&7/16&-63/128&-63/128\\
2&-25/32&7/40&-339/800&-111/200\\
3&-25/16&7/80&-1279/3200&-1839/3200\\
4&-85/32&7/136&-3603/9248&-1347/2312
\end{array}                                                            \tag{90.1}
\]

with \(a^2b^2=-\tfrac5{32}(l^2+1)\) and \(p^2q^2=\tfrac{7}{8(l^2+1)}\)
exactly for \(l\leq4\); the pure terms are the single-mode York values
of (86.1).
\end{theorem}

\begin{proof}
Symmetric combinations of engine evaluations, as in Theorem
\ref{thm:v19v84-pairs}; the stated closed patterns are read off the
certified values.
\end{proof}

\subsection{The self-energy summand}

\begin{theorem}[No momentum-sector cancellation]
\label{thm:v19v87-no-cancel}
With \(v_1=1/(2c_g)\), \(v_l=1/(2c_gl)\), the \(n_1\)-linear summand
from the vacuum of the orthogonal mode \(l\) is, by (87.2),

\[
 \sigma^{\perp}_{1l}
 =\underbrace{2c^{ab}v_1v_l}_{\hbox{config}}
 +\underbrace{\frac{c^{pq}}{8v_1v_l}}_{\hbox{momenta}}
 +\underbrace{\frac{c^{aq}v_1}{2v_l}}_{\hbox{mixed}}
 +\underbrace{\frac{c^{bp}v_l}{2v_1}}_{\hbox{mixed}},                    \tag{90.2}
\]

and for \(c_g=3/4\) its values are
\(-\tfrac{2981}{4608}\), \(-\tfrac{4673}{5760}\), \(-\tfrac{74957}{69120}\), \(-\tfrac{3389}{2448}\)
for \(l=1,\dots,4\), with the configuration part
\(-\tfrac5{18},-\tfrac{25}{72},-\tfrac{25}{54},-\tfrac{85}{144}\) and the
mixed \(a^2q^2\) part \(-\tfrac{63}{256},-\tfrac{339}{800},-\tfrac{3837}{6400},-\tfrac{3603}{4624}\),
both growing linearly in \(l\) with the same sign, while the momentum
part \(\tfrac{63}{512},\tfrac{63}{640},\tfrac{189}{2560},\tfrac{63}{1088}\)
is positive and decreasing.  Hence the isotropic sum of the
York-completed self-energy diverges like the fourth power of the
cutoff, and the momentum sector increases rather than cancels the
configuration divergence.
\end{theorem}

\begin{proof}
Insert (90.1) into (90.2); the mixed term \(c^{aq}v_1/(2v_l)=c^{aq}l/2\)
with \(c^{aq}\) of order one is linear in \(l\).
\end{proof}

\subsection{Interpretation}

\begin{proposition}[Sector labels are not physical energies at fourth order]
\label{prop:v19v87-interpretation}
The fourth-order diagonal of the averaged Hamiltonian constraint on
a fixed slice determines the branch labels \(\tau_e\) of the states, and
its cutoff dependence has been shown to contain a graviton mass term
that no sector of the slice constraint cancels.  In the full theory
the Hamiltonian constraint generates the dynamics and the physical
energy of a graviton is defined through the reduced Hamiltonian in an
internal time, in which the massless pole is protected.  The
divergence found here is therefore a property of the slice
labelling, gauge dependent at fourth order, and not a physical
graviton mass; its removal requires the dynamical framework, which the
programme has not constructed.
\end{proposition}

\begin{proof}
This is an interpretive statement: the computations of V76--V87 are
exact for the slice constraint; the standard protection argument
concerns the on-shell two-point function, which involves the
evolution; no computation here contradicts it.
\end{proof}

\subsection{Executable certificate}

The verifier

\[
 \texttt{scripts/verify\_sm\_v19\_v87\_orthogonal\_momentum\_growth.py}
                                                                    \tag{90.3}
\]

evaluates the York-completed balance for the four orthogonal pairs,
checks the single-mode terms and the configuration cross term
\(-\tfrac5{32}(l^2+1)\), and computes the summands (90.2) with their
parts.  It emits

\[
 \texttt{artifacts/sm\_v19\_v87\_orthogonal\_momentum\_growth.json}. \tag{90.4}
\]

The hashes are

\[
\begin{array}{c|l}
\hbox{object}&\hbox{SHA--256}\\ \hline
\hbox{verifier}&
\texttt{FA10503933E17BD2A448118BA6A821590AD106F67725CD568DB56351209637E5}\\
\hbox{canonical payload}&
\texttt{87BB43CE6F697FB7EBFF292DFAF899428414C8BBB574D3277D832E1B286A1A4A}.
\end{array}                                                     \tag{90.5}
\]

\begin{firewall}[Four pairs, one polarization pattern, slice constraint]
\label{fire:v19v87}
Closed forms in \(l\) for the momentum coefficients are read off four
values and not proved; only one polarization pattern per pair is
used; and the interpretation of Proposition
\ref{prop:v19v87-interpretation} is a statement about what has not
been computed, namely the dynamics.
\end{firewall}

\begin{decision}[V88 acquisition]
\label{dec:v19v87-next}
V88 will go upstream to the dynamics in the only setting where the
programme has exact tools: the plane-symmetric class, where the
Einstein evolution reduces to ordinary differential equations in
\(z\) and \(t\).  It will derive the exact evolution equations for
\(G_{ab}(z,t)\) in the gauge of V79 with unit lapse and zero shift,
expand them to second order about the standing wave, and identify the
reduced Hamiltonian whose spectrum, rather than the slice labels, is
the physical energy at fourth order.
\end{decision}

\begin{assessment}[V87 acquisition]
\label{ass:v19v87}
With the momentum sector included on the constrained basis, the
self-energy summand from orthogonal modes is negative and grows
linearly, so the isotropic sum diverges quartically in both sectors
and no cancellation of the graviton mass term occurs in the averaged
slice constraint.  The programme concludes that the fourth-order
sector labels are gauge-dependent slice quantities and that the
massless-graviton protection lives in the dynamics, which is now the
declared next target.
\end{assessment}

\begin{record}[V87 append record]
\label{rec:v19v87}
V87 is appended to the validated V86 whose installed SHA--256 was
\texttt{CA95C12451CDD05132E61A065516F4EE0CD3001747A973A06664BA94FD05D325}.
After installation only V87 is the active numeric SM note; frozen
archives and unrelated notes are unchanged.
\end{record}


\section{V88: the first exact step of nonlinear dynamics from the time-symmetric standing wave}
\label{sec:v19v88}

V87 concluded that the physical graviton energy must be sought in the
dynamics.  This note takes the first exact step: the second-order
initial acceleration of every component of the slice metric under the
Einstein evolution with unit lapse and zero shift, starting from the
balanced time-symmetric standing wave of V41.  At a time-symmetric
slice the acceleration is \(-2(\operatorname{Ric}-\Lambda g)\), so the
computation needs the Ricci tensor of the perturbed metric to second
order, which the three-dimensional engine supplies exactly, together
with the Ricci tensor of the conformal factor.  The result shows the
backreaction anisotropically: the transverse directions accelerate
inward and the longitudinal direction outward, with covariant trace
\(2\Lambda\), the turnaround of V43, and with a longitudinal component
that immediately leaves the spatial gauge of V79.  This is the
starting point for the reduced dynamics.

\subsection{Second-order Ricci tensor}

\begin{proposition}[Ricci tensor of the standing wave to second order]
\label{prop:v19v88-ricci}
For \(\gamma=\delta+\varepsilon h\), \(h=\cos z\operatorname{diag}(1,-1,0)\):
\(\operatorname{Ric}_1=\tfrac12h\), and the second-order part satisfies
\(\delta^{ij}(\operatorname{Ric}_2)_{ij}-h^{ij}(\operatorname{Ric}_1)_{ij}=R_2=-\tfrac14-\tfrac74\cos2z\),
in agreement with V41.  The conformal factor \(\psi=1+\varepsilon^2w_2\),
\(w_2=\tfrac7{128}\cos2z\), contributes at second order
\(\operatorname{Ric}^{\rm conf}_{ij}=-\partial_i\partial_j\phi-\delta_{ij}\Delta\phi\),
\(\phi=2w_2\).
\end{proposition}

\begin{proof}
The engine computes the Ricci tensor from the Christoffel symbols
exactly; the first-order and trace checks are certified.  The
conformal formula is the linearization of the three-dimensional
conformal transformation of the Ricci tensor.
\end{proof}

\subsection{Initial acceleration}

\begin{theorem}[Second-order initial acceleration]
\label{thm:v19v88-acceleration}
Let \(g_*=\psi^4\gamma\) be the balanced datum with \(K=0\) and
\(\Lambda=-\tfrac{\varepsilon^2}8+O(\varepsilon^4)\), evolved with unit
lapse and zero shift.  Then \(\ddot g=-2(\operatorname{Ric}(g_*)-\Lambda g_*)\)
at \(t=0\), whose first-order part is \(-h\), and whose second-order
part is

\[
 \ddot g^{(2)}_{xx}=\ddot g^{(2)}_{yy}=-\frac34-\frac38\cos2z,
 \qquad
 \ddot g^{(2)}_{zz}=\frac14-\frac14\cos2z,
 \qquad
 \ddot g^{(2)}_{ij}=0\ (i\neq j).                                        \tag{91.1}
\]

The covariant trace satisfies
\(\langle\delta^{ij}\ddot g^{(2)}_{ij}\rangle+\langle|h|^2\rangle=-\tfrac54+1=2\Lambda_2\),
so the volume accelerates at the rate \(\Lambda_2\operatorname{Vol}<0\):
the plane is a turnaround.  The longitudinal component
\(\ddot g_{zz}^{(2)}\neq0\) shows that the evolution leaves the gauge
\(g_{zz}=1\) of V79 at second order.
\end{theorem}

\begin{proof}
Insert Proposition \ref{prop:v19v88-ricci} into the evolution
equation (46.1) at \(K=0\); \(g_*\) enters \(\Lambda g_*\) at second
order through \(\Lambda_2\delta\) only.  The trace identity uses
\(g^{ij}=\delta^{ij}-\varepsilon h^{ij}+\dots\) and
\(\ddot g^{(1)}=-h\).  Values certified.
\end{proof}

The mean parts of (91.1) are the beginning of an anisotropic
Kasner-like backreaction: contraction of the two transverse directions
at rate \(-\tfrac34\varepsilon^2\) and expansion of the longitudinal
direction at rate \(\tfrac14\varepsilon^2\), the same signature as the
exact Kasner exponents \((\tfrac23,\tfrac23,-\tfrac13)\) of V39 with
the roles of the directions determined by the wave's polarization
plane.

\subsection{Executable certificate}

The verifier

\[
 \texttt{scripts/verify\_sm\_v19\_v88\_second\_order\_acceleration.py}
                                                                    \tag{91.2}
\]

computes the Ricci tensor series with the engine, checks
\(\operatorname{Ric}_1=h/2\) and the \(R_2\) consistency, adds the
conformal contribution, evaluates (91.1), and verifies the covariant
trace identity.  It emits

\[
 \texttt{artifacts/sm\_v19\_v88\_second\_order\_acceleration.json}.  \tag{91.3}
\]

The hashes are

\[
\begin{array}{c|l}
\hbox{object}&\hbox{SHA--256}\\ \hline
\hbox{verifier}&
\texttt{202320C221DA1D92CA5EF845F2CFB6A0233B627905B4CE8372DD09151DBBECD8}\\
\hbox{canonical payload}&
\texttt{518F6494FF0254FC99D4A58F4B87B47F1C258A573565600DADBC058878CBE360}.
\end{array}                                                     \tag{91.4}
\]

\begin{firewall}[Initial data of the dynamics only]
\label{fire:v19v88}
Only the second time derivative at the time-symmetric slice is
computed; the evolution over a finite time, the secular growth, the
reduced Hamiltonian, and its spectrum are not.  Unit lapse and zero
shift are a gauge choice whose relation to CMC time is not analyzed.
\end{firewall}

\begin{decision}[V89 acquisition]
\label{dec:v19v88-next}
V89 will integrate the plane-symmetric evolution in closed form to
second order in the amplitude over a finite time, using a
trigonometric-polynomial engine in \(z\) and \(t\) with polynomial
secular terms, verifying \(\dot\tau=|K|^2-\Lambda\) along the solution
and extracting the second-order energy exchange between the wave and
the homogeneous expansion, which is the classical content of the
reduced Hamiltonian at this order.  The V90 audit and ten-version
sweep follow.
\end{decision}

\begin{assessment}[V88 acquisition]
\label{ass:v19v88}
The first exact step of nonlinear dynamics is taken: the second-order
initial acceleration of the standing-wave datum is anisotropic,
transverse inward and longitudinal outward, with covariant trace
\(2\Lambda\), and it exits the spatial gauge of the plane class.  The
dynamical programme announced in V87 has its initial data.
\end{assessment}

\begin{record}[V88 append record]
\label{rec:v19v88}
V88 is appended to the validated V87 whose installed SHA--256 was
\texttt{FB9082D81A79F3EA48279358A0B81AE6800467D1ED41C44EABD3EF64376CAD63}.
After installation only V88 is the active numeric SM note; frozen
archives and unrelated notes are unchanged.
\end{record}


\section{V89: closed-form second-order evolution of the standing wave and its effective stress}
\label{sec:v19v89}

V88 computed the initial acceleration of the balanced standing wave.
This note integrates the plane-symmetric Einstein evolution with unit
lapse and zero shift in closed form through second order in the
amplitude, for all times, with an exact trigonometric-polynomial
engine in \(z\) and \(t\) that carries the secular terms as
polynomials in \(t\).  The solution is a finite expression: the
transverse metric contracts quadratically in time and breathes at
twice the wave frequency, the longitudinal metric breathes without
secular drift, and the mean curvature grows linearly with the exact
rate \(|K|^2-\Lambda\) of V43.  The secular parts satisfy the
Bianchi--I equations with an effective null-dust source along the
propagation axis whose energy density is exactly \(-\Lambda/8\pi\):
the second-order dynamics reproduces the Isaacson stress of the wave,
and the negative cosmological constant of V39--V41 is precisely the
one that cancels it.  This is the first exact statement of energy
exchange between the wave and the homogeneous background in the
programme, and it is the classical content of the reduced Hamiltonian
at second order.

\subsection{Setting}

The metric is \(g=\delta+\varepsilon g_1+\varepsilon^2g_2\) with
\(g_1=\cos z\cos t\operatorname{diag}(1,-1,0)\) the linearized
standing wave (the V41 datum with its time dependence), and
\(g_2=\operatorname{diag}(b_2,c_2,a_2)(z,t)\) with the initial data
\(g_2(0)=4w_2\delta=\tfrac7{32}\cos2z\,\delta\), \(\dot g_2(0)=0\)
fixed by the balanced conformal factor of V41 and time symmetry.  The
evolution equations \(\ddot g_{ij}=-2(\operatorname{Ric}_{ij}-\Lambda g_{ij})-2K_{ik}K^k{}_j+\dots\)
with unit lapse and zero shift, expanded to second order, read

\[
 \ddot b_2-b_2''=-2S_{xx},\qquad
 \ddot c_2-c_2''=-2S_{yy},\qquad
 \ddot a_2=(b_2+c_2)''-2S_{zz},                                       \tag{92.1}
\]

with the sources
\(S_{ij}=\cos^2t\operatorname{Ric}_2[h]_{ij}-2(K_1K_1)_{ij}-\Lambda_2\delta_{ij}\),
\(h=\cos z\operatorname{diag}(1,-1,0)\), \(K_1=\tfrac12\cos z\sin t\operatorname{diag}(1,-1,0)\),
\(\Lambda_2=-\tfrac18\), and primes denoting \(\partial_z\).

\subsection{The closed-form solution}

\begin{theorem}[Second-order plane-symmetric evolution]
\label{thm:v19v89-solution}
The unique solution of (92.1) with the stated initial data is

\[
\begin{aligned}
 b_2=c_2&=-\frac18-\frac{t^2}8+\frac{\cos2t}8+\frac{\cos2z}8+\frac3{32}\cos2z\cos2t,\\
 a_2&=\frac1{16}-\frac{\cos2t}{16}+\frac7{32}\cos2z-\frac{t^2}8\cos2z.
\end{aligned}                                                          \tag{92.2}
\]

The second-order mean curvature
\(\tau_2=-\tfrac12\operatorname{tr}\dot g_2+\tfrac12g_1^{ij}\dot g_{1,ij}\)
is

\[
 \tau_2=\frac t4-\frac{\sin2t}{16}+\cos2z\Bigl(\frac t8-\frac{\sin2t}{16}\Bigr),
                                                                    \tag{92.3}
\]

and it satisfies the exact rate law
\(\dot\tau_2=|K_1|^2-\Lambda_2=\tfrac12\cos^2z\sin^2t+\tfrac18\)
identically in \(z\) and \(t\).  The \(z\)-averaged second-order
Hamiltonian constraint
\(\cos^2t\langle R_2[h]\rangle-\langle|K_1|^2\rangle=2\Lambda_2\)
holds for all \(t\).
\end{theorem}

\begin{proof}
The engine represents fields as finite sums
\(c_{kmn}e^{ikz}e^{imt}t^n\) with Gaussian-rational coefficients and
solves each \(z\)-mode of (92.1) exactly: non-resonant terms by a
polynomial ansatz, resonant terms \(e^{\pm ikt}t^n\) by integrating a
first-order recursion, and the \(k=0\) modes by double integration;
the homogeneous part is fixed by the initial data.  The wave-equation
residuals of (92.1) vanish identically, the initial accelerations
reproduce (91.1), and (92.2)--(92.3) and the two identities are
verified coefficient by coefficient.
\end{proof}

\subsection{Effective stress and the meaning of \(\Lambda\)}

\begin{theorem}[Isaacson stress from the exact second-order dynamics]
\label{thm:v19v89-isaacson}
The \(z\)-averaged accelerations are
\(\langle\ddot b_2\rangle=-\tfrac14-\tfrac12\cos2t\) and
\(\langle\ddot a_2\rangle=\tfrac14\cos2t\).  Their time averages,
\(2\Lambda_2\) transversally and \(0\) longitudinally, are the
linearized Bianchi--I equations
\(\ddot m_{ij}=2\Lambda_2\delta_{ij}+16\pi(T_{ij}-\tfrac12T\delta_{ij})\)
for an effective stress with \(p_x=p_y=0\) and

\[
 8\pi\rho_{\rm eff}=8\pi p_z=-\Lambda_2=\frac18,                        \tag{92.4}
\]

that is, null dust along the propagation axis with energy density
equal to the negative of the cosmological constant.  The averaged
Hamiltonian constraint of Theorem \ref{thm:v19v89-solution} is
\(2\Lambda=-16\pi\rho_{\rm eff}\): the balance condition of V41 is
exactly the cancellation of the wave's Isaacson energy by \(\Lambda\).
\end{theorem}

\begin{proof}
The mean second-order Ricci tensor of a homogeneous perturbation
\(m_{ij}(t)\) in synchronous gauge is \(\tfrac12\ddot m_{ij}\), so the
mean of \(\operatorname{Ric}=\Lambda g\) at second order gives
\(\tfrac12\ddot m_{ij}=\Lambda_2\delta_{ij}-\langle\operatorname{Ric}^{(2)}_{ij}[g_1,g_1]\rangle\),
and the wave's contribution is identified with
\(8\pi(T_{ij}-\tfrac12T\delta_{ij})\).  With \(T=-\rho+p_z\) and the
certified averages, \(T_{xx}-\tfrac12T=\tfrac12(\rho-p_z)=0\) and
\(T_{zz}-\tfrac12T=\tfrac12(\rho+p_z)=\rho\), giving (92.4).  The
Hamiltonian statement is the \(z\)-mean identity of Theorem
\ref{thm:v19v89-solution}, whose left side is \(-\tfrac14\) for all
\(t\) because the potential energy \(\cos^2t\langle R_2\rangle\) and
the kinetic energy \(\langle|K_1|^2\rangle\) exchange with constant
sum.
\end{proof}

\begin{proposition}[Secular structure]
\label{prop:v19v89-secular}
The secular parts of (92.2) are transverse,
\(\langle b_2\rangle_{\rm sec}=\langle c_2\rangle_{\rm sec}=-t^2/8\),
with no longitudinal secular term; the mean volume element is
\(\langle\sqrt g\rangle=1-\varepsilon^2\bigl(\tfrac7{32}+\tfrac{t^2}8+\tfrac1{32}\cos2t\bigr)+O(\varepsilon^4)\),
and \(\langle\tau\rangle=\varepsilon^2(\tfrac t4-\tfrac1{16}\sin2t)\).
The perturbative solution is uniform only for \(t\ll\varepsilon^{-1}\),
after which the mean contraction is of order one.
\end{proposition}

\begin{proof}
Read off (92.2)--(92.3); the volume element uses
\(\sqrt{\det(1+\varepsilon A)}=1+\tfrac\varepsilon2\operatorname{tr}A+\varepsilon^2(\tfrac18(\operatorname{tr}A)^2-\tfrac14\operatorname{tr}A^2)\)
with \(\langle|g_1|^2\rangle=\cos^2t\).
\end{proof}

\subsection{Executable certificate}

The verifier

\[
 \texttt{scripts/verify\_sm\_v19\_v89\_second\_order\_evolution.py}
                                                                    \tag{92.5}
\]

builds the sources from the engine's second-order Ricci tensor,
integrates (92.1) exactly, checks the residuals, the V88 initial
accelerations, the closed forms (92.2)--(92.3), the rate law, the
averaged Hamiltonian constraint, and the effective-stress identities
(92.4).  It emits

\[
 \texttt{artifacts/sm\_v19\_v89\_second\_order\_evolution.json}.     \tag{92.6}
\]

The hashes are

\[
\begin{array}{c|l}
\hbox{object}&\hbox{SHA--256}\\ \hline
\hbox{verifier}&
\texttt{AD525DAD770FC031ECD694E0DA1315EE4E61E69387193D84CFC23941427A2EC8}\\
\hbox{canonical payload}&
\texttt{2073FC4B02DAFA38AFAF997D66329B91171153B89DF364E3DF15BF959BFC135A}.
\end{array}                                                     \tag{92.7}
\]

\begin{firewall}[Second order, one mode, synchronous gauge]
\label{fire:v19v89}
The solution is second order in the amplitude and valid for
\(t\ll\varepsilon^{-1}\); the secular terms signal the breakdown of
the naive expansion and the need for a multiple-scale or reduced
Hamiltonian treatment, which is not done.  Only the single plane
standing wave is evolved, in synchronous gauge; the identification of
the effective stress is the standard one and carries no information
about the fourth-order graviton mass question of V78--V87.
\end{firewall}

\begin{decision}[V90 acquisition]
\label{dec:v19v89-next}
V90 is the compulsory companion audit and the ten-version sweep
V80--V89.  V91 will then use the secular structure of (92.2) to set
up the two-scale reduced dynamics: a slowly contracting Bianchi--I
background driven by the wave's Isaacson stress, with the wave
amplitude and frequency adiabatically renormalized, and compare the
resulting adiabatic invariant with the branch labels of V52--V62.
\end{decision}

\begin{assessment}[V89 acquisition]
\label{ass:v19v89}
The second-order plane-symmetric evolution of the balanced standing
wave is solved in closed form, with the exact rate law for the mean
curvature and the exact averaged Hamiltonian constraint along the
solution.  The dynamics exhibits the wave's Isaacson stress as null
dust with \(8\pi\rho_{\rm eff}=-\Lambda\), identifying the negative
cosmological constant of the torus as the exact counterweight of the
wave energy, and it shows a transverse quadratic contraction with no
longitudinal secular drift.
\end{assessment}

\begin{record}[V89 append record]
\label{rec:v19v89}
V89 is appended to the validated V88 whose installed SHA--256 was
\texttt{59607A9849A893B494A82D7F8ACD908C370A2DD940805F40809C45BA5A1E1046}.
After installation only V89 is the active numeric SM note; frozen
archives and unrelated notes are unchanged.
\end{record}


\section{V90: compulsory companion audit and ten-version sweep after the first exact dynamics}
\label{sec:v19v90}

This is the required audit following V85 and the ten-version sweep
of V80--V89 required every tenth version.  Every active companion
note present in the shared folder was inspected read-only, the
frozen archives of the earlier cycles were searched for prior
versions of the V89 result, and the direction for V91--V99 is fixed
on that basis.

\subsection{Companion snapshot}

\begin{record}[V90 companion snapshot]
\label{rec:v19v90-snapshot}
The files present were

\[
\begin{array}{c|r|r}
\hbox{file and SHA--256}&\hbox{bytes}&\hbox{lines}\\ \hline
\texttt{Renewal\_NS\_Stability\_Research\_Notes\_V26.tex}
&278283&5319\\
\multicolumn{3}{l}{
\texttt{82C887F8C2C69E4F28FAD7C3DC8DE5B9099CB5A4BF431EBA1FFF3E91935978AD}}
\\
\texttt{Renewal\_GRH\_Cofinal\_Visibility\_Attack\_V56.tex}
&546654&16487\\
\multicolumn{3}{l}{
\texttt{33DB47D83CE375E8D3A3E234C6A537D1E89F0E5486F2FFAF0DDFB638ED80E2A4}}
\\
\texttt{Renewal\_Yang\_Mills\_Mass\_Gap\_Research\_Notes\_V34.tex}
&489911&12282\\
\multicolumn{3}{l}{
\texttt{3CF829F57455C43D3CE54387782C3C60EB3497F1B680A11DDCAED13306D8C469}}
\\
\texttt{Renewal\_YM\_Parallel\_Research\_Notes\_V5.tex}
&54174&1351\\
\multicolumn{3}{l}{
\texttt{306F1BB84CFA8A8D47EA569EC17E9545F9867C8D32B548EC9D9C444B4F3C0512}}.
\end{array}                                                     \tag{93.1}
\]

Since V85 the GRH note advanced from V55 to V56 (the Pick-disk
condition and its \(y\)-expansion, two of five coefficient
certificates) and the Yang--Mills mass-gap note reappeared at V34
(symmetrized dyadic blocking with exact permutation covariance, an
exact alias recursion for the blocked propagator, and the Gaussian
anomaly card restated on the fixed-point tower).  The NS note and the
YM parallel note are unchanged since V80.
\end{record}

\begin{decision}[V90 audit decision]
\label{dec:v19v90-audit}
No companion theorem, numerical constant, or executable artifact is
imported.  The GRH advance is a statement about polynomial
inequalities on a disk; the Yang--Mills advance is a lattice blocking
of a Gaussian field in Euclidean signature without constraints.
Neither touches the Lorentzian constrained dynamics that is the
current subject of this note.
\end{decision}

\begin{proof}
The V56 GRH section proves exact factorizations of \(y\)-coefficients
of a Pick expression and certifies two sums of squares; the V34
Yang--Mills section proves a blocking covariance and a linear
recursion for a free propagator.  The objects of V88--V89 are the
second-order Einstein evolution equations of a constrained
gravitational wave on a torus; no shared operator, hypothesis, or
constant appears.
\end{proof}

\subsection{Ten-version sweep V80--V89}

\begin{proposition}[What V80--V89 established]
\label{prop:v19v90-sweep}
The block V80--V89 converted the programme from a plane-symmetric
constraint calculus into a three-dimensional constrained calculus and
then into dynamics.  Its exact results are:

\begin{enumerate}
\item the three-dimensional curvature engine and the orthogonal
second-order cross term \(2ab\cos x\cos z\) without frequency
enhancement (V81);
\item the fourth-order orthogonal cross coefficient
\(-\tfrac5{64}(l^2+1)\), hence a quartically divergent isotropic
configuration self-energy with no angular cancellation (V82);
\item the York-projected momentum sector in three dimensions, and the
erratum that the plane-class momentum completion of V67--V69 violated
the momentum constraint, superseding every momentum-dependent number
of V67--V78 (V83);
\item the corrected collinear pair forms with \(p^2q^2\) values
\(\tfrac{35}{18},\tfrac{35}{64},\tfrac{119}{450},\tfrac{91}{576},\tfrac{37}{350}\),
the amplitude--momentum coupling, and the intrinsic cutoff four (V84);
\item the York polarization wedge \(-\tfrac5{64}\) and the weak
\((1,3)\) mixing \(\sin^22\theta=6498/1477867\) (V86);
\item the orthogonal momentum sector \(p^2q^2=7/(8(l^2+1))\) with
linearly growing mixed coupling, so that no sector of the averaged
slice constraint cancels the graviton mass term (V87);
\item the exact second-order initial acceleration of the standing
wave, transverse inward and longitudinal outward with covariant trace
\(2\Lambda\) (V88);
\item the closed-form second-order evolution (92.2)--(92.3), the exact
rate law \(\dot\tau=|K|^2-\Lambda\) along it, and the identification
of the wave's Isaacson stress as null dust with
\(8\pi\rho_{\rm eff}=-\Lambda\) (V89).
\end{enumerate}

The frozen archives of the eighteen earlier cycles contain no
Isaacson or null-dust identification and no second-order integrated
evolution; the V89 result is new to the programme.
\end{proposition}

\begin{proof}
Items are the certified theorems of the cited sections; the archive
statement is a full-text search of the eighteen frozen
\texttt{V100\_f}\(k\) files for the two terms and for an integrated
second-order solution.
\end{proof}

\begin{proposition}[Direction after the sweep]
\label{prop:v19v90-direction}
The sweep closes the fourth-order slice programme: the averaged slice
constraint has been computed exactly in every sector for the pairs
within reach, and it does not protect the massless graviton.  The
open structural gaps are unchanged in kind but reordered in priority:
the reduced dynamics in an internal time is now first, because V89
shows that at second order it is an exactly solvable two-scale system
(a Bianchi--I background driven by the wave's null-dust stress, with
the wave riding on it); the number-changing terms, the mode-number
limit, and the Standard Model coupling follow.  V91 will construct
the two-scale reduction and its adiabatic invariant; V92--V94 will
carry the wave to fourth order on the contracting background, where
the frequency renormalization of V77 becomes a time dependence; the
V95 audit follows.
\end{proposition}

\begin{proof}
This is a decision, justified by Theorems \ref{thm:v19v87-no-cancel}
and \ref{thm:v19v89-isaacson}: the first excludes the slice
constraint as the carrier of the protection, the second exhibits the
dynamics as tractable.
\end{proof}

\begin{firewall}[Audit and sweep only]
\label{fire:v19v90}
No computation is added in V90.  The reordering of priorities is a
judgement about tractability, not a theorem; the four structural
gaps of V80 remain open and the Standard Model has not entered the
dynamics.
\end{firewall}

\begin{decision}[V91 acquisition]
\label{dec:v19v90-next}
V91 constructs the two-scale reduction of the plane-symmetric system:
a homogeneous Bianchi--I background \(\operatorname{diag}(A,A,C)(T)\)
in the slow time \(T=\varepsilon^2t\) coupled to a wave of slowly
varying amplitude and frequency, derives the averaged equations from
the exact second-order solution (92.2), identifies the adiabatic
invariant of the wave (the ratio of its energy to its frequency in
the contracting background), and checks the reduced equations against
(92.2)--(92.3) to second order.
\end{decision}

\begin{assessment}[V90 acquisition]
\label{ass:v19v90}
The audit is current and imports nothing.  The sweep records that
V80--V89 completed and then closed the fourth-order slice programme,
corrected its own momentum sector, and opened the dynamics with an
exact second-order solution whose Isaacson stress fixes the meaning
of the negative cosmological constant.  The direction for V91--V99 is
the two-scale reduced dynamics.
\end{assessment}

\begin{record}[V90 append record]
\label{rec:v19v90}
V90 is appended to the validated V89 whose installed SHA--256 was
\texttt{40E2A97FF014E5A30C86045EDEA32F6CD3CFCFDC37ECD02BE657930DCD3C5B4F}.
After installation only V90 is the active numeric SM note; the
companions, frozen archives, and all unrelated notes are unchanged.  The
next compulsory companion audit is V95 and the next ten-version sweep
is V100, which will be frozen as \texttt{V100\_f19}.
\end{record}


\section{V91: the exact plane-symmetric system, an independent confirmation of V89, and the two-scale reduction}
\label{sec:v19v91}

This note does three things.  It derives the exact plane-symmetric
Einstein equations with cosmological constant in the diagonal
exponential variables, symbolically and without perturbation theory,
and records the constraints and the wave equation in that form.  It
uses them to confirm that the closed-form second-order solution of V89
satisfies the full Einstein equations, a check independent of the
perturbative engine that produced it.  And it constructs the
two-scale reduction announced in V90: a homogeneous Bianchi--I
background in the slow time \(T=\varepsilon t\), driven by the wave's
null-dust stress and by \(\Lambda\), carrying a wave whose amplitude
and frequency follow the adiabatic invariant.  The reduced system is
solved as an exact rational series to order ten, its constraint is
preserved, it reproduces the secular parts of V89, and it is proved
to collapse transversally within slow time \(\sqrt{128/3}\).  The
slow time is \(\varepsilon t\), not \(\varepsilon^2t\) as anticipated
in the V90 decision.

\subsection{The exact plane-symmetric system}

\begin{proposition}[Plane-symmetric Einstein equations]
\label{prop:v19v91-system}
For \(ds^2=-dt^2+e^{2a}(e^{2b}dx^2+e^{-2b}dy^2)+e^{2c}dz^2\) with
\(a,b,c\) functions of \((z,t)\), the equations
\(\operatorname{Ric}=\Lambda g\) are equivalent to

\[
\begin{aligned}
 &a_t^2+2a_tc_t-b_t^2-e^{-2c}\bigl(3a_z^2-2a_zc_z+2a_{zz}+b_z^2\bigr)=\Lambda,\\
 &a_{tz}+a_ta_z-a_zc_t+b_tb_z=0,\\
 &b_{tt}+b_t(2a_t+c_t)-e^{-2c}\bigl(b_{zz}+b_z(2a_z-c_z)\bigr)=0,
\end{aligned}                                                          \tag{94.1}
\]

together with the two diagonal evolution equations
\(R^x{}_x+R^y{}_y=2\Lambda\) and \(R^z{}_z=\Lambda\).  In the
homogeneous case \(b=0\), \(a=u(t)\), \(c=w(t)\), these reduce to
\(R^x{}_x=\ddot u+\dot u(2\dot u+\dot w)\),
\(R^z{}_z=\ddot w+\dot w(2\dot u+\dot w)\), and
\(G_{tt}=\dot u^2+2\dot u\dot w\).
\end{proposition}

\begin{proof}
The Ricci tensor of the metric is computed symbolically from the
Christoffel symbols; the three displayed forms are the certified
simplifications of \(G_{tt}-\Lambda\), \(-\tfrac12R_{tz}\), and
\(\tfrac12(R^x{}_x-R^y{}_y)\), and the homogeneous forms are their
certified specializations.
\end{proof}

\subsection{Independent confirmation of V89}

\begin{theorem}[V89 solves the full equations to second order]
\label{thm:v19v91-confirm}
With \(b=\tfrac\varepsilon2\cos z\cos t\),
\(a=\tfrac{\varepsilon^2}2(b_2-\tfrac12\cos^2z\cos^2t)\),
\(c=\tfrac{\varepsilon^2}2a_2\) for the closed forms (92.2), and
\(\Lambda=-\varepsilon^2/8\), every component of
\(\operatorname{Ric}-\Lambda g\) vanishes through order
\(\varepsilon^2\).
\end{theorem}

\begin{proof}
Direct symbolic substitution into the exact Ricci tensor of
Proposition \ref{prop:v19v91-system} and series expansion in
\(\varepsilon\); the five independent components \(tt,tz,xx,yy,zz\)
are certified to vanish.  This confirms the perturbative engine of
V88--V89 and the sources (92.1) by an independent route.
\end{proof}

\subsection{The two-scale reduction}

\begin{theorem}[Reduced Bianchi--I system with the wave's null dust]
\label{thm:v19v91-reduced}
Let \(H_\perp=\dot u\), \(H_\parallel=\dot w\),
\(\theta=2H_\perp+H_\parallel\), and let
\(8\pi\rho=\langle b_t^2+e^{-2c}b_z^2\rangle\) be the \(z\)- and
phase-averaged wave energy of the constraint in (94.1); for
\(b=\tfrac\varepsilon2\cos z\cos t\) it equals \(\varepsilon^2/8=-\Lambda\).
The averaged equations are

\[
 \dot H_\perp+H_\perp\theta=\Lambda,\qquad
 \dot H_\parallel+H_\parallel\theta=\Lambda+8\pi\rho,\qquad
 H_\perp^2+2H_\perp H_\parallel=\Lambda+8\pi\rho,                      \tag{94.2}
\]

with \(\rho\propto(a_\perp a_\parallel)^{-2}\), which is the
conservation law of null dust along \(z\) and also the consequence of
the adiabatic invariant \(V\omega A^2\) of the wave equation in (94.1)
on the homogeneous background, \(\omega=e^{-w}=1/a_\parallel\),
\(A\propto1/a_\perp\).  In the slow time \(T=\varepsilon t\), with
\(\Lambda=-\tfrac18\) and \(8\pi\rho(0)=\tfrac18\), the solution with
\(u=w=\dot u=\dot w=0\) at \(T=0\) is the exact rational series

\[
\begin{aligned}
 a_\perp^2&=1-\frac{T^2}8+\frac{T^4}{384}+\frac{T^6}{46080}+\frac{13T^8}{10321920}+\frac{191T^{10}}{3715891200}+\dots,\\
 a_\parallel^2&=1+\frac{T^4}{384}+\frac{T^6}{5760}+\frac{149T^8}{10321920}+\frac{277T^{10}}{232243200}+\dots,\\
 \langle\tau\rangle/\varepsilon&=\frac T4+\frac{T^3}{64}+\frac{T^5}{768}+\frac{7T^7}{61440}+\frac{233T^9}{23224320}+\dots,
\end{aligned}                                                          \tag{94.3}
\]

the constraint of (94.2) is preserved to the order computed, and the
orders \(T^0,T^1,T^2\) coincide with the secular parts of
(92.2)--(92.3).
\end{theorem}

\begin{proof}
The averaged equations are \(R^x{}_x=\Lambda\),
\(R^z{}_z=\Lambda+8\pi\rho\), \(G_{tt}=\Lambda+8\pi\rho\) in the
homogeneous forms of Proposition \ref{prop:v19v91-system}, with the
effective stress of Theorem \ref{thm:v19v89-isaacson}
(\(p_x=p_y=0\), \(p_z=\rho\)).  For the wave equation of (94.1) on the
background, \(\ddot b+\dot b\,d\ln V/dt+e^{-2w}b=0\) for the mode
\(\cos z\), the WKB invariant is \(A^2\omega V\) with \(V=e^{2u+w}\),
giving \(A\propto1/a_\perp\), and
\(8\pi\rho=\tfrac12\omega^2A^2\propto(a_\perp a_\parallel)^{-2}\).  The
series is solved order by order in exact rational arithmetic; the
constraint, the conservation law, and the invariant are certified
identically to order ten.  The comparison with V89 is the certified
check \(a_\perp^2=1-T^2/8+O(T^4)\), \(a_\parallel^2=1+O(T^4)\),
\(\langle\tau\rangle=\varepsilon T/4+O(T^3)\).
\end{proof}

\begin{theorem}[Transverse collapse of the reduced system]
\label{thm:v19v91-collapse}
Along any solution of (94.2) with \(\rho\geq0\), \(\Lambda<0\), and
\(H_\perp(0)=0\),

\[
 \dot H_\perp=\frac\Lambda2-\frac32H_\perp^2-4\pi\rho\leq\min\Bigl(\frac\Lambda2,\,-\frac32H_\perp^2\Bigr),
                                                                    \tag{94.4}
\]

so \(H_\perp(t)\leq\Lambda t/2\) and the maximal solution has finite
lifetime \(t_*\leq4/\sqrt{-3\Lambda}\), at which
\(H_\perp\to-\infty\) and \(a_\perp\to0\) at least as fast as
\((t_*-t)^{2/3}\), unless the solution ceases earlier.  For the
balanced datum \(\Lambda=-\varepsilon^2/8\) this is
\(T_*\leq\sqrt{128/3}\approx6.53\).
\end{theorem}

\begin{proof}
Eliminating \(H_\parallel\) between the first and third equations of
(94.2) gives (94.4), certified as a series identity and valid exactly
by the algebra.  From \(\dot H_\perp\leq\Lambda/2\),
\(H_\perp(t_1)\leq\Lambda t_1/2<0\) for any \(t_1>0\); for
\(t\geq t_1\), \(d(1/H_\perp)/dt=-\dot H_\perp/H_\perp^2\geq3/2\), so
\(1/H_\perp(t)\geq2/(\Lambda t_1)+\tfrac32(t-t_1)\), whose right side
vanishes at \(t_2=t_1-4/(3\Lambda t_1)\); hence \(H_\perp\to-\infty\)
no later than \(t_2\), and \(\ln a_\perp=\int H_\perp\leq\tfrac23\ln(t_2-t)+\hbox{const}\).
Minimizing \(t_2=t_1-4/(3\Lambda t_1)\) over \(t_1\) gives
\(t_1=2/\sqrt{-3\Lambda}\) and \(t_*\leq t_2=2t_1=4/\sqrt{-3\Lambda}\);
with \(\Lambda=-\varepsilon^2/8\), \(t_1=\sqrt{32/3}/\varepsilon\) and
\(t_*\leq\sqrt{128/3}/\varepsilon\).
\end{proof}

The collapse is the Kasner-like anisotropic contraction anticipated
in V88: the transverse directions, in which the wave's null dust
exerts no pressure, contract under \(\Lambda<0\) and the wave's
energy, while the longitudinal direction, supported by \(p_z=\rho\),
first stays and then expands as \(\rho\) grows.  The negative
cosmological constant forced by the torus (V39) thus does more than
balance the wave at one instant: it drives the balanced datum to
collapse.

\subsection{Executable certificate}

The verifier

\[
 \texttt{scripts/verify\_sm\_v19\_v91\_two\_scale\_reduction.py}    \tag{94.5}
\]

computes the exact Ricci tensor symbolically (it uses the
\texttt{.grh\_sympy} library already present in the repository),
certifies the forms (94.1), their homogeneous reductions, the
Isaacson mean energy, and Theorem \ref{thm:v19v91-confirm}; it then
solves (94.2) as an exact rational series to order ten and certifies
the constraint, the conservation law, the invariant, the identity
(94.4), and the agreement with V89.  It emits

\[
 \texttt{artifacts/sm\_v19\_v91\_two\_scale\_reduction.json}.        \tag{94.6}
\]

The hashes are

\[
\begin{array}{c|l}
\hbox{object}&\hbox{SHA--256}\\ \hline
\hbox{verifier}&
\texttt{9BDC7CD77714C11B1D0C536051CC49CDE689636DFAD8B8141CEF91429D903C07}\\
\hbox{canonical payload}&
\texttt{ECF0FC966F2E7F3F008EDE16B69B7B69CA8E7A52BAC8415C20D8C6EA0A8932A7}.
\end{array}                                                     \tag{94.7}
\]

\begin{firewall}[Averaged system, adiabatic regime]
\label{fire:v19v91}
The reduced system (94.2) is the averaged equations with the standard
WKB invariant; its validity requires \(|H|\ll\omega\), which holds for
\(\varepsilon\ll1\) away from the collapse and fails as
\(t\to t_*\), where Theorem \ref{thm:v19v91-collapse} is a statement
about the reduced system only.  The series (94.3) is not proved to
converge to the collapse time.  The fourth-order secular terms of the
full solution have not been computed and compared with (94.3); that is
the V92 test.
\end{firewall}

\begin{decision}[V92 acquisition]
\label{dec:v19v91-next}
V92 will extend the \((z,t)\) engine to fourth order in the amplitude
and compare the \(\varepsilon^4t^4\) secular terms of the mean metric
with the \(T^4\) coefficients \(1/384\) of (94.3), after removing
the phase-averaged contribution of the wave's own square, thereby
testing the two-scale reduction at the first order where it makes a
nontrivial prediction.
\end{decision}

\begin{assessment}[V91 acquisition]
\label{ass:v19v91}
The exact plane-symmetric system is recorded, the V89 solution is
confirmed against the full Einstein equations by an independent
symbolic route, and the two-scale reduction is constructed and solved
to order ten with its constraint preserved.  The reduced system is
proved to collapse transversally within slow time \(\sqrt{128/3}\),
turning the negative cosmological constant of the torus from a
balancing device into the driver of a Kasner-like collapse of the
balanced standing wave.
\end{assessment}

\begin{record}[V91 append record]
\label{rec:v19v91}
V91 is appended to the validated V90 whose installed SHA--256 was
\texttt{758D764D3AED784FA836BE169A28D2BAE6C94BAFEC71F40B4950A8C4B46C0D3B}.
After installation only V91 is the active numeric SM note; frozen
archives and unrelated notes are unchanged.
\end{record}


\section{V92: the exact fourth-order evolution, the failure of the homogeneous two-scale model, and a focusing theorem}
\label{sec:v19v92}

This note carries the plane-symmetric evolution of the balanced
standing wave to fourth order in the amplitude, in the exponential
variables of V91 and with the time-symmetric conformal datum of
V66, as an exact trigonometric-polynomial solution in \((z,t)\).  The
fourth-order Hamiltonian constraint at the initial slice reproduces
the V66 value \(\Lambda_4=-145/2048\) from an independent engine, and
both constraints are preserved identically along the solution.  The
test announced in V91 is then decisive: the transverse secular
coefficient agrees exactly with the homogeneous two-scale model,
while the longitudinal one is four times smaller than predicted, and
the discrepancy is traced to the inhomogeneous secular growth of the
longitudinal metric at the wave's own wavelength, a feature of
standing waves that no homogeneous background can carry.  The
third-order wave has a nonlinear frequency shift and an amplitude
growth that also exceeds the adiabatic value.  Finally, the pointwise
rate law of V43 is turned into a focusing theorem for the full
equations: every geodesic normal to the balanced slice focuses within
proper time \((\pi/2)\sqrt{3/|\Lambda|}\), so the toroidal balanced
datum is a recollapsing universe whose lifetime in wave periods is of
order \(\varepsilon^{-1}\).

\subsection{The fourth-order solution}

\begin{theorem}[Exact fourth-order evolution]
\label{thm:v19v92-solution}
Let the datum be \(b(0)=\tfrac12\operatorname{artanh}(\varepsilon\cos z)\),
\(a(0)=2\ln\psi+\tfrac14\ln(1-\varepsilon^2\cos^2z)\), \(c(0)=2\ln\psi\),
\(\psi=1+\varepsilon^2w_2+\varepsilon^4w_4\), \(w_2=\tfrac7{128}\cos2z\),
\(\langle w_4\rangle=0\), with vanishing time derivatives, evolved by
the exact system of Proposition \ref{prop:v19v91-system}.  Then the
fourth-order Hamiltonian constraint at \(t=0\) fixes

\[
 \Lambda_4=-\frac{145}{2048},\qquad
 w_4=\frac{121}{4096}\cos2z+\frac{977}{131072}\cos4z,                  \tag{95.1}
\]

in agreement with V66, and the solution through order \(\varepsilon^4\)
is a finite sum of terms \(t^n\cos(kz)\cos(mt)\), \(t^n\cos(kz)\sin(mt)\)
with \(n\leq4\), \(k\leq4\), \(m\leq4\).  Its third-order wave is

\[
\begin{aligned}
 b_3&=\Bigl(\frac{61}{512}+\frac5{128}t^2\Bigr)\cos z\cos t-\frac3{128}t\cos z\sin t
     +\frac3{512}\cos z\cos3t\\
    &\quad+\Bigl(\frac5{256}-\frac1{128}t^2\Bigr)\cos3z\cos t-\frac1{128}t\cos3z\sin t
     +\frac{17}{768}\cos3z\cos3t,
\end{aligned}                                                          \tag{95.2}
\]

and the \(z\)-averaged secular parts of the fourth-order metric are

\[
 \langle a_4\rangle_{\rm sec}=-\frac{1555}{32768}-\frac{145}{4096}t^2-\frac{t^4}{384},\qquad
 \langle c_4\rangle_{\rm sec}=\frac{205}{16384}-\frac3{2048}t^2+\frac{t^4}{3072},   \tag{95.3}
\]

the remaining mean terms being of the form \(t^n\cos2t\), \(t\sin2t\),
\(\cos4t\) with \(n\leq2\).  The Hamiltonian and momentum constraints
vanish identically in \((z,t)\) at orders two and four along the
solution.
\end{theorem}

\begin{proof}
Orders one to four of the system are, successively, the free wave
equation for \(b_1\), the equations \(a_{2,tt}-a_{2,zz}=\Lambda_2\)
and \(c_{2,tt}=\Lambda_2+2a_{2,zz}+2b_{1,z}^2\) whose solutions are
the V91 forms, the third-order wave equation with sources bilinear in
\((b_1;a_2,c_2)\), and the fourth-order equations for \(a_4\), \(c_4\)
with sources bilinear in \((a_2,c_2)\) and in \((b_1,b_3)\); each is
solved mode by mode by the engine of V89.  The constraint at \(t=0\)
is a Poisson equation for \(w_4\) whose solvability condition is
(95.1).  All statements are certified coefficient by coefficient.
\end{proof}

\subsection{The two-scale test}

\begin{theorem}[Transverse confirmed, longitudinal refuted]
\label{thm:v19v92-test}
The homogeneous two-scale model of Theorem \ref{thm:v19v91-reduced}
predicts \(t^4\) coefficients \(-\tfrac1{384}\) for \(\langle a_4\rangle\)
and \(\tfrac1{768}\) for \(\langle c_4\rangle\), and \(\tfrac1{64}\) for
the \(t^3\) coefficient of \(\langle\tau_4\rangle\).  The exact
values are \(-\tfrac1{384}\), \(\tfrac1{3072}\), and \(\tfrac5{256}\).
The secular \(t^2\) part of the \(z\)-averaged longitudinal source,
twelve times the \(t^4\) coefficient, decomposes exactly as

\[
 \underbrace{-\langle c_{2,t}^2\rangle}_{-1/128}
 +\underbrace{4\langle b_{1,z}b_{3,z}\rangle}_{5/256}
 \underbrace{-2\langle c_2(2a_{2,zz}+2b_{1,z}^2)\rangle}_{-1/128}
 =\frac1{256},                                                         \tag{95.4}
\]

the other three terms of the source contributing no secular part.
The homogeneous model has only the middle term, with the adiabatic
value \(\tfrac1{64}\) in place of \(\tfrac5{256}\).
\end{theorem}

\begin{proof}
The three secular terms are certified.  The first and third arise
from the secular modulation \(-\tfrac{t^2}{16}\cos2z\) of \(c_2\),
which is the response of the longitudinal metric to the
\(z\)-modulated pressure \(2b_{1,z}^2\) of the standing wave: its
time average \(\tfrac18(1-\cos2z)\) is not homogeneous.  The middle
term is the wave's amplitude growth, \(\alpha/4\) with \(\alpha\)
from Theorem \ref{thm:v19v92-wave}.
\end{proof}

\begin{theorem}[Third-order wave: growth and frequency shift]
\label{thm:v19v92-wave}
The resonant terms of (95.2) are those of
\(\tfrac12(1+\varepsilon^2\alpha t^2)\cos z\cos((1+\varepsilon^2\delta)t)\)
with

\[
 \alpha=\frac5{64},\qquad\delta=\frac3{64},                             \tag{95.5}
\]

against the adiabatic value \(\alpha=\tfrac1{16}\) of the invariant
\(V\omega A^2\) on the homogeneous background; the excess \(\tfrac1{64}\)
and the slaved component \(-\tfrac1{128}t^2\cos3z\cos t\) are the
parametric response of the wave to the modulation \(\cos2z\) of the
longitudinal metric, and \(\delta\) is the nonlinear frequency
renormalization of the standing wave at second order in its
amplitude.
\end{theorem}

\begin{proof}
The coefficients of \(t^2\cos z\cos t\) and \(t\cos z\sin t\) in
(95.2) are \(\alpha/2\) and \(-\delta/2\); certified.  The
\(\cos3z\) terms are non-resonant and are the forced response to the
sources \(-2c_2b_{1,zz}\) and \(-b_{1,z}c_{2,z}\), whose \(\cos z\)
parts cancel and whose \(\cos3z\) parts add.
\end{proof}

\subsection{Consequence for the reduction}

\begin{proposition}[No homogeneous background for a standing wave]
\label{prop:v19v92-inhomogeneous}
The slow background of the balanced standing wave is plane symmetric
but not homogeneous: its longitudinal part carries the secular term
\(-\tfrac{T^2}{16}\cos2z\) at the same order as its homogeneous part
\(-\tfrac{T^2}{16}\) in the transverse metric.  Any reduction to a
Bianchi--I background therefore fails at order \(T^4\) in the
longitudinal sector, while the transverse sector, whose equation
(E2) contains no wave source and whose second-order secular term is
homogeneous, is correctly captured through \(T^4\).  The reduced
dynamics must keep the \(z\)-dependence of the slow fields exactly;
since the slow fields are trigonometric polynomials in \(z\) at every
order, the reduction is the secular part of the \(\varepsilon\)-series
itself.
\end{proposition}

\begin{proof}
Theorems \ref{thm:v19v92-test} and \ref{thm:v19v92-wave}.
\end{proof}

\subsection{Focusing for the full equations}

\begin{theorem}[Focusing of the balanced datum]
\label{thm:v19v92-focusing}
Let \((g,K=0)\) be any time-symmetric solution of the vacuum
constraints with \(\Lambda<0\) on a compact three-manifold, evolved
with unit lapse and zero shift.  Then along every geodesic normal to
the initial slice, for as long as the synchronous coordinates exist,

\[
 \tau(t)\geq\sqrt{3|\Lambda|}\,\tan\Bigl(\sqrt{|\Lambda|/3}\,t\Bigr),   \tag{95.6}
\]

so the normal congruence focuses, with \(\tau\to\infty\) and the
volume element of the congruence tending to zero, within proper time
\(t_f=\tfrac\pi2\sqrt{3/|\Lambda|}\).  For the balanced torus datum
\(\Lambda=-\varepsilon^2/8+O(\varepsilon^4)\) this is
\(t_f=\pi\sqrt6/\varepsilon+O(\varepsilon)\), consistent with the
reduced estimate \(\sqrt{128/3}/\varepsilon\) of Theorem
\ref{thm:v19v91-collapse}, which is smaller, and with the exact
fourth-order mean curvature, whose leading slope \(\varepsilon^2t/8\)
saturates (95.6) at \(t\to0\).
\end{theorem}

\begin{proof}
The rate law of V43, \(\partial_t\tau=|K|^2-\Lambda\) for unit lapse
and zero shift (certified along the second-order solution in V89),
together with \(|K|^2\geq\tau^2/3\), gives
\(\partial_t\tau\geq\tau^2/3+|\Lambda|\) with \(\tau(0)=0\); the
comparison solution of the equality is the right side of (95.6),
which blows up at \(t_f\).  The volume element satisfies
\(\partial_t\ln\sqrt g=-\tau\), so it tends to zero as
\(\tau\to\infty\).  The constants follow from \(|\Lambda|=\varepsilon^2/8\).
\end{proof}

The theorem is a Raychaudhuri statement and contains no perturbation
theory; it says that the toroidal branch of the programme, forced to
\(\Lambda<0\) by the rigidity of V39, is a recollapsing universe with
lifetime of order \(\varepsilon^{-1}\) wave periods.  A continuum on
the torus can only be a cosmological phase whose Hubble time
\(\sim|\Lambda|^{-1/2}\) is long compared with the wave periods, that
is, the hierarchy \(\varepsilon\ll1\) of the balanced datum.

\subsection{Executable certificate}

The verifier

\[
 \texttt{scripts/verify\_sm\_v19\_v92\_fourth\_order\_evolution.py}  \tag{95.7}
\]

integrates orders one to four with the engine of V89, checks the
second order against V91, the third- and fourth-order residuals, the
constraints at both orders, (95.1), the two-scale values and the
budget (95.4), (95.5), the slaved component, and the leading slope of
the mean curvature.  It emits

\[
 \texttt{artifacts/sm\_v19\_v92\_fourth\_order\_evolution.json}.     \tag{95.8}
\]

The hashes are

\[
\begin{array}{c|l}
\hbox{object}&\hbox{SHA--256}\\ \hline
\hbox{verifier}&
\texttt{CD33A0DDE10E9A21344F4FE9F70C56CB19E937762C12E14D330E057EB36C5DC1}\\
\hbox{canonical payload}&
\texttt{54D8AA450B8133796A92D8E79645EFBAB9424EC06056C29BD7587DDD396C0B7C}.
\end{array}                                                     \tag{95.9}
\]

\begin{firewall}[Fourth order, one standing wave, synchronous gauge]
\label{fire:v19v92}
The solution is a formal series through order four, valid for
\(t\ll\varepsilon^{-1}\); nothing is proved about its convergence.
Theorem \ref{thm:v19v92-focusing} holds only while the synchronous
coordinates exist, and the focusing it asserts may be a caustic of
the congruence rather than a curvature singularity.  The frequency
shift \(\delta\) is a property of the single standing wave and is not
the running coupling of V77, which concerns the slice constraint.
\end{firewall}

\begin{decision}[V93 acquisition]
\label{dec:v19v92-next}
The focusing theorem changes the target: the toroidal continuum is
cosmological, and its emergence must be studied as the limit of many
balanced modes with \(\Lambda_2=-\langle|\nabla h|^2\rangle/8\) and a
lifetime \(\sim|\Lambda|^{-1/2}\).  V93 will establish the
\(N\)-mode form of the focusing bound and of the second-order
secular expansion for a general finite-mode balanced datum, using
the three-dimensional engine of V81 for the sources, and will ask
whether isotropic superpositions produce an isotropic
Friedmann-like contraction rather than the anisotropic one of the
single wave.
\end{decision}

\begin{assessment}[V92 acquisition]
\label{ass:v19v92}
The fourth-order plane-symmetric evolution is exact, self-consistent,
and cross-checked against V66.  It confirms the two-scale reduction in
the transverse sector and refutes it in the longitudinal sector,
identifying the inhomogeneous secular modulation of a standing wave
as the obstruction, and it yields the nonlinear frequency shift
\(3/64\) and amplitude growth \(5/64\) of the wave.  The focusing
theorem establishes, without perturbation theory, that the toroidal
balanced datum recollapses within proper time \(\pi\sqrt6/\varepsilon\),
so that the torus branch of the programme is a cosmological phase, not
a stationary continuum.
\end{assessment}

\begin{record}[V92 append record]
\label{rec:v19v92}
V92 is appended to the validated V91 whose installed SHA--256 was
\texttt{3AF43F12D134A310C91E960E0B6FE6858E811AF2D83DD7F70243EE92F7E7BCA9}.
After installation only V92 is the active numeric SM note; frozen
archives and unrelated notes are unchanged.
\end{record}


\section{V93: finite-mode balanced data, additive null dust, and the Friedmann limit}
\label{sec:v19v93}

The focusing theorem of V92 refers to \(\Lambda\) alone, so its
content for a general finite-mode balanced datum on the torus is
fixed once \(\Lambda_2\) is known as a functional of the datum.  This
note certifies, with the three-dimensional engine of V81--V88, that
\(\Lambda_2=-\langle|\nabla h|^2\rangle/8\) for arbitrary sums of
transverse-traceless standing modes, that the space-time-averaged
secular acceleration of the slice metric is additive over modes and
is, mode by mode, the null dust of V89 along the mode's wavevector,
and that isotropic mode sets produce the isotropic Friedmann
contraction of a radiation-dominated universe with negative
cosmological constant.  The single standing wave of V88--V92 is thus
the anisotropic special case of a general statement: the balanced
torus datum evolves, on the slow scale, as a Bianchi--I universe
whose stress is the Isaacson radiation of its modes, and it focuses
within \(\pi\sqrt6/\sqrt{\langle|\nabla h|^2\rangle}\).

\subsection{The cosmological constant of a finite-mode datum}

\begin{theorem}[\(\Lambda_2\) and the focusing time]
\label{thm:v19v93-lambda}
For \(h=\sum_k\varepsilon_k\cos(k\cdot x)\,e_k\) with transverse
traceless polarizations \(e_k\), the balanced datum has

\[
 \Lambda_2=-\frac18\langle|\nabla h|^2\rangle=-\frac1{16}\sum_k|k|^2\varepsilon_k^2|e_k|^2,
                                                                    \tag{96.1}
\]

additively over modes, and by Theorem \ref{thm:v19v92-focusing} the
normal congruence of the datum focuses within proper time

\[
 t_f=\frac{\pi\sqrt6}{\sqrt{\langle|\nabla h|^2\rangle}}\,(1+O(\varepsilon^2)),
                                                                    \tag{96.2}
\]

which in units of the shortest wave period \(2\pi/|k|_{\max}\) is
\(\tfrac{\sqrt6}2|k|_{\max}/\sqrt{\langle|\nabla h|^2\rangle}\) periods.
\end{theorem}

\begin{proof}
\(\Lambda_2=\tfrac12\langle R_2[h]\rangle\) by the mean of the
second-order Hamiltonian constraint, and \(\langle R_2[h]\rangle=-\tfrac14\langle|\nabla h|^2\rangle\)
for transverse traceless \(h\); the identity and its additivity are
certified for eight single modes of four directions and three
frequencies, both polarizations, and for the six-mode isotropic set,
including the vanishing of the space mean of every cross term
\(R_2[h_k,h_{k'}]\), \(k'\neq\pm k\), which contains only the
harmonics \(e^{\pm ik\cdot x\pm ik'\cdot x}\).  (96.2) inserts (96.1)
into \(t_f=\tfrac\pi2\sqrt{3/|\Lambda|}\).
\end{proof}

\subsection{Additive null dust}

\begin{theorem}[Secular acceleration of a finite-mode datum]
\label{thm:v19v93-nulldust}
With unit lapse and zero shift, the space mean of the second-order
metric has the secular part

\[
 \langle g_2\rangle_{\rm sec}=t^2\sum_k\Lambda_2^{(k)}\bigl(\delta-\hat k\hat k\bigr),
 \qquad \Lambda_2^{(k)}=-\frac1{16}|k|^2\varepsilon_k^2|e_k|^2,        \tag{96.3}
\]

that is, each standing mode contributes the null dust of Theorem
\ref{thm:v19v89-isaacson} along its own wavevector, with
\(8\pi\rho^{(k)}=-\Lambda_2^{(k)}\), and the contributions add.
\end{theorem}

\begin{proof}
At second order, \(\ddot g_2=-2R_2[g_1]-2R_1[g_2]+4K_1K_1+2\Lambda_2\delta\)
with \(g_1=\sum_kh_k\cos|k|t\), \(K_1=-\tfrac12\dot g_1\), and
\(\operatorname{tr}K_1=0\).  The space mean of \(R_1[g_2]\) vanishes
because \(R_1\) is a sum of second derivatives.  The constant-in-time
part of \(\langle-2R_2[g_1]+4K_1K_1\rangle\) receives no contribution
from pairs of distinct frequencies, and pairs of equal frequency and
distinct wavevectors have zero space mean (Theorem
\ref{thm:v19v93-lambda}); the diagonal terms give
\(-\langle\operatorname{Ric}_2[h_k]\rangle+\tfrac12|k|^2\langle h_kh_k\rangle\)
per mode, which the engine certifies to equal
\(2\Lambda_2^{(k)}(\delta-\hat k\hat k)\) for every mode of the list,
including the oblique wavevector \((1,1,0)\) with both polarizations,
whose off-diagonal entries are those of \(\hat k\hat k\).  Twice
integrating the constant part gives (96.3).
\end{proof}

\begin{corollary}[Friedmann limit of isotropic data]
\label{cor:v19v93-friedmann}
If the mode set satisfies \(\sum_k\Lambda_2^{(k)}\hat k\hat k=\tfrac13\Lambda_2\delta\),
in particular for the three axes with both polarizations and equal
amplitudes, then \(\langle\ddot g_2\rangle=\tfrac43\Lambda_2\delta\),
that is, the mean scale factor obeys \(\ddot a/a=\tfrac23\Lambda\),
the Friedmann acceleration of a radiation fluid with
\(8\pi\rho=-\Lambda\) and \(p=\rho/3\) in the presence of \(\Lambda\).
In general the anisotropic part of the secular contraction is
\(-t^2\sum_k\Lambda_2^{(k)}(\hat k\hat k-\tfrac13\delta)\).
\end{corollary}

\begin{proof}
Sum (96.3) over an isotropic set; the six-mode case is certified with
\(\Lambda_2=-\tfrac34\) and \(\langle\ddot g_2\rangle=-\delta\).  The
Friedmann identification uses
\(\ddot a/a=\tfrac13\Lambda-\tfrac{4\pi}3(\rho+3p)\) with \(8\pi\rho=-\Lambda\), \(p=\rho/3\).
\end{proof}

\subsection{Interpretation}

\begin{proposition}[The toroidal continuum is a radiation universe at turnaround]
\label{prop:v19v93-interpretation}
Every finite-mode balanced datum on the torus is, on the slow scale
\(\varepsilon t\), a Bianchi--I universe at its moment of time
symmetry, whose matter is the Isaacson radiation of its own
gravitational waves with total density \(8\pi\rho=-\Lambda\), and
whose subsequent contraction is isotropic if and only if the
directional energy distribution \(\sum_k\Lambda_2^{(k)}\hat k\hat k\)
is isotropic.  The continuum limit of the programme, if it exists, is
therefore a cosmological phase: the slow expansion rate is
\(H\sim\sqrt{\langle|\nabla h|^2\rangle}\) and the hierarchy
\(H\ll|k|\) between the Hubble rate and the mode frequencies is the
condition under which the wave description and the averaging are
valid, as in the standard treatment of gravitational radiation in
cosmology.
\end{proposition}

\begin{proof}
Theorems \ref{thm:v19v93-lambda}--\ref{thm:v19v93-nulldust} and
Corollary \ref{cor:v19v93-friedmann}; the hierarchy statement is the
firewall of V91.
\end{proof}

\subsection{Executable certificate}

The verifier

\[
 \texttt{scripts/verify\_sm\_v19\_v93\_multimode\_secular.py}        \tag{96.4}
\]

computes, with the engines of V82 and V88, \(\Lambda_2\),
\(\langle|\nabla h|^2\rangle\), \(\langle\operatorname{Ric}_2\rangle\) and
\(\langle hh\rangle\) for the listed modes, checks (96.1) and the
null-dust form for each, the vanishing cross terms for five pairs, and
the additivity and Friedmann value for the isotropic set.  It emits

\[
 \texttt{artifacts/sm\_v19\_v93\_multimode\_secular.json}.           \tag{96.5}
\]

The hashes are

\[
\begin{array}{c|l}
\hbox{object}&\hbox{SHA--256}\\ \hline
\hbox{verifier}&
\texttt{9BBA18DFAD6F00F5D79DED4CFA66E2DDC000EA56AFAE475ED0303E49C15DBC71}\\
\hbox{canonical payload}&
\texttt{09FAF70E8F721E99099A96C2C702E85BA21D65CCD7651015A739A60456ADE6E8}.
\end{array}                                                     \tag{96.6}
\]

\begin{firewall}[Second order, space means, finitely many modes]
\label{fire:v19v93}
Only the space-time-averaged secular part of the second-order metric
is computed; the inhomogeneous secular terms of V92, which for
several modes include modulations at all differences of wavevectors,
are not, and the fourth-order corrections to (96.1) involve the pair
forms of V82--V87 whose closed forms in the frequency were read off
finitely many values.  The mode list is finite; the statement for
infinite mode sums requires the mode-number limit that the programme
has not constructed.
\end{firewall}

\begin{decision}[V94 acquisition]
\label{dec:v19v93-next}
V94 will examine the fourth-order correction to the cosmological
constant of a finite-mode datum as a bilinear form in the mode
energies, using the certified pair coefficients of V82 and V84, to
determine under what summability condition on \((\varepsilon_k)\) the
ratio \(\Lambda_4/\Lambda_2\) stays bounded as modes are added; this
is the classical form of the mode-number limit and decides whether
the radiation universe of Proposition
\ref{prop:v19v93-interpretation} survives the addition of
short-wavelength modes.
\end{decision}

\begin{assessment}[V93 acquisition]
\label{ass:v19v93}
The cosmological reading of the torus branch is now general: for any
finite-mode balanced datum, \(\Lambda_2\) is minus one eighth of the
mean gradient energy, the secular contraction is the additive null
dust of the modes, isotropic sets contract as a Friedmann radiation
universe with \(\ddot a/a=\tfrac23\Lambda\), and the focusing time is
\(\pi\sqrt6/\sqrt{\langle|\nabla h|^2\rangle}\).
\end{assessment}

\begin{record}[V93 append record]
\label{rec:v19v93}
V93 is appended to the validated V92 whose installed SHA--256 was
\texttt{69B1E27600DE7A806E170F75EDE0570B6A7B5FF462E8736DD87DFC426BE0EA87}.
After installation only V93 is the active numeric SM note; frozen
archives and unrelated notes are unchanged.
\end{record}


\section{V94: the fourth-order cosmological constant of two-mode data and the classical mode-number bound}
\label{sec:v19v94}

V93 identified the toroidal continuum as a radiation universe at
turnaround with \(\Lambda_2=-\langle|\nabla h|^2\rangle/8\).  The
mode-number limit asks whether the next coefficient
\(\Lambda_4\) stays controlled when modes are added.  This note
computes \(\Lambda_4\) exactly, with the three-dimensional
Lichnerowicz engine of V82, for single modes in seven directions and
for twenty-four two-mode pairs covering orthogonal, oblique, and
collinear wavevectors with all polarization combinations.  The
single-mode coefficient is direction independent,
\(\Lambda_4=-\tfrac{145}{2048}|k|^2\); half of the pair coefficients
take the universal value \(-\tfrac1{16}(|k_a|^2+|k_b|^2)\) exactly and
the other half exceed it by a negative interaction, the largest
relative to the gradient weight being \(-\tfrac{19}{144}\).  On the
analytic side, every term of \(\Lambda_4\) is shown to be a space mean
of a monomial of differential order two in four copies of \(h\) with
bounded inverse Laplacians, which gives the rigorous bound
\(|\Lambda_4|\leq C\|h\|_{H^2}^4\) on the torus: the classical
mode-number limit of the fourth-order balanced datum exists for
\(h\in H^2\).

\subsection{Single modes}

\begin{theorem}[Direction independence of the single-mode coefficients]
\label{thm:v19v94-single}
For \(h=\varepsilon\cos(k\cdot x)\,e\) with \(e\) transverse
traceless and \(|e|^2=2\),

\[
 \Lambda_2=-\frac{|k|^2}8\varepsilon^2,\qquad \Lambda_3=0,\qquad
 \Lambda_4=-\frac{145}{2048}|k|^2\varepsilon^4,                        \tag{97.1}
\]

for both polarizations of each of the wavevectors
\((0,0,1),(1,0,0),(0,0,2),(0,0,3),(1,1,0),(1,0,1),(1,1,1)\).
\end{theorem}

\begin{proof}
Certified by the Lichnerowicz series of the engine for the fourteen
modes, after rescaling the rational polarization matrices to
\(|e|^2=2\) (\(\Lambda_2\) scales with \(|e|^2\), \(\Lambda_4\) with
\(|e|^4\)).  The \(|k|^2\) scaling is the dimension of \(\Lambda\); the
independence of the direction and of the polarization is the
content.
\end{proof}

\subsection{Two-mode pairs}

\begin{theorem}[Quartic form of a pair]
\label{thm:v19v94-pairs}
For \(h=\alpha h_a+\beta h_b\) with unit-norm polarizations and
distinct wavevectors,
\(\Lambda_4=\alpha^4\Lambda_4^{(a)}+\beta^4\Lambda_4^{(b)}+\alpha^2\beta^2C_{ab}+\alpha^3\beta D_{ab}+\alpha\beta^3G_{ab}\),
with \(D_{ab}=G_{ab}=0\) except for the resonant collinear pair
\((0,0,1),(0,0,3)\) of equal polarization type, where
\(D_{ab}=-\tfrac{109}{512}\).  For the twelve pairs
\(z_+x_\times\), \(z_\times x_+\), \(z_+(101)_\times\), \(z_\times(101)_+\),
\(z_+(110)_+\), \(z_\times(110)_\times\), \(z_+(111)_+\), \(z_\times(111)_\times\),
\(z_+(002)_\times\), \(z_\times(002)_+\), \(z_+(003)_\times\), \(z_\times(003)_+\),

\[
 C_{ab}=-\frac1{16}\bigl(|k_a|^2+|k_b|^2\bigr)                          \tag{97.2}
\]

exactly, and for the remaining twelve pairs \(C_{ab}\) is (97.2) plus a
negative residual:

\[
\begin{array}{l|c|c|c}
\hbox{pair}&|k_a|^2+|k_b|^2&C_{ab}&\hbox{residual}\\ \hline
z_+x_+&2&-5/32&-1/32\\
z_\times x_\times&2&-17/128&-1/128\\
z_+(101)_+&3&-1503/5120&-543/5120\\
z_\times(101)_\times&3&-9/32&-3/32\\
z_+(110)_\times&3&-19/96&-1/96\\
z_\times(110)_+&3&-15/64&-3/64\\
z_+(111)_\times&4&-193/576&-49/576\\
z_\times(111)_+&4&-13/36&-1/9\\
z_\pm(002)_\pm&5&-95/144&-25/72\\
z_\pm(003)_\pm&10&-2605/2048&-1325/2048
\end{array}                                                            \tag{97.3}
\]

where \(z_\pm(00l)_\pm\) denotes the two equal-type collinear pairs.
Over all twenty-four pairs, \(-\tfrac{19}{144}\leq C_{ab}/(|k_a|^2+|k_b|^2)\leq-\tfrac1{16}\).
\end{theorem}

\begin{proof}
The five coefficients are separated by the six evaluations
\((\alpha,\beta)=(1,0),(0,1),(1,\pm1),(2,\pm1)\); all values are
certified, and \(z_+x_+\) reproduces the V82 value \(-\tfrac5{64}(l^2+1)\)
at \(l=1\).  The polarization labels refer to the rational bases of
the verifier: for the oblique wavevectors the plus type is
\(\hat u\hat u-\hat v\hat v\) and the cross type \(\hat u\hat v+\hat v\hat u\)
for the transverse pair \((u,v)\) listed there.
\end{proof}

The universal value (97.2) is a bilinear form in the mode energies,
\(-\tfrac1{16}\sum_{a\neq b}|k_a|^2\varepsilon_a^2\varepsilon_b^2=\tfrac12\Lambda_2\sum_b\varepsilon_b^2+\tfrac1{16}\sum_a|k_a|^2\varepsilon_a^4\),
that is, a relative correction \(\tfrac12\langle|h|^2\rangle\) to
\(\Lambda_2\); it is not the volume normalization, which would give
\(\tfrac16\langle|h|^2\rangle\).  The residuals are the genuine
mode--mode interactions; for the orthogonal plus--plus family the V82
closed form gives the residual \(-\tfrac1{64}(l^2+1)\), still growing
with the gradient weight.

\subsection{The Sobolev bound}

\begin{theorem}[Fourth-order coefficient of an \(H^2\) datum]
\label{thm:v19v94-sobolev}
There is a constant \(C\) such that for every transverse traceless
\(h\in H^2(T^3)\) the fourth-order Lichnerowicz coefficient of the
balanced datum \(\psi^4(\delta+\varepsilon h)\) with zero-mean
conformal corrections satisfies

\[
 |\Lambda_4|\leq C\|h\|_{H^2}^4.                                          \tag{97.4}
\]

Consequently \(|\Lambda_4/\Lambda_2|\leq8C\|h\|_{H^2}^4/\|\nabla h\|_{L^2}^2\),
and for a sequence of finite-mode data converging in \(H^2\) the
coefficients \(\Lambda_2,\Lambda_4\) converge.
\end{theorem}

\begin{proof}
The engine's fourth-order coefficient is
\(\Lambda_4=\tfrac12\langle R_4+R_2w_2-8D_1w_3-8D_2w_2\rangle\), with
\(w_2=\Delta^{-1}R_2\), \(w_3=\Delta^{-1}(R_3-8D_1w_2)\), where
\(R_n\) is the order-\(n\) part of the scalar curvature of
\(\delta+\varepsilon h\) and \(D_n\) the order-\(n\) part of its
Laplacian acting on a function.  Each \(R_n\) is a finite sum of
monomials with \(n\) factors of \(h\) and exactly two derivatives,
each \(D_n\) a sum of monomials with \(n\) factors of \(h\), two
derivatives, and one function argument; \(\Delta^{-1}\) on
zero-mean functions of the torus is bounded \(H^s\to H^{s+2}\).  On
\(T^3\), \(H^2\subset L^\infty\), \(H^1\subset L^6\), so
\(h\in L^\infty\), \(\partial h\in L^6\), \(\partial^2h\in L^2\), and
\(R_2\), \(R_3\), \(D_1w_2\) are in \(L^2\) by H\"older
(\(\partial h\,\partial h\in L^3\), \(h\,\partial^2h\in L^2\)), whence
\(w_2,w_3\in H^2\subset L^\infty\) with \(\partial w\in L^6\),
\(\partial^2w\in L^2\), all with norms bounded by powers of
\(\|h\|_{H^2}\).  Every monomial in the four means then has H\"older
exponents summing to at most one, and is bounded by
\(C\|h\|_{H^2}^4\).  The consequence uses \(|\Lambda_2|=\|\nabla h\|^2/8\).
\end{proof}

\begin{proposition}[The classical mode-number limit at fourth order]
\label{prop:v19v94-limit}
For a finite-mode datum, \(\Lambda_4\) is the quartic form
\(\sum_a\Lambda_4^{(a)}\varepsilon_a^4+\sum_{a<b}C_{ab}\varepsilon_a^2\varepsilon_b^2+(\hbox{resonant odd terms})\).
If \(|C_{ab}|\leq c(|k_a|^2+|k_b|^2)\) for all pairs, as holds with
\(c=\tfrac{19}{144}\) for the certified ones, then
\(|\Lambda_4|\leq\bigl(\tfrac{145}{2048}+c\bigr)\langle|\nabla h|^2\rangle\langle|h|^2\rangle+|\hbox{odd}|\),
so that \(|\Lambda_4/\Lambda_2|\leq8(\tfrac{145}{2048}+c)\langle|h|^2\rangle\)
up to the resonant terms: the relative fourth-order correction is
controlled by the mean square amplitude alone, an \(L^2\) condition,
sharper than the \(H^2\) condition of Theorem \ref{thm:v19v94-sobolev}.
\end{proposition}

\begin{proof}
Insert the bound into the quartic form and use
\(\sum_a|k_a|^2\varepsilon_a^2=\langle|\nabla h|^2\rangle\),
\(\sum_a\varepsilon_a^2=\langle|h|^2\rangle\) for unit-norm
polarizations; the pair bound is certified for the twenty-four pairs
and is a hypothesis beyond them.
\end{proof}

\subsection{Executable certificate}

The verifier

\[
 \texttt{scripts/verify\_sm\_v19\_v94\_fourth\_order\_pairs.py}      \tag{97.5}
\]

evaluates the Lichnerowicz series for the fourteen single modes and
the six evaluations of each of the twenty-four pairs, checks (97.1),
the orthogonal plus--plus value of V82, and records (97.2)--(97.3)
with the range of the ratio.  It emits

\[
 \texttt{artifacts/sm\_v19\_v94\_fourth\_order\_pairs.json}.         \tag{97.6}
\]

The hashes are

\[
\begin{array}{c|l}
\hbox{object}&\hbox{SHA--256}\\ \hline
\hbox{verifier}&
\texttt{96DC683BCA3289102C952F0C314E4E2B5DD4426A7E9C2D1D39898A2CE1885737}\\
\hbox{canonical payload}&
\texttt{8EAB06A49BCC6457586DB10B45F170DC58ECCD902F9AD82372504D854EB71E9F}.
\end{array}                                                     \tag{97.7}
\]

\begin{firewall}[Configuration sector, finitely many pairs, unspecified constant]
\label{fire:v19v94}
Only time-symmetric data are treated, so the momentum sector of
V83--V87 does not enter.  The universal value (97.2) and the pair
bound with \(c=\tfrac{19}{144}\) are certified for the listed pairs and
are not proved for general wavevectors; the constant \(C\) of (97.4)
is not computed; and the resonant odd terms, which occur for
collinear pairs with frequency ratio three, are not bounded in
general.  Nothing is said about convergence of the series in
\(\varepsilon\) or about the quantum mode sum, whose amplitudes
\(\varepsilon_k^2\sim1/(2c_g|k|)\) are not square summable.
\end{firewall}

\begin{decision}[V95 acquisition]
\label{dec:v19v94-next}
V95 is the compulsory companion audit.  V96 will take up the last
sentence of the firewall: the quantum mode sum with the vacuum
amplitudes of V59--V78 has \(\langle|h|^2\rangle\) logarithmically
divergent and \(\langle|\nabla h|^2\rangle\) quadratically divergent,
so the radiation universe of V93 has a divergent \(\Lambda_2\); V96
will formulate the renormalized mode-number limit as a subtraction
of the vacuum expectation in the operator-level constraint of V52,
and determine what finite quantity replaces \(\langle|\nabla h|^2\rangle\)
in the focusing time.
\end{decision}

\begin{assessment}[V94 acquisition]
\label{ass:v19v94}
The fourth-order cosmological constant of finite-mode balanced data
is now a certified quartic form: direction-independent single-mode
coefficients, a universal bilinear pair value equal to half the mean
square amplitude times \(\Lambda_2\), and negative interaction
residuals bounded by \(\tfrac{19}{144}\) of the gradient weight over
twenty-four pairs.  The rigorous \(H^2\) bound and the certified
\(L^2\)-type pair bound together show that the classical
mode-number limit of the balanced datum exists at fourth order for
square-summable amplitudes with finite gradient energy, and that the
obstruction to the continuum lies in the quantum amplitudes, not in
the classical interactions.
\end{assessment}

\begin{record}[V94 append record]
\label{rec:v19v94}
V94 is appended to the validated V93 whose installed SHA--256 was
\texttt{F907BADAAFA9EA70378F0DAA941B03690F76578815ABE12058E01323750CB518}.
After installation only V94 is the active numeric SM note; frozen
archives and unrelated notes are unchanged.
\end{record}


\section{V95: compulsory companion audit after the cosmological turn}
\label{sec:v19v95}

This is the required audit following V90.  Every active companion
note present in the shared folder was inspected read-only before the
renormalized mode-number limit announced in V94 begins.

\begin{record}[V95 companion snapshot]
\label{rec:v19v95-snapshot}
The files present were

\[
\begin{array}{c|r|r}
\hbox{file and SHA--256}&\hbox{bytes}&\hbox{lines}\\ \hline
\texttt{Renewal\_NS\_Stability\_Research\_Notes\_V26.tex}
&278283&5319\\
\multicolumn{3}{l}{
\texttt{82C887F8C2C69E4F28FAD7C3DC8DE5B9099CB5A4BF431EBA1FFF3E91935978AD}}
\\
\texttt{Renewal\_GRH\_Cofinal\_Visibility\_Attack\_V60.tex}
&580602&17282\\
\multicolumn{3}{l}{
\texttt{E090E736AE1C43D60756DB3E9C7929E0F2A179905B74630ADBE85266036737F4}}
\\
\texttt{Renewal\_Yang\_Mills\_Mass\_Gap\_Research\_Notes\_V34.tex}
&489911&12282\\
\multicolumn{3}{l}{
\texttt{3CF829F57455C43D3CE54387782C3C60EB3497F1B680A11DDCAED13306D8C469}}
\\
\texttt{Renewal\_YM\_Parallel\_Research\_Notes\_V5.tex}
&54174&1351\\
\multicolumn{3}{l}{
\texttt{306F1BB84CFA8A8D47EA569EC17E9545F9867C8D32B548EC9D9C444B4F3C0512}}.
\end{array}                                                     \tag{98.1}
\]

Since V90 the GRH note advanced from V56 to V60 (grouped
sum-of-squares certificates, closure of the two-pair band for all
paddings, and a degree-free exclusion schema with three-pair gates);
its V60 checkpoint records having read this note at V91.  The NS,
Yang--Mills mass-gap, and YM parallel notes are unchanged since V90.
\end{record}

\begin{decision}[V95 audit decision]
\label{dec:v19v95-audit}
No companion theorem, numerical constant, or executable artifact is
imported.  The GRH advance concerns polynomial inequalities on a
disk and their sum-of-squares certificates; it shares with the
present programme only the certificate discipline.  The V94 direction
stands: V96 formulates the renormalized mode-number limit.
\end{decision}

\begin{proof}
The only changed file is the GRH note, whose new sections
V57--V60 prove sum-of-squares identities and case closures for a
Pick-disk condition; none of its objects is a constrained
gravitational datum, an Einstein evolution, or a mode sum.
\end{proof}

\begin{assessment}[V95 acquisition]
\label{ass:v19v95}
The audit is current.  Since V90 the programme has solved the
second- and fourth-order plane-symmetric dynamics exactly, confirmed
and then refuted the homogeneous two-scale model sector by sector,
proved the focusing theorem that makes the toroidal branch a
recollapsing radiation universe, generalized its cosmological
constant, secular contraction, and focusing time to finite-mode
data, and bounded its fourth-order coefficient by Sobolev norms and
by certified pair coefficients.  The remaining obstruction is
quantum: the vacuum amplitudes are not square summable.
\end{assessment}

\begin{record}[V95 append record]
\label{rec:v19v95}
V95 is appended to the validated V94 whose installed SHA--256 was
\texttt{076B3818F4AB32DAB70BEF66321BBD2C2F7518016ADE1C59AD4A5BFE1B295429}.
After installation only V95 is the active numeric SM note; the
companions, frozen archives, and all unrelated notes are unchanged.  The
next compulsory companion audit is V100, which is also the
ten-version sweep and the freeze as \texttt{V100\_f19}.
\end{record}


\section{V96: the \texorpdfstring{$\Lambda=0$}{Lambda = 0} sectors are expanding radiation universes}
\label{sec:v19v96}

The physical space of the operator-level averaged constraint (V52,
V62) was constructed at \(\Lambda=0\): one two-dimensional space per
spectral value \(e>0\) of the normal-ordered graviton energy, with
mean curvature \(\pm\tau_e\), \(\tfrac23\tau_e^2=e\).  Those sectors
are not time symmetric, so the focusing theorem of V92, which needs
\(\Lambda<0\) and \(K=0\), does not apply to them.  This note
evolves the classical datum of such a sector, a constant-mean-curvature
plane-symmetric datum with \(\Lambda=0\) and \(\tau=-3H\) on the
expanding branch, exactly to second order, with coefficients in the
field \(\mathbb Q(H)\), \(H^2=\tfrac1{24}\).  The datum is an
expanding radiation universe: the balance \(\tfrac23\tau^2=E\) is
the Friedmann equation \(3H^2=8\pi\rho\), the second-order
corrections coincide with those of the balanced \(\Lambda<0\) datum,
and the wave redshifts with amplitude \(e^{-a}\) and phase
\(\int e^{-c}dt\).  On the contracting branch the Raychaudhuri bound
focuses the congruence within \(3/\tau_e\).  The finite quantity that
replaces the divergent \(\langle|\nabla h|^2\rangle\) of the vacuum
mode sum is the normal-ordered energy \(e\) itself:
\(H^2=e/6=\tfrac1{12c_g}\sum_k\omega_kn_k\), finite on every state
of finite graviton number.

\subsection{The classical datum of a sector}

\begin{proposition}[Sector datum]
\label{prop:v19v96-datum}
Let the constraint be imposed as in Theorem \ref{thm:v19v52-induced}
at \(\Lambda=0\).  The classical counterpart of the sector \((e,\mp)\)
is the constant-mean-curvature datum \(K=A+\tfrac\tau3g\),
\(\tau=\mp\tau_e\), \(\tfrac23\tau_e^2=e\), whose averaged
Hamiltonian constraint (65.1) reads \(\tfrac23\tau^2=\langle|A|^2-R\rangle\).
For the wave datum at its time-symmetric phase, \(A=0\) and
\(\langle-R_2\rangle=\tfrac14\langle|\nabla h|^2\rangle\), so with
\(H=-\tau/3\),

\[
 3H^2=\frac18\langle|\nabla h|^2\rangle=8\pi\rho_{\rm Isaacson},
 \qquad H^2=\frac e6,                                                  \tag{99.1}
\]

the Friedmann equation of a flat radiation universe.  In the number
basis \(e=\tfrac1{2c_g}\sum_k\omega_kn_k\) after normal ordering, so
\(H^2=\tfrac1{12c_g}\sum_k\omega_kn_k\) is finite on every state of
finite graviton number, while the unsubtracted sum
\(\tfrac1{4c_g}\sum_k\omega_k\) of (55.2) diverges.
\end{proposition}

\begin{proof}
(65.1) with \(\Lambda=0\); \(\langle R_2\rangle=-\tfrac14\langle|\nabla h|^2\rangle\)
is Theorem \ref{thm:v19v93-lambda}; the Isaacson identification is
Theorem \ref{thm:v19v89-isaacson} with \(-\Lambda_2\) replaced by
\(3H^2\).  The last sentence is the definition of the normal-ordered
constraint of Proposition \ref{prop:v19v52-operator}.
\end{proof}

\subsection{Second-order evolution of the expanding datum}

\begin{theorem}[Expanding plane-symmetric datum to second order]
\label{thm:v19v96-evolution}
In the variables of Proposition \ref{prop:v19v91-system} with
\(\Lambda=0\), let \(a=sa_1+s^2a_2\), \(b=sb_1+s^2b_2\), \(c=sc_1+s^2c_2\)
with \(a_1=c_1=Ht\), \(b_1=\tfrac12\cos z\cos t\), \(H^2=\tfrac1{24}\),
and initial data \(a_2(0)=-\tfrac18-\tfrac1{64}\cos2z\),
\(c_2(0)=\tfrac7{64}\cos2z\), \(b_2(0)=0\), all second-order time
derivatives zero at \(t=0\).  Then the balance \(3H^2=\langle b_{1,z}^2\rangle(0)=\tfrac18\)
holds, the Hamiltonian and momentum constraints vanish identically at
orders one and two along the solution, and

\[
\begin{aligned}
 a_2&=-\frac18-\frac{t^2}{16}-\frac1{64}\cos2z\cos2t,\qquad
 c_2=\frac1{32}-\frac{\cos2t}{32}+\frac7{64}\cos2z-\frac{t^2}{16}\cos2z,\\
 b_2&=-\frac H2\,t\cos z\cos t+\frac H4\,t^2\cos z\sin t+\frac H4\cos z\sin t,
\end{aligned}                                                          \tag{99.2}
\]

so that \(a_2,c_2\) coincide with the second-order fields of the
balanced \(\Lambda<0\) datum (V91--V92), \(\langle\tau\rangle=-3Hs+s^2(\tfrac t4-\tfrac1{16}\sin2t)\),
and the wave obeys the redshift law: its resonant part is that of
\(\tfrac12(1-Ht)\cos z\cos(t-\tfrac12Ht^2)\), amplitude
\(e^{-a_1}=(V\omega)^{-1/2}\) and phase \(\int e^{-c_1}dt\).
\end{theorem}

\begin{proof}
At first order the system is linear and homogeneous, satisfied by
the isotropic expansion and the free wave.  At second order the
sources are \(-2a_{1,t}^2-a_{1,t}c_{1,t}=-3H^2\) for \(a_2\),
\(-c_{1,t}^2-2a_{1,t}c_{1,t}+2a_{2,zz}+2b_{1,z}^2\) for \(c_2\), and
\(-3Hb_{1,t}-2c_1b_{1,zz}\) for \(b_2\); since \(-3H^2=\Lambda_2\) of
the balanced datum, the \(a_2\), \(c_2\) equations and data are those
of V92 and the solutions agree.  The \(b_2\) source is resonant and
produces the secular terms of (99.2), which are certified to equal
the expansion of \(\tfrac12(1-Ht)\cos z\cos(t-\tfrac12Ht^2)\) at order
\(H\).  All identities are certified in exact \(\mathbb Q(H)\)
arithmetic.
\end{proof}

The secular parts, \(\langle a_2\rangle=-\tfrac32H^2t^2\) and no
\(t^2\) term in \(\langle c_2\rangle\), are those of the Bianchi--I
radiation reduction (94.2) with \(\Lambda=0\), \(H_\perp=H_\parallel=H\),
\(8\pi\rho=3H^2\) at \(t=0\): the transverse directions decelerate
under the null dust, the longitudinal one does not.  For an
isotropic mode set the reduction is the flat radiation Friedmann
solution \(a(t)^2=1+2Ht\), by Corollary \ref{cor:v19v93-friedmann}.

\subsection{The two branches}

\begin{theorem}[Branches of a sector]
\label{thm:v19v96-branches}
For the full Einstein equations with \(\Lambda=0\), unit lapse and
zero shift, along every normal geodesic \(\partial_t\tau=|K|^2\geq\tau^2/3\).
Hence on the contracting branch \(\tau(0)=\tau_e>0\),
\(\tau(t)\geq\tau_e/(1-\tau_et/3)\) and the congruence focuses within
proper time \(3/\tau_e=\sqrt{3/(2e)}\); on the expanding branch
\(\tau(0)=-\tau_e\), \(\tau\) is non-decreasing and
\(\tau(t)\geq-\tau_e/(1+\tau_et/3)\), so the expansion rate
\(-\tau/3\) is at most the Friedmann rate \(H/(1+3Ht)\) of an
isotropic universe with the same initial rate, and no focusing bound
follows.
\end{theorem}

\begin{proof}
The V43 rate law with \(\Lambda=0\) and \(|K|^2\geq\tau^2/3\); the two
comparison solutions of \(\dot\tau=\tau^2/3\).
\end{proof}

\subsection{Executable certificate}

The verifier

\[
 \texttt{scripts/verify\_sm\_v19\_v96\_expanding\_sector.py}          \tag{99.3}
\]

implements the \((z,t)\) engine of V89 over \(\mathbb Q(H)\),
\(H^2=\tfrac1{24}\), integrates orders one and two, and checks the
balance, the residuals, both constraints, (99.2), the secular
coefficients, and the redshift law.  It emits

\[
 \texttt{artifacts/sm\_v19\_v96\_expanding\_sector.json}.             \tag{99.4}
\]

The hashes are

\[
\begin{array}{c|l}
\hbox{object}&\hbox{SHA--256}\\ \hline
\hbox{verifier}&
\texttt{8ECAD7F1DA8A339B18B94B555214CDAB9351064D9ACF65A067163F4084623517}\\
\hbox{canonical payload}&
\texttt{C177E529C48F6B51EE6DB2EFC56C6DF8043BFB9E27D2C5C04FF8AC1CF6B03E0E}.
\end{array}                                                     \tag{99.5}
\]

\begin{firewall}[Second order, classical datum, normal ordering by fiat]
\label{fire:v19v96}
The second-order Hubble correction, fixed by the third-order
Hamiltonian constraint, is not computed, and nothing is proved about
the expanding branch beyond second order: eternal expansion of the
full plane-symmetric solution is not established here.  The
finiteness of \(H^2=e/6\) rests on the normal ordering of Proposition
\ref{prop:v19v52-operator}, a choice of the model, not a derivation;
the zero-point energy is declared not to gravitate.  The quantum
sector and its classical datum are related by the coherent-state
correspondence of V59--V62, not by a derivation of the classical
dynamics from the constraint operator.
\end{firewall}

\begin{decision}[V97 acquisition]
\label{dec:v19v96-next}
With the gravitational sector of the emergent cosmology understood
at this order, V97 opens the matter sector, the last structural gap:
the plane-symmetric Einstein--scalar system for a massless field
\(\phi(z,t)\) on the torus, derived symbolically as in V91, its
second-order dynamics with the exact engine, and the additive
radiation law \(3H^2=8\pi(\rho_{\rm grav}+\rho_\phi)\) with the
scalar's Isaacson stress.  V98 and V99 will treat the massive scalar,
whose pressure differs from \(\rho/3\), and the gauge field, before
the V100 audit, sweep, and freeze.
\end{decision}

\begin{assessment}[V96 acquisition]
\label{ass:v19v96}
The \(\Lambda=0\) physical sectors of the operator-level constraint
are identified with expanding or contracting flat radiation
universes whose Hubble rate is set by the normal-ordered graviton
energy, \(H^2=e/6\), finite on every state of finite graviton number.
The expanding datum is evolved exactly to second order, with both
constraints preserved, the same second-order corrections as the
balanced torus, and the redshift of the wave; the contracting branch
focuses within \(3/\tau_e\).
\end{assessment}

\begin{record}[V96 append record]
\label{rec:v19v96}
V96 is appended to the validated V95 whose installed SHA--256 was
\texttt{2A04A2E261DAFDA4476F70C1F28C12CE0F9CBB4641A6E62B4875A3910ACBD78B}.
After installation only V96 is the active numeric SM note; frozen
archives and unrelated notes are unchanged.
\end{record}


\section{V97: the matter sector opens: a massless scalar is a graviton polarization in the plane-symmetric cosmology}
\label{sec:v19v97}

The last structural gap of the programme is the coupling of matter
to the emergent geometry.  This note takes the first exact step: the
plane-symmetric Einstein equations with a massless scalar field
\(\phi(z,t)\) on the torus, derived symbolically in the variables of
V91, and the evolution of a mixed graviton--scalar datum to fourth
order with the engine of V92.  The result is a symmetry: the system
is invariant under rotations of the pair \((b,\phi/\sqrt2)\), so the
scalar enters the constraints, the metric equations, and its own wave
equation exactly as a second, abelian polarization of the
gravitational wave.  Every result of V89--V96 on the balance, the
Isaacson stress, the secular contraction, the focusing time, the
Hubble rate, and the redshift transfers verbatim to scalar radiation,
with \(\langle|\nabla h|^2\rangle\) replaced by
\(\langle|\nabla h|^2\rangle+4\langle|\nabla\phi|^2\rangle\) in the
normalization where \(8\pi\) is absorbed into \(\phi\).  At the
massless level and in plane symmetry, the emergent cosmology cannot
tell matter radiation from gravitational radiation.

\subsection{The Einstein--scalar system}

\begin{proposition}[Plane-symmetric Einstein--scalar equations]
\label{prop:v19v97-system}
For the metric of Proposition \ref{prop:v19v91-system} and a scalar
\(\phi(z,t)\) with \(\operatorname{Ric}=\Lambda g+d\phi\otimes d\phi\)
(the massless field with \(8\pi\) absorbed into \(\phi\)), the
equations are

\[
\begin{aligned}
 &a_t^2+2a_tc_t-b_t^2-\tfrac12\phi_t^2-e^{-2c}\bigl(3a_z^2-2a_zc_z+2a_{zz}+b_z^2+\tfrac12\phi_z^2\bigr)=\Lambda,\\
 &a_{tz}+a_ta_z-a_zc_t+b_tb_z+\tfrac12\phi_t\phi_z=0,\\
 &f_{tt}+f_t(2a_t+c_t)-e^{-2c}\bigl(f_{zz}+f_z(2a_z-c_z)\bigr)=0\quad\hbox{for }f=b\hbox{ and }f=\phi,\\
 &a_{tt}+2a_t^2+a_tc_t-e^{-2c}\bigl(a_{zz}+2a_z^2-a_zc_z\bigr)=\Lambda,\\
 &c_{tt}+c_t^2+2a_tc_t-e^{-2c}\bigl(2a_z^2-2a_zc_z+2a_{zz}+2b_z^2+\phi_z^2\bigr)=\Lambda.
\end{aligned}                                                          \tag{100.1}
\]
\end{proposition}

\begin{proof}
Symbolic computation of the Ricci tensor and of \(\Box\phi\) for the
metric, with the six displayed forms certified against
\(G_{tt}-T_{tt}\), \(R_{tz}-\phi_t\phi_z\), \(\tfrac12(E_{xx}-E_{yy})\),
\(\Box\phi\), \(\tfrac12(E_{xx}+E_{yy})\), and \(E_{zz}\), where
\(E=\operatorname{Ric}-d\phi\otimes d\phi\); the \(\Lambda\) terms
are added as in V91.
\end{proof}

\begin{theorem}[The scalar is an abelian polarization]
\label{thm:v19v97-rotation}
The system (100.1) is invariant under the rotations

\[
 b\mapsto\cos\theta\,b-\sin\theta\,\phi/\sqrt2,\qquad
 \phi\mapsto\sqrt2\sin\theta\,b+\cos\theta\,\phi,                     \tag{100.2}
\]

which fix \(a\), \(c\), and \(\Lambda\): the constraints and the two
metric equations are invariant, and the pair of wave equations
rotates into itself.  Consequently, for every solution \((a,b,c)\) of
the vacuum system of V91 and every \(\theta\), the fields
\((a,\cos\theta\,b,c,\sqrt2\sin\theta\,b)\) solve the Einstein--scalar
system with the same \(\Lambda\), and conversely every plane-symmetric
Einstein--scalar solution whose \(b\) and \(\phi\) are proportional is
the rotation of a vacuum solution.
\end{theorem}

\begin{proof}
Substitution of (100.2) with symbolic \(\theta\) into the certified
forms (100.1): the differences vanish identically for the four
invariant equations, and the two wave operators transform by the
rotation matrix of (100.2), all certified.  The invariance is
manifest from the fact that \(b\) and \(\phi/\sqrt2\) enter the four
invariant equations only through the quadratic forms
\(b_t^2+\tfrac12\phi_t^2\), \(b_z^2+\tfrac12\phi_z^2\), \(b_tb_z+\tfrac12\phi_t\phi_z\),
and that the wave operator is linear with coefficients depending on
\(a,c\) only.
\end{proof}

\subsection{The mixed datum to fourth order}

\begin{theorem}[Mixed graviton--scalar datum]
\label{thm:v19v97-mixed}
Let \(\beta^2+\tfrac12\gamma^2=1\), and take the datum
\(b(0)=\beta B_0\), \(\phi(0)=\gamma B_0\), \(B_0=\tfrac12\operatorname{artanh}(\varepsilon\cos z)\),
with \(a(0),c(0)\) the conformal datum of Theorem
\ref{thm:v19v92-solution}, time symmetric.  For \((\beta,\gamma)=(\tfrac13,\tfrac43)\)
the fourth-order evolution, computed from (100.1) without use of the
symmetry, has \(\Lambda_2=-\tfrac18\), \(\Lambda_4=-\tfrac{145}{2048}\),
the same \(a_2,c_2,a_4,c_4\) as the vacuum datum of V92, and
\(b=\beta B\), \(\phi=\gamma B\) with \(B\) the vacuum wave (95.2); both
constraints vanish identically at orders two and four.  In
particular the scalar's third-order wave is
\(\phi_3=\tfrac43b_3\), with the same amplitude growth \(\tfrac5{64}\)
and frequency shift \(\tfrac3{64}\).
\end{theorem}

\begin{proof}
The engine of V92 extended by the scalar sources of (100.1) at
each order; the comparison with the \((\beta,\gamma)=(1,0)\) run of the
same code is certified field by field.  The result is the instance
\(\cos\theta=\tfrac13\), \(\sin\theta=\tfrac{2\sqrt2}3\) of Theorem
\ref{thm:v19v97-rotation}.
\end{proof}

\subsection{Consequences for the emergent cosmology}

\begin{corollary}[Scalar radiation in the balance and in the dynamics]
\label{cor:v19v97-transfer}
For plane-symmetric data with graviton and scalar waves,
\(\Lambda_2=-\tfrac18\bigl(\langle|\nabla h|^2\rangle+4\langle|\nabla\phi|^2\rangle\bigr)\)
for time-symmetric data on the torus, and
\(3H^2=\tfrac18\bigl(\langle|\nabla h|^2\rangle+4\langle|\nabla\phi|^2\rangle\bigr)\)
for the \(\Lambda=0\) expanding datum, where
\(\langle|\nabla h|^2\rangle=8\langle b_z^2\rangle\) in the normalization
\(g_{xx}=e^{2a+2b}\); the scalar's Isaacson stress
is null dust along its wavevector with \(8\pi\rho_\phi=\tfrac12\langle\phi_t^2+\phi_z^2\rangle\),
additive with the graviton's; the secular contraction, the focusing
time \(\tfrac\pi2\sqrt{3/|\Lambda|}\), the redshift law, and the
fourth-order interaction coefficients of the plane class are those
of the vacuum wave of equal weight.
\end{corollary}

\begin{proof}
Theorem \ref{thm:v19v97-rotation} maps every mixed plane-symmetric
solution with proportional \(b,\phi\) to a vacuum solution of weight
\(\beta^2+\tfrac12\gamma^2\); the general non-proportional case at
second order follows from the additivity of Theorem
\ref{thm:v19v93-nulldust}, whose proof uses only the quadratic
structure of the sources, which (100.1) shares.
\end{proof}

The statement is sharper than an analogy.  In the plane-symmetric
sector the massless scalar is dynamically indistinguishable from the
cross polarization would be if it coupled linearly; the two are
distinguished only outside plane symmetry, where the gravitational
cross polarization couples nonlinearly through the off-diagonal
metric while the scalar keeps its abelian coupling, and by mass,
which the scalar can carry and the graviton cannot.  The programme's
matter sector therefore begins with a theorem of indistinguishability
and a precise list of what distinguishes.

\subsection{Executable certificate}

The verifier

\[
 \texttt{scripts/verify\_sm\_v19\_v97\_scalar\_sector.py}             \tag{100.3}
\]

derives (100.1) symbolically, certifies the rotation symmetry
(100.2), runs the fourth-order evolution for the vacuum and the mixed
datum, and checks Theorem \ref{thm:v19v97-mixed}.  It emits

\[
 \texttt{artifacts/sm\_v19\_v97\_scalar\_sector.json}.                \tag{100.4}
\]

The hashes are

\[
\begin{array}{c|l}
\hbox{object}&\hbox{SHA--256}\\ \hline
\hbox{verifier}&
\texttt{6D38EF99DB0154548CD9F0DEDBE866970D73D5F3A6007A987DA53C8B8C8F8E3B}\\
\hbox{canonical payload}&
\texttt{FDD275CF710E9E03212100AEB214FB0B5D9B00715300E5D87A3D7AB41E4C5D98}.
\end{array}                                                     \tag{100.5}
\]

\begin{firewall}[Massless, plane symmetric, classical]
\label{fire:v19v97}
The symmetry is exact only for the massless minimally coupled scalar
in plane symmetry with a diagonal metric; a mass term, a potential, a
nonminimal coupling, a second spatial dependence, or the gravitational
cross polarization break it.  The scalar is classical; its
quantization in the seed of V44 and its place in the operator
constraint of V52 are not constructed.  Nothing here concerns gauge
fields or fermions, which are the content of the Standard Model.
\end{firewall}

\begin{decision}[V98 acquisition]
\label{dec:v19v97-next}
V98 breaks the symmetry by mass: the plane-symmetric Einstein--scalar
system with \(\Box\phi=m^2\phi\), its second-order dynamics for a
massive standing mode of frequency \(\sqrt{k^2+m^2}\), the pressure
\(p_z\neq\rho\) of its Isaacson stress, and the resulting departure of
the secular contraction from the null-dust form, which is the first
place where the emergent cosmology distinguishes matter from
gravitational radiation.
\end{decision}

\begin{assessment}[V97 acquisition]
\label{ass:v19v97}
The matter sector is opened with an exact theorem: in the
plane-symmetric emergent cosmology a massless scalar is an abelian
polarization of the gravitational wave, entering the constraints,
the metric equations, and the wave equation through the rotation
invariant \(b^2+\tfrac12\phi^2\) and its derivatives.  All balance,
dynamics, focusing, and redshift results transfer to scalar
radiation, and the fourth-order mixed datum is certified to be the
rotation of the vacuum one.
\end{assessment}

\begin{record}[V97 append record]
\label{rec:v19v97}
V97 is appended to the validated V96 whose installed SHA--256 was
\texttt{65470AEB05A476D51DC61E9756DD49D05C405FFF84113E94D4EA7025BF99FBF3}.
After installation only V97 is the active numeric SM note; frozen
archives and unrelated notes are unchanged.
\end{record}


\section{V98: mass breaks the symmetry: the massive scalar and the first distinction between matter and gravitational radiation}
\label{sec:v19v98}

V97 showed that a massless scalar is a graviton polarization in the
plane-symmetric emergent cosmology.  This note adds a mass.  The
potential enters the Einstein equations as a field-dependent
cosmological term, \(\operatorname{Ric}=\Lambda g+d\phi\otimes d\phi+Vg\),
\(V=\tfrac12m^2\phi^2\), and the second-order dynamics of a
time-symmetric massive standing mode is computed exactly for three
masses.  The mode's Isaacson stress has energy density
\(\rho=\omega^2/4\), longitudinal pressure \(p_z=k^2/4\), and no
transverse pressure, so \(p_z/\rho=k^2/\omega^2<1\): it is a mixture
of null dust and dust along its wavevector.  The consequence is a
longitudinal secular contraction \(-\tfrac{m^2}{16}t^2\) that vanishes
for massless fields and gravitons: mass is the first property that
the emergent geometry can read off its matter.

\subsection{The massive system}

\begin{proposition}[Plane-symmetric Einstein--massive-scalar equations]
\label{prop:v19v98-system}
With \(V=\tfrac12m^2\phi^2\) and \(8\pi\) absorbed into \(\phi\), the
equations of Proposition \ref{prop:v19v97-system} are modified as
follows: the Hamiltonian constraint acquires \(-V\) on its left side,
the two metric equations acquire \(+V\) on their right sides
(\(\Lambda\to\Lambda+V\)), the momentum constraint and the \(b\) wave
equation are unchanged, and the scalar obeys
\(\phi_{tt}+\phi_t(2a_t+c_t)-e^{-2c}(\phi_{zz}+\phi_z(2a_z-c_z))+m^2\phi=0\).
The Raychaudhuri source is \(R_{tt}=-\Lambda+\phi_t^2-V\).
\end{proposition}

\begin{proof}
Symbolic computation as in V97 with the stress tensor
\(T=d\phi\otimes d\phi-g(\tfrac12|d\phi|^2+V)\), whose trace-reversed
form is \(d\phi\otimes d\phi+Vg\); all six forms and \(R_{tt}\) are
certified.
\end{proof}

\subsection{The massive standing mode}

\begin{theorem}[Second-order dynamics of a massive mode]
\label{thm:v19v98-mode}
Let \(\phi_1=\cos z\cos\omega t\), \(\omega^2=1+m^2\), \(b=0\), with
the time-symmetric conformal datum \(a_2(0)=c_2(0)=2w_2\),
\(\langle w_2\rangle=0\), on the torus.  For \(m^2=0,3,8\)
(\(\omega=1,2,3\)):

\[
 \Lambda_2=-\frac{\omega^2}4,\qquad
 \langle a_2\rangle_{\rm sec}=-\frac{\omega^2+1}{16}t^2+\hbox{const},\qquad
 \langle c_2\rangle_{\rm sec}=-\frac{\omega^2-1}{16}t^2+\hbox{const}=-\frac{m^2}{16}t^2+\hbox{const},
                                                                    \tag{101.1}
\]

both constraints vanish identically along the solution, and the
mean curvature obeys the pointwise rate law
\(\dot\tau_2=-\Lambda_2+\phi_{1,t}^2-V_1\).
\end{theorem}

\begin{proof}
The second-order equations are \(a_{2,tt}-a_{2,zz}=\Lambda_2+V_1\),
\(c_{2,tt}=\Lambda_2+V_1+2a_{2,zz}+\phi_{1,z}^2\), with \(\Lambda_2\)
and \(w_2\) fixed by the constraint at \(t=0\); all statements are
certified by the engine for the three masses.
\end{proof}

\begin{corollary}[Isaacson stress of a massive mode]
\label{cor:v19v98-stress}
The time-averaged stress of the mode is
\(\rho=\tfrac{\omega^2}4\), \(p_z=\tfrac14\), \(p_x=p_y=0\), so that
\(\Lambda_2=-\rho\) is the balance and the secular accelerations of
(101.1) are the Bianchi--I values
\(\langle\ddot g_2\rangle_{xx}=2\Lambda_2+(\rho-p_z)\),
\(\langle\ddot g_2\rangle_{zz}=2\Lambda_2+(\rho+p_z)\).  In particular
\(p_z/\rho=k^2/\omega^2\) and the transverse pressure vanishes: a
massive standing mode is null dust plus dust along its wavevector, in
the proportions \(k^2:m^2\), and only the dust part contracts the
longitudinal direction.
\end{corollary}

\begin{proof}
The averages of \(\tfrac12(\phi_t^2+\phi_z^2)+V\),
\(\tfrac12(\phi_t^2+\phi_z^2)-V\), and \(\tfrac12(\phi_t^2-\phi_z^2)-V\)
for \(\phi_1\) with \(\omega^2=1+m^2\); the Bianchi--I identification
is certified as an identity between (101.1) and these averages for the
three masses.
\end{proof}

\subsection{Interpretation}

\begin{proposition}[What the geometry reads off its matter]
\label{prop:v19v98-reading}
In the plane-symmetric emergent cosmology at second order, the
metric sector depends on the matter content only through the
time-averaged stress of each standing mode, which for gravitons and
massless scalars is null dust and for a massive scalar is null dust
plus dust in the ratio \(k^2:m^2\).  The first quantity of the matter
sector visible to the geometry is therefore the dispersion relation
\(\omega^2=k^2+m^2\) through the pressure ratio \(p_z/\rho=k^2/\omega^2\),
equivalently the longitudinal secular contraction \(-\tfrac{m^2}{16}t^2\)
per unit amplitude.  The strong energy condition holds on average,
\(\langle R_{tt}\rangle=-\Lambda+\tfrac18(2+m^2)>0\) for the mode of
unit amplitude, but not pointwise when \(V>\phi_t^2\), so the
focusing theorem of V92 applies to massive matter only after
averaging.
\end{proposition}

\begin{proof}
Theorems \ref{thm:v19v93-nulldust}, \ref{thm:v19v97-rotation},
\ref{thm:v19v98-mode}; the average of \(R_{tt}\) is
\(\langle\phi_t^2\rangle-\langle V\rangle=\tfrac{\omega^2}4-\tfrac{m^2}8=\tfrac{2+m^2}8\).
\end{proof}

\subsection{Executable certificate}

The verifier

\[
 \texttt{scripts/verify\_sm\_v19\_v98\_massive\_scalar.py}            \tag{101.2}
\]

derives the massive system symbolically, and for \(m^2=0,3,8\)
integrates the second order, checks the constraints, the rate law,
(101.1), and the Bianchi--I identification.  It emits

\[
 \texttt{artifacts/sm\_v19\_v98\_massive\_scalar.json}.               \tag{101.3}
\]

The hashes are

\[
\begin{array}{c|l}
\hbox{object}&\hbox{SHA--256}\\ \hline
\hbox{verifier}&
\texttt{DD675124035729FE20DFBBEC5674EB659DDB39D99EAFC889A02E8E6DE1F4B820}\\
\hbox{canonical payload}&
\texttt{CEC9A64395B58FEFAEC8A350442AEAA033B678BF89B736E93B59B65A14DF72C3}.
\end{array}                                                     \tag{101.4}
\]

\begin{firewall}[Three masses, one mode, second order]
\label{fire:v19v98}
The laws (101.1) are certified for three integer frequencies and
read as general in \(\omega\) by the averaging argument of the
corollary, which is elementary but is not itself a certificate for
non-integer \(\omega\).  The mode is classical and time symmetric;
the coupling of the massive scalar to the graviton (\(b\neq0\)) and
its fourth-order self-interaction are not computed.
\end{firewall}

\begin{decision}[V99 acquisition]
\label{dec:v19v98-next}
V99 treats the gauge field: the plane-symmetric Einstein--Maxwell
system for a potential \(A_x(z,t)\), whose stress is null dust on
average but whose anisotropic part \(T_{xx}-T_{yy}\) sources the
gravitational polarization \(b\) directly, so that a standing
electromagnetic wave generates a gravitational wave at twice its
wavenumber; the second-order dynamics will give the exact secular
conversion rate.  V100 is then the audit, the ten-version sweep
V90--V99, and the freeze as \texttt{V100\_f19}.
\end{decision}

\begin{assessment}[V98 acquisition]
\label{ass:v19v98}
The massive scalar breaks the polarization symmetry of V97 in a
computable way: its standing mode is null dust plus dust in the
ratio \(k^2:m^2\), with \(\Lambda_2=-\omega^2/4\), transverse secular
contraction \(-(\omega^2+1)t^2/16\) and longitudinal
\(-m^2t^2/16\), certified for three masses with both constraints and
the rate law.  The dispersion relation of matter is the first
matter property visible to the emergent geometry.
\end{assessment}

\begin{record}[V98 append record]
\label{rec:v19v98}
V98 is appended to the validated V97 whose installed SHA--256 was
\texttt{21488E2E6FC4B9D0D0B9B73C34D1F01D6D6A4E2FE1245C4D9436B0F9FAF8734F}.
After installation only V98 is the active numeric SM note; frozen
archives and unrelated notes are unchanged.
\end{record}


\section{V99: the gauge field: spin is visible, and a standing electromagnetic wave does not radiate gravitationally at second order}
\label{sec:v19v99}

V97 and V98 placed the scalar in the plane-symmetric emergent
cosmology.  This note places the gauge field.  The plane-symmetric
Einstein--Maxwell system for a potential \(A_x(z,t)\) is derived
symbolically; its stress is traceless and pointwise null dust in the
\((t,z)\) block, but it carries an anisotropic transverse part
\(T_{xx}-T_{yy}\) that sources the gravitational polarization \(b\)
directly.  The second-order dynamics of a standing electromagnetic
wave is computed exactly: the isotropic metric sector \(a_2\) and
\(\Lambda_2\) coincide with those of the massless scalar of equal
energy, the longitudinal metric \(c_2\) is static where the scalar's
grows, and the induced gravitational polarization is
\(b_2=\tfrac14\sin^2z\sin^2t\), bounded, with no secular term: the
standing wave induces a static deformation and a homogeneous
second-harmonic breathing but no growing gravitational wave.  The
emergent geometry thus distinguishes spin one from spin zero at
second order through the polarization sector and through the
inhomogeneous secular terms, while their averaged Isaacson stresses
are identical.

\subsection{The Einstein--Maxwell system}

\begin{proposition}[Plane-symmetric Einstein--Maxwell equations]
\label{prop:v19v99-system}
For the metric of Proposition \ref{prop:v19v91-system} and the
potential \(A_\mu=(0,A(z,t),0,0)\), with \(8\pi\) absorbed and
\(\rho:=\tfrac12e^{-2a-2b}(A_t^2+e^{-2c}A_z^2)\),
\(\pi:=\tfrac12e^{-2a-2b}(A_t^2-e^{-2c}A_z^2)\), the stress is
traceless with \(T_{tt}=e^{-2c}T_{zz}=\rho\),
\(e^{-2a-2b}T_{xx}=-e^{-2a+2b}T_{yy}=-\pi\), and the equations
\(\operatorname{Ric}=\Lambda g+T\), \(d{*}F=0\) are

\[
\begin{aligned}
 &a_t^2+2a_tc_t-b_t^2-e^{-2c}\bigl(3a_z^2-2a_zc_z+2a_{zz}+b_z^2\bigr)-\rho=\Lambda,\\
 &a_{tz}+a_ta_z-a_zc_t+b_tb_z+\tfrac12e^{-2a-2b}A_tA_z=0,\\
 &b_{tt}+b_t(2a_t+c_t)-e^{-2c}\bigl(b_{zz}+b_z(2a_z-c_z)\bigr)+\pi=0,\\
 &a_{tt}+2a_t^2+a_tc_t-e^{-2c}\bigl(a_{zz}+2a_z^2-a_zc_z\bigr)=\Lambda,\\
 &c_{tt}+c_t^2+2a_tc_t-e^{-2c}\bigl(2a_z^2-2a_zc_z+2a_{zz}+2b_z^2\bigr)-\rho=\Lambda,\\
 &A_{tt}+A_t(c_t-2b_t)-e^{-2c}\bigl(A_{zz}-A_z(2b_z+c_z)\bigr)=0.
\end{aligned}                                                          \tag{102.1}
\]

The Maxwell equation does not involve \(a\), and the transverse
equation (E2) has no electromagnetic source.
\end{proposition}

\begin{proof}
Symbolic computation of the Ricci tensor, of \(T=F\cdot F-\tfrac14gF^2\),
and of \(\nabla_\mu F^{\mu x}\); the six forms and the stress
components are certified.  The absence of \(a\) from the Maxwell
equation is the conformal invariance of the Maxwell field in the
transverse conformal factor; the absence of a source in (E2) is the
tracelessness of \(T\) together with \(T_{xx}+T_{yy}=0\).
\end{proof}

\subsection{The standing electromagnetic wave}

\begin{theorem}[Second-order dynamics of a standing electromagnetic wave]
\label{thm:v19v99-wave}
Let \(A_1=\cos z\cos t\), \(b=0\) at first order, with the
time-symmetric conformal datum \(a_2(0)=c_2(0)=2w_2\),
\(\langle w_2\rangle=0\), \(b_2(0)=\dot b_2(0)=0\).  Then
\(\Lambda_2=-\tfrac14\), both constraints vanish identically along the
solution, the rate law \(\dot\tau_2=-\Lambda_2+\rho_1\) holds
pointwise, and

\[
 a_2=-\frac{t^2}8-\frac1{32}\cos2z\cos2t,\qquad
 c_2=-\frac1{32}\cos2z,\qquad
 b_2=\frac14\sin^2z\sin^2t=\frac1{16}(1-\cos2z)(1-\cos2t).            \tag{102.2}
\]
\end{theorem}

\begin{proof}
The second-order sources are \(\Lambda_2\) for \(a_2\),
\(\Lambda_2+2a_{2,zz}+\rho_1\) for \(c_2\), and \(-\pi_1\) for
\(b_2\), with \(\rho_1=\tfrac14(1-\cos2z\cos2t)\) and
\(\pi_1=\tfrac14(\cos2z-\cos2t)\).  The \(c_2\) source vanishes
identically, so \(c_2\) is its initial datum; the \(b_2\) source
contains no resonant term, since \(\pi_1\) has no \(\cos2z\cos2t\)
component, and the particular solution plus the homogeneous fix give
(102.2).  All certified.
\end{proof}

\begin{corollary}[Scalar versus vector radiation]
\label{cor:v19v99-comparison}
For the massless scalar \(\phi_1=\cos z\cos t\) (V98, \(m=0\)) and the
electromagnetic wave \(A_1=\cos z\cos t\), which have the same
averaged Isaacson stress (null dust with \(\rho=\tfrac14\)):
\(\Lambda_2\) and \(a_2\) coincide; the scalar's longitudinal metric
carries the inhomogeneous secular term \(-\tfrac{t^2}{16}\cos2z\)
(the modulation responsible for the failure of the homogeneous
reduction in V92) while the electromagnetic \(c_2\) is static; and
the scalar induces no gravitational polarization while the
electromagnetic wave induces the bounded \(b_2\) of (102.2).
\end{corollary}

\begin{proof}
Certified by running the V98 engine at \(m=0\) alongside the
present one.  The difference in \(c_2\) comes from the Ricci
sources: \(R_{zz}=\phi_z^2\) for the scalar, whose time average
\(\tfrac12\sin^2z\) is modulated, against \(R_{zz}=e^{2c}\rho\) for the
Maxwell field, whose value \(\tfrac14(1-\cos2z\cos2t)\) has an
unmodulated time average.  The difference in \(b\) is the anisotropic
stress \(\pi\), absent for the scalar.
\end{proof}

\subsection{Interpretation}

\begin{proposition}[What the geometry reads off the gauge field]
\label{prop:v19v99-reading}
At second order the emergent geometry distinguishes the spin of its
massless radiation by two signatures invisible in the averaged stress:
the polarization sector, excited by the anisotropic stress of the
vector field and not by the scalar, and the inhomogeneous secular
growth of the longitudinal metric, present for the scalar and absent
for the vector.  A standing electromagnetic wave does not convert into
a growing gravitational wave at second order; conversion requires a
background field or a travelling configuration, whose \(\pi\) has a
resonant component.
\end{proposition}

\begin{proof}
Theorem \ref{thm:v19v99-wave} and Corollary
\ref{cor:v19v99-comparison}; the last sentence is the observation that
\(\pi_1\) for the standing wave has no \(\cos2z\cos2t\) term, so no
source at the resonant wavenumber and frequency of the \(b\)
equation.
\end{proof}

\subsection{Executable certificate}

The verifier

\[
 \texttt{scripts/verify\_sm\_v19\_v99\_maxwell\_sector.py}            \tag{102.3}
\]

derives (102.1) symbolically, integrates the second order, checks the
constraints, the rate law, (102.2), and the comparison with the V98
scalar at \(m=0\).  It emits

\[
 \texttt{artifacts/sm\_v19\_v99\_maxwell\_sector.json}.               \tag{102.4}
\]

The hashes are

\[
\begin{array}{c|l}
\hbox{object}&\hbox{SHA--256}\\ \hline
\hbox{verifier}&
\texttt{B8AC2091F90B7C08CA2D75A104556858CF5879343EF135C51FE18E53941C8A6C}\\
\hbox{canonical payload}&
\texttt{515BEF3BBAD5B439CDC237FABA713D1AD69EBD7872379A4E7774417C86A4DAB3}.
\end{array}                                                     \tag{102.5}
\]

\begin{firewall}[Abelian, one polarization, second order, classical]
\label{fire:v19v99}
Only the abelian field with one transverse component and one
standing mode is treated, to second order; the non-abelian gauge
fields of the Standard Model, the fermions, the Higgs coupling, and
the quantization of the matter fields in the seed are not.  The
non-conversion statement is second order and for a standing wave.
\end{firewall}

\begin{decision}[V100 acquisition]
\label{dec:v19v99-next}
V100 is the compulsory companion audit and the ten-version sweep
V90--V99, after which the note is frozen as \texttt{V100\_f19} and a
new V1 opens.  The sweep will fix the direction of the twentieth
cycle: the emergent geometry is a cosmology whose matter content
enters through averaged stresses, dispersion relations, polarization
sources, and inhomogeneous secular terms; the twentieth cycle must
decide whether the Standard Model's non-abelian and fermionic sectors
can be placed in the same exact framework, or whether the framework
must first be extended beyond plane symmetry.
\end{decision}

\begin{assessment}[V99 acquisition]
\label{ass:v19v99}
The gauge field is placed in the plane-symmetric emergent cosmology
with an exact system and an exact second-order solution.  Its
averaged stress is the null dust of the massless scalar, but the
geometry reads its spin through the induced polarization
\(\tfrac14\sin^2z\sin^2t\) and through the absence of the scalar's
inhomogeneous secular growth; a standing electromagnetic wave does
not radiate gravitationally at second order.
\end{assessment}

\begin{record}[V99 append record]
\label{rec:v19v99}
V99 is appended to the validated V98 whose installed SHA--256 was
\texttt{F708B86181BF89951FB06FCD8F7E740D1EFB23D352462E0C2DD54756D592BD2F}.
After installation only V99 is the active numeric SM note; frozen
archives and unrelated notes are unchanged.
\end{record}


\section{V100: companion audit, ten-version sweep, and close of the nineteenth note cycle}
\label{sec:v19v100}

This is the required audit following V95, the ten-version sweep of
V90--V99, and the closing record of the nineteenth cycle.  Every
active companion note present in the shared folder was inspected
read-only.

\subsection{Companion snapshot}

\begin{record}[V100 companion snapshot]
\label{rec:v19v100-snapshot}
The files present were

\[
\begin{array}{c|r|r}
\hbox{file and SHA--256}&\hbox{bytes}&\hbox{lines}\\ \hline
\texttt{Renewal\_NS\_Stability\_Research\_Notes\_V26.tex}
&278283&5319\\
\multicolumn{3}{l}{
\texttt{82C887F8C2C69E4F28FAD7C3DC8DE5B9099CB5A4BF431EBA1FFF3E91935978AD}}
\\
\texttt{Renewal\_GRH\_Cofinal\_Visibility\_Attack\_V60.tex}
&580602&17282\\
\multicolumn{3}{l}{
\texttt{E090E736AE1C43D60756DB3E9C7929E0F2A179905B74630ADBE85266036737F4}}
\\
\texttt{Renewal\_Yang\_Mills\_Mass\_Gap\_Research\_Notes\_V35.tex}
&520002&12860\\
\multicolumn{3}{l}{
\texttt{6C014195DC1E1A52F7943D3840A669EB813C202A77B90B80F31E251AB1C1BB9D}}
\\
\texttt{Renewal\_YM\_Parallel\_Research\_Notes\_V5.tex}
&54174&1351\\
\multicolumn{3}{l}{
\texttt{306F1BB84CFA8A8D47EA569EC17E9545F9867C8D32B548EC9D9C444B4F3C0512}}.
\end{array}                                                     \tag{103.1}
\]

Since V95 the Yang--Mills mass-gap note advanced from V34 to V35
(alias Ward identity, exact contraction \(\rho\le17/36\) of the
blocked propagator map, geometric convergence of the symmetrized
Gaussian tower off the axes); the NS, GRH, and YM parallel notes are
unchanged since V95.
\end{record}

\begin{decision}[V100 audit decision]
\label{dec:v19v100-audit}
No companion theorem, numerical constant, or executable artifact is
imported.  The Yang--Mills advance is a contraction estimate for a
Euclidean Gaussian blocking; it does not bear on the Lorentzian
constrained dynamics with matter of V96--V99.
\end{decision}

\begin{proof}
The only changed file is the Yang--Mills note, whose V35 section
proves identities of a free-field blocking; no shared object with the
Einstein--matter systems of this cycle appears.
\end{proof}

\subsection{Ten-version sweep V90--V99}

\begin{proposition}[What V90--V99 established]
\label{prop:v19v100-sweep}
The block turned the programme from constraint calculus into
dynamics, and from vacuum into matter.  Its exact results are:

\begin{enumerate}
\item the exact plane-symmetric Einstein system in exponential
variables, the independent confirmation of the V89 solution, and the
homogeneous two-scale reduction with its collapse bound
\(T_*\leq\sqrt{128/3}\) (V91);
\item the exact fourth-order evolution with \(\Lambda_4=-145/2048\)
recovered from the dynamics, the confirmation of the reduction in the
transverse sector and its refutation in the longitudinal sector by
the inhomogeneous secular modulation of a standing wave, the wave's
amplitude growth \(5/64\) and frequency shift \(3/64\), and the
focusing theorem: every time-symmetric datum on a compact slice with
\(\Lambda<0\) focuses within \(\tfrac\pi2\sqrt{3/|\Lambda|}\) (V92);
\item for finite-mode data, \(\Lambda_2=-\langle|\nabla h|^2\rangle/8\),
additive null dust per mode, the Friedmann limit
\(\ddot a/a=\tfrac23\Lambda\) for isotropic sets, and the focusing time
\(\pi\sqrt6/\sqrt{\langle|\nabla h|^2\rangle}\) (V93);
\item the fourth-order cosmological constant of two-mode data as a
certified quartic form, direction-independent single-mode
coefficients, a universal pair value \(-\tfrac1{16}(|k_a|^2+|k_b|^2)\)
for half the pairs, negative residuals for the other half, and the
Sobolev bound \(|\Lambda_4|\leq C\|h\|_{H^2}^4\) (V94);
\item the identification of the \(\Lambda=0\) physical sectors of the
operator-level constraint with expanding or contracting flat radiation
universes, \(H^2=e/6\) from the normal-ordered graviton energy, the
exact second-order expanding datum in \(\mathbb Q(H)\), and the
redshift law of the wave (V96);
\item the plane-symmetric Einstein--scalar system and the theorem that
a massless scalar is an abelian polarization of the gravitational
wave, with the mixed datum certified to fourth order (V97);
\item the massive scalar as null dust plus dust in the ratio
\(k^2:m^2\), with the longitudinal secular contraction
\(-m^2t^2/16\) as the first matter property visible to the geometry
(V98);
\item the plane-symmetric Einstein--Maxwell system, the standing
electromagnetic wave's induced polarization \(\tfrac14\sin^2z\sin^2t\)
without secular growth, and the spin signatures that distinguish
vector from scalar radiation at equal averaged stress (V99).
\end{enumerate}
\end{proposition}

\begin{proof}
Items are the certified theorems of the cited sections.
\end{proof}

\subsection{Position at the close of the cycle}

\begin{assessment}[The nineteenth cycle]
\label{ass:v19v100-cycle}
The cycle began at V32 with the transport of reflection positivity and
the sharp-time no-go, established the toroidal rigidity that forces
\(\Lambda<0\) on time-symmetric data, constructed the operator-level
averaged constraint and its two-branch physical space, computed the
fourth-order interaction structure of the slice constraint through
the three-dimensional engine and corrected its own momentum sector,
and concluded that the graviton mass term of the slice constraint is
not cancelled in any sector.  It then went upstream to the dynamics
and found the emergent continuum to be a cosmology: the balanced torus
recollapses within \(\tfrac\pi2\sqrt{3/|\Lambda|}\), the \(\Lambda=0\)
sectors are radiation universes with Hubble rate set by the graviton
energy, and matter enters through averaged stresses, dispersion
relations, polarization sources, and inhomogeneous secular terms, with
a massless scalar indistinguishable from a graviton polarization and
mass and spin the first distinguishable properties.

What is not done is unchanged in kind: no gauge-invariant
three-dimensional constrained dynamics beyond plane symmetry, no
number-changing terms beyond one block, no renormalized mode-number
limit beyond the normal-ordering choice, no derivation of the
classical dynamics from the constraint operator, and no non-abelian
or fermionic sector.  The continuum Standard Model coupled to
Einstein gravity has not been constructed.
\end{assessment}

\begin{firewall}[Cycle boundary]
\label{fire:v19v100}
The cosmological identifications of V92--V99 are second-order
statements in plane symmetry with classical matter, certified by
exact arithmetic; calling them an emergent Standard Model--Einstein
continuum would ignore the five gaps listed above.
\end{firewall}

\begin{decision}[Post-rollover V1 acquisition]
\label{dec:v19v100-next}
After this file is frozen as
\texttt{Renewal\_SM\_Einstein\_Continuum\_Research\_Notes\_V100\_f19.tex},
the new active V1 of the twentieth cycle will give context and a
compact result ledger, and will open the non-abelian sector in the
only setting where the programme has exact tools: the plane-symmetric
Einstein--Yang--Mills system for an \(SU(2)\) potential with one
transverse colour component, whose self-interaction is the first
non-abelian term the emergent geometry can read.
\end{decision}

\begin{record}[V100 append record]
\label{rec:v19v100}
V100 is appended to the validated V99 whose installed SHA--256 was
\texttt{7FE8006305743C7380E7D5D624FAD3D7E7E6E898D047EB60FE9C65AC2E8A0A68}.
This file closes the nineteenth cycle and is to be preserved as the
\texttt{f19} archive.  Unrelated notes are unchanged.
\end{record}


\end{document}
